\documentclass[a4paper,11pt]{article}
\usepackage[T1]{fontenc}
\usepackage{lmodern}
\usepackage{amsmath,amssymb,amsthm,mathtools}
\usepackage{mathrsfs}
\usepackage[margin=27mm]{geometry}
\usepackage{microtype,enumitem,float}
\usepackage{longtable,array}
\usepackage{tikz}
\usepackage[hidelinks]{hyperref}
\usepackage{bookmark}
\hypersetup{
 pdftitle={Sharp local cyclicity of the complete quadratic graphic I4-1},
 pdfauthor={Haibo Lu},
 pdfsubject={Complete manuscript on exact local cyclicity in the full quadratic germ},
 pdfkeywords={finite cyclicity, quadratic vector fields, DRR graphic, saddle-node, Dulac map, entry-exit}}
\numberwithin{equation}{section}
\newtheorem{theorem}{Theorem}[section]
\newtheorem{proposition}[theorem]{Proposition}
\newtheorem{lemma}[theorem]{Lemma}
\newtheorem{corollary}[theorem]{Corollary}
\theoremstyle{definition}
\newtheorem{definition}[theorem]{Definition}
\theoremstyle{remark}
\newtheorem{remark}[theorem]{Remark}
\newcommand{\R}{\mathbb R}
\newcommand{\N}{\mathbb N}
\newcommand{\eps}{\varepsilon}
\newcommand{\dd}{\,d}
\newcommand{\tauc}{\tau_{\rm cub}}
\DeclareMathOperator{\diag}{diag}
\DeclareMathOperator{\dist}{dist}
\DeclareMathOperator{\sgn}{sgn}
\setlist{nosep}
\setlength{\emergencystretch}{1em}
\raggedbottom
\title{Sharp local cyclicity of the complete quadratic graphic $I_4^1$}
\author{Haibo Lu\\
\small Shanghai Institute of Technology, Shanghai 201418, People's Republic of China\\
\small Corresponding author: \href{mailto:luhaibo1985@gmail.com}{\texttt{luhaibo1985@gmail.com}}\\
\small \href{https://orcid.org/0009-0000-2717-5968}{ORCID: 0009-0000-2717-5968}}
\date{18 October 2026}
\begin{document}
\maketitle

\begin{abstract}
The complete Dumortier--Roussarie--Rousseau graphic $I_4^1$ consists of
two parabolic branches joining a finite saddle--node to a
multiplicity-four vertex at infinity. We prove that its local cyclicity
is exactly three in the full quadratic coefficient germ. One
coefficient neighborhood and one whole-graphic Hausdorff collar give
the upper bound, and every smaller neighborhood and collar contain one
field with three hyperbolic complete cycles. We compare the curvature
of the completed physical return with that of a cubic entry--exit
action, whose quadratic coefficient cancels from the weighted
curvature. Relative estimates near reversibility, separate root scales,
and the ratio of the two physical tail gains preserve the comparison
through root coalescence. At degree loss, comparison of the actual
small-power section maps gives a two-zero bound. A finite coefficient
cover yields one uniform count. The lower construction uses two
coefficients to match a vanishing complete gain to a steep action
transition and preserve three simple roots in one field. Whole-orbit
convergence and a diagonal limit give attainment at the endpoint.
In a fixed analytic slice, the ratio
$p=\alpha_1/\sqrt{-\alpha_0}$ of the translated axis coefficients,
defined for $\alpha_0<0$,
tends to zero from above along three-cycle sequences approaching the
organizing field and graphic. The bound is uniform down to $p=0$, where at most
one cycle occurs, and is sharp below every positive cutoff on $p$.
\end{abstract}

\noindent\textbf{Keywords:} finite cyclicity; quadratic vector field;
DRR graphic; saddle--node; Dulac map; entry--exit asymptotics.

\medskip
\noindent\textbf{2020 Mathematics Subject Classification:}
34C07 (primary); 34C23, 37G15 (secondary).

\medskip
\section{Introduction}
\label{sec:introduction}

Determining how many limit cycles bifurcate from a quadratic graphic
under variation of all quadratic coefficients is a local part of the
Dumortier--Roussarie--Rousseau finiteness programme
\cite[Introduction, author copy pp.~1--2]{DIR02}. We study its graphic
$I_4^1$, whose two parabolic branches join a finite saddle--node to a
nilpotent saddle--node at infinity
\cite[Proposition~4.2.2 and Figure~8, journal pp.~111--113]{DRR94}.

For the present graphic the axis polynomial is cubic. Its incident
roots can coalesce at different rates, and other coefficient limits
reduce its degree. The cyclicity problem therefore asks for a bound on
the return through both vertices throughout these configurations.
We obtain the sharp bound three: one coefficient neighborhood and one
Hausdorff collar admit at most three cycles, and arbitrarily small
neighborhoods and collars admit three hyperbolic cycles in one field.

The return-map approach connects singular-passage geometry with zero
counting. Saddle--node normal forms provide finitely smooth transition
coordinates \cite{DIR02}; transition maps give finite cyclicity for
elementary quadratic graphics surrounding a focus or center
\cite{DGR02}. For graphics through nilpotent singularities, finite
cyclicity has been obtained by controlling the regular transition
derivatives \cite[Theorems~5.4 and~6.3, journal pp.~372 and 388--389]{ZR02}
and nonlinear transitions through a finite saddle--node
\cite[Theorem~3.8, author copy p.~23]{RZ04}. In the center case,
division in the center ideal permits a displacement--Rolle argument
\cite[Sections~5.1 and 5.4--6, author copy pp.~19 and 23--25]{RR08}.
Sharp counts have been obtained for a local
codimension-four nilpotent saddle--node family
\cite[Theorem~1.1, peer-reviewed author version p.~4]{Huzak17JDE}
and for hyperbolic hemicycles under quadratic perturbations
\cite[Definition~1.4 and Theorem~B, published pp.~3 and 7]{MV25}.
For slow--fast canard cycles containing one remote hyperbolic saddle
with nonunit hyperbolicity ratio, whose slow dynamics is regular or
has only finitely many singularities, all in the interior of the slow
segment, Huzak proves an upper bound of three
\cite[Definition~1 and Theorem~2.2, peer-reviewed author version
pp.~7--8]{Huzak22}.

At the infinity passage of $I_4^1$, entry--exit theory describes how an
orbit follows the axis and reaches its outgoing section.
De Maesschalck and Schecter establish the entry--exit relation and its
smooth parameter dependence \cite[Theorem~1.1 and Corollary~1.2]{DMS16}.
Huzak and Kristiansen develop this description for degenerate turning
points, including the $n=1$ geometry associated with $I_4^1$
\cite[printed p.~3 and Theorem~6.1, printed p.~25]{HK26}.
Their strict $n=1$ transition theorem assumes a uniformly negative
axis polynomial on its selected chart
\cite[Theorem~2.4, printed pp.~6--7]{HK26}. These local transition
results provide entry--exit coordinates and regularity. The present
argument also needs the other root and degree-loss charts, the finite
saddle--node closing, and one count for the complete physical return.
The neighboring fixed-parameter graphic $I_2^1$ is studied in
\cite{Lu26I21}. At the $I_4^1$ endpoint its finite-passage multiplier
vanishes; the lower construction below realizes three cycles internally
before taking this endpoint limit.

The organizing field in the original classification coordinates is
\begin{equation}\label{eq:problem-field}
 F_*:\quad
 \dot x=\lambda_*x-y+x^2,\qquad
 \dot y=x+\lambda_*y+2\lambda_*x^2+xy,
 \qquad \lambda_*=\frac{1}{\sqrt3}.
\end{equation}
Its invariant parabola is
\begin{equation}\label{eq:problem-parabola}
 y=\pi_*(x):={x^2\over2}-\lambda_*x-{\lambda_*^2+1\over2}.
\end{equation}
Write
\(\iota_*(x,y)=(x,y,1)/\sqrt{1+x^2+y^2}\),
let \(p_*=(-2\lambda_*,2/3)\), \(q_*=(0,1,0)\), and let
\(\Gamma_-=\overline{\iota_*\{(x,\pi_*(x)):x<-2\lambda_*\}}\) and
\(\Gamma_+=\overline{\iota_*\{(x,\pi_*(x)):x>-2\lambda_*\}}\). The complete
marked graphic is
\begin{equation}\label{eq:positive-word}
 \Gamma_{I4}=\Gamma_-\cup\Gamma_+,
 \qquad q_*\longrightarrow p_*\longrightarrow q_*.
\end{equation}
The two occurrences of \(q_*\) are one marked projective vertex.

Let \(\mathcal Q_2\) be the twelve-dimensional affine space of real planar
vector fields of degree at most two, and let \(d_{\rm H}\) be the Hausdorff
distance for a fixed smooth metric on the closed upper hemisphere.

\begin{theorem}[Exact local cyclicity of \(I_4^1\)]\label{thm:main}
There are a neighborhood $\mathcal V$ of $F_*$ in $\mathcal Q_2$ and
$\eps>0$ such that every $F\in\mathcal V$ has at most three distinct
isolated periodic orbits $\gamma$ with
\[
 d_{\rm H}\bigl(\iota_*(\gamma),\Gamma_{I4}\bigr)<\eps.
\]
Conversely, for every neighborhood $\mathcal W$ of $F_*$ and every
$\delta>0$, some $F\in\mathcal W$ has three geometrically distinct
hyperbolic periodic orbits satisfying this inequality with $\delta$
in place of $\eps$. Thus
\[
 \operatorname{Cycl}_{\mathcal Q_2}(\Gamma_{I4},F_*)=3.
\]
Both assertions concern the complete two-branch marked graphic and
the full twelve-coefficient quadratic germ.
\end{theorem}

\subsection{The mechanism of the proof}

The finite saddle--node supplies a small multiplier $m>0$. Including
the two regular closing maps, cycles are zeros of one displacement in
an affine physical input:
\begin{equation}\label{eq:intro-displacement}
 F_{\theta,m}(x)=x-m\Pi_{\theta,m}(x)=0.
\end{equation}
The cubic entry--exit action reveals the key cancellation: its
positively weighted curvature contains the linear and cubic
coefficients, while the quadratic coefficient cancels. The product of
the two matched axis distances decreases with the incoming quadratic
label, so every curvature zero has the same crossing orientation.
Section~\ref{subsec:representative-cubic} derives the identity and
transfers this orientation to the completed physical map. With a
positive displacement derivative on the upper continuation,
Rolle's theorem then gives at most three displacement zeros.

Uniformity requires more than convergence to the action. The two
nonlinear section maps contribute to physical curvature, and their
individual gains can vanish at different rates. We retain both gains
and their actual section images. Near reversibility, division by the
coefficient distance controls the curvature as its leading term
vanishes. Separate root-collar estimates handle independently
approaching roots; matching actual inward heights also controls the
original ghost tails beyond the growing-action ranges. On the upper
continuation, a bounded complete derivative gives $F'>1/2$ after
$m$ is chosen small.
At degree loss, a comparison of the two actual small-power section
maps instead gives the same upward orientation at every critical
point of the displacement, and hence at most two zeros on that
whole coefficient fibre.

These estimates count one connected interval of actual trajectories.
We fix the finite coefficient cover and overlap reserves, then a
positive unscaled height cap, and finally independent coefficient and
multiplier bounds. Geometric capture places every nearby whole cycle
in the interval already counted. Fixed positive ratio intervals have
a stronger two-cycle bound, forcing three-cycle parameters to
concentrate at $p\to0^+$ in the slice defined below.

The lower construction is independent of this upper count. First fix
$c>1$ near one. The complete organizing gain is
$\exp(-4\pi/\sqrt{c-1})$; choose the action-model parameter
$\epsilon>0$ so that this equals $2\epsilon^{5/2}$.
The central transition has slope of order $\epsilon^{-5/2}$,
compensating for the vanishing gain. Two coefficients then calibrate
the offset and the complete gain in one field, and one small scale
preserves all three simple roots and their physical hyperbolicity.
Their complete orbits converge to both branches of the organizing
parabola. A diagonal choice as $c\downarrow1$ gives the endpoint
lower bound. The analytic affine and positive-time slice transports
the conclusions to all twelve quadratic coefficients.

\subsection{The organizing family and coordinate model}

Put
\[
 A_0=\begin{pmatrix}\sqrt3/2&0\\[2pt]\sqrt3/4&3/4\end{pmatrix},
 \qquad \rho={2\over\sqrt3}.
\]
The affine change
\begin{equation}\label{eq:canonical-change}
 \binom X Y=A_0\binom x y,
 \qquad t_{\rm can}=\rho t
\end{equation}
has positive time factor and sends \eqref{eq:problem-field} to
\begin{equation}\label{eq:base}
 H_0:\qquad \dot X=X-Y+X^2,\qquad \dot Y=X+X^2+XY.
\end{equation}
The parabola becomes $Y=(X^2-1)/2$, the finite vertex is $p=(-1,0)$,
and the point at infinity is $q=(0,1,0)$.  For a quadratic field $F$ in the
original coordinates, its field in the canonical coordinates is
\begin{equation}\label{eq:fixed-coefficient-transport}
 \mathscr C_0(F)(Z)=\rho^{-1}A_0F(A_0^{-1}Z),
 \qquad Z=(X,Y)^{\mathsf T}.
\end{equation}
The map $\mathscr C_0:\mathcal Q_2\to\mathcal Q_2$ is a linear
isomorphism, with inverse
\[
 \mathscr C_0^{-1}(H)(u)=\rho A_0^{-1}H(A_0u).
\]
In particular, it sends the full twelve-coefficient germ at $F_*$ onto the
full twelve-coefficient germ at $H_0$.

Let
\[
 \iota_0(X,Y)={ (X,Y,1)\over\sqrt{1+X^2+Y^2}}
\]
be the canonical affine embedding, and let $S$ be the projective extension
of $A_0$ to the closed upper hemisphere.  Thus
\begin{equation}\label{eq:fixed-projective-map}
 S(U,V,W)=
 {\left(A_0\binom{U}{V},W\right)
  \over\sqrt{\left\lVert A_0\binom{U}{V}\right\rVert^2+W^2}},
 \qquad S\circ\iota_*=\iota_0\circ A_0.
\end{equation}
This is a smooth projective diffeomorphism preserving the upper hemisphere.
Set
\begin{equation}\label{eq:canonical-graphic}
 \Gamma:=S(\Gamma_{I4}).
\end{equation}
It is the complete two-branch marked graphic of $H_0$, with the same
positive-time branch order because $\rho>0$.  For the fixed smooth metric
used in Theorem~\ref{thm:main}, compactness gives a constant
$1\le L_0<\infty$
such that
\begin{equation}\label{eq:fixed-map-lipschitz}
 d_{\rm H}(S K,S K')\le L_0d_{\rm H}(K,K')
\end{equation}
for all nonempty compact subsets $K,K'$ of the hemisphere.  All calculations
below are made for $H_0$ and $\Gamma$; the final proof pulls the resulting
coefficient neighborhood back by \eqref{eq:fixed-coefficient-transport} and
uses \eqref{eq:fixed-map-lipschitz} to transfer the Hausdorff collar.

The canonical model also lies in the organizing family
\begin{equation}\label{intro:eq-i21-family}
 H_c:\quad \dot X=X-Y+X^2,\qquad
 \dot Y=cX+X^2+XY,
 \qquad
 \Gamma_c=\overline{\iota_0\{Y=(X^2-c)/2\}},
\end{equation}
where the closure is taken in the closed upper Poincare hemisphere and
\(\iota_0(X,Y)=(X,Y,1)/\sqrt{1+X^2+Y^2}\). Along the invariant parabola,
\begin{equation}\label{intro:eq-parabola-speed}
 \left.\dot X\right|_{\Gamma_c}
 =\frac{(X+1)^2+c-1}{2}.
\end{equation}
For every fixed \(c>1\), the finite point \(X=-1\) is crossed with positive
speed. At \(c=1\), the same point becomes a finite saddle--node. The two
infinite ends remain the two marked occurrences of one projective vertex;
no equatorial arc is added. This is the geometric change from the fixed \(I_2^1\) problem to the \(I_4^1\) endpoint.

In the original coordinates, direct substitution and \(3\lambda_*^2=1\) give
\begin{align}
 F_*\bigl(y-\pi_*(x)\bigr)
 &=\dot y-\pi_*'(x)\dot x
   =2x\bigl(y-\pi_*(x)\bigr),\notag\\
 \left.\dot x\right|_{y=\pi_*(x)}
 &=\lambda_*x-\pi_*(x)+x^2
   ={1\over2}(x+2\lambda_*)^2.
 \label{eq:parabola-invariance-speed}
\end{align}
Writing \(u=x+2\lambda_*\) and \(w=y-\pi_*(x)\), the complete system is
\begin{equation}\label{eq:finite-saddle-node-coordinates}
 \dot u=\frac{u^2}{2}-w,\qquad
 \dot w=2(u-2\lambda_*)w .
\end{equation}
Its linearization at \((u,w)=(0,0)\) is
\[
 \begin{pmatrix}0&-1\\[2pt]0&-4\lambda_*\end{pmatrix}.
\]
Thus it has one zero and one nonzero eigenvalue, while the restriction to
\(w=0\) has the quadratic degeneracy \(\dot u=u^2/2\). The finite point is
therefore a saddle--node. The multiplicity of the other marked vertex is checked
separately in the infinity chart of Section~\ref{sec:normalization}.
The lower bound used here is constructed internally in
Theorem~\ref{p4:fixed-c}. The complete organizing multiplier
$\exp(-4\pi/\sqrt{c-1})$ vanishes as $c\downarrow1$;
at the endpoint the regular finite passage becomes a saddle--node.
We therefore derive the upper estimate at $c=1$ with the finite
multiplier separated from the actual infinity passage.

% Figure contract: the complete marked graphic and its positive-time return.
% Native LaTeX picture; schematic, not a numerical phase portrait.
\begin{figure}[H]
\centering
\setlength{\unitlength}{1pt}
\begin{picture}(430,190)
  % Panel (a): the two marked branches and the repeated projective vertex.
  \linethickness{.35pt}
  \put(18,164){\line(1,0){144}}
  \put(129,171){\makebox(0,0)[l]{\small equator}}
  \linethickness{.8pt}
  \qbezier(88,164)(23,164)(26,105)
  \qbezier(26,105)(29,50)(88,50)
  \qbezier(88,50)(147,50)(150,105)
  \qbezier(150,105)(153,164)(88,164)
  \put(26,109){\vector(0,-1){14}}
  \put(150,96){\vector(0,1){14}}
  \put(31,101){\makebox(0,0)[l]{\scriptsize (+t)}}
  \put(145,103){\makebox(0,0)[r]{\scriptsize (+t)}}
  \put(88,164){\circle*{3}}
  \put(88,50){\circle*{3}}
  \put(88,179){\makebox(0,0){\(q_*^-=q_*^+=q_*\)}}
  \put(88,39){\makebox(0,0){\(p_*\)}}
  \put(45,87){\makebox(0,0){\(\Gamma_-\)}}
  \put(132,87){\makebox(0,0){\(\Gamma_+\)}}
  \linethickness{.4pt}
  \put(18,127){\line(3,1){32}}
  \put(120,137){\line(3,-1){32}}
  \put(9,129){\makebox(0,0)[r]{\(\Sigma_-\)}}
  \put(159,129){\makebox(0,0)[l]{\(\Sigma_+\)}}
  \put(88,20){\makebox(0,0){\small
    \(q_*^-\xrightarrow{\,\Gamma_-\,}p_*
      \xrightarrow{\,\Gamma_+\,}q_*^+\quad(+t)\)}}
  \put(88,0){\makebox(0,0){\small (a) The complete marked graphic}}

  % Panel (b): all four positive-time first-hit factors of the return.
  \put(255,150){\makebox(0,0){\(\Sigma_-\)}}
  \put(395,150){\makebox(0,0){\(\Sigma_+\)}}
  \put(255,70){\makebox(0,0){\(\eta_{\rm in}\)}}
  \put(395,70){\makebox(0,0){\(\eta_{\rm out}\)}}
  \put(379,150){\vector(-1,0){108}}
  \put(255,138){\vector(0,-1){56}}
  \put(271,70){\vector(1,0){108}}
  \put(395,82){\vector(0,1){56}}
  \put(325,162){\makebox(0,0){\(T_\theta\)}}
  \put(245,110){\makebox(0,0)[r]{\(E_\theta\)}}
  \put(405,110){\makebox(0,0)[l]{\(C_\theta\)}}
  \put(325,58){\makebox(0,0){\(M_m(\eta)=m\eta\)}}
  \put(325,25){\makebox(0,0){\small
    \(\mathcal R_\theta=C_\theta\circ M_m\circ E_\theta\circ T_\theta\)}}
  \put(325,0){\makebox(0,0){\small (b) The positive-time return composition}}
\end{picture}
\caption{The complete marked graphic and the return used to count its nearby
cycles.  In panel~(a), \(\Gamma_-\) runs from \(q_*^-\) to \(p_*\) and
\(\Gamma_+\) from \(p_*\) to \(q_*^+\) in positive time.  The equality
\(q_*^-=q_*^+=q_*\) records two occurrences in the return word of one
repeated projective vertex, not two points.  The branches have a common
tangent there, and no equatorial arc belongs to \(\Gamma_{I4}\).  Panel~(b)
shows the complete composition of positive-time first-hit maps on their
common through-side domain: the infinity transition \(T_\theta\), the regular
entry map \(E_\theta\), the finite saddle--node multiplier
\(M_m(\eta)=m\eta\), and the regular closing map \(C_\theta\).  Its arrows
denote maps, not additional trajectories.  All shapes and distances are
schematic.}
\label{fig:graphic-sections}
\end{figure}



\subsection{The small positive ratio and its sharpness}

The analytic slice in Lemma~\ref{lem:slice} is fixed once for the full
quadratic germ. For the same field $F$, write
$\theta_F=\theta(\mathscr C_0(F))$, let
$P(v)=\mu_1+\mu_2v-v^3$, $R(v)=\mu_3+v+v^2$, and use the root
$\gamma(F)$ of $R$ near zero. Set
\begin{equation}\label{eq:sharp-sector-coordinates}
 \alpha_0(F)=P(\gamma(F)),\qquad
 \alpha_1(F)=P'(\gamma(F)),\qquad
 s(F)=\sqrt{-\alpha_0(F)},\quad p(F)=\alpha_1(F)/s(F),
\end{equation}
where $s,p$ are used for $\alpha_0<0$. The ratio therefore has one
specified normalization throughout the upper and lower arguments.

\begin{theorem}[Uniform small positive ratio and sharpness]
\label{thm:small-sector}
There are an ordinary open neighborhood $\mathcal V_{\rm s}$ of $F_*$
in $\mathcal Q_2$, constants $\eps_{\rm s}>0$ and $p_{\rm s}>0$, such
that every $F\in\mathcal V_{\rm s}$ with
\[
 \alpha_0(F)<0,\qquad 0\le p(F)\le p_{\rm s}
\]
has at most three distinct isolated periodic orbits whose complete
images have Hausdorff distance less than $\eps_{\rm s}$ from
$\Gamma_{I4}$. On its connected admitted physical return interval,
the corresponding isolated displacement zeros have total multiplicity
at most three. On the face $p=0$ the bound is one.

For every neighborhood $\mathcal W$ of $F_*$, every $\delta>0$ and
every $\eta>0$, one field $F\in\mathcal W$ has three geometrically
distinct hyperbolic complete periodic orbits at Hausdorff distance
less than $\delta$ from $\Gamma_{I4}$ and satisfies
\[
 \alpha_0(F)<0<\alpha_1(F),\qquad
 0<p(F)<\min\{p_{\rm s},\eta\}.
\]
The upper bound retains the independently varying raw coefficients;
no restriction on $\gamma/s$, relative root distances, or the actual
finite multiplier compared with a source radius is added.
\end{theorem}

Proposition~\ref{prop:lower-coefficient-face} locates the constructed
triples more precisely: $p\to0^+$ and $\gamma/s\to+\infty$, while
the two positive roots remain separated in the positive-gamma chart.
Thus the sharpness family is retained by the unrestricted small-ratio
estimate. Section~\ref{sec:sharp-small-sector} proves the theorem from
the full regional estimates below.

\subsection{Organization}

Section~\ref{sec:normalization} constructs the physical return and
slice. Sections~\ref{sec:counting-tools}--\ref{sec:small-positive-source-dictionaries}
give the counting calculus and cubic mechanism.
Section~\ref{sec:common-cover-conclusion} assembles the regional
estimates, uniform budgets, capture and concentration.
Section~\ref{sharp:lower-section} constructs the simultaneous lower
bound with its model, passage, tail and coefficient proofs; the
small-ratio and full-germ proofs follow. The remaining regional transition
proofs and coefficient dictionaries are developed in the appendices.
Two illustrative ghost paths are provided in the optional supplement.
\medskip
\setcounter{tocdepth}{1}
\tableofcontents
\medskip

\section{The full quadratic germ and its physical return}
\label{sec:normalization}

We first separate the affine and positive-time changes from the five
essential coefficients. We then construct the return through both
singular points, including the two regular closing maps. The resulting
displacement is written in an affine section height, which will be used
throughout the zero count.

\subsection{Reduction to five parameters}

The following five-parameter family is transverse to the seven affine
and positive-time directions:
\begin{equation}\label{eq:slice}
 H_\theta:\quad
 \begin{cases}
 \dot X=AX-Y+X^2+(\mu_2+\mu_3)XY+\mu_1Y^2,\\
 \dot Y=CX+X^2+XY+\mu_3Y^2,
 \end{cases}\qquad
 \theta=(A,C,\mu_1,\mu_2,\mu_3)
\end{equation}
near \(\theta_0=(1,1,0,0,0)\).

\begin{lemma}\label{lem:slice}
There are neighborhoods \(\mathcal U\subset\mathcal Q_2\) of \(H_0\) and
\(\mathcal W\) of \((I,0,1,\theta_0)\) such that every
\(H\in\mathcal U\) has a unique representation with parameters in
\(\mathcal W\),
\begin{equation}\label{eq:affine-action}
 H(z)=a_{\rm t}L^{-1}H_\theta(Lz+c),
 \qquad (L,c,a_{\rm t},\theta)\ \hbox{near }(I,0,1,\theta_0),
\end{equation}
with \(a_{\rm t}>0\).  The resulting maps
\(H\mapsto \mathfrak A(H):=(L,c,a_{\rm t})\) and
\(H\mapsto\theta(H)\) are
analytic on \(\mathcal U\), and satisfy
\(\mathfrak A(H_0)=(I,0,1)\) and \(\theta(H_0)=\theta_0\).
\end{lemma}
\begin{proof}
The derivative of the coefficient map has columns
\[
 DH_0h-Dh\,H_0,\quad
 h=(1,0),(0,1),(X,0),(Y,0),(0,X),(0,Y),
\]
followed by \(H_0\) and the five vectors
\[
 (X,0),\quad(0,X),\quad(Y^2,0),\quad(XY,0),\quad(XY,Y^2).
\]
In monomial order \((1,X,Y,X^2,XY,Y^2)\) in each component,
the determinant is \(3\). The analytic inverse function theorem
proves the claim; positivity of \(a_{\rm t}\) persists near \(1\).
\end{proof}

For \(\mathfrak A=(L,c,a_{\rm t})\), denote by
\(T_{\mathfrak A}\) the extension of
\(z\mapsto Lz+c\) to the closed hemisphere, obtained by normalizing
\((X,Y,W)\mapsto(L(X,Y)^{\mathsf T}+cW,W)\).  After shrinking
\(\mathcal U\) to have compact closure in the inverse-function
neighborhood, these extensions and their inverses have a common Lipschitz
constant \(1\le L_*<\infty\).  Moreover,
\begin{equation}\label{eq:slice-graphic-motion}
 d_{\rm H}\bigl(T_{\mathfrak A(H)}\Gamma,\Gamma\bigr)\longrightarrow0
 \qquad\text{as }H\longrightarrow H_0,
\end{equation}
because \(\mathfrak A(H)\to(I,0,1)\).

\subsection{The two singular points}

In the chart at the positive \(Y\)-direction, use
\(v=X/Y,\ z=1/Y\). Multiplying the vector field by \(z>0\) preserves
the direction of its orbits in the finite plane and gives
\begin{equation}\label{eq:raw}
 \dot v=D_\theta(v,z)=P_\theta(v)+Q_\theta(v)z,\qquad
 \dot z=-zG_\theta(v,z)=-z\{R_\theta(v)+S_\theta(v)z\},
\end{equation}
where
\begin{equation}\label{eq:polynomials}
 \begin{aligned}
 P(v)&=\mu_1+\mu_2v-v^3,& Q(v)&=-1+Av-Cv^2,\\
 R(v)&=\mu_3+v+v^2,& S(v)&=Cv .
 \end{aligned}
\end{equation}
The physical side is \(z>0\); the equator \(z=0\) is invariant.
We always use the positive-time chart field on this side.
At the organizing parameter \(\theta_0=(1,1,0,0,0)\), the chart is
\begin{equation}\label{eq:infinity-chart-field}
 \dot v=-v^3-(1-v+v^2)z,\qquad
 \dot z=-z(v+v^2+vz).
\end{equation}
The linearization is \(\bigl(\begin{smallmatrix}0&-1\\0&0\end{smallmatrix}\bigr)\).
Eliminating the analytic unit \(1-v+v^2\) gives the nullcline
\(z=-v^3/(1-v+v^2)\), on which
\begin{equation}\label{eq:infinity-multiplicity-four}
 \left.\dot z\right|_{\dot v=0}
 =\frac{v^4}{(1-v+v^2)^2}.
\end{equation}
Here \((1+v)(1-v+v^2)=1+v^3\); the reduced zero has order four,
the multiplicity of the nilpotent singularity at infinity.
The organizing arc has equation
\begin{equation}\label{eq:organizing-arc}
 z_\Gamma(v)=\sqrt{1+v^2}-1
\end{equation}
and is traversed with decreasing \(v\); its expansion
\(z_\Gamma(v)=v^2/2+O(v^4)\) gives the common equatorial tangent.
Choose a small fixed \(\delta>0\) and the sections
\(\Sigma_\pm=\{v=\pm\delta\}\), with small height intervals around this arc.
Take the height intervals small enough that \(-D\) has a positive lower
bound on their closures.
On the fixed small chart, a coefficient shrink gives
\begin{equation}\label{loc:signs}
 Q<-q_0<0,\qquad G_v=1+2v+Cz>g_0>0,\qquad
 P_\theta(\delta)<0,\quad G_\theta(-\delta,z)<0.
\end{equation}

The small zero of the transverse eigenvalue and the first two
equatorial coefficients there are
\begin{equation}\label{eq:center-parameters}
 \gamma=\frac{-1+\sqrt{1-4\mu_3}}2,\qquad
 \alpha_0=P(\gamma),\qquad \alpha_1=P'(\gamma).
\end{equation}
The Jacobian with respect to \((\mu_3,\mu_1,\mu_2)\) at \(\theta_0\)
is \(-1\), so these, together with \(A,C\), are parameter coordinates.
The translated polynomials and their rescalings are introduced in
Section~\ref{sec:coefficient-regimes}. The five raw coefficients
\(A-1,C-1,\gamma,\alpha_0,\alpha_1\) vary independently.

For the finite saddle-node, set \(u=X+1,\ W=Y+u\).
At \(\theta_0\),
\[
 \dot u=u^2-W,\qquad \dot W=u^2+(u-2)W .
\]
The zero set of the second component is an analytic graph
\(W=\psi(u,\theta)\), used here only to locate equilibria.
The first component on this graph has a Weierstrass factor
\[
 e(u,\theta)(u^2+b_1(\theta)u+b_0(\theta)),\qquad
 e(0,\theta_0)=1/2 .
\]
Indeed, at \(\theta_0\) it equals \(u^2(1-u)/(2-u)\).
Differentiation gives
\[
 b_0=C-A+O(|\theta-\theta_0|^2),\qquad
 b_1=A-1+\mu_2+\mu_3+O(|\theta-\theta_0|^2).
\]
Consequently \(K=b_0-b_1^2/4\) is a discriminant coordinate with
\(dK=dC-dA\) at the organizing field.
We also write \(K_{\rm fin}=K\) for this finite-passage discriminant.

The organizing analytic germ has eigenvalues $0,-2$ and a nonzero
quadratic weak coefficient $1/2$. Thus its multiplicity is two, the
case $\mu=1$ of the saddle--node normal form. Fix a finite joint
differentiability order \(N\ge12\) before restricting the coefficients.
The ordinary section maps and their inverses will be used through phase
order three, while division by the discriminant costs at most one
parameter derivative. The intrinsic root-collar kernels use phase order
five, axis order six and selected mixed order seven; these estimates and
the strong projection estimates are obtained from the analytic chart
field. Choose the order for the normal-form theorem now, increasing it
if necessary before choosing any neighborhood. This finite surplus
accommodates coefficient differentiation and Hadamard division, so a
jointly smooth normalizer of all orders is unnecessary.
The finite-smooth multiparameter unfolding normal form
\cite[Theorem~3 and formula~(2.2), October 2000 author copy p.~6]{DIR02} gives,
by a joint \(C^N\) orbital equivalence preserving positive time,
\begin{equation}\label{eq:finite-normalform}
 \xi'=\eps_f+\frac{\xi^2}{1+b\xi},\qquad \eta'=-\eta .
\end{equation}
Here the parameter in the cited theorem is the complete vector
\(\theta-\theta_0\). For \(\mu=1\), its polynomial \(P(\theta,\xi)\)
is the constant \(\eps_f\). The reduction divides by the nonzero strong
rate, which is positive at the organizing parameter and remains so
nearby; it thus gives the stated positive-time equivalence. The theorem
imposes no rank condition on the unfolding coefficients.
The linear calibration
\(t_{\rm nf}=2t,\ \xi=(X+1-Y)/8+O(|(X+1,Y)|^2)\) gives
\(\eps_f=(C-A)/16+O(|\theta-\theta_0|^2)\).
The double-equilibrium loci in the two descriptions coincide.
In parameter coordinates \((K,\vartheta)\), Hadamard division gives
\[
 \eps_f(K,\vartheta)=K h(K,\vartheta),\qquad
 h(K,\vartheta)=\int_0^1\partial_K\eps_f(tK,\vartheta)\,dt,
 \qquad h(\theta_0)=1/16.
\]
Continuity makes \(h\) a positive unit on a smaller coefficient
neighborhood.

For fixed internal sections \(\xi=\pm\rho\), with \(1+b\xi>1/2\),
the passage exists when \(K>0\) and is exactly
\begin{equation}\label{eq:finite-multiplier}
 \eta_{\rm out}=m_{\rm fin}\eta_{\rm in},\quad
 m_{\rm fin}=e^{-T_f},\quad
 T_f=\int_{-\rho}^{\rho}
       \frac{d\xi}{\eps_f+\xi^2/(1+b\xi)}
     =\frac{\pi}{\sqrt{\eps_f}}+O(1).
\end{equation}
We write \(m=m_{\rm fin}\) in the completed return below.
Completing the square in \(\xi^2+b\eps_f\xi+\eps_f\) proves the last
estimate. Differentiating its arctangent and logarithmic expressions
also shows that each fixed finite parameter derivative of \(m\) is
bounded by \(CK^{-M}e^{-c/\sqrt K}\). Thus \(m\) extends flatly
by zero to \(K=0\), to the chosen differentiability order.
If \(K\le0\), an orbit entering the selected channel approaches the
saddle-node or the left equilibrium and cannot exit through
\(\xi=\rho\). The strong sides of this box point inward.

\subsection{The exact displacement}

Let \(T_\theta:\Sigma_+\to\Sigma_-\) be the infinity transition when
it exists in positive time. The rest of the return consists of two
regular transitions and the finite passage
\eqref{eq:finite-multiplier}; see Figure~\ref{fig:graphic-sections}.

\begin{lemma}\label{lem:closing}
There are uniformly regular increasing maps \(C_\theta,E_\theta\)
on fixed section intervals such that the finite-side return is
\(q\mapsto C_\theta(mE_\theta(q))\). Their derivatives and inverse
derivatives are bounded through the chosen finite order on buffered
section intervals.
If \(a_\theta=C_\theta(0)\), \(f_\theta=C_\theta-a_\theta\),
and \(d_\theta=E_\theta\), then in the affine coordinate
\(x=z_{\rm in}-a_\theta\) the cycle equation is
\begin{equation}\label{eq:return}
 \Delta_\theta(x)=x-m\Pi_\theta(x)=0,\qquad
 \Pi_\theta(x)=\mathcal C_{\theta,m}(T_\theta(a_\theta+x)),
\end{equation}
where
\begin{equation}\label{eq:divided-closing}
 \mathcal C_{\theta,m}(q)=d_\theta(q)
       \int_0^1 f_\theta'(tm d_\theta(q))\,dt .
\end{equation}
In the output coordinate \(t\), the equivalent equation is
\begin{equation}\label{eq:inverse-return}
 S_\theta(t)-a_\theta-f_\theta(mt)=0,\qquad
 S_\theta=(d_\theta\circ T_\theta)^{-1}.
\end{equation}
\end{lemma}
\begin{proof}
The parameter-dependent flow-box theorem on the two compact regular
segments gives positively oriented \(E_\theta,C_\theta\), with bounded
derivatives and reciprocal slopes on smaller fixed intervals.
The integral in \eqref{eq:divided-closing} extends
\(f_\theta(md_\theta(q))/m\) to \(m=0\), with positive slope
\(f_\theta'(md_\theta)d_\theta'\). Substitution proves both equations
and preserves isolated zeros and their multiplicities.
\end{proof}

At a zero, boundedness of \(\mathcal C_{\theta,m}\) gives
\(|x|\le Cm\).
In the inverse equation, the nonlinear closing term has an additional
factor \(m^2\):
\begin{equation}\label{eq:closing-quadratic}
 f_\theta(mt)=mf_\theta'(0)t+
 m^2t^2\int_0^1(1-u)f_\theta''(umt)\,du .
\end{equation}
The affine input keeps this small quadratic closing term explicit in
the later derivative comparisons.

The separate value-and-slope normalization
\eqref{eq:reference-normalization} is applied to the actual input height
and completed output height. It retains their two centers and gains;
only the input normalization is affine in the physical incoming height.

\subsection{Order and uniform transverse sections}
\label{sec:localization}

We first bound the through trajectories from above by an invariant graph.
Choose \(z_{\rm r}=\kappa\delta\), with
\(\delta\ll\kappa\ll1\). On
\(z_{\rm r}/2\le z\le2z_{\rm r}\), a coefficient shrink gives
\(|P|\le z_{\rm r}/8,\ Q\le-1/2,\ |G|\le C\delta\).
The orbit through \((\delta,z_{\rm r})\) consequently has
\(D\le-z_{\rm r}/8\) and \(|dz/dv|\le C\delta\).
It crosses the chart inside this band, by a first-exit argument.
Its graph \(\overline z_\theta(v)\) is the upper boundary; all
selected section windows are strictly below it.

\begin{lemma}\label{lem:loc-through}
The heights on \(\Sigma_+\) which reach \(\Sigma_-\) in the chart
below this upper boundary form an interval
\begin{equation}\label{loc:through-interval}
 \mathcal T_\theta=(\Theta_\theta,\overline z_\theta(\delta)] .
\end{equation}
Every such orbit has \(D<0\). Both \(T_\theta\) and
\begin{equation}\label{loc:speed-margin}
 \sigma_\theta(h)=\min_{|v|\le\delta}
          \{-D_\theta(v,z_\theta(v;h))\}
\end{equation}
are strictly increasing functions of the input \(h\).
\end{lemma}
\begin{proof}
Every physical input has \(D<0\) on the right section by
\eqref{loc:signs}. On the nullclines,
\begin{equation}\label{loc:nullcline}
 \dot D|_{D=0}=-zQG,\qquad \dot G|_{G=0}=G_vD .
\end{equation}
A first crossing from \(D<0\) to \(D>0\) requires \(G>0\);
equality would be an equilibrium. The quadrant \(D>0,G>0\)
is forward invariant by these identities, so this orbit cannot
subsequently reach the left section. A first terminal nullcline hit on
\(v=-\delta\) is also impossible: there \(\dot D<0\) by
\eqref{loc:signs}. A through orbit therefore has \(D<0\) everywhere
and solves
\begin{equation}\label{loc:graph-equation}
 z_v=-zG/D .
\end{equation}
Scalar uniqueness strictly orders these graphs. A graph above a through
graph has smaller \(D\), since \(Q<0\), and cannot cross either
\(D=0\) or the invariant upper boundary. It also passes through;
openness proves the interval assertion and ordering proves that \(T\)
increases. The positive difference of the two \(-D\)-values has a
positive minimum on the compact \(v\)-interval, so \(\sigma\) increases.
\end{proof}

The lower endpoint is a stopped limiting passage. A \emph{gap} is the
positive height difference from it on a specified transverse cut,
distinct from the central height.

\begin{lemma}\label{lem:loc-stopped-passage}
There are \(0<r<\delta\) and \(d_0,c_r>0\), independent of the
through orbit and of the sufficiently small coefficient, such that
every through orbit with endpoints in the selected outer windows and
\(\sigma_\theta<d_0\) crosses \(v=r\) and \(v=-r\) once, in order.
On \(r\le |v|\le\delta\) it has \(-D\ge c_r\).
The two outer transitions consequently have uniformly bounded
derivatives and inverse derivatives of any prescribed finite order
on slightly larger fixed domains.
\end{lemma}
\begin{proof}
The incidence set
\(\mathcal E_\theta=\{D=0,\ zG=0\}\) at \(\theta_0\) meets the
small closed chart only at the origin. Compactness puts the whole
nearby incidence set in \(|v|<r/3\).
Choose a neighborhood \(\mathcal O\) of it in \(|v|<2r/3\).
The endpoint lower bounds persist on the selected outer windows.
At an interior minimum of \(-D\) on a through graph,
\begin{equation}\label{loc:minimizer-identity}
 D\,D_v-QzG=0 .
\end{equation}
Transversality excludes a minimum near the fixed outer sections.
As \(\sigma_\theta\to0\), \eqref{loc:minimizer-identity} gives
\(|zG|\le C\sigma_\theta\); hence the minimizing point must enter
\(\mathcal O\). This proves the first assertion for some \(d_0\).
If the outer speed had no positive lower bound, a limiting orbit
would have a point with \(D=0\), \(D\le0\) on both time sides,
and thus \(\dot D=0\). By \eqref{loc:nullcline} this point lies
in \(\mathcal E\), a contradiction. The decreasing coordinate \(v\)
now gives all ordered hits. The hitting-time theorem on these
uniformly transverse compact slabs proves the derivative bounds.
\end{proof}

Fix nested section windows, whose interiors contain the corresponding
intersections of the organizing graphic, and speed bounds
\begin{equation}\label{loc:nested-buffers}
 K_\pm\Subset K_\pm^\sharp\Subset\Sigma_\pm,\qquad
 0<d_{\rm in}<d_{\rm out}<d_0 .
\end{equation}
The inner and outer sets specified by
\[
 h\in K_+,\quad T_\theta(h)\in K_-,\quad \sigma_\theta(h)\le d_{\rm in},
\]
and respectively by the interiors of \(K_\pm^\sharp\) and
\(\sigma_\theta(h)<d_{\rm out}\), are intervals, denoted by
\(J_\theta^{\rm in}\subset J_\theta^{\rm out}\).
All intervening transverse sections have nested windows as well.
The central height
\begin{equation}\label{loc:central-height}
 h_\theta(h)=z_\theta(\gamma(\theta);h)
\end{equation}
is positive and strictly increasing on \(J_\theta^{\rm out}\).
For any fixed raw height cap \(h_*>0\), define the affine physical
receiver
\begin{equation}\label{loc:analytic-receiver}
 I_\theta(h_*)=
 \{x:\ a_\theta+x\in J_\theta^{\rm out},
       \ 0<h_\theta(a_\theta+x)<h_*\}.
\end{equation}
It is an interval, possibly empty. Thus the whole interval between any
finite collection of its inputs satisfies the same input, inverse-output
and speed restrictions. These geometric restrictions define the domain
before any zero estimate is applied, and keep it on the through side of
the stopped boundary.

The estimates in the following sections choose a positive cap
\(h_{\rm out}\) in \eqref{loc:analytic-receiver} with strict buffers.
Fix this cap and the nested windows before the final independent raw
coefficient restrictions and the multiplier bound. Subsequent restrictions
retain these caps and section bounds. After these choices, the following
capture lemma places all nearby complete cycles below a smaller cap
\(0<h_{\rm in}<h_{\rm out}\).

\begin{lemma}\label{lem:loc-capture}
For every fixed set of these admissible nested windows and speed bounds in
\eqref{loc:nested-buffers}, and every \(h_{\rm in}>0\), there
are a coefficient neighborhood \(\mathcal V_0\) of \(\theta_0\) and
\(\eps_0>0\) such that, for every \(\theta\in\mathcal V_0\), every periodic
orbit \(C\) of \(H_\theta\) with
\(d_{\rm H}(\iota_0(C),\Gamma)<\eps_0\) follows the selected
positive-time circuit and has \(K>0\).
It meets the return section once in its prime period.
Its infinity arc has input in \(J_\theta^{\rm in}\), all selected
transverse hits in the inner windows, and
\(0<z(\gamma)<h_{\rm in}\).
\end{lemma}
\begin{proof}
Surround the two singularities by small blocks and the two compact
regular branches by flow boxes, with nested transverse windows.
In the finite normal-form box take weak faces \(\xi=-\rho\) and
\(\xi=\rho\), with transverse windows about the two organizing hits.
The weak speed at these fixed faces is positive after a coefficient
shrink. Thus they are respectively entrance and exit faces. The strong
faces \(\eta=\pm\eta_0\) point inward, since \(\eta'=-\eta\),
and are chosen away from the organizing arc. Remove from the two weak
faces the parts outside the transverse windows. These parts and the
strong faces are compact and disjoint from \(\Gamma\).

The infinity block has transverse faces \(v=\delta\) and
\(v=-\delta\), the invariant equator below, and the invariant graph
\(\overline z_\theta\) above. At the two selected face windows the
speed \(D\) is uniformly negative. Its upper face and the portions
of the transverse faces outside these windows stay a positive distance
from \(\Gamma\). The two regular flow boxes join the finite exit to
the right infinity entrance and the left infinity exit to the finite
entrance, in positive time. Their lateral faces have the same positive
distance property. Choose the Hausdorff radius below all these boundary
distances, which persist for nearby coefficients. A close cycle must
visit both regular branch boxes, because it approximates the complete
graphic. It can leave each visited box only through the selected next
face; the equator cannot be crossed by uniqueness. These entrances and
exits therefore force the finite channel, one regular branch,
\(\Sigma_+\to\Sigma_-\), and the other regular branch in this order.
The finite channel has no exit for \(K\le0\), so \(K>0\);
Lemma~\ref{lem:loc-through} gives \(D<0\) on the infinity passage.

Every intersection of the cycle with the selected connected segment
of \(\Sigma_+\) lies in its inner window and has the same transverse
orientation. A prime periodic orbit of a planar vector field is a
Jordan curve. Along a connected transverse segment its intersections
alternate between entering and leaving its bounded complementary
component. Hence these equally oriented intersections consist of one
point. The same circuit then contains one infinity arc, and its height
is one isolated zero of \eqref{eq:return} when the cycle is isolated.

If the uniform inner bounds failed, take
\(\theta_n\to\theta_0\) and \(C_n\) with
\(d_{\rm H}(\iota_0(C_n),\Gamma)\to0\).
Let \(A_n\) be the unique infinity-block subarc selected above and \(A_0\)
the complete organizing arc in that block.
Every Hausdorff limit of \(A_n\) is contained in \(A_0\).
Conversely every interior point of \(A_0\) is approached by
\(\iota_0(C_n)\), and the selected circuit forces the approaching
points to be on \(A_n\); closure includes the endpoints.
Thus \(d_{\rm H}(A_n,A_0)\to0\).
For a common modulus of continuity \(\omega\),
\[
 \left|\min_{A_n}(-D_{\theta_n})-\min_{A_0}(-D_{\theta_0})\right|
 \le \|D_{\theta_n}-D_{\theta_0}\|_\infty
          +\omega(d_{\rm H}(A_n,A_0)).
\]
Since the minimum on \(A_0\) is zero, \(\sigma_{\theta_n}(A_n)\to0\).
The hits of \(v=\gamma(\theta_n)\) converge to \((0,0)\), the unique
hit of \(v=0\) on \(A_0\), and all transverse hits enter the inner
windows by continuity. This contradiction supplies one coefficient
neighborhood and one Hausdorff radius.
\end{proof}

\subsection{Transferring a slice bound}

The same geometric transfer will be used for the full coefficient cover
and for its ratio-restricted parts.

\begin{lemma}\label{lem:slice-bound-transfer}
Let $\mathcal U_5$ be a neighborhood of $\theta_0$, let
$\mathcal B$ be any subset of the slice parameter space, and fix
$\eps_5>0$ and an integer $N_{\rm cyc}\ge0$.
Suppose every $\theta\in\mathcal U_5\cap\mathcal B$ has at most
$N_{\rm cyc}$ distinct isolated cycles at whole-graphic Hausdorff
distance less than $\eps_5$ from $\Gamma$.
There are one open neighborhood $\mathcal V$ of $F_*$ in
$\mathcal Q_2$ and one $\eps>0$ such that every $F\in\mathcal V$
with $\theta(\mathscr C_0(F))\in\mathcal B$ has at most
$N_{\rm cyc}$ distinct isolated periodic orbits $C$ satisfying
$d_{\rm H}(\iota_*(C),\Gamma_{I4})<\eps$.
The transport preserves the complete
marking, positive time and return-zero multiplicities. The maps and
their inverses also preserve whole-orbit Hausdorff convergence to the
marked graphics as the fields tend to $F_*$.
\end{lemma}
\begin{proof}
For $H=\mathscr C_0(F)$ use the analytic decomposition in
Lemma~\ref{lem:slice} for this field. Enlarge $L_0$ to bound both $S$ and
$S^{-1}$, and use the common bound $L_*$ for
$T_{\mathfrak A(H)}$ and its inverse. Continuity gives an open
neighborhood $\widehat{\mathcal V}$ of $H_0$ on which
\begin{equation}\label{eq:lift-neighborhood}
 \theta(H)\in\mathcal U_5,\qquad
 \eta(H):=d_{\rm H}(T_{\mathfrak A(H)}\Gamma,\Gamma)<\eps_5/2.
\end{equation}
Set
\begin{equation}\label{eq:original-coefficient-neighborhood}
 \mathcal V=\mathscr C_0^{-1}(\widehat{\mathcal V}),
\end{equation}
an ordinary open neighborhood in all twelve coefficients, and take
\begin{equation}\label{eq:original-collar-radius}
 \eps=\frac{\eps_5}{2L_*L_0}.
\end{equation}
The fixed change and $w=LZ+c$ send a cycle $C$ of this same field
to a cycle $\widetilde C$ of $H_{\theta(H)}$, with
\[
 \iota_0(\widetilde C)=T_{\mathfrak A(H)}S(\iota_*(C)).
\]
Writing $d_{\rm original}=d_{\rm H}(\iota_*(C),\Gamma_{I4})$
and $d_{\rm slice}=d_{\rm H}(\iota_0(\widetilde C),\Gamma)$,
we have
\begin{align}
 d_{\rm slice}&\le L_*L_0d_{\rm original}+\eta(H)<\eps_5,
                                      \label{eq:two-map-hausdorff}\\
 d_{\rm original}&\le L_0L_*(d_{\rm slice}+\eta(H)).
                              \label{eq:two-map-hausdorff-reverse}
\end{align}
The affine diffeomorphisms and positive time factors biject cycles,
preserving distinctness, isolation, the two branches and their marked
order. The induced transverse return maps are locally conjugate, so
zero multiplicities are preserved as well. The first inequality applies
the slice bound; both inequalities and $\eta(H)\to0$ give the
convergence assertion as the fields tend to the organizing field.
\end{proof}

\section{Counting on one physical interval}
\label{sec:counting-tools}

The upper bound follows from two elementary counting criteria on a
connected interval. One uses the orientation of displacement critical
points; the other uses the orientation of curvature zeros. Each criterion
counts the whole physical interval, including the regions between the
captured roots.

\subsection{A common coordinate and the counting criteria}

Fix the incoming raw height coordinate in Lemma~\ref{lem:closing} before
normalizing the tails. Absorb the regular maps between the fixed raw ports
and the chosen sections into the actual adapters. All regional estimates
will then concern the displacement
\begin{equation}\label{sharp:upper-displacement}
 F(x)=x-m\Pi(x),\qquad m>0.
\end{equation}
Every local estimate below includes the divided adapter
\eqref{eq:divided-closing}. Its derivatives remain bounded, with positive
first derivative, as \(m\downarrow0\). Apply the local tail estimates
with these adapters and normalize each map by its exact center and first
variation. The relation with the physical map is the identity
\begin{equation}\label{sharp:upper-affine}
 \Pi(a_++c_+w)=a_-+c_-\widehat\Pi(w),\qquad c_\pm>0.
\end{equation}
Thus the signs of second and third derivatives transfer by the positive
factors \(c_-/c_+^2\) and \(c_-/c_+^3\). The affine relation is
essential here: independently chosen nonlinear coordinates need not
preserve curvature. Comparisons obtained in an inverse coordinate are
therefore transferred to the direct displacement before counting its
critical points.

All carriers below are intervals of actual through trajectories in the
nested section windows. Scalar first variation orders them by the central
height. The preimage of an output interval under an increasing passage map
is an interval. Consequently the interval between any finite set of
captured roots stays in the same enlarged physical domain. The estimates
used between roots therefore apply at nonroot points as well.

\begin{lemma}\label{sharp:upper-critical-count}
Let \(F\) be \(C^3\) on a connected carrier \(I\). If every critical
point satisfies \(F''>0\), then \(F\) has at most one critical point,
which is simple and upward, and at most two isolated zeros, counted with
multiplicity. More generally, let \(J\subseteq I\) be a connected interval,
possibly empty or all of \(I\). If \(F'\) has at most \(j\)
zeros on \(J\), counted with multiplicity, where \(0\le j\le2\), and \(F'>0\) on
\(I\setminus J\), then \(F'\) has at most \(j\) zeros on \(I\), and
\(F\) has at most \(j+1\) isolated zeros there, counted with multiplicity.
For \(J=\varnothing\), take \(j=0\): there is no critical point and
at most one simple displacement zero.
\end{lemma}
\begin{proof}
In the first case every zero of \(F'\) is a simple upward crossing.
Two such crossings would require an intervening zero with another
orientation, so there is at most one. In the second case no critical
point lies outside \(J\), whether its complement is empty or lies on
one or both sides. Rolle's theorem with
multiplicities proves both assertions.
\end{proof}

\begin{lemma}\label{sharp:upper-curvature-count}
Let \(F=x-m\Pi\), \(m>0\), be \(C^3\) on a connected carrier \(I\),
and let \(J\subseteq I\) be a connected interval, possibly empty or
all of \(I\). Suppose that every zero of \(\Pi''\) on \(J\) is simple
and has the same crossing orientation, and that \(F'>0\) on
\(I\setminus J\). Then \(\Pi''\) has at most one zero on \(J\),
\(F'\) has at most two zeros on \(I\), counted with multiplicity,
and \(F\) has at most three
isolated displacement zeros on \(I\), counted with multiplicity.
If \(J=\varnothing\), there is no critical point and at most one
simple displacement zero.
\end{lemma}
\begin{proof}
For nonempty \(J\), two simple curvature zeros of the same orientation
would require an intervening zero with another orientation, so there is
at most one. Since \(F''=-m\Pi''\), Rolle's theorem with multiplicities
gives at most two zeros of \(F'\) on \(J\). Apply
Lemma~\ref{sharp:upper-critical-count} on \(I\). For empty \(J\),
\(F'>0\) throughout \(I\). Positive affine changes in
\eqref{sharp:upper-affine} preserve the curvature orientation, even if
their scale ratios degenerate.
\end{proof}

The regional constructions specify the actual intervals \(I,J\) and
establish the hypotheses at every point of those intervals. For each
coefficient fibre, apply Lemma~\ref{sharp:upper-critical-count} or
Lemma~\ref{sharp:upper-curvature-count} to its whole carrier. The later
coefficient cover takes the maximum of these bounds; its overlap choices
are independent of the phase input.

\subsection{Ordinary adapters and divided closing maps}

The next two lemmas control the regular maps appended to a singular
passage. They use ordinary third derivatives of those maps and allow the
singular passage's first gain to tend to zero.

\begin{lemma}\label{phys:lem-c3-adapters}
Let the ordinary adapters below have uniform \(C^3\) norms and positive
first derivatives bounded away from zero on their original buffered
images. All parameters and gains are frozen under the phase derivation
\(D\). The following identities and bounds are uniform in those
parameters.

For an actual raw output \(T\) in such an image, put
\(\ell=\log\psi'\). Then
\begin{equation}\label{phys:eq-observer-log}
 \begin{aligned}
 D\log\frac{(\psi\circ T)'}{T'}&=\ell'(T)DT,\\
 D^2\log\frac{(\psi\circ T)'}{T'}
 &=\ell''(T)(DT)^2+\ell'(T)D^2T .
 \end{aligned}
\end{equation}
Here \(T'>0\) is the exact raw first gain; the formulas impose no lower
bound on it. If the full translated output and its two absolute phase
derivatives are bounded, and \(mT\) stays in the fixed closing buffer,
then
\begin{equation}\label{phys:eq-absolute-closing-two}
 \sum_{j=0}^2\left(
 |D^j(f(mT)-f(0))|
 +|D^j(\log f'(mT)-\log f'(0))|\right)\le C m.
\end{equation}
The constant depends on these absolute output bounds, not on a
normalized raw gain or a ratio of gains.

For \(f(0)=0\), the divided closing map \(C_m(q)=f(md(q))/m\),
extended at zero by \(C_0=f'(0)d\), has jets
\begin{equation}\label{phys:eq-closing-jets}
 \begin{aligned}
 C_m'&=f'(md)d',\\
 C_m''&=m f''(md)(d')^2+f'(md)d'',\\
 C_m'''&=m^2f'''(md)(d')^3
       +3m f''(md)d'd''+f'(md)d''' .
 \end{aligned}
\end{equation}
They are uniformly bounded and the first is a positive unit on each
prescribed compact \(q\)-interval with buffered \(md\)-image.
Finally, for a coefficient-only center \(z_0\), \(0<d\le1\), and a
fixed compact phase interval whose \(z_0+dy\)-image stays buffered,
\begin{equation}\label{phys:eq-observer-normalization}
 N_d(y)=\frac{\psi(z_0+dy)-\psi(z_0)}{d\psi'(z_0)},\qquad
 N_0(y)=y,\qquad \|N_d-\mathrm{id}\|_{C^3}\le C d.
\end{equation}
This last bound concerns phase derivatives, not derivatives of the
center or scale with respect to a parameter.
\end{lemma}
\begin{proof}
The chain rule proves \eqref{phys:eq-observer-log} and
\eqref{phys:eq-closing-jets}; the integral formula
\(C_m=d\int_0^1f'(tmd)\,dt\) gives the extension at zero.
The ordinary \(C^3\) bounds and positive slopes then give
\eqref{phys:eq-absolute-closing-two} by the mean value theorem.
Taylor's integral formula gives
\eqref{phys:eq-observer-normalization}, whose second and third
rows contain factors \(d\) and \(d^2\), respectively.
\end{proof}

\begin{lemma}\label{phys:lem-c3-core-transport}
Let \(M\) be an intrinsic core map with uniformly bounded \(C^4\)
norm on a fixed buffered phase interval. Let \(A,B\) be ordinary
\(C^3\) adapters close to the identity, with uniformly positive
slopes, on buffers containing all compositions below. On a common
compact interval for \(A^{-1}\) and \(M\circ A^{-1}\),
\begin{equation}\label{phys:eq-c3-core-transport}
 \|B\circ M\circ A^{-1}-M\|_{C^3}
 \le C\bigl(\|A-\mathrm{id}\|_{C^3}
              +\|B-\mathrm{id}\|_{C^3}\bigr).
\end{equation}
The fourth derivative is required only of the intrinsic core.
\end{lemma}
\begin{proof}
The inverse formulas through order three and the uniform slope bound
give \(\|A^{-1}-\mathrm{id}\|_{C^3}\le C\|A-\mathrm{id}\|_{C^3}\).
In the third chain-rule row for \(M\circ A^{-1}-M\), the translated
factor \(M'''(A^{-1})-M'''\) is bounded by
\(\|M''''\|_\infty\|A^{-1}-\mathrm{id}\|_\infty\); the other
rows need fewer core derivatives. The ordinary composition of
\(B-\mathrm{id}\) with \(M\circ A^{-1}\) is bounded in \(C^3\)
without translating any fourth derivative of \(B\). These estimates
prove \eqref{phys:eq-c3-core-transport}.
\end{proof}

\subsection{Weighted curvature of a completed passage}

To use an action calculation in the physical input, we must include the
curvature contributed by each section map. The following identity keeps
these contributions separate and allows a positive weight adapted to
the coefficient chart.

\begin{lemma}[Weighted curvature of a completed passage]
\label{neg:lem-completed-curvature}
Let \(u,v>0\) be radial coordinates, with \(P_\pm,B_\pm>0\)
on the two strict through intervals, and suppose
\[
 v_u=\frac{uB_+(u)P_-(v)}{vB_-(v)P_+(u)}(1+\varrho)>0,
 \quad S(-u^2/2)=-v^2/2,\quad
 \Pi=\phi_-\circ S\circ\phi_+^{-1},
\]
Here \(P_\pm,B_\pm\) are \(C^2\), \(\varrho=\varrho(u,v)\)
is \(C^2\) on the product of the two radial intervals, and
\(\phi_\pm\) are the separately normalized actual \(C^3\) section
maps with positive derivatives. Put
\(\mathcal D_+=P_+\partial_u/(uB_+)\),
\(\mathcal D_-=P_-\partial_v/(vB_-)\), and \(R_B=B_+/B_-\).
For any smooth positive weight \(W=W(u,v)\), one has exactly
\begin{align}
 \phi_+'\frac{P_+}{W}\frac{\Pi''}{\Pi'}
 ={}&\frac{P_+' /u-R_B(1+\varrho)P_-'/v}{W}
 -\frac{P_+}{uW}\frac d{du}\log R_B \notag\\
 &-\frac{B_+}{W}
       \{\mathcal D_++(1+\varrho)\mathcal D_-\}\log(1+\varrho)
       \notag\\
 &-\frac{B_+(1+\varrho)}{W}\mathcal D_-\log\phi_-'
 +\frac{P_+}{u^2W}(u\partial_u)\log\phi_+'.
 \label{deg:weighted-exact-curvature-chain}
\end{align}
Primes on each \(\phi_\pm\) use its own quadratic label; primes on
\(P_\pm\) use the displayed radial variable. The derivative \(d/du\)
follows the matched graph. In particular both section terms remain in
the curvature of the actual affine input.
\end{lemma}
\begin{proof}
One has \(S'=R_BP_-(1+\varrho)/P_+\). Logarithmically differentiate
this identity using \(d/db_+=-u^{-1}d/du\), and substitute in
\[
 \phi_+'\frac{\Pi''}{\Pi'}=
 \frac{S''}{S'}+S'\frac{\phi_-''}{\phi_-'}-
                         \frac{\phi_+''}{\phi_+'}.
\]
Multiplication by \(P_+/W\) gives the formula. Its coefficients need
not extend without weights to a stopped endpoint.
\end{proof}



\section{Action curvature and coefficient geometry}
\label{sec:small-positive-source-dictionaries}

We first define the coefficient coordinates and passage scales. The
representative cubic calculation then explains how the limiting action
controls the curvature of the completed physical return.
\subsection{Coefficient geometry and the regime decision map}
\label{sec:coefficient-regimes}

Different relative sizes of the five coefficients in \eqref{eq:slice}
lead to different passage scales. We introduce these scales here and
summarize the resulting cases in Figure~\ref{fig:coefficient-phase-decision}.
The corresponding bounds are proved in the later face propositions.
Before counting zeros, every phase chart is transported to the physical
sections of \eqref{eq:return}.

Translate the zero of the transverse eigenvalue to the origin and put
\begin{equation}
 \xi=v-\gamma,
 \qquad \alpha _0=P(\gamma),
 \qquad \alpha _1=P'(\gamma),
 \qquad q_\gamma=1+2\gamma,
 \qquad \ell(v)=1-Av+Cv^2.
 \label{deg:central-coordinates}
\end{equation}
Then \eqref{eq:raw} becomes
\begin{equation}
 \begin{split}
  \dot \xi&=\alpha _0+\alpha _1\xi-3\gamma \xi^2-\xi^3
                 -\ell(\gamma+\xi)z,\\
  \dot z&=-z\{q_\gamma \xi+\xi^2+C(\gamma+\xi)z\}.
 \end{split}
 \label{deg:central-field}
\end{equation}
Thus \((A,C,\gamma,\alpha _0,\alpha _1)\) are analytic coordinates on
the five-parameter slice near the organizing field. After one common
coefficient shrink, \(q_\gamma\) and \(\ell\) are positive units on the
fixed raw rectangle, and the regular transition maps between fixed sections
and their inverses have uniform bounds through every derivative order fixed
in advance.

For negative opening write
\begin{equation}
 \alpha _0=-s^2,
 \qquad p=\frac{\alpha _1}{s},
 \qquad v=\gamma+sX,
 \qquad z=s^2Y.
 \label{deg:source-dictionary}
\end{equation}
This gives the source field \eqref{ann:source-field} with parameters
\((a,p,\gamma,s)=(-1,p,\gamma,s)\). Two exact compactifications cover
its small-\(p\) coefficient ends. For \(\gamma\ne0\), set
\begin{equation}
 g=|\gamma|,\qquad k=\frac{s}{g^{3/2}},\qquad
 \pi=\frac{p}{\sqrt g},\qquad \mathfrak q=k\pi .
 \label{deg:gamma-chart}
\end{equation}
Here \(g\) is always a coefficient radius, never a section height. In the
cubic chart set
\begin{equation}
 t=s^{1/3},\qquad U=\frac{p}{t},\qquad
 G=\frac{\gamma}{t^2},\qquad \tau_{\rm cub}=t^2.
 \label{deg:cubic-chart}
\end{equation}
We write \(\tau_{\rm cub}\) for this cubic scale throughout. The
far-root, two-saddle and degree-loss arguments each use an unadorned
\(\tau\) for their local scale; any relation such as \(\tau=p^2\)
is stated where that scale is introduced.

\begin{figure}[H]
\centering
\small
\fbox{\begin{minipage}{0.94\linewidth}
\centering
\(\theta\) in the five-parameter slice
\(\longrightarrow (\gamma,\alpha _0,\alpha _1;A,C)\)

\medskip
\begin{tabular}{@{}p{.27\linewidth}@{\(\longrightarrow\)}p{.62\linewidth}@{}}
\shortstack[l]{\(K\le0\) or no\\through passage}
 & stopped/no-passage endpoint; it is not a periodic-orbit zero\\[2pt]
\(K>0\), high central height
 & affine high-height estimate, Corollary~\ref{rew:cor-high-derivative-signed}\\[2pt]
\(\alpha _0\ge0\), low height
 & determinant chart or signed trace chart,
   Lemma~\ref{deg:lem-nonnegative-domain}\\[2pt]
\(\alpha _0<0,\ p\le0\)
 & raw Jacobian contraction, Lemma~\ref{deg:lem-raw-jacobian}\\[2pt]
\(\alpha _0<0,\ 0<p\ll1\)
 & gamma-dominant, cubic, degree-loss, or remote-root face,
   Lemma~\ref{deg:lem-small-source-routing}\\[2pt]
\(\alpha _0<0,\ p\asymp1\)
 & compact positive-ratio passage,
   Lemma~\ref{ann:lem-compact-positive-domains}\\[2pt]
\(\alpha _0<0,\ p\gg1\)
 & signed-trace terminal block plus compact annuli,
   Lemma~\ref{deg:lem-large-source-routing}
\end{tabular}

\medskip
Each row \(\longrightarrow\) zeros of the same exact displacement
\(x-m\Pi(x)\) on one coefficient fibre
\(\longrightarrow\) one bound on each whole phase fibre, then the
maximum across coefficient alternatives
\(\longrightarrow\) one common coefficient neighborhood and one common
Hausdorff collar
\(\longrightarrow\) lift from five to twelve coefficients.
\end{minipage}}
\caption{Global coefficient--phase decision map. The arrows specify the
logical routing, not additional coordinate changes. Strict overlaps between
adjacent phase and coefficient blocks are recorded in
Table~\ref{tab:coefficient-phase-cover}.}
\label{fig:coefficient-phase-decision}
\end{figure}

For later reference, the symbols with global scope are
\(\theta=(A,C,\mu_1,\mu_2,\mu_3)\),
\((\gamma,\alpha _0,\alpha _1)\), the finite multiplier \(m\), and the
actual input of \eqref{eq:return}. Symbols such as \(g,k,\pi,\mathfrak q\),
\(t,U,G\), and the degree-loss exponents are local to the chart in which
they are defined. Actual phase heights carry labels such as \(h\),
\(h_L\), \(V\), \(\widehat V\), \(W\), or \(\Theta\); they are never
identified with a coefficient radius or with a stopped trace gap.

\subsection{The cubic action and the physical curvature}
\label{subsec:representative-cubic}

The compact cubic block illustrates the curvature mechanism behind the
three-zero bound. At the scale \eqref{deg:cubic-chart}, its axis
polynomial is
\[
 P(\xi)=1-U\xi+3G\xi^2+\xi^3,
 \qquad U\ge0,
\]
with \(U,G\) in a fixed compact set. On a strict through interval,
let \(u>0\) and \(-r<0\) be the incoming and outgoing axis points.
The limiting action equation is
\begin{equation}\label{intro:compact-action}
 \int_0^u\frac{\xi}{P(\xi)}\,d\xi
 =\int_0^r\frac{\xi}{P(-\xi)}\,d\xi,
 \qquad b=-u^2/2,\qquad S(b)=-r^2/2.
\end{equation}
On the traversed segment \(P>0\). Differentiating gives
\(r_u=uP(-r)/(rP(u))>0\) and
\begin{equation}\label{intro:compact-action-curvature}
 S'=\frac{P(-r)}{P(u)},\qquad
 \frac{P(u)}{u+r}\frac{S''}{S'}=3-\frac{U}{ur},\qquad
 (ur)_b=-\frac ru-\frac{uP(-r)}{rP(u)}<0.
\end{equation}
These are the cubic specialization of
\eqref{neg:action-curvature}--\eqref{neg:entry-exit-monotonicity}.
The term containing \(G\) cancels. At a curvature zero, \(U=3ur\),
so the derivative of the middle expression with respect to \(b\)
is strictly negative. Thus every curvature zero is a downward crossing,
and there can be at most one on a connected interval.

The actual input is an affine raw section coordinate. To compare its
curvature with the action formula, write the two separate outward maps as
\[
 \Gamma_\pm=a_\pm+c_\pm\phi_\pm,\qquad c_\pm>0,
 \qquad \widehat\Pi=\phi_-\circ S_{\rm act}\circ\phi_+^{-1}.
\]
Here \(S_{\rm act}\) is the actual matched transition in the two normal
section labels. The centers and gains are retained exactly. In these
coordinates the second-derivative chain rule is
\begin{equation}\label{intro:completed-curvature-chain}
 \phi_+'\frac{\widehat\Pi''}{\widehat\Pi'}
 =\frac{S_{\rm act}''}{S_{\rm act}'}
  +S_{\rm act}'\frac{\phi_-''}{\phi_-'}
  -\frac{\phi_+''}{\phi_+'}.
\end{equation}
Each section derivative uses its own quadratic label. The actual gains
do not enter this identity. Returning to the physical affine input
multiplies its second and third derivatives by
\(c_-/c_+^2\) and \(c_-/c_+^3\), so their signs are preserved.

On its bounded normal-section interval,
Proposition~\ref{deg:prop-compact-cubic} proves the weighted identity
\begin{equation}\label{intro:actual-cubic-curvature}
 \phi_+'(b)\frac{P(u)}{u+r}
              \frac{\widehat\Pi''}{\widehat\Pi'}
 =3-\frac{U}{ur}+E(u,r),\qquad
 \|E\|_{C^1_{u,r}}\le C\tau_{\rm cub}(1+|\log\tau_{\rm cub}|).
\end{equation}
It includes both section terms in
\eqref{intro:completed-curvature-chain}. Near a possible zero the
principal derivative with respect to \(u\) has a margin
\(c(1+r_u)\), whereas the error derivative along the matched graph is
at most \(C\tau_{\rm cub}(1+|\log\tau_{\rm cub}|)(1+r_u)\).
The physical input \(\phi_+(-u^2/2)\) decreases with \(u\).
Consequently the actual curvature retains the same downward crossing.
The bounded-action comparison overlaps this description on the same
map. The rest of the admitted interval has bounded physical first
derivative; choosing \(m\) last makes \(F'>1/2\) there.
Lemma~\ref{sharp:upper-curvature-count} then gives three zeros on the
whole interval.

Two estimates make this calculation compatible with singular
coefficient limits. First, the exterior continuations carry the common
first-variation factor \(\tau_{\rm cub}^{2-\lambda_G}\), where
\(\lambda_G=12\tau_{\rm cub}G/(1+2\tau_{\rm cub}G)\).
The common power cancels in the completed first derivative; the actual
normalized section maps have ordinary \(C^3\) bounds. The general
raw cancellation is displayed in \eqref{sharp:upper-tail-ratio}.
Second, near an axis root the normal height has an action exponential.
The proof bounds derivatives on the product of the two separate section
domains, using this exponential to pay the losses in each root distance.
Only then is the estimate restricted to the actual matching graph.
This controls the two root distances independently and includes the
dependence on the moving entrance.

The other coefficient faces use the same physical chain rule with
different positive weights, as recorded in
Lemma~\ref{neg:lem-completed-curvature}. On the reversible face the
principal curvature vanishes. There the exact first integral permits
division by the coefficient distance before taking the limit, giving
the relative estimate \eqref{neg:relative-estimate}.
The coefficient cover selects one of these estimates for each whole
coefficient fibre.

\subsection{The complete physical return in the compact cubic block}

Fix the passive coefficient set and the margins in the raw sections.
Let \(z_*>0\) be the
unscaled height cap supplied by Corollary~\ref{rew:cor-high-derivative-signed};
its high-height estimate is used below after the cubic low-height analysis.

\begin{proposition}[Compact cubic coefficients]
\label{deg:prop-compact-cubic}
Fix \(U_0,G_0<\infty\), the original passive compact set and buffered
sections. For \(0\le U\le U_0\), \(|G|\le G_0\),
\(p=tU,\gamma=t^2G,s=t^3\), put \(\tau=\tau_{\rm cub}=t^2\).
Choose a positive unscaled raw cap \(0<h_c<z_*\) before the final
coefficient and multiplier bounds. On the entire original interval
\[
 I_\theta=\{x:a_\theta+x\in J_\theta^{\rm out},\quad
                 0<H(x)=z_\theta(\gamma;a_\theta+x)<h_c\},
 \qquad F(x)=x-m\Pi_{\theta,m}(x),
\]
the connected lower curvature receiver satisfies
\(\Pi''=0\Rightarrow\Pi'''<0\). Its bounded and growing descriptions
have actual resection overlap. The remaining admitted ghost raw tail has
\(0<\Pi'\le C(d_0)\tau^3\); its upper core and signed high complement
have bounded complete derivative and strictly overlap. After the common
multiplier shrink the same \(F'\) has at most two zeros, and \(F\)
has at most three isolated zeros, counted with multiplicity, on the whole
coefficient fibre. All original window, inverse-image, speed, stop and
height clips and intervening nonroot inputs are retained. Real-root and
fold endpoints are one-sided; the stopped fold continuation is excluded.
\end{proposition}

\begin{proof}
We construct normal sections on which the action controls curvature,
then continue their estimates to the original raw sections. Near the
roots, the two section variables must be estimated independently before
they are matched. The ghost tail and upper continuation require only a
bound on the complete first derivative. These estimates will cover one
connected physical interval to which the counting lemma applies.

In raw coordinates \(v-\gamma=\tau\xi,z=\tau^2\eta\), the exact field is
\begin{equation}\label{deg:compact-cubic-field}
 \dot\xi=-D,\quad\dot\eta=-\eta H,\qquad
 D=\tau P+Q\eta,\quad H=\xi B+C\tau^2(G+\xi)\eta,
\end{equation}
\begin{equation}\label{deg:compact-cubic-polynomials}
 P=1-U\xi+3G\xi^2+\xi^3,\quad B=1+\tau(2G+\xi),\quad
 Q=1-A\tau(G+\xi)+C\tau^2(G+\xi)^2.
\end{equation}
The unique simple negative root is \(-c\), by Descartes' rule, and
\begin{equation}\label{deg:compact-factorization}
 P(\xi)=(\xi+c)[(\xi-\alpha)^2-\Delta],\quad
 c-2\alpha=3G,\quad c(\alpha^2-\Delta)=1,\quad
 2\alpha c=U+c^{-1}.
\end{equation}
The four positive quantities \(c,c^{-1},\alpha,\alpha^{-1}\) are units.
A triple root would force \(P=(\xi+1)^3,U=-3\), outside this block.

Choose \(\ell_0\) so the original positive input window exceeds the
first-lip height at its raw port with a fixed margin. The actual lips
\(\eta=\tau\ell_0P\), with abscissae \(u,-r\), obey
\begin{equation}\label{deg:compact-lip-speed}
 \frac d{dt}\log\frac\eta P=-\xi B+\tau P'(1+Q\ell_0)
                              -C\tau^3(G+\xi)\ell_0P.
\end{equation}
Beyond a fixed positive threshold the signs are negative on the positive
lip and positive on the negative lip; their first hits are unique. At
\(a=\tau u\) the raw height is
\begin{equation}\label{deg:compact-lip-height}
 z=\ell_0[a^3+3G\tau a^2-U\tau^2a+\tau^3].
\end{equation}
Choose \(0<d_{\rm in}<d_0<\delta\), then \(\ell_0\delta\ll1\),
and finally \(\tau\ll d_{\rm in}\). Comparison with the organizing
arc and Lemma~\ref{lem:loc-capture} place captured low-height cycles at
\begin{equation}\label{deg:compact-raw-range}
                         0<a=\tau u<d_{\rm in}.
\end{equation}
The analytic carrier still includes all its original through inputs.
To continue them to the raw sections, set
\(v-\gamma=\sigma q\), \(\sigma=\pm1\), and
\(\mathcal W=z/q^2\). The graph equation becomes
\begin{equation}\label{deg:compact-raw-corridor}
 q\mathcal W_q=\mathcal W
 \frac{B_\sigma(q)+\sigma Cv_\sigma(q)q\mathcal W}
      {\ell(v_\sigma(q))\mathcal W+M(\sigma q)/q^2}-2\mathcal W.
\end{equation}
Here
\[
 v_\sigma(q)=\gamma+\sigma q,\qquad
 B_\sigma(q)=1+2\gamma+\sigma q,\qquad
 M(x)=x^3+3\gamma x^2-spx+s^2,
\]
and \(\ell\) is defined in \eqref{deg:central-coordinates}.
At fixed relative distance from a lip this is a small perturbation of
\(q\mathcal W_q=1-2\mathcal W\); constants around \(1/2\) give a
strict invariant strip to both original raw sections. The negative
continuation crosses its axis root at positive height. The original
windows and scalar order supply one through component.

\paragraph{Normal sections and independent root distances.}
In the bounded normal-section calculation below, \(p\) denotes the
positive axis footpoint; the fixed coefficient \(p=tU\) is represented
by \(U\) in \eqref{deg:compact-cubic-polynomials}.
Choose \(v>0\) with \(P\ge3/4\) on \([-4v,4v]\).
The limiting action exit for \(p=v\) lies in \([v/2,2v]\), because
both densities are \(s(1+O(v))\). Opposite strict action brackets and
Corollary~\ref{neg:cor-rootfree} give an actual comparison orbit,
independently of the return equation, with both footpoints in fixed
subintervals of \((0,3v)\).
Fix normal boxes and independent weak cuts \(p_w,r_w\) below these
footpoints. Scalar order selects the lower receiver, whose footpoints
have fixed positive lower bounds. The strong projections and normal
linearization give
\[
 \begin{gathered}
 \dot p=-\tau P(p),\quad\dot\zeta=-pB(p)\zeta,\quad\zeta_{\rm section}=u_+(p),\\
 J_+(p)=\int_{p_w}^p\frac{\xi B}{P}\,d\xi,\qquad
 J_-(r)=\int_{r_w}^r\frac{sB(-s)}{P(-s)}\,ds.
 \end{gathered}
\]
The projections have integrable kernels
\(C(1+T)^N e^{-(c_*-C\tau)T}\). To retain the higher rows, write
the attracting field as
\(V=(h-qy,-ny-ky^2)\), where \(h=O_{C^6}(\tau)\),
\(n\ge c_*>0\), and \(q,n,k\) have bounded finite derivatives.
Apply the variational induction in the proof of
Lemma~\ref{lem:normal-time-projection} to this attracting field.
The normal forcing always contains a height or a lower normal
variation, while pure tangential forcing has an \(O(\tau)\)
coefficient. Integrating the normal row first gives the same integrable
Volterra kernel. At phase order five it yields tangential rows
\(C(1+T)^Ne^{C\tau T}\) and normal rows
\(C(1+T)^Ne^{-(c_*-C\tau)T}\), using axis-flow order six.
A selected coefficient derivative requires mixed order seven of the
explicit coefficients and adds only polynomial time factors to this
majorant. Taylor division thus gives the required \(C^4\) units
\(u_\pm,u_\pm^{-1}\). The normal gap persists at the fold; these
coordinates are defined on the actual normal sections.

The action exponential must absorb derivative losses as either endpoint
approaches a root. We estimate those losses in five overlapping charts.
For \(\Delta=s_0^2\), \(a_0=\alpha-s_0\), \(X=a_0-p\), use
\((X,s_0/X)\) and \((s_0,X/s_0)\); the singular action integral is
\[
 \frac1{2s_0}\log\frac{\ell(X+2s_0)}{X(\ell+2s_0)},
                    \qquad s_0=0:\quad X^{-1}-\ell^{-1}.
\]
For \(\Delta=-\kappa^2\), use the left, central and right charts
\((X,\kappa/X)\), \(p=\alpha+\kappa z,|z|\le Z_0\), and
\((\widetilde X,\kappa/\widetilde X)\). Their arctangent action is
\(\kappa^{-1}[\arctan((p-\alpha)/\kappa)+\arctan(\ell/\kappa)]\).
Restoring the smooth positive density gives lower bounds
\(c/X-C\), \(cs_0^{-1}(1+|\log(X/s_0)|)-C\), or
\(c_{Z_0}/\kappa-C\). Every required ordinary phase derivative loses
only fixed powers of these independent variables. For example the real
inner loss is bounded by
\(s_0^{-N}z^{-N}e^{-c(1+|\log z|)/(2\tau s_0)}\), uniformly after
\(c/(2\tau_0s_{0,\max})>N+1\). Thus \(e^{-J_+/(2\tau)}\) pays
all root-distance losses, without comparing the two gaps.

\paragraph{Matching through the weak core.}
To compare the two actions, we need a core map defined all the way to
zero height. We construct that map on the whole chord below each
admitted orbit, before choosing an action cutoff. In raw height the
fixed weak-core equation is
\[
 \eta_\xi=f_\tau(\xi,\eta)
 =\frac{\eta(a+c_1\eta)}{T+Q\eta},\qquad
 a=\xi B,\quad c_1=C\tau^2(G+\xi),\quad T=\tau P.
\]
Here \(P\ge p_c>0\), \(Q\ge q_c>0\), and the coefficients are
bounded on a fixed enlarged weak slab. Direct differentiation gives
\[
 \partial_\eta f_\tau
 =\frac{aT+c_1\eta(2T+Q\eta)}{(T+Q\eta)^2},\qquad
 |\partial_\eta f_\tau|
 \le\frac{|a|}{T}+\frac{|c_1|}{Q}\le\frac C\tau,
 \qquad \eta\ge0.
\]
The constants use neither endpoint root distance.

Retain the actual intrinsic weak cuts from the projection construction
above. Write their raw parametrizations as
\(\iota_\pm(w)=(\psi_\pm(w),h_\pm(w))\), where \(w=\zeta_\pm\),
\(\psi_+(0)=p_w\), \(\psi_-(0)=-r_w\), and \(h_\pm(0)=0\).
The projection and normal linearization above, constructed by the
backward integrals of Lemma~\ref{lem:normal-time-projection} with the
displayed higher rows, give uniformly \(C^4\) inverse charts;
\(h_\pm'\) and the chart Jacobians
\(\Delta_\pm=\det D_{(v,\eta)}(\pi_\pm,\zeta_\pm)\),
in the respective radii \(v=\xi,-\xi\), are positive units.
Their axis values are one.
Use reversed physical time in the negative radial chart, so that both
attracting projected equations are
\(\dot\pi_\pm=-\tau P(\pm\pi_\pm)\).
The sections are \(\pi_+=p_w\), \(\pi_-=r_w\), while
\(\psi_-\) records the raw abscissa \(-v\).
Their exact projected axis equations give the transverse factors
\[
 \begin{aligned}
 A_+(w)&:=h_+'(w)-f_\tau(\iota_+(w))\psi_+'(w)
       =\frac{\tau P(p_w)}{D(\iota_+(w))\Delta_+(w)},\\
 A_-(w)&:=h_-'(w)-f_\tau(\iota_-(w))\psi_-'(w)
       =\frac{\tau P(-r_w)}{D(\iota_-(w))\Delta_-(w)}.
 \end{aligned}
\]
The reflection on the negative side reverses the abscissa and the
attracting time orientation, leaving both factors positive.
Since the weak buffers have fixed bounded height, \(c\tau\le A_\pm\le C\).

Fix any already admitted through orbit of the lower receiver, with weak
heights \(z>0\) and \(w>0\). Its core segment, the axis and the two
weak-section arcs from zero to these heights bound a strip. At either
cut the positive-time derivative of
\(\xi-\psi_\pm(h_\pm^{-1}(\eta))\) is \(-DA_\pm/h_\pm'<0\).
The strip lies to the left of the positive arc and to the right of the
negative arc, making these respectively incoming and outgoing.
Scalar uniqueness prevents a trajectory from crossing the
axis or this admitted orbit, and \(-\dot\xi\ge\tau p_c\) gives a
finite passage. Thus every point of \(0\le z_0\le z\) reaches the
negative arc inside this strip. The same actual intrinsic map
\(M_\tau\) is consequently defined on the entire chord, with
\(M_\tau(0)=0\) and \(0\le M_\tau(z_0)\le w\).
Uniqueness identifies these chord maps with each other and with the
local map of the weak-core flow.

Let \(V\) be the raw initial-height variation between the two actual
endpoint abscissae. Its equation \(V_\xi=(\partial_\eta f_\tau)V\)
on a slab of fixed length gives \(e^{-K/\tau}\le V\le e^{K/\tau}\).
Differentiating the actual endpoint equation yields
\[
 M_\tau'(z_0)=V\,
 \frac{A_+(z_0)}{A_-(M_\tau(z_0))}>0.
\]
Absorbing the fixed powers of \(\tau\) in the exponential constant,
\[
 e^{-K/\tau}\le M_\tau'(z_0)\le e^{K/\tau},\qquad
 \frac{M_\tau(z)}z=\int_0^1M_\tau'(tz)\,dt.
\]
The matching of these same intrinsic heights is
\(w_-=M_\tau(w_+)\), \(w_\pm=u_\pm e^{-J_\pm/\tau}\).
The positive \(C^4\) units \(u_\pm,u_\pm^{-1}\) therefore give
\[
 J_+-J_-
 =\tau\log\frac{M_\tau(w_+)}{w_+}
   +\tau\log\frac{u_+}{u_-},\qquad
 |J_+-J_-|\le K+C\tau\le K_0.
\]
Thus the actions differ by a bounded amount on every admitted through
orbit in the original independent section domains. The bound is obtained
before choosing an action cutoff and uses no return equation.

On the fixed weak core \(D\ge\tau p_c+q_c\eta\), the first four
normalized variational equations have bounds \(e^{K_j/\tau}\).
Factor the map just constructed at zero height as
\(M_\tau(z)=m_\tau zQ_\tau(z)\),
\(Q_\tau(0)=1,m_\tau^{\pm1}\le e^{K/\tau}\).
In particular, the exact quotient is
\[
 Q_\tau(z)=m_\tau^{-1}\int_0^1 M_\tau'(vz)\,dv.
\]
Four core variational equations bound its three ordinary derivatives.
Enlarge the exponential
constant to include this quotient, \(m_\tau^{-1}\), and the actual
input and output chord reserves; then choose \(J_*\) larger than
twice that constant with a fixed buffer. Thus
\(w_+=u_+(p)e^{-J_+/\tau}\) lies on the actual core chord and its
image lies in the negative weak image. Since
\(u_-(r)e^{-J_-/\tau}\downarrow0\) as \(r\uparrow c\), it has a
unique actual inverse. Taking logarithms on the product of separate
section domains gives
\begin{equation}\label{rc:strong-matching}
 \begin{gathered}
 J_+-J_-=C_\theta+\tau R(p,r),\qquad C_\theta=\tau\log m_\tau,\\
 R=\log u_+-\log u_-+\log Q_\tau(w_+),\\
 |C_\theta|\le K_0,\qquad\|R\|_{C^3_{p,r}}\le C(1+|\log\tau|).
 \end{gathered}
\end{equation}
Split the positive action exponential to absorb the core bounds and the
losses in the five root charts. Differentiation in the negative section
variable acts directly on \(u_-\); this avoids differentiating an
exponentially large inverse. Since the heights are intrinsic normal
coordinates, there is no moving initial \(\log P\) term.
On the actual matched graph, with \(P_+=P(p),P_-=P(-r)\),
\begin{equation}\label{rc:strong-variation}
 r_p=\frac{pB_+P_-}{rB_-P_+}(1+\rho),\qquad
 1+\rho=\frac{1-\tau P_+R_p/(pB_+)}{1+\tau P_-R_r/(rB_-)},
 \qquad\|\rho\|_{C^2_{p,r}}\le\epsilon_\tau:=C\tau(1+|\log\tau|).
\end{equation}
The polynomial multipliers and their derivatives are bounded; all
remaining denominators are units, including when \(P_+,P_-\) vanish
independently.

\paragraph{Curvature in the physical input.}
Let \(\Gamma_\pm(b)=a_\pm+c_\pm\phi_\pm(b)\), \(c_\pm>0\), be
these actual outward embeddings into the incoming raw height \(h\)
and divided output, with \(b=-p^2/2,S(b)=-r^2/2\).
The fixed normal-height equation
\(d(\eta-\xi^2/2)/d\log\eta=\tau\Phi\) has four bounded height
rows and clock \(O(1+|\log\tau|)\), independently of the axis gap.
For the exterior continuation retain \(\lambda_G=12\tau G/(1+2\tau G)\).
At each frozen bounded base \(q>0\),
\begin{equation}\label{deg:compact-section-integral}
 \int_1^{\delta/(\tau q)}\rho^{-(2-\lambda_G)}
 [\tau q^{-2}\rho^{-2}+\tau Uq^{-1}\rho^{-1}+\tau q\rho]
 (1+|\log\rho|^2)\frac{d\rho}\rho\le C\tau(1+q).
\end{equation}
The actual relative hierarchy yields
\(\phi_\pm=\mathrm{id}+O_{C^3}(\epsilon_\tau)\), including inverse
maps on their actual images. At an ordinary observer with exact small
gain \(a_\tau\asymp\tau^{2-O(\tau)}\), the normalized second and
third rows are \(a_\tau\chi''/\chi'_*,a_\tau^2\chi'''/\chi'_*\).
Only its ordinary \(C^3\) bounds are used. Outgoing and incoming
cocycles at every cut cancel before normalization; the output closing
retains \(f'(md)d'\) and its ordinary third chain rule.
Thus \(\widehat\Pi=\phi_-\circ S\circ\phi_+^{-1}\) is the actual
completed map up to its two retained positive affine gains.
The exact chain rule gives
\begin{equation}\label{deg:compact-finite-curvature}
 \phi_+'(b)\frac{P(p)}{p+r}\frac{\widehat\Pi''}{\widehat\Pi'}
 =3-\frac U{pr}+E(p,r),\qquad\|E\|_{C^1_{p,r}}\le\epsilon_\tau.
\end{equation}
To display the error, put \(W=p+r\),
\(K_{\rm mat}=B_+(1+\rho)/B_-\), and
\(L_\pm=\phi_\pm''/\phi_\pm'\), each in its own quadratic label.
Then
\[
 \begin{aligned}
 E={}&-\frac{(K_{\rm mat}-1)P_-'}{rW}
       -\frac{P_+}{pW}\partial_p\log K_{\rm mat}
       -\frac{K_{\rm mat}P_-}{rW}\partial_r\log K_{\rm mat}\\
    &+\frac{K_{\rm mat}P_-}{W}L_-
       -\frac{P_+}{W}L_+ .
 \end{aligned}
\]
Here \(P_-'\) uses the radial variable \(r\).
Since \(K_{\rm mat}-1=O_{C^2}(\epsilon_\tau)\) and
\(L_\pm=O_{C^1}(\epsilon_\tau)\), every term has the asserted
\(C^1\) bound on the product of the two normal-section domains:
\(P_\pm\) occur as multipliers and the remaining denominators are units.
Near a zero, \(U\ge U_*>0\), and restriction to the matched graph gives
\[
 \begin{aligned}
 \frac d{dp}\left(3-\frac U{pr}\right)
 &=\frac U{pr}\left(\frac1p+\frac{r_p}{r}\right)
 \ge c(1+r_p),\\
 \left|\frac d{dp}E(p,r(p))\right|
 &\le\epsilon_\tau(1+r_p).
 \end{aligned}
\]
Since \(w_p=(\phi_+(-p^2/2))_p<0\), every curvature zero is
simple and downward in the physical input.

On \(J_+\le J_{\rm rf}>J_*+1\),
both actions are bounded. The five singular integrals give a fixed
positive polynomial margin on the whole traversed segment and its
buffered exits. Corollary~\ref{neg:cor-rootfree}, with cubic coefficient
one, gives actual absolute \(C^3\) comparison there. The limiting map has
\(S_0'=P_-/P_+\),
\(S_0''/S_0'=(p+r)(3-U/(pr))/P_+\), and
\((pr)_b=-r/p-pP_-/(rP_+)<0\).
At \(U=3pr\) its \(S_0'''\) has a uniform negative margin; away
from that set the curvature has a nonzero margin. This transfers the
same crossing orientation to the bounded-action region. That comparison
strictly overlaps the high-action estimate for the same map.

On ghost fibres choose \(M\) above all roots and bottlenecks, large
for Lemma~\ref{deg:lem-weighted-cubic-tail}. Its actual growing
receiver \(u\ge2M,\tau u\le d_0\) has matched \(r_0\le r<c\)
and \eqref{deg:compact-curvature}, with error
\(C\tau(1+u)(1+|\log(\tau(1+u))|)^4\). Choose
\(M>2U_0/r_0\), then \(d_0\) with
\(Cd_0(1+|\log d_0|)^4<1/4\), then \(\tau\). Its curvature is
positive throughout this growing range.
To establish actual overlap, take a buffered band
\([2M-\beta,3M+\beta]\), \(0<\beta<M/2\), in the fixed normal
boxes. The lip projection is \(k_\tau(u)=u+O_{C^3}(\tau)\),
its linearized height is \(v_{\rm lip}=\tau\ell_0P(u)V_\tau(u)\),
with \(V_\tau\) a \(C^3\) unit. The same-orbit invariant gives
\begin{equation}\label{deg:compact-resection}
 \int_{k_\tau(u)}^{p_\tau(u)}\frac{qB(q)}{P(q)}\,dq
 =\tau\log\frac{u_+(p_\tau(u))}{v_{\rm lip}(u)},\qquad
 \|p_\tau-\mathrm{id}\|_{C^3}\le C\tau(1+|\log\tau|),\quad p_\tau'>0.
\end{equation}
The divergent bottleneck prefix cancels in this local integral. Its
\(p\)-derivative after subtracting \(\tau u_+'/u_+\) is a unit,
so the implicit equation gives all three rows. The closed band
\([2M+\beta,3M-\beta]\) has \(p_\tau<3M,u>2M\); the inverse
image of \(p=3M\) also has \(u>2M\), leaving no gap. These actual
orbits retain the original windows: their bounded normal labels have
raw gains \(\tau^{2-\lambda_G}\) times units, their reference images
lie inside the original passive organizing-arc buffers, and the raw
corridor continues them across the negative root at positive height.
Scalar order keeps the band in the comparison component. Resection
identifies orbits; each curvature estimate uses its own actual adapters.

\paragraph{The remaining raw ghost tail.}
On a ghost fibre, through trajectories can continue beyond the cutoff
\(\tau u=d_0\) of the growing estimate. The first-lip abscissa differs from the minimum
speed \(\sigma_\theta\); small central height may coexist with
\(a=\tau u>d_0\). Put \(a_b=d_0/2\), and let
\(x_*=\delta-\gamma\) be the actual positive raw port. For every
admitted \(a_b\le a\le x_*\), at fixed coefficients set
\[
 \begin{gathered}
 M_\tau=x^3+3\tau Gx^2-\tau^2Ux+\tau^3,\qquad B_0=1+2\tau G+x,\\
 Q_0=1-A(\tau G+x)+C(\tau G+x)^2,\quad
 \zeta_0=C(\tau G+x),\quad D_0=M_\tau+Q_0z.
 \end{gathered}
\]
The true graph is \(z_x=z(xB_0+\zeta_0z)/D_0\),
\(z(a)=\ell_0M_\tau(a)\). On \([a,x_*]\), \(M_\tau\ge c(d_0)>0\)
and \(Q_0\) is a positive unit. The graph satisfies
\(z_x\le Cx+Cz\), so its positive initial height and Gronwall
give \(z\le C\delta^2\), below the fixed upper chart boundary.
Its numerator is positive after the coefficient shrink, and hence
\(z\ge\ell_0M_\tau(a)>0\). This constructs the regular outward
passage, including its auxiliary zero-flight endpoint. Append the
observer only when the image lies in its original buffer; the auxiliary
endpoint needs no observer value when its image is outside that buffer.
Use \eqref{deg:scalar-polynomial-lip-slope}, with \(\sigma=\tau\),
\(K_{\rm lip}=K\), \(I_0=I\), and \(x_T=x_*\):
\begin{equation}\label{deg:compact-raw-tail-input}
 \Gamma_{+,a}'=-a\chi'(z_{\rm port})\frac{z_{\rm port}}{D_0(x_*,z_{\rm port})}
                   B_0(a)K_{\rm lip}(a)e^{I_0(a)}.
\end{equation}
The original lip transversality makes \(K_{\rm lip}\) a positive
unit. All admitted factors have bounds and reciprocals depending on
\(d_0\), so \(c(d_0)\le-\Gamma_{+,a}'\le C(d_0)\), or
\(c(d_0)\tau^2\le\Gamma_{+,b_+}'\le C(d_0)\tau^2\)
for \(b_+=-a^2/(2\tau^2)\).

We match this outward map to the negative branch by their actual central
heights. Differentiate the positive inward passage directly. For its height
\(\eta_0^+(u)\), \(W=\eta/(\tau P)\),
\(H_+=P'/(\xi B)\), \(\varpi=C\tau^3(G+\xi)P/(\xi B)\),
\begin{equation}\label{deg:compact-raw-tail-inward}
 \partial_u\log\eta_0^+=-\frac{uB_+}{\tau P_+}\Lambda_+,
 \quad\Lambda_+=\frac{K_{\rm lip}(a)}{1+Q_+\ell_0}
       \exp\int_u^0\frac{\xi B}{\tau P}
                         \frac{W(\varpi-Q)}{(1+QW)^2}\,d\xi.
\end{equation}
In the action clock \(dq=-\xi B\,d\xi/(\tau P)\),
\[
 W_q=W\left[-\frac{1+\varpi W}{1+QW}+\tau H_+\right],\qquad W(0)=\ell_0.
\]
On the original raw chart \(Q-\varpi>0\); the braces increase with
\(W\), and at \(\ell_0\) equal
\(-K_{\rm lip}/(1+Q\ell_0)\le-c_1\). Thus
\(W\le\ell_0e^{-c_1q}\) down to a fixed small scaled cut
\(\delta_c\), on the entire original raw range. The integral in
\eqref{deg:compact-raw-tail-inward} becomes the bounded integral
\(\int W(Q-\varpi)/(1+QW)^2\,dq\).
For the negative branch put \(P_-(v)=P(-v)=(c-v)L_-(v)\),
\(B_-(v)=B(-v)\), \(Q_-(v)=Q(-v)\), with
the exact unit \(L_-\) used in
Lemma~\ref{deg:lem-weighted-cubic-tail}. Its inverse inward graph is
\[
 \eta_v=f_-(v,\eta)
 =\frac{\eta\{vB_-+\zeta_-\eta\}}{\tau P_-+Q_-\eta},\qquad
 \zeta_-=-C\tau^2(G-v),\qquad
 \eta(r)=\tau\ell_0P_-(r).
\]
On \(\delta_c\le v\le r<c\), set
\(H_-=P_-'/ (vB_-)\) and
\(\varpi_-=\tau P_-\zeta_-/(vB_-)\).
The compact coefficient block gives bounded \(H_-\),
\(\varpi_-=O(\tau^3)\), and positive-unit \(Q_-\), even at
the simple root. The reflected ratio \(W_-\) has the same action-clock
equation and strict decay. Its complete moving-initial factor is
\[
 K_-=1-\tau H_-(1+\ell_0Q_-)+\ell_0\varpi_->c_K>0.
\]

Choose \(2\delta_c<r_0\). A fixed positive band below every
positive raw-tail label and the fixed negative band
\([\delta_c,r_0]\) supply action margins \(A_+,A_->0\).
Thus both inward graphs reach their central slab with heights
\(\eta_i\le C\tau e^{-c_1A_{\rm in}/\tau}\),
\(A_{\rm in}=\min(A_+,A_-)\). On that slab \(P_\pm\ge p_c>0\).
In the decreasing radial coordinate each graph satisfies
\(\eta_t\le C\tau\eta^2\); first exit keeps \(0<\eta\le2\eta_i\)
and gives an actual finite continuation to \(\xi=0\).
For \(f_+\) use the positive graph, with
\(\zeta_+=C\tau^2(G+v)\). The exact rational difference is
\[
 f_{\pm,\eta}-f_\pm/\eta
 =\frac{\eta\{\zeta_\pm\tau P_\pm-n_\pm Q_\pm\}}
             {(\tau P_\pm+Q_\pm\eta)^2},\qquad n_\pm=vB_\pm.
\]
Its central integral is bounded by \(C\eta_i/\tau^2\).
Together with the bounded noncentral ratio integral this proves the
positive \(\Lambda_+\) bound and, independently, the reflected identity
\begin{equation}\label{deg:compact-reflected-inward}
 \begin{gathered}
 -\partial_r\log\eta_0^-(r)=\frac{rB_-}{\tau P_-}\Lambda_-,\\
 \Lambda_-=
 \frac{K_-(r)}{1+\ell_0Q_-(r)}
       \exp\int_r^0(f_{-,\eta}-f_-/\eta)\,dv,\qquad
 0<c_\Lambda\le\Lambda_\pm,\Lambda_\pm^{-1}\le C_\Lambda.
 \end{gathered}
\end{equation}
No root factor has been inverted in these error bounds.

The ratio memory and the central linear cocycle give
\(\eta_0^-(r)=\tau\ell_0P_-(r)e^{-A_-(r)/\tau}E_-(r)\),
where \(E_-,E_-^{-1}\le C\) and
\(A_-(r)=\int_0^r vB_-/P_-\,dv\).
The simple root makes \(A_-(r)\to\infty\) as \(r\uparrow c\).
Hence the actual negative centre map is strictly decreasing with range
\begin{equation}\label{deg:compact-negative-centre-image}
              \eta_0^-([r_0,c))=(0,\eta_0^-(r_0)].
\end{equation}
Start with an admitted shot \(u_{\rm ref}\) in the fixed resection
band, whose negative match has \(r_{\rm ref}\ge r_0\).
For small \(\tau\), \(\tau u_{\rm ref}<a_b\).
The weighted inward construction, before either port observer is
attached, is strictly decreasing up to \(u_b=a_b/\tau\), since
\(\tau u\le a_b<d_0\). Thus
\(\eta_0^+(u_b)<\eta_0^-(r_{\rm ref})\), and the full negative
image gives a unique \(r_b\ge r_{\rm ref}\).
This auxiliary centre match does not require the formal \(a_b\) port
image to lie in an observer window.
Positive centre-height monotonicity and this full negative range then
give every original admitted \(a\ge a_b\) its unique negative match.
Scalar uniqueness at the same \(\xi=0\) identifies the joined graph
with the original through orbit; its negative output and closing
images therefore retain their original buffers. Matching them
gives \(S'=B_+P_-\Lambda_+/(B_-P_+\Lambda_-)\).
In the root collar the negative output embedding keeps its own endpoint
and exterior gains:
\[
 \begin{gathered}
 \Gamma_{-,b_-}'=d_c\theta_cX^{\beta_-}K_c(X,X^{\beta_-}),\quad X=c-r,\\
 d_c=\tau^{2-\lambda_G}U_c,\quad\theta_c=\tau^{\beta_-},\quad
 \beta_-=\tau P'(-c)/(cB(-c))>0.
 \end{gathered}
\]
Here \(U_c,K_c\) are bounded positive units. On the fixed complement
the ordinary negative exterior calculation gives
\(\Gamma_{-,b_-}'=\tau^{2-\lambda_G}\) times a bounded positive
unit. These estimates overlap on a fixed collar band and give
\(0<\Gamma_{-,b_-}'\le C\tau^2\) on the entire \(r_0\le r<c\),
since \(\tau^{-\lambda_G}\) is a unit. The closing contributes its actual
bounded positive first derivative. Now \(P_+(u)=\tau^{-3}M_\tau(a)
\ge c(d_0)\tau^{-3}\), while \(P_-\le C\). Therefore
\begin{equation}\label{deg:compact-raw-tail-derivative}
 0<\Pi'=\frac{\Gamma_{-,b_-}'}{\Gamma_{+,b_+}'}
           \frac{B_+P_-}{B_-P_+}\frac{\Lambda_+}{\Lambda_-}
                       \le C(d_0)\tau^3.
\end{equation}
The two absolute gains are retained in this ratio. The estimate covers
\(a_b\le a\le d_0\) and every original admitted \(a>d_0\), subject
to all the original window, speed and height restrictions. It therefore
continues the derivative bound beyond the growing estimate's cutoff.

\paragraph{The upper continuation.}
For the upper complement continue the fixed comparison outward to
\(\xi=\pm4v\). The fixed-normal fast invariants give positive-unit
heights there; scalar order preserves these lower bounds for all larger
inputs, and the field bounds them above through a fixed
\(\eta_{\max}>0\). With \(f=\eta H/D\), specialization of
\eqref{deg:scalar-raw-cocycle} gives
\begin{equation}\label{deg:compact-core-variation}
 f_\eta=\frac d{d\xi}\log\frac\eta D+
       \frac{\tau P'+(Q'+C\tau^2(G+\xi))\eta}{D}.
\end{equation}
Endpoint quotients are units, and \(D\ge\tau p_c+q_c\eta\) bounds
the integral even at exponentially small central height. The two outward
first variations have the same exact factor \(\tau^{2-\lambda_G}\);
their ratio and closing give \(0<\Pi'\le C_{\rm core}\).
The comparison orbit and a fixed action neighborhood lie in both the
lower and core regions. The comparison height \(\eta_{\rm comp}\)
is exponentially small.
The signed threshold is
\(H_H=\tau^3K_H(1+\tau U^2)\), \(\eta_H=H_H/\tau^2\).
After all fixed cuts and caps, choose \(\tau_0\) so
\(\eta_{\rm comp}<\eta_H<\eta_{\max}/2\) and
\(\tau^2\eta_{\max}<h_c/2\). Core and signed high then strictly
overlap on the original height axis.

\paragraph{Covering and counting the whole interval.}
The through lemma and increasing \(H,\sigma_\theta\) make
\(I_\theta\) connected. The comparison is admitted: its raw outward
images retain the original buffers, its central speed is
\(O(\tau^3+\tau^2\eta_{\rm comp})<d_{\rm out}\), and \(H<h_c\).
At its transverse positive lip every lower input has \(W<\ell_0\).
On a ghost fibre the fixed lower height of the original positive window
and the original lip reserve give \(W>\ell_0\) at the raw port.
Every input in \(I_\theta\) already has an actual through trajectory;
scalar order below the comparison gives \(W<\ell_0\) at its positive
lip. Thus a hit exists. Every hit has negative transverse derivative,
so it is unique. At the hit the radial derivative of \(W\) is positive,
whereas its input-height derivative is positive by scalar variation.
The implicit hit equation therefore gives \(a_x<0\) in the original
affine input, before any return-zero argument.
On real/fold fibres a through orbit crosses the incident inner root at
positive height, with \(W\to\infty\) on its inner side, giving the
same unique hit. An outer-component lip hit would force
\(\eta\le\tau\ell_0P\to0\) at the next root and a stopped passage.
Thus all lower through inputs have one of the actual descriptions above.
The bounded and growing ones have no gap by
\eqref{deg:compact-resection}; the growing one reaches \(a=d_0\),
and \eqref{deg:compact-raw-tail-derivative} continues every admitted tail.
The whole original carrier has the exact allocation
\begin{equation}\label{deg:compact-physical-cover}
 \begin{split}
 I_{\rm tail}&=\{x\in I_\theta:x\le x_{\rm comp},\ \Delta<0,\ a(x)>d_0\},\\
 I_{\rm lower}&=\{x\in I_\theta:x\le x_{\rm comp}\}\setminus I_{\rm tail},\\
 I_{\rm core}&=\{x\in I_\theta:x>x_{\rm comp},\ H(x)\le H_H\},\\
 I_{\rm high}&=\{x\in I_\theta:H(x)>H_H\},\qquad
 I_\theta=I_{\rm tail}\cup I_{\rm lower}\cup I_{\rm core}\cup I_{\rm high}.
 \end{split}
\end{equation}
They are intervals or empty clips, since \(a(x)\) decreases. Equality
\(a=d_0\) and the comparison belong to the lower piece; \(H=H_H\)
belongs to the core. Closed action dividers are retained, and all estimate
packages have the strict overlaps above. Natural stopped endpoints and
the infinite-time fold branch are excluded.
Choose \(m\) last with
\(m\max(C_{\rm core},C_H,C(d_0)\tau_0^3)<1/2\).
Then \(F'>1/2\) on tail, core and high. On the lower receiver the
exact two-gain composition gives
\[
 F=a_+-a_\theta-ma_-+c_+w-mc_-\widehat\Pi(w),\qquad
 F''=-\frac{mc_-}{c_+^2}\widehat\Pi'',\quad
 F'''=-\frac{mc_-}{c_+^3}\widehat\Pi'''.
\]
With \(J=I_{\rm lower}\), Lemma~\ref{sharp:upper-curvature-count}
applies on the connected \(I_\theta\): the positive factors preserve
the lower downward curvature crossings, and \(F'>1/2\) on
\(I_\theta\setminus J\). It gives at most one curvature zero on \(J\),
at most two zeros of \(F'\), counted with multiplicity, and at most
three isolated zeros of the
same physical \(F\) on \(I_\theta\), counted with multiplicity.
All choices depend only on the fixed coefficient block and section
margins. The estimates allow the two gains and the two root gaps to
vary independently.
\end{proof}



\section{Uniform regional returns and the three-cycle upper bound}
\label{sec:common-cover-conclusion}

The regional estimates apply to the completed physical map. We fix
the coefficient subcover and its reserves, choose a common raw-height
cap, and only then restrict the independent coefficients and finite
multiplier. Each coefficient alternative counts the whole returning
interval; geometric capture completes the upper bound.
\subsection{From the action to the physical return}

The passage estimates below concern the completed map in the physical
input of Section~\ref{sec:counting-tools}. The ordinary section calculus
of Lemma~\ref{phys:lem-c3-adapters} and the curvature identity of
Lemma~\ref{neg:lem-completed-curvature} retain both nonlinear section
maps and their independent gains. We use four comparisons.

On strict passages, Proposition~\ref{prop:source-normalization}
gives ordinary phase \(C^3\) convergence. The common absolute tail
factor cancels in \eqref{sharp:upper-tail-ratio}; after the regular
closing this bounds the upper continuation. Near reversibility,
Theorem~\ref{neg:thm-relative} instead divides the error by the
coefficient distance before comparison. Its root-free margin and
entry/exit buffers are part of its hypotheses.

For coalescing roots, Proposition~\ref{deg:prop-compact-cubic}
estimates the two root distances independently before matching their
actual heights. Its bounded and growing descriptions overlap on the
same map; the remaining original ghost tail is controlled by
\eqref{deg:compact-raw-tail-derivative}. The quadratic counterpart is
Proposition~\ref{bp-tail:prop-actual-raw-tail}. In that case the signed
density during the negative-root dwell is essential. Proposition~\ref{rawpair:prop-original-tails}
identifies both tail estimates with the common physical receiver.

At degree loss, the two saddle scales remain independent. The actual
section comparison in Proposition~\ref{p2:prop-finite-v-fe-actual-budget},
the growing-range matching and the fixed-exit construction of
Proposition~\ref{fx:prop-inverse-phase} cover one returning interval.
Proposition~\ref{p2:prop-physical-degree-count} counts its critical
points. The exact dictionaries of
Appendix~\ref{app:proof-small-source-routing} place whole coefficient
fibres in these ranges. The coefficient cover and raw budgets below
fix their overlap reserves before the common height cap is chosen.

The following bounds apply to the same completed displacement on its
actual connected through interval. The appendices prove the transition
and section estimates used in each region, with the one-sided endpoint
estimates and continuations needed to count the full physical carrier.

\subsection{The signed high-height interval}

\begin{corollary}[Signed high-height interval]
\label{rew:cor-high-derivative-signed}
There exist \(K,z_*,\gamma_*,C_H>0\) such that, for all sufficiently
small raw coefficients and either sign of \(\alpha_0\), the complete
transition on the actual fixed sections satisfies \(0<\Pi'\le C_H\)
throughout the interval of through trajectories
\begin{equation}\label{ann:high-source-interval}
 K(|\alpha_0|+\alpha_1^2)<H:=z(\gamma)<z_*,
 \qquad |\gamma|<\gamma_*.
\end{equation}
After the common choice \(mC_H<1/2\), the complete displacement has
\(\Delta'>1/2\) and hence at most one isolated zero, counted with
multiplicity, on this entire interval and on each of its truncations.
The assertion includes \(\alpha_0=\alpha_1=0\). At a positive
boundary height \(H=K(|\alpha_0|+\alpha_1^2)<z_*\), with the same
coefficient bounds, the derivative bound still holds; this equality
face is assigned to the low-height region.
\end{corollary}
\begin{proof}
Fix the coefficients and a through trajectory in
\eqref{ann:high-source-interval}. Freeze \(r_0=\sqrt H\) and put
\[
 a=\frac{\alpha_0}{H},\qquad
 p=\frac{\alpha_1}{\sqrt H},\qquad s=r_0,
 \qquad v=\gamma+r_0X,\quad z=r_0^2Y.
\]
The reference orbit has \(Y(0)=1\), and
\[
 |a|<K^{-1},\qquad |p|<K^{-1/2}.
\]
Choose \(K\) so large that the closed block with the corresponding
nonstrict bounds lies inside the signed compact strict family of
Proposition~\ref{sharp:upper-high}. Then choose \(z_*\) and
\(\gamma_*\) small enough for its raw-tail estimates. Its complete
derivative bound applies with the frozen scale \(r_0\), including
both signs of \(a\) and \(a=p=0\). The phase-independent constant bounds the derivative of the same
physical map, with \(r_0\) frozen during differentiation.

Scalar first variation makes \(H\) increasing with that input, so
\eqref{ann:high-source-interval} and its truncations are intervals.
Choose \(mC_H<1/2\) last; then \(\Delta'>1/2\) throughout.
The closed block includes the positive equality face, assigned to
low height. For double-zero coefficients every positive height is high.
\end{proof}

\subsection{The compact positive-ratio whole-carrier bound}

\begin{proposition}\label{sharp:upper-annulus}
For each fixed \(0<p_0<P<\infty\), the complete captured carrier
with \(a=-1\) and \(p_0\le p\le P\) has at most two
isolated displacement roots, counted with multiplicity.
\end{proposition}
\begin{proof}
Use the one connected actual receiver and its open preimages from
Lemma~\ref{ann:lem-compact-positive-domains}. Keep the physical
affine input \(x\) throughout. On \(I_E\cup I_M\),
\(\Pi_{xx}<0\), so for the physical multiplier \(m>0\),
\(\Delta''=-m\Pi''>0\). On \(I_H\), the complete derivative
bound of Lemma~\ref{ann:p0-paired-high-derivative} gives
\(\Delta'>1/2\) after the single final choice
\(m<\min\{m_{\rm ad},1/(2C_H)\}\).
Every critical point has \(\Delta''>0\).
Lemma~\ref{sharp:upper-critical-count} gives the whole-carrier
multiplicity bound.
\end{proof}
\subsection{The gamma-dominant and far-root carriers}

\begin{proposition}\label{sharp:upper-gamma}
For a fixed \(\Pi_{\max}<\infty\), the bounded-\(\pi\)
negative-gamma carrier has at most two displacement zeros, and the
bounded-\(\pi\) positive-gamma carrier has at most three. The positive
far-root carrier of \eqref{deg:positive-far-range} has at most two.
All counts hold on the entire captured physical interval.
\end{proposition}
\begin{proof}
Propositions~\ref{deg:prop-negative-gamma} and
\ref{deg:prop-positive-far} give the two-zero counts on their
original connected carriers, with multiplicity, using the covers
\eqref{deg:negative-gamma-cover} and \eqref{deg:far-physical-cover}.
For positive bounded \(\pi\),
Corollary~\ref{deg:cor-positive-quadratic-whole-carrier} gives at most
two zeros of the same physical \(F'\), including the ghost raw tail.
Its paired/high overlap, in \(\eta=gY_0\), is
\begin{equation}\label{sharp:upper-gamma-overlap}
 gK_H(1+\Pi_{\max}^2g)<\eta_{\max}/2,\qquad
 (gk)^2\eta_{\max}<H_{\rm cap}/2.
\end{equation}
The remaining root-free fibres are contractions by
Lemma~\ref{deg:lem-raw-jacobian}. Lemma~\ref{sharp:upper-critical-count}
then gives the three-zero count on the whole original carrier.
\end{proof}

\subsection{The compact cubic block: a common curvature orientation}
\label{sharp:upper-cubic-section}

\begin{proposition}[The complete compact cubic carrier]
\label{sharp:upper-cubic}
Fix finite \(U_0,G_0\). On the coefficient block
\begin{equation}\label{sharp:upper-cubic-block}
 a=-1,\quad p=tU,\quad\gamma=t^2G,\quad s=t^3,
             \qquad0\le U\le U_0,\quad |G|\le G_0,
\end{equation}
the same displacement has at most two critical points and at most three
isolated zeros, counted with multiplicity, on its whole physical carrier
under the common positive unscaled raw cap. All captured target cycles
lie in this carrier after the common coefficient and marked
Hausdorff-collar shrink.
\end{proposition}

\begin{proof}
Put \(\tau=t^2\). Proposition~\ref{deg:prop-compact-cubic} supplies the
actual clipped carrier and its allocation \eqref{deg:compact-physical-cover}.
Its bounded-action and independent root-collar estimates give the same
curvature orientation. The growing range has positive curvature by
\eqref{deg:compact-curvature}, and is identified with the bounded
range by the actual resection
\begin{equation}\label{sharp:upper-cubic-overlap}
 p_\tau(u)=u+O_{C^3}(\tau(1+|\log\tau|)),\quad p_\tau'>0,
                \qquad2M+\beta\le u\le3M-\beta.
\end{equation}
These are estimates for the completed physical map, so
\begin{equation}\label{sharp:upper-cubic-common-crossing}
 \Pi''(x)=0\quad\Longrightarrow\quad\Pi'''(x)<0
\end{equation}
on the connected lower receiver. Real and fold fibres retain their
one-sided through endpoints.

The original ghost tail has the derivative bound
\eqref{deg:compact-raw-tail-derivative}; the core has
\eqref{deg:compact-core-variation}. The construction places
their comparison and high intervals in strict overlap:
\begin{equation}\label{sharp:upper-cubic-height-overlap}
 \eta_{\rm comp}<\tau K_H(1+\tau U^2)<\eta_{\max}/2,
                      \qquad\tau^2\eta_{\max}<h_c/2.
\end{equation}
One final multiplier bound makes \(F'>1/2\) on this complement.
The cover in Proposition~\ref{deg:prop-compact-cubic} includes every
input, retaining the original clips and equality dividers. Its whole-carrier
count, equivalently Lemma~\ref{sharp:upper-curvature-count}, gives
the stated multiplicity bounds. Capture is applied after these analytic
choices, as in Lemma~\ref{lem:loc-capture}.
\end{proof}

\subsection{Degree loss: a common orientation of critical points}
\label{sharp:upper-degree-section}

The compact graph \eqref{deg:degree-graph} and reciprocal graph
\eqref{deg:degree-infinity-graph} share the physical carrier of
Proposition~\ref{p2:prop-unified-carrier-cover}. The reciprocal face
\(j=0\) concerns only the one-sided closure of the estimates.

\begin{proposition}[The complete degree-loss carrier]
\label{sharp:upper-degree}
Fix the compact positive-unit ranges in
Proposition~\ref{deg:prop-degree-drop}, or the ranges of \(d,A_d\) and
\(0<j\le j_0\) in Proposition~\ref{deg:prop-degree-infinity}.
There are independent active coefficient bounds, one fixed raw-height
cap and a final finite-closing multiplier bound for which the complete
displacement has at most one critical point on its connected interval
of actual through inputs in the original nested windows.
Every such point is a simple upward crossing of the same physical
\(F'\). The displacement therefore has at most two isolated zeros,
counted with multiplicity, and no identity interval.
All estimates hold at every retained input, including nonclosed
through orbits and positive divider equalities. They are uniform for
physical \(j>0\) as \(j\downarrow0\).
\end{proposition}

\begin{proof}
Write the return in the affine input of Lemma~\ref{lem:closing} as
\begin{equation}\label{sharp:upper-degree-displacement}
 F_\theta(x)=x-f_\theta\bigl(m\mathfrak T_\theta(a_\theta+x)\bigr)
             =x-m\Pi_{\theta,m}(x).
\end{equation}
Here $\mathfrak T_\theta=d_\theta\circ T_\theta^\infty$ and
$f_\theta=C_\theta-C_\theta(0)$.
Proposition~\ref{p2:prop-unified-carrier-cover} gives its actual
clipped carrier. Proposition~\ref{p2:prop-physical-degree-count}
retains both gains, output translation and ordinary $C^3$ closing,
and proves
\begin{equation}\label{sharp:upper-degree-critical-orientation}
 F_\theta'(x)=0\ \Longrightarrow\ F_\theta''(x)>0,
 \qquad \#\{x:F_\theta'(x)=0\}\le1.
\end{equation}
Its whole-carrier count gives the remaining conclusions.
\end{proof}

\subsection{The determinant passage}

\begin{proposition}\label{sharp:upper-determinant}
Fix \(P_D<\infty\), the passive compact coefficient block and original
window reserves. For independent small \(r_D>0\), \(|\gamma|\), with
\(\alpha_0=r_D^2\), \(|\alpha_1/r_D|\le P_D\), the displacement on
\eqref{deg:det-actual-receiver} has at most one critical point and at most
two isolated zeros, counted with total multiplicity, after a common
coefficient and final multiplier shrink. The raw cap is fixed, positive
and unscaled. The assertion holds for every original window, speed or
height clip and includes saddle ratio one.
\end{proposition}

\begin{proof}
Use the single affine \(F(x)=x-m\Pi(x)\) of Lemma~\ref{lem:closing}.
Lemma~\ref{deg:lem-determinant-preparation} supplies the full stopped
composition, both exact absolute gains and ordinary \(C^3\) divided
closing. On the strict and high sets in \eqref{deg:det-height-cover},
choose \(m\) last so that their complete derivative bounds give
\(F'>1/2\). Thus every critical point lies in the actual stopped range.

There \eqref{sharp:upper-det-absolute} retains every factor of the same
physical derivative:
\[
 \Pi'(x_b+\sigma_+\delta)=\mathcal A_\theta\delta^q e^{R_\theta(\delta)},
 \quad q=\wp-1,\quad
 |R_\theta|+|\delta R_\theta'|\le C\delta^\alpha(1+|\log\delta|).
\]
The quotient of the two gains is a unit, so \(\mathcal A_\theta\)
is bounded above and below; neither absolute gain needs a lower bound.
Let \(C_0\) bound \(|\log\mathcal A_\theta+R_\theta|\).
At a critical point \(m\Pi'=1\), set \(t=-\log\delta>0\). Then
\[
 -qt\ge L:=\log(1/m)-C_0.
\]
For \(L>0\), this requires \(q<0\) and \(|q|\ge L/t\);
in particular \(\wp=1\) produces no critical point. Moreover
\[
 \frac{|\delta R_\theta'|}{|q|}
 \le\frac C L t e^{-\alpha t}(1+t)\le C_1/L.
\]
Choose \(L>2C_1\). The Euler derivative of the logarithm of the full
physical gain satisfies
\(\delta\partial_\delta\log\Pi'<q/2<0\) at every critical point.
Since \(dx/d\delta=\sigma_+>0\),
\begin{equation}\label{sharp:upper-det-critical}
 F'(x)=0\quad\Longrightarrow\quad F''(x)=-m\Pi''(x)>0.
\end{equation}
The actual stopped/strict and low/high cover
\eqref{deg:det-height-cover} retains every clipped and intervening
input. Lemma~\ref{sharp:upper-critical-count} now gives at most
one critical point and two displacement zeros.
\end{proof}

\subsection{The signed trace passage}

\begin{proposition}\label{sharp:upper-trace}
Fix \(0<p_0<p_1<\infty\), the passive compact block and original
window reserves. For \(\alpha_0=r^2a\), \(\alpha_1=rp\),
\(|a|\le a_*\), \(p_0\le|p|\le p_1\), independent small
\(r>0\), \(|\gamma|\), and a final common multiplier shrink,
the displacement on the original receiver
\eqref{deg:trace-actual-receiver} has at most two isolated zeros for
\(p>0\), and at most one for \(p<0\), counted with total multiplicity.
Its positive unscaled raw cap and all window, speed and height clips
are retained. Both opening signs, zero opening and closed trace endpoints
are included.
\end{proposition}

\begin{proof}
Proposition~\ref{deg:prop-stopped-trace} gives \(F'>1/2\) for
\(p<0\) and \(F'=0\Rightarrow F''>0\) for \(p>0\) on the whole
original clipped receiver. Lemma~\ref{sharp:upper-critical-count}
gives the stated multiplicity bounds.
\end{proof}

\begin{corollary}\label{sharp:upper-local-coefficient-ranges}
The nonnegative-opening sector \(\alpha_0\ge0\) has at most
two captured displacement zeros. There is a fixed \(P_*<\infty\)
such that the negative-opening sector
\(\alpha_0=-s^2<0\), \(p_{\rm src}=\alpha_1/s\ge P_*\)
also has at most two. Both assertions concern the whole physical
interval and count multiplicities.
\end{corollary}

\begin{proof}
First choose the signed trace collar of
Proposition~\ref{sharp:upper-trace}, for example with
\(1/2\le|p|\le2\). Choose \(P_D>2/\sqrt{a_*}\).
If \(\alpha_0>0\) and
\(|\alpha_1|\le P_D\sqrt{\alpha_0}\), use
Proposition~\ref{sharp:upper-determinant} with
\(r_D=\sqrt{\alpha_0}\). On the complement put
\(r=|\alpha_1|\). Then
\[
 p=\sgn\alpha_1,\qquad
 a=\alpha_0/\alpha_1^2<P_D^{-2}<a_*.
\]
This is the exact trace scaling of
Lemma~\ref{deg:lem-positive-atlas}; the full physical domains are
retained in Lemma~\ref{deg:lem-nonnegative-domain}. The same choice treats
\(\alpha_0=0\), \(\alpha_1\ne0\), with \(a=0\).
Each alternative bounds the displacement on its whole coefficient fibre.

The double-zero coefficient face is a contraction by
Corollary~\ref{rew:cor-high-derivative-signed} at every positive
height, without introducing a determinant scale.

Finally suppose \(\alpha_0=-s^2<0\) and
\(p_{\rm src}=\alpha_1/s>0\). With \(r=\alpha_1\), the
trace coefficients are exactly
\[
 p=1,\qquad a=-p_{\rm src}^{-2}.
\]
Choose \(P_*>2/\sqrt{a_*}\) before the subsequent compact
positive-source interval is fixed. Then
\(p_{\rm src}\ge P_*\) belongs to
Proposition~\ref{sharp:upper-trace}. Take the maximum of the
whole-fibre bounds across these coefficient alternatives.
\end{proof}


\subsection{Coefficient-face decomposition and passage domains}
\label{sec:coefficient-limits}
\label{sec:parameter-degenerations}

We use the variables of Section~\ref{sec:coefficient-regimes} to cover the
coefficient space independently of the phase input. The exact coordinate
dictionaries place each whole fibre in a regional physical receiver,
where the criteria of Section~\ref{sec:counting-tools} apply. All compact
blocks and auxiliary reserves are fixed before the common raw cap and
final restrictions in Lemma~\ref{lem:common-raw-receiver}.

\label{subsec:deg-exhaustion}

\begin{lemma}[Small-source coefficient decomposition]
\label{deg:lem-small-source-routing}
There are \(p_*,s_*,\gamma_*>0\) such that, for
\[
 \alpha _0=-s^2<0,\qquad
 0<s\le s_*,\qquad |\gamma|\le\gamma_*,\qquad |p|\le p_*,
\]
the actual passage domains have the coefficient cover in
Table~\ref{tab:coefficient-phase-cover} and
Appendix~\ref{app:proof-small-source-routing}.
The gamma, cubic, degree-loss, and remote-root coordinates are related
by the exact dictionaries \eqref{deg:gamma-degree-dictionary},
\eqref{deg:negative-far-dictionary}, and
\eqref{deg:cubic-degree-dictionary}. Their coefficient ranges overlap
strictly and all thresholds are fixed before the final raw-neighborhood
and multiplier shrink. The face \(p\le0\) is covered directly by
Lemma~\ref{deg:lem-raw-jacobian}.
\end{lemma}

The proof is in Appendix~\ref{app:proof-small-source-routing}.

\begin{lemma}[Central height versus signed-trace gap]
\label{lem:large-ratio-height}
Fix \(K_0,L_c>0\).  Suppose
\(\alpha _0=-s^2<0\), \(\alpha _1>0\), and put
\[
 p=\frac{\alpha _1}{s},\qquad r_L=\alpha _1,\qquad
 v-\gamma=r_LX,\qquad z=r_L^2Y_L.
\]
For \(p\ge1\), every through graph satisfying
\(Y_0=z(\gamma)/s^2\le K_0(1+p^2)\) obeys
\(Y_L(L_c)\le B_c\), with \(B_c\) independent of
\(p,s,\gamma\).  If the lower boundary of the through family has
nonnegative one-sided height \(d_L\) on \(X=L_c\), then
\begin{equation}
 0<\frac{\mathfrak g_L}{r_L^2}
 =Y_L(L_c)-d_L\le B_c.
 \label{deg:large-ratio-gap}
\end{equation}
\end{lemma}

The proof is in Appendix~\ref{app:proof-large-ratio-height}.

\begin{lemma}[Large-source trace domains]
\label{deg:lem-large-source-routing}
There exists \(P_*<\infty\) such that, for
\[
 \alpha _0=-s^2<0,\qquad p=\frac{\alpha _1}{s}\ge P_*,
\]
the part below the high-height threshold belongs to a fixed signed
trace block with bounded actual gap. Its terminal and compact strict
intervals have the overlap \eqref{deg:trace-overlap}; the remaining
regular pieces have bounded complete derivative. The construction is
uniform as \(p\to\infty\), without a bound on \(\gamma/s\).
\end{lemma}

The proof is in Appendix~\ref{app:proof-large-source-routing}.

\subsection{A coefficient cover and its raw-height budgets}
\label{sec:coefficient-raw-budgets}

The coefficient alternatives must be fixed before choosing a common
height cap. Their section constructions then give a raw-height cap
independently of the remaining coefficient radii. In all phase
derivatives below, the coefficients and actual finite multiplier are
held fixed.

\begin{lemma}[Coefficient-only cover]\label{lem:pure-coefficient-cover}
Fix the original nested sections, the passive compact coefficient block,
and the finite orders required by the regional preparations. There are
fixed thresholds
\(P_D,p_*,P_*,k_\flat,\Pi_{\rm far},\Pi_{\max},
\mathfrak q_\flat,U_d,U_{\rm cub},G_{\max}\), and fixed overlapping
coefficient rectangles, with the following properties.

For \(\alpha_0\ge0\), the determinant and signed-trace substitutions
\eqref{deg:determinant-chart}--\eqref{deg:trace-chart} cover every
nonzero coefficient pair. The double-zero pair is retained separately.
For \(\alpha_0=-s^2<0\), put \(p=\alpha_1/s\). The ranges
\begin{equation}\label{eq:pure-outer-coefficient-cover}
 p\le0,\qquad 0<p\le p_*,\qquad
 p_*/2\le p\le2P_*,\qquad p\ge P_*
\end{equation}
cover every coefficient fibre. The small-positive range is covered by the overlapping gamma,
compact-cubic, degree-loss and remote-root alternatives in
Table~\ref{tab:coefficient-phase-cover}.
Here the variables are exactly \eqref{deg:gamma-chart} and
\eqref{deg:cubic-chart}. The degree and remote-root substitutions are
\eqref{deg:gamma-degree-dictionary},
\eqref{deg:negative-far-dictionary}, and
\eqref{deg:cubic-degree-dictionary}. Their full independent rectangles
lie in the prescribed coefficient ranges of their preparations.
All thresholds and reserves are chosen before a final raw coefficient
restriction. Neither a phase input, a return zero, nor a height cap is
used to choose a coefficient alternative.
\end{lemma}

\begin{proof}
Fix a signed compact strict source block about \(a=p=0,Y(0)=1\).
Choose its high-height constant \(K_{\rm hi}\), large enough also for
the signed trace block \(1/2\le|p_L|\le3/2\), and put
\(K_{\rm lo}=2K_{\rm hi}\). The trace preparation on an initial
\(|a_L|\le a_0\le1\) uses the fixed central-height range
\(K_{\rm lo}(1+(3/2)^2)\). Its graph barrier fixes the bounded
actual-gap range; its terminal construction then fixes
\(Q,R_0,a_*,\xi_0\). Choose \(P_D>2/\sqrt{a_*}\).
When \(\alpha_0>0\), use \(r_D=\sqrt{\alpha_0}\) if
\(|\alpha_1|\le P_D\sqrt{\alpha_0}\). Otherwise put
\[
 r_L=|\alpha_1|,\qquad p_L=\operatorname{sgn}\alpha_1,
 \qquad a_L=\alpha_0/\alpha_1^2<a_*/4.
\]
The two coefficient descriptions overlap on
\(P_D/2<|\alpha_1|/\sqrt{\alpha_0}<P_D\).
For zero opening and nonzero trace, the second description has
\(a_L=0\).

For negative opening, first choose the bounded positive-quadratic
root-free and ghost reserves, and the independent positive far-root
bounds \(q_0,\lambda_0,\tau_0\) in
\eqref{deg:positive-far-range}. Choose the reciprocal preparation
bounds \(\mathfrak q_-,\Pi_-,\tau_-\), with its positive-unit lower
bound \(\omega_{\min}\), as in
\eqref{deg:negative-far-rectangle}. Put
\begin{equation}\label{eq:pure-routing-reserves}
\begin{gathered}
\mathfrak q_\flat=\tfrac12\min\{q_0,\mathfrak q_-\},\qquad
\Pi_{\rm far}=\max\{\lambda_0^{-1/2},\Pi_-,2\sqrt{12}\},\\
\Pi_{\max}=2\Pi_{\rm far}.
\end{gathered}
\end{equation}
On \(0<\pi\le\Pi_{\max}\), fix the two bounded-gamma preparations,
including the positive ghost raw-tail reserve. Choose \(k_\flat>0\)
so that \(2k_\flat\) lies below their fixed \(k\)-bounds and below
\(k_d(\mathfrak q_\flat)\). The latter is the bound placing the
gamma-degree dictionary in its prescribed compact positive-unit block
for \(\mathfrak q\ge\mathfrak q_\flat\). Set
\(G_{\max}=k_\flat^{-2/3}\). Fix \(U_d\) for the large-\(U\)
cubic-degree dictionary on this \(G\)-block and take
\(U_{\rm cub}>U_d+1\). Fix the compact cubic preparation on
\([0,U_{\rm cub}]\times[-G_{\max},G_{\max}]\).

These substitutions are algebraic. The gamma polynomial
\(t^3+3\epsilon t^2-\mathfrak qt+k^2\) and Vieta's relations give
\eqref{deg:gamma-degree-dictionary}; the positive-unit entries have
compact closure for \(\mathfrak q\ge\mathfrak q_\flat\).
On the negative far face, the same relations give
\(\omega=1-a+ja\), \(\tau=p^2/\omega^2\), and the full rectangle
\eqref{deg:negative-far-rectangle}. The large-\(U\) polynomial
\(y^3+3Gy^2-Uy+1\) gives
\eqref{deg:cubic-degree-dictionary}; its positive units tend uniformly
to one for bounded \(G\). Let \(\tau_d>0\) be the resulting compact
degree bound. Take both degree radial bounds below \(1/2\), including
their bounded-source, enlarged positive-energy, raw-image and relative
curvature reserves. Include the fixed core/high overlap reserve in
\(\tau_0\) and choose once
\begin{equation}\label{eq:pure-small-ratio-reserve}
0<p_*^2<\tfrac12\min\{\tau_0,\tau_d,
                              \tau_-\omega_{\min}^2\}.
\end{equation}
Thus \(\tau=p^2\) and \(\tau=p^2/\omega^2\) already lie in their
respective independent rectangles.

If \(\gamma\ne0\) and \(k\le2k_\flat\), the overlapping
\(\pi\)-ranges are \(\pi\le\Pi_{\max}\) and
\(\pi\ge\Pi_{\rm far}\). On the latter range split at
\(\mathfrak q_\flat\). The large-\(\mathfrak q\) dictionary is in
the compact degree block, and the small-\(\mathfrak q\) dictionaries
are in the positive far or negative reciprocal rectangle according to
the sign of \(\gamma\). If \(\gamma=0\) or \(k\ge k_\flat\),
then \(|G|^{3/2}\le k_\flat^{-1}\), and
\([0,U_{\rm cub}]\cup[U_d,\infty)\) covers \(U\ge0\).
The \(k\), \(\pi\), \(\mathfrak q\), and \(U\) descriptions
have the strict overlaps just chosen, with their equality faces retained.

For the large positive-source range use the exact trace dictionary
\eqref{deg:large-trace-dictionary}. Its signed opening is
\(a_L=-p_L^2/p^2\). Choose
\(P_*>\max\{2,p_*,3/\sqrt{a_*}\}\), increasing it if necessary
for the fixed bounded-gap trace block of
Lemma~\ref{lem:large-ratio-height}. Finally fix the compact positive
interval \([p_*/2,2P_*]\) and its endpoint, middle and high reserves.
The remaining independent radii \(r_D,r_L,g,t,s\) can subsequently
be restricted in raw coordinates. This proves the coefficient cover.
\end{proof}

\paragraph{The small-positive-source subcover.}
\label{sec:cover-index}\label{subsec:cover-index}
The outer overlaps are \(P_D/2<|\alpha_1|/\sqrt{\alpha_0}<P_D\),
\(p_*/2<p<p_*\), and \(P_*<p<2P_*\). With
\(H=z(\gamma)\), the low and high estimates overlap on
\(K_{\rm hi}(|\alpha_0|+\alpha_1^2)<H<
K_{\rm lo}(|\alpha_0|+\alpha_1^2)\), where
\(K_{\rm lo}=2K_{\rm hi}\). Every phase description is restricted
to the common receiver of Lemma~\ref{lem:common-raw-receiver}.

In the first three rows assume \(\gamma\ne0\), \(k\le2k_\flat\),
using \eqref{deg:gamma-chart}. The next two use \eqref{deg:cubic-chart}
when \(\gamma=0\) or \(k\ge k_\flat\), with \(|G|\le G_{\max}\).
\begingroup
\let\vss\vfil
\small
\setlength{\tabcolsep}{4pt}
\renewcommand{\arraystretch}{1.08}
\begin{longtable}{@{}>{\raggedright\arraybackslash}p{.31\linewidth}>{\raggedright\arraybackslash}p{.39\linewidth}>{\raggedright\arraybackslash}p{.23\linewidth}@{}}
\caption{The small-positive-source coefficient subcover. Each regional estimate
counts the whole physical fibre.}\label{tab:coefficient-phase-cover}\\
\hline
Range & Regional estimate & Overlap or boundary\\
\hline
\endfirsthead
\multicolumn{3}{l}{\tablename\ \thetable\ (continued)}\\
\hline
Range & Regional estimate & Overlap or boundary\\
\hline
\endhead
\hline
\endfoot
\(0<\pi\le\Pi_{\max}\)
 & Proposition~\ref{sharp:upper-gamma}: positive gamma \(3\), negative gamma \(2\)
 & \(\Pi_{\rm far}<\pi<\Pi_{\max}\); real, fold, ghost and root-free faces\\
\(\pi>\Pi_{\rm far}\), \(\mathfrak q\ge\mathfrak q_\flat\)
 & \eqref{deg:gamma-degree-dictionary}; Proposition~\ref{p2:prop-physical-degree-count}: \(2\)
 & Far receivers extend past \(\mathfrak q_\flat\)\\
\(\pi>\Pi_{\rm far}\), \(\mathfrak q<\mathfrak q_\flat\)
 & Positive gamma: \eqref{deg:far-negative-budget} and Proposition~\ref{sharp:upper-gamma}; negative gamma: \eqref{deg:negative-far-dictionary} and Proposition~\ref{p2:prop-physical-degree-count}; both \(2\)
 & \(\mathfrak q_\flat<\mathfrak q_-\), \(\Pi_{\rm far}\ge\Pi_-\)\\
\(0\le U\le U_{\rm cub}\)
 & Proposition~\ref{sharp:upper-cubic}: \(3\), including the original raw ghost tail
 & \(k_\flat<k<2k_\flat\), \(U_d<U<U_{\rm cub}\)\\
\(U\ge U_d\)
 & \eqref{deg:cubic-degree-dictionary}; Proposition~\ref{p2:prop-physical-degree-count}: \(2\)
 & \(U_d<U<U_{\rm cub}\)\\
\end{longtable}
\endgroup

The degree-loss phase subcover and its internal dividers are those of
Proposition~\ref{p2:prop-unified-carrier-cover}. Each row selects a
whole-fibre bound; Proposition~\ref{sharp:upper-coefficient-bounds}
will take their maximum after the common raw restrictions.


\begin{proposition}[Regional raw-height budgets]
\label{prop:regional-raw-budgets}
Fix the coefficient rectangles and all section and overlap reserves of
Lemma~\ref{lem:pure-coefficient-cover}. Let \(\mathcal A_j\) be their
finitely many regional alternatives and write
\[
\mathcal V_\delta=\{\max(|A-1|,|C-1|,|\gamma|,
                                  |\alpha_0|,|\alpha_1|)<\delta\}.
\]
There is \(b>0\) such that, for every
\(0<h\le b\) and each \(j\), there are independent raw bounds
\(\delta_j(h)>0\) and \(\bar m_j(h)>0\) for which the actual regional
transition identities and complete derivative and curvature estimates
hold on their stated phase domains for every
\(\theta\in\mathcal A_j\cap\mathcal V_{\delta_j(h)}\) and every
\(0\le m\le\bar m_j(h)\). These domains together cover every input of
\begin{equation}\label{eq:regional-raw-budget-domain}
\{x:a_\theta+x\in J_\theta^{\rm out},\quad
                         0<h_\theta(a_\theta+x)<h\}.
\end{equation}
The bounds restrict \((A,C,\gamma,\alpha_0,\alpha_1)\) in the
independent raw coordinates and \(0\le m\le\bar m_j(h)\).
These additional restrictions impose no relation between active radii
or root gaps.
Each estimate retains its original input and inverse output windows,
positive-time passage, speed and stop clips, admitted equality dividers,
and intervening nonroot inputs. Its raw observers and divided closing
are composed on their full admitted images, with ordinary phase
\(C^3\) bounds and all nonlinear curvature terms retained.
\end{proposition}

\begin{proof}
We first construct a fixed signed high-height cap on the original raw
sections. Fix the enlarged signed strict source block used in the
coefficient lemma. Choose the core cut \(L\) and raw port radius
\(0<d<\delta\), with \(d\le1/100\), as in
Lemma~\ref{ann:p0-fixed-tail-ports}. Let \(a_H,p_H\) bound its
normalized \(|a|,|p|\), and increase \(L\) until
\(p_H/L\le1/40\), \(a_H/L^2\le1/4000\).
The fixed-core continuity bounds \(s_{\rm core},\rho_{\rm core}>0\)
put both entrance ratios in \([1/4,3/4]\) for independent
\(s\le s_{\rm core}\) and
\(|\gamma|,|A-1|,|C-1|\le\rho_{\rm core}\).

Choose a fixed transverse tube from these ports to the original
sections, and let \(\mu_z>0\) be smaller than its port tolerance and
both raw-window margins. With \(T=\log(d/(sL))\), the entrance-memory
estimate of Lemma~\ref{ann:p0-fixed-tail-ports}, with its fixed
\(\lambda=7/5\), gives
\begin{equation}\label{eq:raw-high-entrance-memory}
 |w_\sigma(T)-w_\Gamma(d)|
 \le C\{s+|\gamma|+d|A-1|+d^2|C-1|\}.
\end{equation}
Indeed its initial term is \(O(e^{-\lambda T})=O(s)\), and
integration of the two source terms has orders \(e^{-T}\) and
\(e^{-\lambda T}\); the passive terms have the three orders displayed.
The constant depends on the fixed ports and strict block, not on the
flight length. Choose \(s_H>0\) and independent passive coefficient
bounds so that the right side times \(d^2\), plus \(2d|\gamma|\),
is less than \(\mu_z\). The invariant ratio strip from the same lemma
keeps both complete tails in the chart. Their raw port errors are
therefore below \(\mu_z\), so both full images lie in the original
windows. The regular observers and divided closing can now be
composed on their larger fixed buffers.

At any signed high input freeze \(r_0=\sqrt H\). The transformed
parameters satisfy \(|a|\le K_{\rm hi}^{-1}\),
\(|p|\le K_{\rm hi}^{-1/2}\), and \(Y(0)=1\).
If \(H\le s_H^2\), this is the same admitted strict block.
The two absolute tail variations in
\eqref{sharp:upper-tail-ratio} have the common
\(r_0^{2-\lambda_\gamma}\) factor; its cancellation, followed by
the actual bridges and divided closing, gives a fixed \(C_H\).
The dilation is frozen during differentiation. This proves the signed
high estimate without a restriction on \(\gamma\log r_0\).

Let \(z_{\rm geo}>0\) be a fixed central reserve below the upper
chart boundary constructed before Lemma~\ref{lem:loc-through}, and
bound \(0<\ell_\theta(\gamma)\le L_{\rm raw}\) on the preliminary
raw block. Choose
\begin{equation}\label{eq:regional-fixed-raw-cap}
b=\tfrac14\min\{s_H^2,z_{\rm geo},d_{\rm out}/(4L_{\rm raw})\}>0.
\end{equation}
Every signed high input below \(b\) consequently has its original
raw-domain description. With \(|\alpha_0|<d_{\rm out}/4\), its
speed satisfies
\(\sigma_\theta\le|\alpha_0|+L_{\rm raw}H<d_{\rm out}/2\).

Now fix any \(0<h\le b\). The lower and strict preparations have
already fixed their compact blocks, chords and positive raw-window
reserves. We place their comparison and upper shots below \(h\)
by the following independent raw restrictions. In the determinant and
trace charts it is enough to require
\begin{equation}\label{eq:regional-det-trace-cap-bounds}
r_D^2<\frac{h}{8K_{\rm hi}(1+P_D^2)},\qquad
r_L^2<\frac{h}{4K_{\rm lo}}.
\end{equation}
Their bounded normalized-height complements have the finite strict
domains of \eqref{deg:det-height-cover} and
\eqref{deg:trace-domain-inclusion}; their larger heights are in the
signed high block. On the fixed compact positive-ratio region the
enlarged strict box \([u_L/2,3u_H]\) is placed below \(h/4\) by
\(s^2<h/(12u_H)\). The endpoint admission and its actual matching
are those of Lemma~\ref{ann:lem-compact-positive-domains}.

For a bounded-gamma preparation let \(k_B,\Pi_B\) be its fixed
coefficient bounds and \(E_B>0\) its fixed upper core height. Retain
the action comparison bound and impose
\begin{equation}\label{eq:regional-gamma-cap-bounds}
g<\min\left\{1,\frac{E_B}{4K_{\rm hi}(1+\Pi_B^2)},
                       \frac1{2k_B}\sqrt{\frac h{E_B}}\right\}.
\end{equation}
Then the normalized high threshold is below \(E_B/2\), while
\((gk)^2E_B<h/4\). The actual comparison and upper shots and every
intervening graph retain the outward admission used in
\eqref{deg:negative-gamma-cover} and in the positive paired
construction. The positive ghost raw tail is part of the same original
domain, including \(gku\ge d_0\); its raw derivative estimate is
\eqref{bp-tail:consumed-derivative}.

For the compact cubic block let \(E_c>0\) be its fixed upper core
height. Retain the exponentially small comparison bound and require
\begin{equation}\label{eq:regional-cubic-cap-bounds}
\tau_{\rm cub}<\min\left\{1,
 \frac{E_c}{4K_{\rm hi}(1+U_{\rm cub}^2)},
 \tfrac12\sqrt{h/E_c}\right\}.
\end{equation}
Here \(\tau_{\rm cub}=(-\alpha_0)^{1/3}\). Thus the normalized
threshold is below \(E_c/2\) and
\(\tau_{\rm cub}^2E_c<h/4\). The comparison remains in the actual
lip domain, with the bounded/growing resection and the original raw
ghost continuation retained in \eqref{deg:compact-physical-cover}.

For the positive far-root rectangle, the fixed \(\tau_0\) already
ensures \(\tau K_{\rm hi}(1+\tau)<E_f/2\) and the action comparison,
where \(E_f>0\) is the selected upper shot. Since
\(a=gq\lambda\), \(q\le q_0\), \(\lambda\le\lambda_0\), the
raw restriction
\begin{equation}\label{eq:regional-far-cap-bound}
g<\frac1{2q_0\lambda_0}\sqrt{h/E_f}
\end{equation}
gives \(a^2E_f<h/4\). The full lower output and upper core images
are the admitted images used in \eqref{deg:far-physical-cover}.
No ratio threshold is changed after \(h\) is fixed.

For the degree preparations, the exact source dictionary
\eqref{deg:degree-source} gives
\begin{equation}\label{eq:degree-cap-raw-identities}
\alpha_1=\tau^2\kappa^2\mathcal R\omega,\qquad
s_s^2=\tau^3\kappa^4\mathcal R^2B.
\end{equation}
The actual \(h_L=0\) source graph and its real enlarged chord have
\(Y_0\le C_Y(1-\log(2\tau))\) by
\eqref{bd:positive-height} and the enlarged positive-energy construction
in Lemma~\ref{p2:lem-bounded-source-budget}. Use the degree bounds
\(\kappa\le\kappa_*\), \(\tau\le\tau_*<1/2\) already fixed in
the coefficient cover, and put
\begin{equation}\label{eq:degree-cap-compact-constant}
D_*=\tau_*[1-\log(2\tau_*)],\qquad
C_d=C_Y\kappa_*^2
       \frac{\mathcal R_{\max}B_{\max}}{\omega_{\min}}D_*.
\end{equation}
The function \(t[1-\log(2t)]\) increases on \((0,1/2)\).
Consequently the actual trace \(h^0\) satisfies
\begin{equation}\label{eq:degree-cap-height-raw-bound}
H_\theta(h^0)\le C_Y\alpha_1\tau\kappa^2
       \frac{\mathcal RB}{\omega}[1-\log(2\tau)]
       \le C_d\alpha_1.
\end{equation}
On the reciprocal graph \eqref{deg:degree-infinity-graph}, write
\(a_{\rm rec,*}=\kappa_*A_{d,\max}<1/2\) for its fixed opening bound.
There
\[
\frac{\mathcal RB}{\omega}
=\frac{dA_dj}{1-a+ja}
\le\frac{d_{\max}A_{d,\max}j_0}{1-a_{\rm rec,*}}.
\]
Thus \eqref{eq:degree-cap-height-raw-bound} holds with
\begin{equation}\label{eq:degree-cap-reciprocal-constant}
C_r=C_Y\kappa_*^2
      \frac{d_{\max}A_{d,\max}j_0}{1-a_{\rm rec,*}}D_*
\end{equation}
in place of \(C_d\), uniformly as \(j\downarrow0\). The independent
raw restriction
\begin{equation}\label{eq:degree-cap-alpha-one-bound}
0<\alpha_1<\frac h{4\max\{C_d,C_r\}}
\end{equation}
places the upper source shot and its chosen enlarged chord below
\(h/4\). The fixed positive-label threshold also gives
\(H_\theta(h^0)>K_{\rm hi}(|\alpha_0|+\alpha_1^2)\), so that this
chord overlaps the same signed high continuation.
The fixed \(p_*\) and degree ratios are unchanged.

Raw output admission in these degree charts is independent of a
reciprocal lower bound on a gain. Indeed
\(r_s^2=(\mathcal R/\omega)\alpha_1\le C\alpha_1\); for a fixed
\(|\lambda_\gamma|<\nu<1\), the retained absolute factor obeys
\[
\Omega=r_s^{2-\lambda_\gamma}
       \le C\alpha_1^{1-\nu/2}\longrightarrow0.
\]
The stopped raw trace has a fixed window margin, so the full enlarged
output chord is admitted before its observer is composed. The exact
input buffer, actual matching and every-point allocation are those of
Proposition~\ref{p2:prop-unified-carrier-cover}. Actual reciprocal
maps use \(j>0\); \(j=0\) retains its one-sided normalized meaning.
The raw-Jacobian estimate already holds on the full original raw-window
and inverse-output-window domain and can be restricted to any such cap.

All these descriptions are restrictions of the original physical map
\(\Pi=\mathcal C_{\theta,m}\circ T_\theta(a_\theta+\cdot)\).
For local nonlinear observers write
\(\Pi=\psi_-\circ P\circ\iota\), where
\(\iota=\psi_+^{-1}\) on the actual incoming image. In particular
the physical curvature contains all three terms
\begin{equation}\label{eq:raw-budget-complete-curvature}
\Pi''=\psi_-''(P)(P')^2(\iota')^2
       +\psi_-'(P)P''(\iota')^2
       +\psi_-'(P)P'\iota''.
\end{equation}
Here the derivatives of \(P\) are evaluated at \(\iota(x)\). The regional
estimates include these ordinary observer derivatives, through
order three, before any positive affine normalization
\eqref{sharp:upper-affine}. Restriction leaves those estimates unchanged.
The divided closing has the ordinary \(C^3\) rows
\eqref{phys:eq-closing-jets}, including \(m=0\).

Take the finite minimum of the original raw-domain and error bounds
and the independent restrictions above. Express
\(\tau_{\rm cub}<T\) as \(|\alpha_0|<T^3\),
\(r_D<R\) as \(\alpha_0<R^2\), and
\(r_L<R\) as \(|\alpha_1|<R\). Half this minimum defines
\(\delta_j(h)>0\) in the independent raw coordinates. All other
normalized coordinates remain in the rectangles already fixed by
Lemma~\ref{lem:pure-coefficient-cover}.

Lastly let \(C_j\) bound the complete derivative on the regional
bounded pieces and let \(M_j\) bound the full translated output on
the original closing image. Intersect the ordinary closing-domain
bound, the regional critical-logarithm bound, and the inequalities
\(mC_j<1/2\), \(mM_j<a_{\rm cl}/2\). Their finite positive
minimum gives \(\bar m_j(h)>0\). The enlarged domain and equality
buffers were chosen earlier; their estimates hold at nonclosed inputs
as well. This proves the asserted raw budgets.
\end{proof}

\subsection{The common raw receiver}
\label{sec:conclusion}

We now use the coefficient cover and regional raw budgets to define one
interval on which all regional estimates hold for the complete return.

\begin{lemma}[The common raw receiver]\label{lem:common-raw-receiver}
Fix the coefficient cover and reserves of
Lemma~\ref{lem:pure-coefficient-cover}, and let \(b>0\) be the raw
budget of Proposition~\ref{prop:regional-raw-budgets}. For every fixed
\(0<h_c\le b\), there are a raw coefficient neighborhood and a
fixed \(\bar m>0\) such that all covering regional constructions
and their every-point estimates hold on the same actual interval
\begin{equation}\label{eq:common-raw-carrier}
 I_\theta(h_c)=\{x:a_\theta+x\in J_\theta^{\rm out},\quad
                    0<h_\theta(a_\theta+x)<h_c\}.
\end{equation}
This interval retains the original input and inverse output windows,
positive-time passage, speed and stop clips, admitted equality dividers,
and all intervening nonroot inputs. The normalized regional descriptions
use their exact centres and positive absolute gains and the same ordinary
\(C^3\) divided closing. The coefficients and actual \(m\) are fixed
in every phase derivative. The neighborhood and multiplier bound may
depend on \(h_c\), but the coefficient cover and its reserves do not.
\end{lemma}

\begin{proof}
Fix any \(0<h_c\le b\). Proposition~\ref{prop:regional-raw-budgets}
gives the finitely many independent bounds
\(\delta_j(h_c),\bar m_j(h_c)\) on the full actual regional
domains. Intersect their raw neighborhoods, including
\(|\alpha_0|<d_{\rm out}/4\). Let \(m_{\rm ad}>0\) be the
ordinary closing-domain bound, and choose
\begin{equation}\label{eq:common-multiplier-bound}
\bar m=\tfrac12\min\{m_{\rm ad},\bar m_1(h_c),\ldots,
                          \bar m_N(h_c),(2C_H)^{-1}\}>0.
\end{equation}
The complete displacement \eqref{eq:return} is therefore defined
on every admitted regional image, with its original divided observer
\eqref{deg:det-divided-closing}. All three nonlinear curvature terms
in \eqref{eq:raw-budget-complete-curvature} remain in its estimates.

The original raw admission set is exactly
\[
J_\theta^{\rm out}=
\mathcal T_\theta\cap\operatorname{int}K_+^\sharp
\cap T_\theta^{-1}(\operatorname{int}K_-^\sharp)
\cap\{\sigma_\theta<d_{\rm out}\}.
\]
Scalar uniqueness and first variation order the actual through graphs.
Thus \(\mathcal T_\theta\), both window preimages, the speed clip,
and the central-height clip are intervals; central height is strictly
increasing by \eqref{loc:central-height}. Their intersection
\eqref{eq:common-raw-carrier} is one interval, possibly empty.

The regional lower, core and high sets cover each admitted through
input before clipping. Intersect them with this same interval. Their
actual matching identities, inverse images and every-point estimates
survive. In particular \eqref{deg:det-height-cover},
\eqref{deg:trace-domain-inclusion}, \eqref{deg:negative-gamma-cover},
\eqref{deg:far-physical-cover}, \eqref{deg:compact-physical-cover},
and Proposition~\ref{p2:prop-unified-carrier-cover} retain their
admitted positive equality dividers. The compact positive-ratio sets
remain the monotone height preimages of
Lemma~\ref{ann:lem-compact-positive-domains}. If a clip removes an
overlap, the retained connected interval is still covered by the
adjoining enlarged estimates. The natural stopped endpoints are
excluded. The original compact-cubic and bounded-positive ghost raw
tails retain \eqref{deg:compact-raw-tail-derivative} and
\eqref{bp-tail:consumed-derivative}, respectively.

The fixed cap is independent of the normalized overlap heights.
Proposition~\ref{prop:regional-raw-budgets} places all comparison
shots and upper core endpoints below \(h_c/4\) by later raw
restrictions and covers the high complement up to \(h_c\).
For degree loss it uses
\eqref{eq:degree-cap-height-raw-bound}--\eqref{eq:degree-cap-alpha-one-bound}
on both the compact and reciprocal rectangles. Hence no ratio reserve
or coefficient threshold is selected again when \(h_c\) is reduced.
Set \(h_{\rm out}=h_c\) in \eqref{loc:analytic-receiver}.

Finally \eqref{eq:finite-multiplier}, with
\(\eps_f=h(\theta)K\) and positive \(h\), gives the actual
\(m(\theta)\le\bar m\) on \(K>0\) after a further raw coefficient
shrink. This restriction is independent of \(s,|\gamma|\) and of
all root gaps. For \(K\le0\) the selected finite channel has no
positive-time exit. Thus the receiver and all its regional estimates
are fixed before any zero count or geometric capture.
\end{proof}

One admissible fixed choice for the subsequent argument is
\begin{equation}\label{eq:common-raw-cap}
h_c=b/2>0.
\end{equation}

\subsection{Pairing the admitted raw tails}
\label{subsec:paired-raw-tails}

On a ghost fibre, a very small central height can come from a first lip
whose raw abscissa lies beyond the growing curvature range. The original
through interval retains these inputs. We pair their actual inward
heights at the original centre and compare the complete input and output
gains. The cubic endpoint power and the quadratic
nonpositive long-dwell correction control this ratio as the negative lip
approaches its root.

\begin{proposition}[Paired raw-tail contraction]
\label{rawpair:prop-original-tails}
Fix the original nested raw-section windows, their invariant upper
boundary, the passive compact coefficient neighbourhood and the ordinary
\(C^3\) regular maps and divided closing of Lemma~\ref{lem:closing}.
Use one positive unscaled cap \(h_c\) in the common receiver
\eqref{eq:common-raw-carrier}, chosen before the final coefficient and
multiplier restrictions. In each of the following regions choose the
fixed raw divider \(d_0\) and buffer \(a_b=d_0/2\) in the original
chart.

For the compact cubic coefficients
\eqref{deg:compact-cubic-field}--\eqref{deg:compact-cubic-polynomials},
\(0\le U\le U_0\), \(|G|\le G_0\), retain the strict ghost fibres
\(\Delta<0\) and put \(\tau=\tauc\). For the bounded positive
quadratic coefficients \eqref{deg:quadratic-field}, retain
\(\pi_\flat\le\pi\le\Pi\), \(-\Delta_0\le\Delta(k,\pi)<0\)
and independently small \(g,k>0\); choose \(d_0<a_0\), where
\(a_0\) is the original positive outward reserve of
Corollary~\ref{deg:cor-positive-large-S-outer}.

Let \(x_{\rm comp}\) be the native admitted comparison divider of
the corresponding regional construction. Every original through input
\(x\le x_{\rm comp}\) has a unique first positive lip.
For every such input with respectively
\(a=\tau u\ge a_b\) or \(a=gku\ge a_b\), its actual inward
height at the original centre has a unique negative match, with
respectively \(r_0\le r<c\) or \(r\ge r_c\),
\(0<R=kr<\mathfrak c\). The matched graph is its original through
trajectory, including its original output and closing admissions.
The complete physical derivative satisfies
\begin{equation}\label{rawpair:eq-contraction}
 0<\Pi_{\theta,m}'(x)\le
 \begin{cases}
 C(d_0)\tau^3,&\text{compact cubic},\\
 C(d_0)g^3,&\text{bounded positive quadratic}.
 \end{cases}
\end{equation}
The constants are uniform in the independent coefficient ranges,
the ghost gap, the negative root distance and \(0\le m\le m_0\).
All phase derivatives fix the coefficients and \(m\). In particular
the estimates retain every original admitted \(a\ge d_0\), all its
window, speed and height clips, and the inputs between return roots.
For the actual finite channel \(K_{\rm fin}>0\), one last common
multiplier restriction gives \(F'>1/2\) on both tails.
\end{proposition}

\begin{proof}
Use the original-domain constructions in
Proposition~\ref{deg:prop-compact-cubic} for the cubic field and
Proposition~\ref{bp-tail:prop-actual-raw-tail} for the quadratic
field. We identify their matches with the common physical receiver
and recall the complete-gain estimates needed for the contraction.

The cubic construction first admits a root-free comparison and a
reference pair in a buffered band about \([2M,3M]\). Its bounded
and high-action descriptions overlap through
\eqref{rc:strong-matching} and the actual resection
\eqref{deg:compact-resection}. This construction does not require
the cubic action to diverge as the positive lip label grows.
The quadratic comparison is the admitted pair of
Lemma~\ref{deg:lem-positive-quadratic-paired-variation}.
On its upper comparison interval, both original raw-section derivatives
in \eqref{deg:quadratic-paired-raw-slopes} contain the common factor
\((gk)^{2-\lambda_g}\), where \(\lambda_g=12g/(1+2g)\).
This factor cancels exactly in their quotient, and
\eqref{deg:quadratic-paired-core-ratio} bounds the remaining core ratio.
On the original ghost tail, the negative-root factor is instead
absorbed by the actual centre variation, as computed below.
The first-lip argument in each native proof applies to every original
through input below its comparison: the fixed window floor and scalar
order give the hit, and the strict transverse sign at a noncentral hit
makes it unique. In the affine physical input its label satisfies
\(a_x<0\).

For inward matching, use the actual original-centre maps in
\eqref{deg:compact-raw-tail-inward}--\eqref{deg:compact-reflected-inward}
and
\eqref{bp-tail:positive-inward-variation}--\eqref{bp-tail:negative-inward-absolute}.
Their full-range monotonicity and absolute endpoint constructions give
\[
 \eta_0^-([r_0,c))=(0,\eta_0^-(r_0)]\quad\hbox{and}\quad
 \eta_0^-([r_c,\mathfrak c/k))=(0,\eta_0^-(r_c)]
\]
in the cubic and quadratic cases, respectively. The negative memory
retains its initial polynomial and the simple-root action diverges at
the excluded endpoint. Decreasing positive centre height from the
reference to \(a_b\), then using its full-range monotonicity, puts
every \(a\ge a_b\) in this negative image. The construction and
its strict overlap with the growing range are proved in
\eqref{deg:compact-negative-centre-image} and
\eqref{bp-tail:whole-tail-matching}.

For an input in the common receiver, this matched graph and the
original through graph have the same positive height at the original
centre. Scalar uniqueness identifies them. Thus the matched output and
closing image are the original admitted ones; the auxiliary match at
\(a_b\) is made before observation. The raw outward passages and
window buffers are those of \eqref{deg:compact-raw-tail-input},
\eqref{bp-tail:positive-raw-graph},
Lemma~\ref{deg:lem-weighted-cubic-tail} and
Corollary~\ref{deg:cor-negative-absolute-outer}.

For the cubic gain calculation, use
\eqref{deg:compact-raw-tail-derivative}. Its proof keeps the factor
\(X^{\beta_-}\) in the negative-root output and uses the actual
central variation to compare it with the positive input gain. For the
quadratic calculation, \eqref{bp-tail:negative-consumed-gain} absorbs
the potentially singular negative-root gain into differentiation with
respect to the actual centre height; the positive estimate in the same
proof then gives \eqref{bp-tail:consumed-derivative}. In both cases
these are complete gain quotients, including the ordinary observers
and divided closing. The preceding identification therefore gives
\eqref{rawpair:eq-contraction} on the common receiver, without
estimating either gain ratio anew.

The regional estimates use the original independent coefficient blocks
and root distances. Their intrinsic collar derivatives precede the
ordinary \(C^3\) section chain of
Lemma~\ref{phys:lem-c3-adapters}. Restricting to the common receiver
preserves the original windows, inverse output, speed, stop and height
clips. The quadratic tail is assigned \(a\ge d_0\), whereas
Proposition~\ref{bp-tail:prop-actual-raw-tail} proves its estimate already
for \(a\ge a_b=d_0/2\). It therefore strictly overlaps the curvature
domain \(a<d_0\) on \(a_b\le a<d_0\), with \(a=d_0\) retained in the tail.
Finally choose \(m_*\le m_0\) with
\(m_*C(d_0)\max\{\tau_0^3,g_0^3\}<1/2\).
Then \(F'=1-m\Pi'>1/2\) on both original admitted tails.
\end{proof}

\subsection{Whole-carrier coefficient synthesis}
\label{sharp:upper-section}
\label{sec:whole-carrier-applications}

% Compatibility inputs; the proofs are in the regional owners.
% Historical aggregate compatibility file.
% The degree-loss and compact-cubic proofs formerly collected here now have
% canonical active owners in sharp-upper-degree.tex and sharp-upper-cubic.tex.
% The complete pre-migration aggregate remains available in Git history.

% Historical aggregate compatibility file.
% The determinant and signed-trace proofs now have their canonical active owner
% in sharp-upper-local-regional.tex, which is reached from Section 13.
% The pre-migration aggregate remains available in Git history.


\begin{proposition}\label{sharp:upper-coefficient-bounds}
With the coefficient thresholds of Lemma~\ref{lem:pure-coefficient-cover}
and the common receiver
\(I_\theta(h_c)\) of Lemma~\ref{lem:common-raw-receiver}, the following
whole-carrier bounds hold after one raw coefficient and multiplier
restriction. Isolated zeros are counted with multiplicity.
\[
\begin{array}{c|c}
 \text{coefficient alternative}&\text{number of roots}\\ \hline
 \alpha_0\ge0&2\\
 \alpha_0<0,\ p\le0&1\\
 \alpha_0<0,\ 0<p\le p_*&3\\
 \alpha_0<0,\ p_*/2\le p\le2P_*&2\\
 \alpha_0<0,\ p\ge P_*&2
\end{array}
\]
Here \(p=\alpha_1/\sqrt{-\alpha_0}\) when \(\alpha_0<0\).
The alternatives cover every sufficiently small raw coefficient and each
controls its whole phase fibre. Their maximum is three.
\end{proposition}
\begin{proof}
Use the fixed coefficient cover of Lemma~\ref{lem:pure-coefficient-cover}
and the raw budgets of Proposition~\ref{prop:regional-raw-budgets}.
Lemma~\ref{lem:common-raw-receiver} has already placed every regional estimate
on the same connected interval \(I_\theta(h_c)\), with the final
coefficient and multiplier restrictions in force.

If \(\alpha_0>0\), the determinant chart applies when
\(|\alpha_1|\le P_D\sqrt{\alpha_0}\), and signed trace applies
otherwise. The exact trace dictionary gives
\(0<a_L<a_*/4\). Propositions~\ref{sharp:upper-determinant}
and \ref{sharp:upper-trace} give two zeros on determinant or positive
trace, and one on negative trace. If \(\alpha_0=0\) and
\(\alpha_1\ne0\), the trace opening is zero. At
\(\alpha_0=\alpha_1=0\), every positive height is in the signed high
region of Corollary~\ref{rew:cor-high-derivative-signed}; the final
multiplier bound makes the return contractive. Zero height is stopped.

For \(\alpha_0<0\) and \(p\le0\),
Lemma~\ref{deg:lem-raw-jacobian} gives at most one zero.
For \(0<p\le p_*\), the coefficient alternatives place each whole
fibre in a bounded-gamma, far-root, compact-cubic or degree-loss receiver.
Propositions~\ref{sharp:upper-gamma} and \ref{sharp:upper-cubic}
give at most three.
The remaining original cubic and bounded-positive quadratic ghost tails
have \(F'>1/2\) by Proposition~\ref{rawpair:prop-original-tails},
with both full gains and their original admissions retained.
The positive far-root budget
\eqref{deg:far-negative-budget}, physical cover
\eqref{deg:far-physical-cover}, and the degree-loss carrier and count
of Propositions~\ref{p2:prop-unified-carrier-cover} and
\ref{p2:prop-physical-degree-count} give two in their alternatives.
Their independent rectangles and overlap reserves are those fixed by
Lemma~\ref{lem:pure-coefficient-cover}, before \(h_c\) was selected.
The two three-root arguments bound the zeros of the actual lower \(F'\)
by two and have \(F'>0\) on the overlapping upper continuation.
Lemmas~\ref{sharp:upper-curvature-count} and
\ref{sharp:upper-critical-count} therefore count the same connected
lower/core/high interval. Both nonlinear section maps and all three
terms of the physical curvature identity remain in these estimates.

For \(p\ge P_*\), the exact large-source dictionary
\eqref{deg:large-trace-dictionary} has \(p_L=1\) and
\(a_L=-p^{-2}\); Proposition~\ref{sharp:upper-trace} gives two
uniformly on this whole sector. On \([p_*/2,2P_*]\),
Lemma~\ref{ann:lem-compact-positive-domains} and
Proposition~\ref{sharp:upper-annulus} give two on the whole receiver.
The three positive ranges overlap strictly. At every coefficient one
covering argument counts all of \(I_\theta(h_c)\), including its
intervening nonroot inputs and admitted equality dividers. The maximum
of the five whole-fibre bounds is three. Proposition~\ref{prop:central-bound}
then applies geometric capture to this already counted receiver.
\end{proof}

\subsection{Geometric capture and the upper bound}
\begin{proposition}[Central three-cycle upper bound]\label{prop:central-bound}
There are a coefficient neighborhood \(\mathcal V_5\) and
\(\eps_5>0\) such that every field in the five-parameter family with
coefficient in \(\mathcal V_5\) has at most three isolated periodic
orbits \(C\) with \(d_{\rm H}(\iota_0(C),\Gamma)<\eps_5\).
\end{proposition}
\begin{proof}
Fix the windows, coefficient cover, reserves, \(h_c\), and multiplier
bound of Lemma~\ref{lem:common-raw-receiver}. Choose
\[
 0<h_{\rm in}<\min\{h_c/2,d_{\rm in}/(4L_{\rm raw})\},
 \qquad |\alpha_0|<d_{\rm in}/4.
\]
Every through orbit of central height \(H<h_{\rm in}\) satisfies
\begin{equation}\label{eq:height-capture}
 \sigma_\theta\le-D_\theta(\gamma,H)
 =-\alpha_0-Q_\theta(\gamma)H
 \le|\alpha_0|+L_{\rm raw}H<d_{\rm in}/2.
\end{equation}
Intersect the analytic neighborhood with \(|\alpha_0|<d_{\rm in}/4\)
and call it \(\mathcal V_{\rm an}\).
Only now apply Lemma~\ref{lem:loc-capture} with these original inner
windows, speed bounds and fixed \(h_{\rm in}\). It supplies
\(\mathcal V_0,\eps_0\); set
\[
 \mathcal V_5=\mathcal V_{\rm an}\cap\mathcal V_0,
 \qquad \eps_5=\eps_0.
\]
Every collar cycle is admitted to the inner windows and height cap,
follows the complete marked positive-time circuit with \(K>0\), and
gives one zero of \eqref{eq:return} in \(I_\theta(h_c)\) at its unique
prime return-section hit.

Each covering alternative in Proposition~\ref{sharp:upper-coefficient-bounds}
controls all of \(I_\theta(h_c)\), including intervening nonroot inputs, so
\begin{equation}\label{eq:total-bound}
 N_{\rm cen}\le\max\{2,1,3,2,2\}=3.
\end{equation}
The regional arguments control the connected lower/core/high union.
On coefficient overlaps, either applicable whole-interval bound is used;
local phase counts are not added. Distinct cycles give
distinct zeros and transverse return preserves isolation, proving the
assertion on this one neighborhood and collar.
\end{proof}

\subsection{Concentration of three-cycle parameters}
\label{sec:three-cycle-concentration}

The two-cycle estimate on a compact positive source interval also
locates the coefficient sequences that can carry three cycles. We use
the canonical slice throughout: for a field \(F\) near \(F_*\), put
\(\theta_F=\theta(\mathscr C_0(F))\), with
\(\mathscr C_0\) from \eqref{eq:fixed-coefficient-transport} and
\(\theta\) from Lemma~\ref{lem:slice}. Define
\(\gamma(F),\alpha_0(F),\alpha_1(F)\) by
\eqref{eq:center-parameters} at \(\theta_F\), and, when
\(\alpha_0(F)<0\), set
\[
 s(F)=\sqrt{-\alpha_0(F)},\qquad
 p(F)=\frac{\alpha_1(F)}{s(F)}.
\]
Thus the ratio below refers to this fixed analytic slice.

\begin{proposition}[Uniform two-cycle exclusion on compact positive ratios]
\label{prop:compact-positive-two-cycle-exclusion}
For every fixed \(0<\eta\le P<\infty\), there are an open
coefficient neighborhood \(\mathcal V_{\eta,P}\) of \(F_*\) in
the full twelve-coefficient quadratic germ and
\(\eps_{\eta,P}>0\) such that every
\(F\in\mathcal V_{\eta,P}\) with
\(\alpha_0(F)<0\) and \(\eta\le p(F)\le P\) has at most
two distinct isolated periodic orbits \(C\) satisfying
\[
 d_{\rm H}(\iota_*(C),\Gamma_{I4})<\eps_{\eta,P}.
\]
The return estimate holds on one complete physical receiver, retaining
the original windows and clips and all inputs between its roots.
\end{proposition}

\begin{proof}
Apply Lemma~\ref{ann:lem-compact-positive-domains} and
Proposition~\ref{sharp:upper-annulus} on the enlarged fixed ratio
interval \([\eta/2,2P]\). They give the two-zero multiplicity bound
on one complete physical receiver, including its endpoint, middle and
high ranges. Choose its cap so that
\[
 0<b_{\eta,P}\le\min\{h_c,z_*\},\qquad
 L_{\rm raw}b_{\eta,P}<d_{\rm out}/4,
 \qquad |\alpha_0|<d_{\rm out}/4.
\]
All constants may depend on \(\eta,P\). The ordinary coefficient
restriction in \eqref{eq:finite-multiplier} makes the actual finite
multiplier satisfy the final bound used by that proposition.

Now choose
\(0<h_{\rm in}<\min\{b_{\eta,P}/2,d_{\rm in}/(4L_{\rm raw})\}\)
and \(|\alpha_0|<d_{\rm in}/4\). The height estimate
\eqref{eq:height-capture} and Lemma~\ref{lem:loc-capture} admit every
sufficiently close whole cycle to this counted receiver, at its unique
prime section hit. Lemma~\ref{lem:slice-bound-transfer}, with
\(\mathcal B=\{\alpha_0<0,\eta\le p\le P\}\) and
\(N_{\rm cyc}=2\), gives the full-germ neighborhood and collar.
\end{proof}

\begin{corollary}[Three-cycle concentration]
\label{cor:three-cycle-concentration}
Let \(F_n\to F_*\) in the full quadratic coefficient germ. Suppose
each \(F_n\) has three distinct isolated periodic orbits
\(C_n^1,C_n^2,C_n^3\), all of this same field, and
\[
 \max_{1\le j\le3}
 d_{\rm H}(\iota_*(C_n^j),\Gamma_{I4})\longrightarrow0.
\]
Then \(\alpha_0(F_n)<0\) and \(\alpha_1(F_n)>0\) for all
sufficiently large \(n\), and
\[
 p(F_n)=\frac{\alpha_1(F_n)}{\sqrt{-\alpha_0(F_n)}}
 \longrightarrow0.
\]
\end{corollary}

\begin{proof}
The whole-coefficient bounds of
Proposition~\ref{sharp:upper-coefficient-bounds}, on the common
receiver followed by geometric capture, give at most two cycles when
\(\alpha_0\ge0\), at most one when
\(\alpha_0<0,\alpha_1\le0\), and at most two when
\(\alpha_0<0,p\ge P_*\), for one fixed \(P_*<\infty\).
In the last range the exact trace scaling is
\(r_L=\alpha_1\), \(p_L=1\), \(a_L=-p^{-2}\);
its independent small \(r_L,|\gamma|\) bounds and positive unscaled
cap cover the entire large-ratio sector.
The same-field transport used above puts these exclusions on one full
coefficient neighborhood and collar. The sequence and all three cycle
sequences eventually enter them. Hence eventually
\(\alpha_0(F_n)<0\), \(\alpha_1(F_n)>0\), and
\(0<p(F_n)<P_*\).

Fix \(0<\eta<P_*\). Apply
Proposition~\ref{prop:compact-positive-two-cycle-exclusion} to
\([\eta,2P_*]\). Its neighborhood and collar are fixed for this
\(\eta\), so the sequence and its three cycles eventually enter
them. If \(p(F_n)\ge\eta\) there, all three cycles belong to
one two-cycle sector, a contradiction. Thus \(p(F_n)<\eta\)
eventually. This proves the limit. The argument uses one fixed
\(\eta\) at a time and requires no common collar as
\(\eta\downarrow0\).
\end{proof}

\begin{corollary}[A consequence in a positive-gamma subregion]
\label{cor:positive-gamma-concentration}
Along any subsequence in Corollary~\ref{cor:three-cycle-concentration}
for which \(\gamma(F_n)>0\), if
\[
 \frac{s(F_n)}{\gamma(F_n)^{3/2}}\longrightarrow0,
\]
then
\[
 \liminf\frac{p(F_n)}{\sqrt{\gamma(F_n)}}\ge\sqrt{12}.
\]
\end{corollary}

\begin{proof}
Fix \(0<\pi_0<\sqrt{12}\). The second alternative of
Lemma~\ref{deg:lem-raw-jacobian} supplies a fixed \(k_0>0\)
and a uniformly contracting complete return when
\[
 \frac{s}{\gamma^{3/2}}\le k_0,\qquad
 \frac{p}{\sqrt\gamma}\le\pi_0.
\]
Its fixed-window endpoint quotient bounds the derivative on the
whole physical interval, including intervening inputs. After its
coefficient and multiplier restriction, geometric capture gives one
neighborhood and collar. The subsequence eventually satisfies the
first inequality and enters both; three isolated cycles exclude the
second. Thus its lower limit is at least every
\(\pi_0<\sqrt{12}\), which proves the assertion.
\end{proof}

\section{Three hyperbolic cycles approaching the marked graphic}
\label{sharp:lower-section}

The complete organizing gain tends to zero as \(c\downarrow1\).
We first construct three hyperbolic cycles at fixed \(c>1\) by
balancing this gain against the central transition slope, then pass
diagonally to \(c=1\) and transfer the cycles to the original marked
graphic.

\begin{theorem}\label{sharp:lower-theorem}
For every neighborhood \(\mathcal V\) of \(F_*\) in the full space
\(\mathcal Q_2\) of real quadratic vector fields and every
\(\delta>0\), there is \(F\in\mathcal V\) with three geometrically
distinct hyperbolic periodic orbits \(\gamma_1,\gamma_2,\gamma_3\)
satisfying
\[
 d_{\rm H}\bigl(\iota_*(\gamma_j),\Gamma_{I4}\bigr)<\delta,
 \qquad j=1,2,3.
\]
The cycles follow both marked branches in the positive-time order
\(q_*\to p_*\to q_*\). In particular,
\(\operatorname{Cycl}_{\mathcal Q_2}(\Gamma_{I4},F_*)\ge3\).
\end{theorem}

Use \(H_c,\Gamma_c\) from \eqref{intro:eq-i21-family};
\(H_1=H_0\) and \(\Gamma_1=\Gamma\).
For \(c>1\), this organizing family appears in
\cite[(1.1)--(1.3), p.~2]{Lu26I21}, with its variables \((x,y,C_0)\)
written here as \((X,Y,c)\).

All coefficient norms below are taken in a fixed norm on the twelve
quadratic coefficients, and \(d_{\rm H}\) uses the fixed smooth metric
of this paper.

\subsection{A construction at a fixed limiting parameter}

Use \eqref{eq:slice}, the organizing family \eqref{intro:eq-i21-family},
and its positive speed \(f_c\) in \eqref{intro:eq-parabola-speed}.
Both infinite ends have endpoint \(q=(0,1,0)\).

\begin{theorem}\label{p4:fixed-c}
There is $K_0>0$ such that, for each fixed $c\in(1,1+K_0)$, there is
a continuous family $F_{c,\rho}$ of fields \eqref{eq:slice},
defined for all sufficiently small $\rho>0$, with
$F_{c,\rho}\to H_c$ in quadratic coefficients.
Each $F_{c,\rho}$ has three geometrically distinct hyperbolic periodic
orbits $\gamma_{j,c,\rho}$, $1\le j\le3$, and
\[
 \max_{1\le j\le3}d_{\rm H}(\iota_0(\gamma_{j,c,\rho}),\Gamma_c)
 \longrightarrow0\qquad(\rho\downarrow0).
\]
The coefficient $C_{c,\rho}$ converges to $c$.
\end{theorem}

All estimates below are uniform on the indicated compact phase and
coefficient sets. After $c$ is fixed, their constants may depend on $c$.
Every phase derivative keeps the coefficients fixed.

\subsubsection{A cubic action map with three transverse roots}

Set $k=3/2$ and
\[
 P_\eps(\xi)=(\xi+1)(\xi-\eps)(\xi-k\eps).
\]
For $0<r<1$ let $u_\eps(r)\in(0,\eps)$ be determined by
\begin{equation}\label{p4:action}
 \int_0^{u_\eps(r)}\frac{\xi}{P_\eps(\xi)}\dd\xi
 =\int_0^r\frac{s}{P_\eps(-s)}\dd s.
\end{equation}
Both integrals increase strictly and the left diverges at $u=\eps$,
giving a unique smooth $u_\eps$ with
\begin{equation}\label{p4:uprime}
 u_\eps'(r)=
 \frac{rP_\eps(u_\eps(r))}{u_\eps(r)P_\eps(-r)}>0.
\end{equation}
Define $S_\eps$ by
\[
 b_\eps(r)=-u_\eps(r)^2/2,\qquad
 S_\eps(b_\eps(r))=-r^2/2.
\]
On compact subintervals of $0<r<1$ this defines a smooth increasing
map in the phase $b$.

\begin{lemma}\label{p4:model}
There are a compact interval $J\Subset(0,1)$, three disjoint closed
subintervals $U_j\subset\operatorname{int}J$, and $\eps_*>0$ with
the following properties. Put
\[
 m_\eps=2\eps^{5/2},\qquad
 c_\eps=\frac{\eps^2}{2}-\frac{101}{100}\eps^{5/2}.
\]
For $0<\eps<\eps_*$ the function
\[
 Q_\eps(r)=\eps^{-5/2}
 \{c_\eps+b_\eps(r)-m_\eps S_\eps(b_\eps(r))\}
\]
has one simple zero in each $U_j$. Its endpoint signs and derivative
lower bounds on these intervals can be chosen independently of $\eps$.
Moreover,
\begin{equation}\label{p4:model-error}
 \left\lVert Q_\eps-\Phi\right\rVert_{C^1(J)}
 \le C_J(\sqrt\eps+\eps|\log\eps|),\qquad
 \Phi(r)=\sqrt{\frac{1-r}{8r}}+r^2-\frac{101}{100}.
\end{equation}
\end{lemma}

\phantomsection\label{app:lower-technical-proofs}
\phantomsection\label{app:lower-model-proof}

\begin{proof}[Proof of Lemma~\ref{p4:model}]
The partial fractions of $\xi/P_\eps(\xi)$ have residues
\[
 a_N=-\frac1{(1+\eps)(1+k\eps)},\quad
 a_P=-\frac2{1+\eps},\quad
 a_B=\frac3{1+k\eps},\quad a_N+a_P+a_B=0.
\]
The action difference in \eqref{p4:action} is therefore
\begin{equation}\label{p4:action-difference}
 a_N\log\frac{1+u}{1-r}
 +a_P\log\frac{\eps-u}{r+\eps}
 +a_B\log\frac{k\eps-u}{r+k\eps}.
\end{equation}
Put $\eps=h^2$, $u=h^2-dh^3$ in \eqref{p4:action-difference}.
The coefficient of $\log h$ is exactly
$3a_P+2a_B=-3h^2/((1+h^2)(1+kh^2))$; all remaining terms extend
smoothly to $h=0$, with limit $\log((1-r)/(8r))-2\log d$.
On positive compact $(d,r)$ rectangles the full $C^1_{d,r}$ error is
$O(h+h^2|\log h|)$. The limiting $d$ derivative $-2/d$ and the
implicit function theorem give
\begin{equation}\label{p4:d-error}
 d_\eps(r)=\frac{\eps-u_\eps(r)}{\eps^{3/2}},\qquad
 \left\lVert d_\eps-d_0\right\rVert_{C^1(J)}
 \le C_J(\sqrt\eps+\eps|\log\eps|),\quad
 d_0(r)=\sqrt{\frac{1-r}{8r}}.
\end{equation}
Uniqueness in \eqref{p4:action} identifies this local solution with
$u_\eps$.

The center terms cancel exactly:
\begin{equation}\label{p4:center-cancellation}
 Q_\eps(r)=d_\eps(r)-\frac{\sqrt\eps}{2}d_\eps(r)^2
               +r^2-\frac{101}{100}.
\end{equation}
This and \eqref{p4:d-error} prove \eqref{p4:model-error}.

The values at $r=1/20,1/2,49/50,999/1000$ have signs $+,-,+,-$,
as follows by evaluating
\[
 \frac{1-r}{8r}-\left(\frac{101}{100}-r^2\right)^2
 =-\frac{10000r^5-20200r^3+11451r-1250}{10000r}.
\]
Furthermore
\[
 \Phi'(r)=2r-\frac1{4\sqrt2\,r^{3/2}\sqrt{1-r}}.
\]
Its zeros are exactly the solutions of $r^5(1-r)=1/128$.
The left side increases on $(0,5/6)$ and decreases on $(5/6,1)$;
its maximum is $3125/46656>1/128$.
Thus $\Phi'$ has exactly two zeros. The alternating signs and Rolle's
theorem give exactly three roots of $\Phi$, all simple: a nonsimple
root and the two intervening Rolle points would give a third root
of $\Phi'$.
Choose $J$ containing the four test points and choose disjoint
neighborhoods $U_j$ of the three roots, indexed in increasing $r$ order,
with opposite endpoint
signs and $|\Phi'|\ge2\kappa>0$ on their union.
Reduce $\eps_*$ so that \eqref{p4:model-error} preserves these
signs and gives $|Q_\eps'|\ge\kappa$ there.
\end{proof}

Figure~\ref{fig:lower-action-profile} shows the model profile and
resolves its two roots near \(r=1\).
% Computed scalar profile used by the lower action construction.
% Native TikZ; the preamble loads tikz (no PGFPlots or external data).
% Curves sample Phi(r)=sqrt((1-r)/(8r))+r^2-101/100 directly.
% Main axes: 0.05 <= r <= 1.02, -0.45 <= Phi <= 0.15;
% the graph itself is sampled only for 0.05 <= r < 1.
% Zoom window: 0.965 <= r < 1, -0.012 <= Phi <= 0.007.
% Sampling uses a unit-interval parameter; straight segments avoid spline
% overshoot. PGF floating-point evaluation resolves the shallow right peak.
% Root markers are numerical locations (bisection residual < 5e-15):
% 0.1115973705382644, 0.9786836116269868, 0.9987531118881820.
\begin{figure}[htb]
\centering
\usepgflibrary{fpu}
\pgfmathdeclarefunction{loweractionphi}{1}{%
  \begingroup
  \pgfkeys{/pgf/fpu=true,/pgf/fpu/output format=fixed}%
  \pgfmathparse{sqrt((1-#1)/(8*#1))+#1*#1-1.01}%
  \pgfmathsmuggle\pgfmathresult
  \endgroup
}
\pgfmathdeclarefunction{loweractionphidelta}{1}{%
  \begingroup
  \pgfkeys{/pgf/fpu=true,/pgf/fpu/output format=fixed}%
  \pgfmathparse{sqrt(#1/(8*(1-#1)))-2*#1+#1*#1-.01}%
  \pgfmathsmuggle\pgfmathresult
  \endgroup
}
\begin{tikzpicture}[x=1cm,y=1cm,font=\small]
  % Panel (a): the full profile, with the divergent left tail clipped.
  \begin{scope}
    \draw[line width=.3pt] (0,0) rectangle (7.5,3.5);
    \draw[densely dashed,line width=.35pt] (0,2.625)--(7.5,2.625);
    \foreach \r/\lab in {.1/{0.1},.3/{0.3},.5/{0.5},.7/{0.7},.9/{0.9},1/{1}} {
      \pgfmathsetmacro{\xx}{(\r-.05)/.97*7.5}
      \draw[line width=.3pt] (\xx,0)--(\xx,-.065);
      \node[below,inner sep=2pt] at (\xx,-.065) {$\lab$};
    }
    \foreach \v/\lab in {-.4/{-0.4},-.2/{-0.2},0/{0},.1/{0.1}} {
      \pgfmathsetmacro{\yy}{(\v+.45)/.60*3.5}
      \draw[line width=.3pt] (0,\yy)--(-.065,\yy);
      \node[left,inner sep=2pt] at (-.065,\yy) {$\lab$};
    }
    \begin{scope}
      \clip (0,0) rectangle (7.5,3.5);
      \draw[line width=.75pt]
        plot[domain=0:1,samples=260,variable=\u]
        ({.90*\u/.97*7.5},
         {(loweractionphi(.05+.90*\u)+.45)/.60*3.5})
        plot[domain=0:1,samples=320,variable=\u]
        ({(.90+.04999*\u)/.97*7.5},
         {(loweractionphidelta(.05-.04999*\u)+.45)/.60*3.5});
      \foreach \r in {.1115973705382644,.9786836116269868,.9987531118881820} {
        \fill ({(\r-.05)/.97*7.5},2.625) circle[radius=1.25pt];
      }
    \end{scope}
    \node[above,inner sep=2pt] at (0,3.5) {$\Phi(r)$};
    \node[below,inner sep=2pt] at (7.5,-.42) {$r$};
    \node at (3.75,-.74) {(a) The action profile};
  \end{scope}
  % Panel (b): a separate linear scale resolves both right-hand zeros.
  \begin{scope}[shift={(9.3,0)}]
    \draw[line width=.3pt] (0,0) rectangle (5.2,3.5);
    \draw[densely dashed,line width=.35pt]
      (0,{.012/.019*3.5})--(5.2,{.012/.019*3.5});
    \foreach \r/\lab in {.97/{0.97},.98/{0.98},.99/{0.99},1/{1}} {
      \pgfmathsetmacro{\xx}{(\r-.965)/.035*5.2}
      \draw[line width=.3pt] (\xx,0)--(\xx,-.065);
      \node[below,inner sep=2pt] at (\xx,-.065) {$\lab$};
    }
    \foreach \v/\lab in {-.01/{-0.010},-.005/{-0.005},0/{0},.005/{0.005}} {
      \pgfmathsetmacro{\yy}{(\v+.012)/.019*3.5}
      \draw[line width=.3pt] (0,\yy)--(-.065,\yy);
      \node[left,inner sep=2pt] at (-.065,\yy) {$\lab$};
    }
    \begin{scope}
      \clip (0,0) rectangle (5.2,3.5);
      \draw[line width=.75pt]
        plot[domain=0:1,samples=360,variable=\u]
        ({.03499*\u/.035*5.2},
         {(loweractionphidelta(.035-.03499*\u)+.012)/.019*3.5});
      \foreach \r in {.9786836116269868,.9987531118881820} {
        \fill ({(\r-.965)/.035*5.2},{.012/.019*3.5})
          circle[radius=1.25pt];
      }
    \end{scope}
    \node[above,inner sep=2pt] at (0,3.5) {$\Phi(r)$};
    \node[below,inner sep=2pt] at (5.2,-.42) {$r$};
    \node at (2.6,-.74) {(b) Enlargement near $r=1$};
  \end{scope}
\end{tikzpicture}
\caption{Computed graph of
$\Phi(r)=\sqrt{(1-r)/(8r)}+r^2-101/100$ on $0<r<1$.
Panel~(a) clips the divergent left-end portion as $r\downarrow0$.
The dots mark numerically computed zeros; panel~(b) resolves the two near $r=1$.
The three simple zeros are established analytically in
Lemma~\ref{p4:model}.}
\label{fig:lower-action-profile}
\end{figure}


The lower construction balances the vanishing outer gain
$m_\eps=2\eps^{5/2}$ with the central slope
$S_\eps'(b_\eps(r))\asymp\eps^{-5/2}$ on the chosen phase interval,
as follows from \eqref{p4:d-error}. Their product remains finite and
positive; the calibration below fixes this gain and the return center
in one field.

\subsubsection{The actual passage through infinity}

Choose $K_0>0$ small enough that, for $1<c<1+K_0$,
\begin{equation}\label{p4:epsilon-c}
 \eps=\eps(c)=
 \left(\frac12e^{-4\pi/\sqrt{c-1}}\right)^{2/5}<\eps_*.
\end{equation}
Fix such a $c$ and this $\eps$ for the remainder of the construction.
Set
\begin{equation}\label{p4:path}
\begin{gathered}
 \beta=1-(1+k)\eps,\qquad \gamma=\beta\rho/3,\\
 \alpha_0=-k\eps^2\rho^3,\qquad
 \alpha_1=((1+k)\eps-k\eps^2)\rho^2,\\
 \mu_3=-\gamma-\gamma^2,\quad
 \mu_2=\alpha_1+3\gamma^2,\quad
 \mu_1=\alpha_0-\alpha_1\gamma-2\gamma^3.
\end{gathered}
\end{equation}
The coefficients $\mu$ are fixed by $c,\rho$; $A,C$ will be chosen
near $(1,c)$. In the physical field \eqref{eq:raw}, substitute
$v=\gamma+\rho\xi$, $z=\rho^2\eta$ and the positive time
$t_{\rm new}=\rho t_{\rm raw}$, giving exactly
\begin{equation}\label{p4:central}
\begin{split}
 \xi'&=-\rho P_\eps(\xi)-\ell_\rho(\xi)\eta,\\
 \eta'&=-\eta\{\xi B_\rho(\xi)+\zeta_\rho(\xi)\eta\},\\
 \ell_\rho&=1-A(\gamma+\rho\xi)+C(\gamma+\rho\xi)^2,\\
 B_\rho&=1+2\gamma+\rho\xi,\qquad
 \zeta_\rho=C\rho(\gamma+\rho\xi).
\end{split}
\end{equation}

Enlarge $J\Subset(0,1)$; its entry and exit feet remain separated from
$0$ and the three zeros of $P_\eps$. Choose compact phase intervals $I_+,I_-$
containing their images $-u^2/2,-r^2/2$, with
$b_\eps(J)\Subset\operatorname{int}I_+$ and
$S_\eps(I_+)\Subset\operatorname{int}I_-$.
Choose $L$ large enough that
\begin{equation}\label{p4:L-window}
 3/8<1/2+b/L^2<5/8
\end{equation}
on both phase intervals and on intervals connecting them to $b=0$.
On $\xi=\pm L$ use the same literal coordinate $b=\eta-L^2/2$.
The tail reference uses $b=0$.

Use the same lip height $h_\rho=(k\eps^2\rho)^{12}$ on both sides.
Between the lips use $Q=\rho\log(\eta/h_\rho)$.

\begin{lemma}\label{p4:core}
On the enlarged input interval $I_+$, the positive-time passage
of \eqref{p4:central} from $\xi=L$ to $\xi=-L$ exists for small
$\rho>0$. In the literal phase $b$ its map satisfies
\begin{equation}\label{p4:core-limit}
 \|S_{\eps,\rho;A,C}-S_\eps\|_{C^1_b}
 \le C_\eps\rho(1+|\log\rho|)\longrightarrow0,
\end{equation}
uniformly for $A,C$ in a sufficiently small fixed rectangle about
$(1,c)$, and the entire passage has $\xi'<0$.
\end{lemma}

\phantomsection\label{app:lower-core-proof}

\begin{proof}[Proof of Lemma~\ref{p4:core}]
Use the same lip height $h_\rho=(k\eps^2\rho)^{12}$ on both sides.
On each outer part put $H=\eta-\xi^2/2$ and use $\log\eta$ as a clock.
Direct differentiation gives
\begin{equation}\label{p4:fast-clock}
 \frac{\dd H}{\dd\log\eta}
 =-\frac{\rho\xi P_\eps+\xi\eta(\ell_\rho-B_\rho)
                    -\zeta_\rho\eta^2}
             {\xi B_\rho+\zeta_\rho\eta}.
\end{equation}
For the selected negative $H$ windows,
$\xi=\pm\sqrt{2(\eta-H)}$ stays a positive distance from zero.
On their fixed enlarged neighborhoods the denominator in
\eqref{p4:fast-clock} has the appropriate nonzero sign, while its
right side and the mixed derivatives needed for a first phase
variation are $O_\eps(\rho)$. Its clock length is
$O_\eps(1+|\log\rho|)$.
Integration and the scalar variational equation therefore give
\begin{equation}\label{p4:fast-error}
 |H_{\rm out}-H_{\rm in}|+
 |\partial_{H_{\rm in}}H_{\rm out}-1|
 \le C_\eps\rho(1+|\log\rho|).
\end{equation}
Applying this estimate up to a first exit from the enlarged windows
excludes such an exit for small $\rho$.

Starting at a positive lip foot $u$ and integrating outward in
reverse time gives at $\xi=L$
$b=-u^2/2+O_\eps(\rho(1+|\log\rho|))$ in $C^1_u$.
Its derivative is $-u+o(1)$, bounded away from zero.
Thus it can be inverted on the entire enlarged entry interval.
The negative lip gives in positive time
$b=-r^2/2+O_\eps(\rho(1+|\log\rho|))$ in $C^1_r$.
Near each foot, $P_\eps>0$. Away from its small foot neighborhood,
the limiting height $(\xi^2-u^2)/2$ or $(\xi^2-r^2)/2$
has a positive lower bound. This lower bound dominates the bounded
term $\rho P_\eps$, so $\rho P_\eps+\ell_\rho\eta>0$
throughout both fast parts. The hits of $\xi=\pm L$ are unique
and transverse.

Between the lips use $Q=\rho\log(\eta/h_\rho)$. In the region
$Q\le0$ the exact graph equation is
\begin{equation}\label{p4:thin}
 Q_\xi=\frac{\xi B_\rho+\zeta_\rho\eta}
                 {P_\eps+\ell_\rho\eta/\rho},
 \qquad \eta=h_\rho e^{Q/\rho}.
\end{equation}
The compact traversed $\xi$ interval, with its exit buffer, has
$P_\eps\ge p_\eps>0$.
Replace $e^{Q/\rho}$ for $Q>0$ by a bounded smooth function of
$Q/\rho$ equal to the exponential for $Q\le0$.
This defines an open scalar ODE domain and changes no solution
before its first return to the lip.
For total derivative order at most three in $Q,\xi,A,C$, the
right side differs from $\xi/P_\eps$ by at most
\begin{equation}\label{p4:thin-error}
 C_\eps\{\rho+\rho^{11-j}\},\qquad 0\le j\le3,
\end{equation}
where $j$ is the number of $Q$ derivatives. Indeed,
$h_\rho/\rho=O_\eps(\rho^{11})$, and each $Q$ derivative
costs one power of $\rho$.
The scalar integral equation and its first variational equation converge
on the fixed $\xi$ interval.

The limiting solution starting at $u$ is
$Q_0(\xi)=\int_u^\xi s/P_\eps(s)\dd s$.
It is strictly negative before its negative exit $-r$, except at
the entrance, and $Q_0'(-r)=-r/P_\eps(-r)<0$.
Uniform signs of $Q_\xi$ near entrance and exit, and a strictly
negative upper bound for $Q_0$ between them, exclude an early first
return and persist under \eqref{p4:thin-error}.
Between two buffered negative exit cuts with opposite limiting signs,
the implicit function theorem gives a unique first lip hit. The thin
field and first variations have $O_\eps(\rho)$ error and the exit has
a fixed inverse bound. Composition with \eqref{p4:fast-error} proves
\eqref{p4:core-limit}; the denominator signs and $P_\eps>0$ on the
thin part give $\xi'<0$ everywhere.
\end{proof}


\subsubsection{The two parabolic tails}

Choose a small raw half-width $d>0$. Write $q=|v-\gamma|$,
$v=\gamma+\sigma q$ with $\sigma=\pm1$, and
$\widetilde P_\rho(x)=\rho^3P_\eps(x/\rho)$.
The exact radial graph equation is
\begin{equation}\label{p4:tail}
 z_q=F_{\sigma,\rho}(q,z)=
 \frac{z\{q(1+2\gamma+\sigma q)
                  +\sigma C(\gamma+\sigma q)z\}}
 {\widetilde P_\rho(\sigma q)+
       [1-A(\gamma+\sigma q)+C(\gamma+\sigma q)^2]z}.
\end{equation}
Its initial condition at $q=\rho L$ is
$z=\rho^2(L^2/2+b)$.
The negative tail is traversed outward in positive time; the positive
tail is constructed outward in reverse time.

\begin{lemma}\label{p4:tails}
For the enlarged phase windows the tails reach the raw ports
$v=\gamma\pm d$ uniquely. Their maps $T_{\sigma,\rho}(b)$ have the
decomposition
\begin{equation}\label{p4:tail-decomposition}
 T_{\sigma,\rho}(b)=a_{\sigma,\rho}
             +c_{\sigma,\rho}\phi_{\sigma,\rho}(b),\qquad
 a_{\sigma,\rho}=T_{\sigma,\rho}(0),\quad
 c_{\sigma,\rho}=\partial_bT_{\sigma,\rho}(0)>0.
\end{equation}
Uniformly on a fixed small $A,C$ rectangle,
\begin{equation}\label{p4:tail-limits}
\begin{gathered}
 a_{\sigma,\rho}\to a_{\sigma,0},\qquad
 c_{\sigma,\rho}/\rho^2\to\ell_\sigma>0
                  \quad\text{in }C^1_{A,C},\\
 \phi_{\sigma,\rho}\to\operatorname{id}
                  \quad\text{in }C^1_b.
\end{gathered}
\end{equation}
On every fixed $q$ interval bounded away from zero, the entire
tail family converges to its limiting parabolic graph.
\end{lemma}

\phantomsection\label{app:lower-tails-proof}

\begin{proof}[Proof of Lemma~\ref{p4:tails}]
Since $\gamma=O_\eps(\rho)$ and $q\ge\rho L$,
\[
 \widetilde P_\rho(\sigma q)=O_\eps(q^3+\rho q^2+\rho^2q+\rho^3),
 \qquad |\widetilde P_\rho|/q^2\le C_\eps q.
\]
For $w=z/q^2$ in $[1/4,3/4]$ the denominator in
\eqref{p4:tail} is at least a positive constant times $q^2$
after $d$ is reduced. Its logarithmic-radius equation is
$w_{\log q}=1-2w+O_\eps(q)$, with the same bound for its
first $w$ derivative. The boundary vectors point inward.
Equation \eqref{p4:L-window} starts all phases strictly inside
this band. The band is invariant, the radial speed does not vanish,
and the raw port is reached once.

Writing the right side as
$z(\beta+\zeta z)/(\widetilde P+\ell z)$, where
$\widetilde P=\widetilde P_\rho(\sigma q)$, direct differentiation gives
\[
 F_{zz}=\frac{2\widetilde P(\zeta\widetilde P-\beta\ell)}
                 {(\widetilde P+\ell z)^3}.
 \]
Together with the band bounds this yields
\begin{equation}\label{p4:tail-majorants}
\begin{split}
 |F_z|&\le C_\eps(1+\rho/q+\rho^2/q^2+\rho^3/q^3),\\
 |F_{zz}|&\le C_\eps q^{-2},\\
 |F_A|+|F_C|&\le C_\eps q^2,\qquad
 |F_{zA}|+|F_{zC}|\le C_\eps .
\end{split}
\end{equation}
The first majorant is uniformly bounded since $\rho/q\le1/L$.
For $J=\partial_bz$, $J(\rho L)=\rho^2$ and $J_q=F_zJ$.
Consequently
\begin{equation}\label{p4:J}
 c_\eps'\rho^2\le J\le C_\eps'\rho^2,\qquad
 |\partial_b\log J|
 \le\int_{\rho L}^d|F_{zz}|J\dd q\le C_\eps'\rho.
\end{equation}
This proves the last assertion of \eqref{p4:tail-limits}, since
$\phi(0)=0$ and $\phi'=J(q,b)/J(q,0)$.

At $\rho=0$ equation
\eqref{p4:tail} has a unique graph in the band with
$z_0(q)=q^2/2+O(q^3)$.
It is obtained by starting $z(a)=a^2/2$ and letting $a\downarrow0$.
The inward-band comparison gives $|z/q^2-1/2|\le Cq$
for these initial data. Any two bounded-band limiting graphs have
distance at most $Ca^2$ at $q=a$; the bound for $F_z$
propagates this distance to $q=d$, proving uniqueness.
For the actual reference $b=0$, coefficient convergence on compact
$q>0$ intervals, $z\le Cq^2$, and propagation from a small fixed $a$
give convergence to this graph.

The parameter-independent initial height gives
$z_{A,q}=(\partial_zF)z_A+F_A$ and the analogous equation for $C$.
The majorants \eqref{p4:tail-majorants} give $|z_A|+|z_C|\le Cq^3$;
their integral equations converge to the unique zero-initial-value
equations at $q=0$, proving $C^1_{A,C}$ convergence of $a_{\sigma,\rho}$.
Finally,
\[
 \log(c_{\sigma,\rho}/\rho^2)=
 \int_{\rho L}^dF_z(q,z(q,0))\dd q.
\]
Its $A,C$ derivatives have integrands
$F_{zA}+F_{zz}z_A$ and $F_{zC}+F_{zz}z_C$, bounded by
$C(1+q)$.
Pointwise convergence for $q>0$ and these integrable majorants give
the first-slope limits in \eqref{p4:tail-limits}.
The same bounds show $\ell_\sigma$ and their reciprocals are bounded
on the coefficient rectangle.
\end{proof}


\subsubsection{The finite connection and two coefficient directions}

At $\rho=0,A=1,C=c$ the two limiting tail graphs are the ends of
the organizing parabola, whose raw height is
\begin{equation}\label{p4:raw-organizing}
 z_c(v)=\frac{v^2}{1+\sqrt{1+cv^2}}.
\end{equation}
The finite arc between $v=-d$ and $v=d$ is a compact affine arc.
It includes the two crossings $Y=0,\ X=\pm\sqrt c$.
At these points
\[
 \dot X=c\pm\sqrt c>0,\qquad
 \dot Y=c\pm c^{3/2}\ne0.
\]
Along the entire arc $f_c\ge(c-1)/2$.
Choose a fixed narrow affine tube and buffered raw section windows.
After the coefficient rectangle and $\rho$ are reduced, the
finite-time flow stays in this tube and hits the positive port
transversely. It defines an increasing map $E_\rho$ from negative
to positive raw height. Its mixed derivatives through order two
converge uniformly to those of $E_0$.
The same transverse finite-time estimates give a uniform $C^3$ norm
and slope bounded away from zero.
Lemma~\ref{p4:tails} admits the port heights and all phase images to
these buffered domains before composition.

Define
\begin{equation}\label{p4:outer-data}
 \Delta_\rho=a_{+,\rho}-E_\rho(a_{-,\rho}),\qquad
 \Lambda_\rho=E_\rho'(a_{-,\rho})c_{-,\rho}/c_{+,\rho}.
\end{equation}
These functions converge in $C^1_{A,C}$ by the finite-time
estimates and the normalized slope limits in \eqref{p4:tail-limits}.
The limiting return center varies nontrivially with $A$, while $C$
changes the complete gain along the connected family $A=1$.
The triangular Jacobian below allows both conditions to be imposed
on the same field.

\begin{lemma}\label{p4:transverse-coefficients}
At $(A,C)=(1,c)$,
\begin{equation}\label{p4:Jacobian}
 \Delta_0=0,\quad \partial_C\Delta_0=0,\quad
 \partial_A\Delta_0>0,\qquad
 \Lambda_0=e^{-4\pi/\sqrt{c-1}},\quad
 \partial_C\Lambda_0=2\pi(c-1)^{-3/2}\Lambda_0>0.
\end{equation}
\end{lemma}

\phantomsection\label{app:lower-coefficients-proof}

\begin{proof}[Proof of Lemma~\ref{p4:transverse-coefficients}]
For $A=1$ the whole $C>1$ family has the invariant parabola, so
$\Delta_0(1,C)=0$.
Write $K=c-1$ and
\[
 \mathcal E_c(X)=((X+1)^2+K)^2
 \exp\left[-\frac4{\sqrt K}
                \arctan\frac{X+1}{\sqrt K}\right].
\]
For $W=Y-(X^2-c)/2$, the exact affine equation at $\mu=0,C=c$ is
\[
 \dot W=2XW+(1-A)X^2,\qquad
 \mathcal E_c'/\mathcal E_c=2X/f_c.
\]
The graph variation in the $A$ direction therefore satisfies
\[
 W_{A,X}=(2X/f_c)W_A-X^2/f_c.
\]
The tail bound $z_A=O(q^3)$ gives $W_A=O(X)$ at infinity.
Since $\mathcal E_c\asymp X^4$, the free homogeneous coefficient
vanishes, and the selected variations are
\[
 W_A^-(X)=-\mathcal E_c(X)\int_{-\infty}^X
       \frac{t^2}{f_c(t)\mathcal E_c(t)}\dd t,\qquad
 W_A^+(X)=\mathcal E_c(X)\int_X^\infty
       \frac{t^2}{f_c(t)\mathcal E_c(t)}\dd t.
\]
The positive integrand is $O(|t|^{-4})$ at both ends, so
$W_A^+-W_A^-$ is positive at the positive port. At fixed $v=d$,
\[
 W_z=\frac{d^2-z_c(d)}{z_c(d)^3}>0,
\]
conversion to raw height gives $\partial_A\Delta_0>0$.

Truncate the limiting tails at $v=\pm a$. The outer derivative is
$E_0'$ times the negative-tail first variation divided by the positive
one, tending to $\Lambda_0=E_0'\ell_-/\ell_+$ as $a\downarrow0$.
The parabola endpoints are $X=\pm X_a$, $X_a=a/z_c(a)\to\infty$,
with equal fixed-$v$ conversion factors $W_z$. Thus the derivative
is $\mathcal E_c(X_a)/\mathcal E_c(-X_a)$: its polynomial ratio tends
to one and its arctangent difference to $\pi$, giving
$\Lambda_0(1,c)=e^{-4\pi/\sqrt K}$. For every $C$ in the rectangle,
$\Lambda_0(1,C)=e^{-4\pi/\sqrt{C-1}}$, so differentiation gives
$\partial_C\Lambda_0(1,c)=2\pi K^{-3/2}e^{-4\pi/\sqrt K}$.
\end{proof}


\subsubsection{A single calibration and three simple actual roots}

Equation \eqref{p4:epsilon-c} gives
$\Lambda_0(1,c)=m_\eps$ exactly.
Consider the two equations
\begin{equation}\label{p4:calibration}
 G_\rho(A,C)=
 \bigl(\Delta_\rho-c_{+,\rho}c_\eps,\,
                         \Lambda_\rho-m_\eps\bigr)=0.
\end{equation}
By \eqref{p4:tail-limits},
$c_{+,\rho}\to0$ together with its parameter derivatives.
Thus $G_\rho\to G_0$ in $C^1_{A,C}$ and $G_0(1,c)=0$.
Lemma~\ref{p4:transverse-coefficients} gives
\[
 \det DG_0(1,c)
 =(\partial_A\Delta_0)(\partial_C\Lambda_0)>0.
\]
On a small fixed ball in the previous buffered domains with $C>1$,
$C^1$ convergence makes
$(A,C)\mapsto(A,C)-DG_0(1,c)^{-1}G_\rho(A,C)$ a contraction into
itself for small $\rho$. The resulting coefficient pair satisfies
\begin{equation}\label{p4:AC}
 (A_\rho,C_\rho)\longrightarrow(1,c).
\end{equation}
The calibrated pair is continuous in $\rho>0$, since all actual hits
are transverse.
No estimate uniform in $c\downarrow1$ is needed for this inverse.

Evaluate all maps at this one coefficient pair. Normalize
$E_\rho\circ T_{-,\rho}$ by its exact center and slope:
\[
 E_\rho(T_{-,\rho}(b))=
 E_\rho(a_{-,\rho})+
 E_\rho'(a_{-,\rho})c_{-,\rho}\widehat\phi_{-,\rho}(b).
\]
Here $\widehat\phi_{-,\rho}=N_{c_{-,\rho}}\circ\phi_{-,\rho}$, with
$\psi=E_\rho$ and $z_0=a_{-,\rho}$ in
\eqref{phys:eq-observer-normalization}. The buffered $C^3$ bounds,
$c_{-,\rho}=O(\rho^2)$, and \eqref{p4:tail-limits} give
$\widehat\phi_{-,\rho}(0)=0$ and $C^1_b$ convergence to the identity.

The return equation is
$T_{+,\rho}(b)=E_\rho(T_{-,\rho}(S_{\eps,\rho}(b)))$.
Dividing its difference by the positive $c_{+,\rho}$ and using
\eqref{p4:calibration} gives exactly
\begin{equation}\label{p4:actual-displacement}
 D_\rho(b)=c_\eps+\phi_{+,\rho}(b)
         -m_\eps\widehat\phi_{-,\rho}(S_{\eps,\rho}(b)).
\end{equation}
For $D_\eps=c_\eps+b-m_\eps S_\eps(b)$, the core bound,
\eqref{p4:J}, and the buffered finite connector give, on the common
enlarged interval and uniformly in the coefficient ball,
\begin{equation}\label{p4:normalized-error}
 \left\lVert D_\rho-D_\eps\right\rVert_{C^1_b}
       \le C_\eps\rho(1+|\log\rho|).
\end{equation}
This remains true at \eqref{p4:AC}. Keep $\eps=\eps(c)$ fixed and put
$\mathcal Q_\rho(r)=\eps^{-5/2}D_\rho(b_\eps(r))$.
The chain rule gives the explicit estimate
\begin{equation}\label{p4:actual-error}
 \left\lVert \mathcal Q_\rho-Q_\eps\right\rVert_{C^1(J)}
 \le \eps^{-5/2}\max(1,\left\lVert b_\eps'\right\rVert_{C^0(J)})
                         \left\lVert D_\rho-D_\eps\right\rVert_{C^1_b}
 \longrightarrow0.
\end{equation}
All coefficients, including $A_\rho,C_\rho$, are fixed in this
$r$ derivative. For the fixed $\eps$, choose the same $\rho$
small enough to preserve the three endpoint sign pairs and the
derivative lower bound of Lemma~\ref{p4:model}.
There is exactly one root $r_{j,\rho}$ in each $U_j$, and
$|\mathcal Q_\rho'(r_{j,\rho})|\ge\kappa/2$.

Each root closes one actual circuit. The central passage and both
tails have strictly decreasing $v$ in positive time.
For each $\rho>0$ all flight times are finite: on the thin part
$-\xi'\ge\rho p_\eps$, the fast parts are compact with nonzero
speed, and the outward tail speed is at least a positive constant
times $q^2$ with $q\ge\rho L>0$.
The finite connection lies in a compact affine tube whose spherical
distance from $q$ is positive. The inner section $\xi=L$ tends to
$q$, so for small $\rho$ it has no further hit during the circuit.
Thus the three roots give three different prime periodic orbits.

They are hyperbolic in the actual raw input coordinate.
Indeed $h(r)=T_{+,\rho}(b_\eps(r))$ is positive and
\[
 h_r=c_{+,\rho}\phi_{+,\rho}'b_\eps'(r)<0,\qquad h_r\ne0,
\]
by \eqref{p4:uprime}. If $R_\rho$ is the first return and
$F_\rho(h)=h-R_\rho(h)$, then the exact identity
\[
 F_\rho(h(r))=c_{+,\rho}\eps^{5/2}\mathcal Q_\rho(r)
\]
gives, at each root,
\begin{equation}\label{p4:physical-simple}
 F_{\rho,h}=
 \frac{\eps^{5/2}\mathcal Q_\rho'}
      {\phi_{+,\rho}'b_\eps'}\ne0.
\end{equation}
The common gain $c_{+,\rho}$ cancels; positivity of the first-return
derivative and $R_\rho'\ne1$ prove hyperbolicity.

\subsubsection{Whole-orbit convergence}

The three roots remain in the same compact phase intervals.
On the central passage $|\xi|\le L$ and $0<\eta\le C_{\eps,L}$.
At $\xi=0$ the thin orbit has $Q\le0$, hence all three satisfy
\begin{equation}\label{p4:third-height}
 0<z(\gamma)=\rho^2\eta(0)
 \le\rho^2h_\rho=(k\eps^2)^{12}\rho^{14}\longrightarrow0.
\end{equation}

Fix $0<d_0<d$. On $d_0\le|v-\gamma|\le d$, both tails converge
uniformly to their organizing ends, by Lemma~\ref{p4:tails} and
\eqref{p4:AC}. The finite connection converges uniformly, in affine
coordinates and finite time, to the corresponding organizing arc.
Each of these arcs is parameterized over a full common interval,
with endpoints converging to the corresponding organizing endpoints.
Consequently both directed Hausdorff distances tend to zero on
these pieces.

The remaining tail pieces have $z\le C|v-\gamma|^2$ and
$|v|\le d_0+O(\rho)$.
The central piece has $|v-\gamma|\le\rho L$ and $z\le C\rho^2$.
For these pieces $Y>0$, and the true spherical embedding is
\[
 \iota_0(X,Y)=\frac{(v,1,z)}{\sqrt{1+v^2+z^2}}.
\]
Their distance to $q$ is at most $C(d_0+\rho)$.
The omitted organizing ends have distance at most $Cd_0$ to $q$.
Conversely, \eqref{p4:third-height} provides an actual point tending
to $q$ on each of the three cycles, so every omitted organizing
point is within $C d_0+o(1)$ of that cycle.
First let $\rho\downarrow0$, then $d_0\downarrow0$.
This proves simultaneous whole-orbit Hausdorff convergence and
Theorem~\ref{p4:fixed-c}. The finite part containing $Y=0$
uses the affine embedding throughout.

\begin{corollary}[Stability order]\label{p4:stability-order}
For each fixed $c\in(1,1+K_0)$ and all sufficiently small $\rho>0$,
order the three constructed cycles of $F_{c,\rho}$ by their roots
$r_{1,\rho}<r_{2,\rho}<r_{3,\rho}$. The first and third cycles
are attracting, and the second is repelling.
\end{corollary}
\begin{proof}
The roots of \(\Phi\), in increasing \(r\) order, cross downward,
upward and downward. The two \(C^1\) estimates
\eqref{p4:model-error} and \eqref{p4:actual-error} preserve these
signs for \(\mathcal Q_\rho'\). Since \(b_\eps'<0\) and
\(\phi_{+,\rho}'>0\), identity \eqref{p4:physical-simple} gives
signs \(+,-,+\) for \(1-R_\rho'\). Positivity of the actual
first-return derivative proves the stated stability order.
\end{proof}

\begin{corollary}[Limits of the three return multipliers]
\label{p4:multiplier-limits}
Let $r_1^0<r_2^0<r_3^0$ be the three zeros of $\Phi$ and let
$M_j(c,\rho)>0$ be the derivative of the positive first return at
the cycle $\gamma_{j,c,\rho}$, ordered by increasing $r$.
For $\eps=\eps(c)$, let $r_{j,\eps}$ be the zero of $Q_\eps$ in
$U_j$ and put $\mu_j(c)=m_\eps S_\eps'(b_\eps(r_{j,\eps}))>0$.
Then
\begin{equation}\label{p4:multiplier-successive-limits}
 \begin{gathered}
 \lim_{\rho\downarrow0}M_j(c,\rho)=\mu_j(c),\qquad
 \lim_{c\downarrow1}\mu_j(c)=L(r_j^0),\\
 L(r)=8\sqrt2\,r^{5/2}\sqrt{1-r}.
 \end{gathered}
\end{equation}
The limiting multipliers satisfy
\begin{equation}\label{p4:multiplier-intervals}
 0<L(r_1^0),L(r_3^0)<\tfrac12,
 \qquad \tfrac32<L(r_2^0)<2.
\end{equation}
\end{corollary}
\begin{proof}
For fixed $c$ and $\eps=\eps(c)$, \eqref{p4:actual-error} gives
$r_{j,\rho}\to r_{j,\eps}$. Evaluate all maps at the calibrated
coefficients \eqref{p4:AC}. In the actual raw input coordinate
$h=T_{+,\rho}(b)$, the first return is
\[
 R_\rho=E_\rho\circ T_{-,\rho}\circ
 S_{\eps,\rho}\circ T_{+,\rho}^{-1}.
\]
At $b=b_\eps(r_{j,\rho})$, its derivative is exactly
\begin{equation}\label{p4:multiplier-exact}
 \begin{split}
 M_j(c,\rho)
 &=\frac{E_\rho'(T_{-,\rho}(S_{\eps,\rho}(b)))
          T_{-,\rho}'(S_{\eps,\rho}(b))
          S_{\eps,\rho}'(b)}{T_{+,\rho}'(b)}\\
 &=m_\eps\,
   \frac{\widehat\phi_{-,\rho}'(S_{\eps,\rho}(b))
          S_{\eps,\rho}'(b)}{\phi_{+,\rho}'(b)}>0.
 \end{split}
\end{equation}
The second equality uses the exact calibration
$\Lambda_\rho=m_\eps$. The core and tail $C^1$ limits, including
$\widehat\phi_{-,\rho}\to\operatorname{id}$, give the first limit
in \eqref{p4:multiplier-successive-limits}.

All derivatives keep the coefficients fixed. From
$u_\eps=\eps-\eps^{3/2}d_\eps$ and
$S_\eps(b_\eps(r))=-r^2/2$ we obtain
\begin{equation}\label{p4:balanced-slope-limit}
 \begin{gathered}
 b_\eps'(r)=\eps^{5/2}(1-\sqrt\eps\,d_\eps(r))d_\eps'(r),\\
 m_\eps S_\eps'(b_\eps(r))
 =\frac{-2r}{(1-\sqrt\eps\,d_\eps(r))d_\eps'(r)}
 \longrightarrow L(r)
 \quad\text{uniformly on }J.
 \end{gathered}
\end{equation}
Indeed \eqref{p4:d-error} gives $d_\eps\to d_0$ in $C^1(J)$ and
$d_0'(r)=-1/(4\sqrt2\,r^{3/2}\sqrt{1-r})$ is bounded away from
zero there. By \eqref{p4:model-error}, $r_{j,\eps}\to r_j^0$.
Together with $\eps(c)\to0$, this proves
\eqref{p4:multiplier-successive-limits}.

The bounds needed for \eqref{p4:multiplier-intervals} follow from
exact rational brackets. The sign of $\Phi(r)$ on $(0,1)$ is the
sign of
\[
 \frac{1-r}{8r}-\left(\frac{101}{100}-r^2\right)^2
 =-\frac{10000r^5-20200r^3+11451r-1250}{10000r}.
\]
Direct substitution gives the endpoint signs in the following table.
Since $L(r)^2=128r^5(1-r)$ increases up to $r=5/6$ and decreases
afterwards, rational endpoint evaluation also gives its indicated
bounds throughout each bracket. The endpoints have common denominator
$10^5$:
\[
\begin{array}{c|c|c|c}
 j&(10^5a_j,10^5b_j)&(\operatorname{sgn}\Phi(a_j),
                  \operatorname{sgn}\Phi(b_j))
       &L(r)^2\ (a_j\le r\le b_j)\\ \hline
 1&(11159,11160)&(+,-)&(1/1000,1/400)\\
 2&(97868,97869)&(-,+)&(12/5,5/2)\\
 3&(99875,99876)&(+,-)&(3/20,1/5)
\end{array}
\]
Each bracket contains $r_j^0$, since $\Phi$ has exactly three
roots. Taking positive square roots gives
\eqref{p4:multiplier-intervals}.
\end{proof}


\subsection{Endpoint attainment}

\begin{lemma}[Endpoint attainment]\label{rew:lem-endpoint-attainment}
There are fields \(F_n\) in \eqref{eq:slice}, with
\(F_n\to H_0\) in quadratic coefficients, and three geometrically
distinct hyperbolic periodic orbits \(\gamma_{j,n}\) of each
\(F_n\), \(j=1,2,3\), such that
\[
 \max_{1\le j\le3}
 d_{\rm H}\bigl(\iota_0(\gamma_{j,n}),\Gamma\bigr)\longrightarrow0.
\]
For every fixed choice of admissible windows in
Lemma~\ref{lem:loc-capture}, all three orbits eventually follow the
same marked positive-time circuit \(q\to p\to q\), meeting its
return section once in their prime periods.
Their positive first-return multipliers $M_{j,n}$ satisfy
\[
 M_{j,n}\longrightarrow L(r_j^0),\qquad j=1,2,3,
\]
with the notation of Corollary~\ref{p4:multiplier-limits}.
\end{lemma}
\begin{proof}
Set \(c_n=1+K_0/(2n)\), with \(K_0\) from
Theorem~\ref{p4:fixed-c}. For each \(n\), first fix \(c_n\), then
choose \(0<\rho_n\le1/n\) so that the theorem applies and the same field
\(F_n:=F_{c_n,\rho_n}\) and its three hyperbolic orbits
\(\gamma_{j,n}:=\gamma_{j,c_n,\rho_n}\) satisfy
\begin{equation}\label{rew:eq-endpoint-diagonal-selection}
 \begin{gathered}
 \|F_n-H_{c_n}\|<\frac1n,\qquad
 \max_{1\le j\le3}
 d_{\rm H}\bigl(\iota_0(\gamma_{j,n}),\Gamma_{c_n}\bigr)<\frac1n,\\
 \max_{1\le j\le3}
 |M_j(c_n,\rho_n)-\mu_j(c_n)|<\frac1n.
 \end{gathered}
\end{equation}
This choice may depend on \(c_n\) and controls the actual coefficient
\(C_{c_n,\rho_n}-c_n\). The last condition is available by
Corollary~\ref{p4:multiplier-limits} after \(c_n\) is fixed; its
second limit gives the asserted convergence of
$M_{j,n}=M_j(c_n,\rho_n)$ on this same diagonal.
Continuity of \(H_c\) gives
\begin{equation}\label{rew:eq-endpoint-coefficients}
 \|F_n-H_0\|\le\frac1n+\|H_{c_n}-H_1\|\longrightarrow0.
\end{equation}

Pair the organizing parabolas at the same \(X\); their heights differ
by \((c_n-1)/2\). Bounded \(D\iota_0\) and comparability of the fixed
metric to the spherical metric give one constant \(M_\Gamma\) with
\begin{equation}\label{rew:eq-endpoint-graphics}
 d_{\rm H}(\Gamma_{c_n},\Gamma)\le M_\Gamma(c_n-1)
 \longrightarrow0.
\end{equation}
Both closures have the same endpoint \(q=(0,1,0)\) and no equatorial
arc. The triangle inequality gives
\begin{equation}\label{rew:eq-endpoint-cycles}
 \max_{1\le j\le3}
 d_{\rm H}\bigl(\iota_0(\gamma_{j,n}),\Gamma\bigr)
 \le\frac1n+M_\Gamma(c_n-1)\longrightarrow0.
\end{equation}

The same pairing gives convergence of the marked parts \(X\le-1\)
and \(X\ge-1\), whose cut is \(p_{c_n}=(-1,(1-c_n)/2)\to(-1,0)\).
The positive speed \eqref{intro:eq-parabola-speed} orders them as
\(q\to p_{c_n}\to q\). For any fixed admissible windows, speed bounds,
and inner height, \eqref{rew:eq-endpoint-coefficients} and
\eqref{rew:eq-endpoint-cycles} admit all three cycles to
Lemma~\ref{lem:loc-capture} for all sufficiently large \(n\).
It gives their marked circuit, transverse hits, and prime return.
\end{proof}

\begin{proof}[Proof of Theorem~\ref{sharp:lower-theorem}]
The inverse coefficient isomorphism \(\mathscr C_0^{-1}\) from
\eqref{eq:fixed-coefficient-transport} transports the fields of
Lemma~\ref{rew:lem-endpoint-attainment} to fields converging to \(F_*\),
hence eventually into \(\mathcal V\). The projective inverse \(S^{-1}\) in
\eqref{eq:fixed-projective-map} is Lipschitz on the compact hemisphere,
with constant \(L_-\), and sends \(\Gamma\) to \(\Gamma_{I4}\).
The transported cycles therefore have distances at most
\[
 L_-\max_{1\le j\le3}
 d_{\rm H}\bigl(\iota_0(\gamma_{j,n}),\Gamma\bigr),
\]
Choose \(n\) so that this is less than \(\delta\), the field lies in
\(\mathcal V\), and capture holds. The fixed coordinate change and
positive time factor preserve hyperbolicity, distinctness, and the
marked order \(q_*\to p_*\to q_*\).
\end{proof}

The strengthened diagonal in
\eqref{rew:eq-endpoint-diagonal-selection} realizes
Theorem~\ref{sharp:lower-theorem} with first-return multipliers converging
to \(L(r_1^0),L(r_2^0),L(r_3^0)\). For all sufficiently large \(n\),
these multipliers lie in \((0,1/2),(3/2,2),(0,1/2)\), respectively.
The fixed coordinate change and positive time factor preserve the
return multipliers, giving attracting, repelling and attracting cycles
in order of increasing \(r\).

\subsection{The coefficient face of the same-field triples}

The construction also identifies how the three cycles enter the small
positive ratio sector. Its calibration varies $A,C$ and preserves the
three coefficients prescribed in \eqref{p4:path}. We compute the ratio
directly from those coefficients, independently of the upper estimates.

\begin{proposition}\label{prop:lower-coefficient-face}
The fields in Lemma~\ref{rew:lem-endpoint-attainment} can be chosen with
\[
 \alpha_0<0<\alpha_1,\qquad p\longrightarrow0^+,
 \qquad \gamma/s\longrightarrow+\infty.
\]
In their positive-gamma coordinates
$g=\gamma$, $k=s/g^{3/2}$, $\pi=p/\sqrt g$ and
$\mathfrak q=k\pi$, they satisfy
\[
 k\longrightarrow0,\qquad \pi\longrightarrow5/\sqrt2,
 \qquad \mathfrak q\longrightarrow0.
\]
Their two positive axis roots in the coordinate
$(v-\gamma)/(gk)$ have squared half-gap tending to $1/72$.
In particular they remain on the bounded positive-gamma real-root
side, separated from the fold. The same field and its three complete
hyperbolic cycles meet every prescribed positive upper cutoff on $p$.
\end{proposition}

\begin{proof}
For each fixed $c>1$, write
\[
 \epsilon=\left(\tfrac12e^{-4\pi/\sqrt{c-1}}\right)^{2/5},
 \qquad \beta=1-\tfrac52\epsilon.
\]
Restrict the fixed $c$-range so that $0<\epsilon<1/10$.
Equation~\eqref{p4:path}, with its fixed root ratio $3/2$, gives
\begin{equation}\label{eq:lower-sector-exact-coefficients}
 \gamma=\frac{\beta\rho}{3},\qquad
 \alpha_0=-\frac32\epsilon^2\rho^3,\qquad
 \alpha_1=\left(\frac52\epsilon-\frac32\epsilon^2\right)\rho^2.
\end{equation}
These are the actual coefficients after calibration: the equations
\eqref{p4:calibration} change only $A,C$. Moreover a field $H_\theta$
already in the slice has, by local uniqueness in Lemma~\ref{lem:slice},
slice parameter $\theta(H_\theta)=\theta$ and affine/time parameters
$(I,0,1)$. Consequently after the fixed inverse transport
$\mathscr C_0^{-1}$ these same coefficients define
\eqref{eq:sharp-sector-coordinates} for the original field.

Taking the square root and dividing in
\eqref{eq:lower-sector-exact-coefficients} gives
\begin{equation}\label{eq:lower-sector-ratios}
 s=\sqrt{3/2}\,\epsilon\rho^{3/2},\qquad
 p=\frac{5-3\epsilon}{\sqrt6}\sqrt\rho,\qquad
 \frac\gamma s=\frac{\beta}
 {3\sqrt{3/2}\,\epsilon\sqrt\rho}.
\end{equation}
For $c_n\downarrow1$, the endpoint construction first fixes $c_n$,
then chooses one $\rho_n\le1/n$ for the actual coefficient difference
and all three whole-cycle distances. Thus
$0<p_n\le5/\sqrt{6n}\to0$, while $\gamma_n/s_n\to+\infty$.
The actual coefficient $C_{c_n,\rho_n}$ tends to $c_n$ with the
chosen accuracy; it need not equal $c_n$. No estimate for the inverse
calibration uniform as $c\downarrow1$ is used.

The gamma coordinates computed from the same path are
\begin{equation}\label{eq:lower-sector-gamma-coordinates}
 k=\frac{9\epsilon}{\sqrt2\,\beta^{3/2}},\qquad
 \pi=\frac{5-3\epsilon}{\sqrt{2\beta}},\qquad
 \mathfrak q=\frac{9\epsilon(5-3\epsilon)}{2\beta^2}.
\end{equation}
The translated roots in \eqref{p4:path} are
$-\rho,\epsilon\rho,3\epsilon\rho/2$. Dividing by
$gk=3\epsilon\rho/\sqrt{2\beta}$ gives
\[
 -r=-\frac{\sqrt{2\beta}}{3\epsilon},\qquad
 a=\frac{\sqrt{2\beta}}3,\qquad b=\sqrt{\beta/2}.
\]
Hence
\begin{equation}\label{eq:lower-sector-root-gap}
 \frac{(b-a)^2}{4}=\frac{\beta}{72}\longrightarrow\frac1{72},
 \qquad kr=\frac3\beta\longrightarrow3.
\end{equation}
In particular $\pi\to5/\sqrt2>\sqrt{12}$, on the separated real-root
side, with an escaping negative root at finite further radius.
These formulas describe this constructed family; they impose no
gamma-chart limit on other three-cycle sequences.

For any neighborhood, collar and positive ratio cutoff, take a
sufficiently large $n$ in this same sequence. Coefficient convergence,
the simultaneous Hausdorff limits and $p_n\to0^+$ meet all three
requirements at once. The fixed inverse coordinate change and positive
time preserve the three hyperbolic cycles of that single field.
\end{proof}

The one-cycle estimate at $p=0$ and these strictly positive-ratio
triples concern different coefficient limits. The realization passes
through the unbounded $\gamma/s$ branch. A bounded-$\gamma/s$
restriction would remove this sharpness family.

\section{The sharp small positive ratio sector}
\label{sec:sharp-small-sector}

The estimate down to $p=0$ uses the coefficient cover and original
physical passages directly. It is not obtained by sending the lower
endpoint of a compact-positive-ratio estimate to zero. The preceding
lower construction supplies its sharpness independently.

\begin{proof}[Proof of Theorem~\ref{thm:small-sector}]
Use Proposition~\ref{sharp:upper-coefficient-bounds} on the common
receiver of Lemma~\ref{lem:common-raw-receiver}, with \(h_c=b/2\).
Its fixed reserve \eqref{eq:pure-small-ratio-reserve} gives
\(0<p_{\rm s}\le p_*\) for which every \(0<p\le p_{\rm s}\)
is covered before the independent raw radii and final multiplier are
restricted. The proposition counts the full connected interval with
multiplicity, including both original ghost tails; no further regional
count is needed.

At $p=0$, $\alpha_1=0$.
Lemma~\ref{deg:lem-raw-jacobian} bounds the actual transition
Jacobian on the full original domain for
$\alpha_0\le0,\alpha_1\le0$. With the divided closing and one
further multiplier shrink, $F'>1/2$, so there is at most one
isolated zero and it is simple. These are fixed multiplier bounds;
\eqref{eq:finite-multiplier} makes the actual $m$ satisfy them by
an ordinary coefficient shrink, independently of $s,\gamma$ and
all root distances.

Apply geometric capture with these fixed choices, as in
Proposition~\ref{prop:central-bound}. This gives one slice
neighborhood and one collar for both statements. Use
Lemma~\ref{lem:slice-bound-transfer} with
$\mathcal B=\{\alpha_0<0,\ 0\le p\le p_{\rm s}\}$, and
intersect with its one-cycle version on $p=0$ if necessary. The
resulting full-germ neighborhood and collar are common to the whole
small-ratio sector. The stronger return-multiplicity assertion is the
slice count above, preserved at the corresponding physical sections.

Finally Proposition~\ref{prop:lower-coefficient-face} supplies
same-field hyperbolic triples with coefficient and whole-orbit
Hausdorff convergence and $p_n\to0^+$. Given
$\mathcal W,\delta,\eta$, choose one sufficiently large $n$ so
that its field lies in $\mathcal W$, all three distances are below
$\delta$, and $0<p_n<\min\{p_{\rm s},\eta\}$.
These are the three cycles calibrated by one coefficient pair in
Theorem~\ref{p4:fixed-c}, proving the sharpness clause.
\end{proof}

\subsection{Transfer to the full quadratic germ}

\begin{proof}[Proof of Theorem~\ref{thm:main}]
Proposition~\ref{prop:central-bound} supplies one slice neighborhood
and one complete-graphic collar with at most three isolated cycles.
Apply Lemma~\ref{lem:slice-bound-transfer} with the whole slice
parameter space and $N_{\rm cyc}=3$. This gives the stated upper
bound on one ordinary full twelve-coefficient neighborhood and collar.
Proposition~\ref{prop:lower-coefficient-face}, independently of that
count, gives one field with three hyperbolic complete cycles in every
smaller neighborhood and collar. Together these prove exact cyclicity
three with the asserted marking and quantifiers.
\end{proof}

The counting argument uses one ordered physical interval, a common
orientation of curvature zeros or displacement critical points, and a
bounded complete derivative on its complement. The family-specific
inputs are the cubic cancellation \eqref{neg:action-curvature},
relative division \eqref{neg:relative-estimate}, and cancellation of
the two physical tail gains in \eqref{sharp:upper-tail-ratio}.
The root and degree-loss estimates establish these inputs where the
ordinary through-domain margins vanish. The finite coefficient cover
then gives one neighborhood after fixing the raw cap and overlap reserves.

\appendix
\section{Scalar transitions and physical section calculus}

The regional passages require estimates on normalized scalar maps, on
derivatives with moving initial sections, and on local Dulac maps. We
develop those tools here. Regular adapters and divided closing maps are
handled by the identities of Section~\ref{sec:counting-tools}.

\subsection{Normalized scalar transitions}
Subtracting a reference orbit and dividing by its first variation removes
the affine part of a scalar transition. The remaining equation vanishes
to second order in the normalized phase variable. An integrable bound on that
equation gives uniform endpoint maps even when the unnormalized first
variation has no uniform bound.

\begin{lemma}[Normalized scalar transitions]\label{lem:scalar-normalization}
Let \(k\ge1\). Consider equations
\[
 \eta_r=N_\vartheta(r,\eta),\qquad L\le r\le R_\vartheta,
 \qquad N_\vartheta(r,0)=\partial_\eta N_\vartheta(r,0)=0,
\]
on \(|\eta|\le a\), where \(L\) and \(a>0\) are fixed.
Suppose that, uniformly in the parameters,
\[
 \|N_\vartheta(r,\cdot)\|_{C^{k+1}([-a,a])}\le g(r),
 \qquad g\in L^1([L,\infty)).
\]
Then there is a fixed neighborhood of zero on which all endpoint
maps have bounded \(C^{k+1}\) norms and positive first derivatives
bounded away from zero. Their inverses have the same bounds on a
fixed smaller neighborhood of zero.

Suppose also that \(R_\vartheta\to\infty\) and the kernels converge
in \(C^0([L,M];C^{k+1}_\eta)\) for every fixed \(M>L\).
Then the endpoint maps
and their inverses converge in \(C^{k+1}\) on smaller neighborhoods.
For a chosen parameter derivative \(\partial_\ell\), convergence
in \(C^k\) holds for the differentiated maps if the differentiated
kernels converge in \(C^0([L,M];C^k_\eta)\), have a common integrable
majorant, and
\[
 \|N_\vartheta(R_\vartheta,\cdot)
               \partial_\ell R_\vartheta\|_{C^k}\longrightarrow0.
\]
The convergence is joint in an additional compact parameter family
if these local convergences, \(R_\vartheta\to\infty\), and the
vanishing endpoint term are uniform in that family.
\end{lemma}
\begin{proof}
Taylor's formula gives \(|N_\vartheta(r,\eta)|\le g(r)|\eta|^2/2\).
Comparison with \(y'=g(r)y^2/2\) keeps all solutions from a common
enlarged initial interval inside \(|\eta|<a\). The first variation is
\[
 \exp\left(\int_L^r
          \partial_\eta N_\vartheta(t,\eta(t))\,dt\right)>0.
\]
It and its reciprocal are bounded. Induction in the triangular
variational equations and Gronwall's inequality bound all remaining
rows through order \(k+1\). The slope lower bound gives a common
image interval, and the inverse identity gives the inverse bounds.

Subtracting the integral equations and their variations gives convergence
on \([L,M]\); each remaining tail costs at most \(C\int_M^\infty g\).
Let the parameters tend to their limit and then \(M\to\infty\).
For the parameter derivative, the variation satisfies
\[
 Z_r=\partial_\eta N_\vartheta\,Z+\partial_\ell N_\vartheta,
 \qquad Z(L)=0.
\]
The endpoint derivative adds
\(N_\vartheta(R_\vartheta,\eta(R_\vartheta))
\partial_\ell R_\vartheta\). The same argument proves \(C^k\)
convergence under the endpoint assumption: differentiating a coefficient
\(N_{\eta^h}(r,\eta)\) adds
\(N_{\eta^{h+1}}(r,\eta)Z\), which requires the extra phase order.
Finally, \(C^{k+1}\) convergence of the endpoint maps \(A_\vartheta\) and
\[
 \partial_\ell A_\vartheta^{-1}
 =-\left(\frac{\partial_\ell A_\vartheta}{A_\vartheta'}\right)
     \circ A_\vartheta^{-1}
\]
give \(C^k\) convergence of the differentiated inverses. Each estimate
is uniform in the additional compact family under the stated assumptions.
\end{proof}

For a graph equation \(Y_r=F(r,Y)\), a reference solution \(Y_*\)
and its first variation \(j>0\), the substitution
\(\eta=(Y-Y_*)/j\) gives the exact kernel
\begin{equation}\label{tail:normalized-kernel}
 N(r,\eta)=j\eta^2\int_0^1(1-t)
                  F_{YY}(r,Y_*+tj\eta)\,dt.
\end{equation}
For \(h\ge2\), \(\partial_\eta^hN=j^{h-1}F_{Y^h}(r,Y_*+j\eta)\).
At fixed \(r,\eta\), a parameter derivative is therefore
\begin{equation}\label{deg:det-parameter-phase-derivative}
\begin{split}
 \partial_\ell\partial_\eta^hN
 =j^{h-1}\bigl\{&(h-1)(\partial_\ell\log j)F_{Y^h}
 +\partial_\ell^{\rm exp}F_{Y^h}\\
 &+F_{Y^{h+1}}(Y_{*,\ell}+j_\ell\eta)\bigr\}.
\end{split}
\end{equation}
Here \(\mathrm{exp}\) holds the graph coordinates fixed. The orders
zero and one follow by integrating in \(\eta\), since
\(N(r,0)=N_\eta(r,0)=0\). Thus \(C^3\)-valued parameter
derivatives require \(F_{Y^4}\) and \(N_{\eta^4}\).
Integrating the indicated parameter derivative gives the sequential
parameter secants used below. This requires neither a mixed parameter
derivative nor convergence of \(j\).

A quantitative form controls how close the normalized map is to the
identity when only an upper bound on the first variation is available.
For an actual graph family
write \(J_n=\partial_b^nY\), \(Q_n=J_n/J_1\), and
\(a_n=F_{Y^n}J_1^{n-1}\). Direct differentiation gives
\begin{equation}\label{tail:relative-jet-equations}
 Q_2'=a_2,\qquad Q_3'=3a_2Q_2+a_3,\qquad
 Q_4'=a_2(4Q_3+3Q_2^2)+6a_3Q_2+a_4.
\end{equation}
If, uniformly on an enlarged label interval,
\(Q_n(L)=O(\varepsilon)\) and
\(\int|a_n|\,dr=O(\varepsilon)\), \(2\le n\le4\),
these identities give \(Q_n=O(\varepsilon)\).
Since \(\partial_b\log J_1=Q_2\), the endpoint map in \(b-b_*\),
normalized by its exact value and slope at \(b_*\), and its inverse on a
smaller common interval, are \(\mathrm{id}+O_{C^3}(\varepsilon)\)
for sufficiently small \(\varepsilon\).
This conclusion is independent of the size of \(J_1(b_*)>0\).


\subsection{Moving-section variations and inward continuation}

Moving a section changes both the initial height and the time of
departure. The following identities retain both effects, including the
endpoint terms that appear after using height as the independent variable.

\begin{lemma}[Scalar differentiation on moving sections]
\label{deg:lem-scalar-moving-identities}
Fix the coefficients of a scalar graph \(z_x=f(x,z)\), with initial
curve \((x_l(a),z_l(a))\), and let \(D\) be a phase vector field,
commuting with differentiation at fixed \(x\). Write \(V_j=D^jz\).
Through order three its equations and moving initial rows are
\begin{equation}\label{deg:scalar-moving-bell}
 \begin{gathered}
 (V_1)_x=f_zV_1,\qquad (V_2)_x=f_zV_2+f_{zz}V_1^2,\\
 (V_3)_x=f_zV_3+3f_{zz}V_1V_2+f_{zzz}V_1^3,\\
 B_1=Dz_l-f_lDx_l,\quad B_2=DB_1-f_{z,l}B_1Dx_l,\\
 B_3=DB_2-(f_{z,l}B_2+f_{zz,l}B_1^2)Dx_l,
 \qquad V_j(x_l;a)=B_j .
 \end{gathered}
\end{equation}
For \(f=z(n+\zeta z)/(M+Qz)\) on a positive graph with positive
denominator, put \(\mathscr D=M+Qz\).
The complete fixed-port first derivative, including an observer \(\chi\), is
\begin{equation}\label{deg:scalar-raw-cocycle}
 \begin{gathered}
 f_z=\frac d{dx}\log\frac z{\mathscr D}
          +\frac{M'+(Q'+\zeta)z}{\mathscr D},\\
 D\chi(z_T)=\chi'(z_T)B_1
    \frac{(z/\mathscr D)_T}{(z/\mathscr D)_l}
    \exp\int_{x_l}^{x_T}\frac{M'+(Q'+\zeta)z}{\mathscr D}\,dx .
 \end{gathered}
\end{equation}
For \(z_l=\ell_0M(a)\), \(n=xB\), and
\(b=-a^2/(2\sigma^2)\), this becomes
\begin{equation}\label{deg:scalar-polynomial-lip-slope}
 \begin{gathered}
 \Gamma_b'=\sigma^2\chi'(z_T)\frac{z_T}{\mathscr D_T}B(a)K(a)e^{I(a)},
 \quad I=\int_a^{x_T}\frac{M'+(Q'+\zeta)z}{\mathscr D}\,dx,\\
 K=1-\frac{M'(a)}{aB(a)}(1+\ell_0Q(a))
       +\frac{\ell_0\zeta(a)M(a)}{aB(a)}.
 \end{gathered}
\end{equation}
If instead \(r_\eta=V(r,\eta)\), with moving data
\((\eta_l,r_l)\), write \(P_j=D^jr|_\eta\), \(l_j=D^j\eta_l\),
\(t_j=D^j\eta_T\), where \(r_T\) is fixed. Then
\begin{equation}\label{deg:scalar-inverse-endpoints}
 \begin{gathered}
 P_{1,l}=Dr_l-V_ll_1,\quad
 P_{2,l}=DP_{1,l}-V_{r,l}P_{1,l}l_1,\\
 (P_1)_\eta=V_rP_1,\quad
 (P_2)_\eta=V_rP_2+V_{rr}P_1^2,\\
 t_1=-P_{1,T}/V_T,\quad
 t_2=-\{P_{2,T}+2V_{r,T}P_{1,T}t_1
              +(V_{r,T}V_T+V_{\eta,T})t_1^2\}/V_T .
 \end{gathered}
\end{equation}
For \(I=\int_{\eta_l}^{\eta_T}J(r(a,\eta),\eta)\,d\eta\),
its full rows are
\begin{equation}\label{deg:scalar-density-endpoints}
 \begin{aligned}
 DI={}&J_Tt_1-J_ll_1+\int J_rP_1\,d\eta,\\
 D^2I={}&\int(J_{rr}P_1^2+J_rP_2)\,d\eta
       +(J_rP_1)_Tt_1-(J_rP_1)_ll_1\\
 &+J_Tt_2+J_{\eta,T}t_1^2-J_ll_2
       -(J_{r,l}Dr_l+J_{\eta,l}l_1)l_1 .
 \end{aligned}
\end{equation}
The observer and divided-closing rows use
Lemma~\ref{phys:lem-c3-adapters}, specifically
\eqref{phys:eq-observer-log} and \eqref{phys:eq-closing-jets}.
Coefficient-only gains and affine calibrations are frozen under \(D\).
\end{lemma}
\begin{proof}
Differentiate the graph equation at fixed \(x\); differentiating
\(V_j(x_l(a),a)=B_j(a)\) gives the initial recursion.
Rational differentiation proves \eqref{deg:scalar-raw-cocycle}.
For the inverse graph differentiate twice at fixed \(\eta\), and
then twice the identity \(r(a,\eta_T(a))=r_T\).
Leibniz differentiation gives \eqref{deg:scalar-density-endpoints};
at the lower endpoint the total radial derivative is \(Dr_l\),
whereas at the upper endpoint it is zero. The observer and closing
claims follow from the cited ordinary \(C^3\) identities.
\end{proof}

In the inward passages, the ratio of height to the axis polynomial
decays on the action clock. The next lemma estimates the resulting
integral without discarding the dependence on its terminal clock time.

\begin{lemma}[Stable ratio integrals with moving endpoints]
\label{deg:lem-stable-ratio-ledger}
At fixed action clock \(q\), let a phase field \(D\) commute with
\(\partial_q\). For \(\kappa>0\), consider
\[
 W_q=\mathcal F=W\{- (1+\omega W)/(1+QW)+\kappa H\},
 \qquad W(0)=\ell_0,
 \qquad L=(Q-\omega)W/(1+QW).
\]
Suppose \(Q,\kappa H,\omega\) and their required two clock rows are
uniformly bounded, \(1+QW\) is a positive unit, and both
\(\mathcal F/W\) and \(\mathcal F_W\) are at most \(-c_0<0\)
on \(0\le W\le\ell_0\). For some \(0\le\rho\le\rho_0\),
assume the first three phase rows of \(\kappa H,Q,\omega\),
and mixed rows containing a phase derivative of total order at most three,
are bounded by
\(C\rho e^{\alpha q}\), with \(\alpha<c_0/12\).
There are an integer \(N\) and \(0<c<c_0/2\), depending only on the
stated bounds, such that
\begin{equation}\label{deg:stable-ratio-fixed-clock}
 0<W\le\ell_0e^{-c_0q},\qquad
 |D^jW|+|D^jL|\le C\rho(1+q)^Ne^{-cq},\quad1\le j\le3.
\end{equation}
The needed mixed rows obey these bounds; pure clock rows have the
same exponential bound without \(\rho\).
For a moving \(T\), put \(T_j=D^jT\) and \(I_T=\int_0^T L\,dq\).
All endpoint contributions are
\begin{equation}\label{deg:stable-ratio-terminal-rows}
 \begin{aligned}
 DI_T={}&\int_0^T DL\,dq+L_TT_1,\\
 D^2I_T={}&\int_0^T D^2L\,dq+2(DL)_TT_1+(L_q)_TT_1^2+L_TT_2,\\
 D^3I_T={}&\int_0^T D^3L\,dq+3(D^2L)_TT_1
       +3(DL_q)_TT_1^2+(L_{qq})_TT_1^3\\
 &+3(DL)_TT_2+3(L_q)_TT_1T_2+L_TT_3 .
 \end{aligned}
\end{equation}
If \(v_q=-\kappa P(v)/(vB(v))\), \(H=P'/(vB)\), and
\(W=Y/P(v)\), its exact polynomial memory is
\begin{equation}\label{deg:stable-ratio-polynomial-memory}
 W(q)=\ell_0\frac{P(a)}{P(v(q))}e^{-q+I(q)},\qquad
 Y_T=\ell_0P(a)e^{-T+I_T}=P(v_T)W(T).
\end{equation}
The corresponding rows of \(W(T)\) are the ordinary chain rule
\(DW=W_D+W_qT_1\),
\(D^2W=W_{DD}+2W_{Dq}T_1+W_{qq}T_1^2+W_qT_2\), and its
third row with coefficients \(1,3,3,1,3,3,1\).
Thus the interior integral contributions are \(O(\rho)\); the full
moving-limit rows also retain their terminal contributions, which are at most
\(C(1+T)^Ne^{-cT}(1+\sum_{j=1}^3|T_j|)^3\).
The constants depend only on the stated coefficient bounds,
\(c_0,\ell_0,\rho_0\), independently of the flight length.
\end{lemma}
\begin{proof}
The negative logarithmic rate gives the first-exit bound for \(W\).
Differentiate its equation three times. Each row has diagonal
\(\mathcal F_W\); the explicit coefficient forcing is
\(O(\rho e^{\alpha q}W)\), and other terms are products of lower
rows. Successive variation of constants absorbs at most three
coefficient growth factors into the fixed exponential decay and
proves \eqref{deg:stable-ratio-fixed-clock}. Differentiating these
equations in \(q\) gives the mixed rows. Leibniz differentiation and
the ordinary chain rule give the terminal formulas.
Since \(\partial_q\log P(v)=-\kappa H\), integrating the ratio
equation gives \eqref{deg:stable-ratio-polynomial-memory}.
\end{proof}

When the radial route is fixed, the same decay controls ordinary
derivatives with respect to the initial height, even for arbitrarily
long flights.

\begin{lemma}[Initial-height rows on a frozen ratio route]
\label{deg:lem-fixed-route-height-rows}
Consider the ratio equation and density \(\mathcal F,L\) of
Lemma~\ref{deg:lem-stable-ratio-ledger}, with coefficients and radial
endpoints fixed. On its full prescribed buffer \(0\le W\le\ell_0\),
assume the stated positive-unit denominator, the two negative margins
\(\mathcal F/W,\mathcal F_W\le-c_0<0\), and bounded
\(Q,\omega,\kappa H\). Let \(W(0;w)=w\), and, for a fixed finite
flight \(T\ge0\), put
\[
 \mathcal N_T(w)=\int_0^T L(q,W(q;w))\,dq,\qquad e_w=w\partial_w.
\]
Uniformly in \(T\) and \(0\le w\le\ell_0\), the ordinary extension
at \(w=0\) satisfies
\[
 \mathcal N_T(0)=0,\qquad
 |e_w^j\mathcal N_T(w)|\le Cw\quad(0\le j\le3),\qquad
 \sum_{j=1}^3|\partial_w^j\mathcal N_T(w)|\le C.
\]
The normalized nonlinear unit \(e^{\mathcal N_T}\) and its reciprocal
have bounded ordinary \(C^3\) norms. In the polynomial-memory setting,
\(Y_T=Y_0e^{-T+\mathcal N_T(w)}\), with \(Y_0=P(v_0)w\).
\end{lemma}
\begin{proof}
The buffer and negative logarithmic margin give
\(0\le W(q;w)\le we^{-c_0q}\). Rational height differentiation
bounds \(\mathcal F_{WW},\mathcal F_{WWW}\) and the first three
height derivatives of \(L\). Apply the triangular equations
\eqref{deg:scalar-moving-bell} to \(V_j=\partial_w^jW\), with
\(V_1(0)=1\), \(V_2(0)=V_3(0)=0\). Successive variation of
constants gives \(|V_j(q)|\le Ce^{-c_0q}\), \(1\le j\le3\).
Every term in \(\partial_w^jL(q,W(q;w))\) contains a decaying
variation, so its integral is uniformly bounded. Also \(|L|\le CW\)
gives \(|\mathcal N_T|\le Cw\). The ordinary extension at zero and
the expansion of \((w\partial_w)^j\) give the weighted rows.
Exponentiation and
\eqref{deg:stable-ratio-polynomial-memory} prove the final claims.
\end{proof}

After the long inward flight, a short central segment remains. Factoring
out its linear propagator gives a nonlinear error controlled by the
incoming height, uniformly as that propagator tends to zero.

\begin{lemma}[Compact inward height continuation]
\label{deg:lem-compact-height-continuation}
For a frozen coefficient field, suppose the actual equation on a
fixed segment, oriented from \(b>0\) to zero, is
\(Y_v=A(v)Y+K(v,Y)Y^2\), with
linear propagator \(m(b)=1\), \(0<m(v)\le1\). On its prescribed
height buffer \(0\le Y\le h_0\), assume the ordinary height
derivatives of \(K\) through order three are bounded by
\(C/\kappa_c\), for a local parameter \(\kappa_c>0\).
For \(0<y\le h_0/2\) and \(y/\kappa_c\) sufficiently small,
its actual continuation satisfies
\begin{equation}\label{deg:scalar-compact-propagator}
 Y(0;y)=m(0)y e^{\mathcal N(y)},\qquad
 |(y\partial_y)^j\mathcal N(y)|\le Cy/\kappa_c,\quad0\le j\le3.
\end{equation}
In composition with \(y=y(a)\), with the field, segment and
\(\kappa_c\) fixed, put \(\Lambda_j=D^j\log y\) and
\(e=y\partial_y\). The three rows of \(\mathcal N(y(a))\) are
\(e\mathcal N\,\Lambda_1\),
\(e^2\mathcal N\,\Lambda_1^2+e\mathcal N\,\Lambda_2\), and
\(e^3\mathcal N\,\Lambda_1^3+3e^2\mathcal N\,\Lambda_1\Lambda_2
+e\mathcal N\,\Lambda_3\). No reciprocal bound on \(m(0)\) is used.
Constants depend only on the stated height bounds and fixed segment,
independently of the extracted multiplier.
\end{lemma}
\begin{proof}
Write \(Y=m(v)yu(v;y)\). The integral
equation for \(u\), with \(u(b)=1\), has coefficients and its first
three Euler forcings bounded by \(Cy/\kappa_c\) on
\(1/2\le u\le2\). First exit and three triangular integral
estimates give \(u^{\pm1}\le C\) and
\(|e^j\log u(0;y)|\le Cy/\kappa_c\). The displayed composition
rows follow by the chain rule.
\end{proof}

\subsection{An exterior spectral estimate}

The outward cubic, quadratic and degree-loss passages share a scalar
graph equation. The quadratic coefficient determines a power in its
first variation, while the remaining terms control the error in its
logarithmic derivatives. We use logarithmic height as the independent
variable to estimate these derivatives directly, including the moving
endpoints.

\begin{lemma}[Exterior spectral estimate]\label{deg:lem-spectral-exterior}
Fix a terminal cut \(d>0\), a small lip constant \(\ell_0>0\), compact
bounds on \(A,C\), and positive constants \(c_0,C_0\). Put
\[
 \begin{gathered}
 M=tx^3+\mu_2x^2+\mu_1x+\mu_0,\qquad B=1+b+tx,\\
 Q=1-A(v+tx)+C(v+tx)^2,\quad \zeta=tC(v+tx),\quad |t|\le1.
 \end{gathered}
\]
All coefficients are fixed under \(E=a\partial_a\). On an enlarged
open interval of \(0<a<d/2\), suppose, pointwise in \(a\), that
some \(0<\rho\le\rho_0\) satisfies
\begin{equation}\label{deg:spectral-exterior-hypotheses}
 \begin{gathered}
 c_0\rho\le m(x):=M(x)/x^2\le C_0\rho\quad(a\le x\le2a),\\
 |b|+|v|+|\mu_2|+|t|a+|\mu_1|/a+|\mu_0|/a^2\le C_0\rho .
 \end{gathered}
\end{equation}
For sufficiently small \(|t|d,\rho_0\), depending only on these
fixed bounds, the graph
\begin{equation}\label{deg:spectral-exterior-graph}
 z_x=\frac{z(xB+\zeta z)}{M+Qz},\qquad z(a)=\ell_0M(a)
\end{equation}
reaches \(x=d\) with positive height and denominator. An ordinary
\(C^3\) observer \(\chi\), with \(\chi'>0\), is appended on the
open preimage of its prescribed buffered terminal image. Set
\[
 \nu=\frac{2\mu_2}{1+b},\quad\lambda=2\nu,\quad e=2-\lambda>3/2,
 \quad\Theta=|t|a+|\mu_1|/a+|\mu_0|/a^2,\quad L_\rho=1+|\log\rho|,
\]
\[
 H_\chi=\sup\{|z(\log\chi')'|
       +|z(\log\chi')'+z^2(\log\chi')''|\},\qquad
 \mathcal T=|t|d+|\mu_2|+|\mu_1|/d+|\mu_0|/d^2+H_\chi ,
\]
where the supremum is on that buffer and is finite.
For any coefficient-only \(\sigma>0\), use
\(\beta=-a^2/(2\sigma^2)\) and \(\Gamma(\beta)=\chi(z_d)\).
Then \(\Gamma'>0\), and
\begin{equation}\label{deg:spectral-exterior-rows}
 |E\log\Gamma'+\lambda|+|E^2\log\Gamma'|
 \le C\{\Theta L_\rho+\mathcal T(a/d)^e\}.
\end{equation}
The same rows hold after a positive coefficient-only affine calibration.
For intrinsic terminal height take \(\chi=\mathrm{id}\), so \(H_\chi=0\).
No reciprocal bound on \(\sigma\) or bound on \(\lambda\log a\) is used.
\end{lemma}
\begin{proof}
Set \(W=z/x^2\), \(\tau=\log(x/a)\), \(N=B+\zeta xW\), and
\(D=m+QW\). Then \(W_\tau=F=W(N/D-2)\). On \([a,2a]\),
\(N,Q\) are positive units and \(D\asymp\rho+W\) for bounded \(W\).
Below a small fixed \(w_0\), \(F\ge cW/(\rho+W)\). Starting at
\(W_l=\ell_0m(a)\), the time to reach \(w_0\) is at most
\[
 C\rho\log\frac{w_0}{\ell_0c_0\rho}+Cw_0<\log2,
\]
by choosing \(w_0\) first and then \(\rho_0\). At the boundaries
\(w_0<1/2\), \(w_1=3/4\), the formula \(F(x,w)=1-2w+O_w(\rho)\)
makes the strip invariant up to \(2a\). Beyond it,
\(|m(x)|\le|t|d+C\rho\ll w_0^2\) preserves the strip, gives
positive-unit \(D\) and \(F_W\le-c<0\). Thus the graph reaches
\(d\), with \(c\rho\le W\le w_1\), even if \(m\) changes sign.

The bounded strong graph with \(b=v=\mu_j=0\), keeping \(t,A,C\),
is \(W_*=1/2+O(|t|x)\), by the integral from zero for
\(xU'=-2U+O(|t|x)\), \(U=W-1/2\). Strip contraction gives
\begin{equation}\label{deg:spectral-exterior-window}
 |W(d)-W_*(d)|\le C\rho+C(a/d)^c.
\end{equation}
Where this comparison applies, the original buffer margin and the
permitted cutoff choices put the terminal height in the observer's
domain. Elsewhere we retain the original terminal admission condition.

Put \(r=\log x\), \(\eta=\log z\). At fixed height, \(W_r=-2W\)
and \(r_\eta=V=D/N\). The raw cocycle
\eqref{deg:scalar-raw-cocycle} has the exact density split
\begin{equation}\label{deg:spectral-exterior-density}
 \begin{gathered}
 \frac{2m+xm'+(xQ'+\zeta x)W}{N}=\nu+J,\\
 J=\frac{(3-\nu)tx+\mu_1/x+
          tx[-A+(3-\nu)C(v+tx)]W}{N}.
 \end{gathered}
\end{equation}
Here \(2m+xm'=3tx+2\mu_2+\mu_1/x\); \(\mu_0\) cancels only
in the density. Let
\(I=\int_{\eta_l}^{\eta_T}J\,d\eta\), with
\(\eta_l=\log(\ell_0M(a))\). The near height length is \(O(L_\rho)\);
on the strip \(d\eta\le C\,d\log x\). Thus every partial integral
is bounded by \(C\{\Theta L_\rho+|t|d+|\mu_1|/a\}\).

Use \eqref{deg:scalar-inverse-endpoints} with
\(P_j=E^jr|_\eta\), \(l_j=E^j\eta_l\), \(t_j=E^j\eta_T\),
\(Er_l=1\). The polynomial gives bounded \(l_1,l_2\),
\(V_l,V_{r,l}=O(\rho)\), and \(P_{1,l}=1+O(\rho)>0\).
For \(\mathscr B=EW_l-F_l\), the exact first cocycle gives
\begin{equation}\label{deg:spectral-exterior-inverse-decay}
 \begin{gathered}
 P_1=\frac{-\mathscr B(1+\ell_0Q(a))}{\ell_0N}
       (a/x)^e\left(\frac W{\ell_0m(a)}\right)^\nu e^I,\\
 0<P_1\le C(a/x)^e,\qquad |P_2|\le CP_1 .
 \end{gathered}
\end{equation}
The initial factor equals \(B(a)K(a)>0\), with \(K\) from
\eqref{deg:scalar-polynomial-lip-slope}. Since \(|\nu|\le C\rho\)
and \(c\rho\le W\le w_1\), its height ratio to power \(\nu\) is a
unit for either sign, proving the first bound without bounding
\(a^{-\lambda}\). Near the lip, fixed-height differentiation gives
\(|V_{rr}|\le C(\rho+W)\), whereas \(d\eta=N/D\,d\log x\).
On the strip all these factors are bounded. Hence
\[
 |V_{rr}|\,d\eta\le C\,d\log x,\qquad
 (P_2/P_1)_\eta=V_{rr}P_1.
\]
Its integral is bounded by \(C\int_a^d(a/x)^e\,d\log x\), proving
the second bound. The two fixed-terminal rows then give
\(|t_1|+|t_2|\le C(a/d)^e\).

Rational differentiation at fixed height gives
\(|J_r|+|J_{rr}|\le C(|t|x+|\mu_1|/x)\). In the two full Leibniz
rows \eqref{deg:scalar-density-endpoints}, the lower terms cost
\(C\Theta\), the near integrals cost \(C\Theta L_\rho\), and
\[
 \int_{2a}^d(|t|x+|\mu_1|/x)(a/x)^e\,d\log x
 \le C(|t|a+|\mu_1|/a).
\]
Every upper term retains \((a/d)^e\). Therefore
\begin{equation}\label{deg:spectral-exterior-integral-rows}
 |EI|+|E^2I|
 \le C\{\Theta L_\rho+(|t|d+|\mu_1|/d)(a/d)^e\}.
\end{equation}
The exact polynomial-lip formula becomes
\begin{equation}\label{deg:spectral-exterior-slope}
 \Gamma'=\sigma^2a^{-\lambda}
 B(a)K(a)\ell_0^{-\nu}m(a)^{-\nu}
 \frac{W_T}{m(d)+Q(d)W_T}z_T^\nu\chi'(z_T)e^I .
\end{equation}
Both \(E\) rows of \(\log(BK)\) are \(O(\Theta)\). Also
\(E m,E^2m=O(\Theta)\), \(m(a)\asymp\rho\), and
\(|\nu|,\Theta\le C\rho\), so those of \(-\nu\log m(a)\) cost
\(C\Theta\). At fixed \(d\), the first two height-logarithmic rows
of \(\log[W_T/(m(d)+Q(d)W_T)]\) cost \(C|m(d)|\), and those of
\(\nu\log z_T\) cost \(C|\nu|\). The observer rows are precisely
those defining \(H_\chi\). Their terminal chain rules and the integral
bounds prove \eqref{deg:spectral-exterior-rows}, since
\(a^{-\lambda}\) contributes \(-\lambda\) and zero to the first and
second Euler rows.
\end{proof}


\subsection{Exact identities on moving polynomial sections}
\label{subsec:deg-moving-sections}

For the cubic, quadratic and degree-drop scalings consider
\begin{equation}
 \dot\xi=-\epsilon\{P(\xi)+Q(\xi)Y\},
 \qquad
 \dot Y=-Y\{\xi B(\xi)+\zeta(\xi)Y\},
 \label{deg:tail-field}
\end{equation}
where \(\epsilon>0\) and \(B,Q\) are positive units, with transverse sections
\begin{equation}
 Y=\ell_0P(\xi),\qquad 0<\ell_0\ll1,
 \label{deg:moving-section}
\end{equation}
Put
\[
 A(\xi)=\int^\xi {xB(x)\over P(x)}\,dx,
 \qquad \mathcal D={P(\xi)\over \xi B(\xi)}\partial_\xi .
\]

\begin{lemma}[Exact moving-section identities]\label{deg:lem-moving-lip}
On a positive graph of \eqref{deg:tail-field} with
\(P>0\) and \(P+QY>0\), put \(W=Y/P\). Then
\begin{equation}
 {d\over d\xi}\left\{
  \epsilon\log Y-A(\xi)\right\}
 ={W\over1+QW}\left\{\zeta-Q{\xi B\over P}\right\}.
 \label{deg:exact-action-cocycle}
\end{equation}
If this graph joins the two lips at \(u\) and \(-r\), normalize
\(A(0)=0\) and write \(A_+(u)=A(u)\), \(A_-(r)=A(-r)\).
With \(P_+(u)=P(u)\), \(P_-(r)=P(-r)\), one has
\begin{equation}
 A_+(u)-A_-(r)
 =\epsilon\{\log P_+(u)-\log P_-(r)+R(u,r)\}.
 \label{deg:action-matching}
\end{equation}
where \(\epsilon R\) is the integral of the right-hand side of
\eqref{deg:exact-action-cocycle} from \(u\) to \(-r\) along this
graph. The two endpoint polynomials and the whole integral are retained.

For \(f=dY/d\xi\), its nonlinear derivatives and the inward
action density \(dq=|dA|/\epsilon\), away from \(\xi=0\), are
\begin{equation}
 \partial_Y^nf=(-1)^{n-2}n!
 \frac{PQ^{n-2}(\zeta P-\xi BQ)}{\epsilon(P+QY)^{n+1}},
 \quad n\ge2,\qquad
 |d\xi|=\frac{\epsilon P}{|\xi|B}\,dq.
 \label{deg:bell-density}
\end{equation}
For a moving initial curve \((\xi_l(x),Y_l(x))\), let
\(Y_T(x)\) be its graph value at a fixed terminal abscissa and let
\(\varphi_T\) be a fixed regular terminal adapter. Write
\(J\) for the scalar first variation. Then
\begin{equation}
 \frac d{dx}\varphi_T(Y_T(x))
 =\varphi_T'(Y_T)J(\xi_T,\xi_l)
                 \{Y_l'(x)-f(\xi_l,Y_l)\xi_l'(x)\}.
 \label{deg:moving-first-variation}
\end{equation}
In the label \(b=-x^2/2\), this derivative is multiplied by
\(-1/x\). Coefficient-only affine calibrations leave logarithmic
slope derivatives unchanged. The quantitative bounds for the actual
polynomial passages are established separately below and in
Appendix~\ref{app:moving-section-bookkeeping}.
\end{lemma}

\begin{proof}
Division of the graph equation and integration give the first two
identities; rational differentiation gives \eqref{deg:bell-density}.
The initial variation at fixed abscissa is
\(Y_l'-f_l\xi_l'\); its scalar variational equation gives
\eqref{deg:moving-first-variation}. The remaining statements are
the ordinary chain rule.
\end{proof}

\subsection{A compact-spectrum local Dulac estimate}

The following local Dulac estimate is used in the compact positive-ratio
passage and in the stopped determinant chart. Its logarithmic remainder
bounds remain uniform when the characteristic ratio crosses one.

\begin{lemma}[Uniform local positive Dulac estimate]
\label{lem:uniform-positive-dulac}
Consider a compact analytic family of hyperbolic saddles with fixed compatible
entrance and exit sections.  Suppose the positive tangential and normal rates
\(L,\nu\) are bounded above and away from zero and
\mbox{\(\rho=\nu/L\in[\rho_0,\rho_1]\Subset(0,\infty)\)}. Writing \(\theta\) for
the family parameter, in a positive local coordinate \(H\) which vanishes at
the incoming separatrix, the local Dulac map has the form
\[
 \mathcal D(H)=c_\theta H^\rho e^{\mathcal R_\theta(H)},
 \qquad c_\theta>0.
\]
For every \(0<\epsilon<\min(1,\rho_0)\), there is a uniform constant
\(C<\infty\) such that
\[
 \left|(H\partial_H)^j\mathcal R_\theta(H)\right|
 \le C H^\epsilon(1+|\log H|),
 \qquad 0\le j\le2.
\]
The estimate includes \(\rho=1\), and all local coordinate changes preserve
the positive-time orientation.
\end{lemma}
\begin{proof}
Straighten the invariant axes. A positive time factor gives
\(u'=Lu\). Linearize the one-dimensional stable-axis equation and
then replace \(v\) by \(ve^{h(u)}\), where
\(h'(u)=(b(u,0)-\nu)/(Lu)\). The quotient is removable at zero.
The remaining normal-coefficient defect vanishes on both axes, so
on one common small rectangle the equation is
\[
 u'=Lu,\qquad v'=-\nu v(1+uvG_\theta(u,v)),
 \qquad \rho=\nu/L\in[\rho_0,\rho_1].
\]
These changes use no division by \(\rho-1\), preserve positive
time, and have common finite derivative bounds on the compact family.
Take the entrance \((H,v_0)\) and the exit \(u=u_0\).
All phase derivatives below hold the coefficients and these cuts fixed.
Put \(\mathsf B=-\rho vG_\theta\). The graph equation is
\[
 \partial_u\log v=-\rho/u+\mathsf B(u,v).
\]
After shrinking the rectangle, \(v\) decreases. Its integral
equation first gives \(v(u;H)\le C(H/u)^\rho\), and then
\[
 I_\rho(H):=H^\rho\int_H^{u_0}u^{-\rho}\,du
 \le CH^{\min(1,\rho)}(1+|\log H|).
\]
The bound is uniform through \(\rho=1\); no separate bound for
\((1-\rho)^{-1}\) is taken.

Let \(\mathcal E=H\partial_H\), \(P=\mathcal E\log v\), and
\(Z=\mathcal E^2\log v\), at fixed \(u\). Differentiating the
moving initial condition \(v(H;H)=v_0\), as in
\eqref{deg:scalar-moving-bell}, gives
\[
 \begin{aligned}
 P(H)&=\rho-H\mathsf B(H,v_0),& Z(H)&=O(H),\\
 P_u&=\mathsf B_{\log v}P,
 &Z_u&=\mathsf B_{\log v}Z+\mathsf B_{\log v\log v}P^2.
 \end{aligned}
\]
The two logarithmic height derivatives of \(\mathsf B=-\rho vG\)
are \(O(v)\). Variation of constants therefore bounds \(P>0\) and gives
\( |P(u_0)-\rho|+|Z(u_0)|\le C(H+I_\rho(H))\).
For
\[
 \mathcal R(H)=\log\frac{v(u_0;H)}{v_0(H/u_0)^\rho},
\]
one has \(\mathcal E\mathcal R=P(u_0)-\rho\) and
\(\mathcal E^2\mathcal R=Z(u_0)\); its value is bounded by
the integral of \(\mathsf B\). These prove the stated estimate
for every fixed \(0<\epsilon<\min(1,\rho_0)\).
The endpoint \(H=0\) is a separatrix endpoint, not a through orbit.
\end{proof}

\section{Strong projection and physical section heights}
\label{sec:strong-projection}
\label{subsec:normal-time-projection}

The outgoing passages are parametrized by their projection to the
equatorial axis. We construct this projection for a common normal-time
field, with bounds independent of the length of the logarithmic corridor.
An integrating coordinate then converts the axis action into the height
of the actual orbit on a fixed section.

\begin{lemma}\label{lem:normal-time-projection}
Fix \(0<w_0<1/4\).  Let \(h,c,k\) be \(C^4\) functions on
\(\R\), with
\[
 \|h\|_{C^4}+\|c-1\|_{C^4}+\|k\|_{C^4}\le\epsilon,
\]
and let \(V\) denote the field
\begin{equation}\label{eq:normal-time-field}
 u'=h(u)+c(u)w,\qquad
 w'=w\{1-2h(u)+(k(u)-2c(u))w\}.
\end{equation}
For \(\epsilon\) sufficiently small, this field has a \(C^3\)
projection \(\widehat\pi\) to \(w=0\), satisfying
\(V\widehat\pi=h(\widehat\pi)\) and
\(\widehat\pi(u,0)=u\).  Uniformly on \(\R\times[0,w_0]\),
\[
 \left\|\widehat\pi-u-\tfrac12\log(1-2w)\right\|_{C^3}
 \le C\epsilon.
\]
Writing \(q=e^u\), \(Y=q^2w\), and
\(\pi(q,Y)=e^{\widehat\pi(u,w)}\), we have
\begin{equation}\label{eq:normal-time-footpoint-bounds}
 c_1w\le\log q-\log\pi\le C_1w,\qquad
 1\le q/\pi<2,\qquad \partial_u\widehat\pi\ge\tfrac12.
\end{equation}
There is also a uniformly \(C^3\) function \(J\) in \((u,w)\),
with \(J=O(w)\), for which \(Z=Ye^{-J}\) satisfies
\(V\log Z=1\) when \(Y>0\).  The constants depend on \(w_0\)
and the indicated coefficient bounds.
\end{lemma}

\begin{proof}
Let \(\phi_h^T\) be the axis flow \(u'=h(u)\). Transporting the
backward orbit's base coordinate by this flow defines the footpoint:
\begin{equation}\label{deg:trace-projection}
\begin{split}
 \widehat\pi(u,w)
 &=\lim_{T\to\infty}\phi_h^T(u(-T))\\
 &=u-\int_0^\infty
 D\phi_h^T(u(-T))c(u(-T))w(-T)\,dT,\\
 \pi(q,Y)&=\exp\!\bigl(\widehat\pi(\log q,Y/q^2)\bigr).
\end{split}
\end{equation}
The integral identity follows by differentiation in \(T\).
Write the backward field \(G\), with \(n=1-2h\), \(d_n=k-2c\), as
\[
 \dot U=-h(U)-c(U)W,\qquad
 \dot W=-n(U)W-d_n(U)W^2.
\]
The strip \(0\le W\le w_0\) is invariant. After decreasing \(\epsilon\),
\(n+2d_nW\ge\lambda:=\tfrac12(1-4w_0)>0\) and
\(W(T)\le w e^{-\lambda T}\).
For initial directions \(i,j,k\in\{u,w\}\), put
\(z_i=(U_i,W_i)\), \(M=DG(U,W)\). The variational rows are
\begin{equation}\label{proj:backward-ordinary-jets}
 \begin{aligned}
 \dot z_i&=Mz_i,\\
 \dot z_{ij}&=Mz_{ij}+D^2G[z_i,z_j],\\
 \dot z_{ijk}&=Mz_{ijk}
 +\sum_{\rm cyc}D^2G[z_{ij},z_k]+D^3G[z_i,z_j,z_k].
 \end{aligned}
\end{equation}
The normal diagonal is at most \(-\lambda\), the base diagonal
is \(O(\epsilon)\), and the product of their off-diagonal couplings
is \(O(\epsilon W)\). Integrating the normal row first therefore
leaves an integrable Volterra coupling in the base row. Higher normal
forcings contain \(W\) or a lower normal variation, while pure base
forcings have \(O(\epsilon)\) coefficients. Induction gives
\[
 |D^jU(T)|\le C(1+T)^N e^{C_3\epsilon T},\qquad
 |D^jW(T)|\le C(1+T)^N e^{-(\lambda-C_3\epsilon)T},
 \quad 1\le j\le3.
\]
The scalar axis-flow equations give the same growth bound through order
four. Thus the integrand of \eqref{deg:trace-projection}, through three
ordinary derivatives, is bounded by
\(C(1+T)^Ne^{-(\lambda-C\epsilon)T}\), with an additional \(w\)
in order zero. This proves uniform \(C^3\) dependence; comparison with
\((h,c,k)=(0,1,0)\) gives the stated error. More generally, row \(j\)
uses axis-flow order \(j+1\), giving \(C^4\) under the small \(C^5\)
bounds used in the compact-positive passage. Flow translation gives
\(V\widehat\pi=h(\widehat\pi)\).

Exponentiation gives \(\pi(q,0)=q\), \(V\pi=b(\pi)\), where
\(b(q)=qh(\log q)\). Positivity of the projection integral and
comparison with its unperturbed value prove
\eqref{eq:normal-time-footpoint-bounds}.

Finally set
\begin{equation}\label{eq:normal-time-integrating-coordinate}
 J(q,Y)=\int_0^\infty k(U(T))W(T)\,dT,
 \qquad Z=Ye^{-J}.
\end{equation}
The same majorant gives three ordinary derivatives in \((u,w)\)
and \(J=O(w)\).  Translation along the flow gives \(VJ=kw\).
Since \eqref{eq:normal-time-field} gives
\(V\log Y=1+kw\), we obtain \(V\log Z=1\).
\end{proof}

On a logarithmic corridor with fixed-length collars, extend the same
triple \((h,c,k)\) by smooth retractions with uniformly bounded
fourth derivatives. This preserves the coefficient bounds and
\(V\log Y=1+kw\), independently of corridor length. Each application
checks that its physical passage stays in the unchanged region.

We now relate the divergence of the axis action at a stopped endpoint
to a physical height tending to zero. The next lemma makes this relation
quantitative without dividing by the axis rate at that endpoint.

\begin{lemma}[Physical height from the axis action]
\label{lem:normal-time-action-height}
Use the preceding projection for a field unchanged on
\(Q/2\le q\le2R\), and write \(h(q)=h(\log q)\) in physical
coordinates. Suppose \(b(q)=qh(q)>0\) on \([Q/2,R)\),
\(b\ge b_0>0\) on \([Q/2,Q]\), and
\(\mathcal A(s)=\int_Q^s b(q)^{-1}\,dq\to\infty\) at \(R\).
For \(2Q<R_0<R<\infty\) and every \(R_0\le s<R\), the strong section
\(Y/q^2=w_0\) with footpoint \(s\) has a unique backward crossing of
\(q=Q\), of height \(y(s)>0\). With \(\mathcal D=b(s)\partial_s\),
\begin{equation}\label{eq:normal-time-action-height}
 \log y=-\mathcal A(s)+2\log s+\log U(s)+E_Q(y),\qquad
 E_Q(y)=J(Q,y)+\mathcal A(\pi(Q,y)),
\end{equation}
where \(U,U^{-1}\) and three logarithmic derivatives of \(U\) are
bounded, \(E_Q(0)=0\), and \(\|E_Q\|_{C^3}\le C_Q\).
If also \(b\le P_*\) on \([Q,R_0]\), then
\[
 0<y(s)\le e_0:=C_QR_0^2e^{-(R_0-Q)/P_*}.
\]
Once \(C_Qe_0\) is small, \(y\) decreases to zero at \(R\), and
\begin{equation}\label{eq:normal-time-action-jets}
 \mathcal D\log y=-K,\qquad
 |K-1|+|\mathcal DK|\le C\epsilon+C_Qe_0,\qquad K\ge\tfrac12.
\end{equation}
The constants are independent of \(R\) and of the axis rate at \(R\).
\end{lemma}
\begin{proof}
The footpoint bounds give a unique strong-section abscissa
\(q_s(s)=sU_s(s)\in(s,2s)\), with \(q_s'\asymp1\) and
three bounded logarithmic derivatives of \(U_s\). In backward
normal time both \(w\) and the footpoint decrease. While
\(Q/2\le\pi<R\),
\(u'=h(\pi)+cw+O(\epsilon w)>0\).
The axis reaches \(Q/2\) in finite backward time for each \(s<R\);
there \(q<Q\). Thus the physical orbit crosses \(q=Q\) exactly
once, with \(0<y<w_0Q^2\) and \(Q/2<\pi(Q,y)<Q\).
The entire segment stays in the unchanged corridor.

The first integral \(\log Z-\mathcal A(\pi)\) gives
\eqref{eq:normal-time-action-height}, with the explicit unit
\[
 U(s)=w_0(q_s/s)^2e^{-J(q_s,w_0q_s^2)}.
\]
The fixed-interval lower bound for
\(b\) bounds \(E_Q\) on the whole possible height segment.
Since \(\mathcal D\log(s^2e^{-\mathcal A})=2h-1\le-1/2\),
\(y(s)\le C_QR_0^2e^{-\mathcal A(R_0)}\le e_0\).
Differentiating this identity after the small-height bound gives
\[
 K=\frac{1-2h-\mathcal D\log U}{1-yE_Q'(y)}.
\]
The Euler bounds imply \(\mathcal D\log U=O(\epsilon)\) and
\(\mathcal D^2\log U=O(\epsilon^2)\). Using
\(\mathcal Dy=-Ky\) and the bounded derivatives of \(E_Q\)
proves \eqref{eq:normal-time-action-jets}. Finally the exact
height identity and \(\mathcal A(s)\to\infty\), for each fixed
coefficient fibre with finite \(R\), give \(y(s)\to0\).
\end{proof}

\section{Normalized transitions on the parabolic tails}
\label{sec:source-normalization}

Above the coefficient height scale, the infinity transition consists
of two parabolic tails joined by a bounded connector. Normalization by
each tail's exact value and first variation gives convergent nonlinear
maps, while the common singular power cancels in the ratio of the two
absolute variations. This gives the signed high-height estimate of
Corollary~\ref{rew:cor-high-derivative-signed} for the actual return
\eqref{eq:intro-displacement}.

All constants are uniform on the stated compact parameter sets.
A positive unit is bounded above and away from zero;
\(O_{C^k}\) denotes ordinary phase derivatives through order \(k\)
unless weights are displayed. Choose the finite intrinsic normal-form
regularity before constructing the maps. Exterior observers use
ordinary \(C^3\) bounds on their actual buffered images.

Put \(\alpha _0=s^2a,\ \alpha _1=sp,\ v=\gamma+sX,\ z=s^2Y\).
For \(s>0\), a positive change of time gives
\begin{equation}\label{ann:source-field}
\begin{split}
 X'&=a+pX-Y+A(\gamma+sX)Y-3\gamma X^2-sX^3
                 -C(\gamma+sX)^2Y,\\
 Y'&=-Y\{(1+2\gamma)X+sX^2+Cs(\gamma+sX)Y\}.
\end{split}
\end{equation}
Here \(s\) is a general positive scale. For negative opening one may
take \(s=\sqrt{-\alpha _0}\), whereas normalizing the height of a
reference orbit allows \(a\) to have either sign.

\subsection{The limiting transition}
\label{ann:subsec-source-scattering}

At \(s=\gamma=0\), the equation is
\begin{equation}\label{ann:limiting-source}
 X'=a+pX-Y,\qquad Y'=-XY.
\end{equation}
We call a complete orbit \emph{strict} if it is a graph on \(\R\)
with \(Y>0\) and \(D=Y-a-pX>0\). Positive time then corresponds to
decreasing \(X\). A compact strict family means a compact family of
such graphs, with uniform positive margins for these two inequalities
on each fixed compact \(X\)-interval.

Removing the common quadratic, linear and logarithmic terms at each end
gives finite labels for these complete graphs.

\begin{lemma}\label{ann:lem-renormalized-ends}
On a compact strict family the limits
\begin{equation}\label{eq:renormalized-ends}
 B_\pm=\lim_{X\to\pm\infty}
 \{Y-X^2/2-2pX-(2a-4p^2)\log|X|\}
\end{equation}
exist. They give increasing coordinates on the family and depend
real-analytically on the parameters and the orbit. The corresponding
maps and their inverses have bounded derivatives of every prescribed
finite order on smaller compact strict families.
\end{lemma}
\begin{proof}
Bound \(|a|,|p|\) by \(M\ge1\). On either outward branch the graph
equation \(Y_X=XY/(Y-a-pX)\), with \(r=|X|\), gives
\(Y_r\ge rY/(Y+M(1+r))\). Below \(Y=M(1+r)\),
\((\log Y)_r\ge1/(4M)\) for \(r\ge1\); above \(r=2M\)
this line is a lower barrier. Compactness supplies a common entrance;
after enlarging the cut, \(Y\ge r^2/8\) and \(Y_r=r+O(1)\), uniformly.
Writing \(c=2a-4p^2\) and
\(Y=X^2/2+2pX+c\log|X|+\mathcal B\), direct expansion in a
parabolic corridor gives
\[
 |\mathcal B_X|\le
 C\frac{1+\log|X|+|\mathcal B|}{X^2},
 \qquad |\partial_{\mathcal B}\mathcal B_X|\le C X^{-2}.
\]
Gronwall's inequality gives boundedness and convergence. Beyond a fixed
large cut, the prescribed-limit integral equation is a contraction on a
common complex neighborhood, since its denominator is comparable to
\(X^2\). Holomorphic Picard iteration and finite-interval dependence
give analyticity. The positive scalar first variation and its reciprocal
are bounded, proving the coordinate and inverse assertions.
\end{proof}

Denote the scattering map by
\begin{equation}\label{ann:limiting-scattering-map}
 S_{a,p}(b)=B_-,\qquad b=B_+.
\end{equation}
Its curvature has a particularly simple sign.

\begin{proposition}\label{ann:prop-source-curvature}
On every strict component,
\begin{equation}\label{ann:source-scattering-division}
 S_{a,p}(b)=b-p\mathcal I(a,p,b),\qquad
 \mathcal I_{bb}>0.
\end{equation}
In particular, the curvature of \(S_{a,p}\) has sign
\(-\sgn p\) when \(p\ne0\).
\end{proposition}
\begin{proof}
When \(p=0\), reflection \(X\mapsto-X\) reverses the flow and gives
\(S_{a,0}=\mathrm{id}\). Analytic division in \(p\) defines \(\mathcal I\).
For \(\sigma=\partial_bY>0\) and \(\tau=\partial_b^2Y\), the graph
equation gives
\[
 \sigma_X=-\frac{X(a+pX)}{D^2}\sigma,\qquad
 \tau_X=-\frac{X(a+pX)}{D^2}\tau
              +\frac{2X(a+pX)}{D^3}\sigma^2.
\]
Integrating the equation for \(\tau/\sigma\) from the incoming to the
outgoing end yields
\[
 \frac{S''}{S'}=-2\int_\R\frac{X(a+pX)\sigma}{D^3}\,dX.
\]
To evaluate this integral, use
\[
 D_X=X-p+\frac{X(a+pX)}D,\qquad
 (\sigma/Y)_X=-X\sigma/D^2.
\]
The integral of \((\sigma/D)_X\), followed by that of
\((\sigma/Y)_X\), therefore gives
\begin{equation}\label{ann:source-log-curvature}
 \frac{S''(b)}{S'(b)}
       =-p\int_\R\frac{\sigma(X,b)}{D(X,b)^2}\,dX.
\end{equation}
All boundary terms vanish by \eqref{eq:renormalized-ends}.
Thus
\(\mathcal I_{bb}=S'\int_\R\sigma D^{-2}\,dX>0\), also at \(p=0\)
by continuity.
\end{proof}

\subsection{Normalization at a reference orbit}

Fix a strict reference orbit and label nearby orbits by \(b\), with
reference value \(b_0\). For positive \(s\), let \(h(b)\) be their
height on \(\Sigma_+\), and put
\[
 k(b)=\mathcal C_{\theta,m}\bigl(T_\theta(h(b))\bigr).
\]
Thus \(h-a_\theta-mk=0\) is exactly the periodic-orbit equation.
The maps \(h,k\) include all transitions to the fixed sections. Their
slopes may vanish at different rates, so normalize their values and first
derivatives separately:
\begin{equation}\label{eq:reference-normalization}
 H(b)=b_0+\frac{h(b)-h(b_0)}{h'(b_0)},\quad
 K(b)=b_0+\frac{k(b)-k(b_0)}{k'(b_0)},\quad
 \Phi=K\circ H^{-1}.
\end{equation}
The input \(\beta=H(b)\) is affine in the actual incoming height.
Reversing the sign of the return equation and dividing by \(mk'(b_0)>0\)
gives exactly
\begin{equation}\label{eq:affine-normalized-displacement}
 D(\beta)=c_0+c_1\beta+\Phi(\beta)=0 .
\end{equation}
The constants retain both centers and both slopes. No bound on their
ratio is assumed.

\begin{proposition}[Convergence on the actual sections]
\label{prop:source-normalization}
\label{ann:prop-source-curvature-persistence}
On each compact strict family, \(\Phi\) extends continuously in
phase \(C^3\) to \(s=\gamma=0\), with limit the positive affine
normalization of \(S_{a,p}\) in \eqref{eq:reference-normalization}.
Let \(\Phi^{\rm int}\) denote the normalized transition before the
exterior observers. Its maps and inverses have every prescribed finite
phase order; near \(p=0\), when the division intervals remain strict,
its \(s\)-derivative for \(s\ge0\) and \(p\)-derivative at \(s=0\)
extend in each prescribed \(C^k\), using one extra intrinsic phase
order. Moreover, on common enlarged domains,
\[
 \|\Phi-\Phi^{\rm int}\|_{C^3}=O(s^{2-\nu}),\qquad 0<\nu<1.
\]
Consequently the actual \(s\)-secant extends continuously in
\(C^3\) at zero. At \(p=s=0\), for small \(\gamma\),
\(\Phi=\mathrm{id}\). All assertions are uniform in passive
coefficients and \(0\le m\le m_0\).
\end{proposition}

\begin{proof}
Ordinary ODE dependence controls the compact core. On each outward tail,
set \(X=\sigma r\), \(\sigma=\pm1\); then
\begin{equation}\label{eq:parabolic-tail}
 Y_r=F_\sigma(r,Y)=
 \frac{(1+2\gamma+\sigma sr)rY+
       \sigma Cs(\gamma+\sigma sr)Y^2}
 {\ell_\sigma Y+3\gamma r^2+\sigma sr^3-\sigma pr-a},
 \quad
 \ell_\sigma=1-A(\gamma+\sigma sr)+C(\gamma+\sigma sr)^2 .
\end{equation}
Choose a small raw radius \(\delta\), then a large core radius \(L\),
then the bounds on \(s,\gamma\). The tail ends at \(R_s=\delta/s\),
where uniformly regular maps join \(v=\gamma\pm\delta\) to \(\Sigma_\pm\).

Write
\[
 L_\gamma=1-A\gamma+C\gamma^2,\quad
 W_\gamma=\frac{1-4\gamma}{2L_\gamma},\quad
 \lambda_\gamma=\frac{12\gamma}{1+2\gamma}.
\]
At \(s=p=a=0\), \(W=Y/r^2\) has equilibrium \(W_\gamma\)
and radial exponent \(-2+\lambda_\gamma\).
Fix \(0<\nu<1\), with \(|\lambda_\gamma|<\nu\). A small band about
\(W_\gamma\) is invariant; its logarithmic equation has perturbation
\(O(sr+|p|/r+|a|/r^2)\). Variation of constants gives
\begin{equation}\label{eq:parabolic-corridor}
 W-W_\gamma=O(sr+|p|/r+r^{-2+\nu}),\qquad L\le r\le R_s .
\end{equation}
This includes the logarithmic term for bounded \(a\) and makes the
denominator comparable to \(r^2\), uniformly on the compact strict family
after choosing \(L\).

Let \(Y_*(r)=Y(r,b_0)\) and \(j(r)=Y_b(r,b_0)>0\).
For the rational expression \(F=(\alpha Y^2+\beta Y)/(\ell Y+m)\),
polynomial division gives the exact identity
\begin{equation}\label{tail:rational-division}
 F_{Y^h}=(-1)^h h!\,
 \frac{m(\alpha m-\beta\ell)\ell^{h-2}}
      {(\ell Y+m)^{h+1}},\qquad h\ge2.
\end{equation}
Using the coefficients in \eqref{eq:parabolic-tail} and the corridor gives
\begin{align}
 F_Y-\lambda_\gamma/r
   &=O(s+|p|r^{-2}+r^{-3+\nu}),\label{ann:tail-first-derivative}\\
 |F_{Y^j}|&\le C_j\{sr^{-2j+2}+|\gamma|r^{-2j+1}
                    +|p|r^{-2j}+|a|r^{-2j-1}\},\quad j\ge2.
 \label{ann:tail-higher-derivatives}
\end{align}
At fixed \(r,Y\), their explicit coefficient derivatives satisfy
\begin{equation}\label{tail:rational-parameter-bounds}
 |\partial_s^{\rm exp}F_{Y^h}|\le C_hr^{-2h+2},\qquad
 |\partial_p^{\rm exp}F_{Y^h}|_{s=0}\le C_hr^{-2h+1},\quad h\ge2.
\end{equation}
Integrating \eqref{ann:tail-first-derivative}, using \(sR_s=\delta\), gives
\begin{equation}\label{ann:tail-first-variation}
 c r^{\lambda_\gamma}\le j(r)\le C r^{\lambda_\gamma}.
\end{equation}
The parameter variations, at fixed \(r\), satisfy
\begin{equation}\label{eq:tail-parameter-variations}
 |Y_{*,s}|\le Cr^3,\quad |(\log j)_s|\le Cr,\qquad
 |Y_{*,p}|_{s=0}\le Cr,\quad |(\log j)_p|_{s=0}\le C .
\end{equation}
For the height rows, apply variation of constants to
\(F_s=O(r^2)\), \(F_p|_{s=0}=O(1)\), with propagator
\eqref{ann:tail-first-variation}. For the slope rows,
\((\log j)_{\ell r}=F_{Y\ell}+F_{YY}Y_{*,\ell}\) is respectively
\(O(1)\) and \(O(r^{-2})\).

Subtracting the exact affine part gives
\begin{equation}\label{eq:normalized-difference}
 \eta=(Y-Y_*)/j,\qquad \eta_r=\mathcal N(r,\eta),
\end{equation}
with kernel \eqref{tail:normalized-kernel}. On a bounded
\(\eta\)-interval this remainder can also be written exactly as
\begin{equation}\label{tail:exact-affine-remainder}
 \mathcal N=
 \frac{j\eta^2 m(\alpha m-\beta\ell)}
 {D_*^2(D_*+\ell j\eta)},\qquad D_*=\ell Y_*+m.
\end{equation}
Both denominators are comparable to \(r^2\). The phase and parameter
formula \eqref{deg:det-parameter-phase-derivative}, together with
\eqref{ann:tail-higher-derivatives},
\eqref{tail:rational-parameter-bounds} and
\eqref{eq:tail-parameter-variations}, gives
\begin{align}
 \|\mathcal N\|_{C^k}
 &\le C\{sr^{-2+\nu}+|\gamma|r^{-3+\nu}
                       +|p|r^{-4+\nu}+|a|r^{-5+\nu}\},\label{eq:N-bound}\\
 \|\mathcal N_s\|_{C^k}&\le Cr^{-2+\nu},\qquad
 \|\mathcal N_p\|_{C^k,s=0}\le Cr^{-3+\nu}.
 \label{eq:N-parameter-bound}
\end{align}
The reference and slope are chosen by finite ODE dependence at the fixed
core cut. Thus \(Y_{*,\ell}\), \(j_\ell\) include their motion while
the normalized initial phase stays fixed. Every further phase row adds
\(j/r^2=O(r^{-2+\nu})\), so the bounds include the extra order needed
for the \(C^k\)-valued parameter variation.

These are the integrable majorants of Lemma~\ref{lem:scalar-normalization}.
The moving endpoint contributes
\[
 \|\mathcal N(R_s,\cdot)R_s'\|_{C^k}=O(s^{1-\nu})\longrightarrow0 .
\]
Finite-interval ODE dependence supplies local kernel convergence.
The lemma gives joint convergence, common domains and inverse bounds.

For each exterior map \(R\), its actual interval is
\(z=z_*+d_s\eta\), where \(d_s=s^2j(R_s)=O(s^{2-\nu})\).
Lemma~\ref{phys:lem-c3-adapters} gives on the admitted buffered image
\begin{equation}\label{eq:exterior-normalization}
 \frac{R(z_*+d_s\eta)-R(z_*)}{R'(z_*)d_s}
                  =\eta+O_{C^3}(d_s),
\end{equation}
including inverse and divided closing. With these observers denoted by
\(E_\pm\), the exact composition is
\(\Phi=E_-\circ\Phi^{\rm int}\circ E_+^{-1}\). The intrinsic
\(C^4\) bound and Lemma~\ref{phys:lem-c3-core-transport} give the
claimed \(O_{C^3}(s^{2-\nu})\) error. Its quotient by \(s\) tends to
zero, requiring only phase derivatives of the exterior maps.

At \(p=s=0\) a first integral is
\begin{equation}\label{ann:gamma-first-integral}
\begin{split}
 B_{a,\gamma}(X,Y)
 &=-Y^{-\nu_\gamma}
       \left(\frac{X^2}{2}-\frac{L_\gamma}{1-4\gamma}Y\right)
       -\frac{a}{1+2\gamma}\log Y\,
          \frac{1-e^{-\nu_\gamma\log Y}}{\nu_\gamma\log Y}
       -a\log2,\\
 &\hspace{35mm}\nu_\gamma=\frac{6\gamma}{1+2\gamma},
\end{split}
\end{equation}
with the quotient continued at zero. Its \(Y\)-derivative is positive
on the strict core. Use it as the initial orbit label, shifted to \(b_0\)
at the reference orbit. The same label occurs on the two tails.
Implicit differentiation gives \(Y_{bb}/Y_b=O(r^{-2+\nu})\),
and the analogous higher-order estimates. Therefore both normalized
tail limits are the identity. Together with
\eqref{eq:exterior-normalization} this proves \(\Phi=\mathrm{id}\).
At \(\gamma=0\), the limiting tail labels for general \(p\) are
\eqref{eq:renormalized-ends}, which proves the stated scattering limit.
\end{proof}

During intrinsic parameter integration, hold \(m\) independently
fixed in \([0,m_0]\). Its divided exterior map extends to both signs
of the finite opening; only the actual \(m(\theta)>0\) enters the
cycle equation.

\subsection{Exact division in the perturbation parameters}

On common enlarged strict domains, put
\begin{equation}\label{ann:exact-source-division}
 \begin{gathered}
 \Phi(\beta)=\beta+pU(\beta)+sV(\beta),\qquad
 U=\int_0^1\Phi_p^{\rm int}(a,tp,0,\gamma;\beta)\,dt,\\
 V=\int_0^1\Phi_s^{\rm int}(a,p,ts,\gamma;\beta)\,dt
          +\frac{\Phi-\Phi^{\rm int}}s .
 \end{gathered}
\end{equation}
The two integrals divide successively in \(p\) and \(s\), retaining
the actual exterior remainder. Both secants have continuous bounded
\(C^3\) limits: the first uses only the intrinsic
\(p\)-derivative at \(s=0\); the last quotient is
\(O_{C^3}(s^{1-\nu})\). The strict margins hold throughout both
intrinsic integration intervals, with fixed references and common
inverse images.

\subsection{Cancellation of the absolute tail derivatives}
\label{ann:subsec-high-source}

Use the actual affine incoming coordinate of Lemma~\ref{lem:closing}
and write
\[
 \Pi(x)=\mathcal C_{\theta,m}\bigl(T_\theta(a_\theta+x)\bigr),
 \qquad \Delta(x)=x-m\Pi(x).
\]
The fixed raw ports, their regular bridges, and the divided closing map
are all included in this same transition.

\begin{proposition}\label{sharp:upper-high}
On any fixed compact strict source family with its two outward tails,
the complete transition satisfies \(0<\Pi'(x)\le C\), uniformly for
small positive \(s\), small \(\gamma\), and \(0\le m\le m_0\).
The constant may depend
on the compact family. This includes a small signed \(a\)-block about
\(a=p=0\), with central height near one.

On the negative-opening source family \(a=-1\), there are constants
\(K_H,z_*,C_H>0\) such that, for either sign of small \(\gamma\),
\[
 Y_0>K_H(1+p^2),\qquad 0<s^2Y_0<z_*
 \quad\Longrightarrow\quad 0<\Pi'(x)\le C_H.
\]
Here \(C_H\) is independent of \(p\) and of the tested phase point.
After one common multiplier shrink, \(\Delta'>1/2\) on this whole
high-height interval.
\end{proposition}
\begin{proof}
Fix a compact strict family and a core cut \(L_t\). With
\(y_c=Y(0)\), ordinary first variation on \([-L_t,L_t]\) gives
\[
 A_\sigma=\frac{\partial Y(\sigma L_t;y_c)}{\partial y_c}
 \in[c_A,C_A],\qquad \sigma=\pm.
\]
The coefficient scale \(s\) is held fixed in this differentiation.
Equations \eqref{ann:tail-first-derivative} and
\eqref{ann:tail-first-variation}, before affine normalization, give on
each outward tail
\[
 J_\sigma(R)=(R/L_t)^{\lambda_\gamma}U_\sigma(R),\qquad
 \lambda_\gamma=\frac{12\gamma}{1+2\gamma},\qquad
 0<c\le U_\sigma,U_\sigma^{-1}\le C.
\]
At ports \(v=\gamma\pm d_H\), \(R=d_H/s\), the raw variations are
\begin{equation}\label{sharp:upper-tail-ratio}
 \frac{dz_\sigma}{dy_c}
 =s^{2-\lambda_\gamma}(d_H/L_t)^{\lambda_\gamma}
                                      A_\sigma U_\sigma(R),\qquad
 T'(z_+)=\frac{A_-U_-(R)}{A_+U_+(R)}.
\end{equation}
Here \(T\) is the passage between the raw ports. The common power
cancels; distinct fixed radii add only the bounded factor
\((d_-/d_+)^{\lambda_\gamma}\). Regular bridges are derivative units, and
\(\mathcal C_{\theta,m}'=f_\theta'(md_\theta)d_\theta'\) is a
uniform positive unit by Lemma~\ref{lem:closing}. Thus the derivative
of the complete transition in the actual incoming coordinate is
bounded on the fixed compact strict family.

For the signed block \(|a|,|p|\le1/8\),
\(3/4\le y_c\le5/4\), the half-graphs satisfy
\[
 3/4+r^2/4\le Y(r)\le5/4+r^2,\qquad
 Y-a-\sigma pr\ge39/64.
\]
These follow by scalar comparison; below \(Y/r^2=1/4\) the
logarithmic vector points upward, and above \(3/4\) it points
downward beyond a common fixed cut. Thus this is a compact strict
family with uniform parabolic entrances.

For the negative-opening high-height range, freeze
\(q=Y_0^{-1/2}\) at a selected point and rescale the source field to
\[
 (\bar a,\bar p,\bar\gamma,\bar s)
   =(-Y_0^{-1},pY_0^{-1/2},\gamma,sY_0^{1/2}),\qquad y(0)=1.
\]
Large \(K_H\), followed by small \(z_*\) and a coefficient shrink,
places these parameters in one compact strict source block about
\(y=1+x^2/2\). Apply the preceding bound with scale \(\bar s\).
The same constant \(C_H\) then works at every tested point, and
\(mC_H<1/2\) gives \(\Delta'>1/2\).
The scale is frozen during the phase derivative, with no bound on
\(\gamma\log\bar s\).
\end{proof}

\subsection{Extension through the height cap}

The derivative bound joins the low-height count to the fixed raw cap.
The overlap condition in the next lemma ensures that this continuation
introduces no additional critical point.

\begin{lemma}[Extension through the height cap]
\label{sharp:upper-height-splice}
On one ordered physical carrier, suppose a connected lower interval
has a prescribed critical-point bound and overlaps a bounded-height
interval where the complete transition derivative is uniformly bounded.
Write its scaled height as \(\eta=\tau Y_0\), \(\tau>0\), and
suppose these two intervals cover \(0<\eta\le\eta_{\max}\).
If Proposition~\ref{sharp:upper-high} applies above
\(Y_0=K_H(1+p^2)\) and
\begin{equation}
 \tau K_H(1+p^2)<\eta_{\max},
 \label{sharp:upper-height-splice-condition}
\end{equation}
then a common multiplier shrink adds no critical point on the
complement of the lower interval, up to the captured raw-height cap.
\end{lemma}
\begin{proof}
The bounded derivative gives \(F'>0\) on the bounded-height part
after shrinking \(m\). Put \(Z=z_*/s^2\). If
\(Z\le K_H(1+p^2)\), the displayed inequality puts every captured
height below \(\eta_{\max}/\tau\). Otherwise the intervals
\((0,\eta_{\max}/\tau]\) and \((K_H(1+p^2),Z)\) overlap
below the cap and cover it. The high-height estimate gives \(F'>0\)
on the added interval. These are restrictions of the same completed
map; scalar order retains all inputs between captured roots.
\end{proof}

\section{Degenerating saddle transitions: action coordinates}
\label{sec:negative-action}

We use an action coordinate to describe a large negative
source-energy interval. Let \(\beta\) be
defined by \(Y+\log Y=X^2/2+\beta\) for the unperturbed negative source.
The passage consists of two parabolic tails joined by a thin region near
the equator. Its limiting transition is determined by an action integral,
but its curvature vanishes on the reversible coefficient set. Absolute
convergence alone does not control the curvature near that set. We estimate
the nonlinear part of the actual transition relative to the distance to
that set. Keeping the section maps and their inverses in physical
coordinates allows this estimate to be used in
\eqref{eq:intro-displacement}.

\subsection{The equation and the section coordinates}

The local scales below are related to the raw coefficients in
Section~\ref{sec:coefficient-regimes}.

Fix a scale \(t>0\) on each coefficient fibre, keeping it fixed as
the orbit varies. Start with \eqref{ann:source-field} at \(a=-1\),
write \(p=t\Lambda_1\), \(s=t^3\Lambda_2\),
\(\gamma=t^2k\), and put \(\xi=tX\), \(\eta=t^2Y\).
After this phase scaling and a positive change of time, the same
form \eqref{ann:source-field} holds in $(\xi,\eta)$ with the new
coefficients
\begin{equation}\label{neg:source-dictionary}
 (a_{\rm new},p_{\rm new},\gamma_{\rm new},s_{\rm new})
 =(-t^2,t^2\Lambda_1,t^2k,t^2\Lambda_2).
\end{equation}
Thus \(p_{\rm new}=tp\), \(s_{\rm new}=s/t\), and the equation is
\begin{equation}\label{neg:field}
\begin{split}
 \dot\xi&=-\eta-t^2P_\Lambda(\xi)
                +At^2L(\xi)\eta-Ct^4L(\xi)^2\eta,\\
 \dot\eta&=-\eta\{\xi+t^2\xi(2k+\Lambda_2\xi)
                         +Ct^4\Lambda_2L(\xi)\eta\},\\
 P_\Lambda(\xi)&=1-\Lambda_1\xi+3k\xi^2+\Lambda_2\xi^3,
 \qquad L(\xi)=k+\Lambda_2\xi .
\end{split}
\end{equation}
The passive coefficients \(A,C\) range over a fixed compact
neighborhood of \((1,1)\). Physical tails have \(\Lambda_2>0\);
at \(\Lambda_2=0\) we will use limits of normalized maps.

For the unperturbed negative source,
\(Y+\log Y=X^2/2+\beta\). Hence its rescaled energy is
\begin{equation}\label{neg:affine-energy}
 \widehat b=t^2\beta+t^2\log(2t^2),\qquad
 \eta+t^2\log(2\eta)=\xi^2/2+\widehat b.
\end{equation}
Energy-label differences are exactly \(t^2\) times source-label
differences; the action integral below keeps these transverse labels.

We specify the intervals on which the estimate will be used. Choose
compact intervals
\[
 J\Subset\operatorname{int}\widehat J_+\Subset(-\infty,0),
 \qquad \widehat J_-\Subset(-\infty,0),
\]
and a compact coefficient set. For \(b\in\widehat J_+\), put
\(x=\sqrt{-2b}\) and define \(r>0\) by
\begin{equation}\label{neg:action-equation}
 A_-(r;\Lambda)=A_+(x;\Lambda),\qquad
 A_+(x;\Lambda)=\int_0^x\frac{u}{P_\Lambda(u)}\,du,
 \quad A_-(r;\Lambda)=\int_0^r\frac{u}{P_\Lambda(-u)}\,du.
\end{equation}
We require that the entry and exit lie in fixed compact subintervals
of \((0,\infty)\), that \(-r^2/2\) lies in
\(\operatorname{int}\widehat J_-\) with a fixed margin, and that
\begin{equation}\label{neg:rootfree-margin}
 P_\Lambda\ge p_0>0
\end{equation}
on a common neighborhood of every traversed segment \([-r,x]\).
The common neighborhood extends to both sides of the limiting exit.
Monotonicity of \(A_-\) then gives comparison exit cuts with opposite
action differences. Enlarged intervals provide the smaller common
images needed for composition and inversion.

Choose a fixed \(L\) beyond the entry and exit intervals. The positive
height pieces between the two sections \(\xi=\pm L\) and the lips
are required to lie in fixed transverse neighborhoods. Each outward
tail lies in a parabolic corridor and reaches the prescribed raw
section, followed by its actual regular transitions. The divided
finite closing map is included on its side. Its multiplier \(m\)
is held fixed as an independent passive parameter in
\(0\le m\le m_0\) during parameter division, as in
Section~\ref{sec:source-normalization}; the cycle equation uses its
actual value \(m>0\). These requirements concern the whole enlarged
input and output intervals.

There is a useful energy coordinate for all small signed \(k\).
Put
\begin{equation}\label{neg:ideal-energy}
\begin{split}
 \ell&=1-At^2k+Ct^4k^2,\qquad
 \nu_k=\frac{6t^2k}{1+2t^2k},\\
 \mathcal B_{t,k}(\xi,\eta)
 &=-\frac12\left[
 \eta^{-\nu_k}\left(\xi^2+\frac1{3k}
             -\frac{2\ell\eta}{1-4t^2k}\right)-\frac1{3k}\right].
\end{split}
\end{equation}
The apparent pole is removable, with
\begin{equation}\label{neg:ideal-energy-zero}
 \mathcal B_{t,0}=\eta+t^2\log\eta-\xi^2/2 .
\end{equation}
Direct differentiation of \eqref{neg:field} gives
\(\dot{\mathcal B}_{t,k}=0\) when
\(\Lambda_1=\Lambda_2=0\). Shrink the parameters so that
\(1+2t^2k\), \(1-4t^2k\), and \(\ell\) are positive.
At the two fixed cuts use the separate coordinates
\begin{equation}\label{neg:cut-labels}
 b_\pm=\mathcal B_{t,k}(\pm L,\eta)+t^2\log2 .
\end{equation}
Here \(\eta\) stays in a compact positive interval and
\(\partial_\eta\mathcal B_{t,k}\) is bounded above and away
from zero. These coordinates and their inverses are smooth in
\(t^2,k,b_\pm\). They agree with \eqref{neg:affine-energy} at
\(k=0\); away from that set, their relation is part of the actual
section maps.

Let \(S_{t,\lambda,k}\), \(\lambda=(\Lambda_1,\Lambda_2)\),
be the passage from \(b_+\) to \(b_-\), and let
\(\Gamma_\pm\) be the maps from these separate coordinates to
the actual input and divided output heights. For references
\(b_\pm^*\) in the enlarged intervals, set
\begin{equation}\label{neg:tail-normalizations}
\begin{split}
 c_\pm&=\Gamma_\pm'(b_\pm^*)>0,
 \qquad a_\pm=\Gamma_\pm(b_\pm^*)-c_\pm b_\pm^*,\\
 \phi_\pm(b)&=b_\pm^*
          +\frac{\Gamma_\pm(b)-\Gamma_\pm(b_\pm^*)}{c_\pm},
 \qquad \Gamma_\pm=a_\pm+c_\pm\phi_\pm .
\end{split}
\end{equation}
On the respective raw images
\(\mathcal I_\pm=\Gamma_\pm(\operatorname{int}\widehat J_\pm)\)
and normalized images
\(\mathcal U_\pm=\phi_\pm(\operatorname{int}\widehat J_\pm)\), put
\begin{equation}\label{neg:exact-composition}
 \Pi_{t,\Lambda}=\phi_-\circ S_{t,\lambda,k}\circ\phi_+^{-1}
       :\mathcal U_+\longrightarrow\mathcal U_- .
\end{equation}
If necessary the domain is restricted to the smaller input interval
whose image has the prescribed output margin. The estimates below
ensure that \(J\subset\mathcal U_+\) for small parameters.

To see the role of both affine factors, write the actual cycle
equation as \(h-a_{\rm f}-mh_{\rm out}=0\), where
\(a_{\rm f}=a_\theta=C_\theta(0)\) and
\(h_{\rm out}=\mathcal C_{\theta,m}(T_\theta(h))\) is the divided
output height from Lemma~\ref{lem:closing}. Substitution of
\eqref{neg:tail-normalizations} gives, without a change of zeros,
\begin{equation}\label{neg:physical-displacement}
 D(u)=a_+-a_{\rm f}-ma_-+c_+u-mc_-\Pi_{t,\Lambda}(u)=0 .
\end{equation}
Thus \(u\) is affine in the actual incoming height. The two
positive first derivatives may tend to zero at different rates;
neither their equality nor a bounded ratio is required.

The zero count therefore requires curvature control of the completed map
\eqref{neg:exact-composition}. Lemma~\ref{neg:lem-completed-curvature}
expresses this curvature with both section contributions; the relative
comparison below estimates them separately.

\subsection{A relative estimate near the reversible coefficients}

For a fixed reference \(b_*\in J\), write
\(\mathfrak T f=f-f(b_*)-(b-b_*)f'(b_*)\).
On a fixed compact interval the quotient norm satisfies
\[
 \|[f]\|_{C^3/\operatorname{Aff}}
      \asymp\|\mathfrak T f\|_{C^3}\asymp\|f''\|_{C^1} .
\]
This quotient isolates the derivatives needed for the curvature count.
The affine coefficients remain in the actual displacement
\eqref{neg:physical-displacement}.

\begin{theorem}[Relative convergence on the actual sections]
\label{neg:thm-relative}
Fix the enlarged intervals, the entry and exit margins, the tail and
regular-transition neighborhoods described above. There exist
\(t_0,\lambda_0,k_0,C>0\), depending only on these data and the
compact passive coefficient set, such that the following holds.
For
\[
 0<t\le t_0,\qquad
 0<\rho_{\rm id}:=(\Lambda_1^2+\Lambda_2^2)^{1/2}\le\lambda_0,
 \qquad \Lambda_2\ge0,\qquad |k|\le k_0,
\]
put \(\theta_i=\Lambda_i/\rho_{\rm id}\). On the fixed smaller
interval \(J\), the map \eqref{neg:exact-composition} satisfies
\begin{equation}\label{neg:relative-estimate}
 \left\|\frac{[\Pi_{t,\Lambda}-\mathrm{id}]}{\rho_{\rm id}}
           -\theta_1[U_0]-\theta_2[V_0]\right\|_{C^3/\operatorname{Aff}}
 \le C\{t^2(1+|\log t|)+\rho_{\rm id}+|k|\},
\end{equation}
where
\begin{equation}\label{neg:limiting-rows}
 U_0(b)=-\frac23(-2b)^{3/2},\qquad
 V_0(b)=\frac25(-2b)^{5/2} .
\end{equation}
The coefficient \(k\) is independently signed. The constants are
uniform in the two separate reference values, held fixed during
parameter integration, and in \(0\le m\le m_0\). No positive
lower bound on \(\theta_2\) is needed. At \(\Lambda_2=0\) the
assertion concerns the one-sided normalized limits.
\end{theorem}

\begin{proof}
On each outward tail use \eqref{eq:parabolic-tail} with dictionary
\eqref{neg:source-dictionary} and the separate initial label
\eqref{neg:cut-labels}. With \(L\) fixed before shrink, the enlarged
windows have \(\eta\asymp r^2\) and denominator \(\asymp r^2\)
throughout division. Thus the tail proof of
Proposition~\ref{prop:source-normalization} applies with its exact
\(j\asymp r^{\lambda_\gamma}\), kernel
\eqref{eq:normalized-difference}, and bounds
\eqref{eq:N-bound}--\eqref{eq:N-parameter-bound} through one extra
phase order. It gives bounded \(C^3\)-valued intrinsic \(s\)-rows
and \(p\)-rows at \(s=0\), including the moving-end error
\(O_{C^3}(s^{1-\nu})\). These differentiations concern only the tails.

At \(p=s=0\), \eqref{neg:ideal-energy} is a first integral;
implicit differentiation gives identity normalized intrinsic limits.
As in \eqref{ann:exact-source-division}, integrate the intrinsic tails
first in \(p\) at \(s=0\), then in \(s\) at fixed \(p\), holding
the two reference values separately fixed. Since
\(\partial_{\Lambda_1}=t^2\partial_p\)
and \(\partial_{\Lambda_2}=t^2\partial_s\), the intrinsic maps and
inverses differ from the identity by
\(Ct^2(|\Lambda_1|+\Lambda_2)\) in \(C^3\).

At the tested coefficients, each raw gain for \(\Lambda_2>0\) satisfies
\[
 d\le C(t^2\Lambda_2)^{2-\nu}
       \le Ct^2\rho_{\rm id},\qquad 0<\nu<1.
\]
The ordinary observers, including divided closing, therefore add
\(O_{C^3}(d)\) by Lemma~\ref{phys:lem-c3-adapters}. Thus
\begin{equation}\label{neg:single-tail-estimate}
 \|\phi_\pm-\mathrm{id}\|_{C^3}
 +\|\phi_\pm^{-1}-\mathrm{id}\|_{C^3}
       \le Ct^2(|\Lambda_1|+\Lambda_2).
\end{equation}
At \(\Lambda_2=0\) use the one-sided normalized limit. Inverses
are on fixed smaller common images; no exterior coefficient derivative
is needed.

For the core, work in a small independent \((\lambda,k)\)-block with
strict core margin and construct \(S_{t,\lambda,k}\) and its first
\(\lambda\)-derivatives in \(C_b^3\), keeping \(t>0,k\) fixed.
Choose an integer \(M>20\), and cut each lip at
\(y=\eta/t^2=t^M\). On a lip use
\(H=\eta-\xi^2/2\) and \(\ell_\eta=\log\eta\).
Equation \eqref{neg:field} gives exactly
\begin{equation}\label{neg:lip-equation}
 \frac{dH}{d\ell_\eta}=t^2\Phi_\pm,
 \qquad
 \Phi_\pm=
 \frac{\xi\eta(2k+\Lambda_2\xi+AL-Ct^2L^2)
              -\xi P_\Lambda+Ct^2\Lambda_2L\eta^2}
      {\xi\{1+t^2(2k+\Lambda_2\xi)\}
                              +Ct^4\Lambda_2L\eta} .
\end{equation}
The signed root \(\xi^2=2(\eta-H)\) is bounded away from zero
on each lip, as is the denominator. Thus the required derivatives
of \(\Phi_\pm\) in \(H,\lambda,k\) are bounded. The logarithmic
clock has length \(O(1+|\log t|)\), giving a change
\(O(t^2(1+|\log t|))\) through total order five; the positive-height
pieces to \(\pm L\) add \(O(t^2)\). Equation~\eqref{neg:ideal-energy}
therefore gives
\begin{equation}\label{neg:lip-label-error}
 b=H+O_{C^5}\bigl(t^2(1+|\log t|)\bigr),
\end{equation}
At \(k=0\), the error is \(t^2\log(2\eta)\), largest at
\(\eta=t^{M+2}\); the removable continuation in \(k\) has the same
mixed bound through order five. The positive entry is
\(\sqrt{-2b}+O_{C^5}(t^2(1+|\log t|))\); the negative exit has
the corresponding estimate about \(-\sqrt{-2b}\).

In the thin core put \(Q=t^2\log y\). Its exact graph equation is
\begin{equation}\label{neg:thin-core}
 Q_\xi=
 \frac{\xi(1+2t^2k+t^2\Lambda_2\xi)
               +Ct^6\Lambda_2(k+\Lambda_2\xi)y}
      {P_\Lambda(\xi)+\ell_\Lambda(\xi)y}
 =\frac{\xi}{P_\Lambda(\xi)}+R,
 \quad
 \ell_\Lambda=1-At^2L+Ct^4L^2 .
\end{equation}
Until the first return to \(y=t^M\), one has \(0<y\le t^M\).
The root-free margin and direct differentiation give
\begin{equation}\label{neg:core-mixed-estimate}
 |\partial_Q^j\partial_\lambda^\alpha R|
       \le C\{t^2+t^{M-2j}\},\qquad j+|\alpha|\le5,
\end{equation}
with the same bounds for the required \(\xi\)-derivatives.
The negative powers come only from differentiating
\(y=e^{Q/t^2}\). For dependence on an open domain, replace this
exponential above the lip by
\(t^M\chi((Q-Q_{\rm lip})/t^2)\), where
\(Q_{\rm lip}=Mt^2\log t\), \(\chi(z)=e^z\) for
\(z\le1\), and \(\chi\) is smooth, positive and constant for
\(z\ge2\). This changes no first-exit orbit, preserves
\eqref{neg:core-mixed-estimate}, and keeps the denominator positive.

The two comparison exit cuts from \eqref{neg:action-equation}
have opposite action differences. Starting at the positive lip,
\(Q\) decreases, stays a fixed amount below the lip near
\(\xi=0\), and increases strictly on the negative side away from
zero. The comparison cuts consequently give a unique actual first
exit. At that exit,
\(Q_\xi=-r/P_\Lambda(-r)+o(1)\) is bounded away from zero.
The differentiated integral equation for \eqref{neg:thin-core},
followed by this implicit exit equation, therefore has uniform
derivatives through total order five. With the two lip maps it gives
\begin{equation}\label{neg:compact-mixed-convergence}
\begin{split}
 &\|S_{t,\lambda,k}-S_{0,\lambda,k}\|_{C_b^3}
  +\|\partial_\lambda S_{t,\lambda,k}
                -\partial_\lambda S_{0,\lambda,k}\|_{C_b^3}\\
 &\hspace{20mm}\le C\{t^2(1+|\log t|)+t^{M-10}\},
 \qquad S_{0,\lambda,k}(b)=-r^2/2 .
\end{split}
\end{equation}
No derivative with respect to \(t\) at zero is asserted.

The exact energy and even \(P_{0,k}\) give
\(S_{t,0,k}=S_{0,0,k}=\mathrm{id}\). At \(\lambda=k=0\),
\eqref{neg:action-equation} gives
\[
 \partial_{\Lambda_1}S_{0,0,0}=-2x^3/3=U_0(b),\qquad
 \partial_{\Lambda_2}S_{0,0,0}=2x^5/5=V_0(b).
\]
Its derivatives are smooth on the compact root-free block. The
identity on the coefficient ideal permits the exact division
\[
 S_{t,\lambda,k}-\mathrm{id}
      =\int_0^1\partial_\lambda S_{t,v\lambda,k}\,\lambda\,dv .
\]
Using \eqref{neg:compact-mixed-convergence} in this integral, and
then the smoothness of the limiting action, gives
\begin{equation}\label{neg:relative-core}
 \|S_{t,\lambda,k}-\mathrm{id}-\Lambda_1U_0-\Lambda_2V_0\|_{C^3}
  \le C\rho_{\rm id}
          \{t^2(1+|\log t|)+\rho_{\rm id}+|k|\} .
\end{equation}
Thus division by \(\rho_{\rm id}\) precedes the limit; the even
coefficient \(k\) creates no additional row.

The same construction bounds \(S\) in \(C^4\). Apply
Lemma~\ref{phys:lem-c3-core-transport} with
\eqref{neg:single-tail-estimate} to obtain
\begin{equation}\label{neg:composition-estimate}
 \|\Pi_{t,\Lambda}-S_{t,\lambda,k}\|_{C^3}
                       \le Ct^2\rho_{\rm id} .
\end{equation}
Only the core needs the extra phase derivative for evaluation at
\(\phi_+^{-1}\). Combining with \eqref{neg:relative-core} proves
\eqref{neg:relative-estimate}.
\end{proof}

The outgoing tail must first be estimated in its own intrinsic label. If
an outgoing graph is written as \(\eta(\xi,S(b))\), then
\[
 \partial_b^2\eta(\xi,S(b))
    =\eta_{BB}(\xi,S(b))(S')^2+\eta_B(\xi,S(b))S'' .
\]
The second term need not decay on that tail. Estimating the map in
\(B\) first and then composing retains this contribution.

\subsection{Zeros of the completed displacement}

The action transition has a curvature formula useful beyond the
small coefficient neighborhood. Differentiating
\eqref{neg:action-equation}, with \(x=\sqrt{-2b}\), gives
\begin{equation}\label{neg:action-curvature}
 S_\Lambda'=
     \frac{P_\Lambda(-r)}{P_\Lambda(x)}>0,
 \qquad
 \frac{S_\Lambda''}{S_\Lambda'}
   =\frac{x+r}{P_\Lambda(x)}
                 \left(3\Lambda_2-\frac{\Lambda_1}{xr}\right),
 \quad S_\Lambda=-r^2/2 .
\end{equation}
The terms containing \(k\) cancel. Moreover
\begin{equation}\label{neg:entry-exit-monotonicity}
 \frac{d(xr)}{db}
      =-\frac r x-\frac{xP_\Lambda(-r)}{rP_\Lambda(x)}<0 .
\end{equation}
Thus, away from \(\lambda=0\), every possible curvature zero is
simple and has the same crossing orientation.

\begin{corollary}\label{neg:cor-rootfree}
On a compact family satisfying the section, tail, and root-free
margins above, assume \(\Lambda_2\ge0\) and that the physical
coefficients satisfy \(k\to0\) whenever
\(\rho_{\rm id}\to0\), uniformly on the family. For example,
the condition \(|k|^3\le C\Lambda_2^2\) suffices.
There is a common \(t_0>0\) such that every actual completed
displacement \eqref{neg:physical-displacement}, for
\(0<t\le t_0\), has at most three isolated zeros on its input
interval, counted with multiplicity.

The same conclusion holds for connected coefficient-dependent
intervals \(J_\Lambda\) if their coefficient--label graph has
compact closure and the enlarged input and output intervals, both
exit margins, and all actual transition neighborhoods have the
same uniform margins.
\end{corollary}
\begin{proof}
At \(\rho_{\rm id}=0\), the exact first integral and the separate
normalized tail limits give \(\Pi=\mathrm{id}\). The displacement
is therefore affine: it has at most one zero unless it vanishes
identically, in which case it has no isolated zero. This coefficient
boundary is understood as a normalized limit; actual raw passages
have \(\Lambda_2>0\).

Near the coefficient ideal, put \(x=\sqrt{-2b}\). The limiting
rows satisfy
\[
 U_0''=-2/x,\qquad V_0''=6x,\qquad
 \frac{(\theta_1U_0+\theta_2V_0)''}{2/x}
                   =-\theta_1+3\theta_2x^2 .
\]
If the last expression vanishes, compactness of the positive
\(x\)-interval and \(\theta_1^2+\theta_2^2=1\) give a
positive lower bound on \(\theta_2\). Its derivative in \(b\)
is \(-6\theta_2\). When \(\theta_2\) is smaller than a fixed
positive constant, the expression is instead bounded away from
zero. This proves a uniform alternative of no curvature zero or
one simple downward crossing. The relative \(C^3\) estimate
preserves that alternative after a common shrink of
\(t,\rho_{\rm id},k\).

On the complementary compact coefficient annulus, the tail estimate
still applies as \(p=t^2\Lambda_1\), \(s=t^2\Lambda_2\),
\(\gamma=t^2k\) tend uniformly to zero. The strict core margin permits
the same lip and first-exit argument, giving absolute \(C^3\)
convergence to \(S_\Lambda\) without division through the ideal. Hence
\(\|[\Pi-S_\Lambda]\|_{C^3/\operatorname{Aff}}\to0\).
Equations \eqref{neg:action-curvature} and
\eqref{neg:entry-exit-monotonicity} give the same uniform crossing
alternative on this compact annulus. Finitely overlapping the two
coefficient regions gives one bound on \(t\).

Lemma~\ref{sharp:upper-curvature-count} gives at most one curvature
zero and three displacement zeros since \(D''=-mc_-\Pi''\).
The affine relation \eqref{neg:physical-displacement} preserves
multiplicity. Compactness of the coefficient--label graph supplies
the same constants fibrewise on \(J_\Lambda\).
\end{proof}

\section{Small characteristic ratios, finite middle passages, and fixed exit}
\label{sec:small-characteristic-ratios}

At degree loss, the two incident saddle exponents vanish at different
rates. We compare their small-power maps using an action coordinate,
exact local invariants and backward shooting from a fixed exit. The
degree-loss section joins these descriptions on their actual raw
receivers and includes the divided closing.
\subsection{Two-saddle comparison and action matching}
\label{sec:two-saddle-comparison}

On this coefficient face the two incident saddles lie on different
scales, and both characteristic ratios tend to zero. Their local maps
must be compared after the intervening nonlinear flight. We first derive
its scalar action equation and the actual local section coordinates.
Section~\ref{subsec:deg-degree-drop} then constructs their physical
receiver and follows the passage over the full action range.

\subsection{The two saddle scales}

Let \(0<\kappa,\tau\ll1\), \(\delta=\kappa\tau\). In coordinates
which put the negative saddle at \(u=-1\), the other two equatorial
roots are \(u=a\) and \(u=c\). Write
\begin{equation}\label{fx:pd}
 p_{\rm d}=-\mathcal R(u+1)(u-a)(u-c),
\end{equation}
where
\begin{equation}\label{fx:coefficient-graph}
 a=\kappa A_{\rm d},\qquad a+c=1-3\nu .
\end{equation}
The quantities \(A_{\rm d},\mathcal R,c,c-a,1-3\nu\) and their
reciprocals belong to fixed positive compact sets; \(A,C\) vary in
fixed compact sets. Thus \(a\) tends to zero while \(-1\) and \(c\)
remain separated. The vector field has the form
\[
 \dot u=\delta p_{\rm d}-Kw,\qquad
 \dot w=-w(L+Ew),
\]
with
\begin{equation}\label{fx:coefficients}
\begin{aligned}
 K&=1-\delta\mathcal R A(u+\nu)
              +C\delta^2\mathcal R^2(u+\nu)^2,\\
 L&=u\{1+\delta\mathcal R(u+2\nu)\},\qquad
 E=C\delta^2\mathcal R^2(u+\nu).
\end{aligned}
\end{equation}
On the intervening axis segment put \(h=-\delta p_{\rm d}>0\).
Positive time runs from the cut \(u_+=a(1-\eta_+)\) to
\(u_-=-1+\eta_-\), with fixed \(0<\eta_\pm<1\). The orbit equation is
\begin{equation}\label{fx:graph}
 w_u=\frac{w\{L(u)+E(u)w\}}{h(u)+K(u)w}.
\end{equation}

The tangential eigenvalues at \(a\) and \(-1\) are
\(\delta\ell_+\) and \(-\delta\ell_-\), whereas their normal
eigenvalues are \(-b_+\) and \(b_-\), where
\[
\begin{aligned}
 \ell_+&=\mathcal R(1+a)(c-a),&
 b_+&=a\{1+\delta\mathcal R(a+2\nu)\},\\
 \ell_-&=\mathcal R(1+a)(1+c),&
 b_-&=1-\delta\mathcal R(1-2\nu).
\end{aligned}
\]
Define the small characteristic ratios used in the local maps by
\begin{equation}\label{fx:exponents}
 \chi=\delta\ell_+/b_+,\qquad
 \beta=\delta\ell_-/b_-,\qquad q=\beta/\chi .
\end{equation}
Thus \(\chi\) is the exponent for the inverse positive Dulac map,
whereas \(\beta\) is the exponent for the forward negative Dulac map.
The normal rate at the positive saddle is of order \(\kappa\),
while that at the negative saddle is bounded away from zero. Hence
\(\chi\) is of order \(\tau\), \(\beta\) is of order
\(\kappa\tau\), and \(q=\kappa Q_{\rm d}\) with
\(Q_{\rm d}\) a positive unit. In particular,
\begin{equation}\label{fx:log-smallness}
 \beta|\log a|+q|\log a|\longrightarrow0.
\end{equation}
The inequality \(\beta\ll\chi\) will determine the sign in the
comparison of the two Dulac maps. To transfer that sign to the return,
we also need the nonlinear transition between their local sections.

\subsection{A transition in accumulated action}

Writing \(A_m=L/h\) separates the linear normal equation in
\eqref{fx:graph} from its nonlinear part. The following calculation
also applies to general coefficients for which \(h>0\):
\begin{equation}\label{eq:comparison-middle}
 w_u=A_m(u)w+\frac{E(u)-A_m(u)K(u)}{h(u)+K(u)w}\,w^2,\qquad h>0 .
\end{equation}
Let \(J'=A_mJ,\ J(u_+)=1\), and write \(h_-=h(u_-),J_-=J(u_-)\).
For a fixed \(\lambda>0\), put
\[
 \rho=\frac{h_-J}{J_-h},\qquad
 w=\lambda h\rho V .
\]
The factor \(J/J_-\) removes the linear amplification between the
cuts, and \(\rho(u_-)=1\). This normalization isolates the remaining
nonlinear change of height. In the saddle configuration above the
amplification is exponentially large in \(1/\beta\); its exact value
is given by the partial fractions in \eqref{fx:partial-fractions}.
Direct substitution retains the denominator in the equation for \(V\):
\begin{equation}\label{eq:comparison-action}
 dV=-\frac{\lambda KV^2}{1+\lambda K\rho V}\,d\rho
 +\frac{\lambda\rho(E-Kh'/h)V^2}{1+\lambda K\rho V}\,du .
\end{equation}

For bounded \(V=V(u_+)\), an early negative cut gives the actual
section coordinates before the nonlinear amplification layer. Larger
values of \(V\) require a different description of the same passage.
Section~\ref{sec:fixed-exit-comparison}
uses the coordinate \(\Theta=\beta\log V\) on a bounded nonnegative
interval. For a positive limiting \(\Theta\) in that interval,
\(V\) is exponentially large in \(1/\beta\), and the orbit leaves
the small-height region before the negative saddle. Its height on a
fixed exit section then supplies a regular coordinate. These are two
overlapping parts of the orbit interval for the same coefficients.

\subsection{Actual section coordinates at both stopped endpoints}
\label{subsec:actual-dulac-matching}

For bounded amplified height, the middle passage is close to its linear
normalization up to an early negative cut. We stop there and use exact
saddle invariants to continue to both actual sections. This retains the
nonlinear section maps needed to compare the two small powers.

\begin{lemma}[Uniform actual-section matching]
\label{lem:actual-dulac-matching}
Consider \eqref{fx:pd}--\eqref{fx:coefficients}, with the compact
coefficient conditions stated there. Fix \(V_1<\infty\) and let
\(V\) be the raw middle input in \eqref{eq:comparison-action}.
Choose the local cuts sufficiently small in the order in the proof,
before decreasing \(\kappa,\tau\).
On \(0\le V\le V_1+1\) there is an increasing change
\(\widehat V=V+O_{C^3}(\epsilon_\theta)\), fixing zero, where
\(\epsilon_\theta/\chi\to0\). The two actual section coordinates,
including the finite closing map on the negative side, have the form
\begin{equation}\label{eq:actual-dulac-pair}
 x_{\rm in}=x_b+P_*\Psi(\widehat V^\chi),\qquad
 f_\theta(mt_{\rm out})=b_*+a_*D
       \bigl(\beta^{-\beta}C(\widehat V)^\beta\bigr),
 \qquad P_*,a_*>0,
\end{equation}
for \(m>0\). Put \(e_C(v)=vC'(v)/C(v)>0\) for \(v>0\).
Uniformly for \(0<V\le V_1+1\),
\begin{equation}\label{eq:actual-dulac-bounds}
 \begin{gathered}
 C(0)=0,\quad C'(0)=1,\quad
 \|C-\mathrm{id}\|_{C^3}\le\epsilon_\theta,\qquad
 \left|\frac{V e_C'(V)}{e_C(V)}\right|=o(\chi),\\
 D'>c_D>0,\qquad |D''/D'|\le C_D,\qquad
 \sup_{0\le t\le(V_1+2)^\chi}
       |t\Psi''(t)/\Psi'(t)|\longrightarrow0,\qquad P_*\to0 .
 \end{gathered}
\end{equation}
The maps \(\widehat V,C,D,\Psi\) extend in \(C^3\) to their
zero endpoints; this does not assert ordinary \(C^3\) regularity
of the compositions with the small powers of \(\widehat V\).
The normalized estimates have a uniform one-sided extension as
\(j=1/c\downarrow0\) at the remote-root face. Actual raw section
maps and their positive affine scales are asserted only for \(j>0\).
\end{lemma}

\begin{proof}
We first construct local coordinates in which the saddle invariant is
exact and the derivative bounds remain uniform as its exponent vanishes.
After dividing by the positive normal rate, and reversing time at
the positive saddle, both fields on a fixed small box take the form
\begin{equation}\label{eq:saddle-normal-time}
 \dot x=-\rho x p(x)+c(x)y,\qquad
 \dot y=y\{1+k(x)y\},\qquad \rho\in\{\beta,\chi\}.
\end{equation}
Here \(p(0)=1\), \(p>c_p>0\), and \(p,c,k\) have uniform
bounds through any prescribed finite order. At the negative saddle
\(x=u+1,y=w\), \(c=K/L<0\); at the positive saddle
\(x=(u-a)/a,y=w/a^2\), \(c>0\). Both \(|c|\) and its
reciprocal are bounded. These assertions follow by substituting the
displayed scalings in \eqref{fx:coefficients}; in particular the
positive normal rate is divided out only after scaling by \(a\).

Put \(h=-\rho xp(x)\). Construct the footpoint \(\pi\) by backward
axis transport as in Lemma~\ref{lem:normal-time-projection}, and put
\(J=\int_0^\infty k(X(-T))Y(-T)\,dT\), using bounded extensions.
Here \(Y(-T)\le ye^{-T/2}\), \(h=O_{C^4}(\rho)\), and the
couplings \(c_xY,k_xY^2\) have integrable mass. The same triangular
argument bounds tangential rows by \(C(1+T)^Ne^{C\rho T}\) and
normal rows by \(C(1+T)^Ne^{-(1/2-C\rho)T}\).
For a normalized coefficient parameter \(\vartheta\), its new forcing
is \(-\bigl(h_\vartheta+c_\vartheta Y,k_\vartheta Y^2\bigr)\):
the tangential term adds polynomial factors, while normal terms keep a
decaying height factor. Thus the backward integrals and their mixed rows
through five phase derivatives and one parameter derivative have an
integrable majorant \(C(1+T)^Ne^{-(1/2-C\rho)T}\), after fixing the
finite order and then decreasing \(\rho\). These mixed rows use axis-flow
derivatives through order seven, supplied by the prescribed finite
coefficient bounds.

Consequently \(\dot\pi=h(\pi)\), and \(Z=ye^{-J}\) satisfies
\(\dot Z=Z\). The axis coordinate
\[
 \ell(x)=x\exp\int_0^x\frac{p(s)^{-1}-1}{s}\,ds
\]
has uniform finite-order bounds and positive derivative. Hence
\(\Xi=\ell(\pi)\) gives the exact invariant
\begin{equation}\label{eq:exact-saddle-invariant}
 \dot\Xi=-\rho\Xi,\qquad I=\Xi Z^\rho,\qquad \dot I=0.
\end{equation}
This is a uniformly regular local coordinate change. Since
\(\pi_y(x,y)=-c(x)+o_\rho(1)+O(y)\) and \(|c|\) is bounded below,
choose the height strip first and then \(\rho\) to preserve its sign.
All these bounds concern intrinsic coordinates, before extracting raw
section scales.

Fix a permissible small upper height \(y_{\max}\) in the negative
box. Choose \(x_->0\) small compared with this height, and then
\(\eta_0>0\) sufficiently small compared with both, as quantified
in the output construction below. Finally fix
\(0<\eta_-<\eta_0\) and put \(u_0=-1+\eta_0\).
This cut lies in the negative box, and
strictly before the cut \(u_-\) in positive time. If
\(J_m(u)=\exp\int_{u_+}^u L/h\), put
\[
 \varepsilon_+=\frac{\lambda h_-}{a^2J_m(u_-)},\qquad
 \varepsilon_0=\frac{\lambda h_-J_m(u_0)}{J_m(u_-)}.
\]
The raw heights at \(u_+\) and \(u_0\) are respectively
\(a^2\varepsilon_+V\) and \(\varepsilon_0 M_0(V)\).
Partial fractions \eqref{fx:partial-fractions} imply, for some
fixed \(c_0>0,N<\infty\),
\begin{equation}\label{eq:early-saddle-smallness}
 \varepsilon_++\varepsilon_0
       \le C\beta^{-N}e^{-c_0/\beta},\qquad
 c\le\varepsilon_0^\beta\le C,\qquad
 \varepsilon_+^\chi\le Ce^{-c_0/(2q)} .
\end{equation}
To check uniformity, \(\beta\le Ca\), so every inverse power of
\(a\) is bounded by a power of \(\beta^{-1}\). The logarithmic
terms in the partial fractions are
\(q\log a\), \(\beta\log a\), and \(\beta\log\beta\),
all tending to zero. The action between \(u_0\) and \(u_-\)
has a fixed positive limit, while the full action to \(u_-\)
has the positive limit in \eqref{fx:gamma-endpoint}.
These facts give all three estimates, with the last obtained by
multiplying the logarithm of the first by \(\chi=\beta/q\).

Before \(u_0\), the middle map satisfies
\begin{equation}\label{eq:early-middle-identity}
 M_0(0)=0,\quad M_0'(0)=1,\qquad
 \|M_0-\mathrm{id}\|_{C^3([0,V_1+1])}
       \le C\beta^{-N}e^{-c_0/(2\beta)}.
\end{equation}
Indeed \(\rho'/\rho=(L-h')/h\) has at most one zero, a minimum;
the endpoint actions give
\(\sup_{[u_0,u_+]}\rho\le C\beta^{-N}e^{-c_0/\beta}\).
The ordinary equation is
\[
 V_u=\mathscr F=\frac{\lambda\rho(E-KL/h)V^2}{1+\lambda K\rho V}.
\]
On \([0,V_1+2]\) its denominator is at least one, and the integral
of its \(C^3_V\) norm is at most
\(C\beta^{-N}(1+|\log\beta|)e^{-c_0/\beta}\).
The first-exit and triangular-variation argument of
Lemma~\ref{lem:scalar-normalization} proves the estimate;
\(\mathscr F=O(V^2)\) gives \(M_0'(0)=1\). No ordinary
\(u\)-derivative at the shrinking positive cut is needed.

On the positive middle cut write the coordinate restrictions as
\(\Xi_{+,m}(y),Z_{+,m}(y)\), and let
\(z_+=Z_{+,m}'(0)>0\). Define
\[
 \widehat V=
 \frac{Z_{+,m}(\varepsilon_+ V)}{z_+\varepsilon_+}
 \left(\frac{-\Xi_{+,m}(\varepsilon_+ V)}
                  {-\Xi_{+,m}(0)}\right)^{1/\chi} .
\]
The weak coordinate is negative on this cut. Taylor's formula and
\eqref{eq:early-saddle-smallness} give
\(\widehat V=V+O_{C^3}(\varepsilon_+/\chi)\), with
\(\widehat V'(0)=1\). Its inverse has the same bounds.
Choose a fixed incoming local cut \(x=x_+>0\) in the positive
box. The stopped orbit \(\Xi=0\) intersects it at a height
\(y_+>c>0\), since \(\Xi_y<0\). On a fixed chord above that
height the function
\[
 L_+(y)=-\Xi_{+,\rm out}(y)Z_{+,\rm out}(y)^\chi
\]
is increasing, vanishes at \(y_+\), and has a uniformly bounded
\(C^3\) inverse. The exact invariant gives
\[
 L_+(y)=\Sigma_+\widehat V^\chi,\qquad
 \Sigma_+=-\Xi_{+,m}(0)(z_+\varepsilon_+)^\chi
                    \le Ce^{-c_0/(2q)}.
\]

Continue this chord to the actual positive section. In
\((X,Y)=(u/a,w/a^2)\) the limiting stopped graph is
\(Y=(X^2-1)/2\), \(X\ge1+x_+\). A smaller chord has positive
height and denominator margins to a fixed large cut and then enters a
common parabolic band. Apply the outgoing-tail normalization
\eqref{ann:tail-first-derivative}--\eqref{eq:N-bound} and the
fixed-port admission argument of Lemma~\ref{ann:p0-fixed-tail-ports}.
Both use only this half-family and its common band. If \(d_+\) is the
exact raw derivative at \(y_+\), then
\(p=(x_{\rm in}-x_b)/d_+\) gives
\(L_+(y)=\Phi_+(p)\), \(\Phi_+(0)=0\), with positive derivative
units and uniform \(C^3\) inverse bounds. Retain the absolute gain
\(d_+=\Omega a^{2-\lambda_\gamma}U_+\le C\), where \(U_+\) is
a positive unit.
Set
\[
 \sigma_+=\Phi_+^{-1}(\Sigma_+),\qquad
 \Psi(t)=\frac{\Phi_+^{-1}(\Sigma_+t)}{\sigma_+},\qquad
 P_*=d_+\sigma_+ .
\]
Then \(\sup|t\Psi''/\Psi'|\le C\Sigma_+\), giving the
positive assertions without any bound on unscaled raw derivatives.

At the negative cut \(u_0\), denote the coordinate restrictions by
\(\Xi_{-,m}(w),Z_{-,m}(w)\), and put
\(\xi_0=\Xi_{-,m}(0)>0\), \(z_0=Z_{-,m}'(0)>0\).
For \(v\) in a fixed bounded interval set
\[
 \mathcal G(v)=\frac{Z_{-,m}(\varepsilon_0v)}{z_0\varepsilon_0}
 \left(\frac{\Xi_{-,m}(\varepsilon_0v)}{\xi_0}\right)^{1/\beta},
 \qquad C=\mathcal G\circ M_0\circ(V\mapsto\widehat V)^{-1}.
\]
This is the actual strong-coordinate correction on the early cut.
Its \(C^3\) difference from the identity is
\(O(\varepsilon_0/\beta)\); the cut \(\eta_0\) is fixed.
Together with \eqref{eq:early-middle-identity} this yields
\(C'(0)=1\) and the claimed estimate for \(C-\mathrm{id}\),
with \(\epsilon_\theta=C\beta^{-N}e^{-c_0/(3\beta)}\).
Hadamard division gives \(C(v)=v c(v)\),
\(\|c-1\|_{C^2}\le C\epsilon_\theta\), so
\(|v e_C'/e_C|\le C\epsilon_\theta=o(\chi)\).

On the fixed outer negative cut \(u=-1-x_-\),
\(\Xi_w>0\), so the stopped branch \(\Xi=0\)
meets it at a height \(w_->c x_->0\). On a slightly larger
positive-height chord its invariant
\[
 I_-(w)=\Xi_{-,\rm out}(w)Z_{-,\rm out}(w)^\beta
\]
has \(I_-'>c>0\) and uniformly bounded \(C^3\) norm.
Writing \(r=\beta^{-\beta}C(\widehat V)^\beta\), equality
of the two invariants is precisely
\[
 I_-(w)=\alpha_\theta r,\qquad
 \alpha_\theta=\xi_0(z_0\varepsilon_0\beta)^\beta .
\]
The constant \(\alpha_\theta\) is bounded above and away from
zero by \eqref{eq:early-saddle-smallness}. Also
\(\alpha_\theta\le C\eta_0\), while \(r\le1+o(1)\)
on the fixed amplified-height interval. The choices above can be
made with \(\pi(-x_-,y_{\max})\ge c y_{\max}>0\), and
then with
\[
 (1+\epsilon)\sup_\theta\alpha_\theta
       <\tfrac12\inf_\theta I_-(y_{\max}).
\]
The whole invariant interval consequently lies strictly inside the
output chord, with a uniform buffer. The inverse function theorem gives
\(w=D_0(r)\) on \([0,1+\epsilon]\), with uniform \(C^3\)
and inverse bounds. At \(r=0\) its output is \(w_->0\).
The extension at the excluded endpoint \(\widehat V=0\) is thus
obtained on a positive-height chord.

From this cut the limiting stopped graph is \(w=(u^2-1)/2\).
Its buffered chord has the same finite-cut margins and negative-tail
band. The preceding normalization and admission argument, now on the
negative half-tail, gives
\(T(w)=t_b+\Omega U_-B_-(w)\), with \(B_-'>c\) and uniform
\(C^3\) map and inverse bounds. In particular \(|T''/T'|\le C\),
although its raw slope may vanish.
For the divided finite closing map, the chain rule gives
\[
 \frac{(f_\theta(mT)/m)''}{(f_\theta(mT)/m)'}
   =\frac{T''}{T'}+
      m\frac{f_\theta''(mT)}{f_\theta'(mT)}T' .
\]
The last term is bounded on the fixed output window since
\(0<\Omega\le C\). Normalizing its exact value and positive
slope at a reference point, and composing with \(D_0\), gives
\(D\) and the constants \(b_*,a_*\) in
\eqref{eq:actual-dulac-pair}. This proves the negative assertions.

Finally, in the remote-root chart the coefficients on both local
boxes are regular functions of \(j\). The source-tail scales are
exactly \(\Omega\) and \(\Omega a^{2-\lambda_\gamma}\);
\(|\lambda_\gamma|<1\) keeps both bounded above. No reciprocal
bound on these raw scales was used. The fixed-cut action has no
\(\beta\log j\) term, so every estimate above is uniform as
\(j\downarrow0\), as a one-sided family of physical passages.
\end{proof}


\subsection{The fixed-exit continuation}
\label{sec:fixed-exit-comparison}

For nonzero limiting action $\Theta=\beta\log V$, orbits of
\eqref{fx:pd}--\eqref{fx:graph} leave the small-height region before
the negative saddle. At the fixed exit $u_c\in(-1,u_-)$ the height
$W=w(u_c)$ provides a regular coordinate. Shooting backwards from this
section constructs the scalar inverse phase and its $C^2$ limit without
using the positive return-section map.

\subsection{The action interval and the exit height}

For \(J(u)=\exp\int_{u_+}^u L/h\), put
\begin{equation}\label{fx:rho-gamma}
 \rho=\frac{h(u_-)J(u)}{J(u_-)h(u)},\qquad
 \Gamma_\beta=-\beta\log\rho,\qquad H_\beta=h/\beta .
\end{equation}
The exact partial fractions are
\begin{equation}\label{fx:partial-fractions}
 \frac Lh=-\frac1{\beta(u+1)}
       +\frac1{\chi(a-u)}
       -\frac{1+1/\beta+1/\chi}{c-u}.
\end{equation}
Thus
\begin{equation}\label{fx:gamma-derivative}
 \Gamma_\beta'=-L/H_\beta+\beta H_\beta'/H_\beta,\qquad
 H_\beta=\frac{b_-(u+1)(a-u)(c-u)}{(1+a)(1+c)} .
\end{equation}
On negative compact intervals this tends in \(C^3\) to
\begin{equation}\label{fx:limit-functions}
 H_{0,c}=\frac{(-u)(1+u)(c-u)}{1+c},\qquad
 \Gamma_c(u)=\log\frac{(1+u)(c-u_-)}{(1+u_-)(c-u)} .
\end{equation}
Indeed, \(\Gamma_c'=-u/H_{0,c}\), and both actions vanish at \(u_-\).
The limit action labels the point at which an orbit leaves the
small-height region. From that point to the fixed exit, the limiting
graph equation is \(w_u=u\). These two limits will parametrize the
inverse transition.

At the shrinking positive cut, integration of
\eqref{fx:partial-fractions} gives
\[
 \beta\log J(u_-)=
 \log\frac{(1+u_+)(c-u_-)}{(1+u_-)(c-u_+)}
 +q\log\frac{(a-u_+)(c-u_-)}{(a-u_-)(c-u_+)}
 +\beta\log\frac{c-u_-}{c-u_+}.
\]
By \eqref{fx:log-smallness},
\begin{equation}\label{fx:gamma-endpoint}
 \Gamma_\beta(u_+)\longrightarrow
 \Gamma_{-,c}:=\Gamma_c(0)=
             \log\frac{c-u_-}{c(1+u_-)}>0 .
\end{equation}
Also \((L-h')'=1+o(1)>0\); hence
\(\Gamma_\beta'=-\beta(L-h')/h\) has at most one zero, a maximum.

Choose \(0<\Delta<\frac14\inf_c\Gamma_{-,c}\), and set
\[
 r_1=\Gamma_c^{-1}(\Gamma_{-,c}-\Delta),\qquad
 r_1^\sharp=\Gamma_c^{-1}(\Gamma_{-,c}-3\Delta/4).
\]
Choose fixed cuts \(u_c,u_R\) and \(\epsilon>0\) with
\[
 -1<u_c<u_--2\epsilon,\qquad
 r_1^\sharp(c)<u_R<0,\qquad
 \Gamma_c(u_R)>\Gamma_{-,c}-\Delta/2 .
\]
These choices are uniform in the compact \(c\)-family. Define
\begin{equation}\label{fx:W-limit}
 \mathcal W(r)=\frac{u_c^2-r^2}{2},\quad
 \mathcal J_c=\mathcal W([u_-,r_1]),\quad
 \mathcal J_c^\sharp=\mathcal W([u_--\epsilon,r_1^\sharp]).
\end{equation}
They are positive compact intervals with a uniform strict inclusion.

\subsection{Convergence of the inverse transition}

We first work only with the scalar equation \eqref{fx:graph}, without
assuming any positive return-section map. Let \(\mathcal K_c\) be a
compact interval strictly inside \((0,u_c^2/2)\), with uniform endpoint
margins, and prescribe \(w(u_c;W)=W\) for \(W\in\mathcal K_c\).
Write
\[
 r(W)=-\sqrt{u_c^2-2W},\qquad
 d_{\mathcal K}=\inf_{c,W\in\mathcal K_c}
       \{\Gamma_{-,c}-\Gamma_c(r(W))\}>0.
\]
Increase \(u_R<0\) towards zero if necessary, so that it lies beyond
all \(r(W)\) and
\(\Gamma_c(u_R)>\Gamma_{-,c}-d_{\mathcal K}/4\).
We follow the backward shot through its positive-height part, the thin
transition layer, and the remaining small-height part. The next three
first-exit arguments continue it to \(u_+\) and estimate its phase.
They apply, in particular, to
\(\mathcal J_c^\sharp\) and to any fixed larger chord still inside
\((0,u_c^2/2)\). Define
\begin{equation}\label{fx:extracted-phase}
 \phi_\beta(u;W)=\beta\log\frac{w(u;W)}{h(u)}+\Gamma_\beta(u).
\end{equation}
Its value at \(u_+\) gives the exact inverse phase
\begin{equation}\label{fx:inverse-phase}
 \Psi_\theta(W)=\phi_\beta(u_+;W)-\beta\log\lambda
                 =\beta\log V=:\Theta .
\end{equation}
The differential equation for this phase is
\begin{equation}\label{fx:phase-equation}
 \phi_\beta'=\frac{r(Eh-KL)}{H_\beta(1+Kr)},\qquad r=w/h .
\end{equation}
The phase changes by only a small amount after the orbit has entered
the small-height part. This will allow the leaving point of the limiting
graph to determine the inverse phase.

\begin{lemma}\label{fx:lem-fast-inverse}
For fixed \(\omega_0>0\), the graph meets \(w=\beta\omega_0\)
at a unique point \(\sigma_\beta(W)<u_R\), and, uniformly on
\(\mathcal K_c\),
\[
 \|\sigma_\beta-\mathcal W^{-1}\|_{C^2}
     \le C(\omega_{\rm d}+\beta|\log\beta|)\longrightarrow0 .
\]
Here \(\omega_{\rm d}\) bounds in \(C^3\) the differences
\(H_\beta-H_{0,c},K-1,L-u,E\) on the fixed negative rectangle.
\end{lemma}
\begin{proof}
On that rectangle \(L+Ew\le-\ell<0\).
The inverse graph, \(U_\beta(W;W)=u_c\), satisfies
\[
 \partial_w U_\beta
       =\frac K{L+Ew}+\frac{\beta H_\beta}{w(L+Ew)} .
\]
Between \(\beta\omega_0\) and the fixed upper height, its
right side differs from \(1/u\) in \(L^1_w\) by at most
\(C\omega_{\rm d}+C\int_{\beta\omega_0}^{W_{\max}}\beta\,dw/w\).
Gronwall's inequality and the first two triangular variational
equations give the asserted \(C^2\) bound.
To justify existence, apply this estimate only up to a first exit
from a fixed rectangle about the limiting graphs. Their endpoints
\(r(W)<u_R\) stay strictly inside,
so an exit is impossible. Negative nonzero speed at the selected
height gives uniqueness and transversality.
\end{proof}

\begin{lemma}\label{fx:lem-corner}
Uniformly on \(\mathcal K_c\),
\[
 \|\phi_\beta(u_R;W)-\Gamma_\beta(\sigma_\beta(W))\|_{C^2}
       \le C\beta .
\]
\end{lemma}
\begin{proof}
On this fixed negative rectangle,
\(g_\beta=-(L-h')/H_\beta\ge g_0>0\), and the exact normal equation is
\[
 \beta r'=\frac{r\{-g_\beta+\beta r(E-KH_\beta'/H_\beta)\}}{1+Kr}.
\]
Set \(\zeta=(u-\sigma)/\beta\), so \(r_\zeta=f_\beta\),
\(r(0;\sigma)=\omega_0/H_\beta(\sigma)\). On the fixed bounded
\(r\)-interval, \(f_r\le-a_0<0\),
\(|f_u|+|f_{uu}|\le Cr\), and \(|f_{ur}|+|f_{rr}|\le C\).
The triangular variation argument of
Lemma~\ref{lem:scalar-normalization}, with these decaying parameter
forcings, gives
\[
 |\partial_\sigma^j r|\le C_j(1+\zeta^j)e^{-c\zeta},\qquad j=0,1,2.
\]
A first-exit argument retains the rectangle up to \(u_R\).
For \(G(\sigma,\zeta)=F_\beta(\sigma+\beta\zeta,r)\), where
\(F_\beta\) is the right side of \eqref{fx:phase-equation}, put
\(T=(u_R-\sigma)/\beta\), \(J=\beta\int_0^T G\,d\zeta\).
The moving endpoint must be retained:
\[
 J_\sigma=\beta\int_0^T G_\sigma\,d\zeta-G(\sigma,T),\qquad
 J_{\sigma\sigma}=\beta\int_0^T G_{\sigma\sigma}\,d\zeta
 -2G_\sigma(\sigma,T)+\beta^{-1}G_\zeta(\sigma,T).
\]
The integral terms are \(O(\beta)\); all boundary terms are
exponentially small because \(u_R-\sigma\ge d_0>0\).
Thus \(J=O_{C^2_\sigma}(\beta)\). At the crossing,
\(\phi_\beta-\Gamma_\beta
=\beta\log(\omega_0/H_\beta(\sigma_\beta))=O_{C^2_W}(\beta)\).
Compose with the preceding \(C^2\) crossing map.
\end{proof}

\begin{lemma}\label{fx:lem-flat-part}
For some fixed \(d>0\) and finite integer \(n\),
\[
 \|\phi_\beta(u_+;W)-\phi_\beta(u_R;W)\|_{C^2}
          \le C\beta^{-n}e^{-d/\beta}.
\]
\end{lemma}
\begin{proof}
The strict cuts and preceding estimates give the fixed phase gap
\[
 \phi_\beta(u_R;W)\le\Gamma_{-,c}-3d_{\mathcal K}/4,\qquad
 \inf_{[u_R,u_+]}\Gamma_\beta\ge\Gamma_{-,c}-d_{\mathcal K}/2.
\]
The second bound follows from endpoint convergence and the fact that
\(\Gamma_\beta\) has at most one critical point, a maximum.
Equation~\eqref{fx:partial-fractions} gives bounded \(L/H_\beta\):
at the shrinking cut its only potentially large term satisfies
\(q/(a-u)\le q/(a\eta_+)\le C\). Thus \(F_\beta\) and its first
two \(r\)-derivatives are bounded; ordinary \(u\)-derivatives at
that cut are not needed.

Since \(h,K>0\), the shot continues while it stays in the rectangle.
The identity \(r=e^{(\phi_\beta-\Gamma_\beta)/\beta}\) and the phase
gap give \(r\le Ce^{-d/\beta}\): up to a first loss, the phase change
is at most \(C\int r\,du\), too small to close half the gap.
For its first two phase variations,
\(r_\phi=r/\beta\), \(r_{\phi\phi}=r/\beta^2\).
The scalar triangular variation argument therefore has exponentially
small integrated coefficients times finite powers of \(\beta^{-1}\).
It preserves the bounded \(W\)-variations at \(u_R\) and gives the
claimed \(C^2\) error.
\end{proof}

\begin{proposition}\label{fx:prop-inverse-phase}
The scalar backward-shooting solutions are defined up to \(u_+\),
and, jointly in the compact \(c\)-family, on every chord
\(\mathcal K_c\) above,
\begin{equation}\label{fx:phase-limit}
 \Psi_\theta\longrightarrow
       \Psi_{0,c}=\Gamma_c\circ\mathcal W^{-1}\quad\hbox{in }C^2 .
\end{equation}
More precisely, for the coefficient error \(\omega_{\rm d}\) in
Lemma~\ref{fx:lem-fast-inverse},
\begin{equation}\label{fx:phase-error-budget}
 \|\Psi_\theta-\Psi_{0,c}\|_{C^2(\mathcal K_c)}
 \le C\{\omega_{\rm d}+\beta(1+|\log\beta|)\}.
\end{equation}
Using \(\mathcal K_c=\mathcal J_c^\sharp\), its inverse
\(\mathcal H_\theta\) is defined on
\begin{equation}\label{fx:action-band}
 0\le\Theta\le\Gamma_{-,c}-\Delta
\end{equation}
and satisfies \(0<c_H\le\mathcal H_\theta'\le C_H\) and
\(|\mathcal H_\theta''/\mathcal H_\theta'|\le C_H\).
\end{proposition}
\begin{proof}
The exact derivative \eqref{fx:gamma-derivative}, together with
\(\Gamma_\beta(u_-)=\Gamma_c(u_-)=0\), gives
\(\Gamma_\beta-\Gamma_c=O_{C^3}(\omega_{\rm d}+\beta)\)
on the fixed negative interval.
Combining the three estimates in \eqref{fx:inverse-phase} and composing
with the uniformly \(C^2\) crossing \(\sigma_\beta(W)\) gives
\eqref{fx:phase-error-budget}; the exponentially small remainder
is absorbed into \(O(\beta)\). All derivatives here are in \(W\)
with the coefficients fixed.
For \(r=\mathcal W^{-1}(W)\),
\(\Psi_{0,c}'(W)=1/H_{0,c}(r)>0\) with uniform upper and lower
bounds. The larger intervals in \eqref{fx:W-limit} give a
fixed image margin around \eqref{fx:action-band}.
The uniform inverse function theorem applies. Explicitly,
\[
 u_{*,c}(\Theta)=
 \frac{c(1+u_-)e^\Theta-(c-u_-)}
      {(1+u_-)e^\Theta+(c-u_-)},\qquad
 \mathcal H_{0,c}(\Theta)=\frac{u_c^2-u_{*,c}(\Theta)^2}{2}.
\]
In particular, \(\Theta=0\) has a positive interior height on
the fixed exit section.
\end{proof}

The strict gap between the fixed-exit action endpoint and a negative
axis cut is
\begin{equation}\label{fx:overlap-gap}
 \Gamma_{-,c}-\Gamma_c(-\sigma)
    =\log\frac{c+\sigma}{c(1-\sigma)}>0,\qquad 0<\sigma<1 .
\end{equation}
Proposition~\ref{p2:prop-finite-v-fe-actual-budget} uses this gap to
place the fixed-exit construction and the bounded-height construction
on overlapping actual receivers. It then proves the completed-return
curvature estimate on their union, which is the physical interval used
by the coefficient cover.


\section{The compact positive-ratio passage}
\label{sec:compact-positive}

We next treat the coefficient face \(\alpha _0=-s^2<0\) with
\(p=\alpha _1/s\) in a fixed interval \([p_0,P]\subset(0,\infty)\).
The low-height orbits pass between a positive equatorial saddle and
a negative saddle escaping to \(X=-\infty\), whose characteristic
number tends to zero.

The escaping saddle prevents a comparison on fixed compact source
sections alone. We parametrize its passage by the point where the strong
leaf meets the equator. An axis integral measures the normal expansion
all the way to that saddle. Matching it with the positive Dulac map gives
a curvature estimate in the input of the same return, uniformly as the
negative characteristic number tends to zero. We then join this endpoint
estimate to the compact middle and the high interval, where the derivative
is bounded.

\begin{lemma}[Domains of the compact positive-ratio passage]
\label{ann:lem-compact-positive-domains}
For \(p_0\le p\le P\), sufficiently small \(s>0,|\gamma|\), and one
fixed sufficiently small bound on \(z(\gamma)\), the actual through
interval is covered by the negative-saddle endpoint, a compact family
of strict source passages, and the high-height interval. Consecutive
pieces overlap in actual trajectories. The completed transition has
negative curvature on the endpoint and compact middle, and bounded
first derivative on the high-height interval. All cuts and estimates
are uniform on the fixed \(p\)-interval.
\end{lemma}

\subsection{The equatorial roots and the strong foliation}
\label{ann:subsec-negative-saddle}

With \(a=-1\), write \eqref{ann:source-field} as
\[
 X'=-(M+LY),\qquad Y'=-Y(XB+\zeta Y),
\]
where
\begin{equation}\label{ann:source-coefficients}
\begin{aligned}
 M&=1-pX+3\gamma X^2+sX^3,& B&=1+2\gamma+sX,\\
 L&=1-A(\gamma+sX)+C(\gamma+sX)^2,&
 \zeta&=Cs(\gamma+sX).
\end{aligned}
\end{equation}

\begin{lemma}\label{ann:lem-axis-roots}
The roots of \(M\) are simple and have the order
\(-R<x_+<x_f\), where \(R\to\infty\), \(x_+\to p^{-1}\).
For \(M_-(q)=M(-q)\), \(B_-(q)=B(-q)\), put
\(\kappa=-M_-'(R)\) and \(\rho_-=\kappa/(RB_-(R))\).
Then
\begin{equation}\label{ann:negative-root-identity}
 sR=3\gamma+p/R+R^{-2},
\end{equation}
and
\begin{equation}\label{ann:small-negative-index}
 \kappa\ge cp_0,\qquad
 0<\rho_-\le C(R^{-1}+|\gamma|+sR)\longrightarrow0 .
\end{equation}
At \(x_+\), the rates \(L_+=-M'(x_+)\), \(\nu_+=x_+B(x_+)\)
give \(0<\rho_0\le\rho_+=\nu_+/L_+\le\rho_1<\infty\).
\end{lemma}
\begin{proof}
On compact sets \(M\to1-pX\), yielding \(x_+\) and the negative
sign immediately to its right. Since \(M(0)>0\), the cubic has a
second positive root and a negative root, both escaping uniformly.
Substitution at \(-R\) proves \eqref{ann:negative-root-identity}.
The factorization
\[
 M=s(X+R)(X-x_+)(X-x_f),\qquad sRx_+x_f=1
\]
gives
\(\kappa=x_+^{-1}(1+x_+/R)(1+R/x_f)\ge cp_0\).
Direct differentiation also gives
\(\rho_-=(2p/R+3\gamma+3/R^2)/B_-(R)\).
The positive rates converge on the compact \(p\)-interval.
\end{proof}

Fix \(Q\gg P\), then \(R_0>2Q\), and shrink the coefficients so that
\begin{equation}\label{ann:delta-Q}
 d_Q=Q^{-2}+P/Q+|\gamma|+sR\ll1 .
\end{equation}
In \(q=-X\) the negative tail is
\begin{equation}\label{ann:negative-tail-field}
 q'=M_-(q)+L_-(q)Y,\qquad
 Y'=Y\{qB_-(q)+\zeta_-(q)Y\},
\end{equation}
with
\[
 \begin{aligned}
 M_-&=1+pq+3\gamma q^2-sq^3,& B_-&=1+2\gamma-sq,\\
 L_-&=1-A(\gamma-sq)+C(\gamma-sq)^2,&
 \zeta_-&=Cs(sq-\gamma).
 \end{aligned}
\]
Let \(\mathcal V\) denote \eqref{ann:negative-tail-field}.  Use the
positive normal time \(d\tau_n=qB_-(q)\,dt\), and write
\(V=\mathcal V/(qB_-(q))\).  In \(u=\log q\), \(w=Y/q^2\),
put \(c=L_-/B_-\) and \(k=\zeta_-q/B_-\).  The axis speed is
\begin{equation}\label{ann:negative-slow-speed}
 h_-(q)=\frac{M_-(q)}{q^2B_-(q)}
       =\frac{q^{-2}+p/q+3\gamma-sq}{B_-(q)}=O(d_Q).
\end{equation}
The exact normal-time equations are
\[
 u'=h_-+cw,\qquad w'=w\{1-2h_-+(k-2c)w\}.
\]
The functions \(h_-\), \(c-1\), \(k\) and their first five Euler
derivatives \(D_q^j=(q\partial_q)^j\) are \(O(d_Q)\) on
\(Q/2\le q\le2R\) and fixed logarithmic collars.
Extend this triple consistently as in
Lemma~\ref{lem:normal-time-projection}.  Its hypotheses hold with
\(\epsilon=Cd_Q\); choosing \(Q\) first and then the small
coefficient bounds makes this quantity sufficiently small.  The
projection is therefore constructed directly from the negative-tail
field, independently of either return-map estimate.

It gives a strong projection \(\pi(q,Y)\) for \(V\) with uniform
scaled bounds
on the whole corridor, including the root \(R\).  The induction
through one further order, using the fifth coefficient derivatives,
gives a bounded \(C^4\) norm for \(\log\pi-\log q\) in
\((\log q,Y/q^2)\).  Throughout this
section, the footpoint refers to this normal-time projection.
In particular, its invariance equation is
\begin{equation}\label{ann:strong-projection}
 \pi(q,0)=q,\qquad \pi_*V=b_-(r)\partial_r,\qquad
 b_-(r)=\frac{M_-(r)}{rB_-(r)},\qquad r=\pi(q,Y).
\end{equation}

Use the section
\begin{equation}\label{ann:negative-strong-section}
 Y/q^2=w_0,\qquad 0<w_0<1/4 .
\end{equation}
At the limiting corner, the leaf over \(r\) satisfies
\(Y-q^2/2=-r^2/2\), and meets the section transversely at
\(q=r(1-2w_0)^{-1/2}\).  The uniform projection estimates give a
unique intersection over every \(R_0\le r<R\), of the form
\[
 q_s(r)=rU_s(r),\qquad U_s\asymp1,\qquad q_s'(r)\asymp1,
 \qquad |(r\partial_r)^j\log U_s|\le C\quad(1\le j\le3).
\]
We use this footpoint \(r\) to parametrize the strong section.

Figure~\ref{fig:tails-roots} shows the compact core, long tails
and root configuration.
% Schematic domains and equatorial roots for a=-1, p in [p_0,P].
% The arrows in panel (a) denote positive-time transition maps.
% Panel (b) is an ordered root diagram, not an orbit on the equator.
\begin{figure}[htb]
\centering
\setlength{\unitlength}{1pt}
\begin{picture}(410,182)
 \put(205,172){\makebox(0,0){\small (a) The transition on fixed physical sections}}
 \put(45,148){\makebox(0,0){\(\Sigma_-\)}}
 \put(155,148){\makebox(0,0){\(X=-L\)}}
 \put(255,148){\makebox(0,0){\(X=L\)}}
 \put(365,148){\makebox(0,0){\(\Sigma_+\)}}
 \linethickness{.6pt}
 \put(142,128){\vector(-1,0){85}}
 \put(242,128){\vector(-1,0){74}}
 \put(353,128){\vector(-1,0){86}}
 \put(100,113){\makebox(0,0){\small outgoing tail}}
 \put(205,113){\makebox(0,0){\small compact passage}}
 \put(310,113){\makebox(0,0){\small incoming tail}}
 \put(45,97){\makebox(0,0){\(\dfrac{-\delta-\gamma}{s}\)}}
 \put(365,97){\makebox(0,0){\(\dfrac{\delta-\gamma}{s}\)}}
 \put(205,69){\makebox(0,0){\small (b) Equatorial roots in the same coordinate}}
 \linethickness{.35pt}
 \put(24,42){\vector(1,0){368}}
 \put(400,42){\makebox(0,0){\(X\)}}
 \put(195,39){\line(0,1){6}}
 \put(195,29){\makebox(0,0){\small \(0\)}}
 \put(74,42){\circle*{3}}
 \put(219,42){\circle*{3}}
 \put(338,42){\circle*{3}}
 \put(74,54){\makebox(0,0){\(-R\)}}
 \put(219,54){\makebox(0,0){\(x_+\)}}
 \put(338,54){\makebox(0,0){\(x_f\)}}
 \put(74,18){\makebox(0,0){\small \(-R\to-\infty\)}}
 \put(219,18){\makebox(0,0){\small \(x_+\to p^{-1}\)}}
 \put(338,18){\makebox(0,0){\small \(x_f\to+\infty\)}}
 \put(74,0){\makebox(0,0){\small \(\rho_-\to0\)}}
 \put(219,0){\makebox(0,0){\small \(\rho_0\le\rho_+\le\rho_1\)}}
\end{picture}
\caption{Two consequences of the rescaling \(v=\gamma+sX\),
\(z=s^2Y\), for \(\alpha_0=-s^2\) and \(p\in[p_0,P]\).
The positions and distances are schematic. In (a), a fixed compact
core does not include the tails leading to \(\Sigma_\pm\);
the arrows denote the three factors of the positive-time transition.
In (b), the ordered zeros of \(M\) are shown as \(s,\gamma\to0\).
The negative incident saddle escapes and its characteristic number
vanishes, although the positive incident saddle remains on a compact
scale. The axis arrow gives the orientation of \(X\), not the flow.}
\label{fig:tails-roots}
\end{figure}


\subsection{An integrating coordinate up to the negative saddle}

The linear normal equation along the negative axis is
\(d\log Y/dq=qB_-(q)/M_-(q)\). Its integral suggests the coordinates
\begin{equation}\label{ann:negative-action}
 \mathcal A_-(r)=\int_Q^r\frac{qB_-(q)}{M_-(q)}\,dq,\qquad
 \mathcal D_-=\frac{M_-(r)}{rB_-(r)}\partial_r .
\end{equation}
Then \(\mathcal D_-\mathcal A_-=1\), and
\(\mathcal A_-\to\infty\) at \(R\).
Thus \(\mathcal A_-\) covers the entire interval up to the saddle,
although the upper bound \(R\) for the footpoint tends to infinity.
The next lemma relates this linear integral to the height
of the actual orbit on the fixed cut \(q=Q\).

\begin{lemma}\label{ann:lem-global-action-preparation}
The height \(Y_Q(r)\) at the fixed cut \(q=Q\) of the orbit through
the strong section over \(r\) satisfies
\begin{equation}\label{ann:negative-action-preparation}
 \log Y_Q=c_\theta-\mathcal A_-(r)+\mathcal R_-(r),
 \qquad R_0\le r<R,
\end{equation}
where, for \(j=1,2\),
\begin{equation}\label{ann:negative-action-remainder}
 |\mathcal D_-^j\mathcal R_-|\le C(d_Q+e_{R_0}),\qquad
 e_{R_0}=C_QR_0^2e^{-c(R_0-Q)/P}\longrightarrow0.
\end{equation}
The value of \(\mathcal R_-\) need not be small.
\end{lemma}
\begin{proof}
Apply Lemma~\ref{lem:normal-time-action-height} to the triple just
constructed. The polynomial in \eqref{ann:negative-tail-field}
gives \(M_-(q)/q\ge p_0/2\) on \([Q/2,Q]\) after shrinking the
coefficients, and \(b_->0\) up to its simple root \(R\).
After \(R_0\) is fixed, the same polynomial gives
\(b_-\le CP\) on \([Q,R_0]\). Thus the lemma constructs the
whole actual height interval within \(Q/2<q<2R\), and yields
\begin{equation}\label{ann:global-action-exact}
 \log Y_Q=-\mathcal A_-(r)+2\log r+\log\mathcal U(r)+E_Q(Y_Q),
\end{equation}
with the unit and fixed-cut bounds stated there. In particular,
\begin{equation}\label{ann:negative-small-height}
 0<Y_Q(r)\le e_{R_0}=C_QR_0^2e^{-c(R_0-Q)/P}.
\end{equation}
Choose \(R_0\) large and then shrink the coefficient neighborhood.
Taking \(c_\theta=0\) and
\(\mathcal R_-=2\log r+\log\mathcal U+E_Q(Y_Q)\),
\eqref{eq:normal-time-action-jets} gives
\(\mathcal D_-\mathcal R_-=1-K\) and
\(\mathcal D_-^2\mathcal R_-=-\mathcal D_-K\), proving the
claimed bounds. The lemma also gives \(Y_Q\downarrow0\) at \(R\).
\end{proof}

\subsection{Derivatives on the return sections}

We now continue from the strong section to the actual return section,
retaining the footpoint power and the divided closing. In the following
proof, \(w_0\) keeps the small-lip restriction
\(0<w_0<1/4\) of \eqref{ann:negative-strong-section}; only
\(w_1\) is chosen in \((1/2,1)\).

\begin{lemma}\label{ann:lem-negative-gate}
Let \(\Gamma_-(r)\) denote the resulting coordinate on the negative
return section.
Then
\begin{equation}\label{ann:negative-gate-derivative}
 -\Gamma_-'(r)=\Omega_-(\theta)r^{1-\lambda_\gamma}U_-(r)>0,
 \qquad \Omega_->0,
\end{equation}
where \(U_-\) and its reciprocal are bounded and
\begin{equation}\label{ann:negative-gate-jets}
 |\mathcal D_-^j\log U_-|\le C(d_Q+e_{R_0}),\qquad j=1,2 .
\end{equation}
\end{lemma}
\begin{proof}
Fix \(w_0\), \(w_1\in(1/2,1)\), \(\mu=3/2\), and the small
physical width \(\delta\) before its nested output windows. All phase
derivatives fix coefficients and multiplier. On the strong section the
inverse of \(\upsilon=\widehat\pi(u,w_0)\), \(\upsilon=\log r\), is
\(u=a(\upsilon)=\log q_s(r)\). The projection bounds give
\begin{equation}\label{ann:repair-inverse-start}
 a_1\asymp1,\qquad |a_2|+|a_3|\le C,\qquad 1<q_s/r<2,
 \qquad a_j=\partial_\upsilon^j a.
\end{equation}
Also, for \(h=M_-/(r^2B_-)\),
\begin{equation}\label{ann:repair-action-budget}
 \mathcal D_-=h\partial_\upsilon,\qquad
 |h|+|h_\upsilon|\le Cd_Q.
\end{equation}
The scalar graph and its logarithmic-radius equation are
\begin{equation}\label{ann:negative-scalar-graph}
 Y_q=f=\frac{Y(qB_-+\zeta_-Y)}{M_-+L_-Y},\qquad
 W_u=F(u,W)=W\left\{\frac{B_-+ZW}{D+L_-W}-2\right\},
\end{equation}
where \(W=Y/q^2\), \(D=q^{-2}+p/q+3\gamma-sq\), and \(Z=q\zeta_-\).
On \(w_0\le W\le w_1\), after choosing \(Q\) and shrinking the
coefficients,
\begin{equation}\label{ann:repair-outgoing-strip}
 D+L_-W\ge c,\qquad F(u,w_0)>c,\qquad F(u,w_1)<-c,\qquad F_W\le-\mu.
\end{equation}
The needed coefficient and state derivatives are bounded; explicitly,
\begin{equation}\label{ann:repair-rational-identities}
 qf_Y=F_W+2,\qquad
 F_{WW}=\frac{2D(ZD-B_-L_-)}{(D+L_-W)^3},\qquad
 F_{WWW}=-\frac{6DL_-(ZD-B_-L_-)}{(D+L_-W)^4}.
\end{equation}
Thus the strip is invariant and \(q'>0\). Since
\(\sup_{r<R}q_s(r)<q_H=(\delta+\gamma)/s\), every orbit reaches
this raw exit once, even when \(q_s>R\).

To extract the long-tail power, compare with
\[
 F^0_\gamma(W)=W\left\{\frac{1+2\gamma}{3\gamma+L_\gamma W}-2\right\},
 \quad L_\gamma=1-A\gamma+C\gamma^2,
 \quad w_\gamma=\frac{1-4\gamma}{2L_\gamma}.
\]
Here \(F^0_\gamma(w_\gamma)=0\),
\((F^0_\gamma)'(w_\gamma)=-2+\lambda_\gamma\),
\((F^0_\gamma)''=O(|\gamma|)\), and
\(F-F^0_\gamma=O(sq+p/q+q^{-2})\) in \(C^1_W\).
Scalar contraction, with \(q_0=q_s(r)\), consequently gives
\begin{equation}\label{ann:repair-memory-remainder}
 \left|f_Y-\frac{\lambda_\gamma}{q}\right|
 \le C\left(s+pq^{-2}+q^{-3}
       +\frac{|\gamma|}{q}(q_0/q)^\mu\right).
\end{equation}
The last term retains entrance memory and has integral at most
\(|\gamma|/\mu\), independently of the tail length. Hence
\begin{equation}\label{ann:repair-integrated-remainder}
 I(r)=\int_{q_0}^{q_H}(f_Y-\lambda_\gamma/q)\,dq,\qquad
 |I(r)|\le C(\delta+|\gamma|+p/q_0+q_0^{-2}).
\end{equation}
For the fixed-entry first variation,
\begin{equation}\label{ann:negative-first-variation}
 (J_Y)_q=f_YJ_Y,
\end{equation}
so the exact cocycle is
\begin{equation}\label{ann:negative-free-factor}
 \frac{J_Y(q)}{J_Y(q_0)}=(q/q_0)^{\lambda_\gamma}
 \exp\int_{q_0}^{q}\{f_Y(\widetilde q,Y(\widetilde q))-\lambda_\gamma/\widetilde q\}\,d\widetilde q.
\end{equation}

The moving initial value \(W(a(\upsilon);\upsilon)=w_0\) gives
\begin{equation}\label{ann:repair-moving-variation}
 J:=\partial_\upsilon W=-F(a,w_0)a_1
 \exp\int_a^uF_W\,dv<0,\qquad |J|\le Ce^{-\mu(u-a)}.
\end{equation}
Differentiate its logarithm using the triangular variation argument of
Lemma~\ref{lem:scalar-normalization}. The first integral forcing is
\(F_{WW}J\); in the second row it is
\(F_{WWW}J^2+F_{WW}J\partial_\upsilon\log(-J)\).
The moving lower-limit terms involve only \(a_1,a_2,a_3\) and the
bounded entrance derivatives of \(F\). Integrating the decaying
forcings in order yields
\begin{equation}\label{ann:repair-moving-log-first}
 |\partial_\upsilon\log(-J)|\le C,
\end{equation}
\begin{equation}\label{ann:repair-moving-log-second}
 |\partial_\upsilon^2\log(-J)|\le C.
\end{equation}
No reciprocal of \(J\) is estimated separately.

For raw admission choose \(0<\eta\ll\delta\), then require
\(2sR<\eta\). Every orbit crosses \(sq=\eta\) in the strip.
Contraction on the fixed remaining raw interval, compared with
\(w_\Gamma(x)=(\sqrt{1+x^2}+1)^{-1}\), gives
\[
 |W_H-w_\Gamma(\delta+\gamma)|
 \le C(\eta/(\delta+\gamma))^\mu
 +C_{\eta,\delta}(|A-1|+|C-1|+|\gamma|+sp+s^2).
\]
The raw height is \(z_H=(\delta+\gamma)^2W_H\); its organizing
height shifts by \(O(\delta|\gamma|)\). Choose \(\eta/\delta\)
from the fixed window reserve, then shrink the coefficients. On this
whole admitted image the divided adapter \(\chi=\mathcal C_{\theta,m}\)
has the ordinary \(C^3\) bounds of Lemma~\ref{phys:lem-c3-adapters}.
Set \(\Gamma_-(r)=\chi(z_H(r))\). Its observer identities and the
preceding two rows give
\begin{equation}\label{ann:repair-completed-log-jets}
 |\partial_\upsilon\log(-\Gamma_-')|
 +|\partial_\upsilon^2\log(-\Gamma_-')|\le C.
\end{equation}
Finally separate \(F_W=-2+\lambda_\gamma+(qf_Y-\lambda_\gamma)\)
in the moving cocycle:
\begin{equation}\label{ann:negative-raw-gate-factor}
 -z_H'=\Omega_\theta r^{1-\lambda_\gamma}\widehat U_-,\qquad
 \Omega_\theta=s^{2-\lambda_\gamma}(\delta+\gamma)^{\lambda_\gamma},
\end{equation}
\begin{equation}\label{ann:repair-unit-factor}
 \widehat U_-=F(a,w_0)a_1(q_s/r)^{2-\lambda_\gamma}e^{I(r)},\qquad
 U_-=\chi'(z_H)\widehat U_-.
\end{equation}
The factors \(\widehat U_-,U_-\) are positive units, with
\(\Omega_\theta\) kept separate. Subtracting
\((1-\lambda_\gamma)\upsilon+\log\Omega_\theta\) in
\eqref{ann:repair-completed-log-jets} bounds two ordinary derivatives
of \(\log U_-\). Equation~\eqref{ann:repair-action-budget} then gives
\(\mathcal D_-\log U_-=O(d_Q)\) and
\(\mathcal D_-^2\log U_-=O(d_Q^2)\), proving the assertion.
The same bounds before the adapter give
\(\mathcal D_-^jz_H=O((sr)^{2-\lambda_\gamma}d_Q^j)\), \(j=1,2\).
\end{proof}

At the other saddle the characteristic number remains compact.
Its Dulac estimate must still include the actual incoming coordinate,
because the curvature is taken in the physical affine input.

\begin{lemma}[The fixed-cut height and the actual positive half-gate]
\label{ann:p0-positive-interface}
For fixed \(p\in[p_0,P]\Subset(0,\infty)\), physical windows,
and then fixed \(Q\), there is a common positive radius \(H_*\)
such that the map from the local positive-saddle exit to \(q=Q\)
is \(\Psi_\theta(v)\), with \(\Psi_\theta(0)=0\), bounded
\(C^3\) norm, and derivative bounded above and away from zero.
The actual incoming half-gate \(x=\Gamma_+(H)\), on a slightly
larger interval than \([0,H_*]\), satisfies
\[
 \Gamma_+'>0,\qquad
 \left|\frac{\Gamma_+''}{\Gamma_+'}\right|
 +\left|\frac{\Gamma_+'''}{\Gamma_+'}\right|\le C_Q.
\]
Its exact value-and-slope normalization and its inverse on a smaller
normalized image have uniform \(C^3\) bounds. Its full raw image
is admitted by the preassigned incoming window.
For every fixed \(0<\overline H\le H_*\),
\(y_+(\overline H)=Y_Q^+(\overline H)\) has a uniform positive
lower bound. No uniform lower bound for the raw \(\Gamma_+'\)
or bound for the raw inverse is asserted.
\end{lemma}
\begin{proof}
The fixed rescaled axis segment from the outgoing saddle cut to \(q=Q\)
has a positive speed margin. Its finite hitting map \(\Psi\) therefore
has \(\Psi(0)=0\), positive derivative units and uniform \(C^3\)
bounds; no raw-tail scale enters this map.

For the incoming half-map, follow the stable separatrix backwards. In
the limiting field its denominator \(D=Y+1-pX\) cannot first vanish,
since \(D_X=XY/D-p>0\) near a proposed zero; a quadratic upper
barrier excludes finite blow-up. The compact half-family consequently
enters a common parabolic band at a fixed \(X=L_0\). A buffered chord
has a finite hitting map \(A_\theta(H)=Y(L_0)/L_0^2\) with positive
derivative units and uniform \(C^3\) bounds, including \(H=0\).

On the outward tail \(F_w\le-\lambda<0\), and the state derivatives
through order three are bounded. The relative-jet argument
\eqref{tail:relative-jet-equations}, with
\(0<J_1\le e^{-\lambda t}\), bounds \(J_2/J_1,J_3/J_1\).
Exact value-and-slope normalization thus gives the asserted map and
inverse bounds. The fixed-port argument of
Lemma~\ref{ann:p0-fixed-tail-ports} admits its whole raw chord: that
argument needs only this common parabolic band, retains the entrance
memory and \(O(d|\gamma|)\) port shift, and reaches the original
window through its fixed raw tube. Compose the ordinary bridge on this
admitted image.

Finally Lemma~\ref{lem:uniform-positive-dulac} and \(\Psi\) give
\(y_+(H)=c_+H^{\rho_+}e^{\mathcal R_+(H)}\), with \(c_+\) a
positive unit. Thus for fixed \(0<\overline H<1\),
\(y_+(\overline H)\ge c_Q\overline H^{\rho_1}e^{-M_Q}>0\).
This height bound involves no lower bound on the raw incoming slope.
\end{proof}

\begin{lemma}\label{ann:lem-positive-dulac}
There is a fixed incoming coordinate \(H>0\) at \(x_+\) for which
the height on \(q=Q\) is
\begin{equation}\label{ann:positive-dulac-map}
 Y_Q^+(H)=c_+H^{\rho_+}e^{\mathcal R_+(H)}.
\end{equation}
For \(0<\epsilon<\min(1,\rho_0)\), the Euler derivatives of
\(\mathcal R_+\) through order two are bounded by
\(CH^\epsilon(1+|\log H|)\).
The continuation to the input \(x\) in \eqref{eq:return}
has the form \(x=\Gamma_+(H)\), with
\begin{equation}\label{ann:positive-gate-derivative}
 \Gamma_+'(H)=\Omega_+(\theta)U_+(H)>0,
\end{equation}
where \(U_+\) and its reciprocal are bounded and its first two
Euler logarithmic derivatives have the same vanishing bound.
\end{lemma}
\begin{proof}
Apply Lemma~\ref{lem:uniform-positive-dulac} to the rates in
Lemma~\ref{ann:lem-axis-roots}, then compose with the fixed-cut map
\(\Psi\) above. Hadamard division gives
\(\Psi(v)=\Psi'(0)ve^{j(v)}\), \(j(0)=0\), with bounded \(C^2\)
norm. Its first two Euler rows along the Dulac map are
\(O(H^{\rho_+})\), so this composition preserves the asserted remainder.
For the incoming half-gate set
\(\Omega_+=\Gamma_+'(0)\), \(U_+=\Gamma_+'(H)/\Gamma_+'(0)\).
The preceding relative \(C^3\) bounds give
\(\mathcal E\log U_+,\mathcal E^2\log U_+=O(H)\),
\(\mathcal E=H\partial_H\). The final input differs from raw height
only by the translation \(a_\theta\).
\end{proof}

\subsection{Endpoint curvature and the total count}

\begin{lemma}[Exhaustion of the low central heights]
\label{ann:p0-low-height-exhaustion}
There are fixed \(R_a<R_b\), \(R_a>2Q\), and \(u_L>0\), chosen
before the final coefficient shrink, such that every actual through
orbit with \(0<u:=Y(0)\le2u_L\) is represented uniquely by the
same negative strong footpoint \(R_b<r<R\) and by the positive
Dulac coordinate \(0<H<\overline H\). The identification on the
fixed cut is
\[
 N_\theta(u)=y_\theta(r)=Y_Q^+(H),
 \qquad N_\theta(u):=Y(Q;Y(0)=u).
\]
All actual input and output windows are those already fixed; zero
height is not included as a through orbit.
\end{lemma}
\begin{proof}
Start at \((Q,y)\), \(0<y<w_0Q^2\). Until the strong-section hit,
the footpoint \(p_f=\pi(q,Y)\) increases below \(R\), and the
projection estimates give
\[
 (\log q)'\ge h(\log p_f)+(1-C\epsilon)w>0,\qquad
 q<2p_f<2R,\qquad (\log w)'\ge\nu_0>0.
\]
Thus there is one transverse hit in finite normal time, inside the
unchanged corridor; it may have \(q_s>R\). Scalar first variation
and \(\partial_q(Y-w_0q^2)=q_sF(\log q_s,w_0)>0\) give
\(dq_s/dy<0\), hence \(dr/dy<0\). The reverse construction in
Lemma~\ref{lem:normal-time-action-height} has
\(y_r=-Ky/b<0\), \(K\ge1/2\), and \(y\to0\) at \(R\).
It is therefore a bijection \([R_a,R)\to(0,y(R_a)]\), identical
to the forward construction by uniqueness.

On \(0\le q\le Q\), choose a small signed height strip with
\(M_-+L_-Y>c>0\). Its finite map satisfies
\(N(0)=0\), \(c_Q\le N'\le C_Q\), with common finite-order bounds.
Choose \(\overline H\) for the positive Euler errors, then \(R_a\)
so that \(e_{R_a}<\inf y_+(\overline H)\) and the fixed-cut errors
are small; take \(R_0=R_a\). On every fixed footpoint interval the
exact action-height formula bounds the unit and action, giving
\(\inf y(R_a)>0\). Choose \(R_b>R_a\) with
\(e_{R_b}<\inf y(R_a)\), then require \(R>2R_b\).
The same fixed-interval formula gives
\[
 y(R_b)=R_b^2\mathcal U(R_b)
 e^{-\mathcal A(R_b)+E_Q(y(R_b))}\ge\eta_b>0.
\]
Choose \(2u_L\) below the domain of \(N\) and \(\eta_b/C_Q\).
Then \(N(u)<y(R_b)\) selects a unique \(r>R_b\); it is also below
\(y_+(\overline H)\) and selects a unique positive Dulac preimage.
Equality at \(q=Q\) identifies the orbit, so both previously admitted
gates and the divided closing apply throughout this interval.
\end{proof}

The exhaustion lemma identifies the same orbit in the two endpoint
coordinates. Equating their descriptions of \(Y_Q\) gives
\begin{equation}\label{ann:endpoint-matching}
 \rho_+\log H+\mathcal R_+(H)
       =-\mathcal A_-(r)+\mathcal R_-(r)+c_\theta .
\end{equation}
Set \(\mathcal D=H\partial_H\) and
\(K_-=1-\mathcal D_-\mathcal R_-\ge1/2\).
Then
\begin{equation}\label{ann:rH}
 Hr_H=-\frac{\rho_++\mathcal D\mathcal R_+}{K_-}
                  \frac{M_-(r)}{rB_-(r)}<0 .
\end{equation}
By \eqref{ann:negative-small-height} and
\eqref{ann:positive-dulac-map}, throughout this endpoint
\[
 H\le C_QR_0^{2/\rho_0}e^{-c(R_0-Q)/(P\rho_1)}.
\]
Thus the same choice of \(R_0\) makes all positive Dulac remainders
small on the entire interval \(R_0\le r<R\).

\begin{proposition}\label{ann:prop-endpoint-curvature}
The complete map
\(x=\Gamma_+(H)\mapsto\Pi_\theta(x)=\Gamma_-(r(H))\)
has \(\Pi_\theta''(x)<0\) on this endpoint. Its displacement has
at most two zeros counted with multiplicity.
\end{proposition}
\begin{proof}
Write \(y(H)=Y_Q^+(H)\) and
\(\mathcal O(y(r))=\Gamma_-(r)\). Put
\(b=M_-/(rB_-)\), \(\mathcal D_-=b\partial_r\), and
\(K=-\mathcal D_-\log y\ge1/2\). The action-height jets give \(K=1+O(d_Q+e_{R_a})\) and
\(\mathcal D_-K=O(d_Q+e_{R_a})\). The coefficient quotients are
\begin{align}
 M_-'/(rB_-)&=(p/r+6\gamma-3sr)/B_-=O(d_Q),
 \label{ann:Mprime-quotient}\\
 M_-/(r^2B_-)&=(r^{-2}+p/r+3\gamma-sr)/B_-=O(d_Q),
 \label{ann:M-quotient}\\
 sM_-/(rB_-^2)&=O(d_Q).\label{ann:B-quotient}
\end{align}
Thus \(b'=O(d_Q)\). The identity
\(\mathcal O'=(-\Gamma_-')b/(yK)\), together with
Lemma~\ref{ann:lem-negative-gate}, gives
\[
 e_{\mathcal O}:=1+y\frac{\mathcal O''}{\mathcal O'}
 =-\frac{\mathcal D_-\log(-\Gamma_-')+b'-\mathcal D_-\log K}{K},
 \qquad |e_{\mathcal O}|\le Cd_Q+C_Qe_{R_a}.
\]
No reciprocal negative saddle rate occurs.
Set \(\mathsf A=\rho_++H\mathcal R_+'\). In the actual input
\(x=\Gamma_+(H)\), the full chain rule is
\begin{equation}\label{ann:complete-log-curvature}
 H\partial_H\log\Pi_x'
 =-1+\mathsf A e_{\mathcal O}
      +H\partial_H\log\mathsf A-H\partial_H\log\Gamma_+'.
\end{equation}
Choose \(Q\) first. For this fixed cut choose \(\overline H\)
so that \(|\mathsf A-\rho_+|\le\min(1,\rho_0/2)\) and the last
two terms each have absolute value at most \(1/16\).
Choose \(R_a\) using the positive fixed-cut image margin, as in
Lemma~\ref{ann:p0-low-height-exhaustion}, and then shrink the
coefficients so that
\( |e_{\mathcal O}|\le1/[16(\rho_1+1)]\).
The right side is at most \(-13/16\). Since
\(\Pi_x'>0\), \(H\Gamma_+'>0\),
\[
 \frac{\Pi_{xx}}{\Pi_x'}
 =\frac{H\partial_H\log\Pi_x'}{H\Gamma_+'}<0.
\]
For the physical multiplier \(m>0\), \(\Delta''=-m\Pi''>0\);
Rolle's theorem with multiplicities gives the endpoint count.
\end{proof}

\begin{lemma}[Fixed parabolic tails and the physical sections]
\label{ann:p0-fixed-tail-ports}
Fix the physical sections \(\Sigma_\pm=\{v=\pm\delta\}\), their
nested windows, and \(0<d<\delta\le1/32\), with \(d\le1/100\).
On an enlarged compact strict source family, one can choose a fixed
\(L\), and then sufficiently small bounds on the raw coefficients
and on \(s>0\), so that the two outward graphs reach
\(v=\gamma\pm d\). Their entire images are admitted by uniformly
regular increasing maps
\[
 \mathcal B_{\sigma,\theta}:z(\gamma+\sigma d)
       \longmapsto z(\sigma\delta),\qquad \sigma=\pm1.
\]
These maps and their inverses have uniform phase \(C^3\) bounds on
slightly larger fixed windows. The divided closing map is defined on
the full negative image for \(0\le m\le m_{\rm ad}\), where
\(m_{\rm ad}>0\) is fixed before any zero-counting multiplier bound.
\end{lemma}
\begin{proof}
Choose \(|\gamma|<\min\{d/2,(\delta-d)/2,1/500\}\).
For \(r=Le^t\), \(\xi=sr\), write \(w=Y/r^2\) on either tail.
Its exact equation is
\[
 \dot w=F_\sigma(t,w)
 =w\left(\frac{n_{0,\sigma}+n_{1,\sigma}w}
                   {\ell_\sigma w+d_\sigma}-2\right),
\]
where
\[
\begin{aligned}
 v_\sigma&=\gamma+\sigma\xi,
 &\ell_\sigma&=1-Av_\sigma+Cv_\sigma^2,\\
 d_\sigma&=3\gamma+\sigma\xi-\sigma p/r-a/r^2,
 &n_{0,\sigma}&=1+2\gamma+\sigma\xi,
 &n_{1,\sigma}&=\sigma C\xi v_\sigma.
\end{aligned}
\]
The entrance \(L\) is chosen on the enlarged limiting family first.
Continuous dependence on the fixed core then puts both entrance
values strictly inside \([1/4,3/4]\). Increase \(L\), if necessary,
so that \(|p|/L\le1/40\) and \(|a|/L^2\le1/4000\), and require
\(A,C\in[9/10,11/10]\) and \(sL<d\). The fixed core is retained
when these choices are made. On
\[
 1/4\le w\le3/4,\quad \xi\le1/100,
 \quad |\gamma|\le1/500,
 \quad |p|/r\le1/40,
 \quad |a|/r^2\le1/4000
\]
the denominator is bounded below and the vector field points strictly
inward at both edges. Differentiating this rational quotient, as in
\eqref{ann:repair-rational-identities}, gives
\begin{equation}\label{ann:p0-tail-envelope}
 F_w\le-\lambda,\qquad \lambda=7/5,
 \qquad |F_{ww}|\le10,
 \qquad |F_{ww}|\le240|d_\sigma|.
\end{equation}
Thus the band remains invariant up to
\(T=\log(d/(sL))\), independently of its length.

The invariant band controls the graph throughout the long tail. To
compose with the regular transitions, we must also place its whole image
inside the fixed physical windows.
At the organizing coefficients the exact radial graph is
\[
 w_\Gamma(\xi)=\frac1{\sqrt{1+\xi^2}+1},\qquad
 \xi\partial_\xi w_\Gamma=F^0_\sigma(\xi,w_\Gamma).
\]
On the preceding coefficient box, the mean-value formula gives
\[
 |F_\sigma-F^0_\sigma|
 \le100\bigl(|\gamma|+|A-1|\xi+|C-1|\xi^2
                         +|p|/r+|a|/r^2\bigr).
\]
For \(T=\log(d/(sL))\), contraction gives the uniform error
\[
 \begin{split}
 \mathcal E_\theta(L,T)={}&e^{-\lambda T}
       |w_\sigma(0)-w_\Gamma(sL)|\\
 &+100\int_0^T e^{-\lambda(T-t)}
 \left(|\gamma|+|A-1|sLe^t+|C-1|(sLe^t)^2
       +\frac{|p|}{L}e^{-t}+\frac{|a|}{L^2}e^{-2t}\right)dt.
 \end{split}
\]
In particular \(|w_\sigma(T)-w_\Gamma(d)|\le\mathcal E_\theta\).
This retains the entrance memory; \(1<\lambda<2\) bounds both
decaying-power convolutions. For fixed \(L\) and bounded \(a,p\),
\(\mathcal E_\theta\to0\) uniformly as
\(s,\gamma,A-1,C-1\to0\).
With \(g(v)=\sqrt{1+v^2}-1\), the actual port heights satisfy
\[
 |z_\sigma-g(\gamma+\sigma d)|
       \le d^2\mathcal E_\theta(L,T)+2d|\gamma|.
\]

On the fixed annulus \(d/2\le|v|\le2\delta\), the organizing graph
has \(-D_0(v,g(v))\ge v^2/4\ge d^2/16\).
Choose a fixed tube about this graph, and then shrink the raw
coefficients so that \(-D_\theta\ge d^2/32\) there. The actual
graph equation
\[
 z_v=\frac{z(R_\theta(v)+Cvz)}
                {[1-Av+Cv^2]z-P_\theta(v)}
\]
has bounded state derivatives through order three in the tube.
Gronwall's inequality, with the port error just obtained, gives
admission to that tube by a first-exit argument. Its finite physical
length then gives the asserted regular maps and inverse bounds.
First place their entire images strictly inside the prescribed inner
windows. Next choose \(m_{\rm ad}\) so that
\(m d_\theta(K_-^\sharp)\) lies strictly inside the fixed domain of
\(f_\theta\). Only after these domain assertions do we compose the
divided closing map.
\end{proof}

\begin{lemma}[Actual curvature on a compact positive source block]
\label{ann:p0-middle-actual-curvature}
Fix \(0<p_0\le p\le P\) and a buffered positive central-height
interval \(0<u_-\le u\le u_+\). For \(a=-1\), after a common
coefficient shrink the complete transition has \(\Pi_{xx}<0\) on
this whole block in the actual affine input. The assertion is
uniform for \(0\le m\le m_{\rm ad}\).
\end{lemma}
\begin{proof}
We first construct the limiting strict block. On the two half-graphs
\(X=\sigma r\),
\[
 Y_r=\frac{rY}{Y+1-\sigma pr},\qquad Y(0)=u.
\]
The height increases and is bounded above by
\(u_++1+r^2+2Pr\). The negative denominator is automatically
positive. On the positive half, at a sufficiently small positive
denominator its derivative \(rY/D-p\) is positive, which excludes
a first zero. More quantitatively,
\[
 D\ge\min\{1/2,u_-/(4P^2)\}>0.
\]
These inequalities exclude finite blow-up and give both infinite
half-graphs. A common parabolic entrance is fixed as follows. For
\(Q=\max\{4,16P\}\), \(w=Y/r^2\), and \(t=\log r\), the source
equation satisfies
\[
 0<w\le1/4\ \Longrightarrow\ (\log w)_t\ge2/3,
 \qquad w\ge3/4\ \Longrightarrow\ (\log w)_t\le-2/5.
\]
Its entrance bounds at \(Q\) are uniform. After a fixed additional
logarithmic time both half-graphs enter the strict parabolic band.

The resulting enlarged strict family is admitted on the actual
sections by Lemma~\ref{ann:p0-fixed-tail-ports}. On this compact
family, Proposition~\ref{ann:prop-source-curvature} gives
\(S''=-p\mathcal I_{bb}\) with a uniform negative margin since
\(p\ge p_0\). The positive affine normalization of \(S\) has the
same margin up to bounded positive factors.
Proposition~\ref{prop:source-normalization} with phase order two
therefore gives \(\Phi''<0\) after a common coefficient shrink.

Retain the exact centres and gains of the actual maps \(h,k\) in
\eqref{eq:reference-normalization}, including both bridges and divided
closing. With \(c_h=h'(b_0)>0\), \(c_k=k'(b_0)>0\), one has
\begin{equation}\label{ann:p0-two-actual-adapters}
 \Pi_{xx}(x)=\frac{c_k}{c_h^2}\Phi''\!\left(
 b_0+\frac{x+a_\theta-h(b_0)}{c_h}\right)<0.
\end{equation}
The curvature sign is preserved by the positive prefactor, even when
both absolute gains tend to zero.
\end{proof}

\begin{lemma}[Paired first variations at high central height]
\label{ann:p0-paired-high-derivative}
For the fixed compact positive-ratio region, there are fixed
\(K_H,z_*,C_H>0\) such that the complete actual transition has
\(0<\Pi_x'\le C_H\) on
\[
 u=Y(0)>K_H(1+p^2),\qquad 0<s^2u<z_*.
\]
The constant is independent of the tested phase point. The actual
multiplier restriction \(mC_H<1/2\) gives \(\Delta'>1/2\).
\end{lemma}
\begin{proof}
The derivative assertion is Proposition~\ref{sharp:upper-high} with
\(a=-1\). Its paired absolute variations
\eqref{sharp:upper-tail-ratio} cancel the common
\(s^{2-\lambda_\gamma}\) power before the actual bridges and
ordinary divided closing are composed. For the high range freeze
\(H=s^2u\); the signed normalization in
Corollary~\ref{rew:cor-high-derivative-signed} gives the same
\(C_H\) at every phase point.

For the domain assertion, apply
Lemma~\ref{ann:p0-fixed-tail-ports} to that signed compact source
block, fixing its core, tail entrance, ports and raw-bridge tubes
before shrinking the coefficients and \(z_*\). This admits both
full raw images, including the positive-height equality and
double-zero faces. Shrink \(z_*\) further so that
\[
 \sigma_\theta\le-D_\theta(\gamma,H)
   =H(\ell_\gamma-\bar a)\le2H<d_{\rm out}.
\]
Restrict the receiver to \(H<\min\{h_{\rm out},z_*\}\), and
choose \(m<\min\{m_{\rm ad},1/(2C_H)\}\) last.
\end{proof}

\begin{proof}[Proof of Lemma~\ref{ann:lem-compact-positive-domains}]
Fix the coefficient interval, finite smoothness, physical windows and
closing-domain cap \(m_{\rm ad}\). Lemma~\ref{ann:p0-low-height-exhaustion}
chooses \(Q,\overline H,R_a=R_0,R_b,u_L\) in that order and matches
all \(0<u\le2u_L\) through the fixed-cut equality
\(N(u)=y(r)=Y_Q^+(H)\). Its full raw images are admitted.
Independently fix the core, tail entrance and raw tubes for
Lemma~\ref{ann:p0-paired-high-derivative}, obtaining \(K_H,C_H,z_*\).
Choose \(u_H>\max\{2u_L,K_H(1+P^2)\}\) and apply
Lemma~\ref{ann:p0-middle-actual-curvature} on \([u_L/2,3u_H]\).
Intersect these coefficient bounds, including \(R>2R_b\), without
changing \(Q\).

Choose the unscaled cap and coefficient bounds so that
\[
 0<h_{\rm cap}\le z_*,\qquad
 L_{\rm raw}h_{\rm cap}<d_{\rm out}/4,\qquad
 |\alpha_0|<d_{\rm out}/4,
 \quad 0<\ell_\theta(\gamma)\le L_{\rm raw}.
\]
Then every through orbit with \(H=z(\gamma)<h_{\rm cap}\) satisfies
\(\sigma_\theta\le-D_\theta(\gamma,H)
=-\alpha_0+\ell_\theta(\gamma)H<d_{\rm out}/2\).
Let \(J\) be the original buffered receiver clipped by
\(H<\min\{h_{\rm cap},h_{\rm out}\}\), in the affine coordinate
\(x=z_{\rm in}-a_\theta\). The final global choice takes
\(h_{\rm out}\le h_{\rm cap}\), as in Section~\ref{sec:conclusion}.

The strictly increasing central height
\(u(x)=z_\theta(\gamma;a_\theta+x)/s^2\) gives the relative open
intervals
\[
 I_E=J\cap\{0<u<2u_L\},\qquad
 I_M=J\cap\{u_L<u<2u_H\},\qquad
 I_H=J\cap\{u>u_H\}.
\]
They cover \(J\), including all four internal equality levels, and
have physical overlaps \(u_L<u<2u_L\) and \(u_H<u<2u_H\).
The endpoint estimate holds on the larger \(r>R_a\) domain, the
middle estimate on \([u_L/2,3u_H]\), and the high estimate below
\(z_*\). Original window and speed clips preserve this cover; the
fixed-cut match has its own image margin. Raw overlap lengths need no
positive lower bound.

Uniqueness identifies these as restrictions of the same completed map
\eqref{eq:return}. Proposition~\ref{ann:prop-endpoint-curvature} and
Lemma~\ref{ann:p0-middle-actual-curvature} give \(\Pi_{xx}<0\) on
\(I_E\cup I_M\), while
Lemma~\ref{ann:p0-paired-high-derivative} gives
\(0<\Pi_x'\le C_H\) on \(I_H\). All cuts, caps and estimates
precede the final multiplier restriction and geometric capture.
\end{proof}

\section{Gamma-dominant and compact cubic passages}
\label{sec:gamma-cubic-passages}

When axis roots coalesce, the incoming and outgoing root distances can
tend to zero independently. We estimate the complete return
through these limits in the gamma-dominant and compact cubic charts.
The discriminant separates real, multiple and ghost roots, and the
far-root ranges are retained. In each case the estimates apply to the
same physical input used for zero counting.
\subsection{Gamma-dominant and root-coalescence faces}
\label{subsec:deg-gamma-dominant}

In the chart \eqref{deg:gamma-chart}, negative \(\gamma\) leaves two
separated incident roots as \(k\to0\). For positive \(\gamma\), the
positive roots may coalesce and the negative root escapes. These two
geometries require different comparisons: a physical endpoint power law
on the negative face, and weighted curvature on the positive face.
We first treat the negative face and the positive far-root face, then
the bounded positive face. In each case the displacement is the original
\(F(x)=x-m\Pi_{\theta,m}(x)\) of \eqref{eq:return}, with
\(x=h-a_\theta\). Fix the original passive compact set, nested windows
and their reserves before the regional cuts. Every carrier below is
\[
 I_\theta=\{x:a_\theta+x\in J_\theta^{\rm out},\quad
                  0<H(x):=z_\theta(\gamma;a_\theta+x)<H_{\rm cap}\},
\]
where \(0<H_{\rm cap}<z_*\) is a positive unscaled cap chosen before
the coefficient and multiplier bounds. Thus every set retains the original
input window, inverse output window, stop and speed clips, and all
intervening nonroot inputs. Phase derivatives fix every coefficient and
\(m\); a stopped finite channel contributes no cycle.

\begin{proposition}[Negative quadratic face]
\label{deg:prop-negative-gamma}
Fix \(\Pi_{\max}<\infty\) and a sufficiently small routing bound
\(k_\flat\) from Lemma~\ref{deg:lem-small-source-routing}.
For \(\gamma=-g<0\), \(0<k\le2k_\flat\), and
\(0\le\pi\le\Pi_{\max}\), independently small \(g\) and \(m\)
give the physical endpoint law \eqref{deg:negative-gamma-physical-power}.
The endpoint, core and high intervals cover the original \(I_\theta\),
with strict estimate overlaps and positive dividing heights retained.
Every critical point of the same \(F\) is simple and has \(F''>0\).
Consequently the whole coefficient fibre has at most one critical point
and at most two isolated displacement zeros, counted with multiplicity.
\end{proposition}

\begin{proof}
We first match the two root passages on fixed normal sections, retaining
the middle multiplier in their logarithmic actions. The exterior maps
then turn this match into the physical power law. A comparison orbit
joins its endpoint interval to the bounded core and high-height interval,
so the final critical-point argument applies to one connected carrier.

Put \(\epsilon=gk\). The exact dictionary is
\(v-\gamma=\epsilon\xi\), \(z=\epsilon^2\eta=s^2Y_0\),
\(\eta=gY_0\), \(s=g^{3/2}k\). After positive time scaling,
\[
 \dot\xi=-gP-Q\eta,\qquad \dot\eta=-\eta(H_0+H_1\eta),
\]
\[
 \begin{gathered}
 P=1-\pi\xi-3\xi^2+k\xi^3,\quad
 Q=1+Ag(1-k\xi)+Cg^2(1-k\xi)^2,\\
 H_0=(1-2g)\xi+kg\xi^2,\quad H_1=-Ckg^2(1-k\xi).
 \end{gathered}
\]
Its incident roots \(a>0,-b<0\), their distances from zero,
\(|P'(a)|,P'(-b)\), and normal rates are positive units on the
closed \((\pi,k)\) block. The third root lies beyond a fixed large
cut \(L\). Choose independent fixed normal sections and inner weak
cuts at the two incident roots. The intrinsic construction
\eqref{eq:saddle-normal-time}--\eqref{eq:exact-saddle-invariant}
has kernels \(C(1+T)^Ne^{-(c-Cg)T}\); four differentiated integral
equations give uniformly bounded strong projections and normal coordinates.
With inward footpoint distances \(t,r>0\), put
\(A_+=-H_0(a)/P'(a)\), \(A_-=-H_0(-b)/P'(-b)\),
\(\nu=A_+/A_-\). These are positive units. The anchored normal-time
actions satisfy
\begin{equation}\label{deg:negative-gamma-actions}
 \begin{gathered}
 J_+(t)=-A_+\log t+j_+(t),\quad J_-(r)=-A_-\log r+j_-(r),\\
 |(\delta\partial_\delta)^j(j_\pm(\delta)-j_\pm(0))|\le C\delta,
 \quad j\le3.
 \end{gathered}
\end{equation}
Indeed, in the intrinsic equations
\(\dot\Xi=-\rho\Xi,\ \dot Z=Z,\ \rho=g/A_\pm\),
\(|\Xi|Z^\rho\) is constant and
\(J=-(g/\rho)\log(|\Xi_n|/|\Xi_m(0)|)\); its axis derivative
is \(H_0/P\). Pole subtraction gives the displayed smooth remainders.
On each actual cut write
\(|\Xi_n|=\delta a_n(\delta)\), \(Z_n=b_n(\delta)\),
\(Z_m=z_0wU_m(w)\). The positive units have uniform \(C^3\) bounds;
Taylor division uses the intrinsic fourth derivative. The exact passage is
\[
 w=w_0e^{\Psi_g(w)},\qquad
 w_0=(b_n/z_0)e^{-J/g},\qquad
 \Psi_g=\rho^{-1}\log(|\Xi_m(0)|/|\Xi_m(w)|)-\log U_m(w).
\]
Here \(|\Psi_g|\le Cw/g\), \(|\partial_w^j\Psi_g|\le C/g\),
\(j\le3\). Choose each fixed collar before \(g\), so that
\(w_0\le\delta^{A_{\min}/(2g)}\). Then
\(1-w\Psi_g'(w)\ge1/2\) and \(w\asymp w_0\). Implicit Euler
differentiation gives
\(|(\delta\partial_\delta)^j\log(w/w_0)|\le Cg^{-(j+1)}w\),
\(j\le3\), hence ordinary \(C^3\) extension to zero. Thus,
independently on both sections, \(w_\pm=u_\pm e^{-J_\pm/g}\),
where \(u_\pm\) and their inverses are uniform \(C^3\) units.

Take \(w\) to be the raw height \(\eta\) on each fixed inner
weak cut, where the normal restriction is \(Z_m=z_0wU_m(w)\).
On the core \(P\ge p_c>0\),
the scalar equation is
\(f=\eta(H_0+H_1\eta)/(gP+Q\eta)\).
Choose an actual input chord \(0\le w\le e^{-K_0/g}\), with
\(K_0\) fixed large enough that a first-exit comparison keeps its
whole passage and output inside the two weak-cut buffers.
Here \(|f_{\eta^h}|\le C_hg^{-h}\). The initial-height rows
\(V_j=\partial_w^j\eta\), through order five, satisfy
\[
 V_j'=f_\eta V_j+
 \sum_{h=2}^j f_{\eta^h}\mathcal B_{j,h}(V_1,\ldots,V_{j-h+1}),
\]
where \(\mathcal B_{j,h}\) are the partial Bell polynomials.
The linear kernel is only \(O(g^{-1})\); successive variation of
constants puts the higher inverse powers in the forcing and gives
\mbox{\(\|M_g\|_{C^5}+\mu_g^{-1}\le e^{K/g}\)}, after enlarging \(K\).
For the two fixed physical cuts the linear multiplier is exactly
\begin{equation}\label{deg:negative-core-multiplier}
 g\log\mu_g=\int_{\xi_{i,+}}^{\xi_{i,-}}\frac{H_0}{P}\,d\xi .
\end{equation}
Taylor's integral division
\(N_g(w)=\mu_g^{-1}\int_0^1(1-t)M_g''(tw)\,dt\)
therefore gives the actual middle map
\begin{equation}\label{deg:negative-gamma-middle}
 M_g(w)=\mu_gw(1+wN_g(w)),\qquad
 \mu_g^{\pm1}+\|N_g\|_{C^3}\le e^{K/g}.
\end{equation}
Choose \(t_0\) before \(g\), with its action exceeding all these fixed
exponents. The weak input then lies on the actual positive core chord,
and \(g(t\partial_t)^j\log(1+wN_g(w))=O(t^N)\), \(j\le3\),
for any fixed \(N\). Matching the actual heights, while retaining
\(g\log\mu_g\), gives
\begin{equation}\label{deg:negative-gamma-matching}
 \begin{gathered}
 A_-\log r-A_+\log t=C_\theta+e_+(t)+e_-(r),\qquad |C_\theta|\le C,\\
 |(t\partial_t)^je_+|\le C(t+t^N),\quad
 |(r\partial_r)^je_-|\le Cr,\quad j\le3.
 \end{gathered}
\end{equation}
The negative height is strictly monotone from zero, and
\(A_--r\partial_re_->0\). Reducing \(t_0\) admits a unique actual
output in its independent fixed negative collar. Consequently
\begin{equation}\label{deg:negative-gamma-power}
 \begin{gathered}
 r=A_\theta t^\nu e^{E_\theta(t)},\qquad c\le A_\theta,A_\theta^{-1}\le C,\\
 |(t\partial_t)^jE_\theta|\le Ct^\beta,\quad j\le3,
 \qquad0<\beta<\min(1,\inf\nu).
 \end{gathered}
\end{equation}

Continue both sections to \(\pm L\). Their fixed-normal fast graphs
have limit \(\eta=(\xi^2-u^2)/2\); at the cuts their heights,
denominators and gap derivatives are units. At a fixed small raw cut
\(d\), differences of the bounded \(L\)-labels are
\(O(\epsilon^{2-\lambda_\gamma})\), while the reference graph has
\(z/d^2=1/2+O(d)+o(1)\). The raw corridor first variation is
\(C(d_1/d)^{\lambda_\gamma}\) times a unit. Choosing \(d\) first
and then \(g\) therefore propagates the reference to the original
organizing-arc windows uniformly on the retained passive compact set.
This admits the whole outward chord, including its limiting zero end,
before applying its ordinary observers; it imposes no condition on
\(g\log k\). The exact tail calculation
\eqref{ann:tail-first-derivative}--\eqref{eq:N-bound} yields
\[
 d_\pm=\epsilon^{2-\lambda_\gamma}U_\pm,\qquad
 \lambda_\gamma=-12g/(1-2g),\qquad
 c\le U_\pm,U_\pm^{-1}\le C.
\]
Its first remainder is integrable with majorant
\(C(\epsilon+g\pi R^{-2}+R^{-3+\zeta})\), \(0<\zeta<1\);
the higher normalized kernels give ordinary \(C^3\) embeddings.
Only \(d_-\) includes the divided closing, whose first derivative is
\(f'(md)d'\). With exact centers and slopes retained separately,
\[
 x=x_{\rm st}+d_+\phi_+(t),\qquad
 \Pi(x)=y_{\rm st}+d_-\phi_-(r(t)),\qquad
 \phi_\pm(0)=0,\quad\phi_\pm'(0)=1.
\]
These maps and their inverses are uniformly \(C^3\). Differentiating
this identity, and putting \(q=(x-x_{\rm st})/d_+=\phi_+(t)\), gives
\begin{equation}\label{deg:negative-gamma-physical-power}
 \begin{gathered}
 \Pi'(x)=B_{\theta,m}q^{\nu-1}e^{R(q)},\qquad B_{\theta,m}=(d_-/d_+)A_\theta\nu,\\
 c\le B_{\theta,m},B_{\theta,m}^{-1}\le C,\qquad
 |(q\partial_q)^jR|\le Cq^\beta,\quad j\le2.
 \end{gathered}
\end{equation}
In fact
\(R(\phi_+(t))=E+\log(1+tE'/\nu)+\log\phi_-'(r)
-\log\phi_+'(t)+(\nu-1)\log(t/\phi_+(t))\).
The third Euler row in \eqref{deg:negative-gamma-power} controls the two
Euler derivatives of this expression. The absolute amplitude
\(B_{\theta,m}d_+^{1-\nu}\) need not be a unit.

Choose the comparison \(t_c=t_0/2\); its matched \(r_c\) lies in a
fixed positive interval. On the two outer pieces beyond
\(a-t_c/2\) and \(-b+r_{\min}/2\), its fast height has a fixed
positive margin. The intervening piece has \(P\ge p_c\).
Choose a fixed upper central shot \(\eta_{\rm core}>0\); its limiting
graph is \(\eta_{\rm core}+\xi^2/2\). Both bounding shots have
buffered outward admission as above. Scalar order and \(Q>0\) admit
every intervening shot \(\eta_c\le\eta_0\le\eta_{\rm core}\),
including nonroots, and their raw central speed is at most
\(\epsilon^2(g+Q(0)\eta_{\rm core})<d_{\rm out}\).
For \(D=gP+Q\eta\) and \(f=\eta(H_0+H_1\eta)/D\),
use \eqref{deg:scalar-raw-cocycle}, with residual
\([gP'+(Q'+H_1)\eta]/D\).
On the middle piece the last fraction is bounded by
\(C+Cgk\); on the outer pieces the comparison height bounds
\(D\) below. The endpoint quotients are units, so integration bounds
the core first variation and its inverse. The two outward derivatives
are \(\epsilon^{2-\lambda_\gamma}\) times units on this entire
chord. Their ratio, regular bridges and closing give \(0<\Pi'\le C_c\).
The limiting \(t=0\) orbit stops, whereas every \(t>0\) passes;
Lemma~\ref{lem:loc-through} identifies \(x_{\rm st}\) with the
excluded mother endpoint. The endpoint extends through \(t_0\),
strictly beyond the comparison.

The comparison action gives \(\eta_c\le Cg^{-N}e^{-c/g}\). The
signed high threshold and the ordered final shrink are
\begin{equation}\label{deg:negative-gamma-cover}
 H_H=\epsilon^2gK_H(1+g\pi^2),\qquad
 \eta_c<\eta_H:=gK_H(1+g\pi^2)<\eta_{\rm core}/2,\qquad
 \epsilon^2\eta_{\rm core}<H_{\rm cap}/2.
\end{equation}
Hence the endpoint, closed core and signed high packages strictly overlap
and cover \(I_\theta\) on the original height axis. The equality
\(H=H_H\) belongs to the core; original clips are imposed on all
three. Choose \(m\) last so \(m\max(C_c,C_H)<1/2\).
It remains to control critical points in the endpoint interval,
including parameters for which the exponent ratio approaches one.
On the endpoint,
\(F'=1-e^L\), \(L=\log(mB)+(\nu-1)\log q+R(q)\).
If \(\nu\ge1\), then \(L<0\). Otherwise, for \(z=1-\nu>0\),
a critical point satisfies \(z\log(1/q)\ge M:=\log(1/m)-C\), so
\[
 qL_q=-z+O(q^\beta)\le-z+Ce^{-\beta M/z}<-z/2.
\]
This choice is uniform as \(z\downarrow0\). Since \(q_x=1/d_+>0\),
every actual critical point has \(F''>0\). The core and high add none;
Lemma~\ref{sharp:upper-critical-count} gives the whole-carrier count.
\end{proof}

\begin{proposition}[The positive far-root face]
\label{deg:prop-positive-far}
For \(\gamma=g>0\), there are independent bounds
\begin{equation}\label{deg:positive-far-range}
 0<\lambda=\pi^{-2}\le\lambda_0,\qquad
 0<q=k\pi\le q_0,\qquad 0<\tau=g\pi^2=p^2\le\tau_0
\end{equation}
such that the actual lower interval has strictly negative completed-map
curvature. Its bounded core and signed high complement strictly overlap
and have \(F'>1/2\) after one multiplier shrink. They cover the
original raw-capped \(I_\theta\), including all nonroot inputs and
original clips. The same physical \(F\) has at most one critical point
and at most two isolated zeros, counted with multiplicity, on every fibre.
\end{proposition}

\begin{proof}
The negative incident root now escapes in the scaled coordinate.
We estimate its outward map first below one third of the root radius and then
in a root collar rescaled by that radius. The exact corridor exponent
keeps both estimates uniform as the root escapes; their action weights
make the resulting section correction small in physical curvature.

Put \(\chi=q\lambda^2\), \(a=s/p=\tau\chi\),
\(g=\tau\lambda\), \(s=\tau^{3/2}\chi\),
\(v=g+aU\), \(z=a^2Z\); at the center \(Z=\tau Y_0\).
After positive time scaling the exact field is
\[
 U'=-\tau P_+(U)-Q_+(U)Z,\qquad
 Z'=-Z\{UB_+(U)+Ca(g+aU)Z\},
\]
\begin{equation}\label{deg:far-polynomials}
 \begin{gathered}
 P_+=1-U+3\lambda U^2+\chi U^3,\qquad
 P_-=1+T+3\lambda T^2-\chi T^3,\\
 B_+(U)=1+2g+aU,\quad B_-(T)=1+2g-aT,\\
 Q_+(U)=\ell(g+aU),\quad Q_-(T)=\ell(g-aT),\quad
 \ell(v)=1-Av+Cv^2.
 \end{gathered}
\end{equation}
The moving lips are \(Z=\tau\ell_0P_\pm\), with labels
\(b_+=-U^2/2,b_-=-T^2/2\). The source quartet is
\((-\tau,\tau,\tau\lambda,\tau\chi)\); reflection at a fixed
negative base \(T_b\) gives
\((-\tau/T_b^2,-\tau/T_b,\tau\lambda,-\tau\chi T_b)\)
without changing the positive-time orbit order.
The first positive root is \(A_+=1+O(\lambda+\chi)\).
Writing \(w_*=q\lambda T_*\), the unique negative-side root satisfies
\begin{equation}\label{deg:far-roots}
 w_*=3+q/w_*+q^2\lambda/w_*^2,\quad 3<w_*<4,\qquad
 T_*\asymp(q\lambda)^{-1},\quad \chi T_*\le4\lambda.
\end{equation}
Its slope is \(c_-=2+3\lambda T_*+3/T_*\). The action residues
\(r_+=A_+B_+(A_+)/[-P_+'(A_+)]\asymp1\) and
\(r_-=T_*B_-(T_*)/c_-\asymp\lambda^{-1}\) are retained exactly.
Define
\[
 \begin{gathered}
 J_+(U)=\int_0^U\frac{xB_+}{P_+}\,dx,\qquad
 A_-(T)=\int_0^T\frac{xB_-}{P_-}\,dx,\\
 \mathcal D_\pm=\frac{P_\pm}{xB_\pm}\partial_x,\quad
 H_\pm=\frac{P_\pm'}{xB_\pm},\quad h_\pm=\frac{P_\pm}{x^2B_\pm}.
 \end{gathered}
\]
With \(R=T^{-1},d=\chi T,E=T\partial_T\),
\(H_-=(R+6\lambda-3d)/B_-\),
\(h_-=(R^2+R+3\lambda-d)/B_-\), \(\mathcal D_-=h_-E\).
Repeated differentiation gives
\begin{equation}\label{deg:far-action-bound}
 \frac{|\mathcal D_-^jH_-|}{h_-}
 \le C_j(R+\lambda)^{j-1}(R+d),\qquad
 \int_{T_b}^{T_*}|\mathcal D_-^jH_-|\,dA_-
 \le C_j(T_b^{-1}+\chi T_*),\quad j=1,2,3.
\end{equation}
Here \(dA_-=d\log T/h_-\); the retained factor \(R+d\) avoids
any \(\lambda\log T_*\) loss. Also \(\tau H_-\le1/4\).
The reflected coefficient \(\varpi_-=-C\tau a(g-aT)P_-/(TB_-)\)
and its first four radial rows are \(O(\tau^3\lambda^3)\), while
those of \(Q_--1\) are \(O(g)\). Choose \(\ell_0\) with
\(\ell_0\max(|Q_-|,|\varpi_-|)\le1/4\). The positive root
collar has the same oriented margins on its fixed bounded cuts.

Apply Lemma~\ref{deg:lem-stable-ratio-ledger} to
\(\vartheta=Z/(\tau P)\), \(q_a=(A(x_l)-A(x))/\tau\), with
\(D=\mathcal D_l\) on either entrance, \(\kappa=\tau\), \(\omega=\varpi\),
\(\rho=\tau\), and \(\alpha=0\). At fixed clock, endpoint
transport is \(\mathcal D\) at the running point; the coefficient
rows follow from \eqref{deg:far-action-bound} and the bounded positive
collar. The margins above give the two negative diagonals, and
\(\partial_{q_a}=-\tau\mathcal D\) supplies the clock rows.
Thus \(\vartheta\le\ell_0e^{-q_a/4}\), and the exact memory is
\[
 \vartheta=\ell_0\frac{P(x_l)}{P(x)}e^{-q_a+V},\qquad
 V=\int_0^{q_a}\frac{(Q-\varpi)\vartheta}{1+Q\vartheta}\,dq,
\]
with three bounded action rows. In particular the initial polynomial
factor is retained through the long negative flight.
Fix inner cuts \(U_i>0\), \(0<T_i<T_b\) on a positive core,
with \(A_-(T_b)-A_-(T_i)\ge c_i>0\), and put
\(J_i=J_+(U_i)\), \(A_i=A_-(T_i)\).
At either fixed inner cut the terminal clock is
\(q_i=(A(x_l)-A(x_i))/\tau\), with
\(\mathcal D_lq_i=1/\tau\) and
\(\mathcal D_l^2q_i=\mathcal D_l^3q_i=0\).
All three total rows in \eqref{deg:stable-ratio-terminal-rows}
are bounded: their terminal densities decay as \(e^{-c/\tau}\),
absorbing the factors through \(\tau^{-3}\).
Writing \(V_\pm\) for the complete terminal integrals, the actual
raw weak heights are
\[
 \begin{aligned}
 w_+(U)&=\tau\ell_0P_+(U)e^{-[J_+(U)-J_i]/\tau+V_+(U)},\\
 w_-(T)&=\tau\ell_0P_-(T)e^{-[A_-(T)-A_i]/\tau+V_-(T)}.
 \end{aligned}
\]
The same five-row scalar core construction, with \(g\) replaced by
\(\tau\), gives
\(M(w)=\mu_\theta w(1+wN_\theta(w))\) on its actual weak chord,
with \(\mu_\theta^{\pm1}+\|N_\theta\|_{C^3}\le e^{K/\tau}\).
Its exact first multiplier has
\(\log\mu_\theta=(A_i-J_i)/\tau\).
Fix the cuts, large radial cut and normal height first; enlarge \(K\)
to include the fixed action-bracket reserves, then choose
\(X_0>0\) with \(J_+(A_+-X_0)-J_i>4K+1\), and
\(0<X_1<X_0\). The actual weak memory and three action rows are
\(O(\tau^{-N}e^{-(3K+1)/\tau})\) on \(A_+-X_0\le U<A_+\).
The exact equality \(M(w_+(U))=w_-(T)\) now gives
\begin{equation}\label{deg:far-action}
 J_+(U)-A_-(T)=\tau\{\log P_+(U)-\log P_-(T)+R_0(U,T)\},
 \quad |R_0|+\sum_{1\le i+j\le3}|\mathcal D_+^i\mathcal D_-^jR_0|\le C.
\end{equation}
Here the remainder has the separated extension
\begin{equation}\label{deg:far-separated-remainder}
 R_0(U,T)=V_+(U)-V_-(T)
       +\log\{1+w_+(U)N_\theta(w_+(U))\}.
\end{equation}
It is defined on the product of the actual positive weak chord and
negative inward domain. Its mixed rows with \(i,j>0\) are exactly zero;
the pure rows through order three are those just estimated.
Only then is it restricted to the actual center-image intersection.
Scalar order gives its unique matched
\(T_b\le T<T_*\); at both comparison labels \(A_+-X_j\),
fixed action brackets and \(A_-(T)\to T-\log(1+T)\) put \(T_j\)
in fixed positive compact intervals. Differentiation yields
\[
 T_U=\frac{UB_+P_-}{TB_-P_+}(1+\varrho),\qquad
 1+\varrho=\frac{1-\tau H_+-\tau\mathcal D_+R_0}
                   {1-\tau H_-+\tau\mathcal D_-R_0},\qquad
 |\mathcal D_+^i\mathcal D_-^j\varrho|\le C\tau\quad(i+j\le2).
\]

Let \(\Gamma_\pm(b_\pm)=d_\pm+c_\pm\phi_\pm(b_\pm)\) be the
actual raw embeddings, including the regular bridges and output closing;
the positive center is measured in \(x=h-a_\theta\).
With \(S(b_+)=-T^2/2\),
\(\widehat\Pi=\phi_-\circ S\circ\phi_+^{-1}\) and
\(\Pi(d_++c_+w)=d_-+c_-\widehat\Pi(w)\).
Put \(L_\pm=\mathcal D_\pm\log\Gamma_\pm'\), where the primes
refer to their own \(b_\pm\). Apply \eqref{deg:weighted-exact-curvature-chain} with
\(u=U,v=T,W=(U+T)/(UT)\), retaining both section terms.
In this weight the polynomial part is
\(3\chi UT-1\), apart from \(O(\tau\lambda)\).
The relative correction is \(O(\tau)\); the section term has bounded
weight \(UB_+T/(U+T)\). We now estimate its actual log-slopes.

At the positive root, pole subtraction and the intrinsic normal clock
retain the embedding power \((A_+-U)^{1+\alpha_+}\),
\(\alpha_+=\tau/r_+\). The normal distance kernel and its first
three weighted rows have majorant \(C(W+\tau u)\), whose normal-clock
integral is \(O(\sigma+\tau(1+|\log\tau|))\).
Its fixed normal-to-raw ordinary \(C^3\) observer therefore gives
\[
 |L_+|\le C_0\sigma+C_\sigma X_0+
                         C_\sigma\tau(1+|\log\tau|)^3.
\]
Choose \(\sigma\), then \(X_0\), to make the fixed terms small;
reducing \(X_0\) strengthens the earlier action gap.

For the negative outward graph put
\(D=\tau P_-+Q_-Z\), \(n=RB_-\),
\(\zeta=-Ca(g-aR)\), \(f=Z(n+\zeta Z)/D\).
It starts at the true moving value \(Z_l(T)=\tau\ell_0P_-(T)\)
and ends at the fixed raw port \(R_{\rm out}=\delta_-/a\).
The original observer \(\psi\) is composed only after raw admission.
Use \eqref{deg:scalar-polynomial-lip-slope} with
\(M=\tau P_-\), \(x_l=T\), \(\sigma=1\), and
\(\chi(Z)=\psi(a^2Z)\). It gives
\[
 \Gamma_-'(b_-)=a^2\psi'(a^2Z_T)\frac{Z_T}{D_T}B_-(T)
 [1+\ell_0\varpi_--\tau H_-(1+\ell_0Q_-)]e^{I(T)},
\]
where \(I=\int_T^{R_{\rm out}}[\tau P_-'+(Q_-'+\zeta)Z]/D\,dR\).
This retains the exact moving initial factor and first gain.
The two logarithmic-slope rows use only the ordinary third observer
derivative, by \eqref{phys:eq-observer-log}.

For \(T_b\le T\le T_*/3\), apply
Lemma~\ref{deg:lem-spectral-exterior} in the raw coordinate \(x=aR\),
with starting abscissa \(aT\), terminal cut \(\delta_-\), and
\[
 t=-1,\quad (\mu_2,\mu_1,\mu_0)=(3g,\tau a,\tau a^2),\quad
 b=2g,\quad v=g,\quad \sigma=a.
\]
Its input label is exactly \(-T^2/2\). With the comparison mass
\(\kappa=\tau/T^2+\tau/T+g+aT\), the root bound \(aT_*<4g\)
gives \(m(x)\asymp\kappa\) on \([aT,2aT]\), and all coefficient
bounds in \eqref{deg:spectral-exterior-hypotheses} hold.
The same margin holds for \(T<7T_*/20\), providing a strict overlap
with the root chart beginning at \(T_*/3\).
Here \(\Theta\le\kappa\le C\tau\), the terminal coefficient is
bounded, and \((aT/\delta_-)^{2-\nu}\le C\tau\), where
\(\nu=12g/(1+2g)\) is the lemma's exact spectral exponent.
Thus, for \(E=T\partial_T\), its two rows give
\[
 |E\log\Gamma_-'|+|E^2\log\Gamma_-'|
                       \le C\tau(1+|\log\tau|)^3.
\]
The observer is used on its original admitted image; neither \(a\)
nor the comparison mass is differentiated as a coefficient.

For the rescaled root range below, the corridor from \(R_c=LT_*\) retains
\(y=Z/R^2\), \(Q_g=\ell(g)\),
\(y_g=(1-4g)/(2Q_g)\), \(\nu=12g/(1+2g)\) exactly.
For \(m_R=\tau/R^2+\tau/R+3g-aR\),
\(B=1+2g-aR\), \(Q=\ell(g-aR)\),
\(\zeta_R=-CaR(g-aR)\), \(D_y=m_R+Qy\), the exact equations are
\[
 R y_R=y\left(\frac{B+\zeta_Ry}{D_y}-2\right),\qquad
 R(y-y_g)_R=-\frac{2Q_gy}{D_y}(y-y_g)+S,
\]
\[
 S=\frac y{D_y}\{(B-1-2g)-2(m_R-3g)+
                                      [\zeta_R-2(Q-Q_g)]y\}.
\]
In a fixed band about \(y_g\), the damping is \(\ge k_0>0\),
\(|S|\le C\xi(R)\), \(\xi=aR+\tau/R+\tau/R^2\).
Choose the raw radius first, then \(L\) and the coefficient bounds.
The actual shots enter the inner band, and the two boundary signs give
its strict continuation to \(\delta_-/a\). Variation of constants yields
\[
 \int_{R_c}^{R_{\rm out}}|y-y_g|\frac{dR}{R}
 \le C\left(|y(R_c)-y_g|+\delta_-+\frac\tau{R_c}
                                      +\frac\tau{2R_c^2}\right).
\]
The exact first density is
\(Rf_Z=(Bm_R+2\zeta_Rm_Ry+\zeta_RQy^2)/D_y^2\).
At \((3g,1+2g,Q_g,0,y_g)\) it equals \(\nu\), and its derivative
in \(y\) is \(O(g)\). Hence
\[
 |Rf_Z-\nu|\le C(g|y-y_g|+\xi),\qquad
 j(R)=(R/R_c)^\nu U(R),\qquad c\le U,U^{-1}\le C.
\]
Here \(j\) is the linear propagator. Extracting its exact exponent makes the residual density integrable
uniformly along the long corridor, without a \(g^2\log q\) loss.
Rational division also gives
\[
 |f_{ZZ}|\le C(aR^{-2}+gR^{-3}+\tau R^{-4}+\tau R^{-5}),\quad
 |f_{ZZZ}|\le C(aR^{-4}+gR^{-5}+\tau R^{-6}+\tau R^{-7}).
\]
For \(1/3\le\theta=T/T_*<1\), use the coefficient-only radius
\(T_*\), \(r=R/T_*\), \(W=Z/T_*^2\). The exact field is
\[
 r'=g\widehat P(r)+Q(r)W,\quad W'=Wn(r,W),\quad
 \widehat P=(1-r)[a_0+(a_0+a_1)r+w_*r^2],
\]
where \(a_0=q^2\lambda/w_*^2\), \(a_1=q/w_*\),
\(Q=\ell(g-gw_*r)\),
\(n=r(1+2g-gw_*r)-Cg^2w_*(1-w_*r)W\).
On the actual compact strip \(1/3\le r\le1+C\sigma\), \(n\) is a
positive unit. The selected strong graph \(r=\phi(W)\) at \((1,0)\)
has coefficient denominators
\(jB_*+g\widehat c\ge cj\),
\(B_*=1+2g-gw_*\), \(\widehat c=w_*+a_1+2a_0\).
Its power-series majorant gives a fixed normal interval and uniform
\(C^4\) bounds. Put \(x_-=1-\theta\),
\(W_0=g\ell_0x_-\widehat C(\theta)\), and \(u=\phi(W)-r\).
Then \(u_0=x_-(1+O(g))\), with three weighted rows, and exactly
\[
 \begin{gathered}
 (\log u)_{\log W}=-\alpha_--K(W,u),\quad
 \alpha_-=g\widehat c/B_* =\tau/r_-,\\
 |(W\partial_W)^i(u\partial_u)^jK|\le C(W+gu),\quad i+j\le3.
 \end{gathered}
\]
Taylor division of the rational normal field also gives
\(\alpha_-+K\ge c(g+W)\), so the actual orbit reaches \(W=\sigma\)
once. Its integrated densities are
\(\int W\,d\log W\le\sigma\),
\(\int gu\,d\log W\le Cg(1+|\log g|)\).
Use \eqref{deg:scalar-moving-bell} with \(s=\log W\),
\(z=\log u\), \(f=-\alpha_--K(W,e^z)\), and
\(D=E_*:=\theta(1-\theta)\partial_\theta\) at fixed \(s\).
Its initial recursion retains the moving clock
\(s_0=\log W_0\). The two integrated densities above give
\[
 u_\sigma=g^{\alpha_-}x_-^{1+\alpha_-}e^{A(\theta)},\qquad
 |A|+\sum_{j=1}^3|E_*^jA|\le C[\sigma+g(1+|\log g|)],
\]
and \(r_\sigma'=g^{\alpha_-}x_-^{\alpha_-}e^A
(1+\alpha_--x_-A_\theta)>0\).
The fixed normal continuation to \(r=L\) has height
\(H_g(r_0)=\sigma+(L^2-r_0^2)/2+O_{C^3}(g)\), with
\(-H_g'\) a unit. At the coefficient-only cut \(R_c=LT_*\),
its full initial Bell rows are \(T_*^2H_g^{(j)}\).
The same corridor and rational densities give
\(\int|f_{ZZ}U_1|\le Cg\) and
\(\int(|f_{ZZ}U_2|+|f_{ZZZ}U_1^2|)\le Cg\).
Thus the two Bell ratios are bounded, including the third observer row.
For \(\Psi_*(r_0)=\psi(a^2Z_{\rm out}(r_0))\), the exact gain is
\[
 -\Psi_*'=a^2\psi'(a^2Z_{\rm out})T_*^2[-H_g']
                      (\delta_-/(aLT_*))^\nu U(R_{\rm out}),\qquad
 \Omega_*\asymp(aT_*)^{2-\nu}\delta_-^\nu L^{-\nu}.
\]
After removing this exact gain, \(G=(\Psi_*-d_*)/\Omega_*\) and its
inverse have ordinary \(C^3\) bounds, with \(-G'\) a unit. Hence
\[
 \Gamma_-'(b_-)=\frac{\Omega_*}{T_*^2}
 \frac{-G'(r_\sigma)}\theta g^{\alpha_-}x_-^{\alpha_-}e^A
                         (1+\alpha_--x_-A_\theta),\qquad
 |E_*\log\Gamma_-'|+|E_*^2\log\Gamma_-'|\le C.
\]
Neither \(\Omega_*\) nor \(x_-^{\alpha_-}\) has been replaced by a
unit. Factor \(P_-=(1-T/T_*)C_*(T)\), where
\(C_*=1+(1+T_*^{-1})T+\chi T_*T^2\).
On this range \(\mathcal D_-=[C_*/(T^2B_-)]E_*\), with
\(C_*/(T^2B_-)\le C\lambda\). Combining both ranges proves the actual
consumed estimate
\begin{equation}\label{deg:far-negative-budget}
 |L_-|\le C\{\tau(1+|\log\tau|)^3+\lambda\},
                       \qquad T_b\le T<T_*.
\end{equation}
In the root range, the residual logarithmic derivatives are bounded,
and the factor \(C_*/(T^2B_-)\) makes their action derivatives small.
Together with the positive-section and relative-variation estimates,
this gives
\begin{equation}\label{deg:far-bell}
 |\varrho|+\max_{i+j\le2}|\mathcal D_+^i\mathcal D_-^j\varrho|
 +|L_-|\le C\{\tau(1+|\log\tau|)^3+\lambda\},\qquad
 |L_+|\le C_0\sigma+C_\sigma[X_0+\tau(1+|\log\tau|)^3].
\end{equation}
Therefore the full physical identity is
\begin{equation}\label{deg:far-curvature}
 \frac{UTP_+\phi_+'}{U+T}\frac{\widehat\Pi''}{\widehat\Pi'}
 =3\chi UT-1+E_f,\quad
 |E_f|\le\epsilon_0+C_\sigma[\tau(1+|\log\tau|)^3+\lambda],
 \quad\epsilon_0<1/8.
\end{equation}
Independent \(q_0,\lambda_0,\tau_0\) make this strictly negative.
The exact endpoint powers also remain: \(x_-=T_*-T\) obeys
\(x_-\sim C_\theta(A_+-U)^{(r_++\tau)/(r_-+\tau)}\), while
the two physical embeddings carry powers \(1+\tau/r_\pm\).
The physical exponent is \(r_+/r_-\asymp3\lambda\); only weighted
endpoint rows and pointwise \(C^3\) on the open passage are used.

It remains to establish raw admission and the whole upper complement.
The two comparison shots \(U=A_+-X_j\) have positive-unit heights at
fixed cuts \(L_+>A_++1\), \(L_->\max T_j+1\), since their fast
invariants give \(Z_+(L_+)=\tfrac12[L_+^2-(A_+-X_j)^2]+O(\tau(1+|\log\tau|))\)
and \(Z_-(L_-)=\tfrac12(L_-^2-T_j^2)+O(\tau(1+|\log\tau|))\).
Choose a fixed upper fast shot and \(\eta_{\max}>0\); scalar order
bounds every intervening central shot through that height at both cuts.
Before applying an observer, compare their ratios to the passive
parabolic references
\[
 t(Y_{{\rm par},\pm})_t=Y_{{\rm par},\pm}
 \left\{\frac{1\pm t+Ct^2Y_{{\rm par},\pm}}
 {\pm t+\ell(\pm t)Y_{{\rm par},\pm}}-2\right\},\qquad
 Y_{{\rm par},\pm}(0)=1/2.
\]
Their analytic recursion has denominator \(j+2\), and the original
prepared passive compact set retains their strict raw-window margins.
The actual/reference ratio fields have negative \(y\)-derivative and
difference \(O(g+s^2/t^2+sp/t)\). Logarithmic variation of constants
at each fixed raw radius yields
\[
 |y(\delta_\pm)-Y_{{\rm par},\pm}(\delta_\pm)|
 \le C[(aR_c/\delta_\pm)^{c_0}+g+\tau/R_c+\tau/R_c^2].
\]
Thus all comparison and upper core shots are admitted uniformly in the
independent rectangle. For the entire lower output family, the raw-height
spectral slope \eqref{deg:spectral-exterior-slope}, with identity observer, gives
\(|h_T|\le Ca^{2-\nu}\delta_-^\nu T^{1-\nu}\), hence total variation
\(\le C(aT_*)^{2-\nu}\delta_-^\nu\). The raw version of the root-range
gain gives the same bound. Since \(aT_*\le4g\), the total is
\(O(g^{3/2})\), smaller than the fixed reference buffer. The bounded
positive normal family has gain \(O(a^{2-\nu})\) and the same input
admission. Ordinary observers are composed only on these admitted images;
the original smaller windows, speed and stop clips are then intersected.

On the fixed core, \(D=\tau P+QZ\) and
\(f_Z=d\log(Z/D)/dU+[\tau P'+(Q'+\zeta)Z]/D\).
The inner positive piece has \(P\ge p_c\), the outer positive piece
has the comparison's \(Z\ge z_c\), and the negative fixed slab has
\(P_-\ge p_c\). The quotient is uniformly integrable and endpoint
\(Z/D\) are units, giving bounded core variation and inverse.
The positive tail has \(m_+=\tau/R^2-\tau/R+3g+aR\),
\(B_+=1+2g+aR\), \(Q_+=\ell(g+aR)\),
\(\zeta_{R,+}=CaR(g+aR)\); the same exact \(y_g,\nu\) corridor
applies. Both absolute outward gains are
\(a^{2-\nu}\delta_\pm^\nu\) times units. Their common factor cancels
in the actual \(\Pi'\), so \(0<\Pi'\le C_c\) through
\(\eta_{\max}\). All ordinary observer and divided closing rows are
those of Lemma~\ref{lem:closing}; their third chain rule uses only
\(f''',d'''\), retaining each small gain.

The action gap gives
\(0<\eta_{c,1}<\eta_{c,0}\le C\tau^{-N}e^{-c/\tau}\).
Fixing the shots and caps first, the independent shrink gives
\begin{equation}\label{deg:far-physical-cover}
 \begin{gathered}
 \eta_{c,1}<\eta_{c,0}<\eta_H:=\tau K_H(1+\tau)<\eta_{\max}/2,\\
 H_H=a^2\eta_H=K_H(s^2+(sp)^2),\qquad a^2\eta_{\max}<H_{\rm cap}/2.
 \end{gathered}
\end{equation}
On the original height axis, lower through the \(X_0\) shot, core
starting at the \(X_1\) shot, and high above \(H_H\) cover \(I_\theta\)
with strict overlaps. The positive equality \(H=H_H\) belongs to the
core; scalar order retains all clipped nonroot inputs. Choose \(m\)
last so \(m\max(C_c,C_H)<1/2\). The upper two packages have
\(F'>1/2\); on the lower one \(\Pi''<0\), so every critical point
of the same \(F\) has \(F''>0\). Lemma~\ref{sharp:upper-critical-count}
proves the whole-fibre bound, excluding its natural stopped endpoint.
\end{proof}

\subsection{Cubic and quadratic source compactifications}
\label{subsec:deg-cubic-quadratic}

The chart \eqref{deg:cubic-chart} covers the compact cubic limits of
the small-\(p\) region. Proposition~\ref{deg:prop-compact-cubic}
proves their complete return and whole-fibre count. The estimates below
supply its growing tails, where a compact phase comparison no longer
applies. We then return to the bounded positive quadratic face of
\eqref{deg:gamma-chart}.

On a growing positive tail, the lip label is unbounded while its raw
abscissa remains small. The next estimate controls the logarithmic
derivatives of the actual exterior map in that regime. Keeping its
absolute first gain will later allow composition with the independent
negative section map.

\begin{lemma}[Complete section derivatives on a growing cubic tail]
\label{deg:lem-raw-growing-section}
Fix a sufficiently small terminal cut \(\delta>0\). Let
\[
 \begin{gathered}
 M(x)=x^3+\mu_2x^2+\mu_1x+\mu_0,\qquad B_0(x)=1+b+x,\\
 Q_0(x)=1-A(v+x)+C(v+x)^2,\qquad \zeta_0(x)=C(v+x).
 \end{gathered}
\]
with bounded \(A,C\). Hold all coefficients fixed as the starting
point \(a\) varies. On an enlarged interval \(I\subset(0,d_0]\),
with \(d_0<\delta/2\), suppose uniformly for \(a\in I\)
\[
 |b|+|v|\le C_0a,\qquad |\mu_j|\le C_0a^{3-j}\ (j=0,1,2),
 \qquad M(x)\ge c_0x^3\quad(a\le x\le\delta).
\]
Use the actual positive graph
\begin{equation}
 z_x=\frac{z\{xB_0(x)+\zeta_0(x)z\}}{D_0(x,z)},\qquad
 D_0=M+Q_0z,\qquad z(a)=\ell_0M(a),
 \label{deg:weighted-raw-positive-graph}
\end{equation}
and an increasing regular terminal adapter \(\chi\), with uniformly bounded
\(C^3\) norm and inverse on its enlarged image. For fixed positive
\(\sigma\), label the input by \(b_+=-a^2/(2\sigma^2)\), and write
\(\Gamma_+(b_+)=\chi(z_\delta)\). After choosing \(\ell_0,d_0\)
small, this continuation and every separate affine normalization
\(\phi_+\) satisfy
\begin{equation}
 |(a\partial_a)^j\log\phi_+'|\le Ca(1+|\log a|)^4,
 \qquad j=1,2.
 \label{deg:raw-growing-section-jets}
\end{equation}
The exact positive factor \(\sigma^2\) in the first derivative is
retained; no bound on its reciprocal is required.
\end{lemma}
\begin{proof}
Set \(E=a\partial_a\) and \(m(x)=M(x)/x^2\). On \(a\le x\le2a\),
the coefficient bounds and the polynomial margin give
\[
 c_0a\le m(x)\le C_1a,\qquad
 |b|+|v|+|\mu_2|+a+|\mu_1|/a+|\mu_0|/a^2\le C_1a.
\]
Thus Lemma~\ref{deg:lem-spectral-exterior} applies with \(t=1\)
and comparison mass \(\rho=a\), including its near-lip strip entry.
This mass is not differentiated;
all polynomial coefficients and \(\sigma\) remain fixed. The adapter
is evaluated on its stated enlargement of the generated terminal image.
In the actual raw-section applications,
\eqref{deg:spectral-exterior-window}, with \(\rho=a\), and the
original fixed terminal-image margin admit the whole family after
choosing \(d_0\) within its original cutoff reserve. The stated
ordinary \(C^3\) adapter and inverse bounds give a uniform bound for
its two logarithmic-height rows on the enlarged terminal image.

The common spectral estimate has
\[
 \lambda=\frac{4\mu_2}{1+b}=O(a),\qquad
 \Theta=a+|\mu_1|/a+|\mu_0|/a^2\le C_1a.
\]
Its terminal coefficient \(\mathcal T\) is bounded on the stated
buffer, and \(\mathcal T(a/\delta)^{2-\lambda}\le Ca\), after
the permitted \(d_0\) shrink. Consequently
\[
 |E\log\Gamma_+'+\lambda|+|E^2\log\Gamma_+'|
 \le Ca(1+|\log a|).
\]
The physical gain is given exactly by
\eqref{deg:scalar-polynomial-lip-slope},
\begin{equation}\label{deg:weighted-complete-positive-derivative}
 \Gamma_+'=\sigma^2\chi'(z_\delta)\frac{z_\delta}{D_0(\delta,z_\delta)}
                 B_0(a)J_0(a)e^{I(a)},\qquad J_0=1+O(a),
\end{equation}
where \(J_0=K\) and \(I\) are the exact factors in that formula.
This retains the factor \(\sigma^2\) without a reciprocal bound.
Since \(\lambda=O(a)\), both uncorrected logarithmic rows satisfy
the stronger first-logarithmic-power bound. Frozen affine calibration
leaves them unchanged, proving \eqref{deg:raw-growing-section-jets}.
\end{proof}

\begin{lemma}[Weighted cubic tail on the actual sections]
\label{deg:lem-weighted-cubic-tail}
Fix a compact set of coefficients \(0\le U\le U_0\),
\(|G|\le G_0\), and use the field
\eqref{deg:compact-cubic-field}--\eqref{deg:compact-cubic-polynomials}.
On its root-free positive fibres choose \(M\) sufficiently large,
independently of \(\tauc\), and retain
\[
 u\ge2M,\qquad a=\tauc u\le d_0,\qquad r_0\le r<c.
\]
Here \(c\) is the simple negative root and the two moving sections are
\(Y=\ell_0P\), with \(\eta=\tauc Y\).  The exterior maps are the
actual continuations to the fixed raw sections, followed by the regular
input adapter and the divided finite-side output adapter, respectively.
These continuations use the nested domains constructed by
\eqref{deg:compact-raw-corridor}.
There are constants, uniform on these domains, such that, with
\[
 E=u\partial_u,\quad
 \mathcal D_+=\frac{P(u)}{uB(u)}\partial_u,\quad
 \mathcal D_-=\frac{P(-r)}{rB(-r)}\partial_r,\quad
 \epsilon_u=\tauc(1+u),\quad L_u=1+|\log\epsilon_u|,
\]
the exact action matching has the form
\begin{equation}
 A_+(u)-A_-(r)=\tauc\{\log P(u)-\log P(-r)+R_+(u)-R_-(r)\},
 \label{deg:weighted-split-action}
\end{equation}
where
\begin{equation}
 |R_+|+|R_-|\le C,\qquad
 |E^jR_+|\le C\epsilon_u,\qquad
 |\mathcal D_-^jR_-|\le C\tauc\quad(1\le j\le3).
 \label{deg:weighted-action-jets}
\end{equation}
The relative variation in \eqref{deg:weighted-relative-variation} obeys
\begin{equation}
 |\varrho|+\frac{|\mathcal D_+\varrho|+|\mathcal D_-\varrho|}{u+r}
       \le C\epsilon_uL_u^4.
 \label{deg:compact-rho}
\end{equation}
If \(b_+=-u^2/2\),
\(b_-=-r^2/2\), and \(\phi_\pm\) are the separately affine-normalized
actual section embeddings, then
\begin{equation}
 |E\log\phi_+'|+|E^2\log\phi_+'|
       \le C\epsilon_uL_u^4,
 \qquad
 |\mathcal D_-^j\log\phi_-'|
       \le C\tauc(1+|\log\tauc|)^4\quad(j=1,2).
 \label{deg:weighted-actual-section-jets}
\end{equation}
Every prime on \(\phi_\pm\) is taken in its own label \(b_\pm\).
The full composition \(\Pi=\phi_-\circ S\circ\phi_+^{-1}\), with
\(S(-u^2/2)=-r^2/2\), therefore satisfies
\begin{equation}
 \phi_+'\frac{P(u)}{u+r}\frac{\Pi''}{\Pi'}
       =3-\frac{U}{ur}+E_u,
 \qquad |E_u|\le C\epsilon_uL_u^4.
 \label{deg:compact-curvature}
\end{equation}
The inward action estimates and the two exterior estimates are obtained
separately on their actual domains.
\end{lemma}

\begin{proof}
We first continue the two inward graphs to the original centre and
differentiate their exact height match. We then estimate the two
exterior embeddings separately. The negative embedding requires a
root-collar calculation because its first variation vanishes there.
Substitution in the completed-curvature identity gives the final estimate.

Choose a small fixed \(\delta_c>0\) such that \(P\) is a positive unit
on \([-2\delta_c,2\delta_c]\), and so that
\(2\delta_c<r_0\).  On the root-free positive fibre, put
\[
 \begin{gathered}
 q=\frac{A_+(u)-A_+(v)}\tauc,\quad
 W=\frac{Y}{P(v)},\quad D(v)=\frac{P(v)}{vB(v)},\\
 H(v)=\frac{P'(v)}{vB(v)},\qquad
 \varpi(v)=\frac{C\tauc^3(G+v)P(v)}{vB(v)}.
 \end{gathered}
\]
For \(\delta_c\le v\le u\), graph division gives the exact equations
\begin{equation}
 v_q=-\tauc D(v),\qquad
 W_q=W\left\{-\frac{1+\varpi W}{1+QW}+\tauc H\right\},
 \qquad W(0)=\ell_0.
 \label{deg:weighted-inward-equations}
\end{equation}
In particular, \(\tauc H=O(\epsilon_u)\),
\(Q-1=O(\epsilon_u)\), and
\(\varpi=O(\epsilon_u^3)\).  Fix \(\ell_0\), then \(d_0\),
small enough that the braces and the derivative of the right-hand side
with respect to \(W\), for \(0\le W\le\ell_0\), are both less than
\(-c_1<0\).  A first-exit argument gives
\begin{equation}
 0<W(q)\le\ell_0e^{-c_1q}.
 \label{deg:weighted-kernel-decay}
\end{equation}

For unbounded \(u\), Euler differentiation keeps the moving entrance
under control. At fixed action clock \(q\),
\begin{equation}
 Ev=\kappa(u)D(v),\qquad
 \kappa(u)=\frac{u^2B(u)}{P(u)}\asymp u^{-1}.
 \label{deg:weighted-transport}
\end{equation}
The first three Euler derivatives of \(\kappa\) are \(O(u^{-1})\).
On \([M,u]\), \(D(v)=O(v^2)\), and each further application of
\(D(v)\partial_v\) increases polynomial degree by at most one.
On \([\delta_c,M]\), this vector field and its first four derivatives
are uniformly bounded, including when the positive factor approaches a
double root: they contain \(P\), rather than \(P^{-1}\).
Differentiating \eqref{deg:weighted-transport} inductively therefore gives
\[
 |E^j\log v|\le C_j\quad(1\le j\le4),\qquad
 |E^j(\tauc H)|+|E^j(Q-1)|+|E^j\varpi|
       \le C_j\epsilon_u\quad(0\le j\le3).
\]
Because the transported field contains \(P\), these bounds remain
uniform when the root-free positive minimum tends to zero.
Apply Lemma~\ref{deg:lem-stable-ratio-ledger} with
\(D=E,\kappa=\tauc,\rho=\epsilon_u,\alpha=0\) on
\(\delta_c\le v\le u\).
The needed clock rows follow from \(\partial_q=-\tauc D(v)\partial_v\):
on the growing part its degree growth is controlled by
\(\tauc v\le\epsilon_u\), and on the compact part its coefficients
are bounded. The strict diagonal chosen above gives
\begin{equation}\label{deg:weighted-kernel-jets}
 |E^jW(q)|\le C\epsilon_u(1+q)^Ne^{-c_2q},\qquad1\le j\le3.
\end{equation}

The nonlinear logarithmic correction on this tail is exactly
\[
 V_+(u)=\int_0^{q_c(u)}
       \frac{(Q-\varpi)W}{1+QW}\,dq,
 \qquad q_c(u)=\frac{A_+(u)-A_+(\delta_c)}\tauc.
\]
The fixed-clock and terminal formulas of the same lemma give \(|V_+|\le C\) and
\(|E^jV_+|\le C\epsilon_u\), \(1\le j\le3\): indeed,
\(|E^jq_c|\le C_j/\epsilon_u\) and
\(q_c\ge[A_+(2M)-A_+(M)]/\tauc\ge a_M/\tauc\).
Every boundary term is therefore at most
\(C\tauc^{-N}e^{-c_2a_M/\tauc}\).

For the negative inward tail replace \((P,B,Q,\varpi)\) by
\((P(-v),B(-v),Q(-v),-C\tauc^3(G-v)P(-v)/(vB(-v)))\).
At fixed \(q\), \(\mathcal D_-\) transports to
\(P(-v)[vB(-v)]^{-1}\partial_v\);
its coefficient rows in the braces are \(O(\tauc)\) on
\([\delta_c,c]\), including the root factor. The same lemma with
\(\rho=\tauc\) gives \(|V_-|\le C\),
\(|\mathcal D_-^jV_-|\le C\tauc\). Since \(r\ge r_0>2\delta_c\),
its complete terminal rows are exponentially small.

Continue both tails to \(\xi=0\) using
Lemma~\ref{deg:lem-compact-height-continuation}, with radial coordinate
\(v=|\xi|\) oriented from \(\delta_c\) to zero and
\(\kappa_c=\tauc\). Here \(P\ge p_c>0\), the linear inward
propagator is at most one, and the nonlinear coefficient has ordinary
height derivatives through order three bounded by \(C/\tauc\).
For the incoming heights \(y_\pm=Y(\pm\delta_c)\), the exact memory
\eqref{deg:stable-ratio-polynomial-memory} gives
\(y_\pm\le C\tauc^{-N}e^{-c/\tauc}\); its three
logarithmic phase rows grow at most polynomially in \(\tauc^{-1}\).
The cited continuation and composition rows therefore make the remaining
nonlinear correction exponentially small, with those phase derivatives.
The exact cocycle \eqref{deg:exact-action-cocycle} gives the actual
heights on the common section:
\begin{align*}
 \log Y_0^+(u)&=\log(\ell_0P(u))-A_+(u)/\tauc+R_+(u),\\
 \log Y_0^-(r)&=\log(\ell_0P(-r))-A_-(r)/\tauc+R_-(r).
\end{align*}
with \eqref{deg:weighted-action-jets}.  Equality of these heights is
exactly the physical matching condition.  This proves
\eqref{deg:weighted-split-action}; it also supplies a separated extension
\(R_0(u,r)=R_+(u)-R_-(r)\), whose genuinely mixed derivatives vanish.

The separated height formulas also determine the relative first
variation and its weighted derivatives. Write
\(P_+=P(u)\), \(P_-=P(-r)\), \(B_+=B(u)\), \(B_-=B(-r)\),
and use radial derivatives in
\(H_+=P_+'/(uB_+)\), \(H_-=P_-'/(rB_-)\).
Differentiating \eqref{deg:weighted-split-action} gives
\begin{equation}
 r_u=\frac{uB_+P_-}{rB_-P_+}(1+\varrho),\qquad
 1+\varrho=
 \frac{1-\tauc H_+-\tauc\mathcal D_+R_+}
      {1-\tauc H_--\tauc\mathcal D_-R_-}.
 \label{deg:weighted-relative-variation}
\end{equation}
Since \(\mathcal D_+=(P_+/(u^2B_+))E\), the numerator differs
from one by \(O(\epsilon_u)\), with the same Euler bound.  The
denominator differs from one by \(O(\tauc)\), with two bounded action
variations of that order.  This proves \eqref{deg:compact-rho} (even
with \(L_u^4\) omitted).  Both denominators are positive after shrinking
\(d_0\) and then \(\tauc\).

For the positive exterior map apply
Lemma~\ref{deg:lem-raw-growing-section} with
\[
 \sigma=\tauc,\quad a=\tauc u,\quad
 (\mu_2,\mu_1,\mu_0)=(3\tauc G,-\tauc^2U,\tauc^3),
 \quad b=2\tauc G,\quad v=\tauc G.
\]
For \(u\ge2M\), the coefficient bounds hold and the leading cubic
term gives \(M(x)\ge c x^3\) on \([a,\delta]\), uniformly on the
fixed compact block. In the translated coordinate \(x=v-\gamma\),
the original positive raw port is \(x_T=\delta-\tauc G\).
The proof of the raw lemma is uniform for terminal cuts in a fixed
neighborhood of \(\delta\); \(x_T\) is coefficient-only and is held
fixed under \(E\). Thus it applies at this actual port. The corridor
supplies its enlarged terminal image.
Since \(a\asymp\epsilon_u\), \eqref{deg:raw-growing-section-jets}
is the positive assertion of \eqref{deg:weighted-actual-section-jets}.

Put \(\lambda_G=12\tauc G/(1+2\tauc G)\).
For the negative exterior map, first restrict to
\(r_0\le r\le c-x_0\), with \(x_0>0\) fixed and small.
Here \(P_-(r)=P(-r)\) has a fixed positive margin. Put \(v=-\xi\),
\(B_-(v)=B(-v)\), \(Q_-(v)=Q(-v)\), and use the increasing
height clock \(s=\log\eta\). The fast invariant
\(H=\eta-v^2/2\) satisfies the exact equation
\[
 H_s=\tauc\mathscr F(H,e^s),\qquad
 \mathscr F=
 \frac{v\eta(B_--Q_-)/\tauc
       -C\tauc(G-v)\eta^2-vP_-(v)}
      {vB_-(v)-C\tauc^2(G-v)\eta}.
\]
On a fixed box with \(v\ge r_0/2\), the denominator is a positive
unit, \((B_--Q_-)/\tauc\) is smooth and bounded, and
\(v=\sqrt{2(e^s-H)}\) gives
\( |\partial_H^i\partial_s^j\mathscr F(H,e^s)|\le C\)
for \(i+j\le3\). The initial rows are
\[
 s_l=\log\{\tauc\ell_0P_-(r)\},\qquad
 H_l=-r^2/2+\tauc\ell_0P_-(r).
\]
The ordinary \(r\)-derivatives of \(s_l-\log\tauc\) are bounded.
The moving initial recursion of
\eqref{deg:scalar-moving-bell} therefore gives the three initial
variations of \(-r^2/2\), with errors \(O(\tauc)\).
Choose a fixed small \(\eta_*>0\). First exit and the three
triangular equations on
\(s_l\le s\le\log\eta_*\), whose length is
\(O(1+|\log\tauc|)\), give
\[
 H(\log\eta_*;r)=-r^2/2+
          O_{C^3_r}\bigl(\tauc(1+|\log\tauc|)\bigr).
\]
This is an actual outward passage: initially
\(D=\tauc P_-+Q_-\eta>0\), and at a possible first zero
\(D_s=Q_-\eta>0\). On the fixed positive-height continuation
to a sufficiently large fixed cut \(v=L>c+1\), the graph equation
is an \(O_{C^3}(\tauc)\) perturbation of \(\eta_v=v\). Hence
\[
 \eta_L(r)=\tfrac12(L^2-r^2)+
          O_{C^3_r}\bigl(\tauc(1+|\log\tauc|)\bigr).
\]
The cut and its buffered chord lie in the original corridor
\eqref{deg:compact-raw-corridor}. From this cut to the raw port use
the negative tail with source quartet
\((a,p,\gamma,s)=(-\tauc,\tauc U,\tauc G,\tauc)\).
For a fixed \(0<\nu<1\), \eqref{eq:N-bound} gives
\[
 \int_L^{\delta/\tauc+G}\|\mathcal N\|_{C^3}\,dv
 \le C\tauc\int_L^\infty
 (v^{-2+\nu}+v^{-3+\nu}+v^{-4+\nu}+v^{-5+\nu})\,dv
 \le C\tauc.
\]
After extraction of its exact first slope, the tail and its inverse
are \(\mathrm{id}+O_{C^3}(\tauc)\). Its physical slope is
\(d_{\rm c}=\tauc^{2-\lambda_G}U_{\rm c}\), with
\(U_{\rm c}\) a positive unit. The original ordinary \(C^3\)
observer and divided closing add \(O_{C^3}(d_{\rm c})\) by
\eqref{eq:exterior-normalization}. Calibrating value and slope at a
frozen reference in this compact interval gives
\(\widetilde\phi_-=\mathrm{id}+O_{C^3}(\tauc(1+|\log\tauc|))\)
in \(b_-=-r^2/2\). Since \(\mathcal D_-\) has bounded coefficients
and derivatives here, the two logarithmic slope rows have the same
bound. They are unchanged by a frozen positive affine calibration,
so the original \(\phi_-\) satisfies
\eqref{deg:weighted-actual-section-jets} on this interval.
Near the root, the first variation vanishes and the preceding compact
normalization no longer gives a bounded inverse. We extract its endpoint
power before estimating logarithmic derivatives.

The root equation gives the exact factorization
\[
 P_-(r)=(c-r)L_-(r),\qquad
 L_-(r)=r^2+(U/c+c^{-2})r+c^{-1},\qquad
 L_-(c)=c^2+U+2/c>0.
\]
Thus \(c\) and \(L_-\), with their required finite coefficient
derivatives, are uniform units on a fixed negative collar,
independently of the positive bottleneck. Put \(x=c-r\), \(y=\eta\).
Dividing by the
positive normal rate \(rB(-r)\) gives
\[
 \dot x=-\beta x p(x)-k_1(x)y,\qquad
 \dot y=y\{1+k_2(x)y\},\qquad
 \beta=\frac{\tauc L_-(c)}{cB(-c)}=O(\tauc),
\]
Here the coefficients are exactly
\[
 \begin{split}
 p(x)&=\frac{cB(-c)L_-(c-x)}{L_-(c)(c-x)B(-c+x)},\\
 k_1(x)&=\frac{Q(-c+x)}{(c-x)B(-c+x)},\qquad
 k_2(x)=-\frac{C\tauc^2(G-c+x)}{(c-x)B(-c+x)}.
 \end{split}
\]
In particular \(p(0)=1\), \(p,k_1\) are positive units and
\(k_2=O_{C^7}(\tauc^2)\) on a fixed box.
Use the normal-time construction
\eqref{eq:saddle-normal-time}--\eqref{eq:exact-saddle-invariant} in the
proof of Lemma~\ref{lem:actual-dulac-matching}, with
\(\rho=\beta\), \(h=-\beta xp\), \(c=-k_1\), \(k=k_2\),
and selected coefficient parameter \(\tauc\) at fixed \(U,G,A,C\).
The explicit coefficients have a bounded joint order-seven budget:
\(h=O_{C_x^7}(\beta)\), \(k_1\) is a positive unit, and their
selected \(\tauc\)-rows are bounded. The normal coefficient and its
parameter rows retain their height factors.
The projection at phase order five uses axis-flow order six and the
selected parameter row. For a fixed normal decay rate \(0<\lambda<1\),
its integrands have the majorant
\(C(1+T)^Ne^{-(\lambda-C\beta)T}\), integrable after
\(C\beta<\lambda/2\). Thus the exact coordinates \((\Xi,Z)\)
and invariant \(I=\Xi Z^\beta\) have five phase derivatives and one
intrinsic \(\tauc\)-derivative, as required for the divisions below.
In the following coefficient-dependent units, a subscript
\({\rm c}\) records dependence on \(\tauc\). On the moving curve
\(y=\tauc\ell_0P_-\), these coordinates have the form
\[
 \Xi_m=xu_{\rm c}(x),\qquad Z_m=\tauc xv_{\rm c}(x),\qquad
 I(x)=\theta_{\rm c} x^{1+\beta}A_{\rm c}(x),\quad
 \theta_{\rm c}=\tauc^\beta,\quad A_{\rm c}=u_{\rm c} v_{\rm c}^\beta,
\]
with positive \(C^4\) units \(u_{\rm c},v_{\rm c},A_{\rm c}\).
To check the coefficient differences before taking the root limit, put
\(Y_{\rm c}(x)=\tauc x\ell_0L_-(c-x)\). Then
\[
 \partial_{\tauc}u_{\rm c}(x)=
 \frac{(\partial_{\tauc}\Xi)(x,Y_{\rm c}(x))}{x}
 +\Xi_y(x,Y_{\rm c}(x))\ell_0L_-(c-x),\qquad
 v_{\rm c}=\ell_0L_-(c-x)e^{-J(x,Y_{\rm c}(x))}.
\]
The first numerator vanishes at \(x=0\); its parameter-valued fifth
phase row gives \(C^4\) after Taylor division. In the second term the
moving factor \(x\) has already cancelled, so \(\Xi_y\) needs only
four phase rows. The formula for \(v_{\rm c}\) removes the factor
\(\tauc x\) before differentiation. Integrating these intrinsic
parameter rows gives \(u_{\rm c}-u_0,v_{\rm c}-v_0=O_{C^4}(\tauc)\),
and hence \(A_{\rm c}-A_0=O_{C^4}(\tauc)\). Also
\(\theta_{\rm c}=1+O(\tauc|\log\tauc|)\).

At a fixed outer cut beyond \(-c\), the stopped trajectory has
positive height. The invariant is a regular increasing coordinate
there, including \(I=0\), as in
\eqref{eq:exact-saddle-invariant}. Choose its buffered invariant chord
first, then the smaller lip collar so \(I(x)\) lies strictly inside it.
Choose the upper horizontal boundary with invariant values above that
chord throughout the local strip. This is possible at \(\tauc=0\),
where \(\pi=c-\sqrt{r^2-2\eta}\), and the strict margin persists
for small \(\tauc\). Thus no orbit of the lip collar exits through
that boundary.
The outward speed starts positive. At a possible first zero of
\(\mathscr D=\tauc P_-+Q_-\eta\), one has
\(\dot r=0\) and \(\dot{\mathscr D}=Q_-\dot\eta>0\), excluding
loss of that speed. Increasing height and the invariant chord then give
the actual first hit of the outer cut at positive height, uniformly for
\(0<x\le x_0\). Continue this fixed chord to the
negative raw section, before either the regular output adapter or the
finite closing map. After extracting the intrinsic exact first slope, write
\[
 T(x)=t_{\rm st}+d_{\rm c} F_{\rm c}(I(x)),\qquad
 d_{\rm c}=\tauc^{2-\lambda_G}U_{\rm c}>0,
\]
where \(U_{\rm c}\) is a positive unit and \(F_{\rm c},F_{\rm c}^{-1}\)
have uniform \(C^4\) bounds on a fixed invariant chord. Here the
finite positive-height passage has \(O(\tauc)\) coefficient changes.
On its subsequent tail, \eqref{eq:N-bound}, integrated with the
frozen negative quartet, bounds the nonlinear normalized kernel by
\(O(\tauc)\) through four phase orders. Thus the intrinsic,
separately normalized maps
satisfy \(F_{\rm c}=F_0+O_{C^3}(\tauc)\); no convergence rate for
the extracted slope \(d_{\rm c}\) is used.

Since \(b_-=-r^2/2\), one has \(db_-/dx=r=c-x>0\).
Consequently the exact first derivative is
\[
 T_{b_-}'=d_{\rm c}\theta_{\rm c} x^\beta
                 K_{\rm c}(x,x^\beta),\qquad
 K_{\rm c}(x,z)=
 \frac{F_{\rm c}'(\theta_{\rm c} xzA_{\rm c}(x))
       \{(1+\beta)A_{\rm c}(x)+xA_{\rm c}'(x)\}}{c-x}.
\]
The function \(K_{\rm c}\) is a positive \(C^3\) unit on a fixed
rectangle \(0\le x\le x_0\), \(0\le z\le1\), and
\(K_{\rm c}(0,z)\) is independent of \(z\).
At \(\tauc=0\) the outer graph is
\(\eta(\xi)=(\xi^2-r^2)/2\). Its intrinsic, separately
normalized section map is affine in \(b_-\). The regular factors
therefore cancel exactly at \(z=1\), giving a positive constant
\(k_0\), independent of \(x\), for which
\[
 \|K_{\rm c}(\cdot,1)-k_0\|_{C^2}
       \le C\tauc(1+|\log\tauc|).
\]
The affine limiting map is what makes the logarithmic derivative small;
boundedness of the regular factors alone would not suffice.

To differentiate along \(z=x^\beta\), use
\(E_x=x\partial_x+\beta z\partial_z\). For
\(0<x\le x_0<1\),
\[
 x|1-x^\beta|\le C\beta.
\]
The uniform \(C^3\) bounds for \(K_{\rm c}\), its constancy in
\(z\) at \(x=0\), and the preceding cancellation imply
\[
 |E_x\log K_{\rm c}|+|E_x^2\log K_{\rm c}|
       \le C\tauc(1+|\log\tauc|)
       \quad\hbox{on }z=x^\beta.
\]
Indeed the differences of the \(x\)-derivatives between \(z\)
and \(1\) have factors \(1-z\), accompanied by \(x\) or
\(x^2\); all remaining terms contain \(\beta\).
Since
\[
 \mathcal D_-=-\frac{L_-(c-x)}{(c-x)B(-c+x)}E_x,
 \qquad E_x\log x^\beta=\beta,
\]
we obtain the first two action-logarithmic bounds for
\(T_{b_-}'\), including their one-sided limits as \(x\downarrow0\).
The derivative itself may vanish like \(x^\beta\).

Finally append the original increasing ordinary \(C^3\) output adapter
\(\psi\) and the divided closing map. Put
\[
 \Psi_{\theta,m}=\mathfrak f_{\theta,m}\circ\psi,
 \qquad \mathfrak f_{\theta,m}(q)=f_\theta(mq)/m,
 \qquad \Gamma_-=\Psi_{\theta,m}\circ T.
\]
Admission to the original raw and closing windows precedes these
compositions. Apply \eqref{phys:eq-observer-log} and
\eqref{phys:eq-closing-jets} with \(D=\mathcal D_-\) and primes in
\(b_-\). They bound the two logarithmic-slope corrections using only
the ordinary \(C^3\) norms and positive first derivative of
\(\Psi_{\theta,m}\), uniformly for \(0\le m\le m_0\).
The invariant formula retains the small gain:
\(|DT|+|D^2T|\le Cd_{\rm c}\), and
\(d_{\rm c}\le C\tauc^2\), since \(\lambda_G=O(\tauc)\).
Both additional logarithmic rows are therefore \(O(d_{\rm c})\),
using only the ordinary \(C^3\) adapter and closing bounds of
Lemma~\ref{phys:lem-c3-adapters}. Thus the actual complete output
satisfies \eqref{deg:weighted-actual-section-jets}; affine normalization
fixes its logarithmic rows.

We can now combine the inward match with both exterior estimates.
Apply Lemma~\ref{neg:lem-completed-curvature} with
\(v=r\), \(P_+=P(u)\), \(P_-=P(-r)\), and \(W=u+r\).
In \eqref{deg:weighted-exact-curvature-chain}, put \(R_B=B_+/B_-\).
The first term differs from
\((P'(u)/u+P'(-r)/r)/(u+r)=3-U/(ur)\) by
\(O(\epsilon_u)\).  The second is \(O(\epsilon_u)\), since
\(R_B-1=\tauc(u+r)/B_-\) and
\(r_u=O(P_-/(u^2r))\).  The third is controlled by
\eqref{deg:compact-rho}, and the last by
\eqref{deg:weighted-actual-section-jets}, since
\(P_+/(u^2(u+r))\) is bounded.  For the negative section term use
\[
 \frac{\tauc(1+|\log\tauc|)^4}{u+r}
       \le C\epsilon_u(1+|\log\epsilon_u|)^4.
\]
To see this uniformly, write \(\epsilon_u=\tauc(1+u)\), use
\(1+|\log\tauc|\le1+|\log\epsilon_u|+\log(1+u)\), and note that
\((1+\log(1+u))^4/(1+u)^2\) is bounded.
This proves \eqref{deg:compact-curvature} in the actual affine
input coordinate.
\end{proof}


\begin{proposition}[Positive quadratic face]
\label{deg:prop-positive-quadratic}
Fix \(\Pi<\infty\). For \(\gamma=g>0\),
\(0<k<k_0\), and \(0<\pi\le\Pi\), the actual low-height through
interval on the real/fold side and its fixed adjacent root-free band
has the moving-section descriptions below. Every zero of the complete
map's curvature on this interval is simple and is a downward crossing
in the same increasing physical input. The bounded-height complement
has uniformly bounded complete derivative. On the root-free side
outside this band, the complete return is a contraction after a
coefficient shrink.
The domains and estimates persist through the actual discriminant
\(\Delta=0\), whose limit at \(k=0\) is \(\pi^2=12\).
\end{proposition}

\begin{proof}
We use the curvature principal part \(3-\pi/(uR)\), with the
homogeneous section powers retained. Its zeros require separate
arguments in the joint growing corner and near two stopped endpoints.
The latter calculation uses fixed normal sections when \(b/a\) is
near two. The remaining growing ranges have a strict curvature sign,
and a paired derivative bound supplies the upper complement.

Fix a preliminary \(0<\pi_{\rm pre}<\sqrt{12}\).
All comparison estimates are uniform for
\(\pi_{\rm pre}\le\pi\le\Pi\). The final
\(\pi_\flat\in(\pi_{\rm pre},\sqrt{12})\) will be chosen after
the comparison margin. For \(0<\pi\le\pi_\flat\),
Lemma~\ref{deg:lem-raw-jacobian} gives contraction after the final
coefficient shrink.

In \(x=gk\xi\) the exact polynomial and system are
\begin{equation}
 \begin{split}
 P(\xi)&=1-\pi\xi+3\xi^2+k\xi^3
       =(\mathfrak c+k\xi)\{(\xi-d)^2-\Delta\},\\
 \dot\xi&=-g\{P+Q_+Y\},\qquad
 \dot Y=-Y\{\xi B_++Ckg^3(1+k\xi)Y\},
 \end{split}
 \label{deg:quadratic-field}
\end{equation}
where
\[
 \mathfrak c=3+2kd,\qquad
 \mathfrak c(d^2-\Delta)=1,\qquad
 \pi=2\mathfrak cd-k/\mathfrak c,\qquad
 B_+=1+g(2+k\xi).
\]
Also \(Q_+=1-Ag(1+k\xi)+Cg^2(1+k\xi)^2\), and
the physical coordinates are \(x=gk\xi\), \(z=g^3k^2Y\).
At \(k=0\), \(d=\pi/6\) and
\(\Delta=(\pi^2-12)/36\). The curves \(Y=\ell_0P\) are transverse
by the analogue of \eqref{deg:compact-lip-speed}, and the reciprocal
raw equation \eqref{deg:compact-raw-corridor} joins them to the fixed
sections. On a provisional compact coefficient block about the fold,
and the real block up to \(\pi=\Pi\),
\(a,a^{-1},d,d^{-1},\mathfrak c,\mathfrak c^{-1}\) are bounded
where the corresponding real root is defined.
The original cocycles of Corollaries~\ref{deg:cor-positive-compact-suppliers}
and \ref{deg:cor-negative-inward-common-centre} give a finite
zeroth corrected-action bound \(K_0\). Fix \(r_c>0\) and a finite
action threshold greater than
\(K_0+\sup\widetilde A_-(r_c)+1\).
For a fixed small \(\delta_c>0\), the real factorization gives
\[
 P(v)\le C(a-v),\qquad vB_+(v)\ge c>0,
 \quad a-\delta_c\le v\le a-X_{\rm up}.
\]
Thus \(J_+(a-X_{\rm up})\ge c\log(\delta_c/X_{\rm up})\),
uniformly as \(b-a\to0\), and the fold has stronger
inverse-distance divergence. Choose \(X_{\rm up}=X>0\) to pay
this threshold. Choose the retained ghost width \(\Delta_0>0\)
sufficiently small afterwards so that the same margin holds on
\(-\Delta_0\le\Delta<0\). Then choose \(\pi_\flat\) with
\(0<12-\pi_\flat^2<18\Delta_0\).
Retain \(|\Delta|\le\Delta_0\) and the compact real block up to
\(\pi=\Pi\). Outside this band on the root-free side,
\(\pi<\pi_\flat\) after shrinking \(k\), so it is already a
contraction. In the remainder of the proof,
\begin{equation}
 \pi_\flat\le\pi\le\Pi.
 \label{deg:quadratic-pi-lower}
\end{equation}
Put \(u_{\rm th}=a-X\), \(a=d-\sqrt\Delta\), on the real/fold
side and \(u_{\rm th}=d-X\) on the ghost side.
The comparison \(u_*=a-X/2\), respectively \(d-X/2\), uses
this same \(X\) in Lemma~\ref{deg:lem-positive-quadratic-paired-variation}.
The independent degenerations are \(g\to0\), \(k\to0\), and the
positive root gap.

Fix a translated raw half-width \(0<d_q<\delta/2\), small enough for the
parabolic corridor estimates below, and shrink the coefficient neighborhood
so that \(|\gamma|<d_q/2\).  The cuts \(x=\pm d_q\), equivalently
\(v=\gamma\pm d_q\), lie on the corresponding sides of \(v=\gamma\) and
strictly within the strip bounded by the fixed sections
\(\Sigma_\pm=\{v=\pm\delta\}\).  The uniformly regular segments from these
cuts to \(\Sigma_\pm\), and their inverses, are included in the actual
section adapters used below.

For real or double roots the positive lip interval ends at \(a\).
For \(\Delta<0\), choose \(0<d_{\rm in}<d_0<d_q/2\) and restrict to
the component containing \(u_{\rm th}\) of the preimage of the
enlarged input and output section intervals, with
\[
 u\ge u_{\rm th},\qquad x=gku<d_0.
\]
Take smaller nested intervals with \(x<d_{\rm in}\).
Indeed \(z=\ell_0(x^3+3gx^2-\pi kg^2x+k^2g^3)\) on the lip;
for \(d_{\rm in}\le x\le d_q\) this is separated from the
organizing arc when \(\ell_0d_q\) is small.  The localization lemma
and the common transverse sign therefore place every actual low-height
cycle in the smaller interval.  On the positive continuation use the
relative outer cut \(x=L gku\), and on the negative continuation
\(x=-LgR\), with \(L>1\) fixed large.  At these cuts
\(\mathcal W=1/2-1/(2L^2)+o(1)\); the corridor equation carries
both traces to \(x=\pm d_q\), even across an axis root at positive
height.  This specifies the actual interval used below.  Its endpoint
on a root-free fibre is a finite section boundary, not \(u=\infty\).

Put \(R=kr\),
\[
 \widehat P(R)=k^2P(-R/k)
  =(\mathfrak c-R)\{(R+kd)^2-k^2\Delta\}.
\]
Here \(B_-(R)=1+g(2-R)\).  Define the two actions on fixed inner
cuts by \(J_+'=uB_+/P(u)\) and
\(A_-'=rB_- /P(-r)\), and put
\(\mathcal D_+=P(u)(uB_+)^{-1}\partial_u\).
In either negative coordinate,
\(\mathcal D_-= \widehat P(R)(RB_-)^{-1}\partial_R
=P(-r)(rB_-)^{-1}\partial_r\).
Retaining both logarithmic terms in the moving-lip matching gives
\[
 \widetilde J_+=J_+-g\log P,\qquad
 \widetilde A_-=A_--g\log P(-r),\qquad
 |\widetilde J_+-\widetilde A_-|\le K_0 .
\]
Here \(K_0\) is independent of \(g,k\), and
\(\widetilde A_-(r)\ge c\log(1+r)-C\).
The comparison action was chosen above
\(K_0+\sup\widetilde A_-(r_c)\).
Every retained low-height input then has \(r\ge r_c\), while the
comparison output itself is bounded above independently of \(k\).
We define \(\varrho\) by the exact relative variation
\[
 R_u=\frac{uB_+\widehat P(R)}{RB_-P(u)}(1+\varrho)>0 .
\]
On the growing negative tail let
\(\mathcal D_-=\widehat P(R)(RB_-)^{-1}\partial_R\).  Its logarithmic
action derivative is
\begin{equation}
 H_-={2\mathfrak c-3R\over1+g(2-R)}+O(r^{-1}+k),\qquad
 |\mathcal D_-^jH_-|\le C(R+r^{-1}),\quad j\ge1.
 \label{deg:quadratic-negative-action}
\end{equation}
In a fixed simple-root collar at \(R=\mathfrak c\), the rational
function \(H_-\) and its required derivatives are bounded, while
\(\mathcal D_-\) contains the factor \(\mathfrak c-R\).
Consequently \(\mathcal D_-^jH_-=O(\mathfrak c-R)\) for \(j\ge1\),
and \(dA_-\asymp dR/(\mathfrak c-R)\) there.  Together with the
growing-tail estimate this gives
\(\int|\mathcal D_-^jH_-|\,dA_-\le C\).  In the Bell densities
\eqref{deg:bell-density}, the factor \(P\) cancels with the action
change of variable.  The factor \(\mathfrak c-R\) is retained in the
exact logarithmic cocycle and the last segment is treated in the fixed
simple-root coordinate.  No power \(k^{-C/g}\) is generated.
More explicitly, on \(r_c\le r\le(\mathfrak c-\ell_-)/k\),
\(\mathcal D_-=V(kr,r^{-1})(kr\partial_{kr}
-r^{-1}\partial_{r^{-1}})\), with \(V\) a positive unit.
The integral of \(kr+r^{-1}\) against \(dA_-\asymp dr/r\)
is bounded on this interval.  The remaining fixed collar
\(\mathfrak c-2\ell_-<R<\mathfrak c\) has normal rate
\(RB_-\) bounded below and a simple factor \(\mathfrak c-R\).
The two estimates overlap at a fixed \(R\)-cut, where the
Leibniz boundary terms cancel.  No action integral crosses that root.

Only the root-free case can have \(u\to\infty\).  We first treat the
joint quadratic corner, where both radial action speeds are bounded.
Put \(q=u^{-1}\), \(S=ku\), and \(t=r^{-1}\).  Choose a small
\(\delta_q>0\), fix \(0<\delta<\delta_q/2\), and choose \(M\) large enough for
\(M\delta\gg\Pi\) and the action gap below.  All these constants
precede the choices of \(g,k\).  The homogeneous quadratic field has
first integral
\[
 Y^{-\lambda_g/2}\left(\kappa_gY-{\xi^2\over2}\right),
 \qquad \lambda_g={12g\over1+2g},\qquad
 \kappa_g=\frac{g(1-Ag+Cg^2)}{1-4g}.
\]
Its radial exponent is the same at the two fixed raw sections.  On
\(u,r\ge2M\), \(S,R\le\delta_q\), the exact ratio and variational
equations, after this power is extracted, give, with
\[
 h_+={P(u)\over u^2B_+(u)},\quad
 h_-={\widehat P(R)\over R^2B_-(R)},\quad
 \theta_+=u^{-1}+ku,\quad \theta_-=r^{-1}+kr,
\]
\begin{equation}
 \begin{split}
 \sum_{i+j\le2}|\mathcal D_+^i\mathcal D_-^j\varrho|
 &\le\varepsilon_g(\theta_++\theta_-),\\
 \mathcal D_\pm\log\phi_\pm'&=-\lambda_gh_\pm+G_\pm,\\
 |G_\pm|+|\mathcal D_\pm G_\pm|&\le\varepsilon_g\theta_\pm,
 \end{split}
 \qquad \varepsilon_g=Cg(1+|\log g|)^4.
 \label{deg:quadratic-weighted-estimates}
\end{equation}
The two \(-\lambda_gh_\pm\) terms come from the same homogeneous
power and must be combined before estimating the curvature.
These bounds hold in the joint corner; the range \(S\to\infty\),
where \(h_+\asymp1+S\), will instead use radial derivatives.
Lemma~\ref{deg:lem-joint-inward-product-direct} gives the first
line on the product of the two separate actual domains, including
the mixed second row; restriction to their original central-image
intersection gives the matched graph. Lemma~\ref{deg:lem-quadratic-outer-jets}
gives the section rows with the stronger factor \(Cg(1+|\log g|)\),
on each admitted raw image.
The negative observer includes the divided closing of
Lemma~\ref{lem:closing}; its ordinary \(C^3\) logarithmic-slope
bounds apply on those original buffered images, with its small gain
and \(m\) fixed under the phase derivative.

We shall also use absolute negative-side estimates on the larger range
\(r\ge r_c\), \(0<R<\mathfrak c\).  Write the exact logarithmic
matching as
\[
 J_+(u)-A_-(r)
 =g\{\log P(u)-\log P(-r)+\mathcal R_+(u)-\mathcal R_-(r)\}.
\]
Coefficient-only constants may be included in either \(\mathcal R_\pm\).
The two functions are obtained by continuing each actual lip to the same
central section.  Thus their separated extension and its derivatives
refer to the actual orbit.  The complete negative inward action remainder and the negative
outward embedding give
\begin{equation}
 |\mathcal D_-^j\mathcal R_-|\le Cg\quad(1\le j\le3),\qquad
 |\mathcal D_-^j\log\phi_-'|\le C\varepsilon_g\quad(j=1,2).
 \label{deg:quadratic-negative-absolute}
\end{equation}
Here and below primes on \(\phi_-\) are in the label \(-R^2/2\).
Changing from \(-r^2/2\) only changes its frozen affine slope.
The first estimate is Corollary~\ref{deg:cor-negative-inward-common-centre}
at the original centre, retaining the initial polynomial and every
coefficient-only action constant. In particular,
\[
 N_-=-g\mathcal D_-\log Z_-
 =1-gH_--g\mathcal D_-\mathcal R_-,\qquad
 |\mathcal D_-^j(N_--1+gH_-)|\le Cg^2,\quad0\le j\le2.
\]
Since \(H_-\) is bounded, \(N_-\ge1/2\). The second estimate is
Corollary~\ref{deg:cor-negative-absolute-outer}, with the stronger
\(Cg(1+|\log g|)\) bound, including original-buffer admission,
\(p\)-weighted initial variations, both moving endpoints and the
ordinary \(C^3\) divided closing. Neither estimate enlarges the
intersection of actual central images.

For \(\mathcal T(-u^2/2)=-R^2/2\) and
\(\Pi=\phi_-\circ\mathcal T\circ\phi_+^{-1}\), put
\(B_{\rm rat}=B_+/B_-\) and
\(\mathcal V=\mathcal D_++(1+\varrho)\mathcal D_-\).
Lemma~\ref{neg:lem-completed-curvature}, with
\(v=R\), \(P_-=\widehat P\), and \(W=R+S\), gives
\eqref{deg:weighted-exact-curvature-chain} for the weighted curvature
\begin{equation}
 \mathscr K=\phi_+'\frac{P(u)}{R+S}\frac{\Pi''}{\Pi'}.
 \label{deg:quadratic-full-chain}
\end{equation}
The derivative \(d/du\) follows the matched graph.  In the joint
corner, exact logarithmic differentiation and the preceding estimates
give
\begin{equation}
 \phi_+'{P(u)\over R+ku}{\Pi''\over\Pi'}
 =F_\lambda(u,R)+E(u,R),
 \label{deg:quadratic-curvature}
\end{equation}
where, with \(q=u^{-1}\) and \(t=r^{-1}\),
\begin{equation}
 F_\lambda=(3-\lambda_g)-{N_\lambda\over uR},\qquad
 N_\lambda=(1-\lambda_g)\pi-\lambda_g(t-q),
 \label{deg:quadratic-principal}
\end{equation}
and, in the joint corner \(u,r\ge2M\), \(S,R\le\delta_q\),
\begin{equation}
 |E|+|(\mathcal D_++(1+\varrho)\mathcal D_-)E|
 \le C\varepsilon_g\left(1+{1\over uR}\right).
 \label{deg:quadratic-error}
\end{equation}
To verify \eqref{deg:quadratic-principal}, use
\[
 {P'(u)/u-\widehat P'(R)/R\over R+ku}=3-{\pi\over uR}
\]
and
\[
 {\widehat P/R^2-P/u^2\over R+ku}
 ={ku-R\over u^2R^2}+{\pi\over uR}-1.
\]
The matching rows of Lemma~\ref{deg:lem-joint-inward-product-direct}
apply on the product domain to
\(N_\pm=1-gH_\pm-g\mathcal D_\pm\mathcal R_\pm\),
\(n=1+\varrho=N_+/N_-\), with \(N_-\ge1/2\).

To estimate the curvature error and its derivative along the match,
put \(W=R+S\), \(\mathcal V=\mathcal D_++n\mathcal D_-\), and
\(A_\pm=B_\pm(H_\pm-\lambda_gh_\pm)\). Direct polynomial calculation gives
\begin{align*}
 A_+&=(6-3\lambda_g)+(\lambda_g-1)\pi q
                 +(3-\lambda_g)S-\lambda_gq^2,\\
 A_-&=(6-3\lambda_g)+(1-\lambda_g)\pi t
                 -(3-\lambda_g)R-\lambda_gt^2.
\end{align*}
Since \(B_+-B_-=gW\), the exact error decomposition is
\begin{align*}
 E_B&=-gA_-/B_--(B_+/W)\mathcal V\log(B_+/B_-),\\
 E_n&=-(B_+/W)\{\varrho(H_--\lambda_gh_-)+\mathcal V\log n\},\\
 E_s&=(B_+/W)(G_+-nG_-),\qquad E=E_B+E_n+E_s.
\end{align*}
The first two action rows of \(h_\pm,B_\pm,H_\pm\) are bounded.
Thus \(|E_B|+|\mathcal VE_B|\le Cg\). The total-two matching row gives
the required bound on \(E_n\), including \(\mathcal V^2\log n\).
For the section term, retain the exact derivative
\[
 \mathcal V(G_+-nG_-)
 =\mathcal D_+G_+-(\mathcal Vn)G_--n^2\mathcal D_-G_-.
\]
Lemma~\ref{deg:lem-quadratic-outer-jets} bounds these terms with the
stronger factor \(Cg(1+|\log g|)\). Finally
\((q+S+t+R)/W=1+1/(uR)\), and \(|\mathcal V\log W|\le C\).
These identities give \eqref{deg:quadratic-error} on the actual
matched joint graph. The other coefficient ranges are treated below.

On the joint corner \(q\) and \(t\) are small.  The lower bound
\eqref{deg:quadratic-pi-lower} and a further common shrink give
\[
 N_\lambda=(1-\lambda_g)\pi-\lambda_g(t-q)
 \geq\frac14\pi_\flat .
\]
Put \(z=uR\) and \(a_0=3-\lambda_g\). At a zero of the actual
\(\mathscr K=F_\lambda+E\), write
\(|E|\le\delta_g(1+1/z)\), with \(\delta_g\to0\).
Then
\[
 \frac{N_\lambda-\delta_g}{a_0+\delta_g}
 \le z\le
 \frac{N_\lambda+\delta_g}{a_0-\delta_g}.
\]
All actual curvature zeros therefore lie in a fixed positive compact
\(z\)-interval. The exact derivative there is
\[
 \mathcal VF_\lambda
 =\frac{N_\lambda(h_++nh_-)
          +\lambda_gh_+q-\lambda_gnh_-t}{z}\ge c>0,
\]
after \(M\) is large and \(g\) small. The differentiated error is
absorbed on this compact interval. If
\(w=\phi_+(-u^2/2)\) is the actual affine input, then
\(\mathcal Vw=-P(u)\phi_+'/B_+<0\). Because \(\mathscr K\) is a positive
multiple of \(\Pi''/\Pi'\), every curvature zero on this matched joint
graph has \(\Pi'''(w)<0\). To obtain this orientation on the
entire low-height interval, we next treat the mixed and remaining
growing ranges, followed by the bounded positive range.
There is also a mixed range \(u\ge2M\), \(r_c\le r\le2M\).
If \(S\ge\delta\), the action comparison below already forces
\(R\ge R_0(\delta)>0\), which contradicts \(R\le2Mk\).
On \(S\le\delta\), the positive half of the joint-corner calculation
still applies: its action speed is bounded.  The negative end is now
compact in \(r\), and \eqref{deg:quadratic-negative-absolute} supplies
its absolute estimates.  In \eqref{deg:weighted-exact-curvature-chain} these give
\(\mathscr K=3-\pi/(uR)+E_M\), with
\(|E_M|\le C_M\varepsilon_g/(R+S)\).  Relative to
\(1/(uR)\), it is bounded by
\(C_M\varepsilon_g\,ur/(u+r)\le2MC_M\varepsilon_g\).
A zero would require \(uR\ge c_M>0\), and hence
\(ku=uR/r\ge c_M/(2M)\).  The action comparison below applies
already for \(r\ge r_c\); with this last fixed threshold in place
of \(\delta\), it forces \(R\ge R_0>0\), contrary to
\(R=kr\le2Mk\).  Thus the mixed range contains no zero.
The remaining corners introduce none.  If \(u\ge2M\),
\(R\ge\delta\), and \(S\le\delta\), the same positive estimates
and \eqref{deg:quadratic-negative-absolute}, including its simple-root
collar, give an error \(O_\delta(\varepsilon_g)\) in
\eqref{deg:weighted-exact-curvature-chain}.  The exact polynomial principal part
is \(3-\pi/(uR)\); the choice \(M\delta\gg\Pi\) leaves a
positive margin.  For the last range \(S\ge\delta\), fix
\(0<\overline R<\tfrac12\inf\mathfrak c\).
For \(r\ge r_c\) and \(R\le\overline R\), subtraction of the exact
logarithmic cocycles gives
\begin{equation}
 \widetilde J_+(u)\ge {1-4g\over3}\log{1\over k}-C_\delta,
 \qquad
 \widetilde A_-(r)\le {1-4g\over3}
   \left(\log{1\over k}+\log R\right)+C_\delta.
 \label{deg:quadratic-large-S-action}
\end{equation}
The upper bound for \(\widetilde A_-\) is used only on this interval,
which stays away from its logarithmic singularity at
\(R=\mathfrak c\).  Write \(c_g=(1-4g)/3\), and shrink \(g\) so
that \(c_g\ge c_*>0\).  Since
\(|\widetilde J_+-\widetilde A_-|\le K_0\), the displayed inequalities
give \(c_g\log R\ge-K_0-2C_\delta\).  Together with the alternative
\(R>\overline R\), this proves the uniform bound
\[
 R\ge R_0(\delta):=
 \min\{\overline R,\exp[-(K_0+2C_\delta)/c_*]\}>0.
\]
Consequently \(uR\ge\delta R_0(\delta)/k\to\infty\).
We prove the required estimate on this range in the radial derivative
\(E=u\partial_u\), keeping the conversion to action derivatives
explicit.  Set
\[
 a=gS=gku,\qquad L_a=1+|\log a|,\qquad
 h_+=\frac{3-\pi/u+u^{-2}+S}{1+2g+gS}.
\]
For \(S\ge\delta\), \(a<d_0\), and \(k\) small,
\begin{equation}
 h_+\asymp_\delta1+S,\qquad
 |E^jh_+|\le C_\delta h_+\quad(j=1,2),\qquad
 \frac{h_+}{R+S}\le C_\delta,
 \qquad \mathcal D_+=h_+E.
 \label{deg:quadratic-large-S-speed}
\end{equation}
In particular the first bound in
\eqref{deg:quadratic-weighted-estimates} is not used here.

Corollary~\ref{deg:cor-positive-large-S-inward}, with
\(k_0\le\delta/(2M)\), gives the actual same-centre Euler-three rows
\begin{equation}\label{deg:quadratic-large-S-action-jets}
 |E^j\mathcal R_+|\le C_\delta a,\qquad1\le j\le3.
\end{equation}
Its strict ghost domain retains the moving initial polynomial,
terminal terms and full compact multiplier; the real/fold through
interval is empty here. The action conversion remains \(\mathcal D_+=h_+E\).

Differentiating the separated action matching gives exactly
\[
 1+\varrho=
 \frac{1-gH_+-g\mathcal D_+\mathcal R_+}
      {1-gH_--g\mathcal D_-\mathcal R_-},\qquad
 H_+=\frac{P'(u)}{uB_+(u)}.
\]
Both denominators remain positive.  Since
\(g|E^jH_+|\le C_\delta a\), \(gh_+\le C_\delta a\),
\eqref{deg:quadratic-negative-absolute} and
\eqref{deg:quadratic-large-S-action-jets} imply
\begin{equation}
 |\varrho|+|E\varrho|\le C_\delta a,\qquad
 |\mathcal D_-\varrho|\le C_\delta g,\qquad
 |\mathcal V\log(1+\varrho)|
       \le C_\delta(h_+a+g).
 \label{deg:quadratic-large-S-relative}
\end{equation}
Thus an action derivative of \(\varrho\) need not be small before
division by \(R+S\).

Choose \(d_0<a_0\) from
Corollary~\ref{deg:cor-positive-large-S-outer}. Its actual positive
raw graph, with \(k_0\le\delta/(2M)\), original observer and frozen
\((gk)^2\) first gain, gives the stronger \(C_\delta aL_a\) bound;
hence
\begin{equation}\label{deg:quadratic-large-S-section-jets}
 |E^j\log\phi_+'|\le C_\delta aL_a^4,\qquad j=1,2.
\end{equation}

We now substitute in \eqref{deg:weighted-exact-curvature-chain}.  Its polynomial
term has principal part \(3-\pi/(uR)\).  Since
\(B_{\rm rat}-1=g(S+R)/B_-\) and
\(\widehat P_R/R\) is bounded for \(R\ge R_0\), the remaining
part of that term is \(O_\delta(a)\).  The second term is exactly
\[
 -\frac{gS h_+}{R+S}
 -\frac{gB_+\widehat P(1+\varrho)}{RB_-^2(R+S)},
\]
and is again \(O_\delta(a)\).  The third and the positive section
term are bounded by \eqref{deg:quadratic-large-S-speed},
\eqref{deg:quadratic-large-S-relative}, and
\eqref{deg:quadratic-large-S-section-jets}.  For the complete negative
section, including its root collar and finite closing, use
\eqref{deg:quadratic-negative-absolute} and
\[
 \frac{g(1+|\log g|)^4}{R+S}\le C_\delta aL_a^4.
\]
Indeed \(a=gS\), \(S\ge\delta\), and
\((1+\log S)^4/S^2\) is bounded for \(S\ge1\); the interval
\(\delta\le S\le1\) is compact.  We have proved
\begin{equation}
 \mathscr K=3-\frac{\pi}{uR}+E_S,\qquad
 |E_S|\le C_\delta a(1+|\log a|)^4.
 \label{deg:quadratic-large-S-curvature}
\end{equation}
First decrease \(d_0\) so that this error is small on \(a<d_0\),
then decrease \(g,k\).  The lower bound on \(uR\) gives strict
positivity on the entire large-\(S\) range.

It remains to justify the bounded positive range and its overlap with
the preceding estimates.  On \(u\le3M\), Corollary~\ref{deg:cor-positive-compact-suppliers},
prepared on \(u<4M\), supplies
\( |\mathcal D_+^j\mathcal R_+|\le C_Mg\), \(1\le j\le3\), and
\( |\mathcal D_+^j\log\phi_+'|\le C_Mg(1+|\log g|)\), \(j=1,2\).
Thus \(N_+=1-gH_+-g\mathcal D_+\mathcal R_+\ge1/2\), with
\(|\mathcal D_+^j(N_+-1+gH_+)|\le C_Mg^2\), \(j\le2\),
on its original real/fold interval or strict ghost preimage.

On bounded moving-lip abscissae the ratio
equation, the growing negative estimate
\eqref{deg:quadratic-negative-action}, and the full section maps give
\begin{equation}
 \mathscr K=3-\frac{\pi}{uR}+E_{\rm lip},\qquad
 |E_{\rm lip}|+|\mathcal D_+E_{\rm lip}|
                  +|\mathcal D_-E_{\rm lip}|
       \le\frac{C\varepsilon_g}{R+k},
 \quad
 \mathscr K=\phi_+'\frac{P(u)}{R+ku}\frac{\Pi''}{\Pi'}.
 \label{deg:quadratic-bounded-lip-curvature}
\end{equation}
These action rows consume the bounded positive and absolute
negative suppliers on their separate actual images, with both
moving initial factors, endpoint rows and physical gains retained.

On a real-root fibre factor
\[
 P(\xi)=(\mathfrak c+k\xi)(\xi-a)(\xi-b),\qquad
 0<a\le b,\quad
 \mathfrak c=3+k(a+b),\quad
 \mathfrak cab=1,\quad
 \pi=\mathfrak c(a+b)-kab.
\]
Here \(a,a^{-1},b,\mathfrak c,\mathfrak c^{-1}\) are bounded.
At the simultaneous stopped endpoint \((u,R)=(a,\mathfrak c)\),
the limiting principal value at \(g=k=0\) is
\[
 3-\frac{\pi}{a\mathfrak c}=2-\frac ba.
\]
The endpoint value vanishes at \(b/a=2\); this is distinct from
the double-root condition \(\pi=\sqrt{12}\) at \(k=0\).
Outside a fixed neighborhood of \(b/a=2\), this value has a
nonzero margin.  A possible curvature zero then has
\(R\ge c_F>0\) and \(P(u)+\widehat P(R)\ge c_0>0\).
Since
\[
 \frac{d}{du}
 =\frac{uB_+}{P(u)}
       \{\mathcal D_++(1+\varrho)\mathcal D_-\},
 \qquad
 \frac{uB_+}{P(u)}\le C(1+R_u)
\]
on this zero neighborhood, the error derivative is at most
\(C\varepsilon_g(1+R_u)\), whereas
\(d(3-\pi/(uR))/du\ge c(1+R_u)\).
For \(R\to0\), the principal term instead dominates the error
relatively.  This proves the crossing assertion except in a fixed
joint-root neighborhood of \(b/a=2\).  The double root \(a=b\)
belongs to the separated case.  The adjacent root-free block has the
same endpoint margin: the arctangent action gives its continuation,
and no passage beyond \(a=d\) is added on \(\Delta=0\).

Near \(b/a=2\), the endpoint principal value can vanish while both
action speeds tend to zero. We therefore need errors whose ordinary
derivatives are small in the two independent root distances.
The root gap \(b-a\) is a positive unit in this range.
Use the actual fixed normal sections constructed in
Lemma~\ref{deg:lem-joint-root-physical-resection}, in the positive
\((\xi,gY)\) and negative \((R,gk^2Y)\) root boxes.  Their strong footpoints
are \(u=a-h\) and \(R=\mathfrak c-t\), with \(h,t>0\).
The retained interval is the intersection of the actual matched graph
with \(h<h_0\), \(t<t_0\), where \(h_0,t_0\) are fixed first.
It is one terminal interval because \(t'(h)>0\).
The normal rates in these two boxes are positive units, independently
of \(g,k\).  In this fixed-root calculation, regard \(A_-\) as
a function of \(R\); the identity
\(rB_-\,dr/P(-r)=RB_-\,dR/\widehat P(R)\) preserves its value.
Lemma~\ref{deg:lem-joint-root-physical-resection} supplies the
pole subtractions with residues \(\mathfrak a_\pm\) and uniform
\(C^3\) regular parts. All coefficient-only action constants,
including possible \(\log k\), remain fixed. The normal-footpoint
quotient is \(\nu_{\rm fp}=\mathfrak a_+/\mathfrak a_-\), without
an added \(g\) in either residue.
The exact residue identity is
\begin{equation}
 3-{\pi\over a\mathfrak c}
 -{g\widehat P_R(\mathfrak c)\over
        \mathfrak cB_-(\mathfrak c)}
 ={b-a\over a}\left({\mathfrak a_+\over\mathfrak a_-}-1\right).
 \label{deg:quadratic-residue-identity}
\end{equation}
This identifies the endpoint value of the actual weighted curvature;
it does not assert that every nearby curvature zero lies on the
coefficient resonance itself.

On the actual matched intersection, the same lemma gives
\eqref{deg:joint-root-normal-weighted-matching},
\eqref{deg:joint-root-normal-matching-rows}, and
\eqref{deg:joint-root-normal-outer-rows}. Its coefficient memory
\(C_{g,k}\) and both homogeneous gains are retained without size
or reciprocal bounds; the separated matching errors vanish with
\(h,t\), and \(|\partial_h\varrho|+|\partial_t\varrho|\le C\varepsilon_g\).

Put \(E_*=-g\widehat P_R(\mathfrak c)/
(\mathfrak cB_-(\mathfrak c))\).
The axis correction in the exact curvature identity is
\[
 E^{\rm ax}=
 -\frac{gkP}{u(R+ku)B_+}
 -\frac{g\widehat P_R}{RB_-}
 -\frac{gB_+\widehat P}{R(R+ku)B_-^2}.
\]
It has endpoint value \(E_*\).
The remaining matching and section terms, before estimating, are
\[
 \begin{split}
 E^\varrho={}&
 -\frac{B_+\varrho(\widehat P_R-\widehat P B_{-,R}/B_-)}
                    {RB_-(R+ku)}
 -\frac{B_+\{\mathcal D_+\varrho+
                    (1+\varrho)\mathcal D_-\varrho\}}
                    {(R+ku)(1+\varrho)},\\
 E^{\rm sec}={}&
 \frac{B_+}{R+ku}
  \{\mathcal D_+\log\phi_+'-
                 (1+\varrho)\mathcal D_-\log\phi_-'\}.
 \end{split}
\]
Write \(E_0=E^{\rm ax}-E_*+E^\varrho+E^{\rm sec}\).
Here \(P\asymp h\), \(\widehat P\asymp t\), and all displayed
denominators are positive units with bounded ordinary endpoint derivatives.
The separated quotient calculation in
Lemma~\ref{deg:lem-joint-root-physical-resection} bounds the ordinary
\(h,t\) derivatives of \(\varrho\), \(\mathcal D_+\varrho\)
and \(\mathcal D_-\varrho\) by \(C\varepsilon_g\).
For the section terms, division of the second rows in
\eqref{deg:joint-root-normal-outer-rows} by
\(h\gamma_+\), respectively \(t\gamma_-\), gives
\[
 |\partial_h(\mathcal D_+\log\phi_+')|
 +|\partial_t(\mathcal D_-\log\phi_-')|
 \le C\varepsilon_g .
\]
The axis correction satisfies
\(E^{\rm ax}-E_*=O(g(h+t))\), with ordinary endpoint derivatives
\(O(g)\). Hence the product rule gives
\( |\partial_hE_0|+|\partial_tE_0|\le C\varepsilon_g\)
on the separated product. Restricting to the actual \(t=t(h)\) yields
\begin{equation}
 \mathscr K=3-\frac{\pi}{uR}+E_*+E_0,\qquad
 |E_0|\le C\varepsilon_g(h+t),\qquad
 \left|\frac{dE_0}{dh}\right|
       \le C\varepsilon_g(1+t'(h)),
 \label{deg:quadratic-normal-curvature}
\end{equation}
Since
\[
 \frac{d}{dh}\left(3-\frac{\pi}{uR}\right)
 =-\frac{\pi}{u^2R}
  -\frac{\pi t'(h)}{uR^2}
 \le-c(1+t'(h)),
\]
the actual curvature is strictly decreasing at every possible zero,
even when both action rates tend to zero independently.

To join this interval to the moving lip, use the actual nonlinear
resection of Lemma~\ref{deg:lem-joint-root-physical-resection}.
Equation~\eqref{deg:joint-root-exact-lip-resection} retains the added
\(g\)'s from the two initial polynomials. Its positive-unit residues
give uniform inclusions between suitably nested fixed endpoint
neighborhoods; they do not give a linear comparison of the internal
root distances. Equation~\eqref{deg:joint-root-physical-composition}
identifies the same completed physical map after the two exterior
embeddings. Only positive affine changes of its physical input and
output remain, so its curvature zeros agree. Both inputs increase
with the corresponding positive endpoint distance. Hence every zero
has the same downward orientation on the overlap.
The bounded construction is used on \(u\le3M\).  On \(u\ge2M\),
the mixed range covers \(r\le2M\); for \(r\ge2M\), the remaining
alternatives are \(S,R\le\delta\), \(R\ge\delta,S\le\delta\),
and \(S\ge\delta\).  They are covered respectively by the joint
corner, the absolute negative-tail estimate, and
\eqref{deg:quadratic-large-S-curvature}.  The fixed thresholds are used
with the enlarged overlaps specified above, so these constructions
exhaust the prescribed interval. Every curvature zero has this same
strict downward orientation in the common physical input.

Put \(u_*=a-X_{\rm up}/2\) on the real/fold side and
\(u_*=d-X_{\rm up}/2\) on the ghost side. In the same original
physical input, Lemma~\ref{deg:lem-positive-quadratic-paired-variation}
gives \(0<\Pi_\theta'\le C_*\) on the entire admitted interval
\([\eta_{\rm comp},\eta_{\max}]\), where
\(\eta_{\rm comp}=gZ_+(u_*)\). The bounded comparison cocycle
gives
\(\eta_{\rm comp}<gZ_+(u_{\rm th})<\eta_{\max}\) after a common
coefficient shrink. Hence this interval overlaps the lower lip
interval strictly and retains every intervening input in the
original buffers. The common choice \(mC_*<1/2\) gives
\(F'>1/2\) on its bounded-height complement, uniformly through
\(\Delta=0\).
On \(\Delta=0\), the formal branch
\(u>d\) approaches the equilibrium in infinite time and is not included.
\end{proof}

\subsection{The admitted raw tail on a bounded positive quadratic fibre}

The growing curvature calculation uses \(gku<d_0\). On a ghost
fibre, the original physical interval can also contain first lips beyond
this divider. We keep that tail and compare its complete input and output
gains through their common central height. The negative gain is controlled
by the sign of its density during the long root dwell.

\begin{proposition}[The original bounded-positive ghost raw tail]
\label{bp-tail:prop-actual-raw-tail}
Fix the passive compact set, original raw sections, nested windows,
ordinary \(C^3\) observers and finite closing of Lemma~\ref{lem:closing}.
Fix \(\pi\in[\pi_\flat,\Pi]\) and the retained strict ghost band
\(\Delta(k,\pi)<0\) of \eqref{deg:quadratic-field}. The parameters
\(g,k>0\) approach zero independently. Keep the original master
\[
 F(x)=x-m\Pi_{\theta,m}(x),\qquad x=h-a_\theta,
 \qquad \Pi_{\theta,m}=\mathcal C_{\theta,m}\circ T_\theta(a_\theta+\cdot),
\]
and one positive unscaled cap
\(0<h_c<\min(H_{\rm cap},z_*)\), chosen before the final coefficient
and multiplier bounds. Its original carrier is
\begin{equation}\label{bp-tail:original-carrier}
 I_\theta=\{x:a_\theta+x\in J_\theta^{\rm out},\quad
                         0<H(x)=z_\theta(g;a_\theta+x)<h_c\}.
\end{equation}
Choose the fixed auxiliary \(d_0<a_0\) and its buffer \(a_b=d_0/2\)
inside the original chart, where \(a_0\) is the positive outward
reserve of Corollary~\ref{deg:cor-positive-large-S-outer}.
There are fixed \(g_{\rm tail},k_{\rm tail}>0\);
the full-carrier coefficient rectangle uses the minima of these bounds
and the native lower, paired and high bounds. Every actual lower input whose positive first
lip has raw abscissa \(a=gku\ge a_b\) satisfies
\begin{equation}\label{bp-tail:consumed-derivative}
                   0<\Pi_{\theta,m}'(x)\le C(d_0)g^3.
\end{equation}
This includes every original admitted \(a>d_0\), through all original
window, inverse-image, speed, stop and height truncations, and all
intervening nonroot inputs. The estimate holds uniformly in the independent
coefficient rectangle and \(0\le m\le m_0\). With the actual finite
channel \(K_{\rm fin}>0\), choose \(m_*\) last so the whole retained
raw tail has \(F'>1/2\). No new receiver or first-lip cutoff is imposed.
\end{proposition}

\begin{proof}
Put \(\epsilon=gk\), \(\eta=gY\), so
\(v-g=\epsilon\xi\), \(z=\epsilon^2\eta\). Use \eqref{deg:quadratic-field} in the original
\(\eta=gY\) coordinate, with \(P=1-\pi\xi+3\xi^2+k\xi^3\),
\(B=1+g(2+k\xi)\), \(Q=\ell(g+\epsilon\xi)\),
\(\ell(v)=1-Av+Cv^2\), and \(H_1=Ckg^2(1+k\xi)\).
On the ghost fibre \(P>0\) on the entire positive segment. In the
translated raw positive abscissa \(t=v-g\), define
\[
 \begin{gathered}
 M(t)=t^3+3gt^2-\pi kg^2t+k^2g^3,\quad B_0(t)=1+2g+t,\\
 Q_0(t)=\ell(g+t),\quad\zeta_0(t)=C(g+t),\quad D_0=M+Q_0z.
 \end{gathered}
\]
Thus \(M(\epsilon\xi)=\epsilon^2gP(\xi)\).
Let \(t_*=\delta-g\) be the actual original positive raw port.
The inverse of its incoming passage obeys
\begin{equation}\label{bp-tail:positive-raw-graph}
 z_t=f_0(t,z)=\frac{z[tB_0+\zeta_0z]}{D_0},\qquad
                         z(a)=\ell_0M(a).
\end{equation}
The lip height \(\ell_0\) is fixed using the original input-window
reserve. On \(a_b\le a\le t_*\), independently small \(g,k\) give
\(M(t)\asymp t^3\) on \([a,t_*]\). With \(Q_0\) a positive unit,
\(f_0\le Ct+Cz\); Gronwall gives \(z\le C\delta^2\) inside the
fixed upper chart boundary. The numerator is positive, so
\(z\ge\ell_0M(a)>0\). Hence this is an actual regular
outward flight. Its original input observer \(\chi_+\) is used only
on admitted port images. A nonadmitted zero-flight endpoint at \(a=t_*\)
is an auxiliary graph value and is not assigned an observer value.

Equation~\eqref{deg:scalar-polynomial-lip-slope}, with
\(\sigma=\epsilon\), \(K_l=K\), \(I_+=I\), and \(x_T=t_*\), gives
\begin{equation}\label{bp-tail:positive-full-gain}
 \Gamma_{+,a}'=-a\chi_+'(z_{\rm port})\frac{z_{\rm port}}{D_0(t_*,z_{\rm port})}
                     B_0(a)K_l(a)e^{I_+(a)}.
\end{equation}
The original transverse lip margin makes \(K_l\) a positive unit.
All admitted remaining factors have bounds and reciprocal bounds
depending on \(d_0\). For the original label \(b_u=-u^2/2\),
\begin{equation}\label{bp-tail:positive-raw-slope}
 -\Gamma_{+,a}'\ge c(d_0)a,\qquad
 c(d_0)\epsilon^2\le\Gamma_{+,b_u}'\le C(d_0)\epsilon^2.
\end{equation}
The exact input gain \(\epsilon^2=(gk)^2\) is retained.

We next obtain an absolute first-variation estimate for the positive
inward passage on the whole raw tail. Let
\(\eta_0^+(u)>0\) be its height at the original center \(\xi=0\).
For \(f=\eta(\xi B+H_1\eta)/(gP+Q\eta)\), put
\(W=\eta/(gP)\), \(H_+=P'/(\xi B)\),
\(\varpi=H_1gP/(\xi B)\). Moving-initial differentiation is exact:
\begin{equation}\label{bp-tail:positive-inward-variation}
 \begin{gathered}
 \partial_u\log\eta_0^+=-\frac{uB_+}{gP_+}\Lambda_+(u),\qquad
 \Lambda_+=\frac{K_l(a)}{1+\ell_0Q_+}
                  \exp\int_u^0\left(f_\eta-\frac f\eta\right)d\xi,\\
 f_\eta-f/\eta=\frac{\xi B}{gP}
                 \frac{W(\varpi-Q)}{(1+QW)^2}.
 \end{gathered}
\end{equation}
Indeed \(g\ell_0P'(u)-f(u,g\ell_0P(u))
=-\ell_0uB_+K_l(a)/(1+\ell_0Q_+)\); division by the actual
height cocycle produces the displayed logarithmic variation.
Fix \(v_i>0\) so that \(P\ge p_i>0\) on \([0,2v_i]\).
For \(\xi\ge v_i\) throughout the original raw chart,
\[
 gH_+=\frac{g(6-\pi/\xi+3k\xi)}B\le C(g+t),\qquad
 |\varpi|\le C_{v_i}\{(g+t)^3+g^3k\}.
\]
The fixed chart and lip choices therefore give \(Q-\varpi>0\)
and the same transverse margin at every running point.
The positive action clock \(dq=-\xi B\,d\xi/(gP)\) gives
\[
 W_q=W\left[-\frac{1+\varpi W}{1+QW}+gH_+\right],\qquad W(0)=\ell_0.
\]
The braces increase with \(W\); at \(\ell_0\) they equal
\(-K_l/(1+\ell_0Q)\le-c_1<0\). Thus
\(W\le\ell_0e^{-c_1q}\) down to \(\xi=v_i\).
In raw coordinates this also gives
\(z=M(t)W\le\ell_0M(t)\le C\ell_0\delta^3\), keeping the
entire noncentral inward arc in the original chart.
The noncentral integral in \eqref{bp-tail:positive-inward-variation}
is exactly \(\int W(Q-\varpi)/(1+QW)^2\,dq\), and is bounded.

Every \(a\ge a_b\) has \(u>2v_i\) after independent coefficient
shrinks. The fixed segment \([v_i,2v_i]\) gives
\(A_{\rm in}=\inf\int_{v_i}^{2v_i}\xi B/P\,d\xi>0\).
Consequently its entering height satisfies
\(\eta_i\le Cg e^{-c_1A_{\rm in}/g}\).
On the remaining slab \([0,v_i]\), the linear coefficient
\(\xi B/(gP)\) is nonnegative, and the graph equation has
\(f\ge-Ckg\,\eta^2\). In the decreasing \(\xi\)-direction a
first-exit comparison therefore keeps \(0<\eta\le2\eta_i\).
The denominator remains at least \(gp_i/2\), so the actual
arc reaches \(\xi=0\) in finite positive time. The regular formula
\[
 f_\eta-f/\eta
 =\frac{\eta(H_1gP-\xi BQ)}{(gP+Q\eta)^2}
\]
bounds the remaining logarithmic variational integral by
\(C\eta_i/g^2\), which is exponentially small. This uses only the
original central passage, with no additional observer or endpoint.
Consequently
\begin{equation}\label{bp-tail:positive-inward-absolute}
 0<c\le\Lambda_+,\Lambda_+^{-1}\le C,
 \qquad -\partial_a\log\eta_0^+
             =\frac{aB_0(a)}{M(a)}\Lambda_+>0.
\end{equation}
This proves monotonicity and absolute first variation for every actual
raw-tail input, with no restriction \(a<d_0\).

The negative inward supplier is already on its entire strict interval.
Write \(R=kr\),
\[
 P_-(r)=1+\pi r+3r^2-kr^3,\qquad
 \widehat P(R)=k^2P_-(R/k),\qquad B_-=1+g(2-R).
\]
Corollary~\ref{deg:cor-negative-inward-common-centre} constructs the
actual original-center height \(\eta_0^-(r)=gZ_-(r)\) and retains
\(\log(\ell_0P_-)-A_-/g+\mathcal R_-\). Its exact derivative is
\begin{equation}\label{bp-tail:negative-inward-absolute}
 \begin{gathered}
 -\partial_r\log\eta_0^-=\frac{rB_-}{gP_-}\Lambda_-,\qquad
 \Lambda_-=1-gH_--g\mathcal D_-\mathcal R_-,\qquad
                         1/2\le\Lambda_-\le3/2,\\
 \mathcal D_-=\frac{P_-}{rB_-}\partial_r,\qquad
 r\ge r_c,\quad 0<R<\mathfrak c.
 \end{gathered}
\end{equation}
This includes the arbitrarily close negative root, its actual endpoint
power and all central continuation factors. In particular
\(\eta_0^-\downarrow0\) at \(R\uparrow\mathfrak c\).

Matching does not assume an image intersection or an action estimate
outside its original domain. The original paired comparison has an
actual negative match \(r_*\ge r_c\) at \(u_*\).
Put \(u_b=a_b/(gk)\). The native separate inward estimates give
\(\eta_0^+(u_b)<\eta_0^+(u_*)=\eta_0^-(r_*)\), since their entire
segment has \(a\le a_b<d_0\). Negative
\eqref{bp-tail:negative-inward-absolute} and its zero endpoint
therefore give a unique \(r_b\ge r_*\). This is a match of the
separate actual centre maps, before any raw observer is attached.
It assigns no observer value to a possibly nonadmitted auxiliary
port image at \(a_b\). Set \(R_b=kr_b>0\).
Positive \eqref{bp-tail:positive-inward-absolute}
and negative \eqref{bp-tail:negative-inward-absolute} monotonicity then
give a unique actual negative match for every \(a\ge a_b\):
\begin{equation}\label{bp-tail:whole-tail-matching}
 \eta_0^+(a)=\eta_0^-(r(a)),\qquad R_b\le R(a)<\mathfrak c.
\end{equation}
The order of the actual centre images therefore extends matching beyond
\(gku<d_0\). The absolute output estimate below holds on all
\(r\ge r_c\), without requiring the stronger bound \(R\ge R_b\).
For an input of the original \(I_\theta\), the through trajectory
already reaches both the centre and \(\Sigma_-\). Equality at that
same centre and scalar uniqueness identify the matched negative graph
with this trajectory. Its output-window and closing inverse-image
admission are inherited before an observer is evaluated. The auxiliary
formal shots do not enlarge those observer domains.

We now bound the actual negative output gain at the fixed raw cut
\(t=d_q\), followed by its original regular section bridges.
Since \(R=kr\), the negative spatial radius is \(gR\).
In \(R=kr,W=k^2\eta\), with raw height \(z=g^2W\), the exact
positive-time field, after a positive constant time change, is
\[
 R'=g\widehat P(R)+Q_-(R)W,\qquad
 W'=W\{RB_-(R)-Cg^2(1-R)W\},\qquad
 Q_-(R)=\ell(g(1-R)).
\]
This retains the independent \(k\)-dependence of
\(\widehat P=k^2+\pi kR+3R^2-R^3\).
For the actual raw negative graph let \(t=g-v>0\),
\(\mathcal W=z/t^2\), \(b=gR\), and
\[
 \begin{gathered}
 p=3+\pi/r+r^{-2}-R=\widehat P(R)/R^2,\quad
 m_-(t)=3g-t+\pi kg^2/t+k^2g^3/t^2,\\
 Q_-(t)=\ell(g-t),\quad
 N_-(t,\mathcal W)=1+2g-t-Ct(g-t)\mathcal W,\quad
 D_-=m_-+Q_-\mathcal W.
 \end{gathered}
\]
Its original lip is \(\mathcal W_l=\ell_0gp\).
Corollary~\ref{deg:cor-negative-absolute-outer} proves actual passage
and original-buffer admission on all \(r\ge r_c,0<R<\mathfrak c\),
including \(p\to0\), by its inverse graph and strip climb.

Set \(\nu=6g/(1+2g)\), \(\lambda_g=2\nu\). The exact density
\eqref{deg:quadratic-outer-density} on this negative graph is
\begin{equation}\label{bp-tail:negative-exact-density}
 \begin{gathered}
 I_- =\int_{\log z_l}^{\log z_T}J_-\,d\log z,\\
 J_-=-\frac{t}{N_-}
 \left[3-\nu+\{-A+(3-\nu)C(g-t)\}\mathcal W
                                      -\frac{\pi kg^2}{t^2}\right].
 \end{gathered}
\end{equation}
The output gain contains \(e^{I_-}\), so an upper bound for
\(I_-\) suffices. Choose \(0<\rho_0<\inf\mathfrak c/4\).
For \(R\le\rho_0\), \(p\) has bounds and reciprocal bounds on
\(r\ge r_c\). Its weak \(\log z\)-length is \(O(1+|\log g|)\),
and \(|J_-|\le Cg\) there. This part contributes \(O(g(1+|\log g|))\).
For \(R\ge\rho_0\), choose a fixed small auxiliary height
\(\omega>0\) with
\(\sup|{-A+(3-\nu)C(g-t)}|\omega\le1/4\), then a fixed
\[
             k_{\rm tail}\le\frac{\rho_0^2}{4\Pi}.
\]
It is a separate quadratic-chart bound, fixed before any final raw
coefficient shrink. While \(\mathcal W\le\omega\),
\(t\ge b=gR\) gives \(\pi kg^2/t^2\le\Pi k/\rho_0^2\le1/4\).
After choosing \(g_{\rm tail}\) with \(\nu<1/4\), the bracket
in \eqref{bp-tail:negative-exact-density} is at least \(2\).
Hence \(J_-\le0\) throughout the arbitrarily long root dwell.
No logarithm of \(p\) is estimated by a coefficient bound.
Choose a fixed inner cut \(d_*\ll\omega\) and then shrink \(g\);
on the strong part up to \(d_*\), \(\mathcal W\ge\omega\) and
\(D_-\ge c\mathcal W\). Thus \(d\log z\le C\,d\log t\).
The remaining density
has majorant \(C(t+kg^2/t)\), with integral at most
\(C(d_*+kg^2/b)\le C(d_*+g/r_c)\).
The compact strong continuation from this auxiliary \(d_*\) to the
original fixed raw cut \(d_q\) contributes a fixed bound. Thus on the whole original
negative interval
\begin{equation}\label{bp-tail:negative-density-value}
                              I_-\le C.
\end{equation}
The singular root dwell is controlled by its sign; this estimate imposes
no condition on \(g\log k\).

Here \(T\) denotes \(t=d_q\). Let \(\psi_-\) include the original
regular bridge from this cut to \(\Sigma_-\) and the divided closing,
and choose \(c_-=\psi_-'(z_{\rm ref})>0\) at a coefficient-only
buffered reference. It and its reciprocal are observer units.
The \(s=-1\) specialization of
\eqref{deg:quadratic-outer-exact-slope}, valid algebraically on this
whole negative domain, and the exact label change give
\begin{equation}\label{bp-tail:negative-full-gain}
 \Gamma_{-,b_R}'=c_-g^2\ell_0^{-1-\nu}g^{-\nu}(gR)^{-\lambda_g}
                    p^{-\nu}\mathcal K_-(r),\qquad
 b_R=-R^2/2,\quad0<\mathcal K_-\le C.
\end{equation}
Here \(\mathcal K_-\) is the product \(U_s\) of that formula,
without its initial \(p^{-\nu}\); its upper bound follows from
\eqref{bp-tail:negative-density-value} and the remaining unit factors.
Here \(\mathscr B_-\) is the original moving-initial transverse
factor of \eqref{deg:quadratic-outer-moving-factor}, a negative unit.
The terminal factors are evaluated on their admitted original images.
The power \(p^{-\nu}\), the exact \(\lambda_g\), and the absolute
gain are all retained; \(\mathcal K_-^{-1}\) is not needed.
Changing back to \(b_r=-r^2/2\) gives exactly
\(\Gamma_{-,b_r}'=k^2\Gamma_{-,b_R}'\); no \(k^{-\lambda_g}\)
has been declared a unit.
Using the actual center variation \eqref{bp-tail:negative-inward-absolute},
\begin{equation}\label{bp-tail:negative-consumed-gain}
 0<\frac{d\Gamma_-}{d\log\eta_0^-}
   =\frac{g\widehat P(R)}{B_-\Lambda_-}\Gamma_{-,b_R}'
   \le Cc_-g^3g^{-\nu-\lambda_g}R^{2-\lambda_g}p^{1-\nu}
   \le Cc_-g^3.
\end{equation}
Indeed \(0<p\le C\) on \(r\ge r_c\), \(R<\mathfrak c\), and
\(2-\lambda_g>1,1-\nu>0\). These positive powers pay both
\(R\to0\) and \(p\to0\). Also \(g^{-\nu-\lambda_g}\) is a
unit because \(\nu+\lambda_g=O(g)\), independently of \(k\).
Thus differentiation with respect to the actual centre height absorbs
the possibly singular root factor in the output gain.

Choose \(c_+>0\) as the analogous ordinary input observer slope unit.
Equations~\eqref{bp-tail:positive-full-gain} and
\eqref{bp-tail:positive-inward-absolute} give
\[
 \frac{d\Gamma_+}{d\log\eta_0^+}
   =\frac{-\Gamma_{+,a}'M(a)}{aB_0(a)\Lambda_+(a)}
                           \ge c(d_0)c_+a^3.
\]
On the same actual matched orbit the complete physical derivative is
therefore the exact ratio
\[
 0<\Pi'(x)=
 \frac{d\Gamma_-/d\log\eta_0^-}{d\Gamma_+/d\log\eta_0^+}
 \le C\frac{c_-}{c_+}\frac{g^3}{a^3}\le C(d_0)g^3.
\]
Only the ratio of the original regular observer units is bounded here;
the full exterior gains and their endpoint powers remained in
\eqref{bp-tail:positive-full-gain} and \eqref{bp-tail:negative-full-gain}.
The input center subtraction \(a_\theta\) has zero phase derivative.
The ordinary \(C^3\) divided-closing chain rule of
Lemma~\ref{deg:lem-scalar-moving-identities} is uniform through
\(m=0\) on these original buffered images. All coefficients remain
frozen, independently of \(g,k\), the ghost gap and the phase scale.
\end{proof}


\begin{corollary}[The bounded positive quadratic carrier]
\label{deg:cor-positive-quadratic-whole-carrier}
Fix \(\Pi_0<\infty\), the original nested raw-section windows,
and their ordinary \(C^3\) regular observers and divided finite
closing. The passive compact \((A,C)\)-set is contained in their
prescribed common coefficient neighbourhood, with the original
buffer and observer margins. Use fibres with \(K_{\rm fin}>0\).
For the field \eqref{deg:quadratic-field} there are independent
upper bounds for \(g,k>0\), one common multiplier bound
\(0<m\le m_*\le m_0\), and a fixed raw central-height cap
\(H_{\rm cap}>0\) with the following property.
On each retained real, fold or adjacent ghost fibre,
\(0<\pi\le\Pi_0\), the original enlarged windows admit one
connected physical carrier below \(H_{\rm cap}\). Its lower
interval has at most one curvature zero, which is simple and
crosses downward in the increasing physical input. The actual
\(F(x)=x-m\Pi_\theta(x)\) has at most two zeros of \(F'\), counted
with multiplicity, on the entire carrier. Its overlapping upper
complement has \(F'>0\). The ghost raw tail \(gku\ge d_0\) also
has \(F'>1/2\); it is retained on the original carrier together
with every intervening input. On the remaining root-free fibres,
\(F'>0\).
\end{corollary}

\begin{proof}
We combine the low-height curvature orientation with the paired upper
bound and the raw-tail contraction, all in the original affine input.
The strict overlaps below ensure that these estimates control every
point of the connected carrier before Rolle's theorem is applied.

Use the original affine input and divided output of
Lemma~\ref{lem:closing} and \eqref{sharp:upper-affine}.
Make the comparison and reserve choices of
Proposition~\ref{deg:prop-positive-quadratic}: \(K_0,r_c,X,\Delta_0,\pi_\flat\),
then \(\delta_q,\delta,M\), with \(M\delta\gg\Pi_0\).
Prepare the bounded supplier on \(u<4M\) and retain \(u\le3M\).
Choose \(d_0\) for the small error in
\eqref{deg:quadratic-large-S-curvature}, with
\(0<d_{\rm in}<d_0<\min(a_0,d_q/2)\).
Fix \(0<H_{\rm cap}<z_*\) before coefficient and multiplier bounds,
and apply Proposition~\ref{bp-tail:prop-actual-raw-tail} with
\(a_b=d_0/2\), \(h_c=H_{\rm cap}/2\); include
\(g_{\rm tail},k_{\rm tail}\) in the final rectangle.

The same-centre suppliers give \(N_\pm\ge1/2\), and
\[
 Z_-(r)\downarrow0\quad(kr\uparrow\mathfrak c),\qquad
 Z_-([r_c,\mathfrak c/k))=(0,Z_-(r_c)].
\]
The comparison action margin places \(Z_+(u_{\rm th})\) inside
this image and, by its bounded action, above \(Z_-(r_{\max})\)
for a fixed \(r_{\max}>r_c\). Strict monotonicity therefore gives
every retained lip its unique actual match,
\[
 Z_+(u)=Z_-(r),\qquad
 R_u=\frac{uB_+\widehat P(R)}{RB_-P(u)}\frac{N_+}{N_-}>0.
\]
Real/fold inputs have \(u_{\rm th}\le u<a\), with only \(u<d\)
at the fold; ghost curvature inputs use the original enlarged
raw-window preimage with \(u\ge u_{\rm th},gku<d_0\).
The separate outward suppliers admit both raw images before observers.
Scalar order and Lemma~\ref{lem:loc-through} make the window
preimages and original speed, height and stop clips intervals.

The comparison is admitted to the mother interval: its central
raw height is below \(h_c\), and its central speed is below the
original margin after shrinking coefficients. Below it the positive
port has \(Y/P>\ell_0\) and the centre has \(Y/P<\ell_0\).
Ghost positivity gives an intermediate lip hit. At every hit the
action-clock derivative is strictly negative, so the hit is unique.
Its radial derivative has the opposite sign, whereas scalar first
variation in the original input is positive. Implicit differentiation
therefore gives \(a_x<0\) in that same affine input. This first-hit
construction and the centre-image identification above preserve both
original output and closing clips. Allocate
\[
 I_{\rm tail}=\{x\in I_\theta:x\le x_{\rm comp},\ \Delta<0,
                                      \ a(x)\ge d_0\},\qquad
 I_{\rm low}=\{x\in I_\theta:x\le x_{\rm comp}\}\setminus I_{\rm tail}.
\]
The tail is an initial interval or empty. Its entry bound
\(H\le(gk)^2Cg e^{-c/g}<h_c\) puts its whole admitted part in
Proposition~\ref{bp-tail:prop-actual-raw-tail}, including \(a=d_0\).
That estimate strictly overlaps the native curvature interval on
\(a_b\le a<d_0\); the connected \(I_{\rm low}\) retains the
strict hypothesis \(a<d_0\).

On the bounded range, the two action-three supplier bounds and
bounded rows of \(H_\pm\) give, by the separated quotient rule,
\begin{equation}\label{deg:quadratic-bounded-product-rows}
 \sum_{i+j\le2}|\mathcal D_+^i\mathcal D_-^j(n-1)|\le C_Mg,
 \qquad \mathcal D_+\mathcal D_-n
       =-\frac{(\mathcal D_+N_+)(\mathcal D_-N_-)}{N_-^2}.
\end{equation}
Thus \eqref{deg:quadratic-bounded-lip-curvature} consumes the full
product and section rows. Proposition~\ref{deg:prop-positive-quadratic}
gives the downward crossings on the bounded lip, fixed-normal and
joint-growing packages. The exact resection
\eqref{deg:joint-root-exact-lip-resection} and composition
\eqref{deg:joint-root-physical-composition} identify the same map
on their overlaps, retaining coefficient memory and both gains.

For \(\sigma_1=\pi_\flat/(24M)\), use the positive bounds
\(R_0(\delta),R_0(\sigma_1)\) from
\eqref{deg:quadratic-large-S-action}, and choose
\[
 k_0<\frac{\min\{R_0(\delta),R_0(\sigma_1)\}}{4M},\qquad
 k_0\le\min\{\delta/(2M),k_{\rm tail}\}.
\]
These enforce the proposition's mixed-range exclusion; the remaining
\(R\ge\delta,S\le\delta\) and \(S\ge\delta\) packages have positive
curvature. The closed alternatives cover \(r\ge2M,u\ge2M\),
within the enlarged domains \(S,R<\delta_q\) and
\(S>\delta/2,a<2d_0\); the bounded range \(u<4M\) overlaps at
\(2M\le u\le3M\). Every equality face is retained.
Hence every curvature zero on the connected lower interval is a
simple downward crossing, so there is at most one. Rolle applied
to \(F''=-m\Pi_\theta''\) gives at most two zeros of \(F'\),
with multiplicity; a zero at the curvature crossing is double.

The paired lemma gives \(0<\Pi_\theta'\le C_*\) on the original
buffer interval \([\eta_{\rm comp},\eta_{\max}]\), with
\(\eta_{\rm comp}=gZ_+(u_*)<gZ_+(u_{\rm th})<\eta_{\max}\).
The signed high package applies for
\(\eta_0>gK_H(1+g\pi^2)\), \(H=(gk)^2\eta_0<z_*\).
Shrink independently so
\[
 gK_H(1+g\Pi_0^2)<\eta_{\max}/2,\qquad
 (gk)^2\eta_{\max}<H_{\rm cap}/2.
\]
The paired and high packages strictly overlap, retain the high equality
in the paired package and cover every intervening upper input under
the original clips. Finally take
\(m_*\max(C_*,C_H,C(d_0)g_0^3,C_{\rm contraction})<1/2\).
The tail and upper complement have \(F'>1/2\), adding no critical
point to the same physical \(F\). For \(\pi\le\pi_\flat\),
Lemma~\ref{deg:lem-raw-jacobian} gives \(F'>0\).
\end{proof}


\section{Bounded action, degree loss, and remote roots}
\label{sec:degree-loss-remote-roots}

At degree loss, the incident saddles have different small exponents, so
one limiting action does not describe the whole return interval. We first
connect bounded action to the central source passage, then construct the
growing-action and fixed-exit maps on the same physical sections.
Comparing their actual input and output derivatives gives at most two
zeros of the complete displacement on the entire interval.
\subsection{Bounded negative action and the central source interval}
\label{subsec:bounded-action}

The bounded-footpoint passage and the central source passage use different
sections and phase labels. We connect them by following the same orbits
between those sections. This gives the overlap needed to combine a
negative-curvature estimate with a bounded-derivative estimate.

Let \(\tau>0\), \(\sigma>0\), and
\begin{equation}\label{bd:polynomial}
\begin{split}
 P_\sigma(\xi)
 &=(1+\sigma\xi)(1-\xi/x_*)(1-\sigma j\xi)
   =1-\Lambda_1\xi+3\Lambda_3\xi^2+\Lambda_2\xi^3,\\
 \Lambda_1&=x_*^{-1}-\sigma+\sigma j,\qquad
 \Lambda_2=\sigma^2j/x_*,\qquad
 3\Lambda_3=\sigma(j-1)/x_*-\sigma^2j .
\end{split}
\end{equation}
Here \(x_*\) ranges in a positive compact interval and \(j\) in a
fixed compact subset of \([0,\infty)\).  Physical raw sections require
\(j>0\).  The value \(j=0\) denotes only a one-sided normalized limit.
We take either a positive compact \(j\)-interval, or a sufficiently
small interval \(0\le j\le j_0<1\).  In the latter case
\(\Lambda_3<0\).  In both cases, after decreasing \(\sigma\),
\[
 \Lambda_1\ge\lambda_1>0,\qquad
 0\le\Lambda_2\le C\sigma^2,\qquad |\Lambda_3|\le C\sigma.
\]
No comparison between \(\tau\) and \(\sigma\), or between either of
them and \(j\), is imposed.

The source parameters and physical coordinates are
\begin{equation}\label{bd:source-dictionary}
\begin{gathered}
 (a,p_{\rm s},\gamma_{\rm s},s_{\rm s})
 =(-1,\sqrt\tau\Lambda_1,\tau\Lambda_3,\tau^{3/2}\Lambda_2),\\
 \xi=\sqrt\tau X,\quad\eta=\tau Y,\qquad
 v-\gamma_{\rm s}=e\xi,\quad z=e^2\eta,\quad e=\tau\Lambda_2.
\end{gathered}
\end{equation}
For a source quartet, denote its positive graph denominator by
\[
 \mathcal D_{a,p,\gamma,s}(X,Y)
 =\{1-A(\gamma+sX)+C(\gamma+sX)^2\}Y
       -a-pX+3\gamma X^2+sX^3.
\]
We write \(\mathcal D_{\rm s}\) for the quartet in
\eqref{bd:source-dictionary}; barred subscripts below mean the same
formula with the displayed barred quartet. These denominators determine
where an orbit can be written as a graph.
Writing
\[
 B_\tau=1+\tau(2\Lambda_3+\Lambda_2\xi),\quad
 L_\tau=\Lambda_3+\Lambda_2\xi,\quad
 Q_\tau=1-A\tau L_\tau+C\tau^2L_\tau^2,
\]
the exact graph equation is
\begin{equation}\label{bd:graph}
 \eta_\xi=
 \frac{\eta\{\xi B_\tau+C\tau^2\Lambda_2L_\tau\eta\}}
      {Q_\tau\eta+\tau P_\sigma}.
\end{equation}
Fix \(L>4(1+\sup x_*)\).  On the actual incoming section \(\xi=L\)
use the increasing coordinates
\begin{equation}\label{bd:incoming-phase}
 h_L=\eta+\tau\log(2\eta)-L^2/2,
 \qquad \beta_L=h_L/\tau-\log(2\tau).
\end{equation}
These formulas prescribe initial points; neither expression is asserted
to be constant along a perturbed orbit.  In particular,
\(h_L=0\), or \(\beta_L=b_\tau:=-\log(2\tau)\), is an actual
interior source orbit.

\begin{lemma}[Source ranges on the actual incoming section]
\label{bd:lem-source-ranges}
There is a fixed \(\delta_\sharp>0\) such that the following
constructions are uniform on \(-\delta_\sharp\le h_L\le0\), with
enlarged input and output intervals.
\begin{enumerate}
\item For \(B_0\le\beta_L\le b_\tau\), where \(B_0\) is fixed
large, the source passage exists and its actual central height obeys
\begin{equation}\label{bd:positive-height}
 |Y(0)+\log Y(0)-\beta_L|\le C_L,
 \qquad Y(0)=\beta_L-\log\beta_L+O_L(1).
\end{equation}
\item Every fixed compact \(\beta_L\)-interval reaches a fixed
compact strict source family: on fixed internal source intervals,
\begin{equation}\label{bd:compact-height}
 0<m_B\le Y(0)\le M_B,\qquad
 Y>0,\qquad \mathcal D_{\rm s}\ge c_B(Y+1)>0.
\end{equation}
\item At a negative base \(\beta_L=-B\), \(B\ge B_-\), freeze
\begin{equation}\label{bd:frozen-dictionary}
 q=B^{-1/2},\quad \rho=\sqrt{\tau B},\quad
 (x,y)=(\xi/\rho,\eta/\rho^2),\quad
 (\lambda_1,\lambda_2,k)
       =(\rho\Lambda_1,\rho^3\Lambda_2,\rho^2\Lambda_3).
\end{equation}
For fixed cuts \(x=\pm L_{\rm f}\), \(L_{\rm f}=8\), this orbit
belongs to a common enlarged negative-action family with core
\(|x|\le R_{\rm f}=5\) and \(P_\sigma(\rho x)\ge3/4\).
The input and output energy windows can be taken respectively inside
\([-7/4,-3/8]\) and \([-2,-1/4]\), with strict margins.
\end{enumerate}
All these same-orbit passages reach \(\xi=\pm L\) with
\begin{equation}\label{bd:parabolic-cuts}
 \eta(\pm L)=L^2/2+
 O_L\bigl(\delta_\sharp+\tau|\log\tau|\bigr).
\end{equation}
Their outward continuations to the fixed physical sections have
positive graph denominators and retain the actual regular maps.
\end{lemma}
\begin{proof}
For the first assertion, freeze \(B=\beta_L\ge B_0\), put
\(q=B^{-1/2}\), \(\rho=\sqrt{\tau B}\), and use the source
scaling \((\xi,\eta)=(\rho x,\rho^2y)\).  Its coefficients are
\[
 (\bar a,\bar p,\bar\gamma,\bar s)
 =(-q^2,\sqrt\tau\Lambda_1q,\tau\Lambda_3,
                         \tau^{3/2}\Lambda_2/q).
\]
Set \(H=y+q^2\log(2y)-x^2/2\) and \(R=L/\rho\).
On a provisional quadratic corridor \(y\asymp1+x^2\), direct
subtraction in the exact graph equation gives
\begin{equation}\label{bd:energy-estimates}
\begin{split}
 |H_x|&\le C_L\{|
 \bar p|+|\bar\gamma|(1+|x|)+\bar s(1+x^2)\},\\
 |\partial_HH_x|&\le C_L\left\{
 \frac{|\bar p|}{1+x^2}+
 \frac{|\bar\gamma|}{1+|x|}+\bar s\right\}.
\end{split}
\end{equation}
The terms with \(\bar p=\bar\gamma=\bar s=0\) cancel exactly.
The first integral on \([-R,R]\) is \(O_L(q^2)\), since its three
largest products are
\(L|\Lambda_1|q^2,L^2|\Lambda_3|q^2,L^3\Lambda_2q^2\).
The second integral tends to zero uniformly for
\(B_0\le B\le b_\tau\).  Start at \(x=0\) with
\(H(0,y)\in[1/4,3]\).  The monotonicity of
\(y+q^2\log(2y)\) improves the bounds at any first contact with the corridor boundary:
\[
 (1+x^2)/16<y<4(1+x^2),\qquad \mathcal D_{\rm bar}\ge cy.
\]
Thus both cuts are reached.  The scalar energy variation is a positive
exponential with exponent bounded by the second integral in
\eqref{bd:energy-estimates}; hence the incoming energy map has an
actual inverse on a common enlarged interval.  Its prescribed value at
\(x=R\) is exactly \(1+B^{-1}\log(2/B)\).  Inverting and undoing
the scale gives \eqref{bd:positive-height}.  This also constructs the
orbit \(h_L=0\) and a real larger source chord through it.

For compact \(\beta_L\), work instead in \((X,Y)\) and put
\(H_0=Y+\log(2Y)-X^2/2\), \(R_0=L/\sqrt\tau\).
At the incoming cut, \(H_0=\beta_L+\log2\).  On a provisional
bounded energy region the denominator is at least \(c(Y+1)\): near
\(\xi=0\), \(P_\sigma\ge1/2\), while away from that fixed small
interval bounded energy forces \(Y\ge X^2/4\), which dominates
the bounded polynomial.  The exact energy subtraction gives
\[
 |H_0'|\le C_L\{p_{\rm s}+|\gamma_{\rm s}|(1+|X|)
                +s_{\rm s}(1+X^2)\}(1+|H_0|).
\]
Its coefficient has a bounded integral because
\(p_{\rm s}R_0=L\Lambda_1\),
\(|\gamma_{\rm s}|R_0^2=L^2|\Lambda_3|\), and
\(s_{\rm s}R_0^3=L^3\Lambda_2\).  Gronwall's inequality improves a
fixed sufficiently large energy bound, and proves continuation.
Inverting \(Y+\log(2Y)\) at the centre proves
\eqref{bd:compact-height}.  Choose larger compact incoming labels
first to obtain strict image margins.  The resulting internal energy
coordinate may retain a bounded translation; no identity with
\(\beta_L+\log2\) on another section is used.

For the negative base in \eqref{bd:frozen-dictionary}, the energy
\(H_{\rm f}=y+q^2\log(2y)-x^2/2\) at \(x=L/\rho\) is
\begin{equation}\label{bd:frozen-incoming}
 b_B(\beta)=q^2\beta+q^2\log(2q^2),\qquad
 b_B(-B)=-1-\frac{\log(B/2)}B.
\end{equation}
Also \(\rho^2=-h_L+\tau\log(2\tau)\le-h_L\).
Choose \(\delta_\sharp\) so that
\(P_{\rm f}(x):=P_\sigma(\rho x)\ge3/4\) on
\([-8,8]\).  On the positive quadratic tail the exact subtraction
gives
\[
 |(H_{\rm f})_x|
 \le Cq^2\{\rho\Lambda_1+
             \rho^2|\Lambda_3||x|+\rho^3\Lambda_2x^2\}.
\]
Its integral to \(L/\rho\) is \(O_L(q^2)\).  A first-contact
argument in a quadratic corridor therefore reaches \(x=8\), with
the actual energy in the interior of \([-3/2,-1/2]\), after increasing \(B_-\). This places the orbit in the negative-action
family. Its relative curvature estimate will be proved below.

To obtain phase derivatives on a neighbourhood of these orbits, we
construct the full negative core on the larger input window
\([-7/4,-3/8]\).  Fix the finite derivative budget \(m\) and an
integer \(M>2m+8\).  Follow the positive fast arc to
\(y=q^{M+2}\).  The logarithmic-height equation has
\(dH_{\rm f}/d\log y=O_{C^m}(q^2\rho)\); its clock length is
\(O(|\log q|)\).  It cannot reach \(x=1/2\) first, since that
would imply \(H_{\rm f}\ge-1/8+o(1)\).  Hence its actual lip
coordinate satisfies
\[
 p=\sqrt{-2b}+O_{C^m}(q^2|\log q|),\qquad 1/2<p<2.
\]
In the thin core put \(\mathfrak q=q^2\log(y/q^2)\) and
\[
 Q_{\rm f}(x)=1-Aq^2(k+\lambda_2x)
                        +Cq^4(k+\lambda_2x)^2.
\]
Until its first return to the same lip height, its exact equation is
\begin{equation}\label{bd:frozen-core}
 \mathfrak q_x=
 \frac{x\{1+q^2(2k+\lambda_2x)\}
       +Cq^6\lambda_2(k+\lambda_2x)e^{\mathfrak q/q^2}}
      {P_{\rm f}(x)+Q_{\rm f}(x)e^{\mathfrak q/q^2}}
 =\frac{x}{P_{\rm f}(x)}+O(q^2\rho)+O(q^M).
\end{equation}
The actions \(\int_0^p v/P_{\rm f}(v)\,dv\) and
\(\int_0^r v/P_{\rm f}(-v)\,dv\) have opposite comparison
differences at \(r=p/2\) and \(r=2p\), with fixed margins.
The logarithmic height first decreases, stays a fixed amount below
the lip near zero, then increases strictly on the negative half.
It therefore has one first exit before \(|x|=5\), with
\[
 r=p+O(\rho+q^2|\log q|)>1/4.
\]
The denominator in \eqref{bd:frozen-core} is at least \(3/4\).
Each exponential derivative costs at most two powers of \(q^{-1}\);
the choice of \(M\) and the transverse exit equation give bounded
finite phase derivatives on nested windows.  The negative fast arc
then reaches \(x=-8\), with
\(b_-=b+O_{C^m}(\rho+q^2|\log q|)\) strictly inside
\([-2,-1/4]\).  The exact independent energy coordinates of
Theorem~\ref{neg:thm-relative} differ from \(H_{\rm f}\) by
\(O_{C^m}(|k|)\) at these fixed positive-height cuts.  Slightly larger
windows thus contain those coordinates and their images as well.

The same energy estimates on both outward tails prove
\eqref{bd:parabolic-cuts}.  In raw variables, the quadratic corridor
of \eqref{eq:parabolic-tail} has a denominator comparable to the
squared radial variable.  Choose \(L\) large and
\(\delta_\sharp\) small first.  All the enlarged chords above enter
its interior, so the first-contact argument continues them to the
fixed raw sections, across any intervening axis root at positive
height.  The fixed regular transitions are then included.  These
are actual first hits of the original field throughout.
\end{proof}

\begin{lemma}[The bounded annulus]
\label{bd:lem-annulus}
Fix positive constants \(p_-,r_{\rm a},R_{\rm BD}\), with
\(p_-<\inf x_*\).  Consider one connected family with positive and
negative strong-axis footpoints
\[
 p_-\le p<x_*,\qquad r_{\rm a}/2\le r\le2R_{\rm BD},
 \qquad b_+=-p^2/2,\quad b_-=-r^2/2.
\]
The projections are synchronized to fixed weak sections, and their
outward continuations use the fixed physical sections of
Lemma~\ref{lem:closing}.  On slightly larger such domains, normalize
each actual embedding by its own exact value and positive first
derivative.  The resulting complete map
\(\Pi=\phi_-\circ S\circ\phi_+^{-1}\) satisfies
\begin{equation}\label{bd:annular-curvature}
 \phi_+'\frac{P_\sigma(p)}{p+r}\frac{\Pi''}{\Pi'}
 =3\Lambda_2-\frac{\Lambda_1}{pr}+E,\qquad
 \|E\|_{C^1_{b_+}}\longrightarrow0 .
\end{equation}
The derivative of the error follows the actual matched family.
That family has a uniform positive gap from the positive root, as
proved below; convergence includes the reciprocal coefficient face
\(j=0\). For sufficiently small coefficients this map is strictly concave and its displacement
has at most two zeros.
\end{lemma}
\begin{proof}
The curvature formula is uniform only after the incoming footpoint has
been separated from the positive root. We first obtain this separation
by matching the two normal heights. On fixed
small strong lips the normal rates are positive units, including at
\(p=x_*\), since
\[
 -P_\sigma'(x_*)\ge c_*,\qquad
 P_\sigma(-r)\ge c_-,\qquad rB_\tau(-r)\ge c_- .
\]
The strong projection and its integrating normal coordinate are uniform
through the required finite phase order: the backward tangential Green
factor is at most \(Ce^{C\tau t}\), whereas the normal variations are
\(C(1+t)^Ne^{-ct}\). Their integral therefore defines the projection
and the coordinate satisfying exactly \(\dot z=-pB_\tau(p)z\) on
the positive side, with the corresponding equation in reversed time on the negative side.
Use the original fixed weak cuts \(c_+<p_-\) and \(c_-<r_{\rm a}/2\).
Integration along the projected axis gives their actual normal heights
\[
 w_+(p)=u_+(p)e^{-J_+(p)/\tau},\qquad
 w_-(r)=u_-(r)e^{-J_-(r)/\tau},\qquad u_\pm^{\pm1}\le C,
\]
where the units and their finite phase derivatives are bounded on the
fixed lips, independently of any outer cut, and
\[
 J_+(p)=\int_{c_+}^p\frac{xB_\tau(x)}{P_\sigma(x)}\,dx,
 \qquad J_-(r)=\int_{c_-}^r\frac{xB_\tau(-x)}{P_\sigma(-x)}\,dx .
\]

On the intervening weak interval \(P_\sigma,Q_\tau\ge c>0\), so
\(-\dot\xi=\tau P_\sigma+Q_\tau\eta\ge c\tau\) and its travel
time is at most \(C/\tau\). Put
\(H=\xi B_\tau+C\tau^2\Lambda_2L_\tau\eta\), so
\(d\log\eta/dt=-H\). At a critical height,
\(d^2\log\eta/dt^2=H_\xi(\tau P_\sigma+Q_\tau\eta)>0\).
The standing passive neighborhood has \(C>0\), and
\(H_\xi=1+\tau(2\Lambda_3+2\Lambda_2\xi)
+C\tau^2\Lambda_2^2\eta>1/2\) on this word. Thus there is no
interior height maximum; the buffered weak endpoint heights bound the
whole passage and its logarithmic rate. At each fixed weak section
the integrating normal coordinate and \(\eta\) differ by a positive
unit. For the actual map \(M_\tau\) between those normal coordinates,
\[
 \left|\tau\log\frac{M_\tau(w)}w\right|\le K .
\]
Its positive first multiplier \(M_\tau'(0)\) is retained. Matching
the two heights now bounds \(J_+(p)\) by \(\sup J_-(r)+K+C\tau\).
The latter is uniform for \(r\le2R_{\rm BD}\), whereas the simple
positive root gives
\[
 J_+(p)\ge a_*\log\frac{\ell_*}{x_*-p}-C,
 \qquad a_*>0 .
\]
The logarithmic residue \(x_*B_\tau(x_*)/[-P_\sigma'(x_*)]\)
is bounded below by \(a_*\).
Consequently every actual member has \(p\le x_*-d_{\rm a}\), for
one fixed \(d_{\rm a}>0\). This gap depends only on the fixed lips and middle passage, so the
outer cuts can now be chosen without changing it.

Enlarge the input within \(p<x_*-d_{\rm a}/2\). The limiting
root-free action has a fixed compact output envelope: at \(\sigma=0\),
its negative action grows linearly at infinity. Choose strict lower
and upper exit brackets there, then decrease \(\sigma\) independently
so they and their polynomial margins persist. Denote an enlarged
upper output bound by \(R_{\rm a}\ge2R_{\rm BD}\). We can therefore choose the annular cuts
\[
 L_{\rm a}>4(1+\sup x_*+R_{\rm a}).
\]
The source cut \(L\), its label \(h_L\), and
\(\delta_\sharp(L)\) stay fixed. The positive-height arcs from the
fixed lips to \(\pm L_{\rm a}\) and their small enlarged chords
have uniform height and denominator margins. Independent final
coefficient bounds place \(eL_{\rm a}\) inside the original raw
ports; scalar uniqueness identifies both decompositions there.

Apply the compact coefficient-annulus argument in the proof of
Corollary~\ref{neg:cor-rootfree}, with \(t^2=\tau\),
\(\lambda=(\Lambda_1,\Lambda_2)\), and \(k=\Lambda_3\).
The root-free polynomial margin, strict exit brackets and nonzero
negative exit slope established above are its inputs. The lip and
first-exit estimates give absolute \(C^3\) error
\(O(\tau(1+|\log\tau|))\), including synchronized resection;
no division by the reversible coefficient distance is needed. If
\(S_0(-p^2/2)=-r_0^2/2\) is defined by the exact limiting action,
\begin{equation}\label{bd:annular-first}
 \|S-S_0\|_{C^3}\longrightarrow0,\qquad
 S_0'=\frac{P_\sigma(-r_0)}{P_\sigma(p)} .
\end{equation}
The exact polynomial identity therefore gives
\[
 \frac{P_\sigma(p)}{p+r_0}\frac{S_0''}{S_0'}
 =\frac{P_\sigma'(p)/p+P_\sigma'(-r_0)/r_0}{p+r_0}
 =3\Lambda_2-\frac{\Lambda_1}{pr_0} .
\]

The two outer embeddings may have very different first derivatives.
We estimate each embedding relative to its own exact first derivative.
On a fixed small positive-height lip let \(p(t)\) be the synchronized
axis solution and \(w=\eta-\xi^2/2+p(t)^2/2\).
The exact field gives
\[
 \dot w=\tau\{\xi P_\sigma(\xi)-pP_\sigma(p)\}
 +\xi\eta(Q_\tau-B_\tau)
 -C\tau^2\Lambda_2L_\tau\eta^2.
\]
Since \(|\xi+p|\) is bounded below, this has the form
\(\dot w=\tau a(t)w+\tau\eta d(t)\), with bounded finite
derivatives and exponentially decreasing \(\eta\).  The fixed
weak section is a positive axis distance away and the axis speed is
\(O(\tau)\), so its remaining time is at least \(c/\tau\).
Terminal hitting-time derivatives multiply the same exponentially
small normal height.  Variation of constants, with four phase
derivatives, gives \(w=O_{C^4}(\tau)\); resecting a moving lip
costs \(O_{C^4}(\tau|\log\tau|)\).  On both fixed outer cuts,
therefore,
\begin{equation}\label{bd:annular-tail-data}
 \eta_\pm(\pm L_{\rm a},b)=L_{\rm a}^2/2+b+
                  O_{C_b^4}(\alpha_\tau),\qquad
 \alpha_\tau=\tau(1+|\log\tau|).
\end{equation}

Set \(s=e\xi\), \(Z=e^2\eta\), \(k_0=\Lambda_3\).
The exact raw tail equation is
\[
 Z_s=\frac{n_2Z^2+n_1Z}{\ell_oZ+m_o},\quad
 \begin{aligned}
 \ell_o&=1-A(\tau k_0+s)+C(\tau k_0+s)^2,\\
 m_o&=s^3+3\tau k_0s^2-\tau\Lambda_1es+\tau e^2,\\
 n_1&=s(1+s+2\tau k_0),\quad n_2=C(\tau k_0+s).
 \end{aligned}
\]
The initial data \eqref{bd:annular-tail-data} enter the invariant
corridor \(Z\asymp s^2\), \(\ell_oZ+m_o\asymp s^2\).
Division writes the right side as
\(F=(n_2/\ell_o)Z+q_o+\mathfrak r/(\ell_oZ+m_o)\), where
\[
 |\mathfrak r|\le C(|s|^4+\tau|k_0||s|^3
                       +\tau e|s|^2+\tau e^2|s|).
\]
If \(j\) is bounded away from zero,
\(\tau|k_0||\log e|\to0\), because
\(k_0=O(\sigma)\) and \(e\asymp\tau\sigma^2\).
Thus the first variation \(J_1=Z_b\) is comparable to \(e^2\).
For the reciprocal range \(j\to0\), retain instead \(k_0<0\).
Direct differentiation gives the exact first cocycle
\[
 F_Z=\frac{n_2\ell_oZ^2+m_o(n_1+2n_2Z)}{(\ell_oZ+m_o)^2},
 \qquad \operatorname{sgn}(s)(n_1+2n_2Z)\ge c|s| .
\]
The latter follows from \(Z\asymp s^2\) and the small fixed raw
radius. At the actual common denominator,
the term \(3\tau k_0s^2\) in \(m_o\) contributes nonpositively
to \(\operatorname{sgn}(s)F_Z\) on both outward tails.  The rest gives
\[
 \operatorname{sgn}(s)F_Z\le
 C(1+\tau e/|s|^2+\tau e^2/|s|^3),\qquad 0<J_1\le Ce^2.
\]
The normalization below uses this exact positive derivative and requires
only the displayed upper bound. For \(2\le n\le3\), in both cases,
\[
 \int |F_{Z^n}|J_1^{n-1}\,d|s|
 \le C(eL_{\rm a}^{3-2n}+\tau|k_0|L_{\rm a}^{2-2n}
                       +\tau L_{\rm a}^{1-2n}+\tau L_{\rm a}^{-2n})\le C\tau.
\]
The quantitative normalization following
\eqref{tail:relative-jet-equations}, through order three, now applies
with initial errors from \eqref{bd:annular-tail-data} and
\(\int|F_{Z^n}|J_1^{n-1}=O(\tau)\), \(n=2,3\).
Each tail and its actual-image inverse are consequently
\(\mathrm{id}+O_{C^3}(\alpha_\tau)\) after normalization by
the retained \(J_1(b_*)>0\). At this actual image the ordinary
\(C^3\) observer adds \(O_{C^3}(J_1(b_*))\) by
\eqref{phys:eq-observer-normalization}. Thus
\(\phi_\pm=\mathrm{id}+o_{C^3}(1)\), including when their
individual gains are much smaller than \(e^2\).

Finally the exact chain rule is
\[
 \frac{\Pi''}{\Pi'}=\frac1{\phi_+'}
 \left\{\frac{S''}{S'}+
       \frac{\phi_-''(S)}{\phi_-'(S)}S'
       -\frac{\phi_+''}{\phi_+'}\right\}.
\]
On the compact family, \(S\), its first three derivatives and
its positive first inverse derivative are bounded. Multiply by
\(P_\sigma(p)/(p+r)\) and use \(r-r_0\to0\) in phase
\(C^3\). Both exterior corrections are \(o_{C^1_{b_+}}(1)\).
This proves \eqref{bd:annular-curvature}.  Its principal term is at most
\(3C\sigma^2-\lambda_1/(2(\sup x_*)R_{\rm BD})<0\).
Positive affine changes preserve this sign in the actual input.
Thus \(x-m\Pi(x)\) is strictly convex, proving the two-zero bound.
\end{proof}

\begin{proposition}[Bounded-footpoint domains and derivatives]
\label{bd:prop-bounded-action}
On one ordered through interval ending at the actual cut \(h_L=0\),
suppose the two fixed comparisons described below and the enlarged
annular domains of Lemma~\ref{bd:lem-annulus} use the same actual
physical sections. The union of the central source interval
\(-\delta_2\le h_L\le0\) and the part with negative strong
footpoint \(0<r\le2R_{\rm BD}\) has an overlapping cover by the
bounded annulus and the negative, compact, and positive source-label
intervals constructed below. All endpoints are fixed before choosing
a displacement zero.

The completed map has strictly negative curvature on the annulus and
on the entire negative source interval. On the compact and positive
source intervals its derivative in the actual affine input is bounded
by one constant. After one common multiplier shrink, the displacement
has positive derivative on these latter intervals, including the cut
\(h_L=0\). The same assertions hold on the enlarged intervals used
for the overlaps. The degree-loss passage below verifies all domain
and overlap hypotheses of this proposition.
\end{proposition}
\begin{proof}
Fix the strict source domain and the constants for its finite phase
derivatives. Choose $B_-$ large, then $\delta_\sharp$ small enough
for the relative negative-curvature estimate below. Choose a fixed $0<\delta_1\le\delta_\sharp/8$,
set $p_0=\sqrt{2\delta_1}$, and require $P_\sigma\ge1/2$ on
$[-2p_0,2p_0]$.
The orbit prescribed by \(h_L=-\delta_1\) is constructed by a
compact fast arc to a small lip, followed by the core first-exit
argument in \eqref{bd:frozen-core}.  Its limiting incoming footpoint
is \(p_0\), and the limiting outgoing footpoint is determined by
\[
 \int_0^r\frac{v}{P_\sigma(-v)}\,dv
       =\int_0^{p_0}\frac{v}{P_\sigma(v)}\,dv.
\]
The comparison values \(r=p_0/2,2p_0\) have strict opposite action
differences.  Thus the actual two footpoints lie in
\((p_0/2,2p_0)\) after shrinking \(\tau\).
Choose \(r_{\rm a}<p_0/4\).  A second fixed comparison
\(h_L=-p_{\rm a}^2/2\), with \(0<p_{\rm a}<r_{\rm a}/8\), has
\(r<r_{\rm a}/4\) and \(p>p_{\rm a}/2\).
Scalar order and the actual synchronized projections imply that
every \(r\ge r_{\rm a}/2\) in the bounded annulus has
\(p\ge p_{\rm a}/2\), while every \(r<r_{\rm a}\) lies in
\(-\delta_1<h_L\le0\).  Both comparisons have strict margins,
so these ranges have an open overlap.  The comparison uses order along the same orbits; \(h_L\) need not equal
the footpoint energy on the other section.

Fix \(\delta_1<\delta_2<\delta_\sharp\).  On
\(-\delta_2\le h_L\le0\), divide the actual incoming labels into
\[
 [\beta_{\min},-B_-],\quad[-B_-,B_+],\quad[B_+,b_\tau],
 \qquad \beta_{\min}=-\delta_2/\tau-\log(2\tau).
\]
Each interval is used with slightly larger fixed phase thresholds.
For a negative base, Lemma~\ref{bd:lem-source-ranges} supplies the
independent enlarged windows $\widehat J_+=[-7/4,-3/8]$ and
$\widehat J_-=[-2,-1/4]$, frozen cuts $L_f=8$ and core $R_f=5$.
The actual incoming frozen energy is strictly inside $[-3/2,-1/2]$;
its separate ideal coordinate differs by $O(|k|)$. Choose a fixed
$J\Subset\widehat J_+$ containing these values with a strict buffer.
With the frozen dictionary~\eqref{bd:frozen-dictionary}, put
\[
 \rho_{\rm id}=(\lambda_1^2+\lambda_2^2)^{1/2}.
\]
Use the tail part of the proof of Theorem~\ref{neg:thm-relative}
with \(t=q\), \(\lambda=(\lambda_1,\lambda_2)\) and the frozen
\(k\). Its coefficients are exactly
\((a,p,\gamma,s)=(-q^2,q^2\lambda_1,q^2k,q^2\lambda_2)\),
and its auxiliary cuts are \(v=q^2k\pm d_H\), with fixed \(d_H\).
The enlarged windows above supply its quadratic corridors;
at \(L_f=8\) its fourth-order kernel majorants become small after
a \(q\) shrink. The same sequential intrinsic \(p\)- and
\(s\)-secants, with the two reference labels fixed, give
\begin{equation}\label{bd:intrinsic-relative-tails}
 \|\phi_\pm^{\rm int}-\mathrm{id}\|_{C^3}
 +\|(\phi_\pm^{\rm int})^{-1}-\mathrm{id}\|_{C^3}
 \le Cq^2\rho_{\rm id}.
\end{equation}
For the actual ordinary \(C^3\) observers, including divided closing,
the separate raw gains satisfy, for fixed \(0<\nu<1\),
\begin{equation}\label{bd:relative-observer-gain}
 0<d\le Cs_f^{2-\nu}\le Cs_f
   =Cq^2\lambda_2\le Cq^2\rho_{\rm id},\qquad
 s_f=q^2\lambda_2=e\rho .
\end{equation}
Equation~\eqref{phys:eq-observer-normalization} on the fixed
buffered cut-label chord therefore preserves
\eqref{bd:intrinsic-relative-tails} for the observed tails and their
actual-image inverses. The fourth phase order is used only in the
intrinsic polynomial kernel, before composing with these observers.

On the negative source interval the limiting curvature tends to zero
with the reversible coefficient distance. We therefore compare the
actual core \(S_q\) with its limiting action \(S_0\) before dividing
by that distance. The exact reversible energy gives
\(S_q(0,k)=S_0(0,k)=\mathrm{id}\). Integrating the coefficient
variation in \eqref{neg:compact-mixed-convergence} before taking any
limit yields
\[
 \|S_q-S_0\|_{C^2}
 \le C\rho_{\rm id}q^2(1+|\log q|).
\]
One intrinsic coefficient secant with two phase derivatives produces
the factor \(\rho_{\rm id}\). This makes the error small relative
to the core curvature even as that curvature vanishes.
For \(b\in J\), put \(p=\sqrt{-2b}\), and let \(r_0\) be its
limiting action exit. Equation~\eqref{neg:action-curvature} gives
\[
 \frac{S_0''(b)}{S_0'(b)}
 =\frac{p+r_0}{P_{\rm f}(p)}
       \left(3\lambda_2-\frac{\lambda_1}{pr_0}\right).
\]
The even coefficient \(k\) cancels exactly. Here
\(\rho^2=\tau B\le-h_L\), \(\rho_{\rm id}\asymp\rho\),
\(\lambda_1/\rho_{\rm id}\ge c_1>0\), and
\(\lambda_2/\rho_{\rm id}=O(\rho^2\sigma^2)\).
The compact root-free margins and positive first variation thus give
\(S_0''/\rho_{\rm id}\le-c_0<0\), without a Taylor error in \(k\).

The complete curvature also contains the two section terms:
\[
 \Pi''(\phi_+(b))=
 \frac{\phi_-'(S_q(b))}{\phi_+'(b)^2}
 \left\{S_q''+(S_q')^2\frac{\phi_-''(S_q)}{\phi_-'(S_q)}
                  -S_q'\frac{\phi_+''}{\phi_+'}\right\}.
\]
The separate tail bounds above make the two section terms
\(O(q^2\rho_{\rm id})\), while the prefactor is a positive unit
and differs from one by that amount. Consequently, at the actual
normalized input \(u_*=\phi_+(b)\),
\begin{equation}\label{bd:relative-negative-curvature}
 \frac{\Pi''(u_*)}{\rho_{\rm id}}
 =\frac{S_0''(b)}{\rho_{\rm id}}
       +O\bigl(q^2(1+|\log q|)\bigr)<-c .
\end{equation}
Choose \(B_-\) large enough to absorb this error, then
\(\delta_\sharp\) small enough for the root-free domain and the
fixed relative coefficient margins. No comparison of \(q\) with \(\rho\) is used.
The incoming resection identifies the point $u_*$ at which this sign
is evaluated. For the two retained positive first gains,
\begin{equation}\label{bd:relative-raw-curvature}
 (\Pi^{\rm raw})''=\frac{g_-}{g_+^2}\Pi'',\qquad g_+,g_->0.
\end{equation}
Each negative base is fixed while differentiating, so this is one
strict sign on the whole negative interval, including nonroot inputs.

On the compact middle interval, \eqref{bd:compact-height} and the
larger label windows give a compact strict core and positive variation
units $A_\pm$. Before normalization its two auxiliary raw variations are
\begin{equation}\label{bd:compact-raw-variation}
 \frac{dz_\pm}{dY_0}
 =s_s^{2-\lambda_\gamma}(d_{H,\pm}/L_t)^{\lambda_\gamma}A_\pm U_\pm,
 \quad s_s=\tau^{3/2}\Lambda_2,\quad
 \lambda_\gamma=\frac{12\gamma_s}{1+2\gamma_s},\quad
 A_\pm^{\pm1},U_\pm^{\pm1}\le C.
\end{equation}
The common singular power cancels in $dz_-/dz_+$; distinct fixed
port radii have a bounded ratio. The actual regular bridges and divided
closing have positive derivative units, so $0<\Pi'\le C_B$.
This cancellation gives the derivative bound without a lower bound on a
raw gain or a bound on $\gamma_s\log s_s$.
On the positive interval, \eqref{bd:positive-height} gives, after
choosing \(B_+\) large,
\[
 Y(0)>K(1+p_{\rm s}^2),\qquad
 0<s_{\rm s}^2Y(0)\le C_L\tau^3\sigma^4j^2|\log\tau|<z_*.
\]
Since \(\alpha_0=-s_{\rm s}^2\) and
\(\alpha_1=s_{\rm s}p_{\rm s}\), these are the hypotheses of
Corollary~\ref{rew:cor-high-derivative-signed} on the actual central
section. Its complete derivative bound holds on this entire interval,
including \(h_L=0\) and the enlarged source chord through that orbit.
Take the maximum of the compact and high-height derivative constants,
then choose \(m\) so that \(\Delta'>1/2\) on both intervals.
The annular curvature comes from Lemma~\ref{bd:lem-annulus}, on its
open overlap with the source interval constructed by the two fixed
comparisons. All assertions concern restrictions of one physical
map and every input in those overlaps. Boundary coefficients denote
normalized one-sided limits; no orbit at a collapsed raw section is
asserted.
\end{proof}

\subsection{Degree-loss and remote-root compactifications}
\label{subsec:deg-degree-drop}

At degree loss the three equatorial roots occupy two scales, and the two
incident saddle exponents vanish at different rates. We construct one
through interval before dividing it according to accumulated action.
This ensures that the local estimates can later be combined for the
complete physical displacement.

\begin{proposition}[Degree-loss passage preparation]\label{deg:prop-degree-drop}
Let \(\delta=\kappa\tau\), with \(0<\kappa,\tau\ll1\), and assume
\begin{equation}
 \begin{gathered}
 a=\kappa A_{\rm d},\qquad
 \mathcal R a c=B\kappa,\qquad
 a+c=1-3\nu,\\
 \mathcal R(1-3\nu)=\omega+B\kappa .
 \end{gathered}
 \label{deg:degree-graph}
\end{equation}
The quantities
\[
 B,\ \omega,\ \mathcal R,\ A_{\rm d},\ c,\ 1-3\nu,\ c-a
\]
and their reciprocals are required to lie in prescribed compact
intervals.  Consider the source coefficients
\begin{equation}
 (a_s,p_s,\gamma_s,s_s)=
 \left(-1,\sqrt\tau\frac{\omega}{\sqrt B},
       \tau\kappa\mathcal R\nu,
       \tau^{3/2}\kappa^2\mathcal R\sqrt B\right).
 \label{deg:degree-source}
\end{equation}
On the fixed physical passage, below the common raw-height bound, the
ordered through interval is covered by the bounded-source, growing-action,
finite-middle-height, and fixed-exit ranges constructed below, together
with their high-height continuation. The actual matching maps,
derivative estimates, and strict overlaps are uniform on the displayed
compact coefficient family for positive \(\kappa,\tau\). At
\(\kappa\tau=0\), only the one-sided limits of normalized estimates
are asserted. The finite-\(V\) and fixed-exit ranges retain the original
section maps and finite closing on their actual buffered domains.
\end{proposition}

\begin{proof}
Set
\[
 r_s=\delta\mathcal R,\qquad
 v-\gamma_s=r_su,\qquad z=r_s^2w,\qquad
 \gamma_s=r_s\nu .
\]
The full two-dimensional system, after a positive change of time, is
\begin{equation}
 \dot u=\delta p_d(u)-K_d(u)w,\qquad
 \dot w=-w\{L_d(u)+E_d(u)w\},
 \label{deg:degree-field}
\end{equation}
where \(p_d=p_{\rm d}\) from \eqref{fx:pd} and
\((K_d,L_d,E_d)=(K,L,E)\) from \eqref{fx:coefficients}.
Thus the saddle rates and exponents are exactly those in
\eqref{fx:exponents}. On the present coefficient family,
\begin{equation}
 \chi=\frac{\delta\ell_+}{b_+}=\tau X_d,\qquad
 \beta=\frac{\delta\ell_-}{b_-}=\kappa\tau B_d,\qquad
 \frac{\beta}{\chi}=\kappa Q_d,
 \label{deg:degree-exponents}
\end{equation}
with \(X_d,B_d,Q_d\) positive units.  These identities provide the
uniform spectral gaps used in the two local saddle charts.

Fix \(0<\eta_\pm<1\), set
\(u_+=a(1-\eta_+)\), \(u_-=-1+\eta_-\), and on
\([u_-,u_+]\) write
\[
 h=-\delta p_d
   =\delta\mathcal R(u+1)(a-u)(c-u)>0,\qquad
 A_m={L_d\over h}.
\]
The scalar equation and its partial fractions are
\eqref{eq:comparison-middle} and \eqref{fx:partial-fractions}.
For \(J(u)=\exp\int_{u_+}^uA_m\) and \(J_-=J(u_-)\),
\begin{equation}
 J(u)=
 \left({1+u_+\over1+u}\right)^{1/\beta}
 \left({a-u_+\over a-u}\right)^{1/\chi}
 \left({c-u\over c-u_+}\right)^{1+1/\beta+1/\chi}.
 \label{deg:degree-middle-cocycle}
\end{equation}
Writing
\[
 \Lambda_\beta=
 \log{(1+u_+)(c-u_-)\over(1+u_-)(c-u_+)},
 \qquad \Lambda_\chi=\log a+O(1),
\]
we obtain
\begin{equation}
 \log J_-={\Lambda_\beta+(\beta/\chi)\Lambda_\chi\over\beta}
             +O(1),\qquad J_-\ge e^{c_0/\beta}.
 \label{deg:degree-large-cocycle}
\end{equation}
Indeed \(\Lambda_\beta\ge c_0>0\), while
\((\beta/\chi)|\log a|=O(\kappa|\log\kappa|)=o(1)\).

To measure how the middle passage concentrates near its negative end,
use the normalized action \(\rho=h(u_-)J/(J_-h)\) of
\eqref{fx:rho-gamma}, so
\(\rho(u_-)=1\).
If \(D_m=L_d-h'\), direct differentiation gives
\[
 {d\rho\over\rho}={D_m\over h}\,du,\qquad
 D_m'(u)=1-4\delta\mathcal R(u+\nu)>0.
\]
For one fixed \(u_0\in(u_-,0)\), this implies
\(\rho(u_0)\le e^{-c/\beta}\) and
\(\operatorname {Var}\rho\le2\).  Since
\[
 {h'\over h}={1\over u+1}-{1\over a-u}-{1\over c-u},
\]
the explicit expression \eqref{deg:degree-middle-cocycle} also yields
\begin{equation}
 \int_{u_-}^{u_+}\rho
 \left(|E_d|+\left|{h'\over h}\right|\right)du=o(1).
 \label{deg:degree-small-mass}
\end{equation}
The only noncompact term is bounded by
\(C\beta+C(1+|\log a|)e^{-c/\beta}/a\).  Thus no quotient
\(\nu/\kappa\) has entered the estimate.

We next place this middle transition inside a single interval of
trajectories joining the fixed physical sections.  Put
\begin{equation}
 s_0=\kappa\sqrt B,\qquad
 x_*={a\over s_0}={A_d\over\sqrt B},\qquad
 u=s_0\xi,\qquad w=s_0^2\tau y.
 \label{deg:degree-parabolic-dictionary}
\end{equation}
With
\[
 \Lambda_1=\frac{\omega}{\sqrt B},\qquad
 \Lambda_2=\kappa^2\mathcal R\sqrt B,\qquad
 \Lambda_3=\kappa\mathcal R\nu,
\]
the exact axis polynomial is
\begin{equation}
 P_s(\xi)=(1+s_0\xi)
       \left(1-{\xi\over x_*}\right)
       \left(1-{s_0\xi\over c}\right),
 \label{deg:degree-source-polynomial}
\end{equation}
and the full source system has the form
\begin{equation}
 \dot\xi=-\tau\{P_s+Q_\tau y\},\qquad
 \dot y=-y\{\xi B_\tau+C\tau^3\Lambda_2
                   (\Lambda_3+\Lambda_2\xi)y\},
 \label{deg:degree-source-field}
\end{equation}
where
\[
 B_\tau=1+\tau(2\Lambda_3+\Lambda_2\xi),\qquad
 Q_\tau=1-A\tau(\Lambda_3+\Lambda_2\xi)
          +C\tau^2(\Lambda_3+\Lambda_2\xi)^2.
\]

We construct two boundaries of the relevant through interval on the
same fixed physical sections.  For the upper comparison, prescribe the
actual point \(h_L=0\) of \eqref{bd:incoming-phase}, with
\(\eta=\tau y\).  Put
\(b_\tau=-\log(2\tau)\), \(q_\tau=b_\tau^{-1/2}\).
Freezing the source height scale changes the active quartet to
\begin{equation}
 \left(-q_\tau^2,\sqrt\tau\Lambda_1q_\tau,
       \tau\Lambda_3,{\tau^{3/2}\Lambda_2\over q_\tau}\right)
       \longrightarrow0.
 \label{deg:degree-zero-source}
\end{equation}
Lemma~\ref{bd:lem-source-ranges} constructs this orbit by an actual
incoming-to-central first hit, together with a real larger source
chord, using \eqref{bd:graph} with \(P_\sigma=P_s\).
Choose the fixed \(L>4(1+\sup x_*)\) of that lemma. Its energy
comparison gives, uniformly also on the larger source chord,
\begin{equation}
 \eta(\pm L)={L^2\over2}+O_L(\tau b_\tau),\qquad
 Q_\tau\eta+\tau P_s\ge c\eta .
 \label{deg:degree-zero-cuts}
\end{equation}
The same proof gives
\(\eta\asymp\tau b_\tau+\xi^2\) and the denominator margin
throughout the central source chord.  Thus \(h_L=0\) is an interior physical orbit, with neighbouring
through trajectories on both sides.

To continue to the fixed raw sections, put \(W=w/u^2\) and
\(r=|u|\).  Direct substitution in \eqref{deg:degree-field} gives
\begin{equation}
 rW_r={WA_\pm(r,W)\over
 K_d(\pm r)W-\delta p_d(\pm r)/r^2}-2W,
 \label{deg:degree-reciprocal}
\end{equation}
where
\[
 A_+={L_d(r)+E_d(r)r^2W\over r},\qquad
 A_-=-{L_d(-r)+E_d(-r)r^2W\over r}.
\]
After choosing a fixed \(U>2\max(1,c,|\nu|)+2\), then \(L\), the
strip
\begin{equation}
 {1\over3}<W<{3\over4},\qquad
 K_d(\pm r)W-\delta p_d(\pm r)/r^2>c_1
 \label{deg:degree-parabolic-strip}
\end{equation}
is invariant from \(s_0L\) to \(U\).  Equation
\eqref{deg:degree-zero-cuts} enters its interior.  The same reciprocal
equation in raw coordinates continues it to the fixed sections and
crosses each intervening equatorial root at positive height.

The lower boundary is the broken trajectory
\begin{equation}
 W^s_+(a,0)\cup\{(u,0):-1\le u\le a\}\cup W^u_+(-1,0).
 \label{deg:degree-stopped-curve}
\end{equation}
The invariant-manifold theorem in the two fixed weighted half-boxes
constructs the upper stable and unstable branches.  Their graph
denominators have a positive margin after division by the local
transverse scale.  To attach them to the fixed sections, first leave
the positive box at \(\xi=x_*+\ell_*\), with \(\ell_*>0\) fixed.
Its height and denominator are bounded below by \(c\ell_*\).
The graph equation therefore reaches \(\xi=L\), with
\[
 \eta_+^{\rm st}(L)=(L^2-x_*^2)/2+o(1).
\]
On the negative side leave the weighted box at \(u=-R\), with
\(R=1+\rho_*\), \(\rho_*>0\) fixed.  Its outgoing graph has
\[
 w_-^{\rm st}(-R)=(R^2-1)/2+o(1),\qquad
 K_dw_-^{\rm st}-\delta p_d\ge c(R^2-1).
\]
The barriers \(2(r^2-1)/5\) and \(3(r^2-1)/5\) continue this
graph until it enters \eqref{deg:degree-parabolic-strip}.  Both
branches then reach \(u=\pm U\), and afterwards the fixed raw
sections.  The source cuts \(u=\pm s_0L\) occur before all three equatorial
roots have been passed. The intermediate fixed cuts complete that
continuation before the raw tails are attached.

Between the saddle boxes \(p_d<0\), so
\(\dot u=-(h+K_dw)<0\) for every \(w\ge0\).  On either exterior
side the stopped branch has \(\dot u<0\) away from its saddle and
\(\partial_w\dot u=-K_d<0\).  Scalar order thus keeps every graph
between the stopped and upper comparison graphs in the same positive
denominator domain.  The inclination lemma in the two saddle boxes,
together with Lemma~\ref{lem:loc-through}, identifies the stable
branch as the lower endpoint of the through interval.  Write the two
actual positive-section traces as
\(h_{\rm in}^{\rm st}<h_{\rm in}^0\), where \(h_{\rm in}^0\) is the
trace of \(h_L=0\). There is no cycle at or below
\(h_{\rm in}^{\rm st}\).

The quadratic outer equation contracts towards its organizing corridor
in logarithmic radius.  More explicitly, for a fixed small raw cut
\(v_0\), its deviation at \(\pm v_0\) is bounded by
\[
 C(r_sU/v_0)^\alpha+Cr_s/v_0
       +C(|A-1|+|C-1|),\qquad 1<\alpha<2.
\]
This follows from \(\mu_i=O(r_s^{4-i})\), \(i=1,2,3\), and
variation of constants after \(u=\pm U\), retaining the passive
coefficient differences from the organizing equation. The fixed regular maps
then put both boundary traces, and by scalar order the whole interval
\((h_{\rm in}^{\rm st},h_{\rm in}^0)\), in the prescribed raw section
windows.
This scalar order locates the orbits in the raw windows. We estimate
the derivatives separately on the enlarged chords transported into
those windows.

The two fixed comparisons of Proposition~\ref{bd:prop-bounded-action}
supply the overlapping source interval $-\delta_2\le h_L\le0$ and
bounded strong-footpoint annulus. Their enlarged independent windows
reach the same raw sections by \eqref{bd:parabolic-cuts} and the outer
bridges above.

Let \(r_b\) be the negative strong footpoint where that section is
met.  On the lower part of the through interval let \(r_\ell\) be
the first point on the actual curve \(y/P_s=\ell_0\).
Choose a fixed \(r_c>2\).  A fixed real comparison through orbit
has a positive ratio at \(-r_c\); choose \(\ell_0\) below that
ratio.  Scalar order gives a unique lower interval on which the
initial ratio is smaller than \(\ell_0\).  Before its first hit,
directly from \eqref{deg:degree-source-field},
\begin{equation}
 {d\over dr}\log{y(-r)\over P_s(-r)}
 \ge {c\over\tau(1-s_0r)}
 \label{deg:degree-ratio-monotonicity}
\end{equation}
on the full interval \(r_c\le r<1/s_0\).  A strict through graph
has positive height at the negative root while \(P_s(-r)\to0\).
Hence the hit exists.  Its uniqueness and the positive first
variation at a fixed section give
\begin{equation}\label{deg:degree-phase-order}
 \partial_{h_{\rm in}}r_\ell<0,\qquad
 \partial_{h_{\rm in}}r_b<0,\qquad
 \frac{dr_\ell}{dr_b}>0
\end{equation}
on the common domain, with \(h_{\rm in}\) the actual raw input height.
No equality between the two footpoint coordinates has been used.

We need a coefficient-independent threshold at which a first ratio
hit has sufficiently large action.  Fix a positive weak-section input \(P_0\) in
the positive saddle box.  Its local Dulac formula and
\(\chi\asymp\tau\) give
\(-\tau\log z_+\le C(P_0)\) for inputs \(P\ge P_0\).
The exact core equation consequently gives
\(\log(y/P_s)(-r_c)\ge-C(P_0)/\tau\), unless the first hit
has already occurred.  If there were no hit before \(R\),
\eqref{deg:degree-ratio-monotonicity} would imply
\[
 \log(y/P_s)(-R)\ge
             \{-C(P_0)+c(R-r_c)\}/\tau.
\]
Choosing \(R=R(P_0)\) large is a contradiction.  Thus
\(r_\ell\to\infty\) forces \(P\to0\) uniformly.  Choose
\(P_0\) so that its positive action exceeds the fixed middle
variational exponent by two; then fix
\(r_{\rm ge}^{\rm out}>R(P_0)\) and
\(r_{\rm ge}>r_{\rm ge}^{\rm out}\).  The growing-action estimates below therefore start with a strict lower
action margin.

At \(r_\ell=r_{\rm ge}\), the actual strong projection has a
compact image.  Choose an independent \(R_{\rm BD}\) so that
this image lies strictly below \(2R_{\rm BD}\).  The transverse
inverse and \eqref{deg:degree-phase-order} then give a fixed
\(r_{\rm ge}^{\sharp}>r_{\rm ge}\) whose entire closed interval
\([r_{\rm ge},r_{\rm ge}^{\sharp}]\) lies inside the enlarged
annular resection.  In particular,
\[
 r_b>2R_{\rm BD}\ \Longrightarrow\ r_\ell>r_{\rm ge}.
\]
Only now decrease the upper \(s_0\)-bound so that this whole fixed
overlap lies in the growing-action product range.  Fix \(V_1>1\) and
\(0<\sigma_1<\sigma_2<\sigma_0<1/4\), with
strictly enlarged ranges at every endpoint.
The four passage regions will be restrictions of this ordered through
interval.

Choose one coefficient-only real action offset \(A_{c,0}(\theta)\),
held fixed by every phase derivative, and put
\[
 \mathcal P(r)=P_s(-r)/B_\tau(-r),\qquad
 \mathcal D=(\mathcal P(r)/r)\partial_r,\qquad
 \mathcal A_-(r)=\int_{r_c}^r\frac{vB_\tau(-v)}{P_s(-v)}\,dv-A_{c,0}.
\]
Lemma~\ref{deg:lem-growing-matching-direct} constructs the two
centre maps independently, retaining the exact initial polynomial
and both core multipliers. Its product rows
\eqref{deg:source-growing-product-phase-rows} and matched identity
\eqref{deg:degree-growing-matching} give the actual graph, with
\(|\zeta|+|X\partial_X\zeta|\le C\tau\) on its enlarged image.

Construct one negative section map for both the growing and fixed-exit
ranges. Fix \(u_c\in(-1,u_-)\) with \(|u_c|>2\sigma_0\), retain
its enlarged product buffer \(\sigma_0\le\bar\sigma<|u_c|/2\), and choose
\(\Delta>0\) below one third of the minimum of the gap
\eqref{fx:overlap-gap} on \([\sigma_1,\sigma_2]\), also with
\(\Delta<\frac14\inf_c\Gamma_{-,c}\). Enlarge the intervals of
\eqref{fx:W-limit} once more, with uniform strict margins:
\[
 \mathcal J_c\Subset\mathcal J_c^\sharp
 \Subset\mathcal J_c^{\sharp\sharp}\Subset(0,u_c^2/2).
\]
Apply Lemma~\ref{deg:lem-growing-landing-direct} on the original
enlarged lip range, with its fixed lower threshold and product buffer.
Equations \eqref{deg:landing-direct-value} and
\eqref{deg:landing-direct-gamma-unit} give the actual landing value,
its strict derivative and its positive first-variation unit.
The action supremum in \eqref{deg:landing-direct-first-row} is
$O(\tau(1+|\log\tau|))$, uniformly in independent $\kappa,j$.
Since \(s_0r\le\sigma_0\), these heights lie strictly inside
\[
 \mathcal J_{\rm gr}=[u_c^2/4,3u_c^2/4]
 \Subset\mathcal J_{\rm gr}^{\sharp}=(u_c^2/8,7u_c^2/8).
\]
Choose a fixed compact positive-height interval \(\mathcal I_{\rm out}\)
containing the closures of \(\mathcal J_{\rm gr}^{\sharp}\) and
\(\mathcal J_c^{\sharp\sharp}\) with a strict buffer, uniformly in
\(c\). Thus it contains \(u_c^2/2\), while the fixed-exit chord
stays strictly below that height.

For every \(W\in\mathcal I_{\rm out}\), the limiting graph
\(w(u)=W+(u^2-u_c^2)/2\) has a uniform positive height from
\(u_c\) to a fixed large negative cut \(-U_-\). Ordinary ODE
dependence gives a positive exact denominator, including at \(u=-1\),
and uniform \(C^5\) bounds for this finite passage and its inverse.
Choose \(U_-\) to enter \eqref{deg:degree-parabolic-strip}.
The tail estimates
\eqref{ann:tail-first-variation}--\eqref{eq:N-parameter-bound},
with the exact first cocycle extracted through three variations, give
\[
 \begin{gathered}
 z_-(W)=z_-^{\rm st}+\Omega Z_-(W),\qquad
 \Omega=r_s^{\,2-\lambda_\gamma},\quad
 \lambda_\gamma=\frac{12\gamma_s}{1+2\gamma_s},\\
 Z_-'>c,\qquad \|Z_-\|_{C^3}+\|Z_-^{-1}\|_{C^3}\le C.
 \end{gathered}
\]
Since \(\Omega\to0\) and the stopped trace lies inside its fixed raw
window, the whole chord is in the domain of \(d_\theta\) after shrinking
the coefficients. Compose with that map, put \(t_b=d_\theta(z_-^{\rm st})\),
and fix an interior reference \(W_\circ\). With
\(U_{\rm out}=\Omega^{-1}(d_\theta\circ z_-)'(W_\circ)\),
this defines a single actual map
\begin{equation}\label{deg:fixed-exit-actual-output}
 \begin{gathered}
 \Xi(W)=t_b+\Omega U_{\rm out}D^n(W),\qquad U_{\rm out}^{\pm1}\le C,\\
 (D^n)'>c,\qquad \|D^n\|_{C^3}+\|(D^n)^{-1}\|_{C^3}\le C.
 \end{gathered}
\end{equation}
For compact \(\nu\), \(|\lambda_\gamma\log r_s|\to0\), so
\(\Omega\asymp r_s^2\); the exact power is retained for the remote-root
face. In particular \(\Xi'\asymp\Omega\) and
\(|\Xi''/\Xi'|\le C\) on the entire common chord.

Fix \(r_\circ\) inside the enlarged growing domain and define
\[
 \begin{gathered}
 a_{\rm grow}=\Omega s_0^2U_{\rm out},\qquad
 t_\dagger=\Xi(W_\theta(r_\circ))+a_{\rm grow}r_\circ^2/2,\\
 T(r)=\Xi(W_\theta(r))=t_\dagger-a_{\rm grow}\mathcal B(r).
 \end{gathered}
\]
The fixed-endpoint variation in
Lemma~\ref{deg:lem-growing-landing-direct} gives
\begin{equation}\label{deg:degree-growing-output}
 \mathcal B'(r)=rG(r),\qquad
 G(r)=\Gamma_\theta(r)
              \frac{\Xi'(W_\theta(r))}{\Omega U_{\rm out}}>0.
\end{equation}
Its exact chain rule is
\begin{equation}\label{deg:source-growing-output-log-row}
 \mathcal D\log G=
 \mathcal D\log\Gamma_\theta
 -s_0^2\mathcal P(r)\Gamma_\theta(r)
                  \frac{\Xi''(W_\theta(r))}{\Xi'(W_\theta(r))}.
\end{equation}
Estimates~\eqref{deg:landing-direct-raw-first-row} and
\eqref{deg:source-growing-actual-quotient-estimates} control the two
terms that enter the complete displacement after finite closing.

On the positive side use the fixed local cuts and stopped outer chord
of Lemma~\ref{lem:actual-dulac-matching}. Its backward-integral and tail
construction holds through five phase derivatives. Thus the natural
raw input is
\begin{equation}\label{deg:fixed-exit-natural-input}
 P=\frac{x_{\rm raw}-x_b}{s_{\rm in}},\qquad
 s_{\rm in}=\Omega a^{2-\lambda_\gamma}U_{\rm in},\qquad
 U_{\rm in}^{\pm1}\le C,
\end{equation}
and the normalized map from the fixed outer chord and its inverse
have uniform \(C^5\) bounds. In the positive saddle coordinates
\(x=(u-a)/a,y=w/a^2\), write the invariant restrictions on
\(x=-\eta_+\) as
\[
 A_+(y)=-\Xi_{+,m}(y),\quad A_0=A_+(0)>0,\qquad
 Z_{+,m}(y)=z_+y g_+(y),\quad g_+(0)=1.
\]
Here \(A_0,z_+\) are positive units and \(A_+,g_+\) have uniform
finite derivative bounds. On the outer chord the invariant coordinate
constructed in that lemma is
\(\mathscr L_\theta(P)=-I(P)/A_0=\Phi_+(P)/A_0\); hence
\[
 \mathscr L_\theta(0)=0,\qquad \mathscr L_\theta'>c>0,\qquad
 \|\mathscr L_\theta\|_{C^5}
       +\|\mathscr L_\theta^{-1}\|_{C^5}\le C.
\]
Choose \(p_0>0\) small, independently of the coefficients, so that
\(\mathscr L_\theta(P)\le\vartheta<1\) on \([0,p_0]\).
Equality of the two invariants is exactly
\begin{equation}\label{deg:fixed-exit-middle-invariant}
 \mathscr L_\theta(P)=\frac{A_+(y)}{A_0}
                         \{z_+y g_+(y)\}^{\chi}.
\end{equation}
This equation has a unique small nonnegative solution \(y=y(P)\).
Indeed, on \(0<y\le\chi^2\) the logarithmic derivative of its
right side, multiplied by \(y\), is
\[
 B_\chi(y)=\chi+y\frac{A_+'(y)}{A_+(y)}
                   +\chi y\frac{g_+'(y)}{g_+(y)}
          =\chi+O(y)\ge\chi/2.
\]
The right side of \eqref{deg:fixed-exit-middle-invariant} tends
uniformly to one at \(y=\chi^2\), while its limit at zero is zero.
Thus it covers \([0,\vartheta]\) for
all sufficiently small \(\chi\). Taking logarithms gives
\[
 y=z_+^{-1}\mathscr L_\theta(P)^{1/\chi}
       \exp O(y/\chi+y).
\]
In particular the exponential factor is a positive unit. From
\(dy/d\mathscr L=y/(\mathscr L B_\chi(y))\), differentiation
twice more and the bounded derivatives of \(A_+,g_+\) give
\begin{equation}\label{deg:fixed-exit-flat-height}
 |\partial_P^j y(P)|\le
 C_j\chi^{-j}\mathscr L_\theta(P)^{1/\chi-j},
 \qquad 0\le j\le3,\quad P>0.
\end{equation}
For \(\chi<1/4\) these derivatives have zero limits at \(P=0\).
Their suprema on the fixed interval tend to zero, since
\(\chi^{-j}\vartheta^{1/\chi-j}\to0\).

Define, with its continuous value at zero,
\begin{equation}\label{deg:fixed-exit-exact-weak-coordinate}
 \Phi_\theta(P)=\{z_+y(P)\}^{\chi}
 =\mathscr L_\theta(P)\frac{A_0}{A_+(y(P))}
                          g_+(y(P))^{-\chi}.
\end{equation}
The second expression and \eqref{deg:fixed-exit-flat-height} show
that \(\Phi_\theta-\mathscr L_\theta\to0\) in \(C^3\).
Consequently \(\Phi_\theta'>c>0\), its \(C^3\) norm and inverse
norm are uniform, and \(\Phi_\theta(P)=PU_\theta(P)\) with
\(U_\theta\) a positive unit on \([0,p_0]\).
Only phase derivatives are used here.

Use \(\lambda=\ell_0\) in the normalized middle action
\(w=\ell_0h\rho V\) of \eqref{eq:comparison-action}.
The exact reference value follows directly from that raw height. Put
\[
 \varepsilon_+=\frac{\lambda h_-}{a^2J_-},\qquad
 \Lambda_+^m=(z_+\beta\varepsilon_+)^\chi,
 \qquad
 \Phi_\theta(a_{\rm in})=\Lambda_+^m.
\]
Write \(q_s=\beta/\chi\), \(j=1/c\), \(d=\mathcal Rc\), and
\(\mathscr R=(1-ju_-)/(1-ju_+)\). The logarithms in
\eqref{deg:degree-middle-cocycle} are exactly
\[
 \Lambda_\beta=\log\frac{(1+u_+)\mathscr R}{1+u_-},\qquad
 \Lambda_\chi=\log\frac{(a-u_+)\mathscr R}{a-u_-}.
\]
The reference telescope is
\begin{equation}\label{p2:reference-telescope}
 \begin{gathered}
 \beta\log(z_+\beta\varepsilon_+)
 =-\Lambda_\beta-q_s\Lambda_\chi+\beta\log(\tau^2U_{\rm tel}),\\
 U_{\rm tel}=\frac{z_+\ell_0\beta\delta d}{a^2\tau^2}
                (1+u_-)(a-u_-)(1-ju_+),\qquad U_{\rm tel}^{\pm1}\le C.
 \end{gathered}
\end{equation}
Here \(\delta=\kappa\tau\), \(a=\kappa A_d\), and
\(\beta=\kappa\tau B_d\); hence all factors of \(\kappa\)
cancel in \(U_{\rm tel}\). Since
\(\Lambda_\beta=\Gamma_{-,c}+O(a)\) and
\(\Lambda_\chi=\log a+O(1)\), the error is bounded by
\(C\{a+q_s(1+|\log a|)+\beta(1+|\log\tau|)\}\).
It tends to zero under independent coefficient bounds, giving
\begin{equation}\label{deg:degree-fixed-exit-gap}
 q_s\log\Phi_\theta(a_{\rm in})
 =\beta\log(z_+\beta\varepsilon_+)
 =-\Gamma_{-,c}+o(1).
\end{equation}
Thus \(\Lambda_+^m\to0\) and the reference
\(a_{\rm in}\in(0,p_0)\) exists uniquely. It corresponds to
\(y=\beta\varepsilon_+\), hence to the raw middle input \(V=\beta\).
For every orbit in this local input interval, the exact height dictionary is
\begin{equation}\label{fx:height-dictionary}
 \begin{gathered}
 y=\varepsilon_+V,\qquad
 V=\beta\left(\frac{\Phi_\theta(P)}
                    {\Phi_\theta(a_{\rm in})}\right)^{1/\chi},\\
 \frac{w(u_+)}{h(u_+)}=\lambda\rho(u_+)V.
 \end{gathered}
\end{equation}

We have constructed the through interval, its first ratio hits and its
input, output and centre-height maps. Proposition~\ref{p2:prop-unified-carrier-cover}
uses them to cover every input, including truncated intervals and
boundaries between the four ranges. Proposition~\ref{p2:prop-finite-v-fe-actual-budget}
constructs the last two transitions on the same shooting chord and
buffered closing image.
\end{proof}

The remote positive root requires a different coordinate at infinity.
Its reciprocal keeps the finite passage coefficients regular while the
physical tail scales are retained exactly.

\begin{proposition}[The remote-root face]\label{deg:prop-degree-infinity}
Let \(d,A_d\) range in positive compact intervals and put
\begin{equation}
 \begin{gathered}
 0\le j\le j_0,\qquad a=\kappa A_d,\qquad
 \mathcal R=dj,\qquad B=dA_d,\\
 \omega=d(1-a+ja),\qquad
 j\nu={j-ja-1\over3},\qquad \delta=\kappa\tau .
 \end{gathered}
 \label{deg:degree-infinity-graph}
\end{equation}
For \(j>0\) this is \eqref{deg:degree-graph} with \(c=j^{-1}\).
For \(j>0\), the conclusion of Proposition~\ref{deg:prop-degree-drop}
extends uniformly as \(j\downarrow0\). At \(j=0\) this denotes only the
one-sided closure of the analytic estimates, not an additional physical
coefficient fibre or an asserted inverse raw gate.
\end{proposition}

\begin{proof}
The factorization and translated root are
\begin{equation}
 p_d(u)=d(u+1)(u-a)(1-ju),\qquad
 \gamma_s={\kappa\tau d\over3}(j-ja-1).
 \label{deg:degree-infinity-factor}
\end{equation}
All coefficients of the exact field on bounded \(u\)-intervals are
analytic at \(j=0\).  In the middle calculation the apparently singular
quantities reduce to
\[
 {1-ju\over1+j},\qquad {j\over1-ju},\qquad
 D_m'=1-4\gamma_s-4\kappa\tau dju.
\]
The limiting functions needed in the accumulated-action and fixed-exit
arguments are
\[
 H_{0,j}(u)=-{u(1+u)(1-ju)\over1+j},\qquad
 \Gamma_j(u)=
 \log{(1+u)(1-ju_-)\over(1+u_-)(1-ju)},
\]
and
\[
 \Gamma_{-,j}-\Gamma_j(-\sigma)
 =\log{1+j\sigma\over1-\sigma}>0.
\]
Thus the small-mass, Lambert, and fixed-exit estimates above extend
continuously to \(j=0\).

The bounded-action domains also have regular coordinates.  Put
\(\sigma=s_0=\kappa\sqrt B\), \(x_*=A_d/\sqrt B\).  Then
\[
 P_\sigma(\xi)=(1+\sigma\xi)(1-\xi/x_*)(1-\sigma j\xi),\qquad
 \Lambda_2=\sigma^2j/x_*,\quad
 3\Lambda_3=\sigma(j-1)/x_*-\sigma^2j<0.
\]
The positive root remains simple and the negative compact-output
interval has a positive polynomial margin.  The two fixed phase
comparisons, the compact central-height estimate, and the negative
frozen first-exit construction in
Lemma~\ref{bd:lem-source-ranges} are consequently uniform on the
compact \(j\)-block.  Lemma~\ref{bd:lem-annulus} uses the sign
\(\Lambda_3<0\) on this face: its actual first tail cocycle is only
bounded above by \(Ce^2\), with \(e=\tau\Lambda_2\), and is
retained exactly in the relative normalization.  Thus it requires
no bound on \(\tau|\Lambda_3||\log e|\), even when \(j\) is
arbitrarily small.  The bounded-action estimate is therefore uniform even as the third root
escapes the compact root chart.

It remains to check the unbounded tails.  Their exact source coefficients
are
\[
\begin{array}{c|cccc}
 &a_\pm&p_\pm&\gamma_\pm&s_\pm\\ \hline
 +&-\tau d/A_d&\tau d(1-a+ja)/A_d&\gamma_s&r_sa\\
 -&-\kappa^2\tau dA_d&\kappa\tau d(1-a+ja)&\gamma_s&r_s,
\end{array}
\qquad r_s=\kappa\tau dj.
\]
After the transverse blow-down, the two maps to the fixed raw sections
have first-cocycle scales
\begin{equation}
 r_s^{\,2-\lambda_\gamma}U_-,\qquad
 r_s^{\,2-\lambda_\gamma}a^{\,2-\lambda_\gamma}U_+,
 \qquad \lambda_\gamma={12\gamma_s\over1+2\gamma_s},
 \label{deg:degree-common-scale}
\end{equation}
with normalized phase derivatives bounded.  The common factor
\(r_s^{2-\lambda_\gamma}\) cancels in the complete transition, while
\(|\lambda_\gamma\log a|=o(1)\). With \(q_s=\beta/\chi\),
the exact telescope \eqref{p2:reference-telescope} gives
\[
 q_s\log\Lambda_+^{m}=-\Lambda_\beta+o(1),\qquad
 e^{-C/q_s}\le\Lambda_+^{m}\le e^{-c/q_s},
\]
and contains no term \(\beta\log j\). The common section scales are
\(s_{\rm in}=\Omega a^{2-\lambda_\gamma}U_{\rm in}\) and
\(a_{\rm out}=\Omega U_{\rm out}\), with
\(\Omega=r_s^{2-\lambda_\gamma}\).

The local boxes and normalized outer maps have uniform finite phase
bounds on the compact $j$-block, so the weak-coordinate construction
\eqref{deg:fixed-exit-middle-invariant}--\eqref{deg:fixed-exit-flat-height}
has exact reference
\[
 \Phi_\theta(a_{\rm in})=(z_+\beta\varepsilon_+)^\chi,
 \qquad q_s\log\Phi_\theta(a_{\rm in})=-\Gamma_{-,j}+o(1).
\]
Proposition~\ref{p2:prop-finite-v-fe-actual-budget} uses these regular
functions \(H_{0,j},\Gamma_j\), their strict action gap and the exact
raw scales to extend the last two passage estimates to this face.

For the outer geometry no fixed cut beyond \(|\nu|\) is available.
Choose instead the positive source cut \(\xi=x_*L_t\), so that
\(u/a=L_t\) exactly.  The actual source construction gives
\(1/3<\eta/\xi^2<3/4\) and a positive divided denominator there.
On the negative side continue from \(u=-s_0L\) only to
\(u=-L_t\).  On this interval the regular factorization gives
\[
 |p_d(-r)|\le Cr,\qquad
 |\delta p_d(-r)|/r^2\le C\tau/L,\qquad
 |K_d-1|+|-L_d(-r)/r-1|\le C\kappa\tau.
\]
The equation for \(w/u^2\) therefore has the same invariant strip
and positive denominator.  These are the two finite entry chords
for the bounded source quartets in the table.  Their normalized
tail equations \eqref{eq:parabolic-tail} continue to the fixed raw
sections; the positive tail crosses \(u=c=j^{-1}\), if that root
lies before the section, at positive height.  The stopped branches
are attached by the same weighted local graphs and finite bridges.
Since \(p_d<0\) on \(-1<u<a\), and \(\partial_w\dot u=-K_d<0\),
every intermediate graph remains on the through side of these
branches.  This constructs the entire ordered interval before its
zero bounds are applied.

The same ratio and strong-projection orders retain the bounded/growing
overlap. The positive $\Gamma_j$ gap retains the fixed-exit overlap.
On the real enlarged source chord through the actual cut \(h_L=0\),
\eqref{bd:positive-height} again gives
\(Y_0\asymp b_\tau\to\infty\), while
\(s_s^2Y_0\le C\tau^3\kappa^4j^2b_\tau\to0\).
Hence this chord supplies the overlap of the actual high interval
with the low carrier through $h_L=0$.
All four domains and estimates are uniform for $j>0$;
$j=0$ retains only the normalized one-sided meaning.
\end{proof}

\subsection{Direct negative landing and its first action variation}

A negative ratio hit must be continued to a fixed section before it can
be used in the physical return. The next lemma constructs this landing
and estimates its first action derivative. The terminal section is kept
fixed while the starting lip moves.

\begin{lemma}\label{deg:lem-growing-landing-direct}
Use the exact field \eqref{deg:degree-field} on either the compact
coefficient graph \eqref{deg:degree-graph} or the reciprocal graph
\eqref{deg:degree-infinity-graph}, with physical \(j>0\).
Keep the original sections, compact passive coefficients, and
coefficient-only scales. Put
\[
 \epsilon_d=\kappa\tau,\qquad s_0=\kappa\sqrt B,\qquad
 x_c=-u_c,\qquad \sigma=s_0r .
\]
Choose the original enlarged product buffer
\(0<\sigma_0\le\bar\sigma<x_c/2<1/2\), and fix
\(R_0>r_c\). On the original enlarged range
\begin{equation}\label{deg:landing-direct-domain}
 R\le r\le\bar\sigma/s_0,\qquad R\ge R_0,\qquad
 q_\ell=\frac{\epsilon_d}{\sigma}=\frac{\tau}{\sqrt B\,r},
 \qquad L_\ell=1+|\log q_\ell|,
\end{equation}
the graph from the original lip \(w=\ell_0h\) reaches \(u_c\)
in positive time, and its inverse reaches the original centre
\(u=0\). Both graphs have positive height and denominator.

Write \(W_\theta(r)=w(u_c;r)\). There is an exact positive unit
\(\Gamma_\theta(r)\) such that
\begin{equation}\label{deg:landing-direct-value}
 \left|W_\theta(r)-\frac{x_c^2-\sigma^2}{2}\right|
 \le C\epsilon_d\{1+\sigma L_\ell\},\qquad
 W_\theta'(r)=-s_0^2r\Gamma_\theta(r),
\end{equation}
and
\begin{equation}\label{deg:landing-direct-gamma-unit}
 c\le\Gamma_\theta\le C,\qquad
 |\Gamma_\theta-1|\le C(q_\ell L_\ell+\epsilon_d).
\end{equation}
If \(\Gamma(\sigma)=\Gamma_\theta(\sigma/s_0)\), then, with all
coefficients frozen,
\begin{equation}\label{deg:landing-direct-first-row}
 \begin{aligned}
 |\partial_\sigma\log\Gamma|
 &\le C\left\{\frac{q_\ell}{\sigma}L_\ell+\epsilon_d\right\},\\
 |\mathcal D\log\Gamma_\theta|
 &\le C\left\{\frac{q_\ell}{r}L_\ell+s_0\epsilon_d\right\},
 \qquad \mathcal D=\frac{\mathcal P(r)}r\partial_r .
 \end{aligned}
\end{equation}
Constants are fixed by the original compact ranges, \(R_0,r_c,\ell_0\)
and the section buffers, independently of the later \(R,\kappa,\tau,j\).
In particular the supremum of the action row is
\(O(\tau(1+|\log\tau|))\) as \(\tau\to0\), uniformly in independent
\(\kappa,j\). At fixed \(r,\tau>0\), sending only \(\kappa\to0\)
does not assert that this row tends to zero or that \(\Gamma_\theta\to1\).
\end{lemma}

\begin{proof}
Put \(x=-u\), \(H(x)=h(-x)\), \(k(x)=K(-x)\),
\(b(x)=-L(-x)\), and \(e(x)=-E(-x)\). The exact graph is
\begin{equation}\label{deg:landing-direct-graph}
 w_x=F(x,w)=\frac{w(b+ew)}{H+kw},\qquad
 w(\sigma;\sigma)=\ell_0H(\sigma).
\end{equation}
Indeed the original dictionary gives
\(s_0^2\tau P_s(-r)=h(-s_0r)\) exactly.
In the compact and reciprocal graphs, respectively,
\[
 H=\epsilon_d\mathcal R(1-x)(a+x)(c+x),\qquad
 H=\epsilon_dd(1-x)(a+x)(1+jx).
\]
On \([\sigma,x_c]\), \(a/\sigma\le C/R_0\).
The exact coefficients \eqref{fx:coefficients} therefore give
\begin{equation}\label{deg:landing-direct-coefficients}
 \begin{gathered}
 H\asymp\epsilon_dx,\quad |H'|+|H''|\le C\epsilon_d,\quad
 |H'/H|\le C/x,\qquad k\asymp1,\\
 |k-1|+|k'|\le C\epsilon_d,\qquad |b-x|\le C\epsilon_dx,\\
 |(b/x)'|\le C\epsilon_d,\quad
 |b'/b|\le C/x,\qquad |e|+|e'|\le C\epsilon_d^2 .
 \end{gathered}
\end{equation}
At the reciprocal face these estimates use
\(\epsilon_d\mathcal R=\epsilon_ddj\) and
\(\epsilon_d\mathcal R\nu=\gamma_s\); no bound on \(\nu\) or
\(\log j\) is required. A common small upper bound on \(\tau\)
makes \(q_\ell\) small without comparing any radial parameters.

First admit the graph. On a provisional band \(0<w<Mx^2\),
the denominator and \(b+ew\) are positive, and
\[
 0<F\le Cx+C\epsilon_d^2w,\qquad
 F\ge\frac{cxw}{C\epsilon_dx+Cw}.
\]
A fixed large \(M\) prevents first contact with \(Mx^2\);
the lip is below it since \(H(\sigma)/\sigma^2=O(q_\ell)\).
Thus the graph stays positive and reaches \(x_c\).
On \([\sigma,2\sigma]\), height is increasing and
\[
 \frac{dx}{dw}\le C\frac{\epsilon_d\sigma+w}{\sigma w}.
\]
For fixed small \(m>0\), the time to height \(m\sigma^2\) is at
most \(C\epsilon_dL_\ell+Cm\sigma<\sigma\), after a common
coefficient shrink. A smaller quadratic lower barrier then gives
\begin{equation}\label{deg:landing-direct-corridor}
 0<w\le Cx^2,\quad b+ew\asymp x,\qquad
 w\asymp x^2\quad(2\sigma\le x\le x_c).
\end{equation}

For the inverse to the centre, the same exact coefficients give,
on \([s_0r_c,\sigma]\) while \(w/H\le\ell_0\),
\[
 \frac{d}{dx}\log(w/H)
 =\frac{b+ew}{H+kw}-\frac{H'}H
 \ge c/\epsilon_d-C/x>0.
\]
Here \(\epsilon_d/x\le C\tau/r_c\). Backward continuation remains
below this lip, whose hit at \(\sigma\) is therefore the first hit
after the inner cut. At that cut its height \(w_b\) satisfies
\(w_b\le C\epsilon_d(a+s_0r_c)\le C\kappa^2\tau\).
On \([0,s_0r_c]\), \(H\ge c\epsilon_da>0\), and for
\(0<w\le2w_b\) the exact derivative satisfies
\[
 F\ge-C\frac{\epsilon_d^2}{\epsilon_da}w^2\ge-C\tau w^2.
\]
In the inverse direction this prevents first contact with \(2w_b\).
Positive denominator and scalar uniqueness give a positive centre
height. This argument permits either sign of \(e(0)w\).

Write \(D=H+kw\) and
\[
 N=\frac{H'+(k'+e)w}{D},\qquad
 I_x=\int_\sigma^xN(t,w(t;\sigma))\,dt,\qquad I=I_{x_c}.
\]
Direct differentiation along the graph gives
\(F_w=(\log(w/D))_x+N\).
The near height change and the quadratic corridor imply
\[
 \int_\sigma^{2\sigma}\frac{\epsilon_d}{D}\,dx
 \le Cq_\ell L_\ell,\qquad
 \int_{2\sigma}^{x_c}\frac{\epsilon_d}{D}\,dx\le Cq_\ell.
\]
The remaining part of \(N\) has integral \(O(\epsilon_d)\), so
\(|I_x|\le C(q_\ell L_\ell+\epsilon_d)\).
Keep the complete moving initial factor
\begin{equation}\label{deg:landing-direct-initial-factor}
 C_\sigma=(1+\ell_0k(\sigma))H'(\sigma)
                 -b(\sigma)-\ell_0e(\sigma)H(\sigma).
\end{equation}
The moving initial point contributes the intrinsic derivative
\(\ell_0H'(\sigma)-F(\sigma,\ell_0H(\sigma))\). Combining this
with the Jacobian to the fixed terminal section gives exactly
\begin{equation}\label{deg:landing-direct-exact-variation}
 w_\sigma(x;\sigma)=\frac{w(x;\sigma)}{D(x;\sigma)}C_\sigma e^{I_x},
 \qquad
 \Gamma(\sigma)=\frac{w(x_c;\sigma)}{D(x_c;\sigma)}
                       \frac{-C_\sigma}{\sigma}e^I .
\end{equation}
Here \(-C_\sigma/\sigma=1+O(q_\ell+\epsilon_d)\), and the terminal
ratio is \(1+O(\epsilon_d)\). These identities prove the strict sign,
the unit estimate, and \(|w_\sigma(x;\sigma)|\le C\sigma\).
The value estimate follows by integrating the exact residual
\[
 F-x=\frac{w(b-xk)+ew^2-xH}{D}.
\]
The first two terms cost \(O(\epsilon_d)\); the near \(xH/D\) term
costs \(O(\epsilon_d\sigma L_\ell)\), and the far term
\(O(\epsilon_d)\). Include the original initial value
\(\ell_0H(\sigma)\) to obtain \eqref{deg:landing-direct-value}.

Apply Lemma~\ref{deg:lem-spectral-exterior} to this same intrinsic
landing graph, with starting abscissa \(\sigma\), terminal cut \(x_c\),
and identity observer. Its constant coefficients are
\[
 t=-r_s,\qquad b_{\rm k}=2\gamma_s,\qquad v_{\rm k}=\gamma_s,
 \qquad
 (\mu_2,\mu_1,\mu_0)
   =(3\gamma_s,\epsilon_d\omega,\epsilon_d\kappa B),
 \qquad r_s=\epsilon_d\mathcal R .
\]
Here \(b_{\rm k},v_{\rm k}\) denote the lemma's constant coefficients.
The two exact coefficient graphs give
\[
 H(x)=-r_sx^3+3\gamma_sx^2+\epsilon_d\omega x
                         +\epsilon_d\kappa B,
 \qquad b(x)=x(1+2\gamma_s-r_sx),
\]
so its \(M,Q,\zeta\) are precisely \(H,k,e\) in
\eqref{deg:landing-direct-graph}. On the original enlarged range,
\(2\sigma<x_c\), and \eqref{deg:landing-direct-coefficients} gives
\[
 H(x)/x^2\asymp q_\ell\quad(\sigma\le x\le2\sigma),\qquad
 |\gamma_s|+r_s\le C\epsilon_d,\qquad
 \kappa/\sigma\le C/R_0 .
\]
Thus the lemma's mass bounds hold with \(\rho=q_\ell\), its
\(\Theta\le Cq_\ell\), and, since the observer is the identity,
\(\mathcal T\le C\epsilon_d\). Its smallness conditions follow from
one independent upper bound on \(\tau\) and the original coefficient
bounds, uniformly on the compact and reciprocal graphs. The actual
corridor and terminal height are those already constructed above.

Take the lemma's frozen quadratic-label scale equal to one, so that
\(\beta_\ell=-\sigma^2/2\). By
\eqref{deg:landing-direct-exact-variation}, the terminal slope in this
label is exactly \(\Gamma(\sigma)=-w_\sigma(x_c;\sigma)/\sigma\).
With \(E_\sigma=\sigma\partial_\sigma\),
\eqref{deg:spectral-exterior-rows} gives
\[
 \left|E_\sigma\log\Gamma+\lambda_\gamma\right|
 \le C\{q_\ell L_\ell+\epsilon_d(\sigma/x_c)^e\},\qquad
 \lambda_\gamma=\frac{12\gamma_s}{1+2\gamma_s},\quad
 e=2-\lambda_\gamma>3/2 .
\]
The spectral power \(\sigma^{-\lambda_\gamma}\) has its exact exponent
held fixed in this phase derivative. Since
\(|\lambda_\gamma|\le C\epsilon_d\), \(\sigma<1\), and \(e>1\),
division by \(\sigma\) yields
\[
 |\partial_\sigma\log\Gamma|
 \le C\{(q_\ell/\sigma)L_\ell+\epsilon_d\}.
\]
The original action identity
\(\mathcal D=s_0(\mathcal P/r)\partial_\sigma\), with
\(\mathcal P/r\asymp1\), gives the second row of
\eqref{deg:landing-direct-first-row}. For \(r\ge R_0\),
\(q_\ell=\tau/(\sqrt B\,r)\) and \(s_0\epsilon_d=O(\kappa^2\tau)\)
give its supremum \(C\tau(1+|\log\tau|)\), independently of the later
\(R\) and of the reciprocal parameter.
\end{proof}

\begin{corollary}\label{deg:cor-growing-landing-fixed-product}
On a fixed original band \(0<\sigma_1\le\sigma\le\sigma_2\le\bar\sigma\),
\begin{equation}\label{deg:landing-direct-fixed-product}
 \left\|W_\theta(\sigma/s_0)-\frac{u_c^2-\sigma^2}{2}\right\|_
                  {C^2_\sigma[\sigma_1,\sigma_2]}
 \le C\epsilon_d(1+|\log\epsilon_d|)\longrightarrow0
 \quad\text{as }\epsilon_d\to0 .
\end{equation}
If, independently, the actual scalar inverse phase
\(\Psi_\theta\to\Psi_0\) in \(C^2\) on a fixed buffered interval
containing this band's height image, then
\(\Psi_\theta(W_\theta(\sigma/s_0))\) converges in \(C^2_\sigma\)
to \(\Psi_0((u_c^2-\sigma^2)/2)\), under the hypotheses of that
independent convergence.
\end{corollary}
\begin{proof}
On the fixed band, \eqref{deg:landing-direct-gamma-unit} and
\eqref{deg:landing-direct-first-row} give
\(\Gamma-1,\Gamma_\sigma=O(\epsilon_d(1+|\log\epsilon_d|))\).
Apply these to \(W_\sigma=-\sigma\Gamma\), then use the value
row \eqref{deg:landing-direct-value}. The composition claim is the
ordinary \(C^2\) chain rule on the assumed buffered image; for the
inverse phase use Proposition~\ref{fx:prop-inverse-phase} under its
own limiting hypotheses.
\end{proof}

In coordinates scaled by the starting lip, the fixed terminal section
has a moving coordinate. Its contribution remains in the derivative: if \(x=\sigma t\), \(w=\sigma^2v\), then
\[
 \frac{d}{d\sigma}W_\theta(\sigma/s_0)
 =2\sigma v(x_c/\sigma;\sigma)
  +\sigma^2v_\sigma(x_c/\sigma;\sigma)
  -x_cv_t(x_c/\sigma;\sigma).
\]
The limiting heights lie inside the original buffered
\(\mathcal I_{\rm out}\), including \(u_c^2/2\); the value estimate
places the whole enlarged landing image there after a common shrink.
With the actual raw map \eqref{deg:fixed-exit-actual-output},
\[
 T(r)=\Xi(W_\theta(r)),\qquad
 T'(r)=-s_0^2r\Gamma_\theta(r)\Xi'(W_\theta(r)).
\]
Its raw endpoint, exact power \(\Omega=r_s^{2-\lambda_\gamma}\),
and ordinary \(C^3\) normalized observer bounds are retained.
No coefficient derivative of a moving raw endpoint is asserted.
For the positive unit \(G=\Gamma_\theta\Xi'/(\Omega U_{\rm out})\),
the original chain rule gives the explicit bound
\begin{equation}\label{deg:landing-direct-raw-first-row}
 |\mathcal D\log G|
 \le C\left\{\frac{q_\ell}{r}L_\ell+s_0\epsilon_d
                           +s_0^2\mathcal P(r)\right\}.
\end{equation}
The last term comes from the bounded logarithmic derivative
\(\Xi''/\Xi'\) of the independently constructed output map;
the exact common tail power cancels.
The error in the \(\Gamma_\theta\) derivative tends to zero with
$\tau$; the output-composition error also requires $s_0\to0$.
We now construct the matching independently of this landing estimate.

\subsection{Separate source heights in the growing range}

The positive and negative half-orbits reach the common centre with
exponentially small heights. We construct these two height maps
separately, retaining their linear multipliers and initial polynomial
factors. Equality of their logarithms then gives the matched action
derivative while allowing both centre multipliers to tend to zero.

\begin{lemma}[Actual source matching at degree loss]
\label{deg:lem-growing-matching-direct}
Fix the original compact or reciprocal coefficient ranges, sections,
passive coefficient neighbourhood and buffered positive weak chord
of \eqref{deg:fixed-exit-exact-weak-coordinate}.
For one fixed \(R_0>r_c\), the separate centre-height maps on
\(0<P<p_0\) and \(R_0\le r\le\bar\sigma/s_0\) have the
pure rows \eqref{deg:source-growing-negative-source-rows} and
\eqref{deg:source-growing-positive-source-rows}.
Their product remainder has the total-order-two rows
\eqref{deg:source-growing-product-phase-rows} before matching.
There is a fixed \(R=r_{\rm ge}^{\rm out}>R_0\), chosen before
the independent coefficient bounds, such that every negative lip
with \(R\le r\le\bar\sigma/s_0\) matches one actual input in
the original strict positive chord. On this entire enlarged graph,
\[
 -X\partial_Xr=e_s\frac{\mathcal P(r)}r(1+\zeta),\qquad
 |\zeta|+|X\partial_X\zeta|\le C\tau,
 \qquad e_s=\tau/\chi\asymp1.
\]
All phase derivatives freeze the coefficients and retain the original
centre, action constants, physical scales and raw section windows.
\end{lemma}
\begin{proof}
\emph{The two centre-height maps.}
Use the fixed coefficient field and positive weak chord, with scales
\begin{equation}
 \epsilon=\kappa\tau,\qquad s_0=\kappa\sqrt B,\qquad
 a=\kappa A_d,\qquad x_* =A_d/\sqrt B,\qquad
 r_s=\epsilon\mathcal R,\qquad \gamma_s=\epsilon\mathcal R\nu .
 \label{deg:source-growing-phase-scales}
\end{equation}
Distinguish the positive unit $B$ from $\mathscr B=B_\tau(-r)$.
For the original inner cut $r_c>2$ and fixed $R_0>r_c$, use
\begin{equation}
 R\le r\le\bar\sigma/s_0,\qquad R\ge R_0,\qquad
 0<\bar\sigma<1/2 .
 \label{deg:source-growing-separate-domains}
\end{equation}
Constants are independent of the later $R$. The positive domain is
$0\le P\le p_0$ with its original larger chord buffer.

\smallskip
\noindent\emph{The negative source equation.}

Write \(x=-u=s_0r\), \(w=s_0^2\tau y\), and
\begin{equation}
 P_-(r)=P_s(-r),\qquad K_-(r)=K(-s_0r),\qquad
 e_-(r)=-E(-s_0r).
 \end{equation}
The exact polynomial is
\begin{equation}
 P_-(r)=\frac{\mathcal R}{\sqrt B}
       (1-s_0r)(x_*+r)(c+s_0r)
 \quad\hbox{on the compact graph},
 \label{deg:source-growing-negative-polynomial-compact}
\end{equation}
and
\begin{equation}
 P_-(r)=\frac d{\sqrt B}
       (1-s_0r)(x_*+r)(1+js_0r)
 \quad\hbox{on the reciprocal graph}.
 \label{deg:source-growing-negative-polynomial-remote}
\end{equation}
The identities $\mathcal RA_dc=B$ and $B=dA_d$, respectively,
give exactly the original $P_s(-r)$.
Set
\begin{equation}
 \mathscr B(r)=1+2\gamma_s-\epsilon\mathcal R s_0r,\quad
 \mathfrak d(r)=\frac{P_-(r)}{r\mathscr B(r)},\quad
 \mathcal D=\mathfrak d(r)\partial_r,\quad
 h_1(r)=\frac{P_-'(r)}{r\mathscr B(r)} .
 \label{deg:source-growing-negative-action-operator}
\end{equation}
Use the chosen coefficient-only real action offset \(A_{c,0}(\theta)\)
and define
\begin{equation}
 \mathcal A_-(r)=\int_{r_c}^r\frac{dt}{\mathfrak d(t)}
                       -A_{c,0},\qquad \mathcal D\mathcal A_-=1.
 \label{deg:source-growing-chosen-action-offset}
\end{equation}
Direct differentiation of \eqref{deg:source-growing-negative-polynomial-compact}, \eqref{deg:source-growing-negative-polynomial-remote}, and \eqref{deg:source-growing-negative-action-operator} gives
\begin{equation}
 \begin{gathered}
 P_-\asymp1+r,\quad \mathscr B\asymp1,\quad
 \mathfrak d\asymp1,\quad
 |\mathcal D^m\mathfrak d|\le C_m,\quad
 |\mathcal D^mh_1|\le C_m/r\quad(0\le m\le3),\\
 |K_--1|\le C\epsilon,\quad
 |\mathcal D^mK_-|\le C_m\tau\quad(1\le m\le3).
 \end{gathered}
 \label{deg:source-growing-negative-coefficient-bounds}
\end{equation}
For the polynomial monotonicity used below, put \(\sigma=s_0r\).
On the compact graph,
\[
 P_-'=\frac{\mathcal R}{\sqrt B}
 \{c(1-2\sigma)+2\sigma-3\sigma^2
                         +a(1-c-2\sigma)\},
\]
whereas on the reciprocal graph
\[
 P_-'=\frac d{\sqrt B}
 \{1-2\sigma+j\sigma(2-3\sigma)
                         +a(j(1-2\sigma)-1)\}.
\]
The fixed bound \(\bar\sigma<1/2\) and a common \(a\)-shrink
give \(P_-'\ge c\) in either graph.  Since
\(\mathscr B'=-\epsilon\mathcal R s_0\le0\), it follows that
\begin{equation}\label{deg:degree-negative-polynomial}
 \mathcal P:=P_-/\mathscr B\asymp1+r,\qquad
 \mathcal P'\ge c .
\end{equation}

Indeed, $P_-'$, $P_-''$, $P_-'''$ are bounded by
$C,Cs_0,Cs_0^2$, with $s_0\le\bar\sigma/r$. Also
\begin{equation}
 K_-=1-A(\gamma_s-\epsilon\mathcal R s_0r)
           +C(\gamma_s-\epsilon\mathcal R s_0r)^2.
 \label{deg:source-growing-negative-K}
\end{equation}
Its positive-order action rows are $O(\epsilon s_0)=O(\tau)$.
On the reciprocal graph, $\mathcal R=dj$ and
$\gamma_s=\epsilon d(j-ja-1)/3$ are regular at $j=0$.

From the original rational graph equation one obtains exactly
\begin{equation}
 y_r=\frac{y\{r\mathscr B+e_-s_0\tau y\}}
                 {\tau(P_-+K_-y)}.
 \label{deg:source-growing-negative-scalar-equation}
\end{equation}
Along the inverse passage introduce accumulated action \(q\)
by \(r_q=-\tau\mathfrak d(r)\), starting at \(r\), and put
\(W=y/P_-\).  The source lip is exactly \(W(0)=\ell_0\).
Then
\begin{equation}
 (\log W)_q=
 -\frac{1+\varpi(r)W}{1+K_-(r)W}+\tau h_1(r)
 =-1+\tau h_1(r)+G(r,W),
 \label{deg:source-growing-stable-ratio-equation}
\end{equation}
where
\begin{equation}
 \varpi=e_-s_0\tau\mathfrak d,\qquad
 G(r,W)=\frac{(K_--\varpi)W}{1+K_-W}.
 \label{deg:source-growing-stable-ratio-nonlinearity}
\end{equation}
Here \(|e_-|\le C\epsilon^2\); direct differentiation gives
\begin{equation}
 |\mathcal D^m\varpi|\le C_m\tau,\qquad 0\le m\le3.
 \label{deg:source-growing-stable-coefficient-rows}
\end{equation}
These derivative bounds apply to the exact nonlinear coefficient.

\smallskip
\noindent\emph{Negative source variations.}\par\nopagebreak[4]

The endpoint at the original inner cut is
\begin{equation}
 T(r)=\frac1\tau\int_{r_c}^r\frac{dt}{\mathfrak d(t)}
     =\frac{\mathcal A_-(r)+A_{c,0}}{\tau}.
 \label{deg:source-growing-inward-terminal}
\end{equation}
It satisfies
\begin{equation}
 T\ge c(r-r_c)/\tau,\qquad
 \mathcal DT=1/\tau,\qquad
 \mathcal D^2T=\mathcal D^3T=0 .
 \label{deg:source-growing-inward-terminal-action-rows}
\end{equation}
Apply Lemma~\ref{deg:lem-stable-ratio-ledger} with
\((\kappa,H,Q,\omega,\rho)=(\tau,h_1,K_-,\varpi,\tau)\)
and density \(L=G\). Its phase derivative is the initial action
derivative \(\delta\) at fixed \(q\):
\begin{equation}\label{deg:source-growing-initial-action-flow}
 \delta f(r(q;r))=(\mathcal Df)(r(q;r)),\qquad
 \partial_q f=-\tau\mathcal Df,
\end{equation}
because \(\mathcal A_-(r(q;r))=\mathcal A_-(r)-\tau q\).
In particular \(\delta\) commutes with \(\partial_q\).
The polynomial factors remain positive on the whole route
\(r_c\le v\le r\), with \(s_0v\le\bar\sigma<1/2\), and on
a fixed enlargement at \(r_c\). Equations
\eqref{deg:source-growing-negative-coefficient-bounds} and
\eqref{deg:source-growing-stable-coefficient-rows} give the required
phase and mixed rows with \(\alpha=0\). Since \(K_-\) is close
to one and \(\varpi,\tau h_1\) are small, both
\(\mathcal F/W\) and \(\mathcal F_W\) have a fixed negative
margin on \(0\le W\le\ell_0\). The lemma therefore gives
\begin{equation}\label{deg:source-growing-stable-ratio-corridor}
 0<W(q;r)\le\ell_0e^{-cq},\qquad 0\le q\le T(r),
\end{equation}
and the fixed-clock rows
\begin{equation}
 |\delta^iG|\le C\tau(1+q)^Ne^{-cq}\quad(i=1,2),
 \qquad |G|+|G_q|\le C(1+q)^Ne^{-cq}.
 \label{deg:source-growing-stable-kernel-rows}
\end{equation}
The inverse passage thus reaches the inner cut positively, with positive
denominator, and its lip is the first ratio hit.
Define on this actual passage
\begin{equation}
 J(r)=\int_0^{T(r)}G(r(q;r),W(q;r))\,dq .
 \label{deg:source-growing-nonlinear-action-integral}
\end{equation}
The ratio equation gives exactly
\(G=-K_-W_q+(\tau K_-h_1-\varpi)W\). Integration by parts
therefore isolates the nonlinear mass:
\begin{equation}\label{deg:source-growing-exact-action-integral}
 J=\ell_0K_-(r)-K_-(r_c)W(T;r)+\int_0^T QW\,dq,
 \qquad Q=\tau(K_-h_1-\mathcal DK_-)-\varpi .
\end{equation}
Here \(\mathcal D^iQ=O(\tau)\), \(0\le i\le2\), using
\(\mathcal D^3K_-\) in the second row. The same fixed-clock
estimates give an integrable \(C\tau(1+q)^Ne^{-cq}\) bound for
\(QW\) and its phase rows through order two. Apply the endpoint formulas
\eqref{deg:stable-ratio-terminal-rows} to this density and the
terminal chain rules of Lemma~\ref{deg:lem-stable-ratio-ledger}
to \(K_-(r_c)W(T;r)\), with
\(T_1=1/\tau\), \(T_2=0\). The fixed gap \(R_0-r_c>0\)
makes all endpoint terms \(\tau^{-N}(1+T)^Ne^{-cT}\) flat.
Together with \(K_-(r)-1=O(\epsilon)=O(\tau)\), this proves
\begin{equation}
 |J-\ell_0|+|\mathcal DJ|+|\mathcal D^2J|\le C\tau .
 \label{deg:source-growing-nonlinear-action-rows}
\end{equation}
The constants depend on \(R_0-r_c>0\), not on the later \(R\).

The polynomial-memory identity
\eqref{deg:stable-ratio-polynomial-memory} retains the initial factor:
\begin{equation}
 \int_0^T\tau h_1(r(q;r))\,dq
             =\log P_-(r)-\log P_-(r_c).
 \label{deg:source-growing-polynomial-memory}
\end{equation}
Consequently the actual height at the inner cut is
\begin{equation}
 y_c(r)=\ell_0P_-(r)e^{-T(r)}e^{J(r)} .
 \label{deg:source-growing-inner-cut-height}
\end{equation}
Uniformly for \(r\ge R_0\),
\begin{equation}
 0<y_c(r)\le C(1+r)e^{-c(r-r_c)/\tau}
           \le C_N\tau^N\quad\hbox{for every fixed }N .
 \label{deg:source-growing-inner-cut-flatness}
\end{equation}
Its first two logarithmic rows are
\begin{equation}
 \mathcal D\log y_c=h_1-1/\tau+\mathcal DJ,\qquad
 \mathcal D^2\log y_c=\mathcal Dh_1+\mathcal D^2J .
 \label{deg:source-growing-inner-cut-log-rows}
\end{equation}

We now continue to the original common centre, without
renormalizing it to an image-dependent section.
On \(0\le v\le r_c\), \(P_-(v)\) and \(\mathscr B(v)\)
are positive units.  The exact linear multiplier is
\begin{equation}
 L_-(v)=\exp\left\{-\frac1\tau
               \int_v^{r_c}\frac{t\mathscr B(t)}{P_-(t)}\,dt
                  \right\}\le1,\qquad
 A_{\rm core,-}=\int_0^{r_c}\frac{t\mathscr B(t)}{P_-(t)}\,dt.
 \label{deg:source-growing-negative-core-linear-factor}
\end{equation}
The exact equation on this fixed compact segment is
\(y_v=A_-y+\mathcal K_-y^2\), where
\begin{equation}
 A_- =\frac{v\mathscr B}{\tau P_-},\qquad
 \mathcal K_-(v,y)=\frac{e_-s_0}{P_-+K_-y}
       -\frac{v\mathscr B K_-}{\tau P_-(P_-+K_-y)} .
 \label{deg:source-growing-negative-core-equation}
\end{equation}
On the original fixed small height buffer, all normal derivatives
through order three of \(\mathcal K_-\) are bounded by \(C/\tau\).
Its inward linear propagator is exactly \(L_-\le1\), and
\eqref{deg:source-growing-inner-cut-flatness} gives \(y_c/\tau\ll1\).
Lemma~\ref{deg:lem-compact-height-continuation}, with
\(\kappa_c=\tau\), gives
an actual positive centre hit and
\(y(0;y_c)=L_-(0)y_c\mathcal U_-(y_c)\), where
\begin{equation}
 |\log\mathcal U_-|+
 \sum_{i=1}^2|(y_c\partial_{y_c})^i\log\mathcal U_-|
                                  \le Cy_c/\tau .
 \label{deg:source-growing-negative-core-height-rows}
\end{equation}
This uses the full multiplier, with no reciprocal bound on it, and
holds for either sign of \(e_-(0)\).

Use the same centre-height scale as on the positive side:
\begin{equation}
 Z_-(r)=\frac{w(0;r)}{s_0^2\tau}
 =\ell_0P_-(r)
       \exp\left\{-\frac{\mathcal A_-(r)+A_{c,0}
                             +A_{\rm core,-}}{\tau}\right\}
       e^{J(r)}\mathcal U_-(y_c(r)).
 \label{deg:source-growing-negative-centre-height}
\end{equation}
Choose the coefficient-only logarithmic constant
\begin{equation}
 c_-=\log\ell_0-(A_{c,0}+A_{\rm core,-})/\tau .
 \label{deg:source-growing-negative-coefficient-constant}
\end{equation}
The separate negative remainder is
\begin{equation}
 S_-(r)=\log Z_-(r)+\mathcal A_-(r)/\tau
                            -\log P_-(r)-c_-
       =J(r)+\log \mathcal U_-(y_c(r)).
 \label{deg:source-growing-negative-source-remainder}
\end{equation}
The chain rule, \eqref{deg:source-growing-inner-cut-flatness}--\eqref{deg:source-growing-inner-cut-log-rows}, and \eqref{deg:source-growing-negative-core-height-rows} bound the two
core contributions by \(Cy_c/\tau^2\) and \(Cy_c/\tau^3\).
Combining with \eqref{deg:source-growing-nonlinear-action-rows} gives
\begin{equation}
 |S_-|\le C,\qquad
 |\mathcal DS_-|+|\mathcal D^2S_-|\le C\tau .
 \label{deg:source-growing-negative-source-rows}
\end{equation}
The physical centre height is therefore
\begin{equation}\label{deg:source-growing-physical-negative-height}
 z_-(r)=r_s^2s_0^2\tau e^{c_-}P_-(r)
                  e^{-\mathcal A_-(r)/\tau}e^{S_-(r)}.
\end{equation}
The factor \(e^{S_-}\), its inverse and its first two action
rows are uniformly bounded.


\smallskip
\noindent\emph{Positive source variations.}

The positive weak coordinate from
\eqref{deg:fixed-exit-exact-weak-coordinate} satisfies
\begin{equation}
 \Phi(P)=\{z_+y(P)\}^{\chi},\qquad
 \chi=\tau X_d,\qquad z_+^{\pm1}+X_d^{\pm1}\le C .
 \label{deg:source-growing-positive-exact-coordinate}
\end{equation}
The local invariant gives
\begin{equation}
 0\le y(P)\le C\vartheta^{1/\chi},\quad
 0<\vartheta<1,\quad 0\le P\le p_0,
 \label{deg:source-growing-positive-weak-flatness}
\end{equation}
and \(\Phi'>c>0\) with its stated uniform phase bounds.
Fix \(P_\dagger\in(0,p_0)\), independently of coefficients,
and set \(X=\Phi(P)/\Phi(P_\dagger)\).  On \(P>0\),
\begin{equation}
 \frac{y(P)}{y(P_\dagger)}=X^{1/\chi},\qquad
 \mathcal E:=X\partial_X=\chi^{-1}y\partial_y .
 \label{deg:source-growing-positive-log-phase}
\end{equation}
Use this exact logarithmic derivative.

Construct \(Z_+(P)\) independently by following positive time
from \(u_+=a(1-\eta_+)\), with \(w=a^2y(P)\), to \(u=0\).
In \(v=u/a\), \(v_+=1-\eta_+\), put
\begin{equation}
 h(av)/a^2=\tau H_+(v),\qquad
 L(av)/a=v\mathscr B_+(v),\qquad K_+(v)=K(av).
 \end{equation}
The original exact functions are
\begin{equation}
 H_+(v)=\frac{\mathcal R}{A_d}
        (1+av)(1-v)(c-av)
 \end{equation}
on the compact graph, and
\begin{equation}
 H_+(v)=\frac d{A_d}(1+av)(1-v)(1-jav)
 \end{equation}
on the reciprocal graph.  On \(0\le v\le v_+\),
\(H_+,\mathscr B_+,K_+\) are positive units, with uniform
finite coefficient-function bounds.  In particular the
positive core has no inverse \(\kappa\) scale.
Its exact scalar equation is
\begin{equation}
 y_v=\frac{y\{v\mathscr B_++aE(av)y\}}
                         {\tau H_++K_+y}.
 \label{deg:source-growing-positive-core-equation}
\end{equation}
Its linear multiplier and positive core action are
\begin{equation}
 L_+(v)=\exp\left\{-\frac1\tau
              \int_v^{v_+}\frac{t\mathscr B_+(t)}{H_+(t)}\,dt
                    \right\}\le1,\qquad
 A_{\rm core,+}=\int_0^{v_+}\frac{t\mathscr B_+(t)}{H_+(t)}\,dt.
 \label{deg:source-growing-positive-core-linear-factor}
\end{equation}
Freeze the coefficients and set \(y(v)=\tau^2\eta(v)\),
with initial height \(\eta_{\rm in}=y(P)/\tau^2\) at the same cut
\(v_+\). Exact division gives
\begin{equation}\label{deg:source-growing-positive-core-rescaled-equation}
 \eta_v=\frac{v\mathscr B_+}{\tau H_+}\eta
                         +\widetilde{\mathcal K}_+(v,\eta)\eta^2,
 \qquad
 \widetilde{\mathcal K}_+
 =\frac{\tau aE(av)}{H_++\tau K_+\eta}
 -\frac{v\mathscr B_+K_+}{H_+(H_++\tau K_+\eta)}.
\end{equation}
Choose a fixed small buffer \(0\le\eta\le\eta_b<\inf X_d^2\).
The nonlinear coefficient and its first three ordinary height
derivatives are bounded uniformly. Indeed the
displayed denominators are positive units and each such derivative
introduces a bounded factor \(\tau K_+\).
Then \(\tau^2\eta_b<\chi^2\) lies inside the original height range;
\eqref{deg:source-growing-positive-weak-flatness} puts every original
\(0<P\le p_0\) in a smaller initial subbuffer after a common
\(\tau\)-shrink.
The inward propagator remains exactly \(L_+\le1\), without a
reciprocal reserve. Thus
Lemma~\ref{deg:lem-compact-height-continuation}, with
\(\kappa_c=1\), supplies the actual positive centre hit
and its Euler-height rows. Since
\(\eta_{\rm in}\partial_{\eta_{\rm in}}=y_+\partial_{y_+}\)
for \(y_+=y(P)\), these give
\begin{equation}
 |\log \mathcal U_+(y_+)|+
 \sum_{i=1}^2|(y_+\partial_{y_+})^i\log \mathcal U_+(y_+)|
                                      \le Cy_+/\tau^2 ,
 \label{deg:source-growing-positive-core-height-rows}
\end{equation}
where \(\mathcal U_+\) is the normalized factor at \(v=0\).
The actual central height is
\begin{equation}
 Z_+(P)=C_+(\theta)y(P)\mathcal U_+(y(P)),\qquad
 C_+(\theta)=\frac{a^2}{s_0^2\tau}
                                      e^{-A_{\rm core,+}/\tau}.
 \label{deg:source-growing-positive-centre-height}
\end{equation}
This retains exactly \(a^2/(s_0^2\tau)=A_d^2/(B\tau)\).
The prefactor of the exponential is a coefficient-only unit times
\(\tau^{-1}\).
Scalar first variation gives \(Z_+'(P)>0\) for \(P>0\).

Define on the original independent positive domain
\begin{equation}
 S_+(P)=\log Z_+(P)-\log Z_+(P_\dagger)-\chi^{-1}\log X
 =\log \mathcal U_+(y(P))-\log \mathcal U_+(y(P_\dagger)).
 \label{deg:source-growing-positive-source-remainder}
\end{equation}
The coefficient-only multiplier cancels as an exact reference
difference.  Equations \eqref{deg:source-growing-positive-log-phase} and \eqref{deg:source-growing-positive-core-height-rows} give
\begin{equation}
 |\mathcal E^iS_+|
 \le C\chi^{-i}\tau^{-2}y(P)
 \le C\tau^{-i-2}\vartheta^{1/\chi},\qquad i=1,2 .
 \label{deg:source-growing-positive-source-rows}
\end{equation}
The rows vanish continuously at $P=0$. This positive map remains
independent of the negative image.

\smallskip
\noindent\emph{Admission to the positive input chord.}

Choose any fixed \(p_*\in(0,p_0)\) inside the original counted
weak chord, with its larger original raw chord buffer.
Put \(j(P)=\Phi(P)/\Phi'(P)\).  Its uniform phase bounds give
\(j'(P)=1+O(P)\); choose \(p_*\) small enough that \(j'(P)>3/4\)
on \(0\le P\le p_*\).  The endpoint is fixed before \(R\).
By \(\Phi'>c\), \(\Phi(p_*)\) is bounded below by a fixed
positive constant.  Equations \eqref{deg:source-growing-positive-exact-coordinate}, \eqref{deg:source-growing-positive-core-linear-factor}, and \eqref{deg:source-growing-positive-centre-height} give
\begin{equation}
 -\tau\log Z_+(p_*)\le C_*(p_*) ,
 \label{deg:source-growing-positive-endpoint-action-budget}
\end{equation}
uniformly in \(\kappa,\tau,j\) after a common small bound.
Indeed, $-\tau\log y(p_*)=-(\tau/\chi)\log\Phi(p_*)
+\tau\log z_+$ and $A_{\rm core,+}$ are bounded; the remaining
scale costs $O(\tau(1+|\log\tau|))$.

On the negative side, \eqref{deg:source-growing-negative-centre-height} gives
\begin{equation}
 -\tau\log Z_-(R)\ge
 c(R-r_c)-\tau\log\{\ell_0P_-(R)\}-C\tau .
 \label{deg:source-growing-negative-endpoint-action-budget}
\end{equation}
First choose a fixed \(R=r_{\rm ge}^{\rm out}>R_0\), larger than any
already prescribed original growing threshold, so that
\(c(R-r_c)>C_*(p_*)+2d_*\) for a fixed \(d_*>0\).
Only afterwards decrease the coefficient bounds so that
\(\tau\log(1+R)\) and the fixed terms in \eqref{deg:source-growing-negative-endpoint-action-budget} are below
\(d_*\).  This preserves the full independent coefficient
family.
Choose the counted threshold \(r_{\rm ge}>R\) with a fixed
positive buffer before decreasing the coefficients. It follows that
\begin{equation}
 0<Z_-(R)<e^{-d_*/\tau}Z_+(p_*).
 \label{deg:source-growing-positive-image-inclusion}
\end{equation}
Furthermore \eqref{deg:source-growing-negative-source-remainder}--\eqref{deg:source-growing-negative-source-rows} imply
\begin{equation}
 \mathcal D\log Z_-=-1/\tau+h_1+\mathcal DS_-<0 .
 \label{deg:source-growing-negative-centre-order}
\end{equation}
Thus every actual growing lip in \eqref{deg:source-growing-separate-domains} has centre height
strictly between zero and \(Z_+(p_*)\).
The continuous strictly increasing positive map \(Z_+\)
maps \([0,p_*]\) onto \([0,Z_+(p_*)]\).
There is a unique \(P(r)\in(0,p_*)\) with
\(Z_+(P(r))=Z_-(r)\).
Scalar uniqueness at the common centre identifies the two
actual graph segments with the same orbit.

The local weak chord is the one parameterized by the original
raw input
\begin{equation}
 P=(x_{\rm raw}-x_b)/s_{\rm in},\qquad
 s_{\rm in}=\Omega a^{2-\lambda_\gamma}U_{\rm in},\qquad
 \Omega=r_s^{2-\lambda_\gamma},\quad U_{\rm in}^{\pm1}\le C.
 \label{deg:source-growing-positive-raw-input-scale}
\end{equation}
The independently constructed stopped outer chord and its
local invariant give the actual continuation to this raw
section for all \(0\le P\le p_0\), with its original buffers.
The matched inputs lie in the strict interior \(P<p_*\).
Both local points are on the same invariant level in the
original coordinate box; the monotone positive invariant
coordinate gives the intervening orbit segment there.
Thus every lip has the actual positive-time continuation from raw
input through the positive middle cut and centre to the negative lip.
The separate annular overlap remains the earlier construction.

\smallskip
\noindent\emph{The matched action derivative.}

On the separate original product domain define
\begin{equation}
 R_{\rm prod}(P,r)=\log P_-(r)+S_-(r)-S_+(P).
 \label{deg:source-growing-product-remainder}
\end{equation}
It remains defined whether or not the two images intersect.
Its phase rows satisfy
\begin{equation}
 \begin{gathered}
 |\mathcal ER_{\rm prod}|+|\mathcal E^2R_{\rm prod}|
                  \le C\tau^{-4}\vartheta^{1/\chi},\\
 |\mathcal DR_{\rm prod}|+|\mathcal D^2R_{\rm prod}|\le C,
 \qquad \mathcal E\mathcal DR_{\rm prod}=0 .
 \end{gathered}
 \label{deg:source-growing-product-phase-rows}
\end{equation}
Although \(\log P_-(r)\) can grow, its two action derivatives have
the bounds needed here.
On the actual matched graph the exact identity is
\begin{equation}
 \mathcal A_-(r)=
 \tau\{c_--\log Z_+(P_\dagger)\}
             -e_s\log X+\tau R_{\rm prod}(P,r),
 \qquad e_s=\tau/\chi=1/X_d .
 \label{deg:degree-growing-matching}
\end{equation}
The constant retains all coefficient-only terms; $e_s$ is a positive unit.

Define separate product functions
\begin{equation}
 N_-=1-\tau\{h_1+\mathcal DS_-\},\qquad
 N_+=1+\chi\mathcal ES_+,\qquad n(P,r)=N_+(P)/N_-(r).
 \label{deg:source-growing-separate-action-numerators}
\end{equation}
For a common small \(\tau\) bound,
\begin{equation}
 N_\pm\ge1/2,\quad
 |N_--1|+|\mathcal DN_-|\le C\tau,\quad
 |N_+-1|+|\mathcal EN_+|
                       \le C\tau^{-3}\vartheta^{1/\chi}.
 \label{deg:source-growing-action-numerator-rows}
\end{equation}
Differentiate equality of the two actual centre heights:
\(\mathcal E\log Z_+=N_+/\chi\) and
\(\mathcal D\log Z_-=-N_-/\tau\).
It gives exactly
\begin{equation}
 -\mathcal E r=e_s\mathfrak d(r)n(P,r)
     =e_s\frac{\mathcal P(r)}r(1+\zeta),\qquad \zeta=n-1 .
 \label{deg:source-growing-actual-matching-slope}
\end{equation}
The matched graph is consequently strictly decreasing. To differentiate
the quotient along this graph, both separate variables must be varied:
\begin{equation}
 \mathcal E_{\rm graph}\zeta=
 \mathcal E_+n-e_sn\mathcal D_-n .
 \label{deg:source-growing-actual-quotient-first-row}
\end{equation}
Here \(\mathcal E_+\) and \(\mathcal D_-\) act on the separate
positive and negative variables of the product.
The second pure derivatives in \eqref{deg:source-growing-negative-source-rows}
and \eqref{deg:source-growing-positive-source-rows} bound precisely
the two derivatives in \eqref{deg:source-growing-actual-quotient-first-row}.  Thus
\begin{equation}
 |\zeta|+|\mathcal E_{\rm graph}\zeta|
 \le C\{\tau+\tau^{-3}\vartheta^{1/\chi}\}
 \le C\tau .
 \label{deg:source-growing-actual-quotient-estimates}
\end{equation}
This is uniform on the whole enlarged graph with independent
$\kappa,\tau,j$; the error tends to zero with $\tau$.

The normalized source heights retain the physical conversion
\(z(\gamma_s)=r_s^2s_0^2\tau Z\), the exact raw power
\(\Omega\), the positive and negative tail quartets on the
reciprocal graph
\begin{equation}
 \left(-\frac{\tau d}{A_d},
       \frac{\tau d(1-a+ja)}{A_d},\gamma_s,r_sa\right),
 \qquad
 \left(-\kappa^2\tau dA_d,
       \kappa\tau d(1-a+ja),\gamma_s,r_s\right),
 \label{deg:source-growing-remote-raw-quartets}
\end{equation}
and the passive coefficient restrictions used to construct the raw
chords.  No estimate of \(\gamma_s\log r_s\), no derivative
with respect to a coefficient, and no frozen moving raw endpoint
is used.  At \(j=0\) only the one-sided normalized estimates
are intended; actual raw sections are asserted for \(j>0\).
\end{proof}


\begingroup
\let\degreecarriersubsection\subsection
\let\degreecarriersubsubsection\subsubsection
\let\section\degreecarriersubsection
\let\subsection\degreecarriersubsubsection
\section{One physical interval and four passage regions}

The preceding constructions use different phase coordinates. We now show
that their domains cover one connected interval of actual through
trajectories, with open overlaps. Their derivative estimates can then
be combined for the same affine displacement.

\begin{lemma}[A common entrance and gate choice]
\label{p2:lem-joint-entrance}
One auxiliary choice for the four passage regions, followed by
independent small coefficient bounds, has the following properties.
The incoming raw height $h$ is $x_{\rm raw}$ in
\eqref{deg:fixed-exit-natural-input} and $x_{\rm in}$ in
Lemma~\ref{lem:actual-dulac-matching}; its stopped value is $x_b$.
Thus the canonical weak input is exactly
\begin{equation}\label{p2:joint-raw-input}
 P=q(h)=\frac{h-x_b}{s_{\rm in}},\qquad
 s_{\rm in}=\Omega a^{2-\lambda_\gamma}U_{\rm in}>0,
 \qquad \Omega=r_s^{2-\lambda_\gamma},\quad U_{\rm in}^{\pm1}\le C.
\end{equation}
This is the affine coordinate on the original raw section, distinct
from the displacement coordinate $x=h-a_\theta$; no new input
observer is inserted.

Let $\widehat p_0=p_0$ be a fixed buffered raw cutoff in
\eqref{deg:fixed-exit-middle-invariant}, and choose
$0<\widehat p_*<\widehat p_0$ so that
$\mathfrak j'=(\Phi_\theta/\Phi_\theta')'>3/4$ on
$[0,\widehat p_*]$. Set $q_{\rm ref}=\widehat p_*/4$ and use
$P_\dagger=q_{\rm ref}$ in
\eqref{deg:source-growing-positive-log-phase}. The final negative
thresholds are denoted by
$\widehat r_{\rm ge}^{\rm out}=r_{\rm ge}^{\rm out}$ and
$\widehat r_{\rm ge}=r_{\rm ge}$, with a fixed positive gap. Every
actual first-lip input with $r_\ell\ge r_{\rm ge}$ satisfies
\begin{equation}\label{p2:joint-entrance-buffer}
 0<P<q_{\rm ref}/2<q_{\rm ref}<\widehat p_*<\widehat p_0,
\end{equation}
including first hits outside the growing product range.

The source constant $\delta_\sharp(L)$ is selected once for the full
bounded-source domain and relative negative-curvature estimate. Its gates
are those below in Lemma~\ref{p2:lem-bounded-source-budget}.
The final $r_{\rm ge}$ precedes $R_{\rm BD}$ and its annular
overlap; all fixed gates and phase buffers precede the final coefficient
bounds. No comparison of the independent scales is imposed.
\end{lemma}
\begin{proof}
Fix the original raw windows and a finite $V_1>1$, say $V_1=2$.
Choose the local cuts and ordinary $C^3$ section maps for this value in
Lemma~\ref{lem:actual-dulac-matching}, retaining its enlarged
interval $0\le V\le V_1+1$. Keep these choices throughout.
Use the strict source block and the two comparison orbits of
Proposition~\ref{bd:prop-bounded-action}: fix its derivative constants,
then $B_-$, then $\delta_\sharp(L)$ for both the domain and relative
estimate, and finally $\delta_1,\delta_2$ and the gates below.
Fix $r_c>2$ and choose $\ell_0$ below the ratio of a fixed real
comparison orbit at $-r_c$, as in
\eqref{deg:degree-ratio-monotonicity}.

Fix a source envelope $\bar\sigma_{\rm src}=1/3$ for the separate
centre-map estimates on $r_c\le r\le\bar\sigma_{\rm src}/s_0$ in
Lemma~\ref{deg:lem-growing-matching-direct}; its derivative constants
are fixed before the positive cutoff and negative thresholds below.
The growing and output intervals will be chosen inside this envelope.
The independent stopped outer chord and local invariant give the
positive input map on $[0,p_0]$ in
\eqref{deg:fixed-exit-natural-input}--\eqref{deg:fixed-exit-exact-weak-coordinate}.
Choose the smaller fixed $\widehat p_*$ above so that the positive
action at $P_0=q_{\rm ref}/2$ also exceeds the fixed middle
variational exponent by two. This is possible since
$\Phi_\theta(P)=PU_\theta(P)$ with uniform positive units and
$\tau/\chi$ a positive unit. These choices precede the final negative
threshold.

The full first-hit cutoff following \eqref{deg:degree-phase-order}
provides a fixed $R(P_0)$. Choose $R=r_{\rm ge}^{\rm out}$ larger
than this and every earlier growing threshold, and so that
$c(R-r_c)>C_*(P_0)+2d_*$ for a fixed $d_*>0$ in
\eqref{deg:source-growing-positive-endpoint-action-budget}.
Choose $r_{\rm ge}>R$ with a fixed positive buffer.
The separate centre maps give
\eqref{deg:source-growing-positive-image-inclusion} with endpoint
$P_0$ instead of $p_*$ after independent coefficient shrink.
Their monotonicity and scalar uniqueness therefore construct the
actual product reference $r_\ell=r_{\rm ge}$ with raw input
$0<P_{\rm ge}<P_0$, in the original buffered chord. The local
invariant supplies its positive-time continuation to the raw section,
as in \eqref{deg:source-growing-positive-raw-input-scale}.

On the full first-hit domain of the original through interval,
\eqref{deg:degree-phase-order} gives $r_{\ell,h}<0$.
Thus $r_\ell(h)\ge r_{\rm ge}$ implies
$h\le h_{\rm ge}=x_b+s_{\rm in}P_{\rm ge}$, and the same positive
affine input gives $0<P(h)\le P_{\rm ge}<P_0$.
This proves \eqref{p2:joint-entrance-buffer} throughout the first-hit
domain, including inputs beyond the product range. The extension follows
from order on the full domain; the product derivative estimates remain
restricted to their original range.

Only now choose $R_{\rm BD}$ so the compact strong projection of
$r_\ell=r_{\rm ge}$ lies strictly below $2R_{\rm BD}$.
The transverse inverse and first-hit order give a fixed
$r_{\rm ge}^{\sharp}>r_{\rm ge}$ in the enlarged annular overlap.
The annular proof now fixes its own outer cuts \(L_{\rm a}\), after
this output bound and its limiting-action envelope. The earlier source
cut \(L\), \(h_L\), \(\delta_\sharp(L)\), and literal gates stay fixed;
both passages continue to the same raw ports. The original weak cuts
satisfy \(c_+=p_{\rm a}/8<p_-=p_{\rm a}/2\) and
\(c_-=p_{\rm a}/8<r_{\rm a}/2\), so the annular root-gap construction
uses this same auxiliary tuple. Only final independent coefficient bounds
are reduced to admit the new positive-height tail cuts.
Choose $0<\sigma_1<\sigma_2<\sigma_0<1/4$, the common
negative section $u_c$, its larger
$\sigma_0<\bar\sigma\le\bar\sigma_{\rm src}$ buffer, with
$\bar\sigma<|u_c|/2$, and strict
action gap $\Delta$ as in the construction after
\eqref{deg:degree-phase-order}. After these fixed choices, shrink
the independent coefficients so the fixed annular overlap lies in
the counted product range and $s_0r_{\rm ge}^{\sharp}<\sigma_1$.
All source, raw and shooting buffers retain strict margins.
All these restrictions are independent upper bounds on the coefficients.
After they have been imposed, the finite multiplier will be chosen.
\end{proof}

\begin{proposition}[One actual carrier for the four passage regions]
\label{p2:prop-unified-carrier-cover}
Use the common choice of Lemma~\ref{p2:lem-joint-entrance}, the original field,
through word, sections, raw windows, $L,h_L$ and closing maps.
The source constant $\delta_\sharp(L)$ is the once-selected constant
for the full original bounded-source domain and its relative sign;
$\delta_1,\delta_2$ keep their literal values from that tuple.
Let $I_\theta$ be the connected interval of actual through inputs in
this word, within the original buffered windows and positive graph
denominator domain, with
\[
             0<H_\theta(h):=z(\gamma_s;h)<z_{\rm cap}.
\]
Here $z_{\rm cap}$ is the original common fixed raw-height cap,
chosen within the standing high-height bound. The stopped trace $x_b$
is excluded. Write $h^0$ for the actual trace of $h_L=0$ and
$I_\theta^- = I_\theta\cap\{h\le h^0\}$.

\begin{enumerate}
\item The original through interval $(x_b,h^0)$ has the strictly increasing
shooting coordinate $W_\theta(h)=w(u_c;h)$ with image
$(0,W_\theta(h^0))$, where $W_\theta(h^0)\to u_c^2/2$ uniformly.
Thus every point of the full original enlarged chord
$\overline{\mathcal J_c^{\sharp\sharp}}$ corresponds to an actual through orbit,
with a strict height buffer. Its backward shot and both raw continuations
are that same orbit.
\item Let $h_S$ be the trace of $h_L=-\delta_2$ and let $h_A,h_B$
be the gate traces of $r_b=r_a/2,2R_{\rm BD}$, respectively. Put
\[
 B_\theta=I_\theta^-\cap([h_S,h^0]\cup[h_B,h_A]),\qquad
 h_T=\min\{h_S,h_B\}.
 \tag{PCov1}
\]
Let $I_{\ell,\theta}$ be the original first-ratio-hit domain. On it put
$\sigma(h)=s_0r_\ell(h)$. Wherever $r_\ell\ge r_{\rm ge}$ use
the original raw middle input
\[
 V(h)=\frac{w(u_+;h)}{\ell_0h_{\rm ax}(u_+)\rho(u_+)}
     =\frac{y(u_+;h)}{\varepsilon_+}>0,
 \quad h_{\rm ax}=-\delta p_d,\quad
 \varepsilon_+=\frac{\ell_0h_{\rm ax}(u_-)}{a^2J_-}.
 \tag{PCov2}
\]
All factors in this definition are actual and positive for a physical
coefficient fibre. Define the following subsets of that same carrier:
\[
\begin{aligned}
 G_\theta&=I_\theta^-\cap I_{\ell,\theta}
          \cap\{r_\ell\ge r_{\rm ge},\ \sigma\le\sigma_0\},\\
 D_\theta&=I_\theta^-\cap I_{\ell,\theta}
          \cap\{\sigma\ge\sigma_1,\ 0<V\le V_1\},\\
 E_\theta&=I_\theta^-\cap I_{\ell,\theta}
          \cap\{\sigma\ge\sigma_1,\ V\ge1\}.
\end{aligned}                                                    \tag{PCov3}
\]
The $V$ coordinate in the last two lines is well defined because
$\sigma_1/s_0>r_{\rm ge}$. These are connected, relatively closed
intervals, possibly empty, and
\[
 I_\theta^-=B_\theta\cup G_\theta\cup D_\theta\cup E_\theta.
 \tag{PCov4}
\]
Each of $G_\theta,D_\theta,E_\theta$ has the exact raw input buffer
\[
 0<q(h):=\frac{h-x_b}{s_{\rm in}}<q_{\rm ref}
       <\widehat p_*<\widehat p_0,
 \qquad s_{\rm in}=\Omega a^{2-\lambda_\gamma}U_{\rm in}>0.
 \tag{PCov5}
\]
The compact annular/growing overlap and the original product and
$V$ overlaps are on these actual inputs, with their equality dividers
included. Intersecting with any original closed raw-window or height
truncation preserves the cover; no divider is required to occur.
\item $G_\theta$ lies in the growing product domain, and
$D_\theta,E_\theta$ lie in the enlarged domains of the finite-$V$ and
fixed-exit maps constructed in
Proposition~\ref{p2:prop-finite-v-fe-actual-budget}.
The high-height interval is $I_\theta^+=I_\theta\cap\{h\ge h^0\}$.
Together with (PCov4) it covers the whole $I_\theta$, including every
non-root through input. On all these intervals the physical equation is
\[
 F(h-a_\theta)=h-a_\theta-f_\theta(m\mathfrak T_\theta(h)),
 \qquad \mathfrak T_\theta=d_\theta\circ T_\theta^\infty.
 \tag{PCov6}
\]
\end{enumerate}
The statements about actual raw maps concern $j>0$; the reciprocal
face $j=0$ has only the uniform one-sided normalized interpretation.
\end{proposition}
\begin{proof}
The independently constructed stopped and $h_L=0$ graphs bound the
through interval. Every intermediate graph reaches $u_c\in(-1,u_-)$ with
positive denominator. Scalar first variation makes $W_\theta$ strictly
increasing, with stopped limit zero.
For the upper comparison graph set $r_0=s_0L$, $r=-u$ and $\mathscr W=w/r^2$.
At its initial cut
\[
 \mathscr W(r_0)=\tfrac12+O_L(\tau b_\tau),\qquad
 b_\tau=-\log(2\tau).
\]
On $r_0\le r\le |u_c|$, direct substitution of the original polynomial
in the graph equation gives, on $1/4\le\mathscr W\le3/4$,
\[
 r\mathscr W_r=1-2\mathscr W+\mathscr E,\qquad
 |\mathscr E|\le C\{\kappa\tau+\kappa\tau/r\}
                 \le C\{\kappa\tau+\tau/L\}.             \tag{PCov7}
\]
Indeed $|\delta p_d(-r)|\le C\kappa\tau(r+a)$ and
$a/r\le x_{\max}/L$, while $K_d=1+O(\kappa\tau)$ and
$-L_d(-r)/r=1+O(\kappa\tau)$; the $E_dw$ term is smaller on this
corridor. The regular reciprocal factorization gives the same bounds
uniformly as $j\downarrow0$. Both boundary vectors point inward.
The logarithmic-radius equation contracts, so variation of constants
gives $\mathscr W(|u_c|)\to1/2$. Hence
$W_\theta(h^0)\to u_c^2/2$. The uniform endpoint margins put
$\overline{\mathcal J_c^{\sharp\sharp}}$ strictly inside this shooting
image after shrink. Strict first variation gives a unique through input
for each $W$; backward uniqueness identifies its shot and the two
already constructed raw continuations. This proves (i) independently
of the fixed-exit phase.

The same incoming first variation preserves the ordering of the
actual $h_L$ labels. Thus the central source interval is $[h_S,h^0]$.
The exact gate inverse and centre matching give the annular interval
$[h_B,h_A]$. The fixed orbit $h_L=-\delta_1$ has trace $h_C$ strictly
inside both intervals, because $-\delta_2<-\delta_1<0$ and its actual
gate labels lie strictly inside the annular bands. Consequently
\[
 [h_S,h^0]\cup[h_B,h_A]=[h_T,h^0].                       \tag{PCov8}
\]
The first fixed comparison alone supplies this overlap. The compact overlap also supplies an actual reference with $r_\ell=r_{\rm ge}$ and $r_b<2R_{\rm BD}$;
its trace $h_{\rm ge}$ satisfies $h_B<h_{\rm ge}$.
Every input failing (PCov8) is below $h_T\le h_B$, so the supplied full
original first-hit order gives $r_\ell>r_{\rm ge}$. The unified entrance
then gives (PCov5) on all inputs with $r_\ell\ge r_{\rm ge}$, including
first hits beyond the product range. Their derivative estimates are applied only within the growing product
domain.

For this actual local input $q$, the original positive cut is reached
on the stopped outer chord's local invariant. Its height is positive,
finite and strictly increasing in $h$. Formula (PCov2) is exactly the
original normalization $w=\ell_0h_{\rm ax}\rho V$, so
$V_h>0$, while the full first-hit order gives $\sigma_h<0$.
All sets in (PCov3) are therefore intervals and include their finite
level boundaries. If an input outside $B_\theta$ has
$\sigma\le\sigma_0$, it is in $G_\theta$. Otherwise
$\sigma>\sigma_0>\sigma_1$ and, because $V_1>1$, its positive $V$
satisfies $V\le V_1$ or $V\ge1$. It is in $D_\theta$ or
$E_\theta$. This proves (PCov4) for every through input, whether or not $F$ vanishes.

The compact band $r_\ell\in[r_{\rm ge},r_{\rm ge}^{\sharp}]$
is strictly inside the annular interval and, after the prescribed
$s_0$ bound, the growing product domain. The inequalities
$\sigma_1<\sigma_2<\sigma_0\le\bar\sigma<|u_c|/2$ retain the actual
product overlap. At $\sigma=\sigma_0$ the other passage region has a strict
$\sigma_1$ buffer, and at $V=1,V_1$ the last two regions overlap
because $1<V_1$. If a level is absent, the corresponding truncated interval is empty. The bounded endpoints and annular/growing overlap are included in these
intervals.

On $G_\theta$, $r_\ell\ge r_{\rm ge}>r_{\rm ge}^{\rm out}$ and
$\sigma\le\sigma_0\le\bar\sigma$, so the independently constructed
negative lip and positive chord identify the actual growing graph.
On $D_\theta\cup E_\theta$ the same raw buffer (PCov5) supplies the
original positive local cuts. The construction in Proposition~\ref{p2:prop-finite-v-fe-actual-budget}, using (i) as its
independent shooting geometry, identifies the bounded-$V$ early-cut
passage and the fixed-exit shooting map with precisely these orbits. In particular its action bound places the whole $E_\theta$
inside the enlarged inverse image, with $\Theta=0$ buffered. This last step is the only use of that proposition; the shooting
geometry in part (i) was constructed independently.

Finally the $h^0$ comparison graph has original source height
$Y_0\asymp b_\tau\to\infty$ and physical height
$H_\theta(h^0)=s_s^2Y_0\to0$. The original source dictionary gives
$|\alpha_0|+\alpha_1^2=s_s^2(1+p_s^2)$, with $p_s^2=O(\tau)$.
It follows that
\[
 K(|\alpha_0|+\alpha_1^2)<H_\theta(h^0)<z_{\rm cap}.
\]
Scalar height order puts every retained input at or above $h^0$ in
that same original high-height interval. The real enlarged $h_L=0$
source chord supplies its open overlap with the low carrier. Thus
$I_\theta^+$ includes the equality $h^0$ and all its non-root inputs;
the upper cap is common and fixed, rather than selected at a zero.
Original-time uniqueness gives (PCov6), with no raw gain inverted in
the domain construction.
\end{proof}

\section{Bounded source}

\begin{lemma}[Physical derivative estimates on the bounded-source intervals]
\label{p2:lem-bounded-source-budget}
Keep the original field, $L,h_L$, raw ports and enlarged windows, regular
maps and ordinary $C^3$ divided closing. Keep the unified tuple's literal
\[
 \begin{gathered}
 \delta_1=\min\{\delta_\sharp(L)/8,x_{\min}^2/8192,1/8192\},\\
 \delta_2=2\delta_1,\quad p_{\rm cmp}=\sqrt{2\delta_1},\quad r_a=p_{\rm cmp}/8,\quad p_a=r_a/16,\\
 p_-=p_a/2,\quad c_\pm=p_a/8.
 \end{gathered}
\]
The constant $\delta_\sharp(L)$ was fixed once, sufficiently small for
the bounded-source domain and relative estimate, before these gate
constants were defined. On the actual source word put
\[
 \beta=h_L/\tau-\log(2\tau),\quad b_\tau=-\log(2\tau),\quad
 \beta_{\min}=-\delta_2/\tau-\log(2\tau).
\]
There are fixed thresholds $B_-,B_+$ and constants $C_B,C_H$, uniform in
the original compact and reciprocal coefficient families, such that:
\begin{enumerate}
\item On every actual input with $\beta_{\min}\le\beta\le-B_-$,
the complete divided transition $\Pi$ has $\Pi''<0$ in the original
affine raw input, hence $F''=-m\Pi''>0$ for $m>0$.
\item On $-B_-\le\beta\le B_+$, one has $0<\Pi'\le C_B$.
\item On $B_+\le\beta\le b_\tau$, including $h_L=0$, and on a real
enlarged source chord through that orbit, one has $0<\Pi'\le C_H$.
The same bound holds on its admitted height-dominant continuation.
\end{enumerate}
All assertions hold on enlarged overlap intervals and at every nonclosed
through input. One final choice $m\max\{C_B,C_H\}<1/2$ gives $F'>1/2$
on the last two pieces. The input is always
$x=h-a_\theta$, and
\[
 \Pi(x)=\mathcal C_{\theta,m}(T_\theta(a_\theta+x)),\qquad
 F(x)=x-m\Pi(x)=x-f_\theta(m d_\theta(T_\theta(a_\theta+x))).
 \tag{BS1}
\]
\end{lemma}

\begin{proof}
Apply Proposition~\ref{bd:prop-bounded-action} to the same sections
and divided closing. Its estimates
\eqref{bd:relative-negative-curvature}--\eqref{bd:relative-raw-curvature}
give the first assertion, and \eqref{bd:compact-raw-variation}
gives the second.
The independent frozen windows, $L_f=8$, $R_f=5$ and polynomial
margin $3/4$ are supplied by Lemma~\ref{bd:lem-source-ranges};
both exact raw gains and their common-power cancellation are retained
in (BS1). It remains to check the literal gates and the upper chord.

The final $\delta_\sharp(L)$ is chosen there for both the domain and
relative curvature estimate, after $B_-$ and before the displayed literal
$\delta_1,\delta_2$. Since $\delta_2\le\delta_\sharp/4$,
the central block has enlarged lower chords down to
$h_L=-3\delta_1>-\delta_\sharp$. Compact labels
$[-B_--1,B_++1]$ and positive labels beginning below $B_+$ give actual
overlaps by monotone incoming first variation. A final $\tau$ shrink
makes $\beta_{\min}<-B_-$ and $b_\tau>B_++2$.

For positive labels choose the signed high-height constants
$K,z_*,C_H$ first, then $B_+$ large. Equation~\eqref{bd:positive-height}
gives, through $\beta=b_\tau$ and the enlarged chord below,
\[
 Y_0>K(1+p_s^2),\qquad
 H=s_s^2Y_0\le C_L\tau^3\sigma^4j^2(1+b_\tau)<z_*.
 \tag{BS10}
\]
To realize the upper chord, freeze $B=b_\tau$ in the original
positive-energy construction. For $\beta/B\in[3/4,3/2]$, its incoming
energy $\beta/B+B^{-1}\log(2/B)$ stays strictly inside $[1/4,3]$.
The energy defect is $O_L(B^{-1})$ and its variation integral tends
uniformly to zero; the same first-contact proof and positive inverse
therefore construct this entire chord. It has
$Y_0=\beta-\log\beta+O_L(1)$ and
$\eta(\pm L)=L^2/2+O_L(\tau b_\tau)$ at the unchanged cuts.
The original quadratic corridor and raw-window margins continue both
sides. In particular $0\le h_L\le\tau b_\tau/2$ is a real enlargement
above the $h_L=0$ orbit.

By (BS10), the signed high-height bound used in
Proposition~\ref{bd:prop-bounded-action}, namely
Corollary~\ref{rew:cor-high-derivative-signed}, applies throughout
this chord and its continuation
$K(|\alpha_0|+\alpha_1^2)\le H<z_*$, including the equality face.
It gives the third assertion, with the height scale frozen at each base point.

The full transported output, including reference values, has a common
bound $M_{\rm out}$ on the original section windows. Choose $m$ last
with $m\le m_0$, $mM_{\rm out}<a_{\rm cl}/2$ and
$m\max(C_B,C_H)<1/2$. Then (BS1) gives $F'>1/2$ on the compact,
positive and high pieces, while the negative curvature gives $F''>0$
for $m>0$. The enlarged estimates include threshold equalities and
every nonclosed input. At $m=0$, $F(x)=x$; collapsed sections denote
only normalized one-sided limits.
\end{proof}

\section{Growing exit}

\paragraph{The growing displacement and its curvature.}
The action derivative of the matched map is known, but the finite
closing is evaluated at its full translated output. We first bound that
output and its first two action derivatives, then use the logarithmic
slope to prove convexity in the original affine input.
All phase derivatives fix every coefficient. Use the already constructed
joint auxiliary tuple, the actual first lip $r=r_\ell$, and its original
enlarged domain
\[
 \widehat r_{\rm ge}^{\rm out}\le r\le\bar\sigma/s_0,\qquad
 s_0=\kappa\sqrt B,\qquad 0<\sigma_0<\bar\sigma<1/2 .
 \tag{GC1}
\]
The counted interval has $r\ge\widehat r_{\rm ge}$ and $s_0r\le\sigma_0$.
Write the actual affine raw input as
\[
 h=x_b+s_{\rm in}P,\qquad
 s_{\rm in}=\Omega a^{2-\lambda_\gamma}U_{\rm in}>0,\qquad
 \Omega=r_s^{2-\lambda_\gamma},\qquad U_{\rm in}^{\pm1}\le C .
 \tag{GC2}
\]
The joint entrance construction supplies $0<P<\widehat p_*/2$.
Set
\[
 X=\Phi_\theta(P)/\Phi_\theta(q_{\rm ref}),\qquad
 \mathcal E=X\partial_X=\mathfrak j(P)\partial_P,\qquad
 \mathfrak j=\Phi_\theta/\Phi_\theta'>0,\qquad
 e_s=\tau/\chi ,
 \tag{GC3}
\]
with the denominator held fixed. Let
\[
 \mathcal P(r)=P_s(-r)/B_\tau(-r),\qquad
 \mathfrak d=\mathcal P/r,\qquad \mathcal D=\mathfrak d\partial_r .
\]
The polynomial estimates and the matched action derivative give
\[
 \begin{gathered}
 \mathcal P\asymp1+r,\quad \mathcal P'>0,\quad
 \mathfrak d^{\pm1}\le C,\quad |\mathcal D\log\mathfrak d|\le C,
 \quad e_s^{\pm1}\le C,\quad \mathfrak j'>3/4,\\
 -\mathcal Er=e_s\mathfrak d\,n,\qquad n=1+\zeta,\qquad
 n\ge1/2,\qquad |\zeta|+|\mathcal E\zeta|\le C\tau .
 \end{gathered}
 \tag{GC4}
\]
Here $\mathcal E$ acts on the actual graph, including the negative
variable in the complete quotient row.

Retain the full actual output
\[
 \begin{gathered}
 T(r)=\Xi(W_\theta(r))
      =t_\dagger-a_{\rm grow}\mathcal B(r),\qquad
 a_{\rm grow}=\Omega s_0^2U_{\rm out}>0,\qquad
 \mathcal B'=rG,\quad G>0,\\
 W_\theta'=-s_0^2r\Gamma_\theta,\qquad
 G=\Gamma_\theta\,\Xi'(W_\theta)/(\Omega U_{\rm out}) .
 \end{gathered}
 \tag{GC5}
\]
Thus $T$ includes its coefficient-dependent translation, original
moving endpoint, tails and regular output section map.
The actual output lies in a fixed compact physical section window $J_T$.
The first action derivative of the landing, composed with this same
$\Xi$, satisfies
\[
 \begin{split}
 |\mathcal D\log G|
 &\le C\left\{
 \frac{\tau}{\sqrt B\,r^2}
       \left(1+\left|\log\frac{\tau}{\sqrt B\,r}\right|\right)
       +s_0\kappa\tau+s_0^2\mathcal P(r)\right\}\\
 &\le C\{\tau(1+|\log\tau|)+s_0\}.
 \end{split}
 \tag{GC6}
\]
The ordinary finite closing has $f_\theta(0)=0$, $f_\theta'>c_f>0$,
and uniform ordinary $C^3$ bounds on a buffered interval containing
$mJ_T$ for $0\le m\le\bar m$. It is the original closing after the
output map $d_\theta$ already retained in $T$.
The physical affine displacement is exactly
\[
 F_\theta(h)=h-a_\theta-f_\theta(mT_\theta(h)),\qquad
 T_\theta(h)=T(r(P(h))).
 \tag{GC7}
\]

\begin{proposition}[Convexity and zero count in the growing range]
\label{p2:prop-growing-physical-count}
On the whole counted growing interval, there are common constants
$C_T,C_{\rm cl}$ and a multiplier bound $m_0>0$, selected after the
independent coefficient bounds, such that
\[
 |T|\le C_T,\qquad 0<\mathcal ET\le C_T,\qquad
 |\mathcal E^2T|\le C_T .
 \tag{GC8}
\]
For every actual through input in that interval, whether or not
$F_\theta=0$,
\[
 \begin{gathered}
 \left|\mathcal E^j\{f_\theta(mT)-f_\theta(0)\}\right|
       \le C_{\rm cl}m,\\
 \left|\mathcal E^j\{\log f_\theta'(mT)-\log f_\theta'(0)\}\right|
       \le C_{\rm cl}m,\qquad j=0,1,2,\quad 0\le m\le m_0 .
 \end{gathered}
 \tag{GC9}
\]
For $0<m\le m_0$, the complete physical displacement is strictly convex:
\[
 F_\theta''(h)>0\qquad\hbox{at every point of the growing interval}.
 \tag{GC10}
\]
Consequently it has at most one critical point, that point is simple
and upward, and it has at most two isolated zeros counted with
multiplicity. At $m=0$, $F_\theta'(h)=1$.
The constants are uniform on the full prescribed passive coefficient
family and both coefficient graphs, with the reciprocal estimates
uniform for physical $j>0$ approaching its one-sided boundary.
\end{proposition}

\begin{proof}
\emph{Bounds on the full output and its action derivatives.}
Put $L(r)=-T_r=a_{\rm grow}rG(r)>0$ and
\[
 d_r=\min\{1,(\widehat r_{\rm ge}
              -\widehat r_{\rm ge}^{\rm out})/2\}>0 .
\]
The original lower buffer places every $[r-d_r,r]$ in (GC1).
Since $\mathfrak d^{-1}$ is bounded, (GC6) gives
\[
 |\partial_r\log L|
 \le r^{-1}+\mathfrak d^{-1}|\mathcal D\log G|\le C_L .
\]
Let $D_T$ be the diameter of $J_T$. Using the actual output values,
\[
 D_T\ge T(r-d_r)-T(r)
      =\int_{r-d_r}^rL(v)\,dv
      \ge d_re^{-C_Ld_r}L(r).
 \tag{GC11}
\]
Thus $L\le C$ without dividing by $a_{\rm grow}$, $\Omega$ or a raw gain.
Equation (GC4) now gives
\[
 \mathcal ET=e_s\mathfrak d\,Ln>0,\qquad
 \mathcal E\log(\mathcal ET)
 =-e_sn\{\mathcal D\log\mathfrak d+\mathcal D\log L\}
       +\frac{\mathcal E\zeta}{n},\qquad
 \mathcal D\log L=\mathfrak d/r+\mathcal D\log G .
 \tag{GC12}
\]
Every term on the right is bounded. Hence $\mathcal ET\le C$ and
$|\mathcal E^2T|\le C$; $|T|$ is bounded by the fixed compact $J_T$.
The reference value of the output is included in this bound, so (GC8)
applies at the actual argument of the closing map.

\emph{Composition with the finite closing.}
The full translated $T$ and its first two action rows satisfy (GC8)
on the buffered actual image. The ordinary $C^3$ closing has positive
derivative units there. Lemma~\ref{phys:lem-c3-adapters}, specifically
\eqref{phys:eq-absolute-closing-two}, therefore gives (GC9) for all
through inputs, with constants fixed before $m$.

\emph{Strict convexity in the actual affine input.}
Before finite closing, the exact logarithmic slope is
\[
 \begin{split}
 H_0
 &:=\mathcal E\log
       \frac{\mathfrak j(P)}{e_sG(r)\mathcal P(r)n(P,r)}\\
 &=\mathfrak j'(P)+e_sn
       \left\{\frac{\mathcal P'(r)}r+\mathcal D\log G(r)\right\}
       -\frac{\mathcal E\zeta}{n}.
 \end{split}
 \tag{GC14}
\]
This identity follows from $\mathcal Er=-e_s\mathfrak d n$;
the derivative of the matched quotient is included.
By (GC4), (GC6) and $\mathcal P'>0$,
\[
 H_0\ge 3/4-C\{\tau(1+|\log\tau|)+s_0\}>1/2
 \tag{GC15}
\]
after common independent $\tau,\kappa$ bounds.
Here $s_0=\kappa\sqrt B$, with $B$ a positive unit.
The original passive coefficient set and joint auxiliary tuple are
unchanged.

Since $\mathcal ET=\mathfrak j T_P$ and (GC5) gives
$T_P=a_{\rm grow}e_sG\mathcal P n/\mathfrak j>0$, define
\[
 \mathscr R=\frac{s_{\rm in}}{mf_\theta'(mT)T_P}>0 .
\]
Every factor in its numerator, $a_{\rm grow}$ and $e_s$ is
coefficient-only. Its complete logarithmic row is therefore
\[
 H_m:=\mathcal E\log\mathscr R
    =H_0-m\frac{f_\theta''(mT)}{f_\theta'(mT)}\mathcal ET .
 \tag{GC16}
\]
The absolute output estimate (GC8) makes the closing term at most
$C_m m$, with a common $C_m$ fixed before the multiplier.
Choose last
$m_0\le\min\{\bar m,1,1/[4(1+C_m)]\}$. Then $H_m>1/4$.

Restore the physical input using (GC2). Direct differentiation of the
same displacement (GC7) gives at every through point
\[
 F_\theta'(h)=1-\frac1{\mathscr R},\qquad
 F_\theta''(h)=\frac{\mathscr R_h}{\mathscr R^2}
       =\frac{H_m}{s_{\rm in}\mathfrak j(P)\mathscr R}>0 .
 \tag{GC17}
\]
The two raw gains remain in this expression, so its sign holds even
when either gain is arbitrarily small.
Translation to $h-a_\theta$ leaves these derivatives unchanged.

Lemma~\ref{sharp:upper-critical-count} now gives the critical and
zero counts, with multiplicity; (GC17) excludes an identity interval.

The bound (GC6) uses $r\ge R_0$ and
$s_0^2\mathcal P(r)\le Cs_0$; its logarithm contains no $j$.
The bounds on $\kappa,\tau$ are independent, uniformly on the
prescribed compact $j$ ranges.
\end{proof}

\section{Finite height and fixed exit}
\subsection{The last two passage regions and their physical curvature}

Bounded amplified height is described by two small-power section maps.
For larger height we instead use the action coordinate on a fixed
negative section. We construct both maps on overlapping domains before
comparing their derivatives in the physical input.

Retain the single joint auxiliary choice, its original physical sections
and nested windows, the weak cuts, and the original coefficient family.
Let \(I_\ell\) be the positive-height first-lip domain in the original
through interval. Write
\[
 P={\frac{h-x_b}{s_{\rm in}}},\qquad
 s_{\rm in}=\Omega a^{2-\lambda_\gamma}U_{\rm in}>0,\qquad
 \Omega=r_s^{\,2-\lambda_\gamma},
 \qquad q_s={\frac{\beta}{\chi}}=\kappa Q_d ,
 \tag{FVFE1}
\]
where \(h\) is actual incoming height, \(Q_d,U_{\rm in}\) are positive
units, and no positive lower bound for either raw first gain is imposed.
The variable \(P\) is the canonical raw input of the joint entrance
construction; it is distinct from the saddle ratio \(q_s\).
Use the unchanged raw middle coordinate
\[
 V={\frac{y_+(h)}{\varepsilon_+}},\qquad
 \varepsilon_+={\frac{\ell_0h_-}{a^2J_m(u_-)}},\qquad
 \sigma=s_0r_\ell .
 \tag{FVFE2}
\]
Here \(y_+\) is the actual height \(w/a^2\) at \(u=u_+\).
Every phase derivative fixes all coefficients and all gains.

\begin{proposition}[Physical passage and convexity in the last two ranges]
\label{p2:prop-finite-v-fe-actual-budget}
For the joint auxiliary choice, followed by independent active
coefficient bounds, put
\[
 J_{34}=I_\ell\cap\{\sigma\ge\sigma_1\},\qquad
 D=J_{34}\cap\{0<V\le V_1\},\qquad
 E=J_{34}\cap\{V\ge1\},\qquad V_1>1 .
 \tag{FVFE3}
\]
These are intervals in the same actual incoming height, with every
displayed divider equality included. Each lies in the prescribed
enlarged domain of its corresponding actual transition, and
\(D\cup E=J_{34}\). Every input in these sets, whether
closed or nonclosed, has its original raw trace in
\(0<P<q_{\rm ref}<\widehat p_*<\widehat p_0\).
The finite-\(V\) map uses the original early negative cut on
the fixed enlarged interval \(0<V\le V_1+1\). For larger amplified
heights, the independently constructed fixed-exit map uses the buffered
shooting chord on \(u=u_c\). The two
constructions agree on \(1\le V\le V_1\) by original-time
uniqueness, so their union is the same physical transition on \(J_{34}\).

For the original return write
\[
 \mathcal F(h):=F(h-a_\theta)
       =h-a_\theta-f_\theta(mT(h)).
 \tag{FVFE4}
\]
The map \(T\) includes the actual negative output adapter. For every
\(0<m\le m_0\),
\[
 \mathcal F''(h)>0\qquad(h\in J_{34}).
 \tag{FVFE5}
\]
Thus the whole combined interval has at most one simple upward
critical point and at most two isolated zeros with multiplicity.
At $m=0$, $\mathcal F'=1$. The whole carrier is counted separately
in Proposition~\ref{p2:prop-physical-degree-count}.
\end{proposition}

\begin{proof}
\emph{The shooting chord and its phase coordinate.}
Part (i) of Proposition~\ref{p2:prop-unified-carrier-cover} realizes
$\overline{\mathcal J_c^{\sharp\sharp}}\Subset(0,u_c^2/2)$ as
unique through orbits, identified with the backward shots by uniqueness.
Proposition~\ref{fx:prop-inverse-phase} applies to this entire buffered
chord: its scalar first-exit construction is independent of the
positive raw-section dictionary. Its extracted phase
$\Psi_\theta(W)$ converges in $C^2$ to
$\Gamma_c\circ\mathcal W^{-1}$, whose derivative is
$1/H_{0,c}(r)>0$. Hence its increasing inverse
$W=\mathcal H_\theta(\Theta)$ has bounded first two derivatives
and a positive derivative unit on
\[
 \Theta_-\le\Theta\le\Gamma_- -\Delta,\qquad \Theta_-<0,
 \tag{FVFE8}
\]
strictly inside its inverse image. The negative buffer comes from the
original enlarged chord around $u_-$; the reciprocal functions give
the same uniform gaps.

\emph{Continuation of the shooting chord to the raw input.}
The exact positive invariant, including its local unit functions,
is
\[
 \Phi_\theta(P)=\{z_+y_+(P)\}^{\chi},\qquad
 \Phi_\theta(P)=P U_\theta(P),\qquad
 \Phi_\theta'>c,\qquad U_\theta^{\pm1}\le C.
 \tag{FVFE9}
\]
Use the reference $P_{\rm ref}=a_{\rm in}$ of
\eqref{deg:degree-fixed-exit-gap}:
\[
 \Phi_\theta(P_{\rm ref})=(z_+\beta\varepsilon_+)^\chi,\qquad
 q_s\log\Phi_\theta(P_{\rm ref})=-\Gamma_-+o(1).
 \tag{FVFE10}
\]
The exact identity \eqref{p2:reference-telescope} proves this limit
on both graphs: on the reciprocal graph use its prescribed $d,j$
in place of $\mathcal Rc,1/c$. All factors remain regular at $j=0$,
so its error bound is unchanged and no $\beta\log j$ occurs.
This reference corresponds to $V=\beta$. On the compact shooting chord the limiting
gap from \(\Gamma_-\) has a fixed positive minimum \(d_{\rm fx}\).
The scalar shooting estimate and (FVFE10) give, after shrink,
\[
 \beta\log V\le\Gamma_- -d_{\rm fx}/2,\qquad
 \beta\log y_+\le-d_{\rm fx}/3 .
 \tag{FVFE11}
\]
Since \(\beta\le C\chi\), this implies \(y_+<\chi^2\).
Thus (FVFE9) can be inverted on the original invariant chord, and
\[
 \Phi_\theta(P)\le C e^{-d_{\rm fx}/(3q_s)},\qquad
 0<P\le e^{-d_{\rm fx}/(4q_s)}<\widehat p_0
 \tag{FVFE12}
\]
after a further independent \(\kappa\) bound.
In particular the physical distance from the stopped trace satisfies
\(\sup(h-x_b)\le C e^{-d_{\rm fx}/(4q_s)}\to0\) on the
entire shooting chord, since \(s_{\rm in}\le C\).
The two local points have the same invariant in the same larger
coordinate box. Positivity and monotonicity along that invariant level
give the intervening orbit segment; the independently actual raw
continuation identifies its raw trace with the shot through orbit.
This places the full chord within the independently selected fixed cutoff
$\widehat p_0$, before any zero selection.

\emph{The bounded-\(V\) passage.}
Lemma~\ref{lem:actual-dulac-matching} independently continues the
entire larger interval $0\le V\le V_1+1$ to the early cut
$u_0=-1+\eta_0>u_-$. Equations~\eqref{eq:early-saddle-smallness}
and~\eqref{eq:early-middle-identity} give its positive invariant
at most $Ce^{-c_0/(2q_s)}$. The inversion used in (FVFE12) places the whole interval strictly
inside the raw cutoff.
The assigned set \(D\) is its restriction to \(0<V\le V_1\).
On the negative early cut, the lemma's exact invariant is
\[
 I_-(w)=\alpha_\theta R,\qquad
 R=\beta^{-\beta}C(v)^\beta,\qquad
 v=\widehat V=V+o_{C^3}(\chi).                 \tag{FVFE14}
\]
Its fixed cut choice puts the full $R\le1+o(1)$ interval strictly
inside the outer positive-height chord. At $R=0$ the output is the
stopped height $w_->0$. Thus both actual tails, ordinary $C^3$
observers and finite closing are on their original buffered domains,
with their exact values and positive first gains retained.

\emph{The domains of the two passage maps.}
The joint entrance supplies the actual product reference
\(r_\ell=\sigma_1/s_0\) after the active shrink.
Direct landing followed by the independent inverse shooting phase
gives
\[
 \Theta(h_{\sigma_1})\longrightarrow\Gamma_c(-\sigma_1),
 \qquad
 \Gamma_- -\Gamma_c(-\sigma)
      =\log{\frac{c+\sigma}{c(1-\sigma)}} .
 \tag{FVFE15}
\]
On the reciprocal chart the gap is
\(\log((1+j\sigma)/(1-\sigma))\), uniformly positive on the fixed
product buffer. The retained choice of \(\Delta\) and its strict
buffer imply
\[
 0<g_0\le\Theta(h_{\sigma_1})\le\Gamma_- -2\Delta .
 \tag{FVFE16}
\]
The actual orders \(r_{\ell,h}<0\) and \(V_h>0\) imply that every
member of \(E\) has \(0\le\Theta=\beta\log V\le\Theta(h_{\sigma_1})\).
Thus the entire fourth range lies in the domain (FVFE8) of the inverse
phase map, including inputs whose orbits are not closed. The equality \(V=1\) has the
negative-action buffer there. Since
\(\mathcal H_\theta([0,\Gamma_- -2\Delta])
\Subset\mathcal J_c^\sharp\Subset\mathcal J_c^{\sharp\sharp}\)
uniformly after shrink, its whole closed output-height subset has
a strict actual shooting buffer.
Decrease the independent coefficient bounds further so that
\(\beta\log(V_1+1)<g_0/2\). Then
\(V(h_{\sigma_1})=\exp(\Theta(h_{\sigma_1})/\beta)>V_1+1\),
so the third range retains the full buffered finite-\(V\) passage.
The two constructions identify the same raw inputs and outputs on \(1\le V\le V_1\) by original-time uniqueness.
No physical input is relabelled at \(V=1\).
Their actual \(\sigma\ge\sigma_1\) domain has
\(r_\ell\ge\sigma_1/s_0>r_{\rm ge}\), so the joint entrance's full
first-hit order supplies the sharper raw bound
\(P<q_{\rm ref}<\widehat p_*<\widehat p_0\).

\emph{Physical convexity for bounded \(V\).}
For \(m>0\), the actual-section construction gives exactly
\[
 \begin{gathered}
 h=x_b+P_*K(v),\qquad
 f_\theta(mT(h))=b_*+a_*H(v),\\
 K(v)=\Psi(v^\chi),\quad
 H(v)=D(\beta^{-\beta}C(v)^\beta),\qquad P_*,a_*>0 .
 \end{gathered}\tag{FVFE17}
\]
Here all regular tails and finite closing are already included.
Put \(e_C=vC'/C>0\). Lemma~\ref{lem:actual-dulac-matching}
supplies the stronger normalized estimates below, with \(\Sigma_+\)
the coefficient-only positive invariant scale in its construction:
\[
 \begin{gathered}
 e_C=1+O(\epsilon_\theta),\qquad
 \left|{\frac{ve_C'}{e_C}}\right|\le C\epsilon_\theta,\qquad
 |D''/D'|\le C,\\
 \sup_t|t\Psi''(t)/\Psi'(t)|\le C\Sigma_+,\qquad
 \epsilon_\theta\le C\beta^{-N}e^{-c_0/(3\beta)},\qquad
 \Sigma_+\le Ce^{-c_0/(2q_s)} .
 \end{gathered}
\]
The exact Wronskian row, with \(R=\beta^{-\beta}C(v)^\beta\), is
\[
 \begin{split}
 \mathcal Q(v)
 &=v\left({\frac{K''}{K'}}-{\frac{H''}{H'}}\right)\\
 &=\chi-\beta e_C-{\frac{ve_C'}{e_C}}
       +\chi v^\chi{\frac{\Psi''}{\Psi'}}
       -\beta e_C R{\frac{D''}{D'}}
 \ge\chi/2>0 .
 \end{split}
 \tag{FVFE18}
\]
On this interval \(R\le1+o(1)\). With
\(q_s=\beta/\chi=\kappa Q_d\longrightarrow0\),
the displayed estimates give the relative bound
\[
 \frac{\mathcal Q(v)}{\chi}
 \ge1-Cq_s-C\frac{\epsilon_\theta}{\chi}-C\Sigma_+
 \ge\frac12 .
\]
Indeed, \(\beta=q_s\chi\) implies
\[
 \frac{\epsilon_\theta}{\chi}
 \le Cq_s\beta^{-N-1}e^{-c_0/(3\beta)}\longrightarrow0.
\]
The input-section error has weight \(\chi\Sigma_+\), whereas
the bounded output and closing error has weight \(\beta\).
Thus independent bounds on \(\kappa,\tau\) give this sign for
every \(v>0\), without extending the small-power compositions
in ordinary \(C^3\) to the excluded stopped endpoint.
Since \(dh/dv=P_*K'>0\), differentiation of the exact physical
displacement, without multiplying it by a nonconstant factor, yields
\[
 \mathcal F''(h)=
 {\frac{a_*}{P_*^2}}{\frac{H'}{(K')^2}}{\frac{\mathcal Q(v)}{v}}>0 .
 \tag{FVFE19}
\]

\emph{Physical convexity on the fixed-exit interval.}
The assigned set \(E\) has
\(0\le\Theta\le\Theta(h_{\sigma_1})\), strictly inside the
larger buffered action interval (FVFE8).
Use the common actual map \eqref{deg:fixed-exit-actual-output},
$\Xi(W)=t_b+a_{\rm out}D^n(W)$, $a_{\rm out}=\Omega U_{\rm out}>0$,
on its whole buffered height chord. Its ordinary $C^3$ normalized
bounds and $0<\Omega\le C$ retain the full output in the original
raw and closing windows. For the fixed closing buffer $a_{\rm cl}>0$,
impose on the whole chord
\[
 m_0\left\{|t_b|+
 a_{\rm out}\sup_{\mathcal J_c^{\sharp\sharp}}|D^n|\right\}
 <a_{\rm cl}/2 .
 \tag{FVFE20a}
\]
The braces are uniformly bounded on the original output window,
so this selects $m$ independently of the first gain.

Put \(\widehat H=D^n\circ\mathcal H_\theta\). The exact raw formulas are
\[
 T(h)=t_b+a_{\rm out}\widehat H(\Theta(P)),\qquad
 \Theta(P)=\beta\log\beta+
       q_s\log{\frac{\Phi_\theta(P)}{\Phi_\theta(P_{\rm ref})}}.
 \tag{FVFE21}
\]
Writing \(p=h-x_b=s_{\rm in}P\), retain the raw gain in the two rows
\[
 0<pT_h=
 a_{\rm out}\widehat H' q_s{\frac{P\Phi_\theta'}{\Phi_\theta}}
       \le C a_{\rm out}q_s ,
 \tag{FVFE22}
\]
\[
 p{\frac{T_{hh}}{T_h}}
 =P{\frac{\Phi_\theta''}{\Phi_\theta'}}
       -P{\frac{\Phi_\theta'}{\Phi_\theta}}
       +q_sP{\frac{\Phi_\theta'}{\Phi_\theta}}
                    {\frac{\widehat H''}{\widehat H'}}
 =-1+O(P+q_s)\le-\tfrac12 .
 \tag{FVFE23}
\]
The ordinary $C^3$ closing has derivative units and bounded
pre-Schwarzian on the whole sampled interval. (FVFE22)--(FVFE23) give
\[
 p\left\{{\frac{T_{hh}}{T_h}}
       +m{\frac{f_\theta''(mT)}{f_\theta'(mT)}}T_h\right\}
 \le-\tfrac12+Cma_{\rm out}q_s<-\tfrac14 .
 \tag{FVFE24}
\]
Therefore
\[
 \mathcal F''(h)=
 -mf_\theta'(mT)T_h
 \left\{{\frac{T_{hh}}{T_h}}
       +m{\frac{f_\theta''(mT)}{f_\theta'(mT)}}T_h\right\}>0 .
 \tag{FVFE25}
\]
No lower bound for \(s_{\rm in}\), \(\Omega\), or \(T_h\) was used.

The positive input order makes \(J_{34}\) one interval.
The two convexity estimates concern the identical physical \(F\)
on an overlapping cover, so (FVFE5) holds on their union.
Lemma~\ref{sharp:upper-critical-count} gives the critical and zero
counts with multiplicity, also for every finite collection approaching
the excluded stopped boundary.
The endpoints at the phase boundaries lie inside the enlarged domains
of their respective passage maps.
At \(m=0\) the original completed displacement has derivative one.
The combined reference logarithms, normalized section maps and
first-cocycle powers above are regular at the reciprocal face; none introduces a
\(\beta\log j\) or \(\tau\log j\) restriction. Actual physical maps
are asserted for \(j>0\), with \(j=0\) a one-sided normalized limit.
\end{proof}

\section{Counting zeros on the complete degree-loss interval}

Each local calculation either gives strict convexity of the physical
displacement or excludes critical points by a positive derivative.
Because the local maps agree on overlapping intervals, the same
alternative holds at every point of the connected through interval.

\begin{proposition}[Physical counting on the degree-loss carrier]
\label{p2:prop-physical-degree-count}
Fix the compact positive-unit ranges of the degree-loss graph, or
the reciprocal graph with $0<j\le j_0$. Use the joint auxiliary
choice of the growing entrance and its actual gates. There are
independent active coefficient bounds, a common raw-height cap and
a last finite-closing multiplier bound with the following property.
On the resulting connected interval $I_\theta$ of actual through
inputs on the original section, the complete affine displacement
\[
 F_\theta(x)=x-f_\theta\bigl(m\mathfrak T_\theta(a_\theta+x)\bigr),
 \qquad x=h-a_\theta,
 \tag{PC1}
\]
has at most one critical point. Every critical point is a simple
upward crossing of $F_\theta'$. Consequently it has at most two
isolated zeros, counted with multiplicity, and no identity interval.
The interval $I_\theta$ consists of the same through word in the
original nested windows, with positive passage denominator and
$0<z(\gamma_s)<z_{\rm cap}$ for one fixed physical height cap.
The assertion holds at every actual input of $I_\theta$, without
assuming its orbit is closed. The original field, marked sections,
nested raw windows, $L,h_L$ and closing map define this same return.
Here $\mathfrak T_\theta=d_\theta\circ T_\theta^\infty$ is the
original transported output and $f_\theta=C_\theta-C_\theta(0)$.
Thus $f_\theta(0)=0$ and at $m=0$ the displacement is $F_\theta(x)=x$.
No physical passage at $j=0$ or at a stopped boundary is asserted.
\end{proposition}

\begin{proof}
Proposition~\ref{p2:prop-unified-carrier-cover} supplies the closed
low cover $B_\theta\cup G_\theta\cup D_\theta\cup E_\theta$ and
its overlapping high continuation on the connected actual interval
$I_\theta$. It includes every divider equality and every nonroot input;
only the stopped endpoint is excluded.
All restrictions have the same affine displacement (PC1).

On the annular part of $B_\theta$,
Lemma~\ref{bd:lem-annulus} gives $\Pi''<0$ after the ordinary
$C^3$ divided closing is included. Its principal margin is
$\lambda_1/(2x_{\max}R_{\rm BD})$; the independent active bounds
make the errors smaller than half this margin.
Lemma~\ref{p2:lem-bounded-source-budget} gives the same sign on
negative source inputs and $F'>1/2$ on the compact, positive and
high source pieces. Propositions~\ref{p2:prop-growing-physical-count}
and~\ref{p2:prop-finite-v-fe-actual-budget} give $F''>0$ on
$G_\theta$ and on the whole combined $D_\theta\cup E_\theta$.
These signs are in the original affine input and include both exact
raw gains and the translated output inside the closing map.

Take the finite intersection of the independent active bounds, then
one multiplier bound for all derivative estimates and the full absolute
output envelope in the original closing buffer. Every critical point
of (PC1) consequently has $F''>0$, including those on the dividers.
Lemma~\ref{sharp:upper-critical-count} gives at most one simple
upward critical point and the two-zero bound with multiplicity.
An identity interval would contain a critical point with $F''=0$.
Section~\ref{sec:common-cover-conclusion} assembles the coefficient
cover and Hausdorff capture.
\end{proof}

\endgroup

The growing-action estimates can now be applied on the physical
interval. They control the matched quotient, which determines
curvature before finite closing:
\begin{equation}\label{deg:degree-growing-elasticity}
 \mathcal E\log\frac{\mathfrak j(P)}
    {e_sG(r)\mathcal P(r)n(P,r)}>\frac12,
 \qquad \mathcal E=X\partial_X.
\end{equation}
This is \textnormal{(GC14)--(GC15)}. To compose with the finite
closing, we also need the full output and its first two action derivatives:
\begin{equation}\label{deg:growing-actual-output-action}
 |T|+|\mathcal E T|+|\mathcal E^2T|\le C,
 \qquad \mathcal E=X\partial_X,\quad \mathcal E T>0.
\end{equation}
Estimate \textnormal{(GC8)} gives this bound on the buffered growing
image, and \textnormal{(GC9)} transfers it through the finite closing.


\subsection{A direct Jacobian estimate}
\label{subsec:deg-raw-jacobian}

In the following two sign ranges, the graph Jacobian itself controls the
complete derivative. An exact logarithmic derivative separates the
endpoint factor from an interior term whose integral has a uniform upper
bound, so no action compactification is needed.

\begin{lemma}\label{deg:lem-raw-jacobian}
Assume either
\(\alpha _0\le0,\alpha _1\le0\), or
\(\alpha _0=-s^2<0,\alpha _1=sp>0,\gamma>0\) and
\[
 {s\over\gamma^{3/2}}\le k_0,\qquad
 {p\over\sqrt\gamma}\le\pi_0<\sqrt {12}.
\]
On the physical input interval selected by Lemma~\ref{lem:loc-capture},
the central transition has a uniformly bounded positive derivative.
After the final choice of the finite saddle--node multiplier, the
complete return is a contraction.
\end{lemma}

\begin{proof}
Put \(\xi=v-\gamma\), the local spatial coordinate of
\eqref{deg:central-field}, and set
\[
 M=-P(\gamma+\xi)
 =-\alpha _0-\alpha _1\xi+3\gamma \xi^2+\xi^3,\qquad
 \mathscr D=M+\ell z=-\dot v.
\]
Along a through orbit, regarded as a graph \(z=z(v)\), write
\[
 f={z\{q_\gamma \xi+\xi^2+Cvz\}\over\mathscr D},\qquad
 \nu={6\gamma\over q_\gamma},\qquad
 N=-\alpha _1+(3-\nu)\xi^2.
\]
Direct differentiation gives
\begin{equation}
 f_z={d\over dv}\log {z^{1+\nu}\over\mathscr D}
       +{N+\{\ell'(v)+(1-\nu)Cv\}z\over\mathscr D}.
 \label{deg:raw-log-jacobian}
\end{equation}
If \(\alpha _1\le0\), then \(N\ge|\alpha _1|+2\xi^2\) after shrinking
\(|\gamma|\).  Let \(\kappa=\frac12\inf\ell\).  On
\(\mathscr D\ge\kappa z\), the last quotient in
\eqref{deg:raw-log-jacobian} is bounded below.  On the complement,
\(M<0\),
\[
 z<{|M|\over\ell-\kappa},\qquad
 |M|\le C(\delta+|\gamma|)N.
\]
Choosing the fixed raw width \(\delta\) and then the coefficient
neighbourhood makes the numerator at least \(N/2\) there.  Thus its
integral in the physical direction from the positive to the negative
section is bounded above.

In the second case put \(\xi=(s/\sqrt\gamma)\widehat x\), with
\(\pi=p/\sqrt\gamma\) and \(k=s/\gamma^{3/2}\).  Then
\[
 M=s^2\{1-\pi\widehat x+3\widehat x^2+k\widehat x^3\}.
\]
Where \(N<0\), the quadratic identity
\[
 1-\pi\widehat x+3\widehat x^2
 =3\left(\widehat x-{\pi\over6}\right)^2+1-{\pi^2\over12}
\]
gives a positive multiple of \(1+\widehat x^2\).  On \(N<0\) one has
\(\widehat x^2\le C\pi/k\), hence
\(|k\widehat x|\le C\sqrt{\pi_0k_0}\).
The cubic term is therefore absorbed by decreasing \(k_0\).
Since \(-N\le sp\), the change of variable cancels the prefactor
\(sp\,s/(s^2\sqrt\gamma)=\pi\le\pi_0\), and hence
\[
 \int_{\{N<0\}}{-N\over\mathscr D}\,d\xi
 \le C\int_{\mathbb R}{d\widehat x\over1+\widehat x^2}<\infty .
\]
Where \(N\ge0\), the preceding two-denominator argument applies.  Indeed
\(M<0\) can then occur only at \(\xi=-u\), \(u>3\gamma\), and
\(N\ge cu^2\), \(|M|\le u^3\le\delta u^2\).

Integration of \eqref{deg:raw-log-jacobian} yields
\[
 0<T_\theta'\le
 C{(z^{1+\nu}/\mathscr D)_-\over
       (z^{1+\nu}/\mathscr D)_+}\le C_0.
\]
The two endpoints lie in the fixed positive-height windows, which bound
the quotient in this estimate.
The estimate holds on the full interval
\(\mathcal T_\theta\cap\operatorname{int}K_+^\sharp
\cap T_\theta^{-1}(\operatorname{int}K_-^\sharp)\), including all
intervening captured inputs. The bounded divided closing of
Lemma~\ref{phys:lem-c3-adapters}, followed by a final multiplier
shrink, gives the contraction.
\end{proof}

\section{Nonnegative opening: determinant and trace faces}
\label{sec:nonnegative-opening}
\label{subsec:deg-nonnegative}

The high-height estimate reduces the remaining problem to bounded
rescaled heights. At their lower boundary, the determinant chart has a
hyperbolic saddle, while the trace chart includes its saddle--node limit.
The determinant count requires the absolute first variation, including
both tail gains. The trace count also requires the additional outgoing
saddle. We treat
both passages in the affine input and with the divided closing of
Lemma~\ref{lem:closing}.

\begin{lemma}[Finite coefficient--phase atlas]
\label{deg:lem-positive-atlas}
Fix the signed high-height constant \(K_H\) and put \(K=2K_H\).
For \(\alpha_0\ge0\), below \(z(\gamma)\le K(\alpha_0+\alpha_1^2)\),
the nonzero coefficients have the overlapping descriptions
\begin{align}
 &\alpha_0=r_D^2>0,\qquad \alpha_1=r_DA_D,\qquad |A_D|\le P_D,
 \label{deg:determinant-chart}\\
 &\alpha_0=r_L^2a_L,\qquad \alpha_1=\sigma r_L,\qquad
 0\le a_L\le a_*,\quad \sigma\in\{-1,1\},
 \label{deg:trace-chart}
\end{align}
where \(P_D>2/\sqrt{a_*}\). The determinant and signed-trace
receivers below cover every low-height through input, including regular
continuations, by overlapping stopped and compact strict intervals.
Their cuts precede a displacement zero. If \(\alpha_0=\alpha_1=0\),
every positive-height through input is in the signed high-height range.
\end{lemma}

\begin{proof}
Use \eqref{deg:determinant-chart} when
\(|\alpha_1|\le P_D\sqrt{\alpha_0}\); otherwise set
\(r_L=|\alpha_1|\), so \(a_L<a_*/4\). The face
\(\alpha_0=0\), \(\alpha_1\ne0\), has \(a_L=0\).
These are whole-coefficient alternatives with a strict overlap and no
restriction on either \(\gamma/r_D\) or \(\gamma/r_L\).
Their bounded normalized heights are respectively \(K(1+P_D^2)\)
and \(K(1+a_*)\). Their physical phase covers are constructed in
Lemma~\ref{deg:lem-determinant-preparation} and
Proposition~\ref{deg:prop-stopped-trace}. For the trace limiting field
\(X'=a_L+\sigma X-Y\), \(Y'=-XY\), at \(a_L=0\) its centre
manifold has \(Y=\sigma X+O(X^2)\), hence
\(\dot Y=-\sigma X^2+O(X^3)\); its weak quadratic coefficient is
nonzero. The double-zero coefficient face has no positive low-height
input, and is covered by Corollary~\ref{rew:cor-high-derivative-signed}.
\end{proof}

\subsection{Preparation on the actual determinant sections}

\begin{lemma}[Determinant section preparation]
\label{deg:lem-determinant-preparation}
Fix \(P_D<\infty\), the passive compact coefficient block, and the
original buffered windows. For independent sufficiently small
\(r=r_D>0\) and \(|\gamma|\), \(\alpha_0=r^2\),
\(\alpha_1=rp\), \(|p|\le P_D\), choose a fixed positive unscaled
raw-height cap \(z_*\). On the actual interval
\begin{equation}\label{deg:det-actual-receiver}
 \mathcal I_\theta=\{x:a_\theta+x\in J_\theta^{\rm out},\quad
       0<H_\theta(x)<z_*\},\qquad
 H_\theta(x)=z_\theta(\gamma;a_\theta+x),
\end{equation}
the stopped transition has the exact composition
\eqref{sharp:upper-det-composition}, with the two absolute positive
gains \eqref{sharp:upper-det-scales} and actual cut gap
\eqref{deg:det-actual-gap}. The normalized end maps and their inverses
have common ordinary \(C^3\) bounds and derivative units. The incoming
and outgoing local chords are independent. The compact-spectrum saddle
estimate \eqref{sharp:upper-det-dulac} and the complete absolute
variation \eqref{sharp:upper-det-absolute} hold, including saddle ratio
one. Every phase derivative fixes all coefficients and the actual \(m\).

The low stopped and compact strict receivers, their full regular
continuations, and the signed high receiver cover \(\mathcal I_\theta\)
as in \eqref{deg:det-height-cover}. On the strict and high pieces the
full \(\Pi'\), including divided closing, is uniformly bounded.
Retain intervening inputs, equality dividers and all original clips;
exclude natural stopped endpoints. The parameters \(m,r,\gamma\)
and the two local chord lengths are independent.
\end{lemma}

\begin{proof}
We construct the two section maps with their absolute gains before
estimating the stopped derivative. Keep \(x=h-a_\theta\), \(\Pi=\mathcal C_{\theta,m}\circ
T_\theta(a_\theta+\cdot)\), and all original finite-side factors.
The original domain is
\[
 J_\theta^{\rm out}=\mathcal T_\theta\cap\operatorname{int}K_+^\sharp
 \cap T_\theta^{-1}(\operatorname{int}K_-^\sharp)
 \cap\{\sigma_\theta<d_{\rm out}\}.
\]
Order of the raw passage and central height makes
\eqref{deg:det-actual-receiver} an interval, before capture is invoked.
With \(v=\gamma+rX\), \(z=r^2Y\), a positive time scaling gives
\begin{equation}\label{deg:det-field}
 \begin{split}
 X'&=1+pX-3\gamma X^2-rX^3-\ell(\gamma+rX)Y,\\
 Y'&=-Y\{q_\gamma X+rX^2+Cr(\gamma+rX)Y\},\qquad q_\gamma=1+2\gamma.
 \end{split}
\end{equation}
At \(r=\gamma=0\), the saddle \((0,1)\) has eigenvalues
\((p\pm\sqrt{p^2+4})/2\). Their uniform nonzero moduli give common
saddle boxes, transverse cuts \(X=\pm L_c\), and persistent branch
intersections for independent small \(r,\gamma\).
The incoming gap is exactly
\begin{equation}\label{deg:det-actual-gap}
 \xi_D=Y(L_c)-Y_b^+(L_c)>0.
\end{equation}
For a limiting through graph starting at \(Y(0)>1\), put
\(D=Y-1-pX\). The identity \(\dot D|_{D=0}=-XY\) excludes a
first exit from \(D>0\) on the negative half-axis; reverse time treats
the positive half-axis. Thus every such graph is strict with quadratic
tails. The gap is continuous, strictly increasing, and tends uniformly
to zero at the stopped boundary. A fixed larger strict graph bounds the
cut heights by order. On any bounded central-height range, the complement
\(\xi_D\ge\xi_*/2\) therefore has a finite compact strict cover with
positive finite-cut height and speed margins.

Choose enlarged tail cuts \(\pm L_t\) and buffered auxiliary ports
with \(0<d_\epsilon<d_+\), joined exactly to the original sections.
For \(X=\epsilon R\), \(\epsilon=\pm1\), the outward equation is
exactly \eqref{eq:parabolic-tail} with \((a,s)=(1,r)\):
\begin{equation}\label{deg:det-tail-graph}
 Y_R=F^\epsilon(R,Y):=F_\epsilon(R,Y)|_{a=1,s=r}.
\end{equation}
The enlarged corridor has \(Y\) and its denominator comparable to
\(R^2\) up to \(d_\epsilon/r<d_+/r\). Choose
\(|\lambda_\gamma|<\nu_0<1\). The intrinsic tail estimates with
reference graph \(Y_*\) and first variation \(j\) give
\begin{equation}\label{deg:det-tail-action}
 F_Y=\lambda_\gamma/R+O(r+|p|R^{-2}+R^{-3+\nu_0}),\qquad
 j=(R/L_t)^{\lambda_\gamma}U_\epsilon,\quad U_\epsilon^{\pm1}\le C,
\end{equation}
\begin{equation}\label{deg:det-rational-derivatives}
 |F_{Y^h}|\le C_h\{rR^{-2h+2}+|\gamma|R^{-2h+1}
              +|p|R^{-2h}+R^{-2h-1}\},\qquad h=2,3.
\end{equation}
The first remainder has bounded integral to the actual endpoint.
Use the normalization of Lemma~\ref{lem:scalar-normalization},
\begin{equation}\label{deg:det-normalized-tail}
 \eta=(Y-Y_*)/j,\qquad
 \eta_R=N=\{F(R,Y_*+j\eta)-F(R,Y_*)-jF_Y(R,Y_*)\eta\}/j.
\end{equation}
Its changed kernel weight is
\begin{equation}\label{deg:det-normalized-kernel}
 \|N\|_{C^3}\le C\{rR^{-2+\nu_0}+|\gamma|R^{-3+\nu_0}
                      +P_DR^{-4+\nu_0}+R^{-5+\nu_0}\}.
\end{equation}
Choose \(L_t\) large, then shrink \(r,|\gamma|\) independently.
The lemma with \(k=2\) supplies enlarged chords and smaller inverse
images, remaining in the corridor, with
\begin{equation}\label{deg:det-phase-variations}
 \|\mathcal T_\epsilon\|_{C^3}
 +\|\mathcal T_\epsilon^{-1}\|_{C^3}\le C,\qquad
 0<c\le\mathcal T_\epsilon'\le C.
\end{equation}

At the actual port the raw centre and gain are
\(z_{\rm ref}=r^2Y_*(d_\epsilon/r)\), \(d=r^2j(d_\epsilon/r)>0\), with
\(d\le Cr^{2-\nu_0}\). The fixed-port proof of
Lemma~\ref{ann:p0-fixed-tail-ports} admits the whole raw images:
choose the auxiliary abscissa ratio from the window reserve, then shrink
the coefficients on the remaining compact bridges. That proof uses only
the common quadratic corridors just established. On these images
Lemma~\ref{phys:lem-c3-adapters}
gives the exact ordinary \(C^3\) normalization
\begin{equation}\label{deg:det-adapter-normalization}
 \widehat G(u)=\frac{G(z_{\rm ref}+du)-G(z_{\rm ref})}
 {G'(z_{\rm ref})d}=u+O_{C^3}(d),
\end{equation}
with common inverse bounds. Its divided-closing rows
\eqref{phys:eq-closing-jets} give, including \(m=0\),
\begin{equation}\label{deg:det-divided-closing}
 \|\mathcal C_{\theta,m}\|_{C^3}\le C,
 \qquad 0<c\le\mathcal C_{\theta,m}'\le C.
\end{equation}

Finite composition, with the positive-tail inverse on the incoming
side, now gives the two exact absolute embedding gains
\begin{equation}\label{sharp:upper-det-scales}
 \sigma_\pm=r^{2-\lambda_\gamma}V_\pm,\qquad
 \lambda_\gamma=\frac{12\gamma}{1+2\gamma},\qquad 0<c\le V_\pm\le C,
\end{equation}
and the complete stopped transition
\begin{equation}\label{sharp:upper-det-composition}
 x=x_b+\sigma_+\delta,\qquad
 \Pi(x)=y_b+\sigma_-O_\theta(D_\theta(I_\theta(\delta))),
 \quad I_\theta(\delta)=\xi_D,\quad 0<\delta\le\delta_*<1.
\end{equation}
The exact centres, cut factors, bridges and divided closing remain in
these maps. \(I,O\) vanish at zero, have derivative units, and are
ordinary \(C^3\) with bounded inverses. The incoming and outgoing saddle
chords have separate buffers. Their common power makes the quotient of
the two absolute gains a unit; the individual gains need no positive
lower bound.

The saddle of \eqref{deg:det-field} belongs to a compact analytic
family: its two rates have uniform positive moduli and
\(\wp\in[\wp_0,\wp_1]\Subset(0,\infty)\). Reduce the independent
incoming and outgoing chords separately to common compatible local
sections; their actual gap holonomies are ordinary \(C^3\) with
positive derivative units. Lemma~\ref{lem:uniform-positive-dulac}
therefore gives a local map \(cz^\wp e^{\mathcal R}\), with two
Euler remainder rows uniform through \(\wp=1\). Differentiation
replaces its remainder by
\(\Psi=\mathcal R+\log(1+z\mathcal R'/\wp)\).
The actual holonomies add \(O(z)+O(z^{\wp_0})\) to its value and
first Euler row. Thus, for fixed
\(0<\alpha<\min\{1,\wp_0\}\), the actual gap-to-gap map satisfies
\begin{equation}\label{sharp:upper-det-dulac}
 D_\theta'(z)=C_\theta\wp z^{\wp-1}e^{\Psi_\theta(z)},\qquad
 C_\theta^{\pm1}\le C,\qquad
 |\Psi_\theta|+|z\Psi_\theta'|\le Cz^\alpha(1+|\log z|).
\end{equation}
The estimate remains uniform at ratio one because it uses no division
by \(\wp-1\).

Put \(i_0=I'(0)>0\), \(o_0=O'(0)>0\). Differentiating the
complete composition in its original affine input gives exactly
\begin{equation}\label{sharp:upper-det-absolute}
 \begin{gathered}
 \Pi'(x_b+\sigma_+\delta)
 =\frac{\sigma_-}{\sigma_+}O'(D(I(\delta)))D'(I(\delta))I'(\delta)
 =\mathcal A_\theta\delta^{\wp-1}e^{R_\theta(\delta)},\\
 \mathcal A_\theta=\frac{\sigma_-}{\sigma_+}o_0C_\theta\wp i_0^\wp,
 \quad 0<c\le\mathcal A_\theta\le C,\qquad
 |R_\theta|+|\delta R_\theta'|\le C\delta^\alpha(1+|\log\delta|).
 \end{gathered}
\end{equation}
Here the full remainder is
\[
 R=\Psi(I(\delta))+(\wp-1)\log\frac{I(\delta)}{i_0\delta}
       +\log\frac{I'(\delta)}{i_0}
       +\log\frac{O'(D(I(\delta)))}{o_0}.
\]
The ordinary \(C^3\) bounds give \(I/(i_0\delta)=1+O(\delta)\)
in the required Euler rows, while \(D(I)\) and its first Euler row
are \(O(\delta^{\wp_0})\). Thus every term has the stated bound.
The common power in \eqref{sharp:upper-det-scales} cancels exactly,
without restricting \(\gamma\log r\); at \(\wp=1\) this full
remainder is still used.

Finally put \(S_D=r^2(1+p^2)\). The actual set cover is
\begin{equation}\label{deg:det-height-cover}
 \begin{gathered}
 \mathcal I_{\rm lo}=\{x\in\mathcal I_\theta:H_\theta\le2K_HS_D\},\qquad
 \mathcal I_{\rm hi}=\{x\in\mathcal I_\theta:H_\theta>K_HS_D\},\\
 \mathcal I_\theta=\mathcal I_{\rm lo}\cup\mathcal I_{\rm hi},\qquad
 \mathcal I_{\rm lo}\subset\{0<\xi_D<\xi_*\}\cup
 \{\xi_D\ge\xi_*/2\text{ in compact strict receivers}\}.
 \end{gathered}
\end{equation}
The low bound \(Y(0)\le2K_H(1+P_D^2)\) gives the finite
strict cover above; shrink \(\xi_*\) so its stopped part lies in
\eqref{sharp:upper-det-composition}. On strict cores the finite-cut
variations are units, and the common raw power
\eqref{sharp:upper-det-scales} cancels before the actual bridges and
divided closing bound \(\Pi'\). The signed high-height corollary
bounds the same derivative above.

The physical overlaps are \(\xi_*/2<\xi_D<\xi_*\) and
\(K_HS_D<H_\theta<2K_HS_D\). Gap equality faces belong to the
adjoining receiver, and both height equality faces are low. Window,
speed and cap clips preserve the cover and all intervening inputs;
a removed overlap places the connected retained interval in one
adjoining receiver. All these choices are made before selecting a
displacement root.
\end{proof}

\subsubsection{The stopped trace passage}
\label{subsec:trace-passage}

The trace chart contains a central saddle--node and an additional
equatorial saddle in the same shrinking neighbourhood. The central
normal form alone does not describe the return: its outgoing height
must also pass through this saddle. We include both passages in the
order of Figure~\ref{fig:trace-extra-saddle}.

\begin{proposition}[Uniform stopped-trace estimate]
\label{deg:prop-stopped-trace}
Fix \(0<p_0<p_1<\infty\), the passive compact coefficient block,
and the original buffered raw windows. There are independent bounds
\(r_*,\gamma_*>0\), an opening bound \(a_*>0\), a fixed positive
unscaled raw-height cap \(h_*\), and a final multiplier bound \(m_*\)
such that, for
\[
 \alpha_0=r^2a_L,\qquad \alpha_1=rp_L,\qquad
 0<r<r_*,\quad |\gamma|<\gamma_*,\quad |a_L|\le a_*,\quad
 p_0\le|p_L|\le p_1,
\]
the single original receiver
\begin{equation}\label{deg:trace-actual-receiver}
 \mathcal I_\theta=\mathcal I_\theta(h_*)
 =\{x:a_\theta+x\in J_\theta^{\rm out},\quad
       0<H_\theta(x)<h_*\},\qquad H_\theta(x)=z_\theta(\gamma;a_\theta+x),
\end{equation}
has \(F'>1/2\) everywhere for \(p_L<0\), and
\(F'=0\Rightarrow F''>0\) for \(p_L>0\), where
\(F(x)=x-m\Pi(x)\), \(0<m\le m_*\).
It consequently has at most one or two isolated zeros, respectively,
counted with total multiplicity. These conclusions apply to every
retained interval with the original window, speed and height clips,
including all intervening nonroot inputs and artificial equality
dividers. Both signs of \(a_L\), \(a_L=0\), the closed trace endpoints
and empty receivers are included; natural stopped endpoints are excluded.
The positive trace has a concave terminal receiver overlapping its
compact strict complement. No relation among \(m,r,\gamma\) is imposed.
\end{proposition}

% Figure contract: the stopped trace and its actual return composition.
% Native LaTeX picture; schematic, not a numerical phase portrait.
\begin{figure}[htb]
\centering
\setlength{\unitlength}{1pt}
\begin{picture}(430,238)
  % Panel (a): the outgoing extra saddle for positive central trace.
  \put(108,228){\makebox(0,0){\small (a) Positive-time geometry in $(q,Y)$}}
  \put(18,108){\line(1,0){198}}
  \put(216,108){\makebox(0,0)[l]{\small $q=-X$}}
  \put(38,108){\line(0,1){91}}
  \put(38,202){\makebox(0,0)[r]{\small $\Sigma_Q$}}
  \put(198,108){\line(0,1){91}}
  \put(198,202){\makebox(0,0)[l]{\small $\Sigma_-$}}
  \put(38,96){\makebox(0,0){\small $Q$}}
  \put(198,96){\makebox(0,0){\small $q_H$}}
  % Equator and its broken limiting connection.
  \linethickness{.35pt}
  \put(55,108){\line(1,0){95}}
  \put(150,108){\line(1,0){48}}
  \put(101,115){\makebox(0,0){\scriptsize equator $Y=0$}}
  \put(150,108){\circle*{5}}
  \put(150,95){\makebox(0,0){\scriptsize $e_R$}}
  \put(150,83){\makebox(0,0){\scriptsize extra saddle}}
  % Boundary connection to the saddle and its unstable branch.
  \linethickness{.8pt}
  \put(55,108){\vector(1,0){78}}
  \qbezier(150,108)(162,119)(174,132)
  \qbezier(174,132)(186,146)(198,163)
  \put(168,125){\vector(1,1){10}}
  \put(163,143){\makebox(0,0)[r]{\scriptsize unstable branch}}
  % Both curves increase with q. The strict orbit stays above the
  % unstable separatrix throughout their shared q interval.
  \linethickness{.55pt}
  \qbezier(38,154)(73,155)(108,158)
  \qbezier(108,158)(138,161)(162,171)
  \qbezier(162,171)(182,180)(198,193)
  \put(86,156){\vector(1,0){12}}
  \put(165,172){\vector(2,1){12}}
  \put(104,172){\makebox(0,0){\scriptsize typical strict orbit}}
  \put(38,183){\makebox(0,0)[l]{\scriptsize trace on $\Sigma_Q$}}
  \put(198,214){\makebox(0,0)[r]{\scriptsize fixed raw cut}}
  \put(108,66){\makebox(0,0){\scriptsize the axis segment ends at $e_R$;}}
  \put(108,54){\makebox(0,0){\scriptsize the positive-height orbit passes above it}}
  % Panel (b): actual composition used for the displacement.
  \put(325,228){\makebox(0,0){\small (b) The actual return composition}}
  \put(246,195){\makebox(0,0){\small weak input $t$}}
  \put(246,129){\makebox(0,0){\small weak input $t$}}
  \put(406,195){\makebox(0,0){\small affine input $x$}}
  \put(406,129){\makebox(0,0){\small raw output}}
  \put(278,195){\vector(1,0){30}}
  \put(340,195){\vector(1,0){30}}
  \put(278,129){\vector(1,0){30}}
  \put(340,129){\vector(1,0){30}}
  \put(307,207){\makebox(0,0){\scriptsize $M_d$ (DIR)}}
  \put(371,207){\makebox(0,0){\scriptsize $A_\theta$}}
  \put(307,141){\makebox(0,0){\scriptsize $\mathfrak H_\theta(t)=b_\theta+B_\theta(t)$}}
  \put(371,141){\makebox(0,0){\scriptsize $O_\theta$}}
  \put(311,180){\makebox(0,0){\scriptsize central passage}}
  \put(311,113){\makebox(0,0){\scriptsize actual outer Dulac}}
  \put(325,94){\makebox(0,0){\scriptsize includes $e_R$ and the fixed-section maps}}
  \put(325,78){\makebox(0,0){\scriptsize $\Pi\bigl(A_\theta(M_d(t))\bigr)=O_\theta(\mathfrak H_\theta(t))$}}
\end{picture}
\caption{The stopped trace and the map used in its zero estimate.  In
panel~(a), $q=-X$ increases in positive time.  The broken limiting
connection consists of the equatorial segment from the trace section
$\Sigma_Q$ to the additional equatorial saddle $e_R$, followed by its
unstable upper-half-plane branch toward the fixed raw section $\Sigma_-$.
The thinner curve is a representative strict positive-height orbit;
it remains above the unstable branch on their common $q$-interval.
The label \emph{trace on \(\Sigma_Q\)} marks the incoming height supplied by
the central saddle--node passage.  All distances and slopes are schematic.
Panel~(b) records the two sides of the physical composition: the central
saddle--node passage is the inverse Dumortier--Ilyashenko--Rousseau (DIR)
normal-form map $M_d$, while the outgoing height has the retained offset
$b_\theta$ and is then sent through the actual outer Dulac map $O_\theta$,
which includes the additional saddle and the fixed-section connectors.
Arrows in panel~(b) denote maps, not flow trajectories.}
\label{fig:trace-extra-saddle}
\end{figure}


\begin{proof}
For positive trace, we first estimate the outgoing equatorial passage
in its actual height coordinate. We then join it to the inverse central
passage and the incoming section map. Their logarithmic derivatives,
together with the regular connectors, give concavity near the stopped
orbit. A compact strict comparison and the high-height estimate extend
the count over the whole receiver. For negative trace, a direct
first-variation identity suffices.

Put \(v=\gamma+rX\), \(z=r^2Y\), and
\(\ell(v)=1-Av+Cv^2\), \(q_\gamma=1+2\gamma\). A positive time change gives
\begin{equation}\label{deg:trace-field}
\begin{split}
 X'&=a_L+p_LX-3\gamma X^2-rX^3-\ell(\gamma+rX)Y,\\
 Y'&=-Y\{(1+2\gamma)X+rX^2+Cr(\gamma+rX)Y\}.
\end{split}
\end{equation}
Keep the original affine \(x=h-a_\theta\) and
\(\Pi=\mathcal C_{\theta,m}\circ T_\theta(a_\theta+\cdot)\).
Lemma~\ref{phys:lem-c3-adapters} bounds divided closing through
ordinary order three, including \(m=0\). Phase derivatives fix all
coefficients and the actual multiplier. The two window admissions,
height and speed clips in \eqref{loc:nested-buffers} make
\eqref{deg:trace-actual-receiver} an interval before capture.

Choose the finite normal-form order and an initial \(|a_L|\le a_0\).
Take the signed high-height constant, increase \(K_H\) so that
\(|\bar a|\le a_0/K_H\), \(|\bar p|\le p_1/\sqrt{K_H}\),
\(\bar Y(0)=1\) lies inside its compact source block, and put
\(K_L=2K_H\). Fix \(X=L_c\) beyond the central equilibria.
The barrier estimate \eqref{sharp:upper-trace-barrier} below fixes a
bounded actual gap range before the terminal chords are selected.
Retain each auxiliary cut's exact continuation, centre and gain to
the original port. First take \(p_L\in[p_0,p_1]\).

\paragraph{The outgoing equatorial saddle.}
Set \(q=-X\).  The exact equations are
\begin{equation}\label{deg:trace-negative-field}
\begin{aligned}
 q'&=M(q)+L_-(q)Y,&M(q)&=-a_L+p_Lq+3\gamma q^2-rq^3,\\
 Y'&=Y\{qB(q)+\zeta(q)Y\},&B(q)&=1+2\gamma-rq,\\
 L_-(q)&=\ell(\gamma-rq),&\zeta(q)&=Cr(rq-\gamma).
\end{aligned}
\end{equation}
Fix \(w_0=1/8\), \(w_1=3/4\), and an auxiliary raw width
\(\delta\) inside the original ports, sufficiently small for this
fixed band. Choose a small \(a_0\) and a large \(Q\); then shrink
\(r,|\gamma|\), initially with \(|a_L|\le a_0\).
There is exactly one root \(R>Q\) of \(M\): \(M(Q)>0\),
\(M'(Q)>0\), and \(M'\) has exactly one positive root, beyond
which \(M\) decreases to \(-\infty\).  At \(R\),
\[
 rR^2=p_L+3\gamma R-a_L/R,\qquad
 \kappa=-M'(R)=p_L+rR^2-2a_L/R\ge p_0/2.
\]
Evaluating \(M\) on fixed compact intervals shows that \(R\to\infty\).
The identity \(rR=3\gamma+p_L/R-a_L/R^2\) then gives \(rR\to0\),
uniformly in the retained parameters.  The tangent and normal
eigenvalues are \(-\kappa\) and \(RB(R)>0\).  Thus the root is a
saddle, its incoming equatorial branch points from \(Q\) towards \(R\),
and its characteristic number satisfies
\[
 \beta_R=\frac{\kappa}{RB(R)}\longrightarrow0.
\]

Put
\[
 \begin{gathered}
 h(q)=\frac{M(q)}{q^2B(q)},\quad b(q)=qh(q),\quad
 \mathcal A(s)=\int_Q^s\frac{dq}{b(q)},\\
 \mathcal D=b(s)\partial_s,\qquad
 \epsilon=|a_L|/Q^2+p_1/Q+|\gamma|+rR.
 \end{gathered}
\]
This is the outgoing construction of
Lemmas~\ref{ann:lem-global-action-preparation} and
\ref{ann:lem-negative-gate}, with
\[
 (s,p,1)\quad\hbox{in their axis polynomial replaced by}\quad
 (r,p_L,-a_L).
\]
We verify the changed inputs. The first five Euler derivatives of
\(h=(-a_L/q^2+p_L/q+3\gamma-rq)/B(q)\) are \(O(\epsilon)\)
on \(Q/2\le q\le2R\) and its fixed logarithmic collars; the same
holds for \(L_-/B-1\) and \(\zeta q/B\). Moreover
\(b>0\) up to \(R\), \(|b'|\le C\epsilon\), and
\(\mathcal A\to\infty\) at \(R\). On the fixed entrance interval,
\[
 M(q)/q\ge p_0-2a_0/Q-3|\gamma|Q-rQ^2\ge p_0/2
 \qquad(Q/2\le q\le Q).
\]
After choosing \(R_0>2Q\), the bound \(b\le Cp_1\) on
\([Q,R_0]\) follows by an independent shrink of \(r,|\gamma|\).

For clarity, the two equations used in that construction are
\begin{equation}\label{deg:trace-normal-field}
 \begin{gathered}
 u'=h+cw,\qquad w'=w\{1-2h+(k-2c)w\},\\
 u=\log q,\quad w=Y/q^2,\quad c=L_-/B,\quad k=\zeta q/B,
 \end{gathered}
\end{equation}
in positive normal time \(d\tau_n=qB(q)\,dT\), and
\begin{equation}\label{deg:trace-graph}
 \frac{dw}{d\log q}=F(\log q,w)
 =w\left\{\frac{1+2\gamma-rq+Crq(rq-\gamma)w}
 {L_-(q)w+3\gamma-rq+p_L/q-a_L/q^2}-2\right\}.
\end{equation}
The first is unchanged in form, so
Lemmas~\ref{lem:normal-time-projection} and
\ref{lem:normal-time-action-height} give the strong footpoint and
its actual fixed-cut height on \(R_0\le s<R\):
\begin{equation}\label{deg:trace-height}
 \log y=-\mathcal A(s)+2\log s+\log U(s)+E_Q(y),\qquad
 0<y\le e_0=C_QR_0^2e^{-c(R_0-Q)/p_1}.
\end{equation}
Here \(U\) and its reciprocal and logarithmic derivatives are bounded,
\(E_Q(0)=0\), \(\|E_Q\|_{C^3}\le C_Q\), and
\begin{equation}\label{deg:trace-height-derivative}
 \mathcal D\log y=-K,\qquad
 |K-1|+|\mathcal DK|\le C\epsilon+C_Qe_0,\qquad K\ge\tfrac12.
\end{equation}
The passage stays in \(Q/2<q<2R\), and \(y\downarrow0\) at \(R\).

For the continuation, fix \(|\gamma|<\delta/2\). The strong-section
abscissa is \(q_s(s)=sU_s(s)\), with \(q_s'\asymp1\) and
three bounded logarithmic derivatives of \(U_s\). On
\(w_0\le w\le w_1\), the denominator in
\eqref{deg:trace-graph} is bounded below,
\(F(w_0)>c\), \(F(w_1)<-c\), \(F_w\le-1\), and its required
state and logarithmic-radius derivatives are bounded. Since
\(q_s(R)<q_H=(\delta+\gamma)/r\), the same invariant-strip and
moving-entrance argument as in
\eqref{ann:repair-moving-variation}--\eqref{ann:repair-completed-log-jets}
gives a unique raw exit and bounded logarithmic phase derivatives.
Only their first row is needed here. The new term
\(-a_L/q^2\) has the same integrable weight as the former \(q^{-2}\),
with bounded coefficient; no reciprocal \(\beta_R\) is introduced.

The full raw image is admitted before composing the closing. Indeed,
choose \(0<d_{\rm r}\ll\delta\), then require \(2rR<d_{\rm r}\).
Comparison from \(rq=d_{\rm r}\) to the raw exit with
\(w_\Gamma(\xi)=(\sqrt{1+\xi^2}+1)^{-1}\) gives, uniformly in \(s\),
\[
 |W_H(s)-w_\Gamma(\delta+\gamma)|
 \le C d_{\rm r}/\delta+
 C_{d_{\rm r},\delta}(|A-1|+|C-1|+|\gamma|+rp_L+|a_L|r^2).
\]
This is the raw comparison in the negative-gate proof with the same
changed polynomial. Choose \(d_{\rm r}/\delta\) from the original
window reserve, then shrink the coefficients; the raw height is
\((\delta+\gamma)^2W_H\), with organizing-height translation
\(O(\delta|\gamma|)\). The compact bridge and divided closing are
therefore defined on the whole image. Their ordinary \(C^3\) bounds
give the completed output \(\Gamma_-(s)\) with
\[
 \Gamma_-'<0,\qquad
 |s\partial_s\log(-\Gamma_-')|\le C,\qquad
 |\mathcal D\log(-\Gamma_-')|\le C\epsilon.
\]
The two physical radii \((\delta\pm\gamma)/r\) remain distinct in
their absolute gains.

Define \(O(y(s))=\Gamma_-(s)\). The same logarithmic chain rule as
in Proposition~\ref{ann:prop-endpoint-curvature} now gives
\begin{equation}\label{deg:trace-extra-saddle}
 \begin{split}
 O'(y)&=\frac{(-\Gamma_-')b}{yK}>0,\\
 y\frac{O''}{O'}
 &=-1-\frac{\mathcal D\log(-\Gamma_-')+b'-\mathcal D\log K}{K}
 =-1+O(\epsilon+C_Qe_0)<-\tfrac12.
 \end{split}
\end{equation}
It holds on the actual height interval \(0<y<y(R_0)\), uniformly
as \(\beta_R\to0\); no derivative at the excluded endpoint is used.

\paragraph{The central passage and its connectors.}
At \(a_L=r=\gamma=0\) the central eigenvalues are \(p_L,0\),
with the nonzero weak quadratic coefficient of
Lemma~\ref{deg:lem-positive-atlas}. Apply the finite-smooth
multiparameter orbital normal form of
\cite[Theorem~3 and formula~(2.2), October 2000 author copy p.~6]{DIR02}
to the reversed field, at the preselected joint order \(N\ge12\).
Its positive orbital factor gives the inverse physical passage, with
no scalar-opening rank condition.

The axis branch has \(X_{\rm ax}=-a_L/p_L+O(a_L^2)\), \(Y=0\),
and the interior branch satisfies
\[
 \left.\partial_{a_L}Y_{\rm int}\right|_{a_L=0}
 =\frac{1+2\gamma}
 {\ell(\gamma)(1+2\gamma)+p_LCr\gamma}>c>0.
\]
Their weak separation is \(a_L\) times a positive unit; the strong
rate stays nonzero. A finite compact cover gives common boxes and phase
bounds across \(a_L=0\). The physical boundary is the interior saddle
for \(a_L>0\), the axis saddle for \(a_L<0\), and their common root
at zero. Center this weak root at \(x_+(d)\), with the other at
\(x_+(d)-d\), \(d\asymp|a_L|\). Only signed-side phase derivatives
are needed.

In reversed normal-form time take \(x'=g_d(x)\), \(\eta'=-\eta\),
a fixed strong entrance \(\eta=\eta_0>0\), and a fixed weak exit
\(x=\rho\). The orbit from \((x_+(d)+t,\eta_0)\) reaches it at
height \(\eta_0M_d(t)\). With bounded normal-form coefficient
\(\mathfrak a\), the two weak factors give
\begin{equation}\label{deg:trace-stopped-map}
 \begin{gathered}
 G_d(t)=\frac{t(t+d)}{1+\mathfrak a(x_+(d)+t)},\qquad
 g_d(x_+(d)+t)=G_d(t),\\
 M_d(t)=\exp\!\left[-\int_{x_+(d)+t}^{\rho}\frac{dx}{g_d(x)}\right].
 \end{gathered}
\end{equation}
For small \(d\ge0\), \(t>0\),
\begin{equation}\label{deg:trace-stopped}
 M_d'=\frac{M_d}{G_d},\qquad
 \frac{M_d''}{M_d'}=\frac{1-G_d'}{G_d},\qquad
 1-G_d'\ge\delta_G>0,\quad G_d\asymp t(t+d).
\end{equation}
The quantities \(G_d\) and \(M_d\) are uniformly small on a short
terminal interval: \(G_d\le Ct(t+d_*)\) for \(0\le d\le d_*\),
while the lower action bound \(c\int_t^{\rho/2}du/[u(u+d_*)]\),
with fixed small \(d_*>0\), diverges at zero.
Keeping the two coalescing factors together makes this estimate
uniform as their separation vanishes.

Absorb the fixed factor \(\eta_0\) into the incoming section map.
From the central strong exit to \(q=Q\), the flight has fixed
finite length and \(q'\ge cq>0\).  Ordinary graph variations give
\begin{equation}\label{deg:trace-height-connector}
 \begin{gathered}
 \mathfrak H_\theta(t)=b_\theta+B_\theta(t),\qquad B_\theta(0)=0,\\
 0<c_Q\le B_\theta'\le C_Q,\qquad
 |\mathfrak H_\theta''/\mathfrak H_\theta'|\le C_Q,\qquad
 0\le b_\theta\le C_Q|a_L|.
 \end{gathered}
\end{equation}
Here \(b_\theta=0\) for \(a_L\le0\). For \(a_L>0\), the
interior saddle's strong branch has a uniform spectral gap; its
smooth finite flight to \(Q\) gives the positive offset. Retain this
offset in \(O\), using phase regularity on each signed side.

After \(Q,R_0\) are fixed, \(y_\theta(R_0)\) has a uniform
positive lower bound.  At \(r=\gamma=0\) it is the height of a
finite regular passage on a compact parameter block, and continuity
preserves positivity.  Choose \(y_*>0\) below this bound.  Only
now shrink \(a_*,t_*\) so that
\begin{equation}\label{deg:trace-domain-inclusion}
 \mathfrak H_\theta([0,2t_*])<y_*/2<y_\theta(R_0).
\end{equation}
Thus every strict terminal orbit, including those with positive
offset, lies in the domain of \eqref{deg:trace-extra-saddle}.

For the incoming connector, at \(r=\gamma=0\) the boundary branch
beyond the positive local cut satisfies \(X,Y>0\) and
\(D=Y-a_L-p_LX>0\).  At \(D=0\), reverse time points into
\(D>0\), since \(\dot D=p_LD-XY\).  Also \(Y_X\le2X\)
for large \(Y\) on bounded \(X\)-intervals.  These facts continue
the branch to every fixed positive cut without loss of its graph
or finite-\(X\) blow-up.  At large fixed \(X\),
\[
 \frac{dW}{d\log X}=\frac{W}{W-p_L/X-a_L/X^2}-2W,
 \qquad W=Y/X^2.
\]
Choose the cut so that \(0<p_L/X+a_L/X^2<1/4\).  The right
side is at least \(1/2\) for \(0<W\le1/4\) with \(D>0\),
and is negative for \(W\ge1\).  The compact family enters a
common quadratic corridor at a fixed cut \(L_t\), as does a small
chord of strict incoming orbits for nearby coefficients.
For nearby coefficients the rational tail equation has
\(F_W\le-1\) and bounded \(F_{WW}\). Its first variation decays
exponentially, so the relative-jet argument
\eqref{tail:relative-jet-equations} bounds the pre-Schwarzian uniformly.
The same fixed raw comparison as on the negative side admits the whole
chord into the positive window. After the finite connector and original
bridges, the actual input adapter therefore satisfies
\begin{equation}\label{deg:trace-input-adapter}
 x=A_\theta(M_d(t)),\qquad A_\theta'>0,\qquad
 |A_\theta''/A_\theta'|\le C_Q.
\end{equation}
Its absolute first derivative needs no lower bound.

\paragraph{Curvature and the overlap of actual phase intervals.}
The three constructed maps refer to the same orbit by flow uniqueness:
\(A_\theta\) continues to the actual input, \(\mathfrak H_\theta\)
to \(q=Q\), and \(O_\theta\) through the extra saddle and closing.
Thus
\[
 x=A_\theta(M_d(t)),\qquad \Pi(x)=O_\theta(\mathfrak H_\theta(t)).
\]
All factors are increasing. Put \(H=\mathfrak H_\theta(t)\).
Differentiation in \(t\), keeping the
affine input \(x\) in the displacement, gives
\begin{equation}\label{deg:trace-composition}
 \begin{split}
 G_d\frac{d}{dt}\log\Pi'(x(t))
 ={}&G_d\frac{O''(H)}{O'(H)}H'+G_d\frac{H''}{H'}\\
 &-M_d\frac{A''(M_d)}{A'(M_d)}-(1-G_d').
 \end{split}
\end{equation}
The first term is nonpositive by \eqref{deg:trace-extra-saddle}
and \eqref{deg:trace-domain-inclusion}.  The next two are bounded
by \(C_Q(G_d+M_d)\), whereas the last is at most \(-\delta_G\).
Shorten the terminal interval so that the whole expression is at
most \(-\delta_G/2\).  Since \(G_d>0\) and \(x'(t)>0\),
this proves \(\Pi''<0\) in the original affine input.

Fix a positive local cut \(X=L_c\) beyond the nearby central
equilibria, and let \(Y_b^+(L_c)\) be the incoming boundary
height.  The normalized actual gap is
\begin{equation}\label{deg:trace-actual-gap}
 \xi=Y(L_c)-Y_b^+(L_c)
 =\frac{z(\gamma+rL_c)-r^2Y_b^+(L_c)}{r^2}.
\end{equation}
Thus \(r^2\xi\) is the actual raw gap on this section, with the
same zero, ordering and orientation as \(\xi\).
Order and the compact local connector make \(\xi(t)\) strictly
increasing. More precisely, the connector on the fixed local cuts
gives \(\xi(t)=\psi_\theta(M_d(t))-\psi_\theta(0)\), with
\(0<c\le\psi_\theta'\le C\). Since \(G_d(u)\ge cu^2\),
the action in \eqref{deg:trace-stopped-map} at \(t=t_*/2\) is
at most \(C/t_*\). Consequently
\(\xi_\theta(t_*/2)\ge c\exp(-C/t_*)>0\), uniformly on the
completed coefficient block, including \(d=0\). Choose
\begin{equation}\label{deg:trace-overlap}
 0<2\xi_0<\inf_\theta\xi_\theta(t_*/2).
\end{equation}
Then \(0<\xi\le2\xi_0\) lies in the terminal interval.
The whole lower range \(Y_0=H_\theta(x)/r^2\le K_L\) reaches this
same cut. Indeed its exact scalar graph is
\begin{equation}\label{sharp:upper-trace-barrier}
 Y_X=\frac{Y\{(1+2\gamma)X+rX^2+Cr(\gamma+rX)Y\}}
 {\ell(\gamma+rX)Y-a_L-p_LX+3\gamma X^2+rX^3}.
\end{equation}
The through denominator is positive. Above a fixed \(M_0>K_L+1\)
it is at least \(\ell_0Y/2\), and on \(0\le X\le L_c\)
\(Y_X\le C_{L_c}+Y\). Comparison after the first crossing of \(M_0\)
and Gronwall give the previously fixed bound
\begin{equation}\label{sharp:upper-trace-bounded-gap}
 0<\xi=Y(L_c)-Y_b^+(L_c)\le B.
\end{equation}
Only original through admission is used here. The divider height is
nonnegative and is subtracted on the same section.

On \(\xi\ge\xi_0\) the limiting passage is uniformly separated from
its stopped graph. To obtain this separation before constructing a strict
cover, let \(t_{\rm ref}\) be the unique phase with \(\xi=\xi_0\).
The uniform vanishing of \(M_d(t)\) bounds \(t_{\rm ref}\) away from
zero. The standing normal form gives a nonbroken physical passage at
every such positive phase: for \(a_L>0\) it crosses above the interior
saddle with \(Y_0>a_L\), and for \(a_L\le0\) with \(Y_0>0\).
Take the positive minimum over the compact reference-phase and
coefficient block at \(r=\gamma=0\); scalar order then gives
\[
 Y_0-\max\{a_L,0\}\ge\eta_0>0\qquad(\xi\ge\xi_0).
\]
For \(D=Y-a_L-p_LX\), \(D(0)>0\) and
\(\dot D=p_LD-XY\) excludes a first zero in either relevant half-axis
orientation. The bound \(Y_X\le C+Y\) excludes finite-\(X\) blow-up.
Thus the bounded complement has a compact strict-family cover with
positive core height and speed margins. At large \(|X|\), its
\(W=Y/X^2\) equation points upward below \(1/4\) and downward above
\(1\), so both half-tails enter quadratic corridors at a common fixed
cut. Their actual outward continuations and fixed raw comparison give
both original-window margins. In particular
\(\xi_0\le\xi\le2\xi_0\) is a physical terminal/strict overlap.

For a strict limiting graph, Proposition~\ref{ann:prop-source-curvature}
with \(a=a_L\), \(p=p_L\), \(b=B_+\), and
\(\sigma=\partial_bY>0\) gives
\begin{equation}\label{deg:trace-strict-curvature}
 \frac{S''}{S'}=-2\int_{\mathbb R}\frac{X(a_L+p_LX)\sigma}{D^3}\,dX
              =-p_L\int_{\mathbb R}\frac{\sigma}{D^2}\,dX<0.
\end{equation}
The compact strict cover bounds \(\sigma\) below and \(D\) above
on a common core; \(p_L\ge p_0\) gives a uniform negative margin.

To transfer this margin, use the exact normalization
\eqref{deg:det-normalized-tail}. The changed coefficient \(a_L\)
gives the weight
\[
 \|N\|_{C^3}\le C\{rR^{-2+\nu_0}+|\gamma|R^{-3+\nu_0}
                  +|p_L|R^{-4+\nu_0}+|a_L|R^{-5+\nu_0}\},
 \qquad 0<\nu_0<1.
\]
The argument of Proposition~\ref{prop:source-normalization}, with this
integrable weight and the admitted enlarged chords, gives phase \(C^3\)
convergence independently in \(r,\gamma\). Exterior factors normalized
by their exact value and slope converge to the identity by
\eqref{deg:det-adapter-normalization}--\eqref{deg:det-divided-closing}.
Use the actual maps \(h(b),k(b)\) and their normalization
\eqref{eq:reference-normalization}, with
\(c_h=h'(b_*)>0\), \(c_k=k'(b_*)>0\). Let
\(\beta=b_*+(x+a_\theta-h(b_*))/c_h\). As in
\eqref{ann:p0-two-actual-adapters},
\begin{equation}\label{deg:trace-physical-curvature}
 \Pi''(x)=\frac{c_k}{c_h^2}\Phi''(\beta)<0.
\end{equation}
Both raw first variations are \(r^{2-\lambda_\gamma}\) times
positive units. Their quotient, the actual bridges and divided closing
bound \(\Pi'\) on the strict pieces and their regular finite-flight
ends; the individual gains may vanish.

At every actual point with \(Y_0>K_H\), freeze the dilation at that
point before a phase derivative. The exact normalized parameters are
\begin{equation}\label{sharp:upper-trace-high-dilation}
 (\bar a,\bar p,\bar\gamma,\bar s)=
 (a_L/Y_0,p_L/\sqrt{Y_0},\gamma,r\sqrt{Y_0}),\qquad \bar Y(0)=1.
\end{equation}
The fixed unscaled cap \(h_*\) gives \(\bar s<\sqrt{h_*}\).
The signed compact source block of
Corollary~\ref{rew:cor-high-derivative-signed} receives these parameters.
Its two absolute original-port variations have the same
\(\bar s^{2-\lambda_\gamma}\) power, including every bridge and
\eqref{deg:det-divided-closing}. Their quotient bounds the derivative
of this same physical \(\Pi\) at every high input, including nonroots.

Thus the original set cover is
\[
 \mathcal I_\theta=\{Y_0\le K_L\}\cup\{Y_0>K_H\},\qquad
 \{Y_0\le K_L\}\subset\{0<\xi\le2\xi_0\}\cup\{\xi\ge\xi_0\}.
\]
Each set is restricted to \(\mathcal I_\theta\). The physical overlaps
are \(K_H<Y_0<K_L\) and \(\xi_0<\xi<2\xi_0\). Enlarged chords
receive both positive gap equality faces; the height equality faces are
low. Original window, speed or height clipping preserves this cover,
including all intervening inputs, and excludes the broken \(t=0\).
On the terminal and strict pieces \(\Pi''<0\); on the high and
regular pieces \(\Pi'\le C\). Choose \(m_*\) last so that the latter
have \(F'>1/2\). Every zero of \(F'\) has \(F''=-m\Pi''>0\).
Lemma~\ref{sharp:upper-critical-count} gives the whole positive-trace
bound.

\paragraph{The negative trace.}
Now take \(\alpha_1=-rp\), \(p\in[p_0,p_1]\), and
\(\alpha_0=r^2a_L\), \(|a_L|\le a_*\). Let \(\delta_{\rm raw}\)
bound \(|v-\gamma|\) on the original small raw chart. In the raw coordinate
\(u=v-\gamma\), put
\[
 \begin{gathered}
 \mathcal M=-r^2a_L+rpu+3\gamma u^2+u^3,\qquad
 \mathscr D=\mathcal M+\ell(v)z>0,\\
 \nu=\frac{6\gamma}{1+2\gamma},\qquad N=rp+(3-\nu)u^2\ge rp+2u^2.
 \end{gathered}
\]
The graph equation
\(z_v=f=z\{(1+2\gamma)u+u^2+Cvz\}/\mathscr D\)
satisfies \eqref{deg:raw-log-jacobian} with these \(\mathcal M,N\).
\begin{equation}\label{deg:trace-jacobian}
 f_z=\frac d{dv}\log\frac{z^{1+\nu}}{\mathscr D}
   +\frac{N+\{\ell'(v)+(1-\nu)Cv\}z}{\mathscr D}.
\end{equation}
On \(\mathcal M<0\),
\[
 |\mathcal M|\le a_*r^2+rp|u|+3|\gamma|u^2+|u|^3
 \le C(a_*r+\delta_{\rm raw}+|\gamma|)N.
\]
Let \(\kappa_0=\tfrac12\inf\ell>0\).  Where
\(\mathscr D\ge\kappa_0z\), the last quotient in
\eqref{deg:trace-jacobian} is uniformly bounded below.
On the complement, \(\mathcal M<0\) and
\(z\le|\mathcal M|/(\ell-\kappa_0)\); the displayed estimate
makes its numerator at least \(N/2\) on the standing small raw chart
after the independent coefficient shrink.  It is therefore nonnegative there.
Integrating from the positive raw section to the negative one
reverses the orientation of the interval and gives
\begin{equation}\label{deg:trace-negative-derivative}
 0<T'\le C\frac{(z^{1+\nu}/\mathscr D)_-}
 {(z^{1+\nu}/\mathscr D)_+}\le C_0.
\end{equation}
The endpoint ratio is bounded on the fixed actually admitted raw
windows: endpoint heights are positive units and graph speeds have
fixed positive lower bounds. No lower central speed is used. The
estimate holds on the full receiver \eqref{deg:trace-actual-receiver},
including its speed clip and every intervening nonroot input.
The ordinary divided closing has bounded positive derivative, so
\(0<\Pi'\le C_{\rm neg}\). Add \(m_*C_{\rm neg}<1/2\) to the final
multiplier choice; then \(F'>1/2\) throughout, giving at most one
simple zero. The positive-phase estimates remain uniform at \(a_L=d=0\).
\end{proof}


\begin{lemma}[Nonnegative-opening domains]
\label{deg:lem-nonnegative-domain}
On every coefficient fibre with \(\alpha_0\ge0\), the coefficient
alternatives and phase intervals of
Lemma~\ref{deg:lem-positive-atlas} cover the connected captured
physical interval, with their regular continuations and all points
between captured roots. The compact strict and regular portions have
bounded complete raw derivative. The common high-height interval is
contractive after the multiplier shrink of
Corollary~\ref{rew:cor-high-derivative-signed}, including the
positive-height equality face \(\alpha_0=\alpha_1=0\).
The only remaining stopped portions are the determinant passage of
Lemma~\ref{deg:lem-determinant-preparation} and the signed trace
passages of Proposition~\ref{deg:prop-stopped-trace}.
\end{lemma}

\begin{proof}
For each whole coefficient fibre use the alternative of
Lemma~\ref{deg:lem-positive-atlas}. The actual determinant cover
\eqref{deg:det-height-cover} and signed trace receiver
\eqref{deg:trace-actual-receiver} retain its stopped, strict, regular
and high pieces, including the original clips and every intervening
input. Their two absolute tail gains have a common power, which cancels
before the complete raw derivative is bounded on the strict and regular
pieces. The signed high-height estimate bounds that same derivative
above, including when both opening coefficients vanish. Choosing the
finite multiplier last makes these bounded pieces contractive.
The gap-zero completion is excluded. The whole-fibre sharp counts are
Propositions~\ref{sharp:upper-determinant} and \ref{sharp:upper-trace};
Corollary~\ref{sharp:upper-local-coefficient-ranges} takes the maximum
across coefficient alternatives.
\end{proof}

\section{Quadratic inward and outward section estimates}
\label{app:moving-section-bookkeeping}

We estimate the actual quadratic passages on their original sections,
keeping the coefficient-dependent first multipliers. The inward maps
provide heights on one common centre section; their equality defines the
match. The outward maps carry the two lips to the original raw ports.
Together these estimates control the section terms in physical curvature.
The final paired estimate bounds the complete derivative above a
comparison orbit, where the separate first variations may be large.
We use the scalar variation and continuation formulas of the physical
section toolkit throughout.

\subsection{Quadratic radial embeddings on common raw sections}

For the quadratic passage estimates below, fix the original raw-section and
ordinary \(C^3\) observer reserves and the passive compact set.
Phase derivatives freeze \(g,k,\pi,A,C,m\), with \(g,k>0\) independently
small. In the positive or reflected negative radial coordinate, use
\begin{equation}\label{deg:quadratic-inward-family}
 \begin{aligned}
 P_s(v)&=1-s\pi v+3v^2+skv^3,& B_s(v)&=1+g(2+skv),\\
 Q_s(v)&=1-Ag(1+skv)+Cg^2(1+skv)^2,&
 \zeta_s(v)&=sCkg^3(1+skv),\qquad s=\pm1.
 \end{aligned}
\end{equation}
The exact scalar graph is
\begin{equation}\label{deg:quadratic-inward-scalar-common}
 Y_v=\frac{Y(vB_s+\zeta_sY)}{g(P_s+Q_sY)},
 \qquad x=sgkv,\quad z=g^3k^2Y.
\end{equation}
Decreasing \(v\) follows positive time for \(s=+1\) and its inverse
for \(s=-1\). On each domain where the actual lip continuation
exists, denote its original-centre height by \(Z_s(x_s)\) and define
\begin{equation}\label{deg:quadratic-inward-memory-common}
 \begin{gathered}
 A_s(x_s)=\int_0^{x_s}\frac{vB_s}{P_s}\,dv,\quad
 \mathcal D_s=\frac{P_s(x_s)}{x_sB_s(x_s)}\partial_{x_s},\\
 \log Z_s=\log\{\ell_0P_s(x_s)\}-A_s(x_s)/g+\mathcal R_s,
 \quad H_s=P_s'/(vB_s),\quad N_s=-g\mathcal D_s\log Z_s .
 \end{gathered}
\end{equation}
The polynomial factor and action in this formula retain the entrance
height and the full linear contraction to the centre.
Coefficient-only memory remains in \(\mathcal R_s\); the estimates
below concern its positive-order rows. Write \(J=A_+\) and omit
\(s=+1\) subscripts when only the positive graph is used.

On a fixed compact segment \([0,b]\) with \(P_s\ge p_b>0\),
the scalar graph has the common preparation
\begin{equation}\label{deg:positive-large-S-inward-compact-graph}
 \begin{gathered}
 Y_v=\frac{vB_s}{gP_s}Y+K_s(v,Y)Y^2,\qquad
 K_s=\frac{\zeta_sP_s-vB_sQ_s}{gP_s(P_s+Q_sY)},\\
 m_s(v)=\exp\left[-\frac1g\int_v^b\frac{wB_s(w)}{P_s(w)}\,dw\right]
 \le1,\qquad m_s(b)=1 .
 \end{gathered}
\end{equation}
Fix a prescribed height buffer \(0\le Y\le h_b\) before shrinking
\(g,k\). The positive units \(B_s,Q_s\), the margin \(p_b\), and
rational differentiation give
\( |\partial_Y^jK_s|\le C_b/g\), \(0\le j\le3\), on that buffer.
Lemma~\ref{deg:lem-compact-height-continuation} therefore applies with
\(\kappa_c=g\) whenever the incoming height satisfies
\(0<y\le h_b/2\) and \(y/g\) is small. Each regional proof supplies
these input bounds and the logarithmic phase rows of \(y\); the
continuation keeps the same original centre and uses no reciprocal
bound on \(m_s(0)\).

\begin{lemma}\label{deg:lem-quadratic-outer-jets}
Fix the raw section buffers, then a compact neighborhood of the passive
coefficients \(A,C\), and \(\pi\in[\pi_\flat,\Pi]\).
There are fixed \(M,L,d_q,\delta_q,g_0,k_0>0\) such that the following
holds for independently small \(0<g<g_0\), \(0<k<k_0\).
For either \(s=\pm1\), let
\[
 x_0\ge2M,\quad q=x_0^{-1},\quad S=kx_0\le\delta_q,
 \quad a=gkx_0,\quad b_0=-x_0^2/2 .
\]
The orbit from the lip at \(v=g+sa\) reaches \(v=g+sd_q\)
with positive height. Its raw embedding \(H_s(b_0)\), and any increasing
observer \(\Phi(H_s(b_0))\) on the prescribed buffered raw image for which
\(\|\partial_z\log\Phi'\|_{C^1}\le C\), have positive first derivative. Put
\[
 E=x_0\partial_{x_0}=a\partial_a,\qquad
 h_s=\frac{3-s\pi q+q^2+sS}{1+2g+sgS},\quad D_s=h_sE,
 \quad\lambda=\frac{12g}{1+2g},\quad L_g=1+|\log g|.
\]
After separate affine normalizations of the embeddings, their slopes obey
\begin{equation}\label{deg:quadratic-outer-weighted-jets}
 D_s\log\phi_s'=-\lambda h_s+G_s,
 \qquad |G_s|+|D_sG_s|\le CgL_g(q+S).
\end{equation}
Primes use the lip label \(b_0\); \(g,k,\pi,A,C\) are held fixed.
The two raw passages exist separately. The conclusion does not require
their inward central images to intersect.
\end{lemma}

\begin{proof}
The common spectral estimate separates the exact homogeneous power
from its logarithmic residual. We first check outward existence and
admission to the raw buffers, then apply its derivative bounds with
the original moving initial data and terminal observer.

Write \(x=|v-g|\), \(W=z/x^2\), \(\tau=\log(x/a)\), and set
\begin{align*}
 m_s(x)&=3g+sx-s\pi kg^2/x+k^2g^3/x^2,\\
 Q_s(x)&=1-A(g+sx)+C(g+sx)^2,\\
 N_s(x,W)&=1+2g+sx+sCx(g+sx)W,
 \qquad D=m_s+Q_sW .
\end{align*}
Division of the exact raw field gives
\begin{equation}\label{deg:quadratic-outer-ratio}
 W_\tau=F_s(x,W):=W(N_s/D-2),\qquad
 W(0)=\ell_0gp,\quad p=3-s\pi q+q^2+sS.
\end{equation}
Use the exact signed dictionary for Lemma~\ref{deg:lem-spectral-exterior},
\[
 t=s,\quad b=2g,\quad v=g,\quad
 (\mu_2,\mu_1,\mu_0)=(3g,-s\pi kg^2,k^2g^3),\quad
 a=gS=gkx_0,\quad\sigma=gk.
\]
Choose \(M\) large and \(\delta_q\) small. For \(S\le\delta_q\),
\(m_s(x)\asymp g\) on \([a,2a]\), for both signs, and
\[
 |b|+|v|+|\mu_2|+a+|\mu_1|/a+|\mu_0|/a^2\le Cg,
 \qquad \Theta=g(S+\pi q+q^2)\le Cg(S+q).
\]
Thus the common lemma's intrinsic construction applies with
\(\rho\asymp g\). For the negative sign it gives the original
root-free outward passage and invariant strip even when \(m_-\)
changes sign beyond \(2a\).

\smallskip
\noindent\emph{Common positive admission.}
For \(s=+1\), the same construction applies on the larger domain
\begin{equation}\label{deg:quadratic-positive-admission}
 x_0\ge2M,\qquad a=gkx_0\le a_0,\qquad 0<g<g_0,\quad0<k<k_0,
\end{equation}
without a bound on \(S\). Indeed \(kg/x\le q\) for \(x\ge a\),
so the choice of \(M\) gives
\[
 x+2g\le m_+(x)\le x+4g,\qquad
 6g+a+\pi gq+gq^2\le C(g+a).
\]
These are the two exterior hypotheses with \(\rho=g+a\).
Choose \(d_q\) within that lemma's fixed cut allowance and
\(a_0<d_q/L\), \(L\ge2\), with \(g_0+a_0<\min(\rho_0,1)\).
Shrink \(g_0\) so that \(g_0\delta_q<a_0\).
Its construction gives, for both signs on their respective domains,
\begin{equation}\label{deg:quadratic-outer-domain}
 D\ge c(g+W),\quad 0<W\le w_1,\qquad
 W\ge w_0\quad(x\ge La).
\end{equation}
For the positive sign, in fact \(D\ge c(g+x+W)\).
The limiting graph in \eqref{deg:spectral-exterior-window} keeps
\(t=s\), hence \(m_s^*(x)=sx\). Its positive window is
\[
 |W(d_q)-W_+^*(d_q)|\le C(g+a)+C(a/d_q)^c.
\]
The original regular-section reserves give common neighbourhoods of the
two limiting terminal families whose continuations lie strictly inside
the prescribed raw buffers. Choose \(a_0\), then \(g_0\), inside the
positive reserve. For the negative sign the same window with \(\rho\asymp g\)
and \(a\le g\delta_q\) gives its original buffered image after shrinking
\(g_0\). Thus both regular continuations are evaluated on admitted
images. The positive graph is the inverse of the actual incoming
passage; the negative graph follows positive time, and \(z_x>0\).

For the logarithmic rows on \(S\le\delta_q\), use \(\rho\asymp g\)
for both signs. The coefficient bounds hold on a slightly enlarged open
phase interval, with all coefficients fixed under \(E\) and \(\beta=b_0\).
The admitted regular continuation followed by the observer is the
\(\chi\) of the common lemma; its prescribed ordinary \(C^3\)
logarithmic-slope reserve gives \(\mathcal T\le C\).

Write \(W_l=\ell_0gp\), \(\eta=\log z\), and
\(\nu=6g/(1+2g)=\lambda/2\).
The exact moving factor specializes to
\begin{equation}\label{deg:quadratic-outer-moving-factor}
 \mathscr B_s=EW_l-F_s(a,W_l)
 =\ell_0\left[g(6-s\pi q+3sS)
 -\frac{1+2g+sa+sCa(g+sa)W_l}{1+\ell_0Q_s(a)}\right].
\end{equation}
It is a negative unit. With
\(\mathscr L_s=[2m_s+x m_s'+(xQ_s'+sCx(g+sx))W]/D\),
the exact spectral density is
\begin{equation}\label{deg:quadratic-outer-density}
 \mathscr L_s\,d\tau=\nu\,d\eta+J_s\,d\eta,\qquad
 J_s=\frac{s\{(3-\nu)x-\pi kg^2/x
       +x[-A+(3-\nu)C(g+sx)]W\}}{N_s} .
\end{equation}
Put \(I(a)=\int_{\eta_l(a)}^{\eta_T(a)}J_s\,d\eta\).
For an observer choose a fixed reference \(z_*\) in its buffered image,
put \(c_\theta=\Phi'(z_*)>0\), and \(\widehat\Phi'=\Phi'/c_\theta\).
The logarithmic derivative assumption makes \(\widehat\Phi'\) a positive
unit there. For the raw embedding take \(c_\theta=1\).
The common exact slope, using \(b_0=-a^2/(2g^2k^2)\), is
\begin{equation}\label{deg:quadratic-outer-exact-slope}
 \begin{gathered}
 (\Phi\circ H_s)'(b_0)=d_\theta a^{-\lambda}U_s,
 \quad d_\theta=c_\theta g^2k^2\ell_0^{-1-\nu}g^{-\nu},\\
 U_s=(-\mathscr B_s)(1+\ell_0Q_s(a))p^{-\nu}
       \frac{W_T}{D_T}z_T^\nu\widehat\Phi'(z_T)e^I .
 \end{gathered}
\end{equation}
By \eqref{deg:spectral-exterior-rows}, its first two logarithmic residual
rows satisfy
\[
 |E\log U_s|+|E^2\log U_s|
 \le C\{gL_g(q+S)+a^{2-\lambda}\}\le CgL_g(q+S),
\]
since \(2-\lambda>1\) and \(a=gS\). The spectral factor retains
\(\lambda=12g/(1+2g)\). With \(R_s=E\log U_s\) and \(G_s=h_sR_s\),
\[
 D_s\log\phi_s'=-\lambda h_s+G_s,\qquad
 D_sG_s=h_s(Eh_s)R_s+h_s^2ER_s.
\]
The bounded \(h_s,Eh_s\) give
\eqref{deg:quadratic-outer-weighted-jets}. Separate affine value/slope
calibrations are constant under \(E\) and do not change these rows.
\end{proof}

\begin{corollary}\label{deg:cor-positive-large-S-outer}
Keep the original positive raw section and observer buffers of
Lemma~\ref{deg:lem-quadratic-outer-jets}, and fix \(\delta>0\).
There are \(a_0,g_0,k_0>0\) such that its positive outward passage and
observer are defined on the whole range
\[
 x_0\ge2M,\qquad S=kx_0\ge\delta,\qquad
 a=gS\le a_0,\qquad 0<g<g_0,\quad0<k<k_0.
\]
With the same fixed-coefficient \(E=x_0\partial_{x_0}\), quadratic
label and coefficient-only slope normalization, one has
\begin{equation}\label{deg:positive-large-S-direct-jets}
 |E\log\phi_+'|+|E^2\log\phi_+'|
 \le C_\delta a(1+|\log a|).
\end{equation}
There is no upper bound on \(S\) independent of \(g\), and no comparison
of \(g\) with \(k\). The observer has exactly the same stated
ordinary \(C^3\) logarithmic-slope input as in the lemma.
\end{corollary}

\begin{proof}
Use the common positive admission \eqref{deg:quadratic-positive-admission}
and its original buffered continuation. Apply
Lemma~\ref{deg:lem-spectral-exterior} with the same positive dictionary,
\(\rho=g+a\), \(\sigma=gk\) and terminal cut \(d_q\).
Its coefficient and mass hypotheses were checked before imposing
\(S\ge\delta\); the admitted ordinary \(C^3\) observer again gives
\(\mathcal T\le C\). Now
\[
 g\le a/\delta,\quad
 \Theta=a+\pi gq+gq^2\le C_\delta a,\quad
 L_\rho\le1+|\log a|,\quad
 \lambda=\frac{12g}{1+2g}\le12a/\delta.
\]
With \(e=2-\lambda>3/2\), the spectral rows give
\[
 |E\log\phi_+'+\lambda|+|E^2\log\phi_+'|
 \le C_\delta\{a(1+|\log a|)+(a/d_q)^e\}
 \le C_\delta a(1+|\log a|).
\]
Absorbing \(\lambda=O_\delta(a)\) proves
\eqref{deg:positive-large-S-direct-jets}. The exact quadratic label,
coefficient-only \((gk)^2\) gain and \(a^{-\lambda}\) are retained.
\end{proof}


\begin{corollary}\label{deg:cor-negative-absolute-outer}
Keep the original negative raw section and buffered observers, and fix
\(r_c>0\). For independently small \(g,k>0\), the actual positive-time
negative outward passage is defined for every
\[
 r\ge r_c,\qquad 0<R=kr<\mathfrak c,
 \qquad p=3+\pi/r+r^{-2}-R>0.
\]
Here \(\mathfrak c\) is the negative simple-root endpoint already defined
in the quadratic chart. Put
\[
 E=r\partial_r=R\partial_R,\qquad B=1+g(2-R),\qquad
 h=p/B,\qquad \mathcal A=hE=\mathcal D_-.
\]
With primes in either quadratic label \(-r^2/2\) or \(-R^2/2\),
\begin{equation}\label{deg:negative-absolute-direct-jets}
 |\mathcal A\log\phi_-'|+|\mathcal A^2\log\phi_-'|
 \le Cg(1+|\log g|).
\end{equation}
The ordinary \(C^3\) buffered observer input is the same as in the lemma;
it includes the actual divided finite closing on its original windows.
\end{corollary}

\begin{proof}
Use the exact negative equations of the lemma, with
\(a=gR\), \(q=1/r\), \(W_l=\ell_0gp\).
Here \(p\) may tend to zero at the simple root, so the weak height
interval can be arbitrarily long. The moving initial variation has
a factor \(p\), which will absorb this dwell in the differentiated
integrals. We first construct the graph in logarithmic height.
On the stated range \(p\) is bounded above, \(R\) is bounded,
\(a\le Cg\), and \(m_-(a)=gp\). Also \(m_-'(x)<0\).
To establish the actual graph without presupposing a positive raw
radial denominator, use \(\eta=\log z\), \(\chi=\log x\), and
\[
 \chi_\eta=V=(m_-+Q_-W)/N_-,\qquad W=e^{\eta-2\chi}.
\]
This equation is regular while \(N_-\) is a positive unit.
Initially \(D_l=gp(1+\ell_0Q_-(a))>0\).
At a first zero of \(D=m_-+Q_-W\), its derivative along this graph
would be \(D_\eta=Q_-W>0\). Hence \(D>0\) throughout the graph.
Below a smaller auxiliary positive strip threshold,
\(D\le gp+Q_-W\) and \(V<1/4\); thus \(W_\eta>0\) and
\(W\ge\ell_0gp\). Consequently
\[
 W_\tau=WN_-/D-2W
 \ge\frac{\inf N_-}{\ell_0^{-1}+\sup Q_-}-2W\ge c_0>0.
\]
A fixed logarithmic length therefore reaches that threshold; the fixed
positive margin supplies the remaining climb to the original \(w_0\).
Thus \(x_1\le La\), uniformly even when \(p\) tends to zero.
Choose \(g_0\) so that \(La<d_q\). Thereafter the original strip is
invariant and \(D\ge cW\). The limiting strong graph and contraction
estimate of the lemma give
\[
 |W(d_q)-W_-^*(d_q)|\le Cg+C(La/d_q)^c\longrightarrow0.
\]
This proves first passage and admission to the original observer buffer
before derivative estimates. Its initial height length satisfies
\[
 \eta_1-\eta_l
 =\log\frac{W_1}{\ell_0gp}+2\log\frac{x_1}{a}
 \le C(1+|\log g|+|\log p|).
\]

Write \(Q_-=q_0+q_1x+Cx^2\), where
\(q_0=1-Ag+Cg^2\), \(q_1=A-2Cg\).
At fixed \(\eta\) the exact partial derivatives are
\[
 D_\chi=-x-\pi kg^2/x-2k^2g^3/x^2-(2q_0+q_1x)W,
 \qquad N_{-,\chi}=-x+CgxW.
\]
On a sufficiently small fixed inner raw interval they imply
\begin{equation}\label{deg:negative-inverse-dissipation}
 V_\chi\le-c\Xi,\qquad |V_{\chi\chi}|\le C\Xi,
 \qquad \Xi=x+kg^2/x+k^2g^3/x^2+W.
\end{equation}
Indeed \(D\le gp+Q_-W\le Cg+CW\); the correction
\(-DN_{-,\chi}/N_-^2\) is absorbed by the negative terms in
\(D_\chi/N_-\). The next derivative has the same density.

Let \(X_j=\mathcal A^j\chi|_\eta\),
\(l_j=\mathcal A^j\eta_l\), and \(t_j=\mathcal A^j\eta_T\).
The full moving initial rows are
\begin{align*}
 l_1&=(2p+Ep)/B,\qquad l_2=O(p),\\
 X_{1,l}&=h-V_ll_1=pK_l>0,\\
 X_{2,l}&=\mathcal A X_{1,l}-V_{\chi,l}X_{1,l}l_1=O(p).
\end{align*}
with \(K_l\) a positive unit. Here \(l_1\), \(Eh\) and \(El_1\)
are bounded. The two inverse equations are \eqref{deg:scalar-inverse-endpoints}.
Up to that inner cut, \eqref{deg:negative-inverse-dissipation} gives
\(0<X_1\le Cp\), \(\int\Xi X_1\,d\eta\le Cp\), and
\(|X_2|\le CX_1\), by division by the positive first variation.
On the strong part choose an auxiliary fixed inner cut \(d_*\le d_q\)
so that \(V_\chi/V\le-3/2\), then shrink \(g_0\) to ensure
\(La<d_*\). The decay follows from
\(D\ge cW\), \(W\ge w_0\), and
\(V_\chi/V\le-2+Cg/w_0+Cx\).
Thus \(X_1\le Cp(a/x)^{3/2}\) up to that cut.
The remaining compact strong-graph segment to the original \(d_q\)
has bounded coefficients and finite logarithmic length, and preserves
this bound up to a fixed factor. At the original terminal cut, \eqref{deg:scalar-inverse-endpoints}
gives \(|t_1|+|t_2|\le Cp a^{3/2}\).
Use the exact density \(J_-=J\) from
\eqref{deg:quadratic-outer-density}, with
\(I=\int_{\eta_l}^{\eta_T}J\,d\eta\).
Its two complete rows are \eqref{deg:scalar-density-endpoints},
with \(D=\mathcal A\), \(r=\chi\), and \(Dr_l=h\).
On the initial part \(|\partial_\chi^jJ|\le Cg\), \(j=0,1,2\),
since \(x\le La\le Cg\) and \(kg^2/x\le g/r\).
Its differentiated integrals are at most
\(Cgp(1+|\log g|+|\log p|)\le Cg(1+|\log g|)\).
On the strong part \(d\eta\le C\,d\log x\), and the integrals are
bounded by
\[
 Cp\int_{La}^{d_*}(x+kg^2/x)(a/x)^{3/2}\,d\log x
 \le Cp(a+kg^2/a)\le Cg.
\]
If the actual entry is below \(La\), the fixed logarithmic segment
from \(x_1\) to \(La\) has the same \(Cg\) bound.
The compact final segment contributes \(Cp a^{3/2}\).
The displayed initial terms are \(O(g)\), and all terminal partials
are bounded. Therefore \(|\mathcal AI|+|\mathcal A^2I|\le CgL_g\).

Finally use the exact complete slope factor of the lemma.
Its moving-initial and \(1+\ell_0Q_-(a)\) logarithmic rows are
\(O(g)\). Although \(p^{-\nu}\) is unbounded as \(p\) tends to zero,
its two action logarithmic rows are \(O(g)\), since
\(\mathcal A\log p=Ep/B\) and \(\mathcal A^2\log p=O(p)\).
The terminal \(W_T/D_T\), \(z_T^\nu\) and ordinary \(C^3\) observer
rows are \(O(p a^{3/2})\), by the full terminal formulas.
The extracted \(a^{-\lambda}\) contributes \(-\lambda h=O(g)\)
and \(-\lambda\mathcal Ah=O(g)\).
All coefficient-only slope factors are frozen, including the change
between the two quadratic labels. These estimates prove
\eqref{deg:negative-absolute-direct-jets} on the whole original domain,
including arbitrarily close approach to its negative root endpoint.
\end{proof}

\subsection{The complete negative inward action remainder}

We next continue the negative lip back to the original centre and
estimate its action remainder. This supplies one side of the later
matching equation. The positive and negative heights are constructed
separately; matching is imposed only on the intersection of their images.

\begin{corollary}\label{deg:cor-negative-inward-common-centre}
Use \eqref{deg:quadratic-inward-family} with \(s=-1,x_s=r\),
\(\pi\in[\pi_\flat,\Pi]\), and fixed small \(\ell_0\).
Fix \(r_c>0\). On the entire strict negative interval
\(r\ge r_c\), \(0<R=kr<\mathfrak c\), continue the lip
\(Y=\ell_0P_-(r)\) to \(\xi=0\) along the inverse of the actual
positive-time negative branch. Its height \(Z_-(r)>0\) is defined on
this entire interval and enters the original central height window
after one coefficient shrink. In \eqref{deg:quadratic-inward-memory-common}, its remainder obeys
\begin{equation}\label{deg:negative-inward-central-remainder}
 \sum_{j=1}^3|\mathcal D_-^j\mathcal R_-|\le Cg.
\end{equation}
Changing the central height by a coefficient-only multiplicative scale changes
only a constant in this logarithmic formula. The estimate does not
differentiate coefficients or bound the value of \(\mathcal R_-\) by
\(Cg\). All constants are independent of the distance to the negative
simple root and of any comparison between \(g\) and \(k\).
\end{corollary}
\begin{proof}
Choose an auxiliary fixed \(0<b<r_c/2\). On \([0,r_c]\), the
polynomial is uniformly positive after decreasing \(k\), and
\[
 A_-(r_c)-A_-(b)\ge a_*>0 .
\]
For \(v\ge b\), \(kv<\mathfrak c\), write
\[
 Q=1-Ag(1-kv)+Cg^2(1-kv)^2,\quad
 H=\frac{P_-'(v)}{vB_-(v)},\quad
 \omega=-\frac{Ckg^3(1-kv)P_-(v)}{vB_-(v)} .
\]
Here \(Q\) is a positive unit. With
\(\mathcal D_v=P_-(v)\partial_v/(vB_-(v))\), direct rational
differentiation in \((v^{-1},kv)\) gives
\begin{equation}\label{deg:negative-inward-coefficient-rows}
 \sum_{j=0}^3|\mathcal D_v^jH|\le C,\qquad
 \sum_{j=1}^3|\mathcal D_v^jQ|\le Cg,\qquad
 \sum_{j=0}^3|\mathcal D_v^j\omega|\le Cg^3 .
\end{equation}
The field \(\mathcal D_v\) contains the factor \(P_-(v)\);
none of these bounds divides by its vanishing root factor.

Set \(t=A_-(r)\), \(q=(t-A_-(v))/g\), and
\(W=Y/P_-(v)\). The original negative scalar equation gives exactly
\[
 v_q=-g\frac{P_-(v)}{vB_-(v)},\qquad
 (\log W)_q=-\frac{1+\omega W}{1+QW}+gH,
 \qquad W(0)=\ell_0 .
\]
At fixed \(q\), \(\partial_t v=P_-(v)/(vB_-(v))\), so every
\(\partial_t\) coefficient derivative is precisely \(\mathcal D_v\).
The logarithmic rate and the normal diagonal are at most \(-c<0\)
for the fixed small lip. Lemma~\ref{deg:lem-stable-ratio-ledger}
applies with \(D=\partial_t\), \(\kappa=g\), \(\rho=g\),
\(\alpha=0\), by \eqref{deg:negative-inward-coefficient-rows};
clock differentiation acts as \(-g\mathcal D_v\). It constructs the
actual positive-height inverse passage to \(b\), with radial speed
\(P_-+QY>0\), and gives
\[
 |\partial_t^jW|+|\partial_t^jL|
 \le Cg(1+q)^Ne^{-cq},\quad1\le j\le3,\qquad
 L=(Q-\omega)W/(1+QW).
\]
Retain the exact source polynomial:
\[
 T=(t-A_-(b))/g,\qquad V_b=\int_0^{T(t)}L\,dq,\qquad
 Y_b=\ell_0P_-(r)e^{-T+V_b}=P_-(b)W(T)\le Ce^{-cT}.
\]
Here \(T'=1/g,T''=T'''=0\), and \(T\ge a_*/g\).
The three complete rows \eqref{deg:stable-ratio-terminal-rows}
therefore give \(|\partial_t^jV_b|\le Cg\): each boundary term
is bounded by \(Cg^{-N}(1+T)^Ne^{-cT}\). The retained polynomial
gives
\[
 \mathcal D_-\log Y_b=H(r)-g^{-1}+\partial_tV_b,
 \qquad |\mathcal D_-^jY_b|\le Cg^{-j}Y_b,\quad j\le3.
\]

Reflect the segment from \(-b\) to the original centre by \(v=-\xi\).
Preparation~\eqref{deg:positive-large-S-inward-compact-graph}, with \(s=-1\),
then applies on the fixed margin above. Exponential flatness gives
\(Y_b\le h_b/2\) and small \(Y_b/g\), so
Lemma~\ref{deg:lem-compact-height-continuation} establishes the actual
continuation, \(0<Y\le2Y_b\), and its logarithmic nonlinear unit
\(\mathcal N\). The three composition rows of that lemma are
bounded by \(CY_b/g^{j+1}\le Cg\).
The exact linear multiplier is \(e^{-A_-(b)/g}\), so the artificial
cut cancels and
\(\mathcal R_-=V_b+\mathcal N(Y_b)\) gives
\eqref{deg:negative-inward-central-remainder} at the original centre.
\end{proof}

\subsection{Positive inward passage on the growing range}

When \(S=ku\) grows, the positive action speed is of order
\(1+S\). We therefore estimate the inward remainder in the radial
derivative \(u\partial_u\), keeping the conversion to action
derivatives explicit. This is the estimate used in
\eqref{deg:quadratic-large-S-curvature}.

For \(s=+1\), the factorization is
\[
 P(v)=(\mathfrak c+kv)\{(v-d)^2-\Delta\},\quad
 \mathfrak c=3+2kd,\quad \mathfrak c(d^2-\Delta)=1,\quad
 \pi=2\mathfrak cd-k/\mathfrak c,
\]
with \(d=\pi/6+O(k)\), \(\mathfrak c=3+O(k)\).

\begin{corollary}[Positive large-\(S\) inward action rows]
\label{deg:cor-positive-large-S-inward}
Fix \(0<\pi_\flat\le\Pi<\infty\), a compact set of the passive
coefficients \(A,C\), and \(\delta>0\). Choose the original lip
height \(\ell_0\) and physical collar width \(d_0\) sufficiently small
within their original buffered windows. There are fixed
\[
 v_i,M,g_0,k_0>0,\qquad
 k_0\le\delta/(2M),
\]
such that the following holds for independently small
\(0<g<g_0\), \(0<k<k_0\), \(\pi\in[\pi_\flat,\Pi]\).
On every strict root-free fibre \(\Delta<0\), every lip point
\begin{equation}\label{deg:positive-large-S-inward-domain}
 \xi=u,\qquad Y=\ell_0P(u),\qquad
 S=ku\ge\delta,\qquad a=gS<d_0
\end{equation}
has an actual positive-height, positive-time first passage to \(\xi=0\).
Its original centre height is \(Y_{0,+}=Z_+>0\), with
\(\mathcal R_+\) defined by \eqref{deg:quadratic-inward-memory-common}.
With all coefficients fixed under \({\mathcal E}=u\partial_u\),
\begin{equation}\label{deg:positive-large-S-inward-action3}
 |{\mathcal E}^j\mathcal R_+(u)|\le C_\delta a,\qquad j=1,2,3 .
\end{equation}
The constants are independent of \(g,k\) and of \(-\Delta>0\).
The statement can be restricted to the original strict raw-window
preimages without extending any image or inverse.

On a real or double-root fibre the prescribed positive through lip ends
at the first root \(a_{\rm root}=d-\sqrt{\Delta}\). That endpoint is
bounded uniformly in this coefficient block, and after the same
\(k\)-shrink its through interval has \(ku<\delta\). Thus
\eqref{deg:positive-large-S-inward-domain} is empty on that branch. No passage beyond the first
root, or at the stopped \(\Delta=0\) limit, is included.
\end{corollary}

\begin{proof}
Fix the finite differentiability budget used below first.
Choose \(v_i>0\) so small that
\[
 P(v)\ge p_i>0\qquad(0\le v\le2v_i)
\]
uniformly for the fixed \(\pi\)-interval and small \(k\).
Then fix \(M>2v_i\) beyond all positive roots and turning points of
the retained compact positive factors, and large enough that
\begin{equation}\label{deg:positive-large-S-inward-growing-polynomial}
 P(v)\asymp v^2(3+kv),\qquad v\ge M .
\end{equation}
The constants may depend on the fixed passive compact,
\(\pi_\flat,\Pi,\delta,\ell_0,v_i,M\).
All calculations are also valid on a slightly enlarged phase range
\(S>\delta/2\), \(a<2d_0\); the choices of \(d_0,g_0,k_0\) are made
for that range before restriction to \eqref{deg:positive-large-S-inward-domain}.
This supplies actual phase buffers for the displayed derivatives.

\smallskip
\noindent\emph{Actual passage and value admission.}
For \(\Delta<0\), \(P>0\) on the positive half-axis.
On \(v_i\le v\le u\) put
\begin{equation}\label{deg:positive-large-S-inward-action-variable}
 q=\frac{J(u)-J(v)}g,\qquad
 v_q=-g\frac{P(v)}{vB(v)},\qquad
 W=\frac{Y}{P(v)},\qquad
 H=\frac{P'(v)}{vB(v)},\quad
 \varpi=\frac{\zeta(v)P(v)}{vB(v)} .
\end{equation}
Division of \eqref{deg:quadratic-inward-scalar-common}, retaining the polynomial factor, gives
exactly
\begin{equation}\label{deg:positive-large-S-inward-ratio}
 W_q=W\left\{-\frac{1+\varpi W}{1+QW}+gH\right\},
 \qquad W(0)=\ell_0 .
\end{equation}
Here \(B,Q\) are positive units, \(gH\le C_\delta a\), and
\(|\varpi|\le C_\delta a^3\), as also checked below.
Choose \(\ell_0,d_0,g_0\) so that the logarithmic rate in
\eqref{deg:positive-large-S-inward-ratio} is at most \(-c_0<0\) while \(0<W\le\ell_0\).
A first-exit argument proves
\begin{equation}\label{deg:positive-large-S-inward-value-decay}
 0<W(q)\le\ell_0e^{-c_0q},\qquad
 0\le q\le T(u):=\frac{J(u)-J(v_i)}g .
\end{equation}
Consequently \(P+QY=P(1+QW)>0\). For each strict ghost coefficient
the finite interval \([v_i,u]\) has a positive minimum of \(P\);
the time integral to \(v_i\) is finite. Height is positive by
uniqueness. The same construction continues to \(v_i/2\), on its
fixed positive-\(P\) buffer, before imposing the cut \(v_i\).
Thus the fixed-\(q\) partials at that moving terminal are genuine
nearby actual-flow derivatives.

Since \(Y\le\ell_0P(v)\), the physical height on the entire arc is
at most
\[
 C\ell_0\{g^3k^2+g^3k^2v+g^3k^2v^2+g^3k^3v^3\}
 \le C\ell_0\{g^3k^2+g^2ka+ga^2+a^3\}.
\]
Together with \(0<gkv\le a<d_0\), this admits the arc to a
fixed small original chart. Neither a limiting axis orbit nor a
normalization is used to infer this admission.

On \(M\le v\le\delta/k\), \(P(v)\le C_\delta v^2\) and \(B\ge1\).
Since \(u\ge\delta/k\), it follows that
\begin{equation}\label{deg:positive-large-S-inward-action-gap}
 T(u)\ge\frac{c_\delta\log(1/k)-C_\delta}{g}.
\end{equation}
Choose \(k_0\) so that the numerator has a fixed positive margin.
Then \(T\) is uniformly large for small \(g\), with no comparison
between \(g\) and \(k\). In particular the actual terminal height
\begin{equation}\label{deg:positive-large-S-inward-inner-value}
 y_i(u):=Y(v_i;u)=P(v_i)W(T(u),u)
              \le C e^{-c_0T(u)}
\end{equation}
lies in any fixed sufficiently small compact-entry interval
\(0<y_i<y_*\) after the permitted coefficient shrink.
This places the input in the prescribed compact-continuation buffer
before its derivatives are estimated.

On \([0,v_i]\), preparation~\eqref{deg:positive-large-S-inward-compact-graph}
and exponential flatness give the prescribed height buffer and small
\(y_i/g\). Lemma~\ref{deg:lem-compact-height-continuation} therefore
continues the actual graph to the original centre with \(0<Y\le2y_i\),
for either sign of \(C\).
Retain the exact memory~\eqref{deg:stable-ratio-polynomial-memory},
with \(V(q,u)=\int_0^qL\,ds\), \(L=(Q-\varpi)W/(1+QW)\);
its initial factor is \(P(u)/P(v)\) and \(|V|\le C\ell_0\).

\smallskip
\noindent\emph{Three transported coefficient rows.}
Write
\[
 h(v)=\frac{P(v)}{v^2B(v)},\qquad
 {\mathcal D}_v=\frac{P(v)}{vB(v)}\partial_v=h(v)v\partial_v,\qquad
 b(u)=\frac1{h(u)} .
\]
At fixed elapsed action \(q\), differentiating
\(J(v)=J(u)-gq\) yields
\begin{equation}\label{deg:positive-large-S-inward-transport}
 {\mathcal E}\log v=\frac{h(v)}{h(u)},\qquad
 {\mathcal E}[F(v(q,u))]=b{\mathcal D}_vF .
\end{equation}
The full second and third transported rows are
\begin{equation}\label{deg:positive-large-S-inward-transport3}
\begin{split}
 {\mathcal E}^2F&=b^2{\mathcal D}_v^2F+({\mathcal E} b){\mathcal D}_vF,\\
 {\mathcal E}^3F&=b^3{\mathcal D}_v^3F+3b({\mathcal E} b){\mathcal D}_v^2F+({\mathcal E}^2b){\mathcal D}_vF .
\end{split}
\end{equation}
The variable coefficients of \({\mathcal D}_v\) are differentiated in its
powers. On the initial domain,
\begin{equation}\label{deg:positive-large-S-inward-initial-speed}
 h(u)=\frac{3-\pi/u+u^{-2}+S}{1+2g+gS}
       \asymp_\delta1+S,\qquad
 |{\mathcal E}^jh(u)|\le C_\delta h(u)\quad(j=1,2).
\end{equation}
Hence
\[
 |{\mathcal E} b|+|{\mathcal E}^2b|\le C_\delta b,\qquad
 b\le C_\delta/(1+S),\qquad
 gh(u)\asymp_\delta g+a .
\]

For \(v\ge M\), let \(t=1/v\), \(s=kv\), and
\({\mathcal E}_v=-t\partial_t+s\partial_s\). Then
\begin{equation}\label{deg:positive-large-S-inward-rational-symbols}
 h(v)=\frac{3-\pi t+t^2+s}{1+2g+gs},\qquad
 H(v)=\frac{6-\pi t+3s}{1+2g+gs},\qquad
 \varpi(v)=Cg^3s(1+s)h(v).
\end{equation}
For \(0\le s\le S\), \(t\le1/M\), and \(gS=a<2d_0\),
these rational symbols give
\[
 h(v)\le C_\delta h(u),\qquad
 |{\mathcal E}_v^jh(v)|\le C_\delta h(v)\quad(j\le2).
\]
The first three \({\mathcal E}_v\) rows of \(gH,Q-1,\varpi\) have
size \(C_\delta a\); the \(\varpi\) rows have the sharper
size \(C_\delta a^3\). Expanding \((h(v){\mathcal E}_v)^j\) and
using \(b h(v)\le C_\delta\) proves the same bounds for
the rows in \eqref{deg:positive-large-S-inward-transport3}.

On \(v_i/2\le v\le M\), \(H\) has uniform ordinary finite
derivatives. The coefficients and finite ordinary derivatives of
\({\mathcal D}_v=P/(vB)\partial_v\) are uniformly bounded, also as
\(-\Delta\downarrow0\); they contain \(P\) rather than \(P^{-1}\)
at the bottleneck. The corresponding ordinary rows of \(gH,Q-1\)
are \(O(g)\), and those of \(\varpi\) are \(O(kg^3)\).
Since \(g\le a/\delta\) on the retained domain, this proves
\begin{equation}\label{deg:positive-large-S-inward-coefficient-Euler}
 |{\mathcal E}^j(gH)|+|{\mathcal E}^j(Q-1)|+|{\mathcal E}^j\varpi|
      \le C_\delta a,\qquad 0\le j\le3,
\end{equation}
uniformly through the arbitrary strict ghost gap.
In particular the quotient in \eqref{deg:positive-large-S-inward-transport} and its first
two Euler rows are bounded on both the growing and compact portions.

For the time rows, \(\partial_q=-g{\mathcal D}_v\) on coefficient
functions and \(gh(v)\le C_\delta a\) on the growing portion.
On the compact portion \(g{\mathcal D}_v\) has bounded finite coefficients.
The same rational calculations give bounded first and second
\(q\)-rows and the mixed \(q,{\mathcal E}\) rows required below.
Their rows containing an \({\mathcal E}\) retain the \(C_\delta a\) factor.
No inverse of the bottleneck polynomial is introduced.

\smallskip
\noindent\emph{Normal variations and the moving terminal.}
The exact diagonal is
\begin{equation}\label{deg:positive-large-S-inward-negative-diagonal}
 \mathscr F_W=-\frac{1+2\varpi W+\varpi QW^2}{(1+QW)^2}+gH
 \le-c_1<0.
\end{equation}
Thus Lemma~\ref{deg:lem-stable-ratio-ledger} applies with
\(D=\mathcal E,\kappa=g,\rho=C_\delta a,\alpha=0\), by
\eqref{deg:positive-large-S-inward-coefficient-Euler} and the mixed
clock bounds above. Its fixed-clock estimate is
\begin{equation}\label{deg:positive-large-S-inward-W-Euler}
 |\mathcal E^jW|+|\mathcal E^jL|
 \le C_\delta a(1+q)^Ne^{-cq},\qquad1\le j\le3,
\end{equation}
with the mixed and pure clock rows specified there. The terminal rows are
\begin{equation}\label{deg:positive-large-S-inward-terminal-clock}
 T_1=b/g,\qquad T_2=(\mathcal E b)/g,\qquad
 T_3=(\mathcal E^2b)/g,\qquad
 |T_j|\le C_\delta/(gh(u))\le C_\delta/a\le C_\delta/g .
\end{equation}
By \eqref{deg:positive-large-S-inward-action-gap}, for each fixed \(N\),
\begin{equation}\label{deg:positive-large-S-inward-terminal-absorption}
 g^{-N}(1+T)^Ne^{-cT}\le C_{\delta,N}g\le C_{\delta,N}a.
\end{equation}
Choose \(k_0\) with \(cc_\delta\log(1/k_0)-cC_\delta>0\);
then \(g^{-N}e^{-\beta/g}\) absorbs every terminal loss, using
independent upper bounds on \(g,k\).
The terminal formulas~\eqref{deg:stable-ratio-terminal-rows}, for
\(V_i=V(T(u),u)\), give
\begin{equation}\label{deg:positive-large-S-inward-V-action3}
 |\mathcal E^jV_i|\le C_\delta a,\qquad1\le j\le3.
\end{equation}
The exact memory, including all moving terminal rows, is
\begin{equation}\label{deg:positive-large-S-inward-height-memory}
 y_i=\ell_0P(u)e^{-T+V_i}=P(v_i)W(T),\qquad
 \mathcal E^j\log y_i=\mathcal E^j\log P(u)-T_j+\mathcal E^jV_i.
\end{equation}
The rational symbols bound \(|\mathcal E^j\log P(u)|\le C_\delta\),
hence \(|\mathcal E^j\log y_i|\le C_\delta/g\) for \(1\le j\le3\).
The value estimate uses \(P(v_i)W(T)\), independently of \(P(u)\).
The compact continuation already constructed has exact multiplier
\(m(0)=e^{-J(v_i)/g}\). Its three composition rows from
Lemma~\ref{deg:lem-compact-height-continuation} are at most
\(Cy_i/g^4\le C_\delta a\), by
\eqref{deg:positive-large-S-inward-terminal-absorption}. Consequently
\[
 \log Y_{0,+}=\log\{\ell_0P(u)\}-J(u)/g+V_i+\mathcal N(y_i),
 \qquad \mathcal R_+=V_i+\mathcal N(y_i).
\]
This proves the three Euler bounds at the original centre. The fixed-cut
action cancels exactly; an additive action normalization changes only a
coefficient constant in \(\mathcal R_+\).

Finally, \(d,\mathfrak c\), and every nonnegative
\(\sqrt{\Delta}\) in the fixed \(\pi\)-block are bounded for
small \(k\). The root \(a_{\rm root}\) is therefore bounded by
the fixed \(M\). Its prescribed through lip has
\(ku\le ka_{\rm root}<\delta\) after the \(k\)-shrink, as stated.
The strict ghost proof permits arbitrarily small \(-\Delta>0\)
but does not add a through interval at \(\Delta=0\).
\end{proof}

\subsection{Bounded positive quadratic passages}

On a bounded positive lip interval, both the root gap and the lip's
distance to its first root can vanish. The inward calculation remains
uniform in action coordinates. For the outward calculation we use
logarithmic height through the weak stage; its moving initial factor
absorbs the long dwell before the graph enters the parabolic strip.

\begin{corollary}[Bounded positive quadratic suppliers]
\label{deg:cor-positive-compact-suppliers}
Fix the original positive raw section and its buffered regular observer,
with its original ordinary \(C^3\) logarithmic-slope reserve
\(\|\partial_z\log\Phi'\|_{C^1}\le C\) on the buffered raw image,
a compact neighborhood of the passive coefficients \(A,C\), and
\(\pi\in[\pi_\flat,\Pi]\). Fix \(U<\infty\) and use
\eqref{deg:quadratic-inward-family}--\eqref{deg:quadratic-inward-scalar-common},
with physical coordinates \(x=gkv,z=g^3k^2Y\).
Retain the original actual moving-lip interval \(Y=\ell_0P(u)\),
intersected with \(u\le U\). Its lower threshold satisfies
\(u_{\rm th}\ge u_0>0\). On real or fold fibres this interval has
\(u_{\rm th}\le u<a\), where \(a\) is the first positive root.
On root-free fibres it is the component of the original section-image
preimage specified in the positive quadratic domain construction.
In each case \(P>0\) on the actual inward axis path \([0,u]\).
No right-hand axis continuation is added at the fold.

There are \(g_0,k_0,C_0>0\), uniform on these domains, such that
the following statements hold for independently small
\(0<g<g_0\), \(0<k<k_0\). Put \(L_g=1+|\log g|\), with \(J,\mathcal D_+\) as in
\eqref{deg:quadratic-inward-memory-common}.
The inward lip orbit reaches the original centre \(v=0\), and its
height \(Z_+(u)=Y(0;u)>0\) has
\begin{equation}\label{deg:positive-compact-central-jets}
 |\mathcal D_+^j\mathcal R_+|\le C_0g,\qquad1\le j\le3.
\end{equation}
The outward lip graph, which is the inverse of the physical positive
incoming passage, reaches the original positive raw port at positive
height and inside its original buffer. For the complete embedding
followed by the original increasing ordinary \(C^3\) observer, use
the lip label \(\beta=-u^2/2\) and any coefficient-only positive
affine value and slope calibration \(\phi_+\). Then
\begin{equation}
 \phi_+'>0,\qquad
 |\mathcal D_+^j\log\phi_+'|\le C_0gL_g,\quad j=1,2.
 \label{deg:positive-compact-outer-jets}
\end{equation}
The prime denotes differentiation in \(\beta\).
Both estimates freeze \(g,k,\pi,A,C\). They are uniform as the
positive root gap and the distance of the lip to its first root tend
to zero independently of \(g,k\). They do not require an intersection
of positive and negative inward images.
\end{corollary}

\begin{proof}
The exact positive quadratic factorization is
\[
 P(v)=(\mathfrak c+kv)\{(v-d)^2-\Delta\},\qquad
 \mathfrak c=3+2kd,\quad
 \mathfrak c(d^2-\Delta)=1,\quad
 \pi=2\mathfrak cd-k/\mathfrak c.
\]
It gives uniform bounds for
\(d,d^{-1},\mathfrak c,\mathfrak c^{-1}\) on the retained fold band.
On a real fibre, \(P(a)=0\) implies
\[
 1=\pi a-3a^2-ka^3\le\Pi a,\qquad a\ge\Pi^{-1}.
\]
On a root-free fibre,
\(\pi=2\mathfrak cd-k/\mathfrak c\) bounds \(d\) away from zero.
Thus the original fixed \(X_{\rm up}\) can be chosen small enough
that \(u_{\rm th}=a-X_{\rm up}\), respectively \(d-X_{\rm up}\),
has the asserted uniform positive lower bound. All constants below
are fixed before \(g,k\).

\smallskip
\noindent\emph{The inward action integral and compact continuation.}
Choose fixed \(b<u_0/2\) with \(P\ge p_b>0\) on \([0,b]\), and put
\[
 \begin{gathered}
 q=(J(u)-J(v))/g,\quad W=Y/P(v),\quad H=P'/(vB),\\
 \omega=\zeta P/(vB),\qquad
 \mathcal D_v=P(v)\partial_v/(vB(v)).
 \end{gathered}
\]
Use the ratio equation and exact memory
\eqref{deg:stable-ratio-polynomial-memory}, with \(\kappa=g\)
and \(V(q)=\int_0^q(Q-\omega)W/(1+QW)\,ds\).
At fixed \(q\), \(\mathcal D_+v=P(v)/(vB(v))\), so its repeated
coefficient rows are powers of \(\mathcal D_v\). On \([b,U]\),
this field has bounded finite derivatives, uniformly at the root/fold:
\begin{equation}\label{deg:positive-compact-inward-coefficients}
 |\mathcal D_v^j(gH)|+|\mathcal D_v^j(Q-1)|+|\mathcal D_v^j\omega|
 \le Cg,\qquad0\le j\le3.
\end{equation}
Clock rows use \(-g\mathcal D_v\), and the fixed small lip gives
both negative margins. These are the inputs of the negative inward proof,
now with \(s=+1\), \(D=\mathcal D_+\), \(\kappa=g,\rho=g,\alpha=0\).
The terminal check is
\[
 q_b=(J(u)-J(b))/g\ge\eta/g,\qquad
 \mathcal D_+q_b=1/g,\quad \mathcal D_+^2q_b=\mathcal D_+^3q_b=0;
\]
the lower bound follows on a fixed subinterval of \([b,u_0]\).
Lemma~\ref{deg:lem-stable-ratio-ledger} and its terminal formulas give
\(|\mathcal D_+^jV(q_b)|\le Cg\), \(1\le j\le3\), and
\begin{equation}\label{deg:positive-compact-central-incoming-height}
 \begin{gathered}
 y_b=\ell_0P(u)e^{-q_b+V(q_b)}=P(b)W(q_b)\le Ce^{-c\eta/g},\\
 \mathcal D_+\log y_b=-g^{-1}+H(u)+\mathcal D_+V(q_b),\qquad
 |\mathcal D_+^jy_b|\le C_jg^{-j}y_b\quad(1\le j\le3).
 \end{gathered}
\end{equation}
Preparation~\eqref{deg:positive-large-S-inward-compact-graph} verifies
the compact-continuation inputs: \(P\ge p_b\),
\(|\partial_Y^jK|\le C/g\), \(y_b\le h_b/2\), and small \(y_b/g\).
Lemma~\ref{deg:lem-compact-height-continuation} then gives
\[
 Z_+(u)=e^{-J(b)/g}y_be^{\mathcal N(y_b)},\qquad
 \mathcal R_+=V(q_b)+\mathcal N(y_b).
\]
Its three composition rows are \(O(y_b/g^{j+1})=O(g)\).
The cut action cancels, proving~\eqref{deg:positive-compact-central-jets}.

\smallskip
\noindent\emph{Construction of the outward graph and buffer admission.}
Put
\[
 a=gku,\qquad p=\frac{P(u)}{u^2}>0,\qquad
 E=u\partial_u,\qquad h=\frac p{B(u)},\qquad \mathcal D_+=hE.
\]
In the raw graph set \(W=z/x^2\), \(r=\log x\), and \(\eta=\log z\).
Its exact coefficients are
\[
 \begin{split}
 m(x)&=3g+x-\pi kg^2/x+k^2g^3/x^2,\\
 Q_0(x)&=1-A(g+x)+C(g+x)^2,\\
 N(x,W)&=1+2g+x+Cx(g+x)W,\qquad D_0=m+Q_0W.
 \end{split}
\]
The initial data are the original lip data
\begin{equation}
 x_l=a,\qquad W_l=\ell_0gp,\qquad
 z_l=a^2\ell_0gp=g^3k^2\ell_0P(u).
 \label{deg:positive-compact-raw-initial}
\end{equation}
Use the inverse logarithmic-height equation first:
\begin{equation}
 r_\eta=V=\frac{D_0}N,\qquad W_\eta=W(1-2V).
 \label{deg:positive-compact-inverse-height}
\end{equation}
Here \(N,Q_0\) are positive units on the original small corridor.
At a possible first zero of \(D_0\), one has \(V=0\) and
\[
 \frac{dD_0}{d\eta}=Q_0W>0.
\]
Thus a graph starting with \(D_0>0\) cannot cross to \(D_0\le0\).
Both \(x\) and \(z\) increase on this graph, including at an axis
root crossed at positive height.

Fix \(L=2\), then a sufficiently small auxiliary height \(w_0>0\).
In the box \(a\le x\le La\), \(0<W\le w_0\), the variable
\(x/(gk)\) lies in the fixed interval \([u_0,LU]\). Hence
\[
 |m|+|m_r|+|m_{rr}|\le Cg,\qquad
 0<V\le C(g+w_0),\qquad \tfrac12W\le W_\eta\le W.
\]
Write \(s=\eta-\eta_l\). The length of this weak stage, up to either
possible first exit, is bounded by
\begin{equation}
 T\le C\{L_g+\log^+(1/p)\}.
 \label{deg:positive-compact-weak-length}
\end{equation}
Since \(m(0)=gp\) and
\[
 |m_s|\le Cg(|m|+W),
\]
Gronwall's inequality, together with the lower growth rate for \(W\),
gives
\[
 |m(s)|\le Cgp e^{Cgs}+CgW(s),\qquad
 r(s)-r_l\le Cgp(1+s)e^{Cgs}+Cw_0.
\]
For every fixed \(j\), the elementary bound
\begin{equation}
 p(1+T)^j e^{CgT}\le C_jL_g^j,\qquad 0<p\le p_{\max},
 \label{deg:positive-compact-polynomial-absorption}
\end{equation}
holds after shrinking \(g\).
Indeed \eqref{deg:positive-compact-weak-length} reduces it to
\(g^{-Cg}\le C\) and the boundedness of
\(p^{1-Cg}\{1+\log^+(1/p)\}^j\).
Choose \(w_0\) so \(Cw_0<\frac12\log L\), then \(g\) so that
\(CgL_g<\frac12\log L\).
The preceding displacement estimate excludes an exit at \(x=La\).
Thus the graph reaches \(W=w_0\) at a finite height time and an
abscissa \(a_*\in[a,La]\), for every \(p>0\).

Keep the original raw cut \(d_q\) fixed. Choose an auxiliary cut
\(0<d_*<d_q\) sufficiently small compared with \(w_0\), before
shrinking \(g,k\). After this shrink \(a_*<d_*/2\).
On \(a_*\le x\le d_*\) use
\[
 W_{\log x}=F=W(N/D_0-2).
\]
The strip \(w_0\le W\le w_1\), for a fixed \(w_1>1/2\), is
invariant, with \(D_0\ge cW\).
Moreover \(F_W\le-c\) and \(V_r/V\le-\alpha\), where
\(\alpha=3/2\), because
\(|m|+|m_r|\le C(g+x)\) on this segment and \(g,d_*\) are
sufficiently small compared with \(w_0\).
The limiting strong graph at \(g=k=0\) has
\(W^*(x)=1/2+O(x)\): with \(U_*=W^*-1/2\) its equation is
\(xU_*'=-2U_*+O(x)\), whose bounded solution follows from the
integral equation at zero.
Since \(x\ge a\ge gku_0\), the strip equation differs by \(O(g)\);
variation of constants gives
\[
 |W(d_*)-W^*(d_*)|\le Cg+C(a_*/d_*)^c.
\]
On the remaining fixed segment \([d_*,d_q]\),
\(F=1-2W+O(g+x)\) retains barriers \(W=1/4,3/4\), with
\(D_0,N,V\) positive units and bounded derivatives.
Continuous dependence gives
\begin{equation}
 |W(d_q)-W^*(d_q)|\le Cg+C(a_*/d_*)^c\longrightarrow0
 \label{deg:positive-compact-original-buffer}
\end{equation}
uniformly in the independent scales, because \(a_*\le LgkU\).
The original passive-coefficient and raw-buffer margins admit
these terminal values and the regular continuation to the positive port.

\smallskip
\noindent\emph{The initial action factor and the full terminal variations.}
Let \(P_j=\mathcal D_+^jr|_\eta\), \(j=1,2\), and
\(l_j=\mathcal D_+^j\eta_l\). From the original initial height,
\[
 l_1=H(u),\qquad l_2=\mathcal D_+H(u),\qquad
 |l_1|\le C,\quad |l_2|\le Cp,\qquad
 \mathcal D_+r_l=h,\quad |\mathcal D_+h|\le Cp.
\]
Both moving initial rows are
\begin{equation}
 P_{1,l}=h-V_ll_1\asymp p>0,\qquad
 P_{2,l}=\mathcal D_+(P_{1,l})-V_{r,l}P_{1,l}l_1=O(p),
 \qquad
 V_l=\frac{gp(1+\ell_0Q_0(a))}{N_l}.
 \label{deg:positive-compact-inverse-initial-rows}
\end{equation}
The derivative of \(P_{1,l}\) here includes both moving initial
coordinates. The inverse equations \eqref{deg:scalar-inverse-endpoints} give
\((P_2/P_1)_\eta=V_{rr}P_1\).
On the weak box, direct differentiation with \(W_r=-2W\) gives
\[
 V_r\le Cg,\qquad |V_{rr}|\le C(g+W),\qquad
 0<P_1(s)\le Cp e^{Cgs}.
\]
The integral controlling \(P_2/P_1\) is at most
\[
 CgpT e^{CgT}+Cp e^{CgT}\int W\,d\eta\le C
\]
by \eqref{deg:positive-compact-polynomial-absorption} and
\(\int W\,d\eta\le Cw_0\). Thus \(|P_2|\le CP_1\) there.
On \(a_*\le x\le d_*\), the negative ratio \(V_r/V\) gives
\begin{equation}
 0<P_1(x)\le Cp e^{CgT}(a_*/x)^\alpha,\qquad
 |P_2(x)|\le CP_1(x).
 \label{deg:positive-compact-inverse-decay}
\end{equation}
The second bound follows by integrating \(|V_{rr}|P_1\,d\eta\)
with \(V\) a unit. At \(d_*\) both rows are \(O(a^\alpha)\).
The remaining fixed segment has bounded \(\eta\)-length, since
\[
 \eta_T-\eta_*=2\log(d_q/d_*)+\log(W_T/W_*).
\]
The barriers bound the last logarithm, so the inverse equations
transfer both \(O(a^\alpha)\) rows to the original cut.

At the original fixed terminal \(x=d_q\), both full implicit rows
\eqref{deg:scalar-inverse-endpoints}, with \(D=\mathcal D_+\),
give \(|t_1|+|t_2|\le Ca^\alpha\le Cg\), since \(V_T\) is a unit.
The initial factor \(p\) absorbs the arbitrary \(|\log p|\)
weak dwell. The weak-stage bound \(V_r\le Cg\) suffices, so this
argument also applies in the compact ghost box.

\smallskip
\noindent\emph{The exact moving factor and differentiated density.}
Use \eqref{deg:quadratic-outer-moving-factor}--\eqref{deg:quadratic-outer-density}
with \(s=+1,x_0=u,p=P(u)/u^2\). These are exact identities of
this same field; only their present bounded-domain estimates are used.
Write \(\mathscr B=\mathscr B_+\), \(\mathscr J=J_+\),
\(\nu=6g/(1+2g)\), \(\lambda=2\nu\). Direct rational differentiation gives
\(|\mathcal D_+^j\log(-\mathscr B)|\le Cgp\), \(j=1,2\).
At fixed \(\eta\), \(W_r=-2W\), and direct differentiation gives
\[
 |\partial_r^j\mathscr J|\le C(x+kg^2/x),\quad 0\le j\le2,
 \qquad |\mathscr J_\eta|\le CxW.
\]
Let
\[
 I(u)=\int_{\eta_l(u)}^{\eta_T(u)}
            \mathscr J(r(u,\eta),\eta)\,d\eta.
\]
Its two complete action rows are \eqref{deg:scalar-density-endpoints},
with \(D=\mathcal D_+\), \(Dr_l=h\), and the fixed terminal \(r_T\).
On the weak box, \(x+kg^2/x\le Cg\).
The integral bounds just established give
\[
 CgpT e^{CgT}\le CgL_g
\]
for every integral in \eqref{deg:scalar-density-endpoints};
here \(P_1^2\le CP_1\) and \(|P_2|\le CP_1\).
On \(a_*\le x\le d_*\), the common bound is
\[
 Cp e^{CgT}\int_{a_*}^{d_*}
 (x+kg^2/x)(a_*/x)^\alpha\,d\log x
 \le Cp e^{CgT}(a_*+kg^2/a_*)\le Cg .
\]
This uses \(\alpha>1\), \(a_*\asymp a\), and
\(kg^2/a=g/u\le Cg\).
On the remaining fixed strong segment the integrals are bounded
by \(Ca^\alpha\).
The initial boundary terms are \(O(g)\); the full terminal terms
are \(O(a^\alpha)\) by
\eqref{deg:scalar-inverse-endpoints}.
Hence
\begin{equation}
 |\mathcal D_+I|+|\mathcal D_+^2I|\le CgL_g.
 \label{deg:positive-compact-density-bound}
\end{equation}
At artificial splits the full Leibniz terms cancel in this single integral.

\smallskip
\noindent\emph{The complete slope on the original observer image.}
Use~\eqref{deg:quadratic-outer-exact-slope} with \(s=+1,x_0=u\),
including the regular continuation in \(\Phi\) on the admitted
image~\eqref{deg:positive-compact-original-buffer}. Normalize its
coefficient-only reference slope as in that formula; the original
ordinary \(C^3\) logarithmic-slope bound still holds. The initial rows are
\[
 \mathcal D_+\log p=Ep/B=O(1),\qquad
 \mathcal D_+^2\log p=O(p),\qquad
 \mathcal D_+\log a=h,\quad \mathcal D_+^2\log a=\mathcal D_+h=O(1).
\]
Thus the powers \(p^{-\nu}\) and \(a^{-\lambda}\) contribute
\(O(g)\); the rows of \(\log(1+\ell_0Q_0(a))\) and
\(\log(-\mathscr B)\) are \(O(gp)\).
The full terminal formulas~\eqref{deg:scalar-inverse-endpoints} and
the ordinary \(C^3\) observer bound give \(O(a^\alpha)\) for both
logarithmic rows of the terminal factors. Combining these estimates
with~\eqref{deg:positive-compact-density-bound} proves
\eqref{deg:positive-compact-outer-jets}, uniformly as \(p\to0\).
The exact slope retains \(d_\theta=c_\theta(gk)^2\ell_0^{-1-\nu}g^{-\nu}\);
frozen affine calibration preserves its logarithmic rows.
\end{proof}

\subsection{Separate inward jets in the joint quadratic corner}

The derivative of the curvature error involves second derivatives of
the relative matching slope. To control these when both lips grow,
we construct three action derivatives of each inward remainder on its
own domain. Their quotient then has the required mixed derivatives on
the product, before equality of the two actual central heights is imposed.

\begin{lemma}\label{deg:lem-joint-inward-product-direct}
Use the original field \eqref{deg:quadratic-field}, its moving lips
\(Y=\ell_0P\), the original central section \(\xi=0\), and a fixed
compact passive \(A,C\) neighborhood. Fix the original sufficiently
small \(\delta_q>0\), \(\pi\in[\pi_\flat,\Pi]\), and a sufficiently
large fixed \(M\). For independently small \(g,k>0\), restrict the
positive side to strict root-free fibres \(\Delta<0\). On each separate
original admissible lip domain with
\[
 u,r\ge2M,\qquad ku,kr\le\delta_q,
\]
the positive inward passage and the inverse negative passage reach
the same original central section at actual positive heights
\(Z_+(u)\), \(Z_-(r)\). Use \eqref{deg:quadratic-inward-family}--\eqref{deg:quadratic-inward-memory-common}
with \(x_+=u,x_-=r\), \(\theta_s=x_s^{-1}+kx_s\).
The separate central remainders satisfy
\[
 \sum_{j=1}^3|\mathcal D_s^j\mathcal R_s|\le Cg\theta_s.
\]
Their values retain any original coefficient-only logarithmic constants;
no small zeroth remainder row is asserted.

On the product of these separate actual domains define
\[
 N_\pm=-g\mathcal D_\pm\log Z_\pm
       =1-gH_\pm-g\mathcal D_\pm\mathcal R_\pm,\qquad
 H_\pm=P_\pm'/(xB_\pm),\qquad n=N_+/N_-,\quad\varrho=n-1 .
\]
Then \(N_\pm\ge1/2\) and
\begin{equation}\label{deg:joint-inward-product-direct-jets}
 \sum_{i+j\le2}|\mathcal D_+^i\mathcal D_-^j\varrho|
                 \le Cg(\theta_++\theta_-).
\end{equation}
On the original actual subset \(Z_+(u)=Z_-(r)\), this \(n\) is exactly
the relative variation in
\[
 R_u=\frac{uB_+(u)\widehat P(R)}{RB_-(r)P_+(u)}\,n,
 \qquad R=kr,\quad \widehat P(R)=k^2P_-(R/k).
\]
The admissible central-image intersection may be empty. All coefficients
are frozen under the phase operators; constants are independent of the
ghost gap and of any comparison between \(g\) and \(k\).
\end{lemma}

\begin{proof}
Choose a fixed small \(0<b<M\), with \(\Pi b<1/4\), so
\(P_s\ge p_b>0\) on \([0,b]\) after the independent \(k\)-shrink.
On the positive strict ghost path and the negative joint interval,
\(P_s>0\); \(B_s,Q_s\) are positive units. Decreasing \(v\)
follows the positive passage or the inverse negative passage,
respectively, in~\eqref{deg:quadratic-inward-scalar-common}.
For \(v\ge b\), put \(t=v^{-1}\), \(V=kv\), \(\vartheta=t+V\), and
\[
 \begin{gathered}
 D_s=\frac{3-s\pi t+t^2+sV}{1+2g+sgV}
                  (-t\partial_t+V\partial_V),\\
 h_s=P_s/(v^2B_s),\quad H_0=\frac6{1+2g},\quad\ell_g=1-Ag+Cg^2.
 \end{gathered}
\]
These variables lie in a fixed compact rectangle and \(h_s\le C_b\),
including at the ghost bottleneck. The exact formulas
\[
 H_s-H_0=\frac{s[-\pi t+(3-gH_0)V]}{B_s},\qquad
 \omega_s=\frac{\zeta_sP_s}{vB_s}
                 =sCg^3V(1+sV)h_s
\]
give, without dividing by \(P_s\),
\begin{equation}\label{deg:joint-inward-weighted-coefficients}
 \begin{aligned}
 |D_s^j(H_s-H_0)|&\le C\vartheta &&(0\le j\le3),\\
 |D_s^j(Q_s-\ell_g)|&\le CgV,\quad
 |D_s^j\omega_s|\le Cg^3V &&(0\le j\le3),\\
 |D_s^j\vartheta|&\le C\vartheta &&(1\le j\le3).
 \end{aligned}
\end{equation}
Use the ratio equation and exact memory of
Lemma~\ref{deg:lem-stable-ratio-ledger}, with
\[
 \tau=(A_s(x_s)-A_s(v))/g,\quad
 T=(A_s(x_s)-A_s(b))/g,\quad W=Y/P_s(v),\quad W(0)=\ell_0,
\]
\[
 f=(Q_s-\omega_s)W/(1+Q_sW),\qquad
 I_b=\int_0^{T(x_s)}f\,d\tau .
\]
At fixed \(\tau\), \(\mathcal D_{x_s}v=P_s(v)/(vB_s(v))\), so
phase rows transport to powers of \(D_s\); clock rows use \(-gD_s\).
They commute, and \(|\vartheta_\tau|\le Cg\vartheta\) gives
\(\vartheta(v(\tau))\le\theta_s e^{Cg\tau}\).
Thus~\eqref{deg:joint-inward-weighted-coefficients} verifies all
phase and mixed coefficient inputs with
\(D=\mathcal D_{x_s}\), \(\kappa=g\), \(\rho=g\theta_s\),
\(\alpha=Cg<c_0/12\). The fixed small lip gives the two negative
margins. The lemma constructs the actual graph to \(b\) and bounds
its fixed-clock rows by \(Cg\theta_s(1+\tau)^Ne^{-c\tau}\).
Each strict fibre has finite action and flight time; their lengths
need no uniform upper bound.

The moving clock has \(\mathcal DT=1/g\), higher rows zero, and
\[
 T\ge c_b\log(x_s/b)/g\ge a_b/g,\qquad
 g^{-N}(1+T)^Ne^{-cT}\le Cg/x_s
\]
for every fixed \(N\), after shrinking \(g\). The first lower bound
absorbs \(x_s\); the second absorbs every fixed power of \(g^{-1}\).
Hence the full terminal formulas~\eqref{deg:stable-ratio-terminal-rows}
give \(\sum_{j=1}^3|\mathcal D^jI_b|\le Cg\theta_s\).
Their exact entrance memory is
\[
 \begin{gathered}
 y_b=\ell_0P_s(x_s)e^{-T+I_b}=P_s(b)W(T)\le Ce^{-cT},\\
 \Lambda_1=\mathcal D\log y_b=H_s(x_s)-1/g+\mathcal DI_b,\\
 \Lambda_j=\mathcal D^j\log y_b=D_s^{j-1}H_s+\mathcal D^jI_b
 \quad(j=2,3).
 \end{gathered}
\]
Thus \(|\Lambda_1|\le C/g\) and
\(|\Lambda_2|+|\Lambda_3|\le C\theta_s\).
On \([0,b]\), preparation~\eqref{deg:positive-large-S-inward-compact-graph}
checks the fixed positive margin, the \(C/g\) height bounds,
\(y_b\le h_b/2\), and small \(y_b/g\).
Lemma~\ref{deg:lem-compact-height-continuation} therefore gives the
actual centre height, with multiplier \(e^{-A_s(b)/g}\), and
\[
 \log Z_s=\log\{\ell_0P_s(x_s)\}-A_s(x_s)/g+I_b+\mathcal N_s(y_b).
\]
Its three composition rows are bounded by \(Cy_b/g^4\le Cg/x_s\).
This proves the weighted remainder estimate on the original centre
and admits its height to the original low-height window. Concatenation
at the same lips retains the original raw images.

\smallskip
\noindent\emph{Separated product and actual matching.}
For the product estimate let \(N_0=1-gH_0\) and \(d_s=N_s-N_0\).
The weighted coefficient and remainder rows give
\[
 \sum_{j=0}^2|\mathcal D_s^jd_s|\le Cg\theta_s,\qquad
 N_s\ge1/2,\qquad
 \varrho=\frac{d_+-d_-}{N_-}.
\]
The common homogeneous term \(N_0\) cancels in this quotient.
The quotient rule through order
two bounds all separated product rows; in particular the genuinely
mixed row is
\[
 \mathcal D_+\mathcal D_-n
 =-\frac{(\mathcal D_+N_+)(\mathcal D_-N_-)}{N_-^2}.
\]
This proves \eqref{deg:joint-inward-product-direct-jets} on the
product before imposing a central-image intersection.
Equality of the actual central heights
then gives
\[
 r_u=\frac{uB_+P_-(r)}{rB_-P_+(u)}\,\frac{N_+}{N_-},
\]
which is precisely the asserted relative variation after \(R=kr\).
Only the original central-image intersection is used; no actual image
or physical carrier is enlarged.
\end{proof}

\subsection{Physical resection at two simple quadratic roots}

Near \(b/a=2\), the endpoint value of the curvature principal part
can vanish at the same time as both action speeds. Fixed normal sections
replace the lip's endpoint powers by pure action residues. This yields
errors that vanish with each independent footpoint distance. The actual
resection below then transfers the resulting curvature orientation back
to the moving-lip description of the same physical map.

\begin{lemma}[Joint simple-root resection on the physical sections]
\label{deg:lem-joint-root-physical-resection}
Use the positive quadratic field and its original raw section buffers.
Restrict the real-root coefficients to a fixed neighborhood of
\(b/a=2\), so that \(b-a\) is a positive unit. All phase derivatives
freeze \(g,k,\pi,A,C\). Write
\[
 P(u)=(\mathfrak c+ku)(u-a)(u-b),\qquad
 \widehat P(R)=(\mathfrak c-R)(R+ka)(R+kb),
\]
\[
 B_+(u)=1+g(2+ku),\qquad B_-(R)=1+g(2-R),
\]
\[
 \mathfrak a_+=\frac{aB_+(a)}{(\mathfrak c+ka)(b-a)},\qquad
 \mathfrak a_-=
 \frac{\mathfrak cB_-(\mathfrak c)}
      {(\mathfrak c+ka)(\mathfrak c+kb)}.
\]
The residues and their reciprocals are bounded. Keep all
coefficient-only constants in the original actions
\(A_+=J_+\), \(A_-=A_-(R)\). Their exact pole subtractions are
\[
 A_+(a-h)=-\mathfrak a_+\log h+j_+(h),\qquad
 A_-(\mathfrak c-t)=-\mathfrak a_-\log t+j_-(t),
\]
where \(j_s-j_s(0)\) have uniform \(C^3\) bounds.

There are fixed normal-height sections in \(y_+=gY\) and
\(y_-=gk^2Y\), and fixed \(h_0,t_0,g_0,k_0>0\), with the
following properties for independently small \(g,k>0\).
The normal-section axis footpoints are \(u=a-h\) and
\(R=\mathfrak c-t\). Each normal section has its own actual outward
map to the original raw port and observer image.
On the actual intersection of the matched graph with
\(0<h<h_0\), \(0<t<t_0\),
\begin{equation}
 \begin{gathered}
 A_+(a-h)-A_-(\mathfrak c-t)=C_{g,k}+E_+(h)+E_-(t),\\
 |(h\partial_h)^jE_+|\le C\varepsilon_g h,\qquad
 |(t\partial_t)^jE_-|\le C\varepsilon_g t,\quad 0\le j\le3.
 \end{gathered}
 \label{deg:joint-root-normal-weighted-matching}
\end{equation}
where \(E_+(0)=E_-(0)=0\) and
\(\varepsilon_g=Cg(1+|\log g|)^4\).
No bound on \(C_{g,k}\) is required.
The quotient \(n\) of the two separated central action derivatives
has
\begin{equation}
 \begin{gathered}
 \sum_{i+j\le2}
 |\mathcal D_+^i\mathcal D_-^j(n-1)|
 \le C\varepsilon_g(h+t),\\
 \sum_{Z\in\{n,\mathcal D_+n,\mathcal D_-n\}}
 (|\partial_hZ|+|\partial_tZ|)\le C\varepsilon_g .
 \end{gathered}
 \label{deg:joint-root-normal-matching-rows}
\end{equation}
where
\[
 \mathcal D_+=\frac{P(u)}{uB_+(u)}\partial_u,\qquad
 \mathcal D_-=\frac{\widehat P(R)}{RB_-(R)}\partial_R.
\]
The normal-section embeddings use labels
\(\beta_+=-u^2/2\), \(\beta_-=-R^2/2\). After separate
coefficient-only positive affine value and slope calibrations,
their complete maps satisfy
\begin{equation}
 \begin{gathered}
 \phi_s=\mathrm{id}+O_{C^3}(\varepsilon_g),\\
 |\mathcal D_+^j\log\phi_+'|\le C\varepsilon_g h,\qquad
 |\mathcal D_-^j\log\phi_-'|\le C\varepsilon_g t,\quad j=1,2.
 \end{gathered}
 \label{deg:joint-root-normal-outer-rows}
\end{equation}
These statements use only the original ordinary \(C^3\)
logarithmic-slope reserve on the terminal observers, including the
divided finite closing on the negative side.

Let \(x=a-u_{\rm lip}\), \(y=\mathfrak c-R_{\rm lip}\) on the
original lips. The actual resection to the normal sections is
increasing and satisfies
\begin{equation}
 \begin{gathered}
 \mathfrak a_+\log h-(\mathfrak a_++g)\log x=C_++O(x+h),\\
 \mathfrak a_-\log t-(\mathfrak a_-+g)\log y=C_-+O(y+t),
 \qquad |C_s|\le C .
 \end{gathered}
 \label{deg:joint-root-exact-lip-resection}
\end{equation}
The remainders have the corresponding finite logarithmic derivative
bounds. The negative constant contains no unbounded pure \(\log k\).
If \(F_s\) denotes this resection in the internal labels, then
\begin{equation}
 S_{\rm fp}=F_-\circ S_{\rm lip}\circ F_+^{-1},\qquad
 \phi_{{\rm lip},s}=\phi_{{\rm fp},s}\circ F_s
 \label{deg:joint-root-physical-composition}
\end{equation}
on the actual common image, up to the retained positive affine
calibrations of the physical coordinates. Thus both constructions
describe the same completed physical map. No affine identification
of the two internal root-distance coordinates is asserted.
\end{lemma}

\begin{proof}
\emph{The normal sections and their exact invariant.}
On the positive side reverse time and use \(y=gY\); on the
negative side multiply time by \(k\) and use \(y=gk^2Y\).
Division by the positive axis normal rate gives, on both root boxes,
\begin{equation}
 v'=gF_s(v)+c_s(v)y,\qquad y'=y\{1+\kappa_s(v)y\},
 \label{deg:joint-root-local-normal-field}
\end{equation}
where
\[
 \begin{array}{lll}
 F_+=P/(vB_+),&
 c_+=Q_+/(vB_+),&
 \kappa_+=Ckg^2(1+kv)/(vB_+),\\
 F_-=\widehat P/(vB_-),&
 c_-=Q_-/(vB_-),&
 \kappa_-=-Cg^2(1-v)/(vB_-),
 \end{array}
\]
and \(Q_-=1-Ag(1-v)+Cg^2(1-v)^2\).
At \(r_+=a\), \(r_-=\mathfrak c\), the axis rates are exactly
\(gF_s'(r_s)=-g/\mathfrak a_s\).
On fixed boxes bounded away from zero, \(gF_s\) has \(C^4\) norm
\(O(g)\), \(c_s=1/v+O_{C^4}(g)\), and
\(\kappa_s=O_{C^4}(g^2)\), uniformly in \(k\).

Apply the backward-integral construction
\eqref{eq:saddle-normal-time}--\eqref{eq:exact-saddle-invariant}
to \(h=gF_s\). The coefficient bounds just established and
\(Y(-T)\le Cy e^{-T/2}\) give differentiated kernels
\(C(1+T)^Ne^{-(1/2-Cg)T}\) through order three.
This supplies \(\pi_s\) and \(Z_s=ye^{-\mathcal B_s}\), with
\(\pi_s'=gF_s(\pi_s)\), \((\log Z_s)'=1\).
The same estimates for the difference from
\(v'=y/v,\ y'=y\) give
\begin{equation}
 \begin{gathered}
 \pi_s(v,y)=\sqrt{v^2-2y}+O_{C^3}(g),\qquad
 \|\mathcal B_s\|_{C^3}\le Cg^2,\\
 \mathcal B_s=O(g^2y),\qquad \partial_v\pi_s\ge\tfrac12.
 \end{gathered}
 \label{deg:joint-root-local-projection}
\end{equation}
Consequently \(\log Z_s-A_s(\pi_s)/g\) is constant along each orbit.

Choose fixed small heights \(y_{0,s}>0\). On \(y=y_{0,s}\) write
\(\pi_s(N_s(U),y_{0,s})=U\). Then
\[
 N_s(U)=\sqrt{U^2+2y_{0,s}}+O_{C^3}(g),\qquad
 Z_{N,s}(h)=y_{0,s}e^{-\mathcal B_s(N_s(r_s-h),y_{0,s})},
\]
\[
 |\partial_h^j\log Z_{N,s}(h)|\le Cg^2,\qquad 1\le j\le3.
\]
These are the fixed normal sections and their actual axis footpoints.
The physical orbit from a sufficiently small lip point reaches its
normal section inside the unchanged box: the normal height grows
at a fixed rate, the projected axis footpoint remains on the inner
side of the root, and its physical displacement is \(O(y_{0,s})\).

\emph{The local amplitudes and the middle first multipliers.}
Fix inner cuts \(\bar v_s=r_s-\delta_s\) in the two boxes.
The invariant gives the exact implicit relation
\[
 \log y_{\rm in}+e_s(y_{\rm in})
 =\log Z_{N,s}(h)-\{A_s(r_s-h)-A_s(\bar v_s)\}/g,
\]
\[
 e_s(y)=-\mathcal B_s(\bar v_s,y)
 -\{A_s(\pi_s(\bar v_s,y))-A_s(\bar v_s)\}/g .
\]
Here \(e_s(0)=0\), and its required height derivatives have
polynomial \(g^{-1}\) bounds, since the inner cut is a fixed
distance from its root. Solving this scalar identity gives
\begin{equation}
 \begin{gathered}
 y_{{\rm in},s}(h)=\alpha_{y,s}h^{\mathfrak a_s/g}
 e^{-[j_s(h)-j_s(0)]/g}u_s(h),\\
 u_s(h)=\frac{Z_{N,s}(h)}{Z_{N,s}(0)}
 e^{-e_s(y_{{\rm in},s}(h))}.
 \end{gathered}
 \label{deg:joint-root-local-pure-residue}
\end{equation}
with
\[
 \alpha_{y,s}=Z_{N,s}(0)
 e^{[A_s(\bar v_s)-j_s(0)]/g},\qquad
 \alpha_{y,s}^{\pm1}\le e^{K/g}.
\]
The coefficient-only constants cancel in
\(A_s(\bar v_s)-j_s(0)\); thus this local bound permits a
global \(\log k\) term in the negative action. Passing to \(Y\)
on the positive cut and \(\widehat Y=k^2Y\) on the negative cut
adds only the coefficient-only factor \(g^{-1}\).

Choose preliminary \(h_0,t_0\) so the positive residue powers dominate these
fixed exponential constants. Three ordinary endpoint derivatives
introduce only polynomial \(g^{-1}\) factors and inverse endpoint
powers. Their nonlinear terms are then bounded by arbitrarily high
fixed endpoint powers after shrinking \(g\). The remaining phase
part of \(\log u_s\) is \(O_{C^3}(g^2)\).
Consequently the first three logarithmic endpoint rows of
\(g\{\log u_s(h)-\log u_s(0)\}\) are bounded by
\(C\varepsilon_gh\).

For the negative growing bridge, take \(R_0=\bar v_-\), the
already fixed inner cut with a margin below the root and away from
zero, and decrease \(k_0\)
so that \(kr_c<R_0\). Retain its exact first multiplier
\[
 \widehat m=\exp\left\{-\frac1g\int_{kr_c}^{R_0}
       \frac{RB_-}{\widehat P}\,dR\right\},\qquad
 \widehat m/k^2\le e^{K/g}k^{c_0/g-2}\le e^{K/g}.
\]
Indeed \(\widehat P\asymp R^2\) on \(Kk\le R\le R_0\), giving
the action lower bound \(c_0\log(1/k)-K\). Choose \(g_0<c_0/4\).
No reciprocal bound for this collapsing multiplier is used.
Put \(W=\widehat Y/\widehat P\) and
\(q=(A_-(R_0)-A_-(R))/g\). Coefficients and both radial cuts are
fixed under the incoming height variation; the flight
\(T=(A_-(R_0)-A_-(kr_c))/g\) may be arbitrarily long.
On \(kr_c\le R\le R_0\), the fixed root margin gives
\(\widehat P\asymp(R+k)^2\) and \(k/R\le1/r_c\). Hence
\[
 Q_--1=O(g),\qquad
 H=\widehat P_R/(RB_-)=O(1),\qquad
 \omega=-Cg^3(1-R)\widehat P/(RB_-)=O(g^3).
\]
The exact ratio equation has \(\kappa=g\). On the unchanged fixed
small buffer \(0\le W\le\ell_0\), these bounds keep
\(1+Q_-W\) a positive unit and both
\(\mathcal F/W,\mathcal F_W\le-c_1<0\), after decreasing \(g_0\).
At the fixed initial cut \(\widehat P(R_0)\) is a positive unit.
The residue formula \eqref{deg:joint-root-local-pure-residue},
with its factor \(g^{-1}\), places
\(w(t)=\widehat Y_0(t)/\widehat P(R_0)\le\ell_0/2\)
after the same independent coefficient shrink. Thus
Lemma~\ref{deg:lem-fixed-route-height-rows} applies on this full buffer.
Retaining the terminal polynomial in its exact cocycle gives
\[
 \widehat Y_T=\widehat Y_0\widehat m
             e^{\mathcal N_T(w)},\qquad
 |\partial_w^j\mathcal N_T|\le C\quad(1\le j\le3),\qquad
 |(w\partial_w)^j\mathcal N_T|\le Cw\quad(0\le j\le3).
\]
The normalized nonlinear unit therefore has bounded ordinary \(C^3\)
rows in the incoming height at the fixed cut, uniformly in \(g,k,T\).
Its three endpoint rows follow by the Euler composition formula of
Lemma~\ref{deg:lem-compact-height-continuation}, with
\(\Lambda_j=(t\partial_t)^j\log w\): the flat \(w\) absorbs the
polynomial \(g^{-1}\) and endpoint powers from the local residue.
These are fixed-route height rows; phase and moving-terminal rows on
other routes remain those of Lemma~\ref{deg:lem-stable-ratio-ledger}.
Conversion to \(Y\) contributes \(k^{-2}\) only to
\(\widehat m/k^2\), with no reciprocal gain estimate.

Fix the prescribed height buffers of the final compact inward
continuations and their coefficient bounds before the final choice
of \(h_0,t_0\). Enlarge \(K_{\rm tot}\) to include the regular
local action parts and the preceding bound for \(\widehat m/k^2\).
For any fixed endpoint power \(n\) and finite polynomial loss
\(g^{-M}\), choose \(t_0<1\) so that
\(\mathfrak a_{\min}\log(1/t_0)>K_{\rm tot}+2c_*\), with
\(c_*>0\). Then decrease \(g_0\); the residue and its three
ordinary endpoint derivatives give
\[
 g^{-M}e^{K_{\rm tot}/g}t^{\mathfrak a_{\min}/g-j}
 \le Ce^{-c_*/g}t^n,\qquad 0<t\le t_0,\quad 0\le j\le3.
\]
In particular the converted input satisfies \(Y_T\le h_b/2\)
and small \(Y_T/g\) in its already prescribed buffer \(h_b\),
uniformly in independently small \(k\). The positive input has
the same bound after the analogous final choice of \(h_0\).
Lemma~\ref{deg:lem-compact-height-continuation} now applies to the
actual final inward segments: their linear propagators contract,
and their nonlinear height bounds are polynomial in \(g^{-1}\).
Its three Euler composition rows and the displayed superpower
bound control the endpoint rows on those original centre windows.
No reciprocal linear multiplier is introduced in this admission.
Equality of the two exact original central heights therefore gives
\eqref{deg:joint-root-normal-weighted-matching}; all coefficient-only
first multipliers remain in \(C_{g,k}\).
Put \(d_+=h\), \(d_-=t\). On the fixed root boxes,
\(\mathcal D_s=-d_s\gamma_s(d_s)\partial_{d_s}\), where
\(\gamma_s\) are positive units with bounded \(C^2\) norms.
The corrected actions \(A_+-E_+\) and \(A_-+E_-\) in
\eqref{deg:joint-root-normal-weighted-matching} have action derivatives
\[
 \mathcal N_+=1-\mathcal D_+E_+,
 \qquad \mathcal N_-=1+\mathcal D_-E_-,
 \qquad n=\mathcal N_+/\mathcal N_- .
\]
The three weighted rows of \(E_s\) give
\[
 |\mathcal D_s^j(\mathcal N_s-1)|
     \le C\varepsilon_gd_s\quad(0\le j\le2),\qquad
 |\partial_{d_s}\mathcal N_s|
  +|\partial_{d_s}(\mathcal D_s\mathcal N_s)|
     \le C\varepsilon_g .
\]
For the ordinary derivatives, divide the second and third action
rows of \(E_s\) by the single factor \(d_s\gamma_s\).
After shrinking the fixed boxes, \(\mathcal N_s\ge1/2\).
The separated quotient calculation in the proof of
Lemma~\ref{deg:lem-joint-inward-product-direct} now applies with
errors \(\varepsilon_gh\) and \(\varepsilon_gt\); its mixed row is
\(O(\varepsilon_g^2ht)\). Differentiating the quotient and its first
action rows with the displayed ordinary bounds proves
\eqref{deg:joint-root-normal-matching-rows} on the product, before
restriction to the actual match and without comparing \(h\) and \(t\).

\emph{The actual fixed-normal exterior embeddings.}
Put \(\sigma_+=gk\), \(\sigma_-=g\). At a normal section,
\[
 x_l=\sigma_sN_s(U),\qquad z_l=\sigma_s^2y_{0,s}.
\]
The ratio \(W_l=y_{0,s}/N_s(U)^2\) is a positive unit, even
when the normal chord lies beyond its axis root.
Use the signed coefficients of~\eqref{deg:quadratic-outer-ratio},
writing \(N_s^{\rm raw}\) for its numerator and \(D_s=m_s+Q_sW\).
As in the bounded positive outward construction,
\((D_s)_\eta=Q_sW>0\) at \(D_s=0\), so the inverse-height
graph retains its positive denominator.
Choose a fixed auxiliary \(d_*<d_q\) small compared with the lower
bound for \(W_l\). On \([x_l,d_*]\), the singular terms in
\(m_s-sx\) are \(O(g)\). The same invariant-strip argument, starting
now at \(x_l\) with no weak stage, gives
\[
 w_*/2\le W\le w^*,\qquad D_s\ge cW,\qquad
 V_r/V=-2+O((g+x)/w_*)\le-3/2,
\]
\[
 |W(d_*)-W_s^*(d_*)|\le Cg+C(x_l/d_*)^c .
\]
The fixed parabolic continuation therefore admits the graph to its
original raw buffer and observer.

Here \(\partial_U\eta_l=\partial_U^2\eta_l=0\), and the inverse
initial rows are bounded with positive first row. Applying
\eqref{deg:scalar-inverse-endpoints} as above gives
\[
 0<P_1\le C(x_l/x)^{3/2},\quad |P_2|\le CP_1,\qquad
 |\partial_U\eta_T|+|\partial_U^2\eta_T|\le C\sigma_s^{3/2}.
\]
The exact first-variation cocycle and density split, with
\(\nu=6g/(1+2g)\), \(\lambda=2\nu\), give
\[
 J_z=\frac{W_T/D_T}{W_l/D_l}(z_T/z_l)^\nu e^{I_s}.
\]
The density's first two fixed-height radial derivatives are bounded
by \(C(x+kg^2/x)\). Thus~\eqref{deg:scalar-density-endpoints},
with \(D=\partial_U\) and both lower clock rows zero, bounds its two
complete rows by
\[
 C(\sigma_s+kg^2/\sigma_s)+C\sigma_s^{3/2}\le Cg.
\]

Since \(z_l\) is phase-constant and \(\beta=-U^2/2\), the full
moving first variation is
\[
 \Gamma_s'(\beta)=
 \sigma_s^2\frac{N_s(U)N_s'(U)}U
 \frac{W_lN_l^{\rm raw}}{D_l}J_z\Phi'(z_T).
\]
Removing only the coefficient-only observer gain gives the exact
form
\begin{equation}
 \Gamma_s'(\beta)=
 c_s^\ast\sigma_s^{2-\lambda}y_{0,s}^{-\nu}
 \frac{N_s(U)N_s'(U)}U N_l^{\rm raw}
 \frac{W_T}{D_T}z_T^\nu\widehat\Phi'(z_T)e^{I_s}.
 \label{deg:joint-root-normal-exact-physical-slope}
\end{equation}
The initial factor equals \(1+O_{C^2}(g)\), because
\eqref{deg:joint-root-local-projection} gives
\(N_s(U)N_s'(U)/U=1+O_{C^2}(g)\).
The terminal factors have two logarithmic rows
\(O(\sigma_s^{3/2})\), using just the original ordinary \(C^3\)
logarithmic-slope observer bound. Thus both ordinary logarithmic
rows of the slope are \(O(g)\).
Subtracting its exact value and calibrating its exact positive
slope at a fixed reference point gives
\(\phi_s=\mathrm{id}+O_{C^3}(g)\).
The action fields contain the factors \(h,t\), proving
\eqref{deg:joint-root-normal-outer-rows}.
The two physical gains
\(\sigma_+^{2-\lambda}\), \(\sigma_-^{2-\lambda}\) are retained
coefficient-only factors in
\eqref{deg:joint-root-normal-exact-physical-slope}; their ratio
and their reciprocals are not estimated.

\emph{Resection of the lip and the same physical curvature.}
On either lip the local normal coordinate is
\[
 y_{{\rm lip},+}=g\ell_0P(a-x),\qquad
 y_{{\rm lip},-}=g\ell_0\widehat P(\mathfrak c-y).
\]
It is \(g\) times a positive unit times the root distance.
The projected lip distance is
\(\widehat x=x(1+O(g))\), respectively
\(\widehat y=y(1+O(g))\).
Applying the exact invariant between the lip and its normal-section
intersection, and subtracting the two pole expansions, gives
\eqref{deg:joint-root-exact-lip-resection}.
The coefficient terms are \(g\log g\), logarithms of local
positive units, and the bounded projection correction
\(\mathfrak a_s\log(\widehat x/x)\).
The negative comparison is in \(gk^2Y\), so it has no pure
\(\log k\). The common growing-bridge multiplier cancels from
this local comparison.
The finite logarithmic derivative estimates follow by differentiating
this invariant on the positive interior of the lip.
In particular
\[
 h=A_+x^{1+g/\mathfrak a_+}e^{O(x+h)},\qquad
 t=A_-y^{1+g/\mathfrak a_-}e^{O(y+t)},\qquad A_s^{\pm1}\le C .
\]
The increasing maps give uniform inclusions between the nested
endpoint neighborhoods;
the lip exponents are \(\mathfrak a_s/g+1\), whereas the normal
residue quotient has no added \(g\).

Both resection coordinates belong to the same orbit, so
\eqref{deg:joint-root-physical-composition} identifies their completed
maps after the retained positive affine physical calibrations.
The normal physical input is increasing in \(h\), since
\(d\{-(a-h)^2/2\}/dh=a-h>0\); the lip physical input is
increasing in \(x\), and the resection is increasing.
Consequently the same completed physical map has the same curvature
zeros and crossing directions in the two descriptions.
\end{proof}

\subsection{Paired first variation on positive quadratic fibres}

Above a comparison orbit the complete derivative is bounded through
the ratio of two variations. Their separate core factors need not be
units, and both exterior gains contain a small common power of \(gk\).
The next lemma retains these factors until their cancellation in the
physical derivative.

\begin{lemma}[Paired first variation above a quadratic comparison]
\label{deg:lem-positive-quadratic-paired-variation}
Fix the original raw sections, nested buffers, ordinary \(C^3\)
regular observers and the divided finite closing of
Lemma~\ref{lem:closing}.
Use the polynomial and physical coordinates of
\eqref{deg:quadratic-field}, and write \(a=d-\sqrt\Delta\) when
\(\Delta\ge0\). The passive compact
\((A,C)\)-set lies inside their prescribed common coefficient
neighbourhood, with the original buffer and observer margins.
Fix a finite coefficient bound \(\Pi\). In the original domain
construction choose \(X\) and the retained ghost width jointly for
the prescribed finite comparison action threshold, and then fix
\(\pi_\flat\). For \(\pi\in[\pi_\flat,\Pi]\), there are fixed
\(\eta_{\max},C_*>0\) and independent upper bounds for \(g,k\)
with the following property on the retained real, fold and ghost fibres.
Use coefficient fibres with the actual \(K_{\rm fin}>0\) and the
original \(0<m\le m_0\) finite-closing admission.
Choose the comparison lip at
\[
 u_* = a-X/2\quad(\Delta\ge0),\qquad
 u_* = d-X/2\quad(\Delta<0),
 \qquad \eta_{\rm comp}=gZ_+(u_*).
\]
The comparison lies strictly inside the enlarged low lip interval
beginning at \(a-X\), respectively \(d-X\). The actual completed
transition in the original affine input satisfies
\[
 0<\Pi'\le C_*
 \qquad\text{throughout}\qquad
 \eta_{\rm comp}\le\eta_0=gY(0)\le\eta_{\max}.
\]
This entire interval is admitted to the original section buffers.
The constants are independent of the root gap and of any relation
between \(g\) and \(k\). After the common choice \(mC_*<1/2\),
\(F'=1-m\Pi'>1/2\) on it.
\end{lemma}

\begin{proof}
Write
\[
 \begin{gathered}
 \eta=gY,\qquad \epsilon_0=gk,\qquad D=gP+Q\eta,\\
 Q=1-Ag(1+k\xi)+Cg^2(1+k\xi)^2,
 \qquad B=1+g(2+k\xi).
 \end{gathered}
\]
The exact field is
\[
 \dot\xi=-D,\qquad
 \dot\eta=-\eta\{\xi B+H_1\eta\},
 \qquad H_1=Ckg^2(1+k\xi),
 \qquad z=\epsilon_0^2\eta.
\]
By \eqref{ann:source-field}, this is the source field with coefficient quartet
\((-g,g\pi,g,\epsilon_0)\); \(\epsilon_0\) is held fixed under
every phase derivative.

\smallskip
\noindent\emph{The actual comparison and the common cuts.}
Use the jointly chosen fixed \(X>0\). The abscissa \(u_*\) ranges over a
fixed positive compact interval, and there is a fixed \(\epsilon>0\)
such that
\[
 P\ge p_*>0\qquad\text{on }[0,u_*+\epsilon].
\]
The factorization gives this margin. The comparison action choice
of Proposition~\ref{deg:prop-positive-quadratic} first fixes \(X\)
by \(J_+(a-X)\ge c\log(\delta_c/X)\), then the ghost width;
the fold has the stronger inverse-distance divergence.

Corollaries~\ref{deg:cor-positive-compact-suppliers} and
\ref{deg:cor-negative-inward-common-centre} give the original inward
cocycles, without changing the centre,
\[
 \log Z_+=\log(\ell_0P)-J_+/g+\mathcal R_+,
 \qquad
 \log Z_-=\log(\ell_0P_-)-A_-/g+\mathcal R_-.
\]
Their zeroth rows on the bounded comparison and the entire strict
negative side are bounded. Indeed their ratio equations give
\(W\le\ell_0e^{-c\tau}\), hence
\(\left|\int (Q-\omega)W/(1+QW)\,d\tau\right|\le C\ell_0\),
and the original compact continuation adds \(O(Y_b/g)\).
The original action constants remain, \(N_-\ge1/2\), and
\(Z_-\downarrow0\) at its simple root. Choose fixed \(r_{\max}>r_c\)
with \(A_-(r_{\max})\) above the bounded comparison action, then shrink
\(k\). Monotonicity and the zeroth cocycles give the unique actual match
\[
 r_*\in[r_c,r_{\max}],\qquad Z_+(u_*)=Z_-(r_*),
 \qquad \eta_{\rm comp}=gZ_+(u_*).
\]

On either comparison lip, \(z/x^2=\ell_0gp\), with \(p\) a
positive unit: on the negative lip \(r_*\) is bounded. Use the
original inverse logarithmic-height graph until it reaches a fixed
small strip height \(w_0\). Its bounded-comparison weak length is
\(O(1+|\log g|)\). The positive weak displacement estimate is
\(O(g(1+|\log g|))+O(w_0)\) in \(\log x\).
Choose \(w_0\) first so small that this strip is entered before
\(\xi=u_*+\epsilon/2\), and then decrease \(g\).
This uses a fixed positive margin of \(P\), and is uniform through
the fold. The negative weak stage reaches its fixed strip by
\(|\xi|\le L_1r_*\), with \(L_1\) fixed.

The outward constructions of
Corollaries~\ref{deg:cor-positive-compact-suppliers} and
\ref{deg:cor-negative-absolute-outer} continue both comparison
traces. Choose one fixed \(L\) larger than those entrance bounds
and the bounded positive factors. On the two coefficient-dependent
but phase-fixed cuts \(\xi=\pm L\), the comparison has
\[
 c_L\le\eta(\pm L)\le C_L,
 \qquad D(\pm L)\asymp\eta(\pm L).
\]
Increasing \(L\), the strip contraction places both values of
\(\eta/L^2\) inside the same small parabolic band. The original
outward continuation, including its weak-stage initial factors,
lands strictly inside the original raw buffers. No raw section or
observer window has been replaced.

\smallskip
\noindent\emph{One upper reference and all intervening inputs.}
Fix \(\eta_{\max}>0\). At \(g=k=0\) the core graph with centre
height \(\eta_{\max}\) is
\(\eta=\eta_{\max}+\xi^2/2\), with strictly positive height and
speed on every fixed core. Ordinary dependence on \([-L,L]\)
constructs its actual upper reference for small independent \(g,k\).
Choose \(L\) large enough that its two values of \(\eta/L^2\)
belong to the same parabolic band as the comparison.

The inverse-height first-exit calculation continues this reference
to an auxiliary raw cut \(d_*<d_q\), where strip contraction gives
\[
 |W(d_*)-W^*(d_*)|\le Cg+C(\epsilon_0L/d_*)^c\longrightarrow0.
\]
The regular segment to the original port has the comparison's limit
\(W^*(d_q)\); the prescribed passive and raw-buffer margins admit it.

Scalar uniqueness orders the graphs between the comparison and
this upper reference. They stay between the two finite positive
core graphs and between their admitted outward tails. Thus every
\(\eta_0\in[\eta_{\rm comp},\eta_{\max}]\) has an actual through
graph on the entire word and remains in the original convex raw
and observer windows. Its central raw height is
\(z(\gamma)=\epsilon_0^2\eta_0\); after the common coefficient
shrink the original low-height and speed restrictions are satisfied.
The whole interval, including nonroot inputs, is consequently used.

\smallskip
\noindent\emph{The paired core variation.}
For
\[
 f=\frac{d\eta}{d\xi}
   =\frac{\eta(\xi B+H_1\eta)}D
\]
use \eqref{deg:scalar-raw-cocycle}, with residual
\([gP'+(Q'+H_1)\eta]/D\).
On \([-L,u_*+\epsilon]\), \(P\) has a fixed positive lower
bound after decreasing \(k\); on the remaining positive core,
the comparison has already entered its invariant strip. Hence
\[
 D_{\rm comp}\ge c g\quad\text{on }[-L,L],\qquad
 \eta/D\le C
\]
on every graph above it. In the part where \(P\) could be negative,
\(\eta\ge c>0\), so \(D\ge c\eta\) after decreasing \(g\).
Since \(P'\) is bounded and
\(Q'=O(gk)\), \(H_1=O(kg^2)\), the remaining integral in the
identity is bounded in absolute value on the full core. The endpoint
quotients \((\eta/D)(\pm L)\) are positive units.

Set
\[
 A_\pm=\frac{\partial\eta(\pm L;\eta_0)}{\partial\eta_0}>0.
\]
Integration gives
\begin{equation}
 \frac{A_-}{A_+}
 =\frac{(\eta/D)(-L)}{(\eta/D)(L)}
   \exp\left\{\int_L^{-L}
          \frac{gP'+(Q'+H_1)\eta}D\,d\xi\right\}
 \in[c,C].
 \label{deg:quadratic-paired-core-ratio}
\end{equation}
This estimate controls the ratio of the two core variations;
separate unit bounds for \(A_\pm\) are not required.
For comparison with a moving lip phase \(u\), the exact incoming
factor is
\[
 \frac{d\eta_0}{du}
  =\eta_0\left\{\frac{P'}P-\frac{uB_+}{gP}
                     +\mathcal R_+'\right\}
  =-\frac{uB_+\eta_0}{gP}N_+.
\]
The core-end derivatives are \(A_\pm\,d\eta_0/du\), so this
complete moving-source factor cancels. The common cuts are phase-fixed.

\smallskip
\noindent\emph{Both original tails and the physical derivative.}
In the source coordinates \((\xi,\eta)\), with quartet
\((-g,g\pi,g,\epsilon_0)\), the displayed parabolic rational
calculation \eqref{ann:tail-first-derivative} gives, on either
admitted tail,
\[
 F_{\sigma,\eta}-\lambda_g/r
   =O(\epsilon_0+g\pi r^{-2}+r^{-3+\nu}),
 \qquad \lambda_g=\frac{12g}{1+2g},\qquad 0<\nu<1.
\]
Its remainder has bounded integral from \(L\) to
\(d_\sigma/\epsilon_0\). The parabolic bands just constructed
justify this application even though the centre height may be
exponentially small. Consequently the
two actual tail first variations are
\[
 J_\sigma=(d_\sigma/(\epsilon_0L))^{\lambda_g}U_\sigma,
 \qquad U_\sigma,U_\sigma^{-1}\le C.
\]
At the original raw ports the complete core-plus-tail variations are
exactly
\begin{equation}
 \frac{dz_\sigma}{d\eta_0}
   =\epsilon_0^{2-\lambda_g}(d_\sigma/L)^{\lambda_g}
       A_\sigma U_\sigma.
 \label{deg:quadratic-paired-raw-slopes}
\end{equation}
All terminal factors lie on the admitted original images.
The original bridges and divided closing, with \(K_{\rm fin}>0\),
\(m\le m_0\), have positive derivative units by Lemma~\ref{lem:closing}.

The common \(\epsilon_0^{2-\lambda_g}=(gk)^{2-\lambda_g}\)
factor cancels in the ratio of the two raw variations. The ratio
\((d_-/d_+)^{\lambda_g}\), the two \(U\)-factors, and the
paired core ratio are bounded units. Composing the original regular
bridges and divided closing therefore proves \(0<\Pi'\le C_*\)
in the same original affine input on the entire stated interval.
Choose \(mC_*<1/2\) only after these constants and buffers have
been fixed. This gives the stated strict sign of \(F'\).
\end{proof}

\section{Coefficient dictionaries and passage domains}
\label{app:coefficient-dictionaries}

We express each coefficient limit in a chart where the corresponding
passage estimate is uniform. The substitutions are exact, so the estimates
apply to the original coefficient fibres. We then choose overlapping
chart ranges in the order required by Section~\ref{sec:coefficient-limits}.
All constants precede the common raw cap and final coefficient
restrictions in Lemma~\ref{lem:common-raw-receiver}.

\subsection{Proof of Lemma~\ref{deg:lem-small-source-routing}}
\label{app:proof-small-source-routing}

\begin{proof}
If \(p\le0\), use Lemma~\ref{deg:lem-raw-jacobian}. For \(p>0\),
first fix the positive-quadratic ghost and root-free reserves, and then
the independent far-root rectangles.

Fix a small number \(\mathfrak q_{\rm far}>0\) for which
Proposition~\ref{deg:prop-positive-far} applies.  On the negative
\(\gamma\)-side, Proposition~\ref{deg:prop-degree-infinity} and the
dictionary below provide numbers \(\mathfrak q_- >0\) and
\(\Pi_->1\).  Put
\begin{equation}
 \mathfrak q_\flat=\frac12\min(\mathfrak q_{\rm far},\mathfrak q_-),\qquad
 \Pi_{\rm far}=\max\{\lambda_0^{-1/2},\Pi_-,2\sqrt{12}\},
 \qquad \Pi_{\max}=2\Pi_{\rm far},
 \label{deg:small-constant-1}
\end{equation}
where \(\lambda_0\) is the constant in
\eqref{deg:positive-far-range}. Apply the whole-carrier estimate of
Corollary~\ref{deg:cor-positive-quadratic-whole-carrier} and
Proposition~\ref{deg:prop-negative-gamma} on
\(0<\pi\le\Pi_{\max}\).

For \(\mathfrak q=k\pi\ge \mathfrak q_\flat\), the degree-loss dictionary given below is
uniform once \(k\le k_d(\mathfrak q_\flat)\).  Choose
\(k_\flat>0\) so that \(2k_\flat\) is smaller than this number and than
the \(k\)-bounds in both compact-\(\pi\) carriers, including the
whole positive-carrier bound \(k_0\le k_{\rm tail}\).  Put
\begin{equation}
 G_{\max}=k_\flat^{-2/3}.
 \label{deg:small-Gmax}
\end{equation}
For this fixed \(G_{\max}\), choose \(U_d\) so that the large-\(U\)
dictionary below belongs to Proposition~\ref{deg:prop-degree-drop}, and
then choose \(U_{\rm cub}>U_d+1\).  Apply Proposition~\ref{deg:prop-compact-cubic} on
\(0\le U\le U_{\rm cub}\), \(|G|\le G_{\max}\). The bounded
positive-quadratic \(k\to0\) face remains in the first compactification;
it is distinct from this compact cubic rectangle.

We first check that the two degree-loss substitutions reproduce the
source coefficients exactly.  In the \(\gamma\)-dominant
chart \eqref{deg:gamma-chart}, let
\(\epsilon=\operatorname {sgn}\gamma\) and put \(t=k\xi\).
The exact polynomial is
\begin{equation}
 F_{\epsilon,\mathfrak q,k}(t)=t^3+3\epsilon t^2-\mathfrak q t+k^2.
 \label{deg:gamma-degree-polynomial}
\end{equation}
For \(\mathfrak q\ge \mathfrak q_\flat\) and \(k\) small it has roots
\(-r<t_a<t_c\), all three simple.  Vieta's formulas are
\[
 t_a+t_c=r-3\epsilon,\qquad
 rt_at_c=k^2,\qquad
 \mathfrak q=r(r-3\epsilon)-\frac{k^2}{r}.
\]
Define
\begin{equation}
 \begin{gathered}
 \tau=p^2,\qquad
 \kappa=\frac{k^2}{\mathfrak q^{3/2}},\qquad
 \mathcal R=\frac{r}{\sqrt{\mathfrak q}},\qquad
 B=\frac{\mathfrak q}{r^2},\qquad \omega=\frac{\sqrt{\mathfrak q}}{r},\\
 \nu=\frac{\epsilon}{r},\qquad
 a=\frac{t_a}{r},\qquad c=\frac{t_c}{r},\qquad
 A_d=\frac{B^{3/2}}{c}.
 \end{gathered}
 \label{deg:gamma-degree-dictionary}
\end{equation}
These identities imply \eqref{deg:degree-graph} exactly.  Moreover
\[
 (\Lambda_1,\Lambda_2,\Lambda_3)
 =(1,\kappa^2,\epsilon/\pi^2),
\]
so
\[
 (\sqrt\tau\Lambda_1,\tau\Lambda_3,\tau^{3/2}\Lambda_2)
 =(p,\gamma,s).
\]
For \(\mathfrak q\in[\mathfrak q_\flat,\infty]\), all positive-unit quantities in
Proposition~\ref{deg:prop-degree-drop} have compact closure:
\(r/\sqrt \mathfrak q\to1\) and
\((B,\omega,c,\nu,A_d)\to(1,1,1,0,1)\) as \(\mathfrak q\to\infty\).

On the negative-\(\gamma\) far face \(\mathfrak q\downarrow0\),
\(\pi\to\infty\), use the roots of
\[
 F_{-,\mathfrak q,k}(t)=t^3-3t^2-\mathfrak q t+k^2.
\]
The exact regular coordinates of Proposition
\ref{deg:prop-degree-infinity} are
\begin{equation}
 \begin{gathered}
 a=\frac{t_a}{r},\qquad j=\frac{r}{t_c},\qquad
 \kappa=a,\qquad \mathcal R=j,\qquad d=B=A_d=1,\\
 \nu=-\frac1r,\qquad \omega=1-a+ja,\qquad
 \tau=\frac{gt_c}{a}.
 \end{gathered}
 \label{deg:negative-far-dictionary}
\end{equation}
Vieta's formulas give
\[
 t_a+t_c=r+3,\qquad rt_at_c=k^2,\qquad
 \mathfrak q=r(r+3)-\frac{k^2}{r},\qquad
 \pi^2=\frac{t_c}{a}(1-a+ja)^2,\qquad
 \tau=\frac{p^2}{\omega^2}.
\]
Hence
\[
 r\sim \mathfrak q/3,\quad t_c\to3,\quad a\sim3/\pi^2,\quad
 j\sim \mathfrak q/9,\quad \tau\sim p^2.
\]
Thus \((j,\kappa,\tau)\to0\), the positive-unit entries stay in
compact ranges, and \(j\nu=-1/t_c=(j-ja-1)/3\) stays bounded.
The same source coefficients are recovered exactly. Fix
\(\mathfrak q_-\), then \(\Pi_-\), so that
\(\omega\ge\omega_{\min}>0\), and let \(\tau_-\) be the
reciprocal radial bound. The rectangle
\begin{equation}
 0<\mathfrak q\le \mathfrak q_-,\qquad \pi\ge\Pi_-,\qquad
 p^2\le\tau_-\omega_{\min}^2
 \label{deg:negative-far-rectangle}
\end{equation}
lies within the range of Proposition~\ref{deg:prop-degree-infinity}.
The estimate therefore holds throughout this rectangle, independently
of the relative rates of its parameters.

For the complementary cubic chart \eqref{deg:cubic-chart}, the large
\(U\) polynomial
\[
 y^3+3Gy^2-Uy+1
\]
has three simple roots \(-r<t_a<t_c\) when
\(U\ge U_d(G_{\max})\).  Set
\begin{equation}
 \begin{gathered}
 \tau=p^2,\qquad \kappa=U^{-3/2},\qquad
 \mathcal R=\frac{r}{\sqrt U},\qquad
 B=\frac{U}{r^2},\qquad \omega=\frac{\sqrt U}{r},\\
 \nu=\frac{G}{r},\qquad a=\frac{t_a}{r},\qquad
 c=\frac{t_c}{r},\qquad A_d=\frac{B^{3/2}}{c}.
 \end{gathered}
 \label{deg:cubic-degree-dictionary}
\end{equation}
Again Vieta's formulas imply \eqref{deg:degree-graph}, and
\[
 (\Lambda_1,\Lambda_2,\Lambda_3)=(1,U^{-3},GU^{-2})
\]
recovers \((p,\gamma,s)=(tU,t^2G,t^3)\) exactly.  Uniformly for
\(|G|\le G_{\max}\), all positive units in
\eqref{deg:cubic-degree-dictionary} tend to one and \(\nu\to0\).

The compact-degree positive units have compact closure. Let
\(\tau_d>0\) be their radial bound, and let the positive far-root
bound \(\tau_0\) include its fixed core/high splice reserve. Choose once
\begin{equation}\label{deg:small-ratio-choice}
 0<p_*^2<\tfrac12\min\{\tau_0,\tau_d,
                                  \tau_-\omega_{\min}^2\}.
\end{equation}
Then the degree and positive far-root scales \(\tau=p^2\), and the
negative far-root scale \(\tau=p^2/\omega^2\), lie in their required
ranges. These radial restrictions must be chosen explicitly, since shrinking
the raw coefficients alone does not control the chart ratios.

The chosen chart ranges now cover every coefficient alternative. If \(\gamma\ne0\) and
\(0<k\le2k_\flat\), then:
\begin{itemize}
 \item for \(0<\pi\le\Pi_{\max}\), use
 Propositions~\ref{deg:prop-positive-quadratic} and
 \ref{deg:prop-negative-gamma}, according to the sign of \(\gamma\);
 \item for \(\pi>\Pi_{\rm far}\) and \(\mathfrak q\ge \mathfrak q_\flat\), use
 Proposition~\ref{deg:prop-degree-drop} through
 \eqref{deg:gamma-degree-dictionary};
 \item for \(\pi>\Pi_{\rm far}\), \(\mathfrak q<\mathfrak q_\flat\), and \(\gamma>0\),
 use Proposition~\ref{deg:prop-positive-far};
 \item in the same range with \(\gamma<0\), use
 Proposition~\ref{deg:prop-degree-infinity} through
 \eqref{deg:negative-far-dictionary}.
\end{itemize}
The compact and far \(\pi\)-ranges, and the two \(\mathfrak q\)-ranges, overlap
strictly.  If \(\gamma=0\), or if \(k\ge k_\flat\), then
\[
 |G|^{3/2}=\frac{|\gamma|^{3/2}}{s}\le k_\flat^{-1}.
\]
For \(U\le U_{\rm cub}\), use
Proposition~\ref{deg:prop-compact-cubic}; for \(U\ge U_d\), use
Proposition~\ref{deg:prop-degree-drop} through
\eqref{deg:cubic-degree-dictionary}.  These intervals also overlap.

The remaining radii \(g=|\gamma|\), \(t=s^{1/3}\), and \(s\)
can be restricted independently. This gives eventual \(s_*,\gamma_*>0\);
the common raw cap is fixed before their final reduction and the choice
of \(m\). The closed alternatives include every equality face, and the bands in \(k,\pi,\mathfrak q,U\) overlap strictly. These choices depend only on the coefficients and are made before any
phase point is selected.
\end{proof}

\subsection{Proof of Lemma~\ref{lem:large-ratio-height}}
\label{app:proof-large-ratio-height}

\begin{proof}
At the central section,
\[
 Y_L(0)=p^{-2}Y_0\le K_0(1+p^{-2})\le2K_0.
\]
The graph equation obtained directly from
\eqref{deg:central-field} is
\begin{equation}
 \frac{dY_L}{dX}=
 \frac{Y_L\{q_\gamma X+r_LX^2+
       Cr_L(\gamma+r_LX)Y_L\}}
  {\ell(\gamma+r_LX)Y_L+p^{-2}-X+3\gamma X^2+r_LX^3}.
 \label{eq:large-ratio-graph}
\end{equation}
The denominator is positive on a through graph. At low height it may
approach zero, so we estimate growth only after the height has crossed
a fixed positive threshold.  Choose \(M_0>2K_0+1\) so
that, whenever \(Y_L\ge M_0\), the denominator is at least
\(\ell_0Y_L/2\).  On that region
\[
 Y_L'\le C_{L_c}+Cr_L(|\gamma|+r_LL_c)Y_L
       \le C_{L_c}+Y_L
\]
after a raw coefficient shrink.  Apply this inequality after the first crossing of \(M_0\), or
equivalently to \(\max(M_0,Y_L)\).  Gronwall's
inequality gives
\[
 Y_L(L_c)\le(M_0+C_{L_c}L_c)e^{L_c}=:B_c.
\]
Strict scalar order gives \(Y_L(L_c)>d_L\); since \(d_L\ge0\),
\eqref{deg:large-ratio-gap} follows.  No derivative lower bound for the
map from \(X=0\) to \(X=L_c\) has been used.
\end{proof}

\subsection{Proof of Lemma~\ref{deg:lem-large-source-routing}}
\label{app:proof-large-source-routing}

\begin{proof}
Fix the constant \(K_0\) in Corollary~\ref{rew:cor-high-derivative-signed}.
Its complete derivative bound applies on the connected high-height interval
\begin{equation}
 Y_0>K_0(1+p^2),\qquad 0<s^2Y_0<z_*.
 \label{deg:large-high-height}
\end{equation}
It remains to cover the lower-height complement by the stopped-trace
and strict source intervals.

Use the positive signed trace chart
\begin{equation}
 \alpha _0=-r_L^2U_0d,\qquad
 \alpha _1=r_Lp_L,\qquad
 d=\frac{p_L^2}{U_0p^2},\qquad
 \frac{s^2}{r_L^2}=\frac{p_L^2}{p^2},
 \label{deg:large-trace-dictionary}
\end{equation}
where \(p_L\) and \(U_0\) are positive units.  In the notation of
Proposition~\ref{deg:prop-stopped-trace}, the signed trace parameter is
\[
 a_L^{\rm tr}:=-U_0d=-\frac{p_L^2}{p^2},
 \qquad \alpha _0=r_L^2a_L^{\rm tr}.
\]
The stopped-trace proposition therefore applies with the signed
parameter \(a_L^{\rm tr}\), which includes the positive factor \(d\).
By increasing \(P_*\),
\(a_L^{\rm tr}\) lies in its signed trace block. For an orbit in
the complementary range \(Y_0\le K_0(1+p^2)\),
\[
 Y_L(0)=\frac{s^2Y_0}{r_L^2}
 \le K_0p_L^2(1+p^{-2})\le C.
\]
The height lemma also applies when \(p_L\) ranges over this compact
positive interval: its coordinates are related to the
\(r_L=\alpha_1\) coordinates by
\(X=p_LX_{\rm old}\), \(Y_L=p_L^2Y_{\rm old}\).
Its Gronwall bound is uniform for the corresponding cuts
\(X_{\rm old}=L_c/p_L\) in a fixed compact interval.
Lemma~\ref{lem:large-ratio-height} therefore gives a fixed number
\(B_0\) such that
\begin{equation}
 0<\frac{\mathfrak g_L}{r_L^2}\le B_0.
 \label{deg:large-bounded-gap}
\end{equation}

Take \(\xi_*=\xi_0\) from \eqref{deg:trace-overlap}. The compact annulus
\(\xi_*\le\mathfrak g_L/r_L^2\le B_0\) is covered by finitely many
strict source intervals, overlapping the terminal range
\(0<\mathfrak g_L/r_L^2\le2\xi_*\). On these intervals the coefficients satisfy
\[
 a_L^{\rm tr}=-\frac{p_L^2}{p^2}\in[-P_*^{-2}C,0],
 \qquad p_L\in[1/2,3/2],
\]
and all orbit segments retain strict positive margins. Proposition
\ref{deg:prop-stopped-trace}, with
\eqref{deg:trace-domain-inclusion} and \eqref{deg:trace-overlap}, supplies
their every-point physical cover. Proposition~\ref{sharp:upper-trace}
controls the derivative on the connected union of these intervals. The choices are
\begin{align*}
 &K_0,\ L_c,\ B_c,\ B_0;
   \quad \text{terminal interval},\ \xi_*;\\
 &P_*,\ \text{finite strict cover},\ L_t;
   \quad \text{common raw cap},\ \text{raw coefficient neighbourhood},\ m.
\end{align*}
The strict margins ensure that every later shrink preserves the domains
and overlaps already chosen.
\end{proof}

\section{Notation and reading index}
\label{app:notation-index}

The original DRR field and graphic are \(F_*,\Gamma_{I4}\);
the canonical endpoint is \(H_0,\Gamma\); the
neighboring organizing family is \(H_c,\Gamma_c\). The full quadratic germ is
\(\mathcal Q_2\), while \(\theta=(A,C,\mu_1,\mu_2,\mu_3)\) denotes the
five-parameter transverse slice. In the positive-infinity chart,
\(v=X/Y\) and \(z=1/Y\); \(h\) denotes a central height only where that
local section is fixed. The exact return is written
\(\mathcal P_\theta(x)=m_\theta\Pi_\theta(x)\), where \(m_\theta\) is the
finite through multiplier and \(\Pi_\theta\) is the completed divided
transition; \(F_\theta=x-\mathcal P_\theta\) is the displacement.

The local finite saddle--node through multiplier is
\(m_{\rm fin}=e^{-T_f}\) in \eqref{eq:finite-multiplier}, written \(m\)
in the completed return \eqref{eq:return}. In the fixed-\(c\) lower construction,
\(\Lambda_\rho\) is the complete external gain, including the two tail
gains and the finite connection, as defined in \eqref{p4:outer-data}.
The action-model target is \(m_\eps=2\eps^{5/2}\); the two-coefficient
calibration \eqref{p4:calibration} imposes \(\Lambda_\rho=m_\eps\).

In the positive-infinity chart, the polynomials are
\[
 P(v)=\mu _1+\mu _2v-v^3,\qquad
 R(v)=\mu _3+v+v^2
\]
as in \eqref{eq:polynomials}.  The symbol \(\gamma\) denotes the small
root of the transverse-rate equation \(R(\gamma)=0\), equivalently the
root selected in \eqref{eq:center-parameters}; it is not, in general, a
root of the equatorial drift \(P\).  The two symbols
\(\alpha _0,\alpha _1\) are the value and first derivative of that drift,
\[
 \alpha _0=P(\gamma),\qquad \alpha _1=P'(\gamma),
\]
not transverse coefficients.  The ratios \(p,\pi,\mathfrak q,U,G\) are
local coefficient dictionaries introduced at their first use in the
coefficient-regime and degree-loss sections.

The scale \(\tauc=t^2\) is reserved for the cubic coefficient chart in
\eqref{deg:cubic-chart}.  An unadorned \(\tau\) is local and must be read
from its defining display: in the coefficient charts it denotes the stated
small coefficient scale (for example \(p^2\) or \(gt_c/a\)), whereas in a
variation integral or a reparametrized orbit it can be a dummy or clock
variable.  No global identification of these uses is intended.  The source
variable and the physical section coordinate are not identified unless the
displayed affine adapter says so.

Finally, \(\Pi_\theta\) in \eqref{eq:return} is the completed, divided
transition: the infinity map composed with the regular section adapters and
the finite-side closing factor after the factor \(m\) has been extracted.
It is the positive transition appearing in the exact displacement
\(F=x-m\Pi_\theta(x)\); it is not an unrestricted Poincare map outside the
chosen physical sections.

The coefficient dictionaries define their ratios before use. The common
raw receiver and capture retain the original physical coordinate throughout;
positive affine changes transfer the derivative signs used for counting.


\section*{Funding}

This research received no specific grant from any funding agency in the
public, commercial, or not-for-profit sectors.

\section*{Declaration of competing interests}
The author declares no known competing financial interests or personal
relationships that could have appeared to influence the work reported in this
article.

\section*{Data and code availability}
The complete LaTeX source accompanies this manuscript. Supporting
symbolic checks and independent multiprecision shooting examples retain
their frozen coefficients and trajectory data in the accompanying
reproduction material. The numerical examples are supplementary and are
not premises of the analytic argument.

\section*{Declaration of generative AI and AI-assisted technologies}
During preparation of this manuscript the author used OpenAI Codex for
mathematical exploration and proof drafting, simulated technical readings,
organization and language editing, LaTeX and code assistance, literature
searches, and symbolic and internal-consistency checks. These tools do not
constitute independent human peer review or formal verification. The author
is responsible for the mathematical claims and the final manuscript.

\begingroup
\sloppy
\bibliographystyle{plain}
\bibliography{references}
\endgroup
\end{document}
